\documentclass[11pt,openany]{book}
\setcounter{secnumdepth}{3}
\usepackage[margin=1.1in]{geometry}
\usepackage[T2A]{fontenc}
\usepackage[utf8]{inputenc}
\usepackage[english,russian]{babel}
\usepackage{amsmath,amssymb,amsthm}
\usepackage{graphicx}
\usepackage{tikz}
\usetikzlibrary{arrows.meta,decorations.markings,decorations.pathmorphing,decorations.pathreplacing,calc}
\input{figures/icons}
\usepackage[unicode,colorlinks=true,linkcolor=blue!60!black,citecolor=blue!60!black]{hyperref}
\usepackage{microtype}
\usepackage{empheq}
\usepackage{array}
\usepackage{booktabs}
\usepackage{multirow}
\usepackage{longtable}
\usepackage{circuitikz}
\definecolor{keyyellow}{rgb}{1.0,0.97,0.78}
% key equations: \begin{empheq}[box=\keybox]{equation} ... \end{empheq} -- light-yellow box, number stays outside
\newcommand{\keybox}[1]{{\setlength{\fboxsep}{7pt}\colorbox{keyyellow}{#1}}}
\usepackage[most]{tcolorbox}
% key statements (propositions etc.): the same light-yellow background, breakable across pages
\newtcolorbox{keyblock}{enhanced,breakable,colback=keyyellow,colframe=keyyellow,boxrule=0pt,arc=0pt,left=6pt,right=6pt,top=2pt,bottom=2pt,boxsep=0pt}
\renewcommand{\chaptermark}[1]{\markboth{\small\chaptername~\thechapter.\ #1}{}}
\renewcommand{\sectionmark}[1]{\markright{\small\thesection.\ #1}}
\renewcommand{\topfraction}{0.95}
\renewcommand{\textfraction}{0.05}
\renewcommand{\floatpagefraction}{0.75}

% part pages: title at the top third, and text may follow on the same page (used for the part introductions)
\makeatletter
\renewcommand\part{\if@openright\cleardoublepage\else\clearpage\fi\thispagestyle{plain}\if@twocolumn\onecolumn\@tempswatrue\else\@tempswafalse\fi\null\vspace{0.18\textheight}\secdef\@part\@spart}
\renewcommand\@endpart{\par\vspace{3em}\if@tempswa\twocolumn\fi}
\makeatother
\newtheorem{theorem}{Теорема}[chapter]
\newtheorem{proposition}{Утверждение}[chapter]
\newtheorem{lemma}{Лемма}[chapter]
\theoremstyle{remark}
\newtheorem{remark}{Замечание}[chapter]
% exercises: \begin{exercise}[\lvlB] ... \end{exercise}, numbered "Задача 4.3"; difficulty marks
\theoremstyle{definition}
\newtheorem{exercise}{Задача}[chapter]
\newcommand{\lvlA}{$\star$}
\newcommand{\lvlB}{$\star\star$}
\newcommand{\lvlC}{$\star\star\star$}
% answer to an exercise in Appendix «Ответы»: \answer{ex:NN:k}{text}
\newcommand{\answer}[2]{\par\noindent\textbf{\ref{#1}.}~#2\par\smallskip}
\newcommand{\dd}{\mathrm{d}}
% neuron type code: first four letters of the primary mediator
\newcommand{\nt}[1]{\textsf{\textbf{#1}}}
\newcommand{\mV}{\,\text{мВ}}
\newcommand{\ms}{\,\text{мс}}
\newcommand{\mM}{\,\text{мМ}}
\newcommand{\um}{\,\text{мкм}}
\newcommand{\ion}[2]{\mathrm{#1}^{#2}}
\newcommand{\Na}{\ion{Na}{+}}
\newcommand{\K}{\ion{K}{+}}
\newcommand{\Cl}{\ion{Cl}{-}}
\newcommand{\Ca}{\ion{Ca}{2+}}
\newcommand{\conc}[1]{[\mathrm{#1}]}

\title{\Huge Компьютер на 20 ваттах\\[16pt] \Large Как вычисляет мозг: количественный подход для инженеров}
\hypersetup{pdftitle={Компьютер на 20 ваттах. Как вычисляет мозг: количественный подход для инженеров},pdfauthor={К. Кладько}}
\author{\Large К.~Кладько}
\date{}

\begin{document}
\frontmatter
\maketitle

\tableofcontents
\mainmatter

\chapter{Типы нейронов}
\label{sec:neurontypes}

\section{Нейрон как чёрный ящик}

Допустим, вам дали мешок с восемьюдесятью шестью миллиардами нейронов~\cite{Azevedo2009} и попросили разложить их по кучкам. Как бы вы это сделали?

Инженер, которому дали мешок незнакомых микросхем, не стал бы сортировать их по форме корпуса. Он спросил бы: сколько у детали входов, сколько выходов и что она выдаёт на выход. Нарисуем нейрон так же --- как блок на схеме (рис.~\ref{fig:blackbox}).

\begin{figure}[t]
\centering
\begin{tikzpicture}[>=Stealth,x=1cm,y=1cm]
% the box
\node[draw,very thick,minimum width=2.6cm,minimum height=2.6cm,fill=blue!6] (N) at (0,0) {};
\node at (0,0.55) {\small нейрон};
\node at (0,-0.15) {$\displaystyle u=\sum_i w_i x_i$};
\node at (0,-0.8) {\small $y=\Theta(u-\theta)$};
% inputs
\foreach \k/\y/\lab in {1/1.05/{x_1},2/0.55/{x_2},3/0.05/{x_3}}{
  \draw[->,thick,blue!60!black] (-4.2,\y) -- (-1.3,\y);
  \node[left] at (-4.2,\y) {$\lab$};
  \node[above,blue!60!black] at (-2.1,\y-0.06) {\scriptsize $w_{\k}$};}
\node at (-4.45,-0.4) {$\vdots$};
\node at (-2.1,-0.4) {$\vdots$};
\draw[->,thick,blue!60!black] (-4.2,-1.05) -- (-1.3,-1.05);
\node[left] at (-4.2,-1.05) {$x_n$};
\node[above,blue!60!black] at (-2.1,-1.11) {\scriptsize $w_{n}$};
\draw[decorate,decoration={brace,amplitude=5pt},gray!60!black] (-4.9,-1.2) -- (-4.9,1.2);
\node[left,align=right,gray!40!black] at (-5.1,0) {\small $n\sim10^3$--$10^5$\\[-2pt]\small входов};
% single output wire, then fan-out
\draw[very thick,red!70!black] (1.3,0) -- (3.2,0);
\foreach \x in {1.6,2.2,2.34,2.9}{\draw[thick,red!70!black] (\x,0) -- (\x,0.4);}
\node[above right,font=\tiny\itshape,gray!30!black,inner sep=1pt] at (1.42,0.45) {щёлк\dots\ щёлк-щёлк\dots};
\node[below,red!70!black,align=center] at (2.25,-0.05) {\scriptsize один выход $y$};
\fill[red!70!black] (3.2,0) circle (0.06);
\foreach \y in {1.3,0.65,0,-0.65,-1.3}{
  \draw[->,thick,red!70!black] (3.2,0) -- (5.0,\y);}
\draw[decorate,decoration={brace,amplitude=5pt},gray!60!black] (5.3,1.45) -- (5.3,-1.45);
\node[right,align=left,gray!40!black] at (5.5,0) {\small копии $y$\\[-2pt]\small на $10^3$--$10^4$\\[-2pt]\small других нейронов};
\end{tikzpicture}
\caption{Нейрон глазами инженера. Много входов $x_i$, каждый со своим весом $w_i$; внутри --- взвешенная сумма $u$ и сравнение с порогом $\theta$ ($\Theta$ --- ступенька: $1$, если аргумент положителен, и $0$ иначе); один выход $y$ --- последовательность импульсов, размноженная ветвлением на тысячи получателей.}
\label{fig:blackbox}
\end{figure}

На выходе блока --- не число, а последовательность коротких импульсов напряжения: «щёлк», пауза, «щёлк-щёлк», длинная пауза. Каждый импульс одинаковой высоты, так что вся информация --- в том, \emph{когда} они происходят. Внутри блока, в первом грубом приближении, стоит сумматор и компаратор: входы складываются с весами, и если сумма перевалила за порог, блок выдаёт импульс. В следующей главе мы заменим его моделью, которая честно работает в непрерывном времени.

Выход один, но провод ветвится, и каждая ветка несёт \emph{ту же самую} копию сигнала. В электронике это называется разветвлением по выходу, fan-out. У нейрона коры он порядка нескольких тысяч; логический вентиль в микросхеме обычно нагружен на единицы других вентилей.

\emph{Что} же передаётся по выходному проводу к получателю? Импульс добегает до конца каждой ветки и там превращается в химический сигнал: ветка выбрасывает в узкую щель между клетками щепотку определённого вещества, \emph{нейромедиатора}, которое меняет входной ток получателя. Вот и признак для сортировки: \emph{какой медиатор выбрасывает выход нейрона}.

\section{Один выход --- один тип сигнала}

Сортировка имеет смысл, только если у каждого нейрона тип выхода один. Так ли это? Опыт говорит, что почти так.

\begin{keyblock}
\begin{proposition}[Один нейрон --- один тип выхода]
\label{prop:dale}
У почти каждого нейрона один главный медиатор, и выбрасывает он его во всех ветках своего выхода. Поэтому тип нейрона задаётся его главным медиатором~\cite{Strata1999}.
\end{proposition}
\end{keyblock}

Договоримся обозначать тип нейрона \emph{первыми четырьмя буквами английского названия его главного медиатора}. Так, нейрон, главный медиатор которого глутамат, --- это \nt{GLUT}-нейрон. Все типы собраны в таблице~\ref{tab:types}.

\begin{table}[t]
\centering
\footnotesize
\setlength{\tabcolsep}{3pt}
\renewcommand{\arraystretch}{1.15}
\begin{tabular}{>{\raggedright}p{1.0cm}>{\raggedright}p{2.05cm}>{\raggedright}p{1.75cm}>{\raggedright}p{1.6cm}>{\centering}p{0.7cm}>{\raggedright}p{1.55cm}>{\raggedright}p{1.8cm}>{\raggedright\arraybackslash}p{3.2cm}}
\toprule
Тип & Главный медиатор & Шина & Скорость & Знак & Нейронов в мозге & Выходов у нейрона & Роль в вычислении \\
\midrule
\nt{GLUT} & глутамат & адресная & быстро и медленно & $+$ & $\sim8\cdot10^{10}$ & $\sim10^4$ & главный положительный вес \\
\nt{GABA} & GABA & адресная & быстро и медленно & $-$ & $\sim3\cdot10^{9}$ & $10^3$--$10^4$ & главный отрицательный вес \\
\nt{GLYC} & глицин & адресная & быстро & $-$ & $10^6$--$10^7$ & $10^2$--$10^3$ & отрицательный вес в спинном мозге и стволе; второй ключ для одного из рецепторов глутамата \\
\nt{ACET} & ацетилхолин & обе & быстро и медленно & $\pm$ & $\sim10^6$ & мышца: $10$--$10^3$; мозг: $\sim10^5$ & выход на мышцу; в мозге --- скорость обучения (гипотеза) \\
\nt{DOPA} & дофамин & широковещ. & медленно & $\pm$ & $\sim5\cdot10^{5}$ & $10^5$--$10^6$ & ошибка предсказания награды $\delta$ \\
\nt{SERO} & серотонин & широковещ. & медленно (один тип рецептора быстрый) & $\pm$ & $\sim3\cdot10^{5}$ & $\sim10^5$ & горизонт планирования (гипотеза) \\
\nt{HIST} & гистамин & широковещ. & медленно & $+$ & $\sim6\cdot10^{4}$ & нет оценки & глобальный сигнал «бодрствуй» \\
\nt{NORA} & норадреналин & широковещ. & медленно & $\pm$ & $\sim5\cdot10^{4}$ & $\sim10^5$ & коэффициент усиления (гипотеза) \\
\nt{ADRE} & адреналин & широковещ. & медленно & $\pm$ & $\sim10^4$ & нет оценки & немного нейронов; роль в мозге неясна \\
\nt{ASPA} & аспартат & адресная & быстро & $+$ & нет оценки & нет оценки & слабый положительный вес; роль спорна \\
\bottomrule
\end{tabular}
\caption{Типы нейронов, по убыванию числа нейронов в мозге. \emph{Шина}: адресная --- медиатор выбрасывается в узкую щель к определённому получателю; широковещательная --- от малой группы нейронов-источников на большие области (раздел~\ref{sec:buses}). \emph{Скорость}: быстро --- через ионотропные рецепторы, миллисекунды; медленно --- через метаботропные, сотни миллисекунд и дольше. \emph{Знак} --- обычный; окончательно его задаёт получатель (раздел~\ref{sec:receiver}), и $\pm$ означает, что встречаются оба знака. \emph{Нейронов в мозге} --- сколько нейронов этого типа во всём мозге человека, по порядку величины; всего нейронов около $8.6\cdot10^{10}$~\cite{Azevedo2009}. Почти все \nt{GLUT}-нейроны --- около $7\cdot10^{10}$ --- это мелкие входные нейроны мозжечка; числа \nt{GLUT} и \nt{GABA} получены из этого подсчёта и доли \nt{GABA}-нейронов в коре~\cite{Hendry1987}. Из малочисленных типов прямо подсчитаны только \nt{HIST}~\cite{Airaksinen1991} и главная из двух групп \nt{DOPA}-нейронов~\cite{Pakkenberg1991}, так что для \nt{DOPA} это нижняя граница; для \nt{SERO} --- сумма двух частичных подсчётов~\cite{Baker1990,Baker1991}. Числа для \nt{NORA}, \nt{GLYC}, \nt{ACET} и \nt{ADRE} --- грубые оценки по порядку величины, без прямого подсчёта. \emph{Выходов у нейрона} --- сколько у одного нейрона синаптических окончаний, то есть мест выброса медиатора на ветвях выходного провода (рис.~\ref{fig:blackbox}), по порядку величины. Для широковещательных типов окончания прямо не считали: у \nt{DOPA} число выведено из того, какую долю своей мишени покрывает ветвление одного выходного провода~\cite{Matsuda2009}, у остальных оценки ещё грубее. У \nt{ACET} два случая: нейрон, управляющий мышцей, и широковещательный нейрон в мозге.}
\label{tab:types}
\end{table}

% the pie is a table-type float with a figure caption, so it can never float ahead of Table~\ref{tab:types}
\begin{table}[tb]
\makeatletter\def\@captype{figure}\makeatother
\centering
\definecolor{vizglut}{HTML}{2A78D6}
\definecolor{vizgaba}{HTML}{E34948}
\definecolor{vizrest}{HTML}{8A8986}
\begin{tikzpicture}[>=Stealth]
% pie: GLUT 96.4 %, GABA 3.6 % (13.0 degrees), the rest 0.006 % (a hairline at 0 degrees)
\fill[vizglut] (0,0) -- (13:1.9) arc (13:360:1.9) -- cycle;
\fill[vizgaba] (0,0) -- (0:1.9) arc (0:13:1.9) -- cycle;
\draw[white,line width=1pt] (0,0) -- (13:1.9);
\draw[vizrest,line width=1.2pt] (0,0) -- (0:1.9);
\node[white,align=center,font=\small] at (190:0.95) {\nt{GLUT}\\$96.4\,\%$};
\draw[gray!60!black] (6.5:1.75) -- (2.05,2.1);
\node[anchor=west,font=\small] at (2.05,2.1) {\nt{GABA} $3.6\,\%$};
% zoom box for the remaining types
\draw[densely dashed,vizrest] (1.9,0) -- (3.1,1.65);
\draw[densely dashed,vizrest] (1.9,0) -- (3.1,-2.2);
\draw[vizrest,rounded corners=3pt] (3.1,-2.2) rectangle (10.1,1.65);
\node[anchor=west,font=\scriptsize] at (3.2,1.4) {все остальные вместе: $\approx0.006\,\%$; логарифмическая шкала};
\foreach \code/\lg/\y/\val/\lx in {GLYC/6.477/1.0/{$10^6$--$10^7$}/4.05,ACET/6.0/0.6/{$\sim10^6$}/3.0,DOPA/5.699/0.2/{$\sim5\cdot10^5$}/2.699,SERO/5.477/-0.2/{$\sim3\cdot10^5$}/2.477,HIST/4.778/-0.6/{$\sim6\cdot10^4$}/1.778,NORA/4.699/-1.0/{$\sim5\cdot10^4$}/1.699,ADRE/4.0/-1.4/{$\sim10^4$}/1.0}{
  \node[anchor=east,font=\scriptsize] at (4.35,\y) {\nt{\code}};
  \fill[vizrest] (4.45,\y-0.13) rectangle ({4.45+\lg-3},\y+0.13);
  \node[anchor=west,font=\scriptsize] at ({4.5+\lx},\y) {\val};}
\draw[vizrest,thick] (7.45,1.0) -- (8.45,1.0);
\draw[vizrest,thick] (7.45,0.92) -- (7.45,1.08);
\draw[vizrest,thick] (8.45,0.92) -- (8.45,1.08);
\draw[gray!60!black] (4.45,-1.72) -- (8.45,-1.72);
\foreach \e in {3,4,5,6,7}{
  \draw[gray!60!black] ({4.45+\e-3},-1.72) -- ({4.45+\e-3},-1.66);
  \node[below,font=\scriptsize] at ({4.45+\e-3},-1.72) {$10^{\e}$};}
\end{tikzpicture}
\caption{Доли типов нейронов в мозге по оценкам таблицы~\ref{tab:types}. Круг почти целиком занят \nt{GLUT}; \nt{GABA} --- узкий сектор, а все остальные типы вместе --- волосок на границе. Справа этот волосок в увеличении, по логарифмической шкале числа нейронов; для \nt{GLYC} столбик доходит до середины диапазона $10^6$--$10^7$, отрезок показывает весь диапазон. У \nt{ASPA} оценки нет, и на рисунке его нет.}
\label{fig:typepie}
\end{table}

Насколько неравны эти кучки, видно на рис.~\ref{fig:typepie}: из каждой тысячи нейронов $964$ --- \nt{GLUT}, $36$ --- \nt{GABA}, а на все остальные типы вместе приходится около шести нейронов на сто тысяч.

У GABA название и так состоит из четырёх букв. Вещества, которые почти никогда не бывают главным медиатором, --- нейропептиды, газы, липиды, пурины, --- своего кода не получают; о них --- в приложении~\ref{sec:othermediators}. Слово «почти» в утверждении~\ref{prop:dale} важно. Многие нейроны выбрасывают вместе с главным медиатором вспомогательные вещества~\cite{Yao2023}. Некоторые выбрасывают и второй быстрый медиатор, даже противоположного знака: часть \nt{DOPA}-нейронов выбрасывает ещё и GABA~\cite{Tritsch2012}. А у некоторых нейронов разные медиаторы выходят из разных веток одного и того же выхода~\cite{Zhang2015}. Всё это измерено, так что «один нейрон --- один медиатор» --- правило с исключениями, а не закон. Но как первое деление оно выдерживает проверку: в клеточном атласе мозга мыши, где нейроны разложены по работающим в них генам на пять с лишним тысяч групп, нейроны по-прежнему делятся на крупные классы прежде всего по главному медиатору~\cite{Yao2023}.

У этого правила есть следствие, которое сразу отличает мозг от искусственной нейронной сети. В искусственной сети один и тот же нейрон может иметь положительный вес к одному получателю и отрицательный к другому: веса $w_{ij}$ --- просто числа в матрице. В мозге знак действия в основном определяется медиатором, а медиатор у нейрона один. Значит, если нейрон $j$ возбуждает одного получателя, он возбуждает и всех остальных. Весь \emph{столбец} матрицы весов, выходящий из нейрона $j$, одного знака:
\begin{empheq}[box=\keybox]{equation}
\operatorname{sign} w_{ij} \;=\; s_j \quad\text{для всех получателей } i, \qquad s_j\in\{+1,-1\}.
\label{eq:dalesign}
\end{empheq}
Нейроны делятся на возбуждающие ($s_j=+1$) и тормозящие ($s_j=-1$), и это свойство самого нейрона, а не связи. Если сети нужна отрицательная связь от возбуждающего нейрона, её приходится делать через посредника: возбуждающий нейрон возбуждает тормозящий, а тот тормозит цель: отрицательный вес с задержкой на одно звено.

Медиаторов известно больше сотни, но почти вся работа мозга держится на полудюжине, и они делятся на два совершенно разных класса.

\section{Две шины: адресная и широковещательная}
\label{sec:buses}

В компьютерных сетях есть два способа доставить сообщение. Первый --- \emph{адресный}, unicast: пакет идёт от одного отправителя к определённому получателю, быстро, и больше его никто не видит. Второй --- \emph{широковещательный}, broadcast: отправитель передаёт одно сообщение сразу всем узлам сети. Нейроны пользуются обоими способами, и медиаторы делятся ровно по этому признаку (рис.~\ref{fig:phone_radio}).

\begin{figure}[t]
\centering
\begin{tikzpicture}[>=Stealth,scale=0.95]
% ---- unicast: point-to-point wiring
\node[above] at (2.0,3.3) {\small \textbf{Адресная}: глутамат и GABA};
\foreach \i in {0,...,4}{
  \fill[blue!60!black] (0.2+0.9*\i,2.5) circle (0.1);
  \fill[blue!60!black] (0.2+0.9*\i,0.6) circle (0.1);}
\foreach \a/\b in {0/0,0/1,1/1,1/3,2/2,3/2,3/4,4/4,2/0}{
  \draw[->,blue!60!black] (0.2+0.9*\a,2.38) -- (0.2+0.9*\b,0.72);}
\fill[red!70!black] (1.1,1.55) circle (0.09);
\pic[scale=1.2] at (1.75,1.9) {letter};
\draw[->,red!70!black,thick] (1.1,1.55) -- (0.28,0.7);
\draw[->,red!70!black,thick] (1.1,1.55) -- (1.95,0.72);
\node[align=center,below] at (2.0,0.2) {\small миллиарды нейронов,\\[-2pt]\small у каждого свои адресаты,\\[-2pt]\small время: миллисекунды};
% ---- broadcast: one small source, global targets
\begin{scope}[xshift=7.2cm]
\node[above] at (1.8,3.3) {\small \textbf{Широковещательная}: дофамин, \dots};
\foreach \i in {0,...,8}{\fill[blue!60!black] (-0.2+0.5*\i,2.75) circle (0.08);}
\fill[orange!85!black] (1.8,0.35) ellipse (0.35 and 0.18);
\pic[scale=1.5] at (2.7,0.75) {mast};
\foreach \x in {-0.2,0.3,0.8,1.3,1.8,2.3,2.8,3.3,3.8}{
  \draw[->,orange!85!black] (1.8,0.5) .. controls (1.8,1.3) and (\x,1.6) .. (\x,2.62);}
\node[align=center,below] at (1.8,0.1) {\small тысячи --- сотни тысяч нейронов,\\[-2pt]\small одно сообщение всем,\\[-2pt]\small время: сотни миллисекунд и дольше};
\end{scope}
\end{tikzpicture}
\caption{Два способа передачи. Слева --- адресная шина, похожая на почту с адресом на конверте; красным выделен один нейрон и два его получателя. Справа --- широковещательная, как радиостанция: небольшая группа нейронов-источников, и одно и то же сообщение получают все.}
\label{fig:phone_radio}
\end{figure}

\emph{Адресная шина} --- это глутамат и GABA. Их выбрасывает подавляющее большинство нейронов. Каждый выход ведёт к определённым клеткам, медиатор действует только в узкой щели контакта, и действует быстро: за миллисекунду открываются ионные каналы получателя, за несколько миллисекунд всё кончается. Это шина \emph{данных}: по ней передаётся то, что мозг вычисляет.

\emph{Широковещательная шина} --- дофамин, серотонин, норадреналин и отчасти ацетилхолин. Их выбрасывают крошечные группы нейронов, но выход каждого из них ветвится невероятно широко. Медиатор часто выбрасывается не в узкую щель, а в межклеточное пространство, где он расплывается и действует на всех, кто рядом. Действует медленно: не открывает каналы напрямую, а запускает внутри получателя цепочку химических реакций. По ней идут не данные, а \emph{управляющие параметры}, общие для больших участков сети, которые меняют режим её работы: коэффициент усиления, скорость обучения, сигнал «это было хорошо». Поэтому такие медиаторы называют \emph{модуляторами}. Долго считалось, что каждый такой канал --- одно число на весь мозг. Современные измерения уточнили картину: канал --- не один провод, а небольшой жгут параллельных линий. Нейроны-источники одного медиатора делятся на подсистемы, у которых свои получатели и своё сообщение, а сколько медиатора выбросить на месте, подстраивают ещё и местные сигналы~\cite{Ren2018,Poe2020}. Число таких линий --- единицы или десятки, а не миллиарды, так что канал остаётся широковещательным.

Если вам нужна одна картинка из этой главы --- вот она. Мозг --- огромная сеть с адресной передачей данных, над которой висит несколько широковещательных каналов с управляющими сигналами. Программист узнает знакомое устройство: вычисление плюс небольшой набор управляющих переменных, каждая из которых общая для большого участка сети.

\section{\nt{GLUT}: положительный вес}

Глутамат --- главный возбуждающий медиатор мозга. Когда выход нейрона выбрасывает глутамат, у получателя открываются каналы, через которые внутрь входит натрий, напряжение на мембране получателя растёт, и его шанс выдать собственный импульс увеличивается. На языке рис.~\ref{fig:blackbox} это вход с \emph{положительным} весом, $w_i>0$.

Кто выдаёт глутамат? Почти все, кто передаёт данные на большие расстояния: от трёх четвертей до четырёх пятых нейронов коры и большинство нейронов, связывающих разные области мозга. Сеть мозга --- это в первую очередь сеть глутаматных связей.

Одна инженерная деталь. После выброса концентрация глутамата в щели взлетает в тысячи раз, примерно до миллимоля, и спадает с постоянной времени около миллисекунды~\cite{Clements1992}: глутамат расплывается из щели, а специальные белки-насосы откачивают его обратно в клетки. Зачем? Затем же, зачем в линии связи после каждого символа сигнал возвращают к нулю. Если медиатор не убрать, следующий импульс придёт на «занятую» линию, и импульсы с интервалом в несколько миллисекунд сольются.

\section{\nt{GABA}: отрицательный вес}

Представьте сеть, в которой все веса положительны. Если в среднем каждый импульс порождает больше одного следующего импульса, активность растёт экспоненциально и через несколько шагов охватывает всю сеть. Электронщик узнает петлю положительной обратной связи с усилением больше единицы: схема уходит в насыщение. Чтобы этого не случилось, нужны отрицательные веса. Главный их источник --- гамма-аминомасляная кислота, сокращённо GABA.

GABA открывает у получателя каналы, пропускающие ионы хлора. Равновесный потенциал хлора лежит около потенциала покоя или немного ниже, поэтому открытые хлорные каналы удерживают мембрану около покоя и не дают поднять напряжение до порога. Это вход с \emph{отрицательным} весом, $w_i<0$. GABA-нейроны --- от пятой части до четверти нейронов коры~\cite{Hendry1987}. Их выходы короткие, и тормозят они своих ближайших соседей.

Торможение в мозге делает гораздо больше, чем просто вычитание. Оно работает как отрицательная обратная связь, которая держит всю сеть в рабочей точке. Считается, что в нормально работающей коре возбуждающий и тормозящий токи, приходящие на нейрон, почти точно уравновешивают друг друга: такой режим предсказан теорией~\cite{vanVreeswijk1996} и виден в измерениях токов в живой коре~\cite{Haider2006}, и нейрон всё время сидит около порога. Схема, стабилизированная сильной отрицательной обратной связью около неустойчивой точки, быстро и чутко реагирует на малые сигналы --- старый инженерный приём.

\section{\nt{DOPA}: сигнал ошибки}

Первый широковещательный канал --- дофамин. Дофаминовых нейронов у человека порядка полумиллиона, примерно один на сто семьдесят тысяч; столько подсчитано в главной из двух их групп~\cite{Pakkenberg1991}, так что это нижняя граница. Их выходы расходятся к тем областям, которые выбирают действия.

Что передаёт этот канал? Ответ дают опыты, в которых регистрируют импульсы дофаминовых нейронов у животного, получающего награду~\cite{Schultz1997}. Животному время от времени неожиданно дают каплю сока, и дофаминовые нейроны отвечают на неё короткой пачкой импульсов. Потом перед соком начинают зажигать лампочку, всегда за секунду до сока. Когда животное выучивает эту связь, нейроны начинают отвечать пачкой на \emph{лампочку}, а на сам сок, пришедший точно в срок, не отвечают совсем. А если после лампочки сок не дают, то в момент, когда он должен был прийти, нейроны \emph{замолкают}, и их частота падает ниже обычной.

Правило, которое объясняет все три случая (рис.~\ref{fig:rpe}): нейроны сообщают не о том, насколько хорошо, а о том, насколько \emph{лучше или хуже, чем ожидалось}.

\begin{figure}[t]
\centering
\begin{tikzpicture}[>=Stealth,x=3.4cm,y=1cm]
\foreach \y/\lab in {2.8/{неожиданный сок},1.4/{сок после лампочки},0/{лампочка, но сока нет}}{
  \draw[gray!70!black] (0,\y) -- (3,\y);
  \draw[densely dotted,gray!60!black] (0,\y+0.3) -- (3,\y+0.3);
  \node[left,align=right,font=\small] at (-0.03,\y+0.35) {\lab};}
% row 1: juice without a cue -> burst at the juice
\draw[very thick,orange!85!black] plot[domain=0:3,samples=120] (\x,{2.8+0.3+0.75*exp(-((\x-2.08)/0.07)^2)});
\pic[scale=0.8] at (2,2.58) {juice};
% row 2: cue then juice -> burst at the cue, nothing at the juice
\draw[very thick,orange!85!black] plot[domain=0:3,samples=120] (\x,{1.4+0.3+0.75*exp(-((\x-1.08)/0.07)^2)});
\pic[scale=0.85] at (1,1.15) {lamp}; \pic[scale=0.85] at (2,1.15) {juice};
% row 3: cue, juice omitted -> burst at the cue, dip when the juice was due
\draw[very thick,orange!85!black] plot[domain=0:3,samples=120] (\x,{0.3+0.75*exp(-((\x-1.08)/0.07)^2)-0.28*exp(-((\x-2.1)/0.12)^2)});
\pic[scale=0.85] at (1,-0.25) {lamp}; \pic[scale=0.85] at (2,-0.25) {nojuice};
\node[right,font=\scriptsize] at (2.36,0.1) {провал};
\node[right,font=\scriptsize] at (2.13,3.75) {сюрприз!};
\node[left,font=\scriptsize] at (0.97,2.25) {всплеск на лампочку};
\node[above,font=\scriptsize] at (2,1.75) {ничего};
\draw[->,gray!70!black] (0,-0.75) -- (3.05,-0.75) node[right,black,font=\small] {$t$};
\draw[gray!60!black] (1,-0.8) -- (1,-0.7) node[below=2pt,font=\scriptsize] {лампочка, $1\,\text{с}$};
\draw[gray!60!black] (2,-0.8) -- (2,-0.7) node[below=2pt,font=\scriptsize] {сок, $2\,\text{с}$};
\end{tikzpicture}
\caption{Частота импульсов \nt{DOPA}-нейронов в трёх опытах (схематично, по~\cite{Schultz1997}). Пунктир --- ровная фоновая частота. Лампочка --- сигнал, капля --- сок, перечёркнутая капля --- сок, который обещали, но не дали. Отклик --- ошибка предсказания~\eqref{eq:rpe}.}
\label{fig:rpe}
\end{figure}

Программист, знакомый с обучением с подкреплением~\cite{SuttonBarto2018}, узнает эту величину сразу. Агент, который учится выбирать действия, хранит оценку $V(s)$ --- сколько награды он ожидает, находясь в состоянии $s$. После каждого шага он сравнивает ожидание с тем, что получил:
\begin{empheq}[box=\keybox]{equation}
\delta_t \;=\; r_t \;+\; \gamma\,V(s_{t+1}) \;-\; V(s_t) ,
\label{eq:rpe}
\end{empheq}
и поправляет оценку: $V(s_t)\leftarrow V(s_t)+\alpha\,\delta_t$. Здесь $r_t$ --- полученная награда, $\gamma<1$ --- коэффициент дисконтирования (насколько будущая награда дешевле немедленной), $\alpha$ --- скорость обучения. Величина $\delta_t$ называется \emph{ошибкой временных разностей}.

Она ведёт себя в точности как дофаминовые нейроны. Неожиданный сок: $V$ ещё ноль, $r>0$, $\delta>0$. Выученный сок: награду предсказала лампочка, $V(s_t)$ перед соком уже равна $r$, и $\delta=0$. Сок не пришёл: $V(s_t)=r$, а $r_t=0$, $\delta<0$. А всплеск переезжает на лампочку, потому что именно в момент лампочки ожидание скачком выросло: $\gamma V(s_{t+1})-V(s_t)>0$. Шаги $t,t+1,\dots$ здесь --- условность алгоритма: в организме часов нет, и ту же величину в непрерывном времени мы получим в главе~\ref{sec:dofa}.

Значит, дофамин --- это широковещательная рассылка скалярного сигнала ошибки $\delta$ всем, кто должен по нему учиться. В первом приближении --- один провод на весь мозг и одно число в каждый момент времени; насколько этот провод на самом деле жгут, разбирается в главе~\ref{sec:dofaalt}.

\section{Остальные широковещательные каналы}

Для остальных модуляторов есть лишь рабочие гипотезы, которые переводят их роли на тот же язык глобальных параметров~\cite{Doya2002}. Держите их именно как гипотезы.

\emph{\nt{NORA}, норадреналин: коэффициент усиления.} Норадреналиновых нейронов ещё меньше, чем дофаминовых, по грубой оценке несколько десятков тысяч, а их выходы расходятся практически по всему мозгу. Они резко активируются, когда происходит что-то неожиданное или важное. Норадреналин меняет крутизну передаточной характеристики нейронов-получателей: слабые входы становятся слабее, сильные --- сильнее. Измеряется это так: фоновые импульсы получателя подавляются сильнее, чем ответ на вход, и отношение сигнала к шуму растёт~\cite{Foote1975NE}. В простой модели это описывают одним числом: если нейрон отвечает частотой $f=F\bigl(g\cdot u\bigr)$ на входной ток $u$, то норадреналин увеличивает коэффициент $g$~\cite{AstonJones2005} (подробнее --- в главе~\ref{sec:nora}). Сеть с большим $g$ работает «контрастнее» и быстрее принимает решения; с маленьким --- «размыто» и больше перебирает варианты, как агент обучения с подкреплением в режиме поиска.

\emph{\nt{ACET}, ацетилхолин: скорость обучения.} Ацетилхолин управляет тем, насколько сеть слушает текущий вход, а насколько --- собственные сохранённые данные. При высоком уровне ацетилхолина входы с органов чувств усиливаются, а внутренние, «вспоминающие» связи ослабляются, и одновременно облегчается изменение весов~\cite{Hasselmo2006}. Грубо говоря, это переключатель «запись/чтение» и вместе с ним --- скорость обучения $\alpha$.

\emph{\nt{SERO}, серотонин: горизонт планирования.} Что передаёт серотонин, понятно хуже всего. Одна из гипотез связывает его с коэффициентом дисконтирования $\gamma$ из~\eqref{eq:rpe}: высокий уровень серотонина соответствует готовности подождать большую награду, а не хватать маленькую сейчас. Серотониновые нейроны работают ровно и медленно, несколько импульсов в секунду в бодрствовании, и почти замолкают в одной из фаз сна~\cite{JacobsAzmitia1992}.

Все шесть медиаторов собраны в таблице~\ref{tab:transmitters}.

\begin{table}[t]
\centering
\small
\begin{tabular}{>{\raggedright}p{2.4cm}>{\raggedright}p{3.0cm}>{\raggedright}p{2.0cm}>{\raggedright\arraybackslash}p{5.2cm}}
\toprule
Тип нейрона & Доля нейронов & Шина & Роль в вычислении \\
\midrule
\nt{GLUT} (глутамат) & большинство & адресная & положительный вес, $w>0$ \\
\nt{GABA} (GABA) & 20--25\% в коре & адресная & отрицательный вес, $w<0$; стабилизация \\
\nt{DOPA} (дофамин) & $\sim 6\cdot10^{-6}$ & широковещ. & ошибка предсказания $\delta$ \\
\nt{NORA} (норадреналин) & $\sim 6\cdot10^{-7}$ & широковещ. & коэффициент усиления $g$ (гипотеза) \\
\nt{ACET} (ацетилхолин) & малая & обе & скорость обучения $\alpha$; «запись/чтение» (гипотеза) \\
\nt{SERO} (серотонин) & $\sim 3\cdot10^{-6}$ & широковещ. & дисконтирование $\gamma$ (гипотеза) \\
\bottomrule
\end{tabular}
\caption{Главные медиаторы мозга на языке вычислений. Правый столбец --- сильное упрощение: надёжен для глутамата, GABA и дофамина, для остальных это гипотезы.}
\label{tab:transmitters}
\end{table}

\section{Знак задаёт получатель}
\label{sec:receiver}

Я немного вас обманул, говоря: глутамат --- положительный вес, GABA --- отрицательный. Ни одна молекула не несёт в себе знака. Молекула --- это просто ключ. Что произойдёт, когда её поймают, решает \emph{замок} --- рецептор на клетке-получателе.

Лучший пример --- дофамин. У него есть рецепторы двух семейств~\cite{Beaulieu2011}. На одних он облегчает возбуждение клетки, на других --- затрудняет. В области, которая выбирает действия, одни нейроны несут рецепторы первого типа, другие --- второго, и один и тот же дофаминовый сигнал одновременно усиливает одну группу и ослабляет другую. Это как один широковещательный пакет, который разные узлы сети интерпретируют по-разному.

\begin{keyblock}
\begin{proposition}[Смысл сообщения задаёт получатель]
\label{prop:receptor}
Знак и скорость действия медиатора определяются не самим медиатором, а рецептором, с которым он связывается. Рецепторы бывают двух родов. \emph{Ионотропные} сами являются ионными каналами и открываются за доли миллисекунды: это прямое управление током, как затвор транзистора. \emph{Метаботропные} запускают внутри клетки цепочку химических реакций и действуют за сотни миллисекунд и дольше: это скорее запись в регистр настроек, которая меняет последующую работу клетки. Адресная шина говорит в основном через первые, широковещательная --- в основном через вторые.
\end{proposition}
\end{keyblock}

Почему тогда вообще можно сказать «глутамат возбуждает»? Потому что почти все рецепторы глутамата, которые встречаются на практике, возбуждающие, а почти все рецепторы GABA --- тормозные. Это не закон природы, а очень надёжное соглашение, и правило~\eqref{eq:dalesign} на нём и держится.

\section{Что унести с собой}

Подавляющее большинство нейронов --- адресная шина данных с весами одного знака на нейрон; ничтожное меньшинство --- широковещательные каналы, передающие на большие участки мозга несколько общих управляющих сигналов. Около двух нейронов из каждых ста тысяч --- и от них зависит, как учится и в каком режиме работает вся остальная сеть. Для инженера это не удивительно: в любом компьютере регистров управления несравнимо меньше, чем ячеек памяти, и всё же без них машина не работает.

В следующей главе мы проверим, что может вычислить сеть из одних \nt{GLUT}-нейронов, самого многочисленного типа, без всяких часов.

\section{Задачи}

\begin{exercise}[\lvlA]
\label{ex:01:1}
\emph{Кучки и провода.} По таблице~\ref{tab:types} (середина диапазона по логарифмической шкале; \nt{ACET} --- $10^5$ выходов): (а)~если нейрон --- человек, сколько населений Земли составят \nt{GLUT}-нейроны и город какого размера --- все широковещательные? (б)~Во сколько раз доля окончаний \nt{DOPA}, \nt{SERO}, \nt{NORA}, \nt{ACET} больше доли их нейронов?
\end{exercise}

\begin{exercise}[\lvlA]
\label{ex:01:2}
\emph{Возврат к нулю.} Глутамат в щели спадает как $c_0e^{-t/\tau_c}$, $c_0=1\mM$, $\tau_c=1\ms$; получатель замечает концентрацию выше $10$~мкМ. (а)~Сколько времени линия «занята» после одного выброса? (б)~Насосы частично заблокированы, $\tau_c=10\ms$. До какой частоты выбросов линия ещё успевает освободиться?
\end{exercise}

\begin{exercise}[\lvlB]
\label{ex:01:3}
\emph{Куда переезжает всплеск.} В опыте рис.~\ref{fig:rpe} после лампочки идут состояния $s_0\to\dots\to s_{K-1}$ с шагом $\Delta$, $T=K\Delta$; награда $r$ приходит на последнем переходе, момент лампочки непредсказуем. Найдите выученные $V(s_k)$ и ответ $\delta$ на лампочку. Положив $\gamma=e^{-\Delta/\tau_\gamma}$, найдите $\tau_\gamma$ при $T=1$~с, $\Delta=0.1$~с, $\gamma=0.95$.
\end{exercise}

\begin{exercise}[\lvlB]
\label{ex:01:4}
\emph{Моделирование: обучение по ошибке.} Промоделируйте цепочку задачи~\ref{ex:01:3} при $K=10$, $\gamma=0.95$, $r=1$ с правилом $V(s_t)\leftarrow V(s_t)+\alpha\delta_t$. В каком испытании $n^*$ ответ на лампочку впервые обгоняет ответ на сок при $\alpha=0.05$, $0.1$, $0.2$, $0.5$? Как $n^*$ зависит от $\alpha$?
\end{exercise}

\begin{exercise}[\lvlB]
\label{ex:01:5}
\emph{Проектирование: рассылка одного числа.} Число $\delta$ надо доставить всем $10^{10}$ нейронам; каждый получатель, включая ретрансляторы, должен получить $10$ окончаний от предыдущего слоя, звено добавляет $1\ms$. Сколько нейронов и звеньев в самом экономном дереве ретрансляторов при $10^4$ выходах на нейрон (\nt{GLUT}) и при $10^{5.5}$ (\nt{DOPA})?
\end{exercise}

\begin{exercise}[\lvlB]
\label{ex:01:6}
\emph{Почему равновесие чуткое.} В сети $80\,\%$ \nt{GLUT} и $20\,\%$ \nt{GABA}, каждый связан с каждым весами $J_E/N$ и $-J_I/N$; $\tau\dot r=-r+Wr+h$. При $J_E=10$ единственное ненулевое собственное число $W$ равно $0.5$. Найдите отношение тормозящего тока к возбуждающему и на сколько процентов достаточно ослабить торможение, чтобы сеть ушла в лавину.
\end{exercise}

\begin{exercise}[\lvlC]
\label{ex:01:7}
\emph{Исследование: \nt{SERO} --- это $\gamma$?} По задаче~\ref{ex:01:3} ответ \nt{DOPA}-нейронов на лампочку равен $re^{-T/\tau_\gamma}$. Задержки $T=1,\dots,5$~с по $n$ предъявлений, $\ln\delta$ измеряется с разбросом $0.5$. Сколько нужно $n$, чтобы отличить $\tau_\gamma=2$~с от $2.4$~с (уровень $0.05$, мощность $0.8$)? Какой механизм, кроме $\gamma$, дал бы тот же спад с $T$?
\end{exercise}

\chapter{Что вычисляют \nt{ГЛУТ}-нейроны}
\label{sec:glutcomputer}

\section{Правила игры}

В главе~\ref{sec:neurontypes} мы увидели, что подавляющее большинство нейронов мозга --- \nt{ГЛУТ}: они только возбуждают, их веса только положительны. Возьмём сколько угодно таких нейронов и спросим: что может вычислить такая сеть, если разрешать себе только то, что бывает в живом организме?

А в живом организме нет тактового генератора, который переключал бы все элементы одновременно. Каждый нейрон работает сам по себе, в непрерывном времени: импульсы приходят и уходят когда придётся, с разными задержками. Есть \emph{входы} --- чувствительные нейроны, которые срабатывают от света, звука или давления, --- и \emph{выходы} --- нейроны, которые управляют мышцами. Между ними --- сеть, и никто не говорит ей, когда начинать и когда заканчивать вычисление. Для электронщика это \emph{асинхронная} схема с элементами, задержки которых заранее точно не известны.

Модель нейрона возьмём самую простую из тех, что честно работают в непрерывном времени. Напряжение $V$ на мембране нейрона в покое равно нулю. Каждый пришедший импульс от нейрона $j$ скачком поднимает его на $w_j\ge0$, после чего напряжение само стекает обратно к нулю с постоянной времени $\tau$:
\begin{empheq}[box=\keybox]{equation}
\tau\,\frac{\dd V}{\dd t} \;=\; -V \;+\; \tau\sum_j w_j \sum_k \delta\bigl(t-t_j^{(k)}-d_j\bigr), \qquad w_j\ge 0 .
\label{eq:glutneuron}
\end{empheq}
Здесь $t_j^{(k)}$ --- моменты импульсов нейрона $j$, $d_j$ --- задержка на пути от него, $\delta$ --- дельта-функция, которая и означает мгновенный скачок. Когда $V$ достигает порога $\theta$, нейрон выдаёт импульс, $V$ сбрасывается в ноль, и примерно миллисекунду нейрон не может сработать снова. Типичные числа: $\tau$ --- от долей миллисекунды до двадцати миллисекунд у разных нейронов, задержки --- от долей миллисекунды до десятков миллисекунд. Это \emph{интегрирующий нейрон с утечкой}: RC-цепь с порогом. Это сознательное упрощение: у нейронов коры ветви входов сами выдают местные импульсы~\cite{Smith2013}, и один нейрон ведёт себя скорее как небольшая многослойная сеть (глава~\ref{sec:synthesis}). Для вопросов этой главы хватает и одной RC-цепи.

Главное отличие от логического вентиля --- утечка: нейрон помнит импульс не вечно, а примерно $\tau$, и складывается только то, что пришло в пределах этого окна.

\section{ИЛИ и И}

Начнём с вентилей. Если один импульс поднимает напряжение выше порога, $w\ge\theta$, нейрон срабатывает на каждый импульс любого входа. Это \emph{ИЛИ}.

С И интереснее. Пусть $w<\theta<2w$: одного импульса мало, нужны два. Но теперь вопрос «пришли ли оба» не имеет смысла, пока не сказано, \emph{в пределах какого времени}. Пусть первый импульс пришёл в момент $0$, а второй --- через время $\Delta$. К приходу второго от первого осталось $we^{-\Delta/\tau}$, и нейрон сработает, если $we^{-\Delta/\tau}+w\ge\theta$. Отсюда окно совпадения:
\begin{empheq}[box=\keybox]{equation}
\Delta \;\le\; \Delta_{\max} \;=\; \tau\,\ln\frac{w}{\theta-w} .
\label{eq:window}
\end{empheq}
Это не логическое И, а \emph{детектор совпадений}: И во времени (рис.~\ref{fig:coincidence}). При $\theta=1.5w$ окно равно $\tau\ln2\approx0.7\tau$. У нейрона коры с $\tau=10\ms$ это около семи миллисекунд. У нейронов слуховой системы, где $\tau$ меньше миллисекунды, окно сжимается до долей миллисекунды.

\begin{figure}[t]
\centering
\begin{tikzpicture}[>=Stealth,x=0.9cm,y=1.6cm]
\foreach \dx/\lab/\c [count=\i] in {0.5/{совпали}/red!70!black, 2.6/{не совпали}/blue!60!black}{
 \begin{scope}[xshift={(\i-1)*7.2cm}]
  \draw[->,gray!70!black] (0,0) -- (6.3,0) node[below left,black] {\small $t$};
  \draw[->,gray!70!black] (0,0) -- (0,1.75) node[left,black] {\small $V$};
  \draw[densely dashed,gray!60!black] (0,1.5) -- (6.0,1.5) node[right,black] {\small $\theta$};
  \draw[very thick,\c] (0,0) -- (1,0) -- (1,1)
      plot[domain=1:1+\dx,samples=30] (\x,{exp(-(\x-1)/1.5)});
  \pgfmathsetmacro{\v}{exp(-\dx/1.5)+1}
  \draw[very thick,\c] (1+\dx,{exp(-\dx/1.5)}) -- (1+\dx,{min(\v,1.5)});
  \ifdim\v pt<1.5pt
    \draw[very thick,\c] plot[domain=1+\dx:6,samples=40] (\x,{\v*exp(-(\x-1-\dx)/1.5)});
  \else
    \draw[very thick,\c] (1+\dx,1.5) -- (1+\dx,0) -- (6,0);
    \draw[->,thick,\c] (1+\dx,1.55) -- (1+\dx,1.72) node[above,font=\scriptsize] {импульс};
  \fi
  \draw[->,thick] (1,-0.35) -- (1,-0.08); \draw[->,thick] (1+\dx,-0.35) -- (1+\dx,-0.08);
  \draw[decorate,decoration={brace,amplitude=3pt,mirror}] (1,-0.42) -- (1+\dx,-0.42) node[midway,below=3pt] {\scriptsize $\Delta$};
  \node[\c] at (3.5,-0.95) {\small \lab};
 \end{scope}}
\end{tikzpicture}
\caption{Детектор совпадений, порог $\theta=1.5w$. Слева второй импульс пришёл через $\Delta<\Delta_{\max}$, сумма достигла порога, нейрон сработал. Справа $\Delta>\Delta_{\max}$: от первого импульса почти ничего не осталось, и нейрон промолчал.}
\label{fig:coincidence}
\end{figure}

Другой способ получить И --- через длительность, а не через точное время. Пока горит свет, чувствительный нейрон выдаёт импульсы часто и ровно; если порогу мало одного такого входа, а двух хватает, нейрон работает, пока активны оба. Так сеть вычисляет по \emph{уровням}, и точное время прихода импульсов неважно. Большая часть того, что ниже, будет именно так.

\section{Чего сделать нельзя}

Теперь попробуем сделать НЕ: нейрон, который работает, когда вход молчит. Не получится: без входных импульсов в~\eqref{eq:glutneuron} нет ни одного скачка, и напряжение стоит на нуле, ниже порога. А сеть из многих нейронов? Тоже нет, и это доказывается сразу для всех сетей.

Удобнее всего говорить о частотах. Пусть $r_i$ --- частота импульсов нейрона $i$, $I_i$ --- приток от входов, а $f_i$ --- его передаточная характеристика: как частота зависит от суммарного притока. У интегрирующего нейрона эта характеристика неубывающая: больше приток --- чаще импульсы. Частоты в сети, пришедшей в равновесие, удовлетворяют уравнениям
\begin{equation}
r_i \;=\; f_i\Bigl(\sum_j w_{ij}\,r_j + I_i\Bigr), \qquad w_{ij}\ge 0 .
\label{eq:ratenet}
\end{equation}

\begin{keyblock}
\begin{proposition}[Монотонность]
\label{prop:monotone}
Пусть сеть~\eqref{eq:ratenet} состоит только из \nt{ГЛУТ}-нейронов и начинает с полной тишины. Если усилить входы, то установившаяся частота ни одного нейрона не уменьшится. Всё, что вычисляет такая сеть, --- \emph{монотонные} функции входов.
\end{proposition}
\end{keyblock}

Доказательство. Начнём с тишины, $r^{(0)}=0$, и будем подставлять частоты в правую часть~\eqref{eq:ratenet}: $r^{(k+1)}=F\bigl(r^{(k)}\bigr)$. Правая часть $F$ не убывает ни по одному аргументу, потому что веса неотрицательны, а $f_i$ не убывают. Поскольку $r^{(1)}\ge0=r^{(0)}$, по индукции $r^{(k+1)}\ge r^{(k)}$: частоты только растут и приходят к наименьшему решению~\eqref{eq:ratenet}. Если входы $I$ заменить на большие $I'$, то на каждом шаге $r^{(k)}(I)\le r^{(k)}(I')$, и в пределе тоже. Настоящая сеть приходит в равновесие не шагами, а непрерывно, но для неё верно то же самое: при положительных связях неравенство между двумя прогонами, выполненное вначале, сохраняется всё время.

Для отдельных импульсов утверждение верно не буквально: лишний входной импульс может заставить нейрон сработать чуть раньше, и из-за миллисекунды невосприимчивости пропадёт его следующий импульс. Но эффект слабый и ненадёжный, вычислять на нём нельзя. Для частот монотонность строгая.

Следствия те же, что у любой схемы без инверторов. Нет НЕ. Нет исключающего ИЛИ: добавление второго входа не может погасить выход. И самое неприятное: \emph{активность нельзя остановить активностью}, можно только убрать то, что её поддерживает.

\section{Два провода: включение и выключение}

Выход из тупика подсказывает сам организм. Во многих органах чувств раздражитель передаётся двумя наборами нейронов: в зрении одни клетки отвечают на усиление света, другие --- на ослабление~\cite{Schiller1992}. Назовём их \emph{включающими} и \emph{выключающими}. Вместо одного провода $x$ сеть получает пару $(x,\bar x)$, и отсутствие, которое она не умеет вычислить, ей сообщает сам датчик.

С парой проводов НЕ бесплатно: достаточно поменять провода местами. И и ИЛИ делаются каждое двумя нейронами:
\begin{empheq}[box=\keybox]{equation}
\begin{aligned}
x\wedge y \;&\to\; \bigl(\;x\wedge y,\;\; \bar x\vee\bar y\;\bigr),\\
x\vee y \;&\to\; \bigl(\;x\vee y,\;\; \bar x\wedge\bar y\;\bigr).
\end{aligned}
\label{eq:dualrail}
\end{empheq}
Справа все операции монотонны, так что утверждение~\ref{prop:monotone} не нарушено.

Для асинхронной схемы у двухпроводного кода есть второе достоинство, даже более важное. Пара проводов бывает в трёх состояниях: «да», «нет» и, когда оба молчат, \emph{«ответа пока нет»}. Получатель узнаёт, что вычисление закончено, не по часам, а по самому сигналу: в каждой паре заработал ровно один провод. Между раздражителями датчики молчат, и сеть возвращается в тишину, готовая к следующему разу. В асинхронной электронике на этом приёме строят схемы, которые работают правильно при любых задержках элементов~\cite{Sparso2001}.

\begin{keyblock}
\begin{proposition}[Асинхронное вычисление]
\label{prop:dualrail}
Если каждый вход приходит парой проводов «включение---выключение», то сеть из \nt{ГЛУТ}-нейронов может вычислить любую логическую функцию входов. Ответ тоже приходит парами проводов, а его готовность видна по тому, что в каждой паре заработал один провод. Часы для этого не нужны. Цена --- примерно вдвое больше нейронов.
\end{proposition}
\end{keyblock}

Пример --- исключающее ИЛИ: «изменился ровно один из двух раздражителей». Прямой провод ответа --- $(x\wedge\bar y)\vee(\bar x\wedge y)$, обратный --- $(x\wedge y)\vee(\bar x\wedge\bar y)$. Это шесть нейронов, и ни один из них не нуждается в торможении (рис.~\ref{fig:dualxor}).

\begin{figure}[t]
\centering
\begin{tikzpicture}[>=Stealth,x=1cm,y=1cm,
  nrn/.style={circle,draw,very thick,fill=blue!6,minimum size=0.75cm,inner sep=0pt,font=\small},
  inp/.style={font=\small}]
\node[inp] (X)  at (0,3.3) {$x$};
\node[inp] (Xb) at (0,2.2) {$\bar x$};
\node[inp] (Y)  at (0,1.1) {$y$};
\node[inp] (Yb) at (0,0)   {$\bar y$};
\node[nrn] (A1) at (3.2,3.3) {$2$};
\node[nrn] (A2) at (3.2,2.2) {$2$};
\node[nrn] (A3) at (3.2,1.1) {$2$};
\node[nrn] (A4) at (3.2,0)   {$2$};
\node[nrn] (O)  at (7.2,2.75) {$1$};
\node[nrn] (Ob) at (7.2,0.55) {$1$};
\foreach \s/\t in {X/A1,Yb/A1,Xb/A2,Y/A2,X/A3,Y/A3,Xb/A4,Yb/A4}{\draw[->,thick,blue!60!black] (\s) -- (\t);}
\draw[->,thick,blue!60!black] (A1) -- node[pos=0.45,above,sloped,font=\scriptsize,black] {$x\wedge\bar y$} (O);
\draw[->,thick,blue!60!black] (A2) -- node[pos=0.45,below,sloped,font=\scriptsize,black] {$\bar x\wedge y$} (O);
\draw[->,thick,blue!60!black] (A3) -- node[pos=0.45,above,sloped,font=\scriptsize,black] {$x\wedge y$} (Ob);
\draw[->,thick,blue!60!black] (A4) -- node[pos=0.45,below,sloped,font=\scriptsize,black] {$\bar x\wedge\bar y$} (Ob);
\draw[->,very thick,red!70!black] (O) -- (8.6,2.75) node[right,black] {$x\oplus y$: «да»};
\draw[->,very thick,red!70!black] (Ob) -- (8.6,0.55) node[right,black] {$\overline{x\oplus y}$: «нет»};
\node[below,font=\small,gray!40!black] at (3.2,-0.55) {И: порог $2$};
\node[below,font=\small,gray!40!black] at (7.2,-0.55) {ИЛИ: порог $1$};
\node[below,font=\small,gray!40!black] at (0,-0.55) {пары проводов};
\end{tikzpicture}
\caption{Исключающее ИЛИ на двухпроводном коде, только из \nt{ГЛУТ}-нейронов. Число в кружке --- порог в единицах веса одного входа. Четыре нейрона И узнают четыре сочетания входов, два нейрона ИЛИ собирают их в пару проводов ответа.}
\label{fig:dualxor}
\end{figure}

\section{Задержки вместо часов}

Для вычислений со временем достаточно задержек в проводах и детекторов совпадений. Лучший пример --- то, как мозг определяет, откуда пришёл звук.

Звук от источника, расположенного сбоку, приходит в ближнее ухо раньше, чем в дальнее. Разница зависит от направления: от нуля, когда источник прямо впереди, до $0.22\,\text{м}/343\,\text{м/с}\approx0.6\ms$ у человека, когда он точно сбоку. Мозг различает разницу около десяти \emph{микросекунд}~\cite{KlumppEady1956,Grothe2010}, хотя сами нейроны срабатывают с точностью не лучше долей миллисекунды. Как?

Проведём от левого уха провод вдоль ряда детекторов совпадений слева направо, а от правого --- справа налево (рис.~\ref{fig:itd}). Сигнал бежит по проводу со скоростью $v$. До детектора, стоящего на расстоянии $x$ от левого края ряда длины $L$, левый сигнал идёт время $x/v$, правый --- $(L-x)/v$. Если в правое ухо звук пришёл позже на $\delta t$, два сигнала встретятся одновременно в той точке, где
\begin{empheq}[box=\keybox]{equation}
\frac{x}{v} \;=\; \frac{L-x}{v} + \delta t
\quad\Longrightarrow\quad
x \;=\; \frac{L}{2} + \frac{v\,\delta t}{2} .
\label{eq:itd}
\end{empheq}
Сработает только тот детектор, к которому оба сигнала пришли в пределах окна~\eqref{eq:window}. Номер сработавшего детектора и есть направление на источник. Разница во времени превратилась в \emph{место}.

\begin{figure}[t]
\centering
\begin{tikzpicture}[>=Stealth,
  nrn/.style={circle,draw,very thick,fill=blue!6,minimum size=0.7cm,inner sep=0pt}]
\foreach \k in {0,...,6}{\node[nrn] (D\k) at (1.4*\k,0) {\scriptsize $2$};}
\node[nrn,fill=red!15] at (D4) {\scriptsize $2$};
\draw[->,thick,blue!60!black] (-1.4,0.9) node[left,black] {\small левое ухо} -- (9.2,0.9);
\draw[->,thick,orange!85!black] (9.8,-0.9) node[right,black] {\small правое ухо} -- (-0.8,-0.9);
\foreach \k in {0,...,6}{
  \draw[->,blue!60!black] (1.4*\k,0.9) -- (D\k.north);
  \draw[->,orange!85!black] (1.4*\k,-0.9) -- (D\k.south);}
\draw[->,thick,red!70!black] (D4.north east) ++(0.1,0.05) -- ++(0.5,0.45) node[above right,font=\small,black] {сработал};
\node[font=\small,gray!40!black] at (4.2,-1.5) {задержка растёт $\longrightarrow$ \hspace{2.2cm} $\longleftarrow$ задержка растёт};
\pic[scale=1.4] at (-2.3,1.75) {speaker};
\node[font=\scriptsize,align=center] at (-2.3,2.35) {источник слева};
\end{tikzpicture}
\caption{Определение направления на звук. Сигналы от двух ушей бегут навстречу друг другу; каждый кружок --- детектор совпадений с порогом $2$. Если звук пришёл в правое ухо позже, сигналы встречаются правее середины, по формуле~\eqref{eq:itd}. Выход сети --- номер сработавшего нейрона. Здесь источник слева: в левое ухо звук пришёл раньше.}
\label{fig:itd}
\end{figure}

Здесь нет ни часов, ни торможения, ни даже двухпроводного кода: сеть не отвечает «да» или «нет», а указывает на один нейрон из многих. Такая схема из линий задержки и детекторов совпадений, по-видимому, действительно работает в слуховом стволе птиц~\cite{Carr1990}. Микросекундная точность получается не из точности каждого нейрона, а из того, что детекторов много и каждый смотрит на свою задержку. У млекопитающих ряда таких задержек не нашли: там ту же задачу, по-видимому, решают детекторы совпадений, настроенные точно рассчитанным во времени торможением от \nt{ГЛИЦ}-нейронов~\cite{Brand2002,Grothe2010}. Как торможение сужает окно совпадения, мы разберём в главе~\ref{sec:gaba} на примере \nt{ГАМК}; глицин тормозит так же, открывая у получателя хлорные каналы.

\section{Память}

Компьютеру нужна память. В такой сети память --- это активность, которая поддерживает сама себя. Возьмём группу \nt{ГЛУТ}-нейронов, сильно связанных друг с другом, и опишем её одной частотой $r$ (рис.~\ref{fig:recurrentgroup}а). В равновесии $r=f(wr+I)$, где $w$ --- сила обратной связи группы на саму себя, а $I$ --- приток извне. Решим это уравнение графически, пересекая кривую $f(wr+I)$ с прямой $r$ (рис.~\ref{fig:bistable}).

\begin{figure}[t]
\centering
% (b): tau dr/dt = -r + f(w r + I), f as in fig:bistable, w=1, I=0.6; Euler dt=0.001 tau.
\begin{tikzpicture}[>=Stealth,
  nrn/.style={circle,draw,thick,fill=blue!6,minimum size=0.5cm,inner sep=0pt}]
\begin{scope}
\node[font=\small] at (-0.5,1.55) {(а)};
\foreach \k in {1,...,6}{\node[nrn] (n\k) at ({1.4+1.0*cos(60*\k)},{1.0*sin(60*\k)}) {};}
\foreach \a/\b in {1/2,2/3,3/4,4/5,5/6,6/1,1/4,2/5,3/6}{\draw[<->,blue!60!black] (n\a) -- (n\b);}
\foreach \k in {2,3,4}{\draw[->,thick,green!45!black] ($(n\k)+(-0.95,0)$) -- (n\k);}
\node[font=\small,green!45!black] at (-0.35,-0.45) {$I$};
\node[font=\Large] at (3.1,0) {$\approx$};
\node[nrn,minimum size=0.85cm,font=\small] (R) at (4.35,-0.1) {$r$};
\draw[->,thick,blue!60!black] (R) edge[loop above,looseness=6] node[above,font=\small] {$w$} (R);
\draw[->,thick,green!45!black] (3.55,-0.95) node[left,font=\small] {$I$} -- (R);
\end{scope}
\begin{scope}[xshift=6.3cm,y=2.6cm,x=1.05cm]
\node[font=\small] at (-0.6,1.25) {(б)};
\draw[->,gray!70!black] (0,0) -- (7.3,0) node[below left,font=\small,black] {$t/\tau$};
\draw[->,gray!70!black] (0,0) -- (0,1.12) node[left,font=\small,black] {$r$};
\foreach \t in {0,2,4,6}{\draw[gray!60!black] (\t,-0.02) -- (\t,0.02); \node[below,font=\scriptsize] at (\t,0) {\t};}
\draw[densely dotted,gray!60!black] (0,0.5) -- (7,0.5) node[right,font=\scriptsize,black,align=left] {порог\\записи};
\fill[green!45!black,opacity=0.25] (1,-0.2) rectangle (2,-0.13);
\fill[gray!60!black,opacity=0.35] (1,-0.29) rectangle (1.5,-0.22);
\draw[very thick,blue!60!black] plot coordinates {(0.0,0.002) (0.1,0.002) (0.2,0.002) (0.3,0.002) (0.4,0.002) (0.5,0.002) (0.6,0.002) (0.7,0.002) (0.8,0.002) (0.9,0.002) (1.0,0.002) (1.1,0.083) (1.2,0.165) (1.3,0.242) (1.4,0.313) (1.5,0.378) (1.6,0.437) (1.7,0.491) (1.8,0.539) (1.9,0.583) (2.0,0.623) (2.1,0.643) (2.2,0.665) (2.3,0.688) (2.4,0.71) (2.5,0.732) (2.6,0.753) (2.7,0.773) (2.8,0.792) (2.9,0.81) (3.0,0.826) (3.2,0.855) (3.4,0.88) (3.6,0.9) (3.8,0.917) (4.0,0.931) (4.4,0.953) (4.8,0.967) (5.2,0.977) (5.6,0.984) (6.0,0.988) (6.5,0.992) (7.0,0.994)};
\draw[very thick,gray!60!black,densely dashed] plot coordinates {(0.0,0.002) (0.1,0.002) (0.2,0.002) (0.3,0.002) (0.4,0.002) (0.5,0.002) (0.6,0.002) (0.7,0.002) (0.8,0.002) (0.9,0.002) (1.0,0.002) (1.1,0.083) (1.2,0.165) (1.3,0.242) (1.4,0.313) (1.5,0.378) (1.6,0.358) (1.7,0.336) (1.8,0.313) (1.9,0.291) (2.0,0.269) (2.2,0.228) (2.4,0.191) (2.6,0.159) (2.8,0.133) (3.0,0.11) (3.2,0.091) (3.4,0.076) (3.6,0.063) (3.8,0.052) (4.0,0.043) (4.4,0.03) (4.8,0.021) (5.2,0.015) (5.6,0.011) (6.0,0.008) (6.5,0.006) (7.0,0.004)};
\node[font=\scriptsize,blue!60!black,anchor=west] at (3.3,0.8) {приток $1\tau$: записано};
\node[font=\scriptsize,gray!50!black,anchor=west] at (2.3,0.3) {приток $0.5\tau$: погасло};
\end{scope}
\end{tikzpicture}
\caption{Группа, сильно связанная сама с собой. (а)~Много \nt{ГЛУТ}-нейронов, возбуждающих друг друга, ведут себя как один нейрон с частотой $r$, который возбуждает сам себя с силой $w$; извне приходит приток $I$. (б)~Частота группы во времени при $w=1$ и той же кривой $f$, что на рис.~\ref{fig:bistable}. Приток $I=0.6$ подаётся на время, отмеченное полоской под осью. Если он держится $1\tau$, группа успевает перейти неустойчивую точку $r=0.5$, разгорается и остаётся гореть, когда приток кончился. Если только $0.5\tau$, она не дотягивает до этой точки и гаснет обратно: чтобы записать бит, приток должен длиться дольше примерно $0.7\tau$.}
\label{fig:recurrentgroup}
\end{figure}

\begin{figure}[t]
\centering
\begin{tikzpicture}[>=Stealth,x=5cm,y=3.2cm]
\draw[->,gray!70!black] (0,0) -- (1.1,0) node[below left,black] {\small $r$};
\draw[->,gray!70!black] (0,0) -- (0,1.1);
\draw[thick,gray!60!black] (0,0) -- (1.05,1.05) node[right,black] {\small $r$};
\draw[very thick,blue!60!black] plot[domain=0:1.05,samples=80] (\x,{1/(1+exp(-(\x-0.5)/0.08))});
\draw[very thick,red!70!black,densely dashed] plot[domain=0:1.05,samples=80] (\x,{1/(1+exp(-(\x+0.3-0.5)/0.08))});
\fill[blue!60!black] (0.002,0.002) circle (0.018);
\draw[blue!60!black,fill=white,thick] (0.5,0.5) circle (0.018);
\fill[blue!60!black] (0.998,0.998) circle (0.018);
\node[below right,font=\small] at (0.02,0.0) {выкл.};
\node[right,font=\small] at (0.52,0.46) {неустойчиво};
\node[below,font=\small] at (0.99,0.96) {вкл.};
\node[anchor=east,font=\small,red!70!black] at (0.19,0.62) {$f(wr+I)$};
\node[anchor=west,font=\small,blue!60!black] at (0.45,0.1) {$f(wr)$};
\end{tikzpicture}
\caption{Память из группы \nt{ГЛУТ}-нейронов с сильной обратной связью. Сплошная кривая --- без внешнего притока: два устойчивых пересечения с прямой $r$, «выключено» и «включено», и неустойчивое между ними. Короткий приток $I$ (штриховая кривая) оставляет только верхнее.}
\label{fig:bistable}
\end{figure}

Если обратная связь сильная, пересечений три: нижнее --- тишина, верхнее --- группа горит на полную, среднее неустойчиво. Получился элемент с двумя устойчивыми состояниями, триггер. Записать в него единицу легко: короткий приток поднимает кривую, нижнее пересечение исчезает, и группа разгорается --- и остаётся гореть, когда приток кончился (рис.~\ref{fig:recurrentgroup}б). Слишком короткий приток не успевает довести группу до неустойчивой точки, и она гаснет обратно. Такая самоподдерживающаяся активность, которая держится секунды после того, как раздражитель исчез, действительно наблюдается в мозге и считается основой кратковременной памяти~\cite{Wang2001}. Но это не единственный способ помнить секунды. Запись можно хранить и в медленной переменной самих синапсов: после пачки импульсов облегчение, как у синапса-«тире» главы~\ref{sec:code}, держится около секунды, и сеть между короткими вспышками может почти молчать --- не триггер, а заряженный конденсатор~\cite{Mongillo2008}. Такие «тихие» промежутки при удержании в памяти наблюдаются~\cite{Stokes2015}, и хранение без постоянных импульсов дешевле по энергии. Какой из двух способов кора использует в каком случае, обсуждается.

А как стереть? По утверждению~\ref{prop:monotone} активностью --- никак; остаётся что-то \emph{убрать}. Первый способ --- убрать поддержку: если группа горит только вместе с притоком от другой группы, память гаснет вместе с ним. Второй --- усталость. У настоящих синапсов при долгой работе запас медиатора истощается~\cite{Tsodyks1997}, а у нейронов нарастают медленные токи, снижающие возбудимость. Обратная связь слабеет, верхнее пересечение исчезает, и группа гаснет сама через секунды. Получается память, которая включается сигналом, а выключается только временем.

\section{Последовательности}

Вместо часов можно иметь \emph{таймер, запускаемый событием}. Соединим группы \nt{ГЛУТ}-нейронов в цепочку: первая возбуждает вторую, вторая --- третью, и так далее. Внешний сигнал зажигает первую группу, и волна активности пробегает по цепочке, каждое звено через одну-две задержки после предыдущего и только один раз: сразу после импульса нейроны невосприимчивы, а волна уже ушла дальше.

Такая цепочка отмеряет время от события: сработало звено номер $k$ --- значит, с момента запуска прошло $k$ задержек. Её выходы можно подать на мышцы и получить последовательность движений, которая разворачивается сама. У певчих птиц во время пения отдельные \nt{ГЛУТ}-нейроны одного из отделов мозга срабатывают строго по очереди, каждый в свой момент песни, --- очень похоже на такую цепочку~\cite{Hahnloser2002}.

В отличие от часов, цепочка молчит, пока её не запустили, отрабатывает один раз и отмеряет время только для тех, кто к ней подключён. Общего такта у сети по-прежнему нет.

\section{Чего не хватает}

Из одних \nt{ГЛУТ}-нейронов, без торможения, получается многое. Но трёх вещей не хватает, и все три из-за утверждения~\ref{prop:monotone}.

\emph{Выбор.} Сеть должна выбирать одно направление, одно действие. Если два раздражителя почти равны, в сети рис.~\ref{fig:itd} сработают оба выхода, и подавить слабый в пользу сильного нечем.

\emph{Сброс.} Стереть память сигналом нельзя.

\emph{Устойчивость.} В сети, где связи возникают не по плану инженера, петли положительной обратной связи неизбежны, и активность в них может нарастать лавиной (глава~\ref{sec:neurontypes}). Единственный тормоз --- та же усталость, медленная и грубая.

Все три беды лечатся одним лекарством --- нейроном, который гасит импульс импульсом. Это \nt{ГАМК}, и нужен он не для того, чтобы вычислять, а для того, чтобы выбирать, сбрасывать и держать вычисление под контролем.

\section{Задачи}

\begin{exercise}[\lvlA]
\label{ex:02:1}
\emph{Ряд детекторов для звука.} Сигналы от ушей бегут по проводам рис.~\ref{fig:itd} со скоростью $v=5\,\text{м/с}$; наибольшая разница времён прихода --- $0.22\,\text{м}/343\,\text{м/с}$.
\par\noindent(а)~Какой длины $L$ должен быть ряд, чтобы покрыть все направления? С каким шагом ставить детекторы, чтобы соседние отвечали разницам $\delta t$, отличающимся на $10$~мкс, и сколько их нужно?
\par\noindent(б)~Детекторы --- нейроны~\eqref{eq:glutneuron} с $\tau=0.5\ms$ и $\theta=1.5w$, на каждый вход приходит по одному импульсу. Сколько детекторов сработает на один звук? Что это значит для утверждения «номер сработавшего детектора и есть направление» и как получатель может извлечь из ряда точность в $10$~мкс?
\end{exercise}

\begin{exercise}[\lvlB]
\label{ex:02:2}
\emph{Неравные входы смещают окно.} Детектор совпадений~\eqref{eq:glutneuron} получает по одному импульсу от двух входов с весами $w_1$ и $w_2$, $w_1,w_2<\theta<w_1+w_2$; второй импульс приходит на $\delta=t_2-t_1$ позже первого ($\delta$ может быть отрицательным).
\par\noindent(а)~Выведите окно совпадения: множество $\delta$, при которых нейрон срабатывает. Проверьте, что при $w_1=w_2$ получается~\eqref{eq:window}.
\par\noindent(б)~Найдите центр окна $c$ и вычислите окно и $c$ при $w_1=1$, $w_2=0.8$, $\theta=1.5$, $\tau=0.5\ms$. Какой вход должен прийти раньше, чтобы ответ был наиболее надёжным, и почему?
\par\noindent(в)~Во всём ряду рис.~\ref{fig:itd} вход от левого уха на $20\,\%$ слабее правого. Как это сказывается на найденном направлении, и сравните ошибку с разрешением около десяти микросекунд. При каком $\theta$ смещение исчезает и чем за это приходится платить?
\end{exercise}

\begin{exercise}[\lvlB]
\label{ex:02:3}
\emph{Монотонность в непрерывном времени.} Частоты подчиняются уравнениям $\tau\,\dd r_i/\dd t=-r_i+f_i\bigl(\sum_jw_{ij}r_j+I_i(t)\bigr)$, где $f_i$ не убывают и удовлетворяют условию Липшица с константой $\ell$.
\par\noindent(а)~Пусть $w_{ij}\ge0$ при $i\ne j$ (знак $w_{ii}$ любой). Докажите: если $r(0)\le r'(0)$ и $I(t)\le I'(t)$ покомпонентно, то $r(t)\le r'(t)$ при всех $t>0$. Указание: рассмотрите $z=r'-r$ и $\sum_i\max(-z_i,0)$.
\par\noindent(б)~Выведите отсюда утверждение~\ref{prop:monotone} для сети, стартующей из тишины, без итераций по шагам.
\par\noindent(в)~Знаки весов произвольны. Докажите, что замена $r_i\to s_ir_i$, $s_i=\pm1$, переводит систему в систему из пункта~(а) тогда и только тогда, когда в каждом цикле графа связей (без учёта направления стрелок, без петель $i\to i$) чётное число отрицательных связей. Примените это к двум группам с взаимным торможением и к петле «\nt{ГЛУТ} возбуждает \nt{ГАМК}, \nt{ГАМК} тормозит \nt{ГЛУТ}» из главы~\ref{sec:gaba}. Что говорит результат о том, какая из них может колебаться?
\end{exercise}

\begin{exercise}[\lvlB]
\label{ex:02:4}
\emph{Проектирование: двухпроводный сумматор.} Полный сумматор складывает три бита $x,y,z$: выдаёт бит суммы $S=x\oplus y\oplus z$ и бит переноса $C$ --- единицу, если единиц не меньше двух.
\par\noindent(а)~Докажите: нейрон с неотрицательными весами, получающий пары проводов $(x_i,\bar x_i)$, на полных входах вычисляет любую пороговую функцию $\bigl[\sum_ic_ix_i\ge\theta\bigr]$ с весами $c_i$ \emph{любого} знака.
\par\noindent(б)~Постройте двухпроводный полный сумматор (пары $(S,\bar S)$ и $(C,\bar C)$) не более чем из восьми \nt{ГЛУТ}-нейронов, указав веса и пороги. Указание: начните с нейронов $g_k=[x+y+z\ge k]$.
\par\noindent(в)~Входы приходят асинхронно: пока какая-то пара молчит, её значение ещё не пришло. Докажите, что в вашей схеме ни один провод ответа не сработает неверно ни при каком порядке прихода входов.
\par\noindent(г)~Сколько нейронов потребовала бы схема из элементов И и ИЛИ, как на рис.~\ref{fig:dualxor}?
\end{exercise}

\begin{exercise}[\lvlB]
\label{ex:02:5}
\emph{Моделирование: сколько длится запись.} Возьмите модель рис.~\ref{fig:recurrentgroup}: $\tau\,\dd r/\dd t=-r+f(wr+I(t))$, $f(u)=1/\bigl(1+e^{-(u-0.5)/0.08}\bigr)$, $w=1$; до записи группа в нижнем устойчивом состоянии $r_0$; приток $I$ подаётся на время $T$.
\par\noindent(а)~Интегрируя методом Эйлера (шаг $10^{-4}\tau$), воспроизведите рис.~\ref{fig:recurrentgroup}б и найдите $r_0$.
\par\noindent(б)~Найдите наименьшую длительность записи $T_{\min}(I)$ при $I$ от $0.25$ до $2$ и постройте её график. Чему равна $T_{\min}(0.6)$?
\par\noindent(в)~Докажите, что при любом $I$ $T_{\min}>\tau\ln\bigl[(1-r_0)/(1-r_u)\bigr]$, где $r_u=0.5$ --- неустойчивое равновесие, и что при $I\to\infty$ $T_{\min}$ стремится к этой границе.
\par\noindent(г)~Найдите $I_c$, ниже которого бит не записывается ни при какой длительности, и проверьте численно, как растёт $T_{\min}$ при $I\to I_c$. Какой смысл имеют оба предела для электронщика, привыкшего к триггерам?
\end{exercise}

\begin{exercise}[\lvlA]
\label{ex:02:6}
\emph{Цепочка как таймер.} Звено цепочки из раздела о последовательностях --- группа из $100$ \nt{ГЛУТ}-нейронов; задержка звена $2\ms$ со случайным разбросом $0.2\ms$, независимым от звена к звену.
\par\noindent(а)~Сколько нужно звеньев и нейронов, чтобы отмерить $1\,\text{с}$? Сколько импульсов тратится на один отсчёт?
\par\noindent(б)~Каков разброс момента, когда срабатывает последнее звено, и его относительная величина? То же для $100\ms$. Как относительная ошибка зависит от отмеряемого интервала?
\par\noindent(в)~Допустим, измерение показывает, что относительная ошибка таймера не зависит от интервала и равна $5\,\%$. Какой источник разброса в цепочке дал бы такой результат?
\end{exercise}

\begin{exercise}[\lvlC]
\label{ex:02:7}
\emph{Исследование: триггер или конденсатор.} В разделе о памяти остался открытым вопрос, как удерживать бит секунды: самоподдерживающейся активностью группы или медленной переменной синапсов, которую изредка освежают короткие пачки.
\par\noindent(а)~Группа из $N=100$ нейронов удерживает бит $10\,\text{с}$. Триггер: все нейроны работают с частотой $20$ импульсов в секунду. Конденсатор: облегчение $u$ спадает с постоянной времени $\tau_F=1\,\text{с}$, бит читается, пока $u\ge u_{\max}/2$, и каждое освежение --- пачка из трёх импульсов у каждого нейрона. Найдите частоту освежений, число импульсов и энергию при стоимости импульса $10^{-10}$~Дж (глава~\ref{sec:code}).
\par\noindent(б)~На $200\ms$ всю группу глушат общим сигналом. Что происходит с битом в каждой модели?
\par\noindent(в)~Открытая часть. Постройте обе модели в виде уравнений для частот (для второй добавьте уравнение для $u$, умножающего вес обратной связи) и предложите опыт с короткими неспецифическими импульсами на группу, который различает две модели по отклику, а не по фоновой активности. Какие измерения внутри одной и той же модели мешают сделать вывод, и какие результаты допускают смешанный механизм?
\end{exercise}

\chapter{\nt{ГЛУТ} и \nt{ГАМК} вместе}
\label{sec:gaba}

\section{Правила игры}

Сеть из одних \nt{ГЛУТ}-нейронов не умеет выбирать, сбрасывать память и удерживать активность от лавины, и всё из-за монотонности (утверждение~\ref{prop:monotone}). Добавим в неё второй по численности тип нейронов --- \nt{ГАМК}. Правила прежние: только то, что бывает в живом организме, и никаких часов.

Модель нейрона та же, что~\eqref{eq:glutneuron}, но импульс от \nt{ГАМК}-нейрона не поднимает напряжение получателя, а опускает. Веса теперь бывают двух знаков, и знак задаёт отправитель, правило~\eqref{eq:dalesign}.

Несколько параметров схемы. \nt{ГАМК}-нейронов в коре от пятой части до четверти~\cite{Hendry1987}. Выходы у них короткие: они тормозят соседей, а не далёкие области. Торможение приходит через миллисекунду и длится несколько миллисекунд. Без входа \nt{ГАМК}-нейрон молчит: чтобы затормозить кого-то, он сам должен получить возбуждение, обычно от \nt{ГЛУТ}-нейронов той же сети.

Поэтому отрицательную связь от \nt{ГЛУТ}-нейрона $A$ к нейрону $B$ приходится делать через посредника: $A$ возбуждает \nt{ГАМК}-нейрон, а тот тормозит $B$. Смена знака стоит одного нейрона и одной задержки, около миллисекунды.

Для торможения важно не только, \emph{от кого} идёт связь, но и \emph{куда} она садится: место посадки во многом определяет, что связь делает~\cite{Markram2004}. \nt{ГЛУТ}-выходы садятся в основном на входные ветви получателя, его \emph{дендриты}, и несут данные: каждый такой вход --- одно слагаемое суммы. Выход на тело нейрона действует на сумму целиком: так многие быстрые \nt{ГАМК}-нейроны делят ответ получателя и задают, когда ему сработать. Выход на то место, где рождается импульс получателя, --- последнее слово, запрет на весь выход сразу. А выход на окончание чужого провода меняет вероятность выброса $p$ одной входной линии, не трогая остальных~\cite{Rudomin1999}; так, в частности, работают ворота боли в главе~\ref{sec:pain}. За пределами мозга выходы садятся на мышечные волокна (глава~\ref{sec:muscles}) и на органы (глава~\ref{sec:chemoutput}), а некоторые не садятся никуда и выбрасывают медиатор в объём ткани или в кровь (главы~\ref{sec:serocomputer} и~\ref{sec:chemoutput}).

\section{НЕ и запрет}

Можно ли теперь сделать НЕ? Почти. Нейрон, который работает при молчащем входе, должен откуда-то брать возбуждение. В живом мозге оно есть: на каждый нейрон постоянно приходит фоновая активность от тысяч других. Пусть фон держит нейрон $Y$ в работе, а вход $x$ через \nt{ГАМК}-посредника его тормозит. Это НЕ.

Полезнее, однако, \emph{запрет}: «$a$, но не при $b$»:
\begin{empheq}[box=\keybox]{equation}
y \;=\; a \wedge \bar b .
\label{eq:veto}
\end{empheq}
Нейрон $Y$ получает возбуждение от $a$ и торможение от \nt{ГАМК}-нейрона, которого возбуждает $b$. Фон здесь не нужен: возбуждение даёт сам $a$. Из запретов и ИЛИ собирается всё. Например, исключающее ИЛИ: $x\oplus y=(x\wedge\bar y)\vee(\bar x\wedge y)$ --- два нейрона-запрета, два \nt{ГАМК}-посредника и один нейрон ИЛИ.

\begin{keyblock}
\begin{proposition}[Полнота \nt{ГЛУТ}+\nt{ГАМК}]
\label{prop:gabacomplete}
Сеть из \nt{ГЛУТ}- и \nt{ГАМК}-нейронов может вычислить любую логическую функцию входов, и каждый бит передаётся одним проводом. Датчики «включения---выключения» из главы~\ref{sec:glutcomputer} для этого больше не нужны. Если функция должна давать единицу при всех молчащих входах, нужен источник фоновой активности.
\end{proposition}
\end{keyblock}

Доказательство: запрет~\eqref{eq:veto} вместе с ИЛИ функционально полон, если есть постоянная единица, потому что НЕ $x$ --- это запрет «единица, но не при $x$». А функции, которые при всех молчащих входах дают ноль, собираются из запретов и ИЛИ и без единицы.

\section{Направление движения}

Запрет во времени, как и И во времени в главе~\ref{sec:glutcomputer}, даёт детектор направления движения.

Пусть два соседних датчика, $A$ и $B$, стоят на расстоянии $\Delta x$ друг от друга. Выходной нейрон $Y$ получает возбуждение от $B$ и торможение от $A$ через \nt{ГАМК}-посредника, с задержкой $d$ (рис.~\ref{fig:direction}). Пятно света, бегущее со скоростью $v$ от $A$ к $B$, проходит $A$ в момент $0$, и торможение приходит к $Y$ в момент $d$; $B$ оно проходит в момент $\Delta x/v$, и тогда же приходит возбуждение. Если эти моменты совпадают с точностью до окна~\eqref{eq:window}, торможение гасит возбуждение, и $Y$ молчит. Это происходит при скорости около
\begin{empheq}[box=\keybox]{equation}
v^* \;=\; \frac{\Delta x}{d} .
\label{eq:vstar}
\end{empheq}
При движении от $B$ к $A$ возбуждение от $B$ приходит в момент $0$, а торможение от $A$ --- только в момент $\Delta x/v+d$, когда $Y$ уже сработал.

\begin{figure}[t]
\centering
\begin{tikzpicture}[>=Stealth,
  nrn/.style={circle,draw,very thick,fill=blue!6,minimum size=0.9cm,inner sep=0pt},
  gab/.style={nrn,fill=red!10},
  inh/.style={-{Bar[width=7pt]},thick,red!70!black}]
\node[nrn] (A) at (0,2) {$A$};
\node[nrn] (B) at (3,2) {$B$};
\node[gab] (G) at (0.9,0.6) {\small \nt{ГАМК}};
\node[nrn] (Y) at (3,-0.6) {$Y$};
\draw[->,thick] (A) -- (G);
\draw[inh] (G) -- node[below left=-2pt] {\scriptsize задержка $d$} (Y);
\draw[->,thick] (B) -- (Y);
\draw[->,very thick,red!70!black] (Y) -- (4.6,-0.6) node[right,black] {\small выход};
\draw[->,thick,orange!85!black] (-0.6,2.9) -- (3.6,2.9) node[right,black] {\small молчит};
\draw[<-,thick,orange!85!black] (-0.6,3.4) -- (3.6,3.4) node[right,black] {\small отвечает};
\foreach \yy/\dd in {2.9/-1,3.4/1}{
  \foreach \k in {1,2,3}{\draw[orange!70!yellow,line width=0.8pt] (1.5+\dd*0.12+\dd*0.13*\k,\yy+0.1-0.1*\k+0.1) -- ++(\dd*0.25,0);}
  \fill[yellow!70!orange,draw=orange!70!black] (1.5,\yy) circle (0.14);}
\node[font=\scriptsize,align=center] at (1.5,3.95) {пятно света};
\end{tikzpicture}
\caption{Детектор направления движения. $B$ возбуждает $Y$, $A$ тормозит его через \nt{ГАМК}-посредника с задержкой $d$. При движении от $A$ к $B$ со скоростью около $\Delta x/d$ торможение гасит возбуждение; при обратном движении оно опаздывает. Жёлтое пятно с чёрточками --- свет, бегущий над датчиками.}
\label{fig:direction}
\end{figure}

В сетчатке избирательность к направлению движения, по-видимому, создаётся именно так: асимметричным торможением со стороны, откуда движение не должно вызывать ответ~\cite{Barlow1965}.

\section{Вычитание или деление}

До сих пор торможение у нас вычитало. На самом деле оно устроено тоньше.

Синапс открывает канал, то есть включает проводимость. У возбуждающего синапса равновесный потенциал лежит далеко выше порога, примерно на $70\mV$ выше покоя: открытый канал тянет напряжение вверх. У тормозящего \nt{ГАМК}-синапса канал пропускает хлор, равновесный потенциал которого лежит около покоя, и открытый канал тянет напряжение не вниз, а \emph{к покою}. Если отсчитывать напряжение от покоя, установившееся напряжение мембраны с проводимостью утечки $g_L$, возбуждающей $g_E$ и тормозящей $g_I$ равно
\begin{empheq}[box=\keybox]{equation}
V_\infty \;=\; \frac{g_E\,E_E}{g_L+g_E+g_I} ,
\label{eq:shunt}
\end{empheq}
где $E_E\approx70\mV$ --- равновесный потенциал возбуждающего синапса. Это закон Ома для делителя напряжения: $g_E$ тянет к $E_E$, а $g_L$ и $g_I$ --- к нулю.

$g_I$ стоит не в числителе со знаком минус, а в знаменателе: торможение не \emph{вычитает} из возбуждения, а \emph{делит} его (рис.~\ref{fig:shunt}). Такое торможение называют шунтирующим: открытые хлорные каналы закорачивают мембрану. Для сети это регулировка усиления. Если $g_I$ пропорциональна суммарной активности соседей, каждый нейрон отвечает не на абсолютную силу своего входа, а на её долю в общей активности. Такое деление на общую активность встречается во многих отделах мозга и считается одной из его стандартных операций~\cite{Carandini2012}: одна и та же сеть работает и при слабом, и при сильном входе.

Оговорка: формула~\eqref{eq:shunt} --- про нейрон, который не выдаёт импульсов. У срабатывающего нейрона напряжение всё время держится около порога, ток через тормозящую проводимость почти постоянен, и на \emph{частоту} импульсов шунт действует в основном вычитанием~\cite{HoltKoch1997}. Частота делится, если нейрону одновременно добавить уравновешенные возбуждающую и тормозящую фоновые проводимости, как в дрожащем, уравновешенном режиме живой коры~\cite{Chance2002}. Так что деление мозг получает не от одного шунта, а от торможения вместе с фоновым шумом~\cite{Carandini2012}.

\begin{figure}[t]
\centering
\begin{tikzpicture}[>=Stealth,x=1.25cm,y=0.05cm]
\draw[->,gray!70!black] (0,0) -- (5.4,0) node[right,black] {\small $g_E/g_L$};
\draw[->,gray!70!black] (0,0) -- (0,68) node[above,black] {\small $V_\infty$, мВ};
\foreach \v in {20,40,60}{\draw[gray!60!black] (-0.05,\v) -- (0.05,\v) node[left=3pt,font=\scriptsize] {\v};}
\foreach \x in {1,2,3,4,5}{\draw[gray!60!black] (\x,-1.5) -- (\x,1.5) node[below=5pt,font=\scriptsize] {\x};}
\draw[very thick,blue!60!black] plot[domain=0:5,samples=60] (\x,{70*\x/(1+\x)});
\draw[very thick,red!70!black] plot[domain=0:5,samples=60] (\x,{70*\x/(3+\x)});
\draw[thick,densely dashed,gray!50!black] plot[domain=0.56:5,samples=60] (\x,{70*\x/(1+\x)-25});
\node[right,font=\small,blue!60!black] at (5.02,58.3) {без торможения};
\node[right,font=\small,red!70!black] at (5.02,45.5) {делит: $g_I=2g_L$};
\node[right,font=\small,gray!40!black] at (5.02,32.5) {вычитает $25\mV$};
\end{tikzpicture}
\caption{Шунтирующее торможение делит напряжение, а не вычитает из него. Синяя кривая --- установившееся напряжение~\eqref{eq:shunt} без торможения, красная --- с тормозящей проводимостью $g_I=2g_L$. Красная --- та же синяя, растянутая по горизонтали втрое: вход как будто поделили на $1+g_I/g_L=3$. Штриховая --- торможение, вычитающее постоянные $25\mV$: кривая опускается целиком, наклон не меняется.}
\label{fig:shunt}
\end{figure}

Есть и второе следствие. Постоянная времени мембраны $\tau_{\text{эфф}}=C/(g_L+g_E+g_I)$, и торможение её укорачивает. А окно совпадения~\eqref{eq:window} пропорционально $\tau$, так что торможение, пришедшее вместе с возбуждением, \emph{сужает окно}. Схема, в которой вход возбуждает нейрон напрямую и одновременно, с задержкой в миллисекунду, тормозит его через \nt{ГАМК}-посредника, --- одна из самых распространённых в мозге: нейрон успевает ответить только на то, что пришло в первую миллисекунду-другую~\cite{Pouille2001}.

\section{Выбор}

Теперь главное, чего не хватало \nt{ГЛУТ}-сети: выбор одного из нескольких. Возьмём две группы \nt{ГЛУТ}-нейронов с входами $I_1$ и $I_2$. Каждая через своих \nt{ГАМК}-посредников тормозит другую с силой $\beta$ (рис.~\ref{fig:wta}). Для частот $r_1,r_2$ получаем
\begin{equation}
\tau\,\frac{\dd r_1}{\dd t} = -r_1 + \bigl[I_1-\beta r_2\bigr]_+ ,\qquad
\tau\,\frac{\dd r_2}{\dd t} = -r_2 + \bigl[I_2-\beta r_1\bigr]_+ ,
\label{eq:wta}
\end{equation}
где $[u]_+=\max(u,0)$: частота не бывает отрицательной.

\begin{figure}[t]
\centering
\begin{tikzpicture}[>=Stealth,
  nrn/.style={circle,draw,very thick,fill=blue!6,minimum size=0.9cm,inner sep=0pt},
  gab/.style={nrn,fill=red!10,minimum size=0.75cm},
  inh/.style={-{Bar[width=7pt]},thick,red!70!black}]
\node[nrn] (E1) at (0,0) {$r_1$};
\node[nrn] (E2) at (4,0) {$r_2$};
\node[gab] (G1) at (1.4,1.1) {};
\node[gab] (G2) at (2.6,-1.1) {};
\draw[->,thick] (E1) -- (G1); \draw[inh] (G1) -- (E2);
\draw[->,thick] (E2) -- (G2); \draw[inh] (G2) -- (E1);
\draw[->,thick,blue!60!black] (-1.6,0) node[left,black] {$I_1$} -- (E1);
\draw[->,thick,blue!60!black] (5.6,0) node[right,black] {$I_2$} -- (E2);
\end{tikzpicture}
\caption{Выбор из двух. Две группы \nt{ГЛУТ}-нейронов (синие) тормозят друг друга через \nt{ГАМК}-посредников (красные).}
\label{fig:wta}
\end{figure}

Пока работают оба нейрона, система~\eqref{eq:wta} линейна, и её поведение задают два собственных числа: $-(1+\beta)/\tau$ для одинаковых отклонений обеих частот и $-(1-\beta)/\tau$ для противоположных. При слабом торможении, $\beta<1$, оба числа отрицательны: оба нейрона работают, торможение лишь приглушает слабого. При сильном торможении, $\beta>1$, второе число становится положительным: любая разница между $r_1$ и $r_2$ нарастает, пока один нейрон не замолчит совсем.

\begin{keyblock}
\begin{proposition}[Победитель забирает всё]
\label{prop:wta}
Если взаимное торможение сильнее единицы, $\beta>1$, состояние, в котором работают обе группы, неустойчиво. Сеть приходит в состояние, где работает одна группа, а другая молчит. Победитель с частотой $r_1=I_1$ удерживает победу, пока $I_2<\beta I_1$.
\end{proposition}
\end{keyblock}

Мы проверили это на модели. При $\beta=0.5$ и входах $1.0$ и $0.9$ обе группы работают, с частотами $0.73$ и $0.53$. При $\beta=2$ и тех же входах вторая группа молчит совсем. У такого выбора есть память: если потом поменять входы местами, прежний победитель остаётся победителем, потому что $1.0<2\cdot0.9$.

С сотней групп всё так же: каждая тормозит остальных, выживает одна.

\section{Сброс памяти}

Память из главы~\ref{sec:glutcomputer} выключалась только усталостью. Добавим к ней \nt{ГАМК}-нейрон, который возбуждается сигналом «сброс» и тормозит группу. Торможение опускает кривую $f(wr+I)$ на рис.~\ref{fig:bistable}, верхнее пересечение исчезает, и группа гаснет --- и остаётся в нижнем состоянии, когда сброс кончился. Теперь память записывается одним сигналом и стирается другим, как у триггера со входами установки и сброса.

Ещё красивее соединить две такие группы взаимным торможением, как на рис.~\ref{fig:wta}, и дать каждой обратную связь на саму себя. Сигнал на вход одной из групп переводит ячейку в её состояние, и ячейка помнит последнее решение. Это прямой аналог защёлки из двух перекрёстно соединённых элементов ИЛИ-НЕ.

\section{Устойчивость и ритм}

Остаётся третья беда \nt{ГЛУТ}-сети --- лавина. Возьмём группу \nt{ГЛУТ}-нейронов с обратной связью на себя $w_{EE}$ и группу \nt{ГАМК}-нейронов, которых первая возбуждает с силой $w_{IE}$ и которые тормозят первую с силой $w_{EI}$. Для частот $E$ и $I$ получаем
\begin{equation}
\tau_E\frac{\dd E}{\dd t} = -E + w_{EE}E - w_{EI}I + h, \qquad
\tau_I\frac{\dd I}{\dd t} = -I + w_{IE}E ,
\label{eq:ei}
\end{equation}
где $h$ --- внешний вход~\cite{WilsonCowan1972}. Без торможения при $w_{EE}>1$ активность растёт без предела: каждый импульс порождает больше одного следующего. С торможением установившаяся частота
\begin{equation}
E^* \;=\; \frac{h}{1-w_{EE}+w_{EI}w_{IE}}
\label{eq:eistar}
\end{equation}
конечна, если знаменатель положителен. Но этого мало: равновесие должно ещё и притягивать. Из линейного анализа~\eqref{eq:ei} получаются два условия.

\begin{keyblock}
\begin{proposition}[Условия устойчивости]
\label{prop:ei}
Равновесие~\eqref{eq:eistar} устойчиво, если торможение
\begin{enumerate}
\item[(i)] \emph{достаточно сильное}: $w_{EI}\,w_{IE} > w_{EE}-1$;
\item[(ii)] \emph{достаточно быстрое}: $\tau_I < \tau_E/(w_{EE}-1)$.
\end{enumerate}
Если нарушено (i), активность нарастает лавиной. Если выполнено (i), но нарушено (ii), торможение запаздывает, и активность начинает колебаться.
\end{proposition}
\end{keyblock}

Условие~(i) --- это определитель матрицы системы~\eqref{eq:ei}, условие~(ii) --- её след. Смысл второго понятен любому, кто настраивал регулятор: запоздавшая обратная связь не гасит отклонение, а раскачивает его.

\begin{figure}[t]
\centering
\begin{tikzpicture}[>=Stealth]
\draw[->,gray!70!black] (0,0) -- (10.9,0) node[below left,black] {\small $t$, мс};
\draw[->,gray!70!black] (0,0) -- (0,3.3) node[left,black] {\small $E$};
\foreach \t in {0,50,100,150}{\draw[gray!70!black] (\t*0.07,0.06) -- (\t*0.07,-0.06) node[below,font=\scriptsize] {\t};}
\input{figures/ei_traces}
\node[font=\small,gray!40!black,anchor=west] at (0.9,3.15) {без торможения: лавина};
\node[font=\small,blue!60!black,anchor=west] at (7.4,1.25) {быстрое торможение};
\node[font=\small,red!70!black,anchor=west] at (7.4,2.55) {медленное торможение};
\end{tikzpicture}
\caption{Частота \nt{ГЛУТ}-группы с обратной связью $w_{EE}=1.5$, $\tau_E=10\ms$, после включения входа. Без торможения (штрих) активность растёт без предела. С торможением силы $w_{EI}w_{IE}=2$ и $\tau_I=2\ms$ (синяя кривая) она выходит на уровень $E^*=h/(1-1.5+2)=0.67h$. С тем же торможением, но медленным, $\tau_I=30\ms$ (красная кривая), активность колеблется.}
\label{fig:ei}
\end{figure}

Считается, что кора работает именно в таком режиме: её собственная обратная связь достаточно сильна, чтобы без торможения уйти в лавину, и держит её в рабочей точке только быстрое торможение~\cite{Tsodyks1997isn}.

У этого режима есть парадоксальная подпись, по которой его можно узнать снаружи~\cite{Tsodyks1997isn}. Добавим в~\eqref{eq:ei} внешний вход $h_I$ прямо \nt{ГАМК}-группе. Установившиеся частоты станут
\begin{equation}
E^* \;=\; \frac{h-w_{EI}\,h_I}{D},\qquad
I^* \;=\; \frac{w_{IE}\,h+(1-w_{EE})\,h_I}{D},\qquad
D \;=\; 1-w_{EE}+w_{EI}w_{IE} .
\label{eq:isn}
\end{equation}
Если $w_{EE}>1$, коэффициент при $h_I$ во второй формуле отрицателен: подтолкнёшь \nt{ГАМК}-нейроны --- и их частота \emph{упадёт}. Причина в петле. Лишнее торможение приглушает \nt{ГЛУТ}-группу, та меньше возбуждает \nt{ГАМК}-нейроны, и они теряют больше, чем получили. При числах рис.~\ref{fig:ei} ($w_{EE}=1.5$, $w_{EI}w_{IE}=2$, $D=1.5$) каждая единица добавленного входа снижает $I^*$ на треть единицы. Эту подпись и нашли в коре: когда ответ нейронов зрительной коры подавляется раздражителем вокруг, тормозящий ток, который они получают, не растёт, а падает вместе с возбуждающим~\cite{Ozeki2009}. Позже ту же подпись нашли в разных областях и слоях коры~\cite{Sanzeni2020}.

Колебания при нарушении условия~(ii) в мозге тоже встречаются. Петля «\nt{ГЛУТ} возбуждает \nt{ГАМК}, \nt{ГАМК} тормозит \nt{ГЛУТ}» с задержкой в несколько миллисекунд порождает ритм с периодом порядка $12$--$50\ms$ (частота $20$--$80$\,Гц), и такие быстрые ритмы в коре хорошо известны~\cite{Cardin2009,Buzsaki2012}. Это не часы компьютера: ритм местный и возникает только тогда, когда сеть активна. Но в течение нескольких циклов он делит время на окна, в которых нейроны могут сработать.

\section{Итог}

\begin{table}[t]
\centering
\small
\begin{tabular}{>{\raggedright}p{3.3cm}>{\raggedright}p{4.4cm}>{\raggedright\arraybackslash}p{4.4cm}}
\toprule
& только \nt{ГЛУТ} & \nt{ГЛУТ} + \nt{ГАМК} \\
\midrule
НЕ & только через датчики «включения---выключения» & запрет $a\wedge\bar b$; НЕ при фоновой активности \\
провода на бит & два & один \\
направление движения & нет & запрет с задержкой \\
регулировка усиления & нет & деление: шунт вместе с фоновым шумом \\
окно совпадения & $\sim\tau$ & сужается торможением \\
выбор одного из многих & нет & взаимное торможение, $\beta>1$ \\
сброс памяти & только усталостью & сигналом через \nt{ГАМК} \\
устойчивость & только усталостью & быстрая отрицательная обратная связь \\
ритм & не нужен & возникает при медленном торможении \\
\bottomrule
\end{tabular}
\caption{Что добавляет \nt{ГАМК} к сети из \nt{ГЛУТ}-нейронов.}
\label{tab:gaba}
\end{table}

\nt{ГАМК} почти не добавляет сети новых \emph{вычислений} --- всё это можно было вычислить и с двухпроводными датчиками, --- а добавляет \emph{управление} (таблица~\ref{tab:gaba}): выбор, сброс, регулировку усиления и устойчивость. Возбуждение в мозге несёт данные, а торможение решает, каким данным пройти, когда и насколько громко. Для каждого четвёртого-пятого нейрона коры это очень разумное разделение труда.

\section{Задачи}

\begin{exercise}[\lvlA]
\label{ex:03:1}
\emph{Шунт в числах.} Постоянная времени мембраны в покое $C/g_L=20\ms$, $E_E=70\mV$.
\par\noindent(а)~По~\eqref{eq:shunt} найдите $V_\infty$ при $g_E=g_L$ и $g_E=4g_L$, без торможения и с $g_I=2g_L$. Сравните с торможением, которое при $g_E=g_L$ вычитает столько же, сколько шунт, и проверьте на втором значении $g_E$, что шунт делит, а не вычитает.
\par\noindent(б)~Найдите $\tau_{\text{эфф}}$ и окно совпадения~\eqref{eq:window} при $\theta=1.5w$ для $g_E=g_L$ без торможения и с $g_I=2g_L$.
\par\noindent(в)~Датчики детектора направления стоят на расстоянии $\Delta x=50\um$, задержка торможения $d=25\ms$. При какой скорости~\eqref{eq:vstar} детектор молчит?
\end{exercise}

\begin{exercise}[\lvlB]
\label{ex:03:2}
\emph{Вывод условий устойчивости.}
\par\noindent(а)~Найдите матрицу линейной системы~\eqref{eq:ei}, её след и определитель, и выведите из них оба условия утверждения~\ref{prop:ei}. Почему при $w_{EE}<1$ условие~(ii) не нужно?
\par\noindent(б)~На границе условия~(ii) при выполненном~(i) собственные числа чисто мнимые. Найдите частоту колебаний на границе и их период при числах рис.~\ref{fig:ei}. Найдите собственные числа при $\tau_I=2\ms$ и $\tau_I=30\ms$ и объясните обе кривые рисунка. Что в модели \texttt{figures/make\_ei.py} удерживает колебания при $\tau_I=30\ms$ от неограниченного роста?
\par\noindent(в)~Выведите~\eqref{eq:isn} и проверьте, что при числах рис.~\ref{fig:ei} $\dd I^*/\dd h_I=-1/3$. При каком условии на $w_{EE}$ подтолкнутые \nt{ГАМК}-нейроны снижают частоту?
\end{exercise}

\begin{exercise}[\lvlB]
\label{ex:03:3}
\emph{Ритм из задержки.} По задаче~\ref{ex:03:2} период колебаний модели~\eqref{eq:ei} выходит в десятки миллисекунд и больше. Заменим медленное торможение запаздывающим: \nt{ГАМК}-посредник отвечает мгновенно, но петля имеет задержку $d$:
\[
\tau_E\frac{\dd E}{\dd t}=-E(t)-G\,E(t-d)+h,\qquad G=w_{EI}w_{IE}>0 .
\]
\par\noindent(а)~Докажите, что при $d=0$ равновесие устойчиво при любом $G$, а при $G<1$ --- при любой задержке.
\par\noindent(б)~При $G>1$ найдите критическую задержку $d_c$, при которой возникают колебания, и их период на пороге. Докажите, что период лежит между $2d_c$ и $4d_c$ и стремится к $4d_c$ при $G\to\infty$.
\par\noindent(в)~Вычислите $d_c$ и период при $\tau_E=10\ms$ для $G=2$ и $G=5$ и сравните с периодами $12$--$50\ms$ быстрых ритмов коры. Проверьте результат, проинтегрировав уравнение методом Эйлера с буфером прошлых значений.
\end{exercise}

\begin{exercise}[\lvlB]
\label{ex:03:4}
\emph{Моделирование: выбор с памятью.} Проинтегрируйте~\eqref{eq:wta} методом Эйлера ($\tau=1$, шаг $10^{-3}$).
\par\noindent(а)~Воспроизведите числа раздела «Выбор»: частоты при $\beta=0.5$ и $\beta=2$ для входов $1.0$ и $0.9$ и сохранение победы после перестановки входов.
\par\noindent(б)~При $\beta=2$ и $I_1=1$ медленно (со скоростью $10^{-3}$ за $\tau$) поднимайте $I_2$ от $0$ до $3$ и опускайте обратно. При каких $I_2$ происходят переключения? Сравните с утверждением~\ref{prop:wta} и найдите ширину петли гистерезиса.
\par\noindent(в)~Стартуя из $r_1=r_2=0$ при $I_1=1+\varepsilon$, $I_2=1$, измерьте время решения (пока проигравший не упадёт ниже $1\,\%$ победителя) при $\varepsilon=10^{-1},\dots,10^{-4}$. Выведите из линейного анализа, на сколько растёт время при уменьшении $\varepsilon$ в десять раз, и сравните с моделированием. Что это значит для выбора между почти равными вариантами?
\end{exercise}

\begin{exercise}[\lvlB]
\label{ex:03:5}
\emph{Проектирование: детектор направления на полосу скоростей.} Схема рис.~\ref{fig:direction}, но детектор должен молчать при движении от $A$ к $B$ с любой скоростью, для которой время пробега $\Delta x/v$ лежит в пределах $5$--$20\ms$. Импульс \nt{ГАМК}-посредника, пришедший к $Y$ через задержку $d_k$ после $A$, запрещает $Y$ срабатывать на отрезке $[d_k,\,d_k+T_I]$, $T_I=5\ms$; возбуждение от $B$ приходит к $Y$ без задержки, и одного импульса $B$ достаточно, чтобы $Y$ сработал.
\par\noindent(а)~Сколько нужно \nt{ГАМК}-посредников и с какими задержками? Докажите, что меньше нельзя.
\par\noindent(б)~Проверьте, что при движении от $B$ к $A$ с любой скоростью $Y$ срабатывает.
\par\noindent(в)~Что делает схема при движении от $A$ к $B$ со скоростью вне полосы? Как изменится ответ (а), если $T_I$ удвоить?
\end{exercise}

\begin{exercise}[\lvlB]
\label{ex:03:6}
\emph{Запреты и фон.}
\par\noindent(а)~Докажите вторую половину утверждения~\ref{prop:gabacomplete}: без фоновой активности сеть из \nt{ГЛУТ}- и \nt{ГАМК}-нейронов вычисляет в точности логические функции с $f(0,\dots,0)=0$. Для полноты постройте схему из одного \nt{ГАМК}-посредника на каждый вход, одного нейрона-запрета на каждый набор входов с $f=1$ и одного нейрона ИЛИ, и сосчитайте нейроны для чётности трёх входов $x\oplus y\oplus z$.
\par\noindent(б)~Пусть \nt{ГАМК}-нейрону разрешено получать входы $x,y,z$ напрямую. Постройте схему чётности трёх входов из двух нейронов: одного \nt{ГАМК} и одного \nt{ГЛУТ}, указав веса и пороги.
\par\noindent(в)~Входы приходят импульсами в разные моменты. Какая гонка сигналов возникает в схеме (б) и как её устранить задержками? Можно ли её устранить, если входы разнесены во времени на неизвестную величину?
\end{exercise}

\begin{exercise}[\lvlC]
\label{ex:03:7}
\emph{Проектирование: сколько посредников нужно для выбора из ста.} Каждая из $n=100$ групп \nt{ГЛУТ}-нейронов должна тормозить каждую другую с силой $\beta$ и не тормозить себя: нужна эффективная матрица связей $W_{ij}=-\beta$ при $i\ne j$, $W_{ii}=0$. \nt{ГАМК}-посредники линейны: частота посредника $m$ равна $\sum_jA_{mj}r_j$, $A\ge0$, и он тормозит группу $i$ с весом $B_{im}\ge0$. Кроме того, группы могут возбуждать друг друга напрямую, $P\ge0$.
\par\noindent(а)~Докажите, что эффективная матрица равна $P-BA$, и что при $P_{ii}\ge0$ одного посредника хватает для любой желаемой $W$.
\par\noindent(б)~Самовозбуждение групп запрещено, $P_{ii}=0$. Докажите, что наименьшее число посредников --- наименьшее $k$ с $\binom{k}{\lfloor k/2\rfloor}\ge n$, и найдите его для $n=100$. Указание: сопоставьте группе $i$ множество посредников, которые её тормозят, и воспользуйтесь теоремой Шпернера о наибольшем семействе подмножеств, ни одно из которых не содержится в другом.
\par\noindent(в)~Прямые возбуждающие связи тоже запрещены, $P=0$. Сколько нужно посредников теперь?
\par\noindent(г)~Сравните три ответа. Какой ценой (задержкой, точностью, числом связей) оплачена экономия в (а) и (б)?
\end{exercise}

\begin{exercise}[\lvlC]
\label{ex:03:8}
\emph{Исследование: когда кора становится сетью, стабилизированной торможением.} Замените в \texttt{figures/make\_ei.py} линейную передаточную функцию \nt{ГЛУТ}-группы на квадратичную $f(u)=[u]_+^2$, оставив \nt{ГАМК}-группу линейной: $\tau_I\,\dd I/\dd t=-I+w_{IE}E+h_I$; параметры рис.~\ref{fig:ei}, $w_{EI}=2$, $w_{IE}=1$. Работайте с копией файла: сам скрипт перезаписывает \texttt{figures/ei\_traces.tex}.
\par\noindent(а)~Эффективная сила обратной связи теперь $w_{EE}f'(u^*)$ и зависит от входа $h$. Найдите аналитически $h_c$, выше которого знак $\dd I^*/\dd h_I$ становится отрицательным, и $h$, при котором нарушается условие~(ii).
\par\noindent(б)~Проверьте (а) моделированием: при $h$ от $0$ до $2$ дайте \nt{ГАМК}-группе добавку $h_I=0.01$ и измерьте изменение $I^*$. Что происходит, если резко включить большой вход $h$ из тишины?
\par\noindent(в)~Открытая часть. Модель предсказывает, что парадоксальный ответ появляется только выше некоторого уровня входа. Предложите опыт, который проверяет это предсказание, и обсудите, что следует из того, что подпись~\eqref{eq:isn} находят в коре и при слабой, фоновой активности. Какие свойства модели надо изменить, чтобы она это воспроизвела?
\end{exercise}

\chapter{Азбука Морзе: что мы знаем о языке, на котором говорят \nt{GLUT}- и \nt{GABA}-нейроны}
\chaptermark{Азбука Морзе}
\label{sec:code}

\section{Алфавит из одного знака}

В азбуке Морзе два знака, точка и тире, и паузы трёх длин между ними. У нейрона, как мы видели в главе~\ref{sec:neurontypes}, алфавит ещё беднее: знак один. Импульс длится около миллисекунды, и все импульсы одного нейрона одинаковы по высоте и форме. Поэтому всё сообщение --- в паузах, то есть в последовательности моментов $t_1,t_2,t_3,\dots$ Это азбука без тире, в которой смысл несёт только ритм точек (рис.~\ref{fig:morse}).

\begin{figure}[t]
\centering
\begin{tikzpicture}[>=Stealth,x=0.08cm,y=1cm]
% Morse letter: dot, dash, dot
\node[left,font=\small] at (-6,2.55) {Морзе:};
\fill[black] (0,2.45) rectangle (6,2.65);
\fill[black] (14,2.45) rectangle (32,2.65);
\fill[black] (40,2.45) rectangle (46,2.65);
\node[right,font=\small,gray!40!black] at (52,2.55) {точка, тире, точка --- смысл в длительностях и паузах};
% stimulus onset
\node[left,font=\small] at (-6,0.4) {нейрон:};
\draw[->,thick,green!45!black] (0,-0.55) -- (0,-0.05);
\node[below,font=\scriptsize,green!45!black] at (0,-0.55) {раздражитель};
\draw[gray!70!black] (0,0) -- (152,0) node[right,black] {\small $t$};
% spikes: single ones blue, the burst red
\foreach \s in {14,26,41,88,112,131}{\draw[very thick,blue!60!black] (\s,0) -- (\s,0.8);}
\foreach \s in {60,63,66,69}{\draw[very thick,red!70!black] (\s,0) -- (\s,0.8);}
% latency
\draw[<->,thick] (0,1.1) -- node[above,font=\scriptsize] {$t_1$: время первого импульса} (14,1.1);
% burst
\draw[decorate,decoration={brace,amplitude=4pt},red!70!black] (58,0.95) -- (71,0.95) node[midway,above=4pt,font=\scriptsize] {пачка --- «тире»};
% counting windows
\foreach \a/\b/\n in {0/50/3,50/100/5,100/150/2}{
  \draw[decorate,decoration={brace,amplitude=4pt,mirror},gray!50!black] (\a+1,-0.2) -- (\b-1,-0.2) node[midway,below=4pt,font=\small] {$n=\n$};}
\node[font=\scriptsize,gray!40!black] at (75,-1.05) {частотный код: число импульсов $n$ в каждом окне};
\foreach \x in {0,50,100,150}{\draw[gray!60!black] (\x,-0.06) -- (\x,0.06);}
\end{tikzpicture}
\caption{Азбука Морзе и азбука нейрона. Одну и ту же последовательность импульсов можно прочитать по-разному: сосчитать импульсы в окнах по $50\ms$ (раздел~\ref{sec:rate}), измерить задержку $t_1$~\eqref{eq:latency} или заметить пачку --- «тире»~\eqref{eq:burst}.}
\label{fig:morse}
\end{figure}

Какие паузы вообще возможны? Снизу их ограничивает невосприимчивость: после импульса нейрон около миллисекунды не может сработать снова, так что больше нескольких сотен импульсов в секунду не бывает. Сверху не ограничивает ничто: нейрон может молчать минутами. У \nt{GLUT}-нейронов коры частоты лежат в пределах от долей импульса до нескольких десятков импульсов в секунду, и большинство нейронов большую часть времени работает у нижней границы.

В главах~\ref{sec:glutcomputer} и~\ref{sec:gaba} мы попутно принимали то один, то другой способ записи чисел импульсами. Теперь спросим прямо: на каком языке \nt{GLUT}- и \nt{GABA}-нейроны говорят друг с другом? Полного ответа нет, этот язык не расшифрован. Но кое-что о нём известно твёрдо, и почти всё известное можно получить, рассуждая как инженер связи: ёмкость канала, шум, приёмник, цена передачи.

\section{Сколько можно передать}

Начнём с верхней границы. Пусть получатель различает моменты импульсов с точностью $\Delta t$: сдвиги меньше $\Delta t$ для него неразличимы. Разрежем время на клетки длины $\Delta t$, меньшие времени невосприимчивости, так что в каждой клетке либо один импульс, либо ни одного (рис.~\ref{fig:cells}). При средней частоте $r$ импульс попадает в клетку с вероятностью $p=r\Delta t$. Больше всего информации поток несёт, когда клетки независимы, и тогда каждая клетка даёт энтропию $-p\log_2p-(1-p)\log_2(1-p)\approx p\log_2(e/p)$ бит при $p\ll1$. Разделив на $\Delta t$, получим предельную скорость передачи и её долю на один импульс~\cite{MacKay1952}:
\begin{empheq}[box=\keybox]{equation}
R_{\max} \;\approx\; r\,\log_2\frac{e}{r\,\Delta t}\ \ \text{бит/с},
\qquad
\frac{R_{\max}}{r} \;\approx\; \log_2\frac{e}{r\,\Delta t}\ \ \text{бит на импульс}.
\label{eq:spikeentropy}
\end{empheq}
При $r=10$ импульсов в секунду и $\Delta t=1\ms$ это около восьми бит на импульс и восьмидесяти бит в секунду.

\begin{figure}[t]
\centering
\begin{tikzpicture}[>=Stealth,x=0.5cm,y=1cm]
% time cut into cells of width Delta t; a cell holds one impulse or none
\foreach \k in {0,...,23}{\draw[gray!55] (\k,0) rectangle (\k+1,0.7);}
\foreach \k in {2,7,8,15,21}{
  \fill[red!10] (\k,0) rectangle (\k+1,0.7);
  \draw[very thick,red!70!black] (\k+0.5,0.05) -- (\k+0.5,0.65);}
\foreach \k in {0,...,23}{
  \pgfmathtruncatemacro{\bit}{(\k==2)||(\k==7)||(\k==8)||(\k==15)||(\k==21)}
  \ifnum\bit=1 \def\bitcol{red!70!black}\else \def\bitcol{gray!60!black}\fi
  \node[font=\scriptsize,text=\bitcol] at (\k+0.5,-0.25) {\bit};}
\draw[<->,gray!60!black] (0,0.95) -- (1,0.95) node[midway,above,font=\scriptsize,black] {$\Delta t$};
\node[font=\small,anchor=east] at (-0.3,0.35) {импульсы};
\node[font=\small,anchor=east] at (-0.3,-0.25) {биты};
\draw[decorate,decoration={brace,amplitude=5pt,mirror},gray!60!black] (0,-0.55) -- (24,-0.55)
  node[midway,below=6pt,font=\scriptsize,black,align=center] {получатель, который только считает: «$n=5$»;\\ получатель, который видит клетки: какие именно $5$ из $24$ --- это $\log_2\binom{24}{5}\approx15$ бит};
\end{tikzpicture}
\caption{Откуда берётся ёмкость~\eqref{eq:spikeentropy}. Время нарезано на клетки длины $\Delta t$, в каждой клетке импульс либо есть ($1$), либо нет ($0$): поток импульсов --- это двоичная строка. Импульс попадает в клетку с вероятностью $p=r\Delta t$. Чем мельче клетки, тем длиннее строка и тем больше бит; счётчик импульсов теряет почти всё, что записано в том, \emph{где} стоят единицы.}
\label{fig:cells}
\end{figure}

Биты на импульс зависят от точности времени лишь логарифмически: при точности $10\ms$ останется около пяти бит на импульс. Но если получатель только считает импульсы за секунду, то при $r=10$ он получит меньше двух бит за всю секунду (раздел~\ref{sec:rate}).

\begin{keyblock}
\begin{proposition}[Ёмкость импульсного канала]
\label{prop:capacity}
Поток импульсов со средней частотой $r$, прочитанный с точностью времени $\Delta t$, несёт не больше~\eqref{eq:spikeentropy}. Почти вся эта ёмкость сидит во \emph{времени} импульсов: получатель, который только считает импульсы, извлекает из того же потока в десятки раз меньше, чем тот, кто различает их моменты с точностью до миллисекунды.
\end{proposition}
\end{keyblock}

Это верхняя граница, а не измерение. Измерения на чувствительных нейронах, для которых известен раздражитель, дают от одного до нескольких бит на импульс; в лучших случаях это до половины границы, вычисленной для их точности~\cite{Rieke1997}. Значит, по крайней мере на входе в нервную систему время импульсов используется, а не пропадает.

\section{Шумный канал}

Ёмкость~\eqref{eq:spikeentropy} посчитана без шума. О том, где он возникает, твёрдо известны три факта.

\emph{Сам генератор импульсов точен.} Если в нейрон коры, вынутый из мозга вместе с кусочком ткани, много раз подать один и тот же ток, быстро меняющийся во времени, импульсы повторяются с точностью около миллисекунды~\cite{Mainen1995}. Если ток постоянный, моменты импульсов от раза к разу расползаются: точное время порождают быстрые перепады входа, а не сам нейрон.

\emph{Синапс ненадёжен.} Импульс, добежавший до \nt{GLUT}-синапса, выбрасывает порцию медиатора лишь с некоторой вероятностью $p$. На синапсах гиппокампа больше половины импульсов до получателя просто не доходят, и причина --- именно случайность выброса, а не сбои проведения~\cite{AllenStevens1994}. Для связиста это канал со стираниями; несколько синапсов между двумя нейронами отчасти спасают дело.

\emph{В живой коре поток импульсов выглядит случайным.} Интервалы между импульсами разбросаны примерно так же, как у случайного пуассоновского потока, а при повторении одного и того же раздражителя число импульсов за окно меняется так, что его дисперсия близка к среднему~\cite{Softky1993}.

Как точный элемент порождает случайный поток? Нейрон получает тысячи входов, каждый ненадёжен, а возбуждение и торможение уравновешены, как в главе~\ref{sec:gaba}. Напряжение мембраны дрожит у порога, и моменты импульсов задают эти дрожания. Шум ли это или сообщение, которое мы не умеем прочитать, --- во многом открытый вопрос. Часть «шума» уже оказалась сообщением: в зрительной коре значительная доля разброса следует за движениями самого животного~\cite{Stringer2019}. Трудность здесь принципиальная: хорошо сжатое сообщение для того, кто не знает кода, неотличимо от случайного потока.

\section{Частота: медленно, но надёжно}
\label{sec:rate}

Самый простой код --- \emph{частотный}: число записано в том, сколько импульсов нейрон выдаёт за окно $T$. Ни стирания, ни дрожание моментов такому коду не страшны. Но у него есть цена. Если поток пуассоновский, число импульсов $n$ за окно имеет среднее $rT$ и разброс $\sqrt{rT}$:
\begin{empheq}[box=\keybox]{equation}
\frac{\sigma_n}{\bar n} \;=\; \frac{1}{\sqrt{r\,T}} .
\label{eq:rateerror}
\end{empheq}
Чтобы узнать частоту с точностью в десять процентов, нужно сто импульсов, при двадцати импульсах в секунду --- пять секунд. Соседние уровни различимы, если они отстоят на пару разбросов, поэтому до частоты $r$ за время $T$ различается около $\sqrt{rT}$ уровней: при $r=10$ и $T=1\,\text{с}$ это три-четыре уровня, меньше двух бит.
Мозг так медленно не работает. Человек узнаёт животное на картинке, показанной на мгновение, и ответ в мозге появляется примерно через $150\ms$~\cite{Thorpe1996}. По грубой оценке, за это время сигнал проходит около десятка последовательных ступеней обработки, по $10$--$20\ms$ на ступень. При обычных частотах коры каждый нейрон успевает выдать на ступени ноль или один импульс, и его частота в таком окне не определена. Выходов два: взять много нейронов или смотреть на время.

\emph{Много нейронов.} Пусть $N$ нейронов несут одно и то же число, и получатель складывает их импульсы. Если их шумы независимы, в~\eqref{eq:rateerror} вместо $rT$ встаёт $NrT$: сто нейронов при $r=20$ за $15\ms$ дают тридцать импульсов и точность около двадцати процентов. Пространство заменяет время. Однако шумы соседних нейронов коры не независимы: коэффициент корреляции чисел импульсов у пары соседей обычно около $0.1$~\cite{Zohary1994}. Если он равен $c$, дисперсия среднего по $N$ нейронам равна
\begin{empheq}[box=\keybox]{equation}
\sigma^2_N \;=\; \sigma^2_1\,\frac{1+(N-1)\,c}{N}
\;\xrightarrow[N\to\infty]{}\; c\,\sigma^2_1 .
\label{eq:correlated}
\end{empheq}

\begin{keyblock}
\begin{proposition}[Предел усреднения]
\label{prop:pooling}
При коррелированном шуме усреднение по группе уменьшает ошибку не больше чем в $1/\sqrt{c}$ раз. При $c=0.1$ это втрое, и такой выигрыш достигается уже на нескольких десятках нейронов; добавлять больше бесполезно.
\end{proposition}
\end{keyblock}

Не всякая корреляция одинаково вредна. Ограничивает точность только та часть общего шума, которая \emph{похожа на сам сигнал}, то есть сдвигает всю группу так же, как её сдвинуло бы изменение измеряемой величины~\cite{MorenoBote2014}. Сколько такого шума в мозге на самом деле, ещё выясняют.

Есть и другой способ использовать много нейронов: записать число в том, \emph{какой} нейрон сработал, как в определителе направления на звук (рис.~\ref{fig:itd}). Это код \emph{местом}, самый надёжный из известных: у чувствительных нейронов номер провода говорит, какая точка кожи тронута или какой высоты звук, и это установлено многократно. Он дорог по числу нейронов, зато быстр: на сообщение уходит одна задержка.

\section{Время: быстро, но от чего отсчитывать}

Второй выход --- писать число во время импульса. Возьмём нейрон~\eqref{eq:glutneuron} и подадим на него с момента $t=0$ постоянный вход, который поднял бы напряжение до уровня $U$. Напряжение растёт как $V=U\bigl(1-e^{-t/\tau}\bigr)$ и достигает порога $\theta$ в момент
\begin{empheq}[box=\keybox]{equation}
t_1 \;=\; \tau\,\ln\frac{U}{U-\theta}, \qquad U>\theta .
\label{eq:latency}
\end{empheq}
Сильный вход --- ранний импульс, слабый --- поздний, подпороговый --- никакого. При $\tau=10\ms$ вход $U=4\theta$ даёт первый импульс через $2.9\ms$, $U=2\theta$ --- через $6.9\ms$, $U=1.2\theta$ --- через $17.9\ms$. Один-единственный импульс несёт число.

Но от какого момента отсчитывать $t_1$? Часов нет, и получатель не знает, когда начался раздражитель. Есть три ответа.

\emph{Относительная задержка.} Два нейрона видят один и тот же раздражитель, и разность их задержек $t_1^A-t_1^B$ зависит только от их входов, а не от момента начала. Сравнивают моменты детекторы совпадений с задержками (рис.~\ref{fig:itd}). Если $N$ нейронов выдали по одному импульсу, сам \emph{порядок} их прихода несёт до $\log_2N!$ бит: для десяти нейронов около двадцати двух бит, тогда как ответ «сработал или нет» дал бы не больше десяти. В сетчатке показано, что по относительным задержкам первых импульсов её выходных нейронов картинку можно восстановить быстро и надёжно~\cite{Gollisch2008}.

\emph{Залп как начало слова.} Резкое изменение раздражителя заставляет сработать почти одновременно множество нейронов. Такой залп сам служит отметкой начала, как длинная пауза между словами в азбуке Морзе, и получатель отсчитывает задержки от него.

\emph{Местный ритм.} В главе~\ref{sec:gaba} петля \nt{GLUT}--\nt{GABA} порождала местный ритм с периодом $12$--$50\ms$. Это не часы, но внутри его цикла нейрон с сильным входом успевает сработать раньше, чем нейрон со слабым, и~\eqref{eq:latency} превращается в код \emph{фазой}: число записано в том, в какой части цикла пришёл импульс. Такой код наблюдали: нейроны, отвечающие за определённое место в пространстве, срабатывают всё раньше в цикле медленного местного ритма по мере того, как животное проходит через «своё» место~\cite{OKeefe1993}. Насколько широко мозг пользуется фазой, пока неизвестно.

\section{Код выбирает получатель}

В последовательности импульсов есть и частота, и времена, и порядок. Что из этого будет прочитано, решает получатель, и решает одним параметром --- своей постоянной времени $\tau$.

Усредним уравнение~\eqref{eq:glutneuron} по времени. Если на нейрон приходят независимые потоки с частотами $r_j$, то среднее напряжение и его относительное дрожание равны
\begin{empheq}[box=\keybox]{equation}
\bar V \;=\; \tau\sum_j w_j\,r_j ,
\qquad
\frac{\sigma_V}{\bar V} \;\approx\; \frac{1}{\sqrt{2\,r_{\text{вх}}\,\tau}} ,
\label{eq:reader}
\end{empheq}
где второе равенство написано для одинаковых весов, а $r_{\text{вх}}$ --- суммарная частота всех входов. Когда $\tau$ велика, напряжение почти не дрожит и повторяет взвешенную сумму частот: получатель --- \emph{интегратор}, он читает частоты и забывает времена. Когда $\tau$ мала, напряжение успевает упасть между импульсами, и порог достигается только тогда, когда несколько импульсов пришли в пределах окна~\eqref{eq:window}: получатель --- \emph{детектор совпадений}, он читает времена (рис.~\ref{fig:reader}). Тысяча входов по пять импульсов в секунду при $\tau=10\ms$ дают дрожание около десяти процентов, почти чистое интегрирование. Если торможение, как в главе~\ref{sec:gaba}, укоротило $\tau$ до $2\ms$, дрожание вырастает до двадцати с лишним процентов, и нейрон начинает замечать совпадения.

\begin{figure}[t]
\centering
\begin{tikzpicture}[>=Stealth,x=0.115cm,y=1cm]
% common input: the same spike train for both receivers
\foreach \s in {6,21,22,38,52,55.5,59,62.5,66,69.5,73,76.5,80,83.5,87,90.5,94,97.5}{\draw[thick,gray!40!black] (\s,5.25) -- (\s,5.65);}
\draw[gray!70!black] (0,5.25) -- (105,5.25);
\node[left,font=\small] at (-1,5.45) {вход};
\node[above,font=\scriptsize,gray!40!black] at (13.5,5.65) {редко};
\node[above,font=\scriptsize,gray!40!black] at (21.5,6.0) {пара};
\node[above,font=\scriptsize,gray!40!black] at (75,5.65) {часто};
% axes and thresholds
\foreach \y/\th in {2.7/2.5,0/1.5}{
  \draw[gray!70!black] (0,\y) -- (106,\y) node[right,black] {\small $t$};
  \draw[densely dashed,gray!60!black] (0,\y+0.6*\th) -- (104,\y+0.6*\th) node[right,black,font=\small] {$\theta$};
}
\node[anchor=south west,font=\small] at (0,4.3) {$\tau=10\ms$: интегратор};
\node[anchor=south east,font=\small] at (104,1.0) {$\tau=2\ms$: детектор совпадений};
% integrator: tau=10 ms, theta=2.5w
\begin{scope}[yshift=2.7cm]
\foreach \a/\b/\v in {0/6/0.000,6/21/1.000,21/22/1.223,22/38/2.107,38/52/1.425,52/55.5/1.351,55.5/59/1.952,59/62.5/2.376,62.5/66/0.000,66/69.5/1.000,69.5/73/1.705,73/76.5/2.201,76.5/80/0.000,80/83.5/1.000,83.5/87/1.705,87/90.5/2.201,90.5/94/0.000,94/97.5/1.000,97.5/104/1.705}{\draw[very thick,blue!60!black] plot[domain=\a:\b,samples=14] (\x,{0.6*\v*exp(-(\x-\a)/10)});}
\foreach \s/\p/\q in {6/0.000/1.000,21/0.223/1.223,22/1.107/2.107,38/0.425/1.425,52/0.351/1.351,55.5/0.952/1.952,59/1.376/2.376,62.5/1.674/2.500,62.5/2.500/0.000,66/0.000/1.000,69.5/0.705/1.705,73/1.201/2.201,76.5/1.551/2.500,76.5/2.500/0.000,80/0.000/1.000,83.5/0.705/1.705,87/1.201/2.201,90.5/1.551/2.500,90.5/2.500/0.000,94/0.000/1.000,97.5/0.705/1.705}{\draw[very thick,blue!60!black] (\s,0.6*\p) -- (\s,0.6*\q);}
\foreach \s in {62.5,76.5,90.5}{\draw[->,thick,red!70!black] (\s,1.58) -- (\s,1.95);}
\end{scope}
% coincidence: tau=2 ms, theta=1.5w
\begin{scope}[yshift=0cm]
\foreach \a/\b/\v in {0/6/0.000,6/21/1.000,21/22/1.001,22/38/0.000,38/52/1.000,52/55.5/1.001,55.5/59/1.174,59/62.5/1.204,62.5/66/1.209,66/69.5/1.210,69.5/73/1.210,73/76.5/1.210,76.5/80/1.210,80/83.5/1.210,83.5/87/1.210,87/90.5/1.210,90.5/94/1.210,94/97.5/1.210,97.5/104/1.210}{\draw[very thick,blue!60!black] plot[domain=\a:\b,samples=14] (\x,{0.6*\v*exp(-(\x-\a)/2)});}
\foreach \s/\p/\q in {6/0.000/1.000,21/0.001/1.001,22/0.607/1.500,22/1.500/0.000,38/0.000/1.000,52/0.001/1.001,55.5/0.174/1.174,59/0.204/1.204,62.5/0.209/1.209,66/0.210/1.210,69.5/0.210/1.210,73/0.210/1.210,76.5/0.210/1.210,80/0.210/1.210,83.5/0.210/1.210,87/0.210/1.210,90.5/0.210/1.210,94/0.210/1.210,97.5/0.210/1.210}{\draw[very thick,blue!60!black] (\s,0.6*\p) -- (\s,0.6*\q);}
\foreach \s in {22}{\draw[->,thick,red!70!black] (\s,0.98) -- (\s,1.35);}
\end{scope}
\foreach \x in {0,50,100}{\draw[gray!60!black] (\x,-0.06) -- (\x,0.06) node[below,font=\scriptsize] {\x\,мс};}
\end{tikzpicture}
\caption{Один и тот же вход --- два разных получателя. Каждый импульс поднимает напряжение на $w$, после чего оно стекает, как в модели нейрона~\eqref{eq:glutneuron}; красные стрелки --- импульсы получателя. Медленный получатель ($\tau=10\ms$, $\theta=2.5w$) срабатывает там, где импульсов \emph{много}, и не замечает пары. Быстрый ($\tau=2\ms$, $\theta=1.5w$) срабатывает только на пару в пределах окна~\eqref{eq:window}, а частый, но не совпадающий поток пропускает.}
\label{fig:reader}
\end{figure}

Измерения показывают, что в активной коре проводимость мембраны возрастает в несколько раз по сравнению с покоем~\cite{Destexhe2003}. Значит, $\tau$ там короче, и нейроны коры в рабочем режиме ближе к детекторам совпадений, чем к интеграторам.

\begin{keyblock}
\begin{proposition}[Код задаёт получатель]
\label{prop:reader}
Одна и та же последовательность импульсов для медленного получателя несёт частоту, для быстрого --- времена, для объёма, в который выброшен медиатор, --- уровень концентрации, сглаженный за секунды (глава~\ref{sec:serocomputer}). Какой код используется, определяется не отправителем, а постоянной времени получателя, и её подстраивает торможение.
\end{proposition}
\end{keyblock}
\section{Точки и тире: синапс как фильтр}

До сих пор синапс был у нас проводом с весом. На самом деле его ответ зависит от предыдущих импульсов, и это даёт языку нейронов второй знак.

\emph{Истощение: синапс передаёт изменения.} Каждый выброс тратит долю $U$ запаса медиатора, готового к выбросу~\cite{Tsodyks1997}, а запас восстанавливается с постоянной времени $\tau_{\text{вос}}$ порядка сотен миллисекунд. При частоте $r$ амплитуда ответа на импульс устанавливается приблизительно на уровне $A(r)=A_0/(1+U r\tau_{\text{вос}})$, где $A_0$ --- ответ отдохнувшего синапса. Ток, который получатель получает в единицу времени, равен
\begin{empheq}[box=\keybox]{equation}
r\,A(r) \;=\; \frac{A_0\,r}{1+U r\,\tau_{\text{вос}}}
\;\xrightarrow[r\to\infty]{}\; \frac{A_0}{U\tau_{\text{вос}}} .
\label{eq:depression}
\end{empheq}
При $U=0.5$ и $\tau_{\text{вос}}=0.8\,\text{с}$ насыщение начинается уже с $1/(U\tau_{\text{вос}})=2.5$ импульса в секунду. Если частота выросла с пяти до двадцати, установившийся ток вырастет всего на треть. Зато первые импульсы на новой частоте приходят, пока запас ещё прежний, и ток на мгновение подскакивает вчетверо, а затем за $\tau_{\text{вос}}/(1+Ur\tau_{\text{вос}})\approx90\ms$ спадает (рис.~\ref{fig:depression}).

Такой синапс передаёт не частоту, а её \emph{изменение}, по существу --- относительное изменение $\Delta r/r$~\cite{Abbott1997depression}. Для электронщика это фильтр верхних частот, поставленный прямо в провод.

\begin{figure}[t]
\centering
\begin{tikzpicture}[>=Stealth,x=12.5cm,y=1cm]
% input rate
\draw[->,gray!70!black] (0,2.6) -- (0.9,2.6) node[right,black] {\small $t$};
\draw[->,gray!70!black] (0,2.6) -- (0,4.3) node[above,black] {\small $r$};
\draw[very thick,blue!60!black] (0,2.9) -- (0.2,2.9) -- (0.2,3.8) -- (0.8,3.8);
\node[left,font=\scriptsize] at (0,2.9) {$5$};
\node[left,font=\scriptsize] at (0,3.8) {$20$};
\node[right,font=\small,blue!60!black] at (0.4,4.1) {частота на входе синапса};
% output current
\draw[->,gray!70!black] (0,0) -- (0.9,0) node[right,black] {\small $t$};
\draw[->,gray!70!black] (0,0) -- (0,2.45) node[left,black] {\small $rA$};
\draw[densely dashed,gray!60!black] (0,0.667) -- (0.8,0.667);
\draw[very thick,red!70!black] (0,0.5) -- (0.2,0.5) -- (0.2,2.0)
   plot[domain=0.2:0.8,samples=60] (\x,{0.667+(2.0-0.667)*exp(-(\x-0.2)/0.089)});
\node[left,font=\scriptsize] at (0,0.5) {$1.7$};
\node[left,font=\scriptsize] at (0,2.0) {$6.7$};
\node[right,font=\scriptsize] at (0.8,0.667) {$2.2$};
\node[right,font=\small,red!70!black] at (0.28,1.6) {ток получателю: подскок и спад за $\approx90\ms$};
\draw[gray!60!black] (0.2,-0.08) -- (0.2,0.08); \node[below,font=\scriptsize] at (0.2,0) {$0.2\,\text{с}$};
\draw[gray!60!black] (0.8,-0.08) -- (0.8,0.08); \node[below,font=\scriptsize] at (0.8,0) {$0.8\,\text{с}$};
\end{tikzpicture}
\caption{Истощающийся синапс по формуле~\eqref{eq:depression} при $U=0.5$, $\tau_{\text{вос}}=0.8\,\text{с}$; ток в единицах $A_0$ за секунду, штриховая линия --- установившийся ток: он вырос лишь с $1.7$ до $2.2$, а в момент скачка --- до $6.7$. Синапс сообщает получателю о том, что частота изменилась, а не о том, какая она.}
\label{fig:depression}
\end{figure}

\emph{Облегчение: пачка как тире.} У других синапсов вероятность выброса мала, но после каждого импульса на десятки миллисекунд растёт. Одиночный импульс через такой синапс чаще всего теряется. А \emph{пачка} --- несколько импульсов подряд с интервалами в несколько миллисекунд --- проходит почти наверняка. Даже без облегчения вероятность, что пропадут все $k$ импульсов пачки, равна
\begin{empheq}[box=\keybox]{equation}
P_{\text{потери}} \;=\; (1-p)^k ,
\label{eq:burst}
\end{empheq}
при $p=0.2$ это $0.8$ для одиночного импульса и $0.41$ для пачки из четырёх; облегчение уменьшает её ещё сильнее. Так у нейрона появляется второй знак: одиночный импульс --- точка, пачка --- тире. Истощающийся синапс лучше всего передаёт первый импульс после паузы; облегчающийся пропускает пачки и глушит одиночные точки. То, что пачки передаются надёжнее, измерено; что мозг действительно пользуется ими как отдельным знаком, --- правдоподобная гипотеза~\cite{Lisman1997}.

\section{Как говорят \nt{GABA}-нейроны}

Самые многочисленные из них в коре --- быстрые~\cite{Tremblay2016}. Их импульсы короче, чем у \nt{GLUT}-нейронов, они могут выдавать их ровно и подолгу с частотой в сотни в секунду, почти не уставая. Их синапсы надёжнее: связь с каждым получателем состоит из многих мест выброса, и импульс почти никогда не теряется. Их выходы садятся рядом с тем местом получателя, где рождается его импульс, так что они управляют не тем, как получатель складывает входы, а тем, сработает ли он и когда.

Сами они получают возбуждение от многих соседних \nt{GLUT}-нейронов и сообщают о суммарной активности вокруг --- ровно то, что нужно для деления из главы~\ref{sec:gaba}. Наконец, быстрые \nt{GABA}-нейроны связаны друг с другом и легко срабатывают вместе; именно они задают местный ритм, о котором шла речь выше~\cite{Cardin2009}.

На языке азбуки Морзе \nt{GABA}-нейроны отвечают не за буквы, а за \emph{паузы}. Они режут непрерывный поток \nt{GLUT}-импульсов на окна, внутри которых получатель может сработать, сужают окна совпадения и приглушают весь поток, когда он становится слишком громким.

\begin{keyblock}
\begin{proposition}[Разделение ролей]
\label{prop:punctuation}
\nt{GLUT}-нейроны несут содержание сообщения: \emph{что} и \emph{сколько}. Быстрые \nt{GABA}-нейроны несут его разметку: \emph{когда} можно говорить и \emph{насколько громко}; ещё один класс, тормозя тормозящих, решает, \emph{для кого} ворота открыты. Отдельные части этого разделения измерены; как общее правило это рабочая гипотеза.
\end{proposition}
\end{keyblock}

Есть и другие, более медленные \nt{GABA}-нейроны, выходы которых садятся на отдельные входы получателя. Они работают как запрет~\eqref{eq:veto} из главы~\ref{sec:gaba}: закрывают один поток \nt{GLUT}-импульсов, не трогая остальных.

Деление на быстрые и медленные \nt{GABA}-нейроны --- старая картина, и она неполна. Сейчас в коре различают по меньшей мере три крупных класса~\cite{Tremblay2016}, и третий устроен неожиданно: он тормозит не \nt{GLUT}-нейроны, а другие \nt{GABA}-нейроны, прежде всего те, что закрывают входы~\cite{Pfeffer2013}. Торможение торможения --- двойное отрицание. Такой нейрон \emph{открывает} ворота, не посылая получателю ни одного возбуждающего импульса. Включают его сигналы из других областей коры и широковещательные каналы, так что он работает как разрешающий вход целого участка сети: пока он активен, запреты сняты, и поток проходит~\cite{Pi2013}. В приложении~\ref{sec:othermediators} этот приём назван растормаживанием.
\section{Экономный язык}

Последнее, что известно о языке нейронов твёрдо, --- он дорог. Импульс вместе с работой синапсов, которую он вызывает, обходится примерно в $10^9$ молекул АТФ (оценка для коры грызунов~\cite{Attwell2001}), а одна молекула даёт около $10^{-19}$ джоуля. Значит, импульс стоит порядка $E_{\text{имп}}\sim10^{-10}$ джоуля. Весь мозг потребляет около $20$ ватт, и если на передачу сигналов уходит половина, $P\sim10$ ватт, то средняя частота одного нейрона
\begin{empheq}[box=\keybox]{equation}
\bar r \;\approx\; \frac{P}{N\,E_{\text{имп}}}
\;\sim\; \frac{10\ \text{Вт}}{8.6\cdot10^{10}\cdot10^{-10}\ \text{Дж}}
\;\approx\; 1\ \text{импульс в секунду}.
\label{eq:energy}
\end{empheq}
Более подробные оценки для коры дают ещё меньше~\cite{Lennie2003}. Код мозга обязан быть \emph{разреженным}: быстро работать может лишь малая доля нейронов.

Из~\eqref{eq:spikeentropy} видно, что это не так уж плохо. При точности $1\ms$ импульс нейрона, работающего с частотой $100$ в секунду, несёт не больше пяти бит, а импульс нейрона, работающего раз в секунду, --- до одиннадцати. Если платить приходится за импульсы, выгодно говорить редко. Кора и правда меняет точность на энергию: при нехватке пищи нейроны сохраняют частоту импульсов, но тратят меньше на синаптические токи, и их ответы становятся менее точными~\cite{Padamsey2022}.

В азбуке Морзе частые буквы короткие: самая частая буква английского текста E --- одна точка. Код нейронов, насколько известно, следует тому же правилу в другой форме: то, что ожидаемо, стоит мало импульсов~\cite{Barlow1961}. Нейрон при постоянном раздражителе постепенно снижает частоту, датчики главы~\ref{sec:glutcomputer} отвечают на перемены, а не на уровни, а \nt{DOPA}-нейроны главы~\ref{sec:neurontypes} сообщают только об ошибке предсказания награды.

\begin{keyblock}
\begin{proposition}[Экономия]
\label{prop:economy}
Энергия ограничивает среднюю частоту нейрона величиной порядка одного импульса в секунду. Поэтому язык \nt{GLUT}-нейронов разрежен и тратит импульсы на новости: на то, что изменилось или не было предсказано. Энергетический предел измерен; что код мозга подстроен под наибольшее число бит на джоуль, --- хорошо обоснованная гипотеза.
\end{proposition}
\end{keyblock}

\section{Итог}

\begin{table}[t]
\centering
\footnotesize
\setlength{\tabcolsep}{4pt}
\renewcommand{\arraystretch}{1.15}
\begin{tabular}{>{\raggedright}p{2.6cm}>{\raggedright}p{2.3cm}>{\raggedright}p{3.0cm}>{\raggedright}p{2.0cm}>{\raggedright\arraybackslash}p{3.6cm}}
\toprule
Способ записи & Что несёт & Кто читает & Время на сообщение & Что известно \\
\midrule
номер сработавшего нейрона (код местом) & какое значение & любой получатель; детекторы совпадений & одна задержка & установлено у чувствительных нейронов и в определении направления на звук \\
частота одного нейрона & величину & интегратор с большой $\tau$; объём (гл.~\ref{sec:serocomputer}) & сотни мс -- секунды & установлено; для ступени обработки в $10$--$20\ms$ слишком медленно \\
частота группы & величину & нейрон с тысячами входов & $10$--$50\ms$ & установлено; выигрыш ограничен корреляциями~\eqref{eq:correlated} \\
время первого импульса, порядок & величину, порядок & детекторы совпадений с задержками & одна задержка & установлено на входе нервной системы; в коре обсуждается \\
фаза в местном ритме & величину, положение & нейроны с общим ритмом & цикл $12$--$50\ms$ и медленнее & наблюдается в отдельных областях; общая роль --- гипотеза \\
пачка («тире») & «важно»: надёжная доставка & облегчающийся синапс & $\sim10\ms$ & надёжность измерена; отдельный знак --- гипотеза \\
изменение частоты & новость & истощающийся синапс & $\sim0.1\,\text{с}$ & установлено \\
импульсы быстрых \nt{GABA}-нейронов & когда и насколько громко & все соседние \nt{GLUT}-нейроны & $1$--$30\ms$ & ритм и деление измерены; «пунктуация» как правило --- гипотеза \\
\bottomrule
\end{tabular}
\caption{Способы, которыми \nt{GLUT}- и \nt{GABA}-нейроны записывают сообщения импульсами. Правый столбец отделяет измеренное от предполагаемого.}
\label{tab:code}
\end{table}

Таблица~\ref{tab:code} собирает то, что мы знаем; видно из неё и то, чего мы не знаем. Мы не знаем, насколько кора пользуется точным временем импульсов, а не только частотами групп. Мы не знаем, есть ли у нейронов «слова» --- устойчивые сочетания импульсов многих нейронов, которые читаются как целое. И мы не знаем, одинаков ли язык в разных частях мозга. Азбуку Морзе можно выучить за вечер, потому что её таблица известна. Таблицу языка нейронов ещё предстоит составить, и составлять её будут, скорее всего, инженеры связи.

\section{Задачи}

\begin{exercise}[\lvlA]
\label{ex:04:1}
\emph{Ёмкость в числах.} $\Delta t=1\ms$. (а)~Сколько бит на джоуль даёт нейрон с частотой $1$ импульс в секунду при цене импульса из~\eqref{eq:energy}? (б)~При $r=10$ во сколько раз меньше границы~\eqref{eq:spikeentropy} извлекает получатель, который только считает импульсы за секунду?
\end{exercise}

\begin{exercise}[\lvlB]
\label{ex:04:2}
\emph{Оптимальная частота.} Нейрон тратит мощность $P_0+rE_{\text{имп}}$, где $P_0$ --- вторая половина из $20$~Вт, поделённая на все нейроны, $E_{\text{имп}}=10^{-10}$~Дж. При какой частоте $r$ число бит на джоуль $R_{\max}(r)/(P_0+rE_{\text{имп}})$ по~\eqref{eq:spikeentropy} с $\Delta t=1\ms$ наибольшее?
\end{exercise}

\begin{exercise}[\lvlB]
\label{ex:04:3}
\emph{Какая корреляция вредна.} $N=1000$ нейронов, дисперсия $\sigma^2=1$, попарная корреляция $c=0.1$. Вычислите информацию Фишера $\mathcal I=f'^{\mathsf T}\Sigma^{-1}f'$ ($\Sigma$ --- матрица ковариаций) для одинаковых наклонов $f'=\mathbf 1$ и для наклонов $\pm1$ с $\mathbf 1^{\mathsf T}f'=0$. Во сколько раз они различаются?
\end{exercise}

\begin{exercise}[\lvlB]
\label{ex:04:4}
\emph{Моделирование: детектор совпадений.} Нейрон~\eqref{eq:glutneuron} получает $1000$ пуассоновских входов по $5$ импульсов в секунду, $w=1$; порог $\theta$ таков, что свободное напряжение пересекает его раз в секунду. Моделью найдите, сколько одновременных входов $m$ зажигают нейрон с вероятностью $1/2$ при $\tau=10$ и $2\ms$.
\end{exercise}

\begin{exercise}[\lvlB]
\label{ex:04:5}
\emph{Проектирование: детектор относительных перемен.} Подберите синапс~\eqref{eq:depression}: установившийся ток при $40$ импульсах в секунду не больше чем на $10\,\%$ выше, чем при $10$, а после скачка до $40$ ток выходит на новый уровень с постоянной времени не больше $50\ms$. Каково наименьшее $\tau_{\text{вос}}$ и при каком $U$?
\end{exercise}

\begin{exercise}[\lvlA]
\label{ex:04:6}
\emph{Время первого импульса и пачки.} (а)~По~\eqref{eq:latency} найдите вход, при котором первый импульс приходит ровно через одну постоянную времени. От чего он зависит? (б)~При вероятности выброса $p=0.2$ сколько импульсов нужно в пачке, чтобы потерять все с вероятностью ниже $5\,\%$?
\end{exercise}

\begin{exercise}[\lvlC]
\label{ex:04:7}
\emph{Исследование: сколько бит в расположении.} Нейрон --- независимые клетки~\eqref{eq:spikeentropy}, $r=10$, $\Delta t=1\ms$. (а)~Разложите энтропию «слова» из $20$ клеток на энтропию числа импульсов и остаток. (б)~Смоделируйте оценку энтропии по частотам слов из $N=30$ и $10^4$ записей. Насколько она занижена?
\end{exercise}

\chapter{Что вычисляют \nt{СЕРО}-нейроны}
\label{sec:serocomputer}

\section{Первый широковещательный канал}

До сих пор все нейроны книги говорили по адресной шине: \nt{ГЛУТ} и \nt{ГАМК} посылали импульсы определённым получателям, и мы выяснили, что такая сеть вычисляет и на каком языке она говорит. Теперь переходим ко второй шине из раздела~\ref{sec:buses} --- широковещательной. Первый её канал --- \nt{СЕРО}. Как и раньше, разрешим себе только то, что бывает в живом организме.

В этой главе появятся три устройства, которыми потом будут пользоваться все широковещательные каналы: нейрон, который при ровном фоновом притоке работает сам, пока его не остановили; провод, который на самом деле объём ткани; и медленная концентрация вместо отдельных импульсов. \nt{ДОФА}, \nt{НОРА} и \nt{АЦЕТ} в главах~\ref{sec:dofa}--\ref{sec:acet} будут собраны из тех же деталей.

С точки зрения инженера у \nt{СЕРО}-нейрона четыре важных отличия.

\emph{Первое: знак выбирает получатель.} У серотонина есть рецепторы обоих знаков, и правило~\eqref{eq:dalesign} здесь не действует: знак определяет не отправитель, а рецептор получателя (утверждение~\ref{prop:receptor})~\cite{BarnesSharp1999}.

\emph{Второе: \nt{СЕРО}-нейрон работает сам, если есть ровный фон.} В бодрствовании он ровно, несколько раз в секунду, выдаёт импульсы, и толчок на каждый импульс ему не нужен: это ритмоводитель. Но ему нужен постоянный фоновый приток. В срезе мозга, где этого притока нет, большинство таких клеток молчит; ровный ритм возвращается, если подать норадреналин через $\alpha_1$-рецепторы, а их блокада его гасит~\cite{VandermaelenAghajanian1983}. В организме этот фон приходит от \nt{НОРА}. Для схемы это генератор, которому нужно питание: пока оно есть, он работает сам, пока его не остановит торможение.

\emph{Третье: он медленный.} Почти все рецепторы серотонина метаботропные: они не открывают канал сами, а запускают внутри клетки цепочку реакций. Ответ медленный: тормозящий ток через главный рецептор нарастает и спадает в пределах примерно секунды~\cite{CourtneyFord2016} --- в десятки раз дольше, чем ответ быстрого синапса \nt{ГЛУТ}.

\emph{Четвёртое: он выбрасывает медиатор не в провод, а в объём.} В областях, куда идут выходы \nt{СЕРО}-нейронов, серотонин часто выходит не в узкую щель, а в межклеточное пространство и расплывается вокруг~\cite{Bunin1998}. Получатель чувствует концентрацию и не может знать, от какого отправителя пришла молекула.

Отдельные импульсы при таких временах уже неважны, поэтому описывать сеть будем частотами $r_j$ и концентрациями. Концентрация $c$ серотонина в объёме, куда выбрасывают медиатор нейроны $j$, растёт от выброса и убывает, потому что медиатор откачивают обратно:
\begin{empheq}[box=\keybox]{equation}
\tau_c\,\frac{\dd c}{\dd t} \;=\; -c \;+\; \kappa\sum_j r_j ,
\qquad
r_i \;=\; f_i\bigl(b_i + s_i\,g\,c\bigr), \quad s_i=\pm1 .
\label{eq:seroneuron}
\end{empheq}
Это ещё один способ прочитать поток импульсов, дополняющий утверждение~\ref{prop:reader}: объём не видит ни отдельных импульсов, ни их порядка, а только частоту, усреднённую за время $\tau_c$. Первое уравнение --- та же RC-цепь, что и мембрана: $\tau_c$ --- время жизни медиатора в объёме; прямо оно не измерено, и мы принимаем его порядка секунды, $\kappa$ --- сколько медиатора даёт один импульс. Второе описывает получателя: $b_i$ --- его собственный приток, у ритмоводителя положительный (в него входит и ровный фоновый вход от \nt{НОРА}); $g$ --- чувствительность рецептора; знак $s_i$ выбирает получатель; $f_i$ --- неубывающая передаточная характеристика.

\section{НЕ за один нейрон}

Возьмём ритмоводитель с тормозящим рецептором, $s=-1$. Пока серотонина вокруг нет, он работает на своём собственном притоке $b$. Когда серотонин появляется, приток уменьшается на $gc$, и если $gc>b$, нейрон замолкает. Это НЕ: выход есть, когда входа нет. На него ушёл один нейрон (рис.~\ref{fig:serogates}). Утверждение~\ref{prop:monotone} здесь неприменимо: оно требовало положительных весов.

Если на тот же ритмоводитель действует серотонин из двух разных объёмов, он молчит, когда серотонин есть хотя бы в одном. Это ИЛИ-НЕ, а из ИЛИ-НЕ можно собрать любую логическую функцию. Двухпроводный код \nt{ГЛУТ}-сети здесь не нужен: отсутствие сигнала вычисляют нейроны, которые работают, пока их не остановили.

\begin{figure}[t]
\centering
\begin{tikzpicture}[>=Stealth,
  nrn/.style={circle,draw,very thick,fill=blue!6,minimum size=0.95cm,inner sep=0pt},
  pace/.style={nrn,double,double distance=1.2pt},
  inh/.style={-{Bar[width=7pt]},thick,red!70!black}]
\node[pace] (N1) at (0,0) {};
\draw[inh] (-1.6,0) node[left,black] {$c$} -- (N1);
\draw[->,very thick,red!70!black] (N1) -- (1.3,0) node[right,black] {$\bar c$};
\node[below=0.55cm] at (0,0) {\small НЕ};
\node[pace] (N2) at (5,0) {};
\draw[inh] (3.4,0.45) node[left,black] {$c_1$} -- (N2);
\draw[inh] (3.4,-0.45) node[left,black] {$c_2$} -- (N2);
\draw[->,very thick,red!70!black] (N2) -- (6.3,0) node[right,black] {$\overline{c_1\vee c_2}$};
\node[below=0.55cm] at (5,0) {\small ИЛИ-НЕ};
\end{tikzpicture}
\caption{Логика из \nt{СЕРО}-нейронов. Двойной кружок --- ритмоводитель: при ровном фоновом притоке он работает сам, пока его не затормозят. Линия с чертой --- тормозящий рецептор, $s=-1$; $c$, $c_1$, $c_2$ --- концентрации серотонина в объёмах вокруг получателя. Слева --- НЕ, справа --- ИЛИ-НЕ.}
\label{fig:serogates}
\end{figure}

\begin{keyblock}
\begin{proposition}[Полнота \nt{СЕРО}-логики]
\label{prop:serocomplete}
Если в сети есть ритмоводители с тормозящими рецепторами и выбросы в разные объёмы не смешиваются, сеть~\eqref{eq:seroneuron} может вычислить любую логическую функцию, и каждый бит передаётся одним проводом.
\end{proposition}
\end{keyblock}

Логически \nt{СЕРО}-сеть сильнее \nt{ГЛУТ}, но условие «выбросы в разные объёмы не смешиваются» в мозге не выполняется (раздел~\ref{sec:volume}).

\section{Отрицательная обратная связь}

У серотониновых нейронов есть тормозящие рецепторы на них самих~\cite{Sprouse1987}. Серотонин, который они выбросили, возвращается к ним же и притормаживает их. Для инженера это знакомая схема: усилитель с отрицательной обратной связью (рис.~\ref{fig:serofeedback}).

Что здесь измерено, а что --- модель. В самом ядре шва, где сидят \nt{СЕРО}-нейроны, серотонин тормозит соседей не через общий объём, а точка-точка, через синапсы: переносчик быстро убирает его и не даёт накопиться вне щели. Одиночный выброс даёт лишь короткую паузу в разряде, а средняя частота от этого не меняется~\cite{CourtneyFord2016}. Поэтому петля ниже, замкнутая через общий объём, --- модель: мы считаем, что делала бы такая обратная связь, если бы она работала на уровне средних частот.

\begin{figure}[t]
\centering
\begin{tikzpicture}[>=Stealth,
  nrn/.style={circle,draw,very thick,fill=blue!6,minimum size=0.95cm,inner sep=0pt},
  pace/.style={nrn,double,double distance=1.2pt},
  blk/.style={draw,very thick,rounded corners=2pt,fill=orange!10,minimum height=0.9cm,minimum width=2.2cm,align=center,font=\small},
  inh/.style={-{Bar[width=7pt]},thick,red!70!black}]
\node[pace] (N) at (0,0) {};
\node[blk] (V) at (3.6,0) {объём\\$\tau_c$};
\draw[->,very thick] (N) -- node[above] {\small $r$} (V);
\draw[->,very thick] (V) -- (6.0,0) node[right] {\small $c$ --- всем получателям};
\draw[inh] (V.south) -- ++(0,-0.9) -| node[pos=0.25,below] {\small $-g\,c$} (N.south);
\draw[->,thick,blue!60!black] (-1.7,0) node[left,black] {\small $b$} -- (N);
\end{tikzpicture}
\caption{\nt{СЕРО}-нейрон с обратной связью через собственный медиатор: концентрация $c$ идёт ко всем получателям и тормозит сам нейрон. В модели это регулятор уровня; в мозге такая роль петли --- гипотеза.}
\label{fig:serofeedback}
\end{figure}

Посчитаем, что делает эта петля. Пусть передаточная характеристика линейна, $f(u)=au$ при $u>0$. В равновесии $c=\kappa r$ и $r=a(b-gc)$. Подставив одно в другое, получим
\begin{empheq}[box=\keybox]{equation}
r \;=\; \frac{a\,b}{1+G}, \qquad G \;=\; a\,g\,\kappa .
\label{eq:feedback}
\end{empheq}
Величина $G$ --- усиление петли. Это та же формула, что и для любого усилителя с отрицательной обратной связью, и она даёт те же два выигрыша.

\emph{Устойчивость к разбросу.} Пусть возбудимость нейрона $a$ изменилась на долю $\varepsilon$. Без обратной связи на ту же долю изменилась бы и частота. С обратной связью, при $G\gg1$, частота $r\approx b/(g\kappa)$ от $a$ почти не зависит: её изменение в $1+G$ раз меньше. При $G=9$ разброс подавляется вдесятеро.

\emph{Быстрота.} Подставим $r=a(b-gc)$ в первое уравнение~\eqref{eq:seroneuron}: получим $\tau_c\,\dd c/\dd t=-(1+G)\,c+\kappa ab$. Концентрация подходит к своему уровню с постоянной времени $\tau_c/(1+G)$ --- в $1+G$ раз быстрее, чем без обратной связи.

Вот первое вычисление, которое могла бы выполнять \nt{СЕРО}-система: держать общий уровень серотонина в мозге на заданной отметке, как термостат держит температуру. Это гипотеза: в опыте автоторможение пока видно только как короткие паузы, не меняющие средней частоты.

\section{Из импульсов --- в уровень}

Второе вычисление спрятано в первом уравнении~\eqref{eq:seroneuron}. Нейроны выдают отдельные импульсы, а получатель чувствует концентрацию, которая сглаживает их за время $\tau_c$. Если частота источника меняется, концентрация следует за ней с запаздыванием:
\begin{equation}
c(t) \;=\; \frac{\kappa}{\tau_c}\int_{-\infty}^{t} e^{-(t-t')/\tau_c}\,\sum_j r_j(t')\,\dd t' .
\label{eq:lowpass}
\end{equation}
Это скользящее среднее активности источников за последние несколько $\tau_c$ --- фильтр нижних частот (рис.~\ref{fig:lowpass}). Так простейший цифро-аналоговый преобразователь пропускает импульсы через RC-цепь и превращает их частоту в уровень. \nt{СЕРО}-система делает то же с сотнями тысяч нейронов сразу.

Эта величина заодно и память. После того как источники замолчали, концентрация держится ещё время порядка $\tau_c$. Это не бит, а плавно тающий след.

\begin{figure}[t]
\centering
\begin{tikzpicture}[>=Stealth,x=1.45cm,y=1cm]
\draw[gray!70!black] (0,3.2) -- (8.1,3.2);
\node[left,font=\small] at (-0.05,3.4) {импульсы};
\node[above,font=\scriptsize,gray!40!black] at (1,3.65) {$2$ в секунду};
\node[above,font=\scriptsize,gray!40!black] at (3.5,3.65) {$8$ в секунду};
\node[above,font=\scriptsize,gray!40!black] at (6.5,3.65) {$2$ в секунду};
\draw[->,gray!70!black] (0,0) -- (8.3,0) node[right,black] {\small $t$};
\draw[->,gray!70!black] (0,0) -- (0,2.75) node[left,black] {\small $c$};
\foreach \s in {0.136,0.529,0.760,1.223,1.714,1.748,1.755,2.663,2.701,2.734,3.414,3.493,3.719,3.800,3.928,3.948,4.074,4.327,4.420,4.589,4.728,4.736,4.914,5.025,5.205,5.220,6.224,6.544,7.178}{\draw[thick,blue!60!black] (\s,3.2) -- (\s,3.6);}
\foreach \a/\b/\v in {0.000/0.136/0.000,0.136/0.529/1.000,0.529/0.760/1.675,0.760/1.223/2.330,1.223/1.714/2.466,1.714/1.748/2.509,1.748/1.755/3.425,1.755/2.663/4.402,2.663/2.701/2.775,2.701/2.734/3.672,2.734/3.414/4.553,3.414/3.493/3.306,3.493/3.719/4.055,3.719/3.800/4.235,3.800/3.928/4.905,3.928/3.948/5.316,3.948/4.074/6.211,4.074/4.327/6.476,4.327/4.420/6.028,4.420/4.589/6.493,4.589/4.728/6.483,4.728/4.736/6.642,4.736/4.914/7.589,4.914/5.025/7.351,5.025/5.205/7.579,5.205/5.220/7.331,5.220/6.224/8.221,6.224/6.544/4.012,6.544/7.178/3.914,7.178/8.000/3.076}{\draw[very thick,orange!85!black] plot[domain=\a:\b,samples=10] (\x,{0.25*\v*exp(-(\x-\a))});}
\foreach \s/\p/\q in {0.136/0.000/1.000,0.529/0.675/1.675,0.760/1.330/2.330,1.223/1.466/2.466,1.714/1.509/2.509,1.748/2.425/3.425,1.755/3.402/4.402,2.663/1.775/2.775,2.701/2.672/3.672,2.734/3.553/4.553,3.414/2.306/3.306,3.493/3.055/4.055,3.719/3.235/4.235,3.800/3.905/4.905,3.928/4.316/5.316,3.948/5.211/6.211,4.074/5.476/6.476,4.327/5.028/6.028,4.420/5.493/6.493,4.589/5.483/6.483,4.728/5.642/6.642,4.736/6.589/7.589,4.914/6.351/7.351,5.025/6.579/7.579,5.205/6.331/7.331,5.220/7.221/8.221,6.224/3.012/4.012,6.544/2.914/3.914,7.178/2.076/3.076}{\draw[very thick,orange!85!black] (\s,0.25*\p) -- (\s,0.25*\q);}
\draw[thick,densely dashed,black] plot[domain=0:2,samples=30] (\x,{0.25*2*(1-exp(-\x))})
  plot[domain=2:5,samples=40] (\x,{0.25*(8-6.271*exp(-(\x-2)))})
  plot[domain=5:8,samples=40] (\x,{0.25*(2+5.688*exp(-(\x-5)))});
\foreach \x in {0,2,5,8}{\draw[gray!60!black] (\x,-0.05) -- (\x,0.05) node[below=3pt,font=\scriptsize] {\x\,с};}
\end{tikzpicture}
\caption{Из импульсов --- в уровень. Сверху --- импульсы \nt{СЕРО}-нейронов: две секунды редко, три секунды часто, потом снова редко. Снизу --- концентрация серотонина~\eqref{eq:lowpass} при $\tau_c=1\,\text{с}$. Штриховая линия --- то же, если бы импульсы шли идеально ровно.}
\label{fig:lowpass}
\end{figure}

\section{Провод --- это объём}
\label{sec:volume}

Вернёмся к условию утверждения~\ref{prop:serocomplete}. Серотонин, выброшенный поблизости разными отправителями, складывается в одну концентрацию.

\emph{Провод здесь --- не волокно, а объём ткани.} Два сигнала остаются разными, только если их выбрасывают в области, разнесённые дальше, чем успевает расплыться медиатор. За время $\tau$ молекула с коэффициентом диффузии $D$ уходит на
\begin{empheq}[box=\keybox]{equation}
\ell \;\sim\; \sqrt{D\,\tau}\, .
\label{eq:diffusionlength}
\end{empheq}
Извилистость межклеточных щелей замедляет диффузию небольшой молекулы примерно в $2.6$ раза~\cite{Sykova2008}, а коэффициент диффузии самого серотонина в ткани не измерен; по размеру молекулы оценим его в несколько единиц на $10^{-6}$~см$^2$/с. Если сигнал живёт $\tau\sim1$~с (тоже оценка), то $\ell\sim\sqrt{3\cdot10^{-6}}$~см $\approx20$~мкм. Адрес в такой сети --- это \emph{место}: получатель узнаёт, от кого сигнал, только по тому, где он сам находится.

Настоящие \nt{СЕРО}-нейроны устроены так, что этих отдельных адресов почти не остаётся. Выход каждого из них раскидан по огромным участкам мозга: ветви одной клетки доходят до многих далёких друг от друга областей~\cite{GagnonParent2014}. Мест выброса на клетку по оценке около $10^5$ (таблица~\ref{tab:types}); прямо их не подсчитали. «Провода» разных \nt{СЕРО}-нейронов перекрываются и смешиваются. Совсем в один провод они всё же не сливаются: \nt{СЕРО}-нейроны делятся на несколько подсистем, которые шлют выход в разные области и по-разному отвечают на одно и то же событие~\cite{Ren2018}. Но таких подсистем единицы, а не сотни тысяч.

\begin{keyblock}
\begin{proposition}[Число независимых проводов]
\label{prop:volumewires}
В сети с объёмной передачей число независимых сигналов ограничено не числом нейронов, а числом непересекающихся областей размером $\ell\sim\sqrt{D\tau}$, в которые они выбрасывают медиатор. Если выход каждого нейрона раскидан по большому объёму, вся сеть передаёт лишь несколько общих переменных.
\end{proposition}
\end{keyblock}

Поэтому логические схемы утверждения~\ref{prop:serocomplete} в мозге собрать не из чего. А регулятору~\eqref{eq:feedback} и фильтру~\eqref{eq:lowpass} отдельные провода не нужны: им хватает одной общей величины. Фильтр прямо следует из измеренного выброса в объём в областях назначения~\cite{Bunin1998}; регулятор уровня --- гипотеза.

\section{Горизонт планирования}
\label{sec:horizon}

Для чего может пригодиться одна общая медленная величина? Хороший кандидат --- параметр, который должен быть одинаков для всей сети и меняться не быстрее, чем за секунды или минуты. В главе~\ref{sec:neurontypes} такой параметр уже появлялся: коэффициент $\gamma$ в формуле ошибки~\eqref{eq:rpe}, который говорит, насколько будущая награда дешевле немедленной. В непрерывном времени его место занимает \emph{горизонт} $\tau_\gamma$: награда $R$, до которой ждать время $D$, стоит сейчас $R\,e^{-D/\tau_\gamma}$.

Горизонт решает, стоит ли ждать. Пусть можно сразу взять маленькую награду $r_0$ или подождать большую $R$. Ждать выгодно, пока
\begin{empheq}[box=\keybox]{equation}
D \;<\; \tau_\gamma\,\ln\frac{R}{r_0} .
\label{eq:patience}
\end{empheq}
При $\tau_\gamma=2\,\text{с}$ и втрое большей отложенной награде стоит ждать до $2.2\,\text{с}$; если горизонт вдвое длиннее, --- до $4.4\,\text{с}$. Терпение пропорционально горизонту.

Формула~\eqref{eq:patience} опирается на экспоненциальное обесценивание. Измеренное обесценивание на задержках в секунды ближе к гиперболе $R/(1+D/\tau_\gamma)$: так у обезьян падают и ценность отложенной награды в выборе, и ответ \nt{ДОФА}-нейронов на её предвестник~\cite{KobayashiSchultz2008}. Для гиперболы условие~\eqref{eq:patience} заменяется на $D<\tau_\gamma\,(R/r_0-1)$: зависимость от отношения наград уже не логарифмическая, а линейная, и при тех же числах граница ожидания сдвигается почти вдвое. Гипербола получается из той же экспоненты, если задержка случайна (задача~\ref{ex:05:7}), а пропорциональность терпения масштабу времени сохраняется.

По гипотезе, горизонт и задаёт \nt{СЕРО}~\cite{Doya2002}. Для такой роли у него всё есть: один общий уровень на всю сеть, который меняется за секунды. Есть и прямые опыты: если искусственно включить серотониновые нейроны, животное дольше ждёт отложенную награду, прежде чем бросить ожидание~\cite{Miyazaki2014}, и дольше пытается добыть награду там, где она стала редкой~\cite{Lottem2018}. Как именно концентрация серотонина превращается в $\tau_\gamma$ внутри сети, неизвестно. В следующей главе горизонт войдёт в ошибку, которую рассылает \nt{ДОФА}.

\section{Итог}

\begin{table}[t]
\centering
\small
\begin{tabular}{>{\raggedright}p{3.6cm}>{\raggedright}p{4.3cm}>{\raggedright\arraybackslash}p{4.3cm}}
\toprule
& \nt{ГЛУТ} & \nt{СЕРО} \\
\midrule
время реакции & миллисекунды & доли секунды --- секунда \\
знак действия & только $+$ & $+$ и $-$, выбирает получатель \\
без входа & молчит & работает сам при ровном фоновом притоке \\
адресация & точка --- точка & объём, адрес --- место \\
НЕ & только через датчики «выключения» & один нейрон \\
память & самоподдерживающаяся активность, гаснет от усталости & концентрация, тает за время $\tau_c$ \\
главные вычисления & совпадения, задержки, логика, последовательности & сглаживание; регулирование уровня и горизонт $\tau_\gamma$ (гипотезы) \\
\bottomrule
\end{tabular}
\caption{Что вычисляют сети из одних \nt{ГЛУТ}-нейронов и из одних \nt{СЕРО}-нейронов.}
\label{tab:glutvssero}
\end{table}

Как логический элемент (таблица~\ref{tab:glutvssero}) \nt{СЕРО}-нейрон богаче \nt{ГЛУТ}, но как система связи он --- широковещательная рассылка: медленная и почти без адресов. Поэтому он не пытается быть процессором, а сглаживает и рассылает несколько общих величин, о которых шла речь в главе~\ref{sec:neurontypes}, --- по гипотезе, в том числе горизонт планирования. Ритмоводитель, объём и медленная концентрация встретятся нам ещё трижды: у \nt{ДОФА}, \nt{НОРА} и \nt{АЦЕТ}.

\section{Задачи}

\begin{exercise}[\lvlA]
\label{ex:05:1}
\emph{Провод --- это объём: оценка.} Примите для серотонина в ткани $D=3\cdot10^{-6}$~см$^2$/с (раздел~\ref{sec:volume}).
(а)~По~\eqref{eq:diffusionlength} найдите $\ell$ для сигналов, живущих $\tau=0.1\,\text{с}$, $1\,\text{с}$ и $1$~мин.
(б)~Сколько непересекающихся кубиков со стороной $\ell(1\,\text{с})$ помещается в $1\,\text{мм}^3$ ткани? Это верхняя граница числа независимых проводов утверждения~\ref{prop:volumewires} в таком объёме.
(в)~За какое время медиатор расплывётся диффузией на $1$~мм и на $5$~см? Чем тогда обеспечена широковещательность \nt{СЕРО} --- диффузией или ветвлением выходного провода с $\sim10^5$ мест выброса (таблица~\ref{tab:types})? Что остаётся на долю диффузии?
\end{exercise}

\begin{exercise}[\lvlA]
\label{ex:05:2}
\emph{Регулятор уровня: вывод.}
(а)~Из~\eqref{eq:seroneuron} при $f(u)=au$ выведите~\eqref{eq:feedback} и покажите, что относительная чувствительность $S=(a/r)\,\partial r/\partial a$ равна $1/(1+G)$ точно, а не только при $G\gg1$.
(б)~При $G=9$ возбудимость $a$ выросла на $20\,\%$. На сколько процентов изменится $r$ (без линеаризации)?
(в)~Пусть $f(u)=a\varphi(u)$ с произвольной гладкой неубывающей $\varphi$. Докажите, что в окрестности равновесия $S=1/(1+G)$ и постоянная времени равна $\tau_c/(1+G)$, где $G=a\,\varphi'(u^*)\,g\kappa$, $u^*=b-gc^*$. Где на сигмоиде $\varphi$ регулятор перестаёт работать?
\end{exercise}

\begin{exercise}[\lvlB]
\label{ex:05:3}
\emph{Медленный рецептор в петле.} Рецепторы \nt{СЕРО} метаботропные, и их ответ запаздывает.
(а)~Опишите рецептор RC-звеном: $\tau_r\,\dd u/\dd t=-u+b-gc$, $r=au$, а объём --- первым уравнением~\eqref{eq:seroneuron}. Выведите характеристическое уравнение замкнутой петли, её собственную частоту $\omega_n$ и коэффициент затухания $\zeta$. При $\tau_c=1\,\text{с}$ и $\tau_r=0.3\,\text{с}$ найдите $G$, при котором $\zeta=1/\sqrt2$, и перерегулирование переходного процесса при $G=9$.
(б)~Замените RC-звено чистой задержкой: $r(t)=a\bigl(b-gc(t-T)\bigr)$. Докажите, что при $G>1$ петля устойчива тогда и только тогда, когда
$T<T^*=\tau_c\arccos(-1/G)/\sqrt{G^2-1}$, а при $G\le1$ --- при любой задержке. Вычислите $T^*$ при $G=2$ и $G=9$.
(в)~Какое подавление разброса $1/(1+G)$ реально доступно регулятору, если задержка в петле --- сотни миллисекунд?
\end{exercise}

\begin{exercise}[\lvlB]
\label{ex:05:4}
\emph{Дробовой шум уровня.} Пусть импульсы всех источников образуют пуассоновский поток с частотой $\lambda$; по~\eqref{eq:lowpass} каждый импульс в момент $t_k$ добавляет к $c$ слагаемое $(\kappa/\tau_c)e^{-(t-t_k)/\tau_c}$.
(а)~Докажите, что $\langle c\rangle=\kappa\lambda$ и $\operatorname{Var}c=\kappa^2\lambda/(2\tau_c)$, так что коэффициент вариации $\mathrm{CV}=1/\sqrt{2\lambda\tau_c}$. (Разбейте время на интервалы $\dd t$, в каждом из которых импульс есть независимо с вероятностью $\lambda\,\dd t$.)
(б)~Найдите CV для $\lambda=2$ и $8$ импульсов в секунду при $\tau_c=1\,\text{с}$ (рис.~\ref{fig:lowpass}).
(в)~Для строго периодических импульсов найдите отношение максимума к минимуму установившейся «пилы» $c(t)$ при тех же частотах.
(г)~Каков CV, если все $3\cdot10^5$ \nt{СЕРО}-нейронов по $2$ импульса в секунду выбрасывают медиатор в один объём? Какое $\tau_c$ понадобилось бы одному нейрону с частотой $4$ импульса в секунду, чтобы получить $\mathrm{CV}=10\,\%$? Почему уровень делают многие нейроны, а не один медленный?
\end{exercise}

\begin{exercise}[\lvlB]
\label{ex:05:5}
\emph{Моделирование: обратная связь против шума.} Для этой главы отдельной модели в \texttt{models/} нет; напишите её на Python. $N=50$ ритмоводителей, $i$-й выдаёт пуассоновские импульсы с частотой $r_i=a_i[b-gc]_+$, где $a_i$ равномерно распределены в $[2.8,\,5.2]$ импульса в секунду; концентрация --- по~\eqref{eq:seroneuron}, каждый импульс увеличивает $c$ на $\kappa/\tau_c$, $\kappa=1/N$, $\tau_c=1\,\text{с}$. Интегрируйте по Эйлеру с шагом $1\ms$ в течение $300\,\text{с}$, первые $20\,\text{с}$ отбросьте. Прогоны: (i)~$b=1$, $g=2.25$ ($G=9$); (ii)~$b=0.1$, $g=0$ --- без обратной связи, тот же средний уровень; (iii)~оба варианта с $a_i$, умноженными на $1.2$. Измерьте средний уровень $c$, его CV, время, за которое автокорреляция $c$ падает в $e$ раз, и частоты отдельных нейронов.
(а)~Линеаризуйте систему вокруг равновесия и покажите, что с обратной связью $\mathrm{CV}^2=1/(2\tau_cN\bar ab)$: при том же среднем уровне обратная связь снижает CV в $\sqrt{1+G}$ раз. Предскажите все измеряемые величины.
(б)~Сравните модель с предсказанием.
(в)~Выравнивает ли обратная связь частоты отдельных нейронов? Объясните, что именно регулирует общий объём.
\end{exercise}

\begin{exercise}[\lvlB]
\label{ex:05:6}
\emph{Проектирование: XOR на \nt{СЕРО}-логике.} Вентиль --- один ритмоводитель с пороговой характеристикой: $r=r_0$, если $b-g\sum_kc_k>0$, и $r=0$ иначе; он выбрасывает медиатор в свой объём, $\tau_c\,\dd c/\dd t=-c+\kappa r$. Логическая единица --- $c_1=\kappa r_0$, ноль --- $c=0$.
(а)~Покажите, что вентиль вычисляет ИЛИ-НЕ своих входов и может управлять следующими вентилями с тем же $g$, если $G'=g\kappa r_0/b>1$. Почему один выход можно подать на сколько угодно входов?
(б)~Соберите XOR$(A,B)$ из пяти вентилей ИЛИ-НЕ. Сколько отдельных объёмов нужно и на каком расстоянии их надо держать при $\tau_c=1\,\text{с}$ (задача~\ref{ex:05:1})?
(в)~Вход вентиля пересёк порог в момент $0$. Найдите, через какое время выход пересечёт порог $b/g$ при переключении $0\to1$ и $1\to0$. Покажите, что сумма этих времён минимальна при $G'=2$.
(г)~Найдите время установления XOR на самом длинном пути при $\tau_c=1\,\text{с}$ для $G'=4$ и $G'=2$.
\end{exercise}

\begin{exercise}[\lvlB]
\label{ex:05:7}
\emph{Терпение при неизвестной задержке.}
(а)~Выведите~\eqref{eq:patience}.
(б)~Животное соглашается ждать втрое большую награду, если ожидание не длиннее $6\,\text{с}$. Какой горизонт $\tau_\gamma$ это означает? До скольких секунд оно будет ждать, если включение \nt{СЕРО}-нейронов удвоит горизонт?
(в)~Пусть задержка $D$ случайна и распределена экспоненциально со средним $\bar D$. Покажите, что ждать выгодно при $\bar D<\tau_\gamma\,(R/r_0-1)$, и сравните с фиксированной задержкой при $\tau_\gamma=2\,\text{с}$, $R=3r_0$.
(г)~Докажите, что при любом распределении задержки со средним $\bar D$ ценность ожидания не меньше, чем при фиксированной задержке $\bar D$. Что это значит для поведения там, где задержки непредсказуемы?
\end{exercise}

\begin{exercise}[\lvlC]
\label{ex:05:8}
\emph{Как концентрация может задавать горизонт.} В разделе~\ref{sec:horizon} открыт вопрос, как уровень серотонина превращается в $\tau_\gamma$. В~\eqref{eq:ctd} горизонт входит только слагаемым $-V/\tau_\gamma$. Пусть \nt{ДОФА}-нейрон получает оценку $V$ двумя путями: прямым возбуждающим с весом $k_1$ и тормозящим через \nt{ГАМК}-посредника, который сглаживает сигнал с постоянной времени $\tau_d$, с весом $k_2$.
(а)~Покажите, что для медленно меняющегося $V$ сумма путей равна $(k_1-k_2)V+k_2\tau_d\,\dd V/\dd t$, и найдите $k_1$, $k_2$, воспроизводящие $\dd V/\dd t-V/\tau_\gamma$, при $\tau_d=0.1\,\text{с}$, $\tau_\gamma=2\,\text{с}$.
(б)~Пусть \nt{СЕРО} умножает прямой вес на $m$: $k_1\to mk_1$. Найдите $\tau_\gamma(m)$ и чувствительность $\dd\ln\tau_\gamma/\dd\ln m$ при $m=1$. На сколько процентов надо изменить $m$, чтобы удвоить горизонт? При каком $m$ горизонт становится бесконечным?
(в)~Рассогласование весов на долю $\mu$ даёт относительную ошибку горизонта $\approx\mu\tau_\gamma/\tau_d$, а конечное $\tau_d$ --- ложную ошибку $\approx\tau_d/\tau_\gamma$ (задача~\ref{ex:06:6}). Покажите, что обе ошибки нельзя сделать меньше $\sqrt\mu$ одновременно. Что это значит, если веса совпадают с точностью $1\,\%$?
(г)~Открытая часть. Критик учится по тому $\delta$, который реально вычисляет схема. Объясните, почему рассогласование весов тогда не ошибка, а просто другой горизонт, и почему высокая чувствительность из~(б) может быть полезна. Предложите механизм, в котором $\tau_\gamma$ зависит от $c$ без тонкой настройки, и опыт, который отличит его от механизма~(б).
\end{exercise}

\chapter{Нейронные сети, которые умеют учиться: добавляем \nt{ДОФА}}
\chaptermark{Сети, которые учатся}
\label{sec:dofa}

\section{Правила игры}

Во всех схемах предыдущих глав веса ставил инженер. В организме инженера нет. Веса должны устанавливаться сами, по ходу работы, и так, чтобы сеть делала что-то полезное. Это и называется обучением.

К прежним правилам добавляется ещё одно. Синапс меняет свой вес сам, и знать он может только то, что происходит \emph{рядом с ним}: активность отправителя $x$, активность получателя $y$ и те вещества, которые до него доходят. Никто не может прислать синапсу персональную поправку.

Это исключает главный метод искусственных нейронных сетей --- обратное распространение ошибки. Там ошибка выхода проходит сеть в обратную сторону по тем же весам, в отдельную фазу после прямого прохода, и каждый вес получает свою поправку. Для этого нужны обратные провода, точно повторяющие прямые, и общий такт. Ни того ни другого в мозге не найдено, по крайней мере в той форме, что в искусственных сетях~\cite{Lillicrap2020,WhittingtonBogacz2019}.

Изменение веса будем записывать как медленный непрерывный процесс:
\begin{equation}
\frac{\dd w}{\dd t} \;=\; \eta\,\Phi(\text{то, что известно синапсу}),
\label{eq:plasticity}
\end{equation}
где $\eta$ --- скорость обучения, настолько малая, что за время одного импульса вес почти не меняется. Вся глава --- о том, какой может быть функция $\Phi$.

\section{Где хранится вес}
\label{sec:weightstore}

Что такое синапс, который «меняет свой вес»? У связи две стороны: окончание провода отправителя (\emph{пресинаптическая} сторона) и участок мембраны получателя с рецепторами (\emph{постсинаптическая} сторона), а между ними --- узкая щель. У \nt{ГЛУТ}-связей мембрана получателя обычно сидит на \emph{шипике} --- маленьком выступе входной ветви получателя. Вес связи удобно разложить на три множителя~\cite{delCastillo1954}:
\begin{empheq}[box=\keybox]{equation}
w \;=\; N_s\,p\,A_1 .
\label{eq:quantal}
\end{empheq}
Здесь $N_s$ --- число мест выброса между двумя нейронами, $p$ --- вероятность, что импульс выбросит порцию медиатора в одном месте (та самая $p$ из главы~\ref{sec:code}), $A_1$ --- ответ получателя на одну порцию, то есть сколько рецепторов её ловит. Множитель $p$ живёт у отправителя, $A_1$ --- у получателя, а $N_s$ требует обеих сторон: новые места выброса вырастают, старые исчезают, и это продолжается всю жизнь.

\emph{Решает получатель.} У типичного \nt{ГЛУТ}-синапса один из рецепторов глутамата проводит ток, только если совпали два условия: глутамат пришёл от отправителя, и мембрана получателя уже возбуждена~\cite{Nowak1984}. Это правило Хебба, встроенное в одну молекулу, --- детектор совпадений из главы~\ref{sec:glutcomputer}, поставленный на вход самого синапса. Сработав, рецептор впускает в получателя кальций, и кальций запускает изменение веса. Получатель --- единственное место, где видны сразу и вход, и выход, поэтому решение принимается у него.

\emph{Исполнить решение может любая сторона.} Самое изученное долговременное усиление идёт у получателя: в мембрану встраиваются новые рецепторы~\cite{Malinow2002}, и шипик растёт~\cite{Matsuzaki2004} --- растёт $A_1$. Но получатель может изменить и $p$: он выбрасывает вещество, которое идёт через щель назад, к отправителю, --- обратный канал из приложения~\ref{sec:othermediators}, --- и тот начинает выбрасывать меньше медиатора~\cite{WilsonNicoll2001}. А кратковременные изменения веса, истощение и облегчение из главы~\ref{sec:code}, --- это изменения $p$, которые отправитель ведёт сам по своей недавней активности.

Для инженера регистр веса разделён между двумя микросхемами. Правило обучения исполняет получатель, а записать результат он может и у себя, и у отправителя, через обратный канал. Поэтому слово «локально» в правилах этой главы означает не одну сторону связи, а всю связь целиком.

\section{Учиться без учителя}

Самое простое, что может сделать синапс, --- усилиться, когда отправитель и получатель активны вместе:
\begin{equation}
\frac{\dd w}{\dd t} \;=\; \eta\,x\,y .
\label{eq:hebb}
\end{equation}
Это правило Хебба~\cite{Hebb1949}. Оно локально и не требует часов. Но у него та же беда, что у \nt{ГЛУТ}-сети в главе~\ref{sec:glutcomputer}: положительная обратная связь. Сильный вес делает $y$ больше, большой $y$ ещё усиливает вес, и веса растут без предела. Лекарство --- отрицательная обратная связь по весу: пусть вес ещё и убывает пропорционально $y^2$. Для нейрона с входами $\mathbf x$ и линейным выходом $y=\mathbf w\cdot\mathbf x$ получаем
\begin{empheq}[box=\keybox]{equation}
\frac{\dd\mathbf w}{\dd t} \;=\; \eta\,y\,\bigl(\mathbf x - y\,\mathbf w\bigr) .
\label{eq:oja}
\end{empheq}
Правило по-прежнему локально: синапс знает свой вход $x_i$, выход $y$ и вес $w_i$.

\begin{keyblock}
\begin{proposition}[Что находит правило Хебба]
\label{prop:oja}
Правило~\eqref{eq:oja} приводит вектор весов к единичной длине вдоль направления, в котором входы меняются сильнее всего, --- к главному собственному вектору матрицы корреляций входов $C=\langle\mathbf x\mathbf x^{\mathsf T}\rangle$~\cite{Oja1982}. Нейрон учится выдавать самую выразительную одну величину, в которую можно сжать его входы.
\end{proposition}
\end{keyblock}

Доказательство. Усредним~\eqref{eq:oja} по входам: $\dd\mathbf w/\dd t=\eta\bigl(C\mathbf w-(\mathbf w^{\mathsf T}C\mathbf w)\,\mathbf w\bigr)$. Неподвижные точки --- собственные векторы $C$ единичной длины. Если $\mathbf w$ близок к собственному вектору $\mathbf e_k$ с числом $\lambda_k$, малая добавка вдоль другого вектора $\mathbf e_j$ растёт как $e^{\eta(\lambda_j-\lambda_k)t}$. Устойчива только та точка, для которой все $\lambda_j<\lambda_k$, то есть главный вектор.

Есть и вариант правила Хебба, который смотрит на время импульсов. Если импульс отправителя пришёл за $s$ миллисекунд \emph{до} импульса получателя, вес растёт на $A_+e^{-s/\tau_+}$; если через $s$ миллисекунд \emph{после} --- уменьшается на $A_-e^{-s/\tau_-}$, с $\tau_\pm$ порядка двадцати миллисекунд. Такое правило измерено на синапсах \nt{ГЛУТ}-нейронов~\cite{BiPoo1998}. Оно усиливает связи, которые \emph{могли быть причиной} ответа, и ослабляет опоздавшие (рис.~\ref{fig:stdp}). Прогоните через сеть много раз одну последовательность возбуждений, и в ней сами вырастут связи от каждого звена к следующему --- цепочка из главы~\ref{sec:glutcomputer}, собранная без инженера.

\begin{figure}[t]
\centering
\begin{tikzpicture}[>=Stealth,x=0.07cm,y=1.3cm]
\draw[->,gray!70!black] (-66,0) -- (68,0) node[right,black] {\small $s$};
\draw[->,gray!70!black] (0,-1.2) -- (0,1.35) node[above,black] {\small изменение веса};
\draw[very thick,blue!60!black] plot[domain=0.5:60,samples=50] (\x,{exp(-\x/20)});
\draw[very thick,red!70!black] plot[domain=-60:-0.5,samples=50] (\x,{-0.6*exp(\x/20)});
\draw[blue!60!black,thick] (0.5,0) -- (0.5,1);
\draw[red!70!black,thick] (-0.5,0) -- (-0.5,-0.6);
\foreach \x in {-40,-20,20,40}{\draw[gray!60!black] (\x,-0.04) -- (\x,0.04);}
\node[below,font=\scriptsize] at (20,-0.04) {$20\ms$};
\node[above,font=\scriptsize] at (-20,0.04) {$-20\ms$};
\node[align=left,font=\small,blue!60!black,anchor=west] at (22,0.95) {отправитель раньше:\\ связь усиливается};
\node[align=right,font=\small,red!70!black,anchor=east] at (-12,-0.95) {отправитель позже:\\ связь ослабевает};
\end{tikzpicture}
\caption{Правило Хебба по времени. По горизонтали --- $s=t_{\text{получ}}-t_{\text{отпр}}$, на сколько импульс получателя отстал от импульса отправителя; $\tau_\pm=20\ms$, $A_-=0.6A_+$. Окно шириной в несколько десятков миллисекунд --- того же порядка, что и окно совпадения~\eqref{eq:window}.}
\label{fig:stdp}
\end{figure}

Но все правила этого раздела учат тому, что \emph{часто}, а не тому, что \emph{полезно}. Синапс, которому известны только $x$ и $y$, не знает, принесло ли действие сок или боль. Чтобы учиться на результате, синапсу нужен третий сигнал.

\section{Третий множитель}

Третий сигнал приносит \nt{ДОФА}: дофаминовые нейроны рассылают широковещательно одно число --- ошибку предсказания награды $\delta$~\eqref{eq:rpe} (глава~\ref{sec:neurontypes}). Пусть изменение веса пропорционально произведению трёх множителей: активности отправителя, активности получателя и $\delta$.

Трудность в том, что награда приходит не тогда, когда сеть выбирала действие, а через сотни миллисекунд или секунды. К приходу $\delta$ произведение $xy$ давно равно нулю, и синапсу нужна память о том, что он участвовал. Самая простая память --- знакомая нам RC-цепь:
\begin{empheq}[box=\keybox]{equation}
\tau_e\,\frac{\dd e}{\dd t} \;=\; -e \;+\; x\,y ,
\qquad
\frac{\dd w}{\dd t} \;=\; \eta\,\delta(t)\,e(t) .
\label{eq:threefactor}
\end{empheq}
Величину $e$ называют \emph{следом}~\cite{Fremaux2016}. Совместная активность оставляет в синапсе метку, которая тает за время $\tau_e$. Если за это время приходит дофамин, вес меняется со знаком $\delta$; если нет, метка исчезает бесследно (рис.~\ref{fig:trace}). На \nt{ГЛУТ}-синапсах измерено: дофамин, пришедший в течение секунды-двух после совместной активности, превращает её в усиление связи, а пришедший позже --- нет~\cite{Yagishita2014}. Окна пластичности длиной в секунды нашли и там, где третьим сигналом служит не дофамин, а сильное возбуждение самого получателя~\cite{Bittner2017}.

\begin{figure}[t]
\centering
\begin{tikzpicture}[>=Stealth,x=7cm,y=1cm]
\fill[orange!12] (0.7,-0.2) rectangle (0.8,5.2);
% rows: x*y, delta, traces, weights
\foreach \y/\lab in {4.4/{$x\,y$},3.1/{$\delta$},1.4/{след $e$},0/{вес $w$}}{
  \draw[gray!70!black] (0,\y) -- (1.42,\y);
  \node[left,font=\small] at (-0.02,\y+0.3) {\lab};}
\draw[very thick,blue!60!black] (0,4.4) -- (0,4.9) -- (0.2,4.9) -- (0.2,4.4);
\node[right,font=\scriptsize,blue!60!black] at (0.21,4.8) {отправитель и получатель активны вместе};
\draw[very thick,orange!85!black] (0,3.1) -- (0.7,3.1) -- (0.7,3.7) -- (0.8,3.7) -- (0.8,3.1) -- (1.4,3.1);
\node[right,font=\scriptsize,orange!85!black] at (0.81,3.6) {пришёл дофамин};
% traces, each in units of its own maximum
\draw[very thick,blue!60!black] plot[domain=0:0.2,samples=20] (\x,{1.4+1.1*(1-exp(-\x))/(1-exp(-0.2))})
  plot[domain=0.2:1.4,samples=60] (\x,{1.4+1.1*exp(-(\x-0.2))});
\draw[thick,densely dashed,gray!50!black] plot[domain=0:0.2,samples=30] (\x,{1.4+1.1*(1-exp(-\x/0.05))/(1-exp(-4))})
  plot[domain=0.2:1.4,samples=80] (\x,{1.4+1.1*exp(-(\x-0.2)/0.05)});
\node[right,font=\scriptsize,blue!60!black] at (1.0,2.15) {$\tau_e=1\,\text{с}$};
\node[right,font=\scriptsize,gray!40!black] at (0.3,1.65) {$\tau_e=50\ms$};
% weights
\draw[very thick,blue!60!black] (0,0.25) -- (0.7,0.25) -- (0.8,0.8) -- (1.4,0.8) node[right,font=\scriptsize] {$\tau_e=1\,\text{с}$: вес вырос};
\draw[thick,densely dashed,gray!50!black] (0,0.2) -- (1.4,0.2) node[right,font=\scriptsize,gray!40!black] {$\tau_e=50\ms$: без изменений};
% time axis marks
\foreach \x/\l in {0/0,0.2/0.2,0.7/0.7,1.4/1.4}{\draw[gray!60!black] (\x,-0.05) -- (\x,0.05) node[below=3pt,font=\scriptsize] {\l\,с};}
\draw[<->,thick,gray!50!black] (0.2,5.35) -- node[above,font=\scriptsize,black] {задержка награды $0.5\,\text{с}$} (0.7,5.35);
\end{tikzpicture}
\caption{След по правилу~\eqref{eq:threefactor}. Совместная активность длится $0.2\,\text{с}$, дофамин приходит через полсекунды после её конца. Следы показаны в долях своего максимума. Медленный след ($\tau_e=1\,\text{с}$) к этому моменту ещё жив, быстрый ($\tau_e=50\ms$) давно растаял.}
\label{fig:trace}
\end{figure}

\begin{keyblock}
\begin{proposition}[Обучение широковещательным сигналом]
\label{prop:threefactor}
Правило~\eqref{eq:threefactor} меняет только те синапсы, которые участвовали в работе сети за последние $\tau_e$, в сторону, заданную общим сигналом $\delta$. Широковещательный сигнал вместе с локальным следом даёт адресную поправку. Задержка между действием и результатом допустима, пока она не больше $\tau_e$.
\end{proposition}
\end{keyblock}

\section{Почему это работает: шум как поиск}
\label{sec:dofagradient}

Почему правило~\eqref{eq:threefactor} ведёт к лучшему? Ответ даёт шум из главы~\ref{sec:code}. Пусть выход нейрона $y=f(u)+\xi$, где $u=\sum_jw_jx_j$, а $\xi$ --- случайное дрожание с нулевым средним и дисперсией $\sigma^2$. Награда $R$ зависит от $y$. Возьмём в качестве следа не просто $x_jy$, а его случайную часть $x_j\xi$ и в качестве третьего множителя --- ошибку $\delta=R-\bar R$, где $\bar R$ --- ожидаемая награда. При малом шуме $R\approx R\bigl(f(u)\bigr)+R'\,\xi$, поэтому
\begin{empheq}[box=\keybox]{equation}
\bigl\langle \delta\,\xi\,x_j \bigr\rangle \;=\; \sigma^2\,R'\,x_j \;=\; \frac{\sigma^2}{f'(u)}\,\frac{\partial\langle R\rangle}{\partial w_j} .
\label{eq:perturbation}
\end{empheq}
Средний шаг идёт по градиенту ожидаемой награды~\cite{Williams1992}. Никто не вычислял производных: сеть случайно отклонилась, дофамин сообщил, стало ли лучше, и участвовавшие синапсы сдвинулись в выгодную сторону (рис.~\ref{fig:perturbation}). Шум, который в главе~\ref{sec:code} мешал передавать сообщения, здесь становится поиском.

\begin{figure}[t]
\centering
\begin{tikzpicture}[>=Stealth,x=2cm,y=3cm]
\draw[->,gray!70!black] (0,0) -- (4.2,0) node[below left,font=\small,black] {выход $y$};
\draw[->,gray!70!black] (0,0) -- (0,1.2) node[left,font=\small,black] {награда $R$};
% reward hill
\draw[very thick,blue!60!black] plot[domain=0:4,samples=60] (\x,{max(0,1-((\x-3)/2.2)^2)});
\node[font=\scriptsize,blue!60!black,anchor=south] at (3,1.02) {лучший выход};
% noise around f(u)
\draw[thick,gray!60!black] plot[domain=0.6:2.4,samples=50] (\x,{0.28*exp(-((\x-1.5)/0.3)^2/2)});
\draw[densely dashed,gray!60!black] (1.5,0) -- (1.5,0.535);
\node[below,font=\small] at (1.5,0) {$f(u)$};
\node[font=\scriptsize,gray!50!black] at (1.5,-0.2) {дрожание $\xi$};
% expected reward
\draw[densely dotted,gray!60!black] (0.5,0.535) node[left,font=\scriptsize,black] {$\bar R$} -- (2.1,0.535);
% two random deviations
\fill[green!45!black] (2.1,0.833) circle (0.035);
\draw[very thick,green!45!black] (2.1,0.535) -- (2.1,0.833);
\node[font=\scriptsize,green!45!black,anchor=west,align=left] at (2.18,0.6) {$\xi>0$: лучше,\\ $\delta>0$};
\fill[red!70!black] (0.9,0.089) circle (0.035);
\draw[very thick,red!70!black] (0.9,0.535) -- (0.9,0.089);
\node[font=\scriptsize,red!70!black,anchor=east,align=right] at (0.83,0.27) {$\xi<0$: хуже,\\ $\delta<0$};
% resulting step
\draw[->,very thick,orange!85!black] (1.55,0.4) -- (1.95,0.4) node[right,font=\scriptsize,black] {шаг};
\end{tikzpicture}
\caption{Шум как поиск~\eqref{eq:perturbation}. Выход нейрона дрожит вокруг $f(u)$. Случайное отклонение вправо (зелёное) дало больше ожидаемого, $\delta>0$; влево (красное) --- меньше, $\delta<0$. В обоих случаях произведение $\delta\,\xi$ положительно, и вес сдвигается вправо, туда, где награда растёт. В среднем шаг пропорционален наклону $R'$: на вершине горки отклонения в обе стороны одинаково плохи, и шаги гасят друг друга.}
\label{fig:perturbation}
\end{figure}

\begin{keyblock}
\begin{proposition}[Спуск по градиенту без обратных проводов]
\label{prop:gradient}
Если в сети есть случайные отклонения активности, а широковещательный сигнал равен ошибке предсказания награды и в каждой обстановке в среднем равен нулю, правило~\eqref{eq:threefactor} в среднем меняет веса по градиенту ожидаемой награды. Обратные провода и отдельная фаза обучения для этого не нужны.
\end{proposition}
\end{keyblock}

Теперь понятно, почему дофамин сообщает не о награде, а об \emph{ошибке} её предсказания. Если бы третьим множителем была сама награда $R\ge0$, среднее в~\eqref{eq:perturbation} не изменилось бы, но появились бы две беды. Во-первых, к полезной части прибавилось бы слагаемое $\bar R\,\xi x_j$: в среднем ноль, но сильный шум. Во-вторых, что хуже, настоящий синапс не умеет отделять в следе $x_jy$ случайную часть от неслучайной. При данных входах $x_jf(u)$ --- одно и то же число от попытки к попытке; если $\delta$ в этой обстановке в среднем равна нулю, эта часть следа ничего не меняет, а при $\delta\ge0$ она раз за разом усиливает всё, что сеть делала, --- правильное и неправильное. Поэтому $\bar R$ должно быть ожиданием \emph{в данной обстановке}, а не средним по всей жизни. Вычислять его --- дело критика, о котором ниже.

Мы проверили это на модели. Две группы \nt{ГЛУТ}-нейронов выбирают одно из двух действий через взаимное торможение, как на рис.~\ref{fig:wta}; веса от сигнала «начали» к группам сначала равны, а шум решает, кто победит. Действие $A$ приносит сок с вероятностью $0.8$, действие $B$ --- с вероятностью $0.2$, и сок приходит через полсекунды после выбора. При $\tau_e=1\,\text{с}$ и $\delta=R-\bar R$ доля выборов $A$ за пятьсот попыток выросла с $0.65$--$0.8$ в первой сотне до $0.96$--$1.0$ в последней, а веса остались в пределах $0.85$--$1.35$. При $\tau_e=50\ms$, меньше задержки награды, сеть не научилась ничему: доля выборов $A$ осталась около половины. При $\delta=R$ без вычитания ожидания сеть тоже стала выбирать $A$, но вес победителя за те же пятьсот попыток вырос в тысячу раз и продолжал расти; а при той же $\delta=R$ и вдесятеро большей скорости обучения сеть в одном из трёх прогонов навсегда закрепила свой первый случайный выбор --- худший.

\section{Цена одного провода}

Сигнал $\delta$ один на всех, а шум, который он оценивает, --- сумма дрожаний всех $N$ нейронов, влияющих на результат. Каждый синапс видит в $\delta$ свою полезную часть и чужой шум $N-1$ соседей. Чтобы выделить своё, приходится усреднять по многим попыткам, и время обучения растёт пропорционально $N$~\cite{Werfel2005}. Обратное распространение такой платы не требует.

\begin{keyblock}
\begin{proposition}[Цена широковещания]
\label{prop:broadcastcost}
Время обучения по одному общему сигналу растёт пропорционально числу нейронов, чей шум влияет на результат. Такое обучение годится для небольших цепей, выбирающих из немногих вариантов, но не для настройки миллиардов весов.
\end{proposition}
\end{keyblock}

Мозг, по-видимому, так и делит работу: выходы \nt{ДОФА}-нейронов идут прежде всего к областям, которые выбирают действия (глава~\ref{sec:neurontypes}), а огромная \nt{ГЛУТ}-сеть коры, где весов подавляющее большинство, учится в основном местными правилами, без внешнего учителя. Раньше их сводили к правилам Хебба~\eqref{eq:oja}. Сейчас всё больше данных за то, что у коры есть собственная цель --- \emph{предсказывать} свой следующий вход, и тогда ошибка вычисляется на месте, как разница между предсказанным и пришедшим~\cite{Keller2018}: учитель встроен в саму сеть. Окна пластичности длиной в секунды, где третьим сигналом служит сильное возбуждение самого получателя~\cite{Bittner2017}, показывают, что местный сигнал ошибки у синапсов действительно бывает. Насколько это общее правило коры, --- гипотеза.

\section{Откуда берётся ошибка: критик в непрерывном времени}

Остаётся вопрос, как сами \nt{ДОФА}-нейроны вычисляют $\delta$. В программах обучения с подкреплением ошибку предсказания считают шагами $t,t+1,\dots$ В организме шагов нет, поэтому запишем её в непрерывном времени~\cite{Doya2000}. Пусть $r(t)$ --- поток награды, а $V(t)$ --- ожидание всей будущей награды, в которой более далёкая ценится меньше:
\begin{equation}
V(t) \;=\; \int_t^{\infty} e^{-(s-t)/\tau_\gamma}\,r(s)\,\dd s .
\end{equation}
Продифференцировав, получим $\dd V/\dd t=V/\tau_\gamma-r$. Невязка этого равенства и есть ошибка предсказания:
\begin{empheq}[box=\keybox]{equation}
\delta(t) \;=\; r(t) \;+\; \frac{\dd V}{\dd t} \;-\; \frac{V}{\tau_\gamma} .
\label{eq:ctd}
\end{empheq}
Постоянная $\tau_\gamma$ --- горизонт планирования, непрерывный аналог коэффициента $\gamma$; по гипотезе, его задаёт \nt{СЕРО} (раздел~\ref{sec:horizon}), и тогда~\eqref{eq:patience} --- правило, сколько стоит ждать. Проверим три случая (рис.~\ref{fig:ctdcases}). Загорелась лампочка, обещающая сок: $V$ скачком выросло, $\dd V/\dd t$ большое, всплеск. Пришёл обещанный сок: $r>0$, но ожидание в тот же момент исчерпалось, $\dd V/\dd t\approx-r$, и всплеска нет. Сок не пришёл: ожидание исчезло без награды, $\dd V/\dd t<0$, провал.

\begin{figure}[t]
\centering
\begin{tikzpicture}[>=Stealth,x=1.15cm,y=1cm]
\foreach \col/\ttl/\lampon/\juice in {0/{неожиданный сок}/0/1, 1/{обещанный сок}/1/1, 2/{сок не пришёл}/1/0}{
 \begin{scope}[xshift=\col*4.6cm]
  \node[font=\small] at (1.5,3.55) {\ttl};
  \foreach \yy in {2.4,1.2,0}{\draw[gray!60!black] (0,\yy) -- (3.1,\yy);}
  % reward r
  \ifnum\juice=1
    \fill[green!45!black] (1.95,2.4) rectangle (2.05,2.85);
    \pic[scale=0.55] at (2.0,3.1) {juice};
  \else
    \pic[scale=0.55] at (2.0,3.1) {nojuice};
  \fi
  % expectation V and error delta
  \ifnum\lampon=1
    \pic[scale=0.55] at (1.0,3.1) {lamp};
    \draw[very thick,blue!60!black] (0,1.2) -- (1,1.2) -- (1,1.45)
      plot[domain=1:2,samples=20] (\x,{1.2+0.25*exp((\x-1)/1.5)}) -- (2,{1.2+0.25*exp(1/1.5)}) -- (2,1.2) -- (3.1,1.2);
    \draw[very thick,red!70!black] (0,0) -- (0.95,0) -- (0.95,0.5) -- (1.05,0.5) -- (1.05,0) -- (3.1,0);
    \ifnum\juice=0
      \draw[very thick,red!70!black] (1.95,0) -- (1.95,-0.45) -- (2.05,-0.45) -- (2.05,0);
      \node[font=\scriptsize,red!70!black,anchor=west] at (2.1,-0.35) {провал};
    \else
      \node[font=\scriptsize,gray!50!black,anchor=north,align=center] at (2.0,-0.08) {$r$ и спад $V$\\ гасят друг друга};
    \fi
  \else
    \draw[very thick,blue!60!black] (0,1.2) -- (3.1,1.2);
    \draw[very thick,red!70!black] (0,0) -- (1.95,0) -- (1.95,0.5) -- (2.05,0.5) -- (2.05,0) -- (3.1,0);
  \fi
  \ifnum\col=0
    \node[font=\small,anchor=east] at (-0.1,2.55) {$r$};
    \node[font=\small,anchor=east] at (-0.1,1.35) {$V$};
    \node[font=\small,anchor=east] at (-0.1,0.1) {$\delta$};
  \fi
 \end{scope}}
\end{tikzpicture}
\caption{Три случая для ошибки~\eqref{eq:ctd}; сверху вниз --- награда $r$, ожидание $V$ и ошибка $\delta$. Слева ожидания нет, и сок целиком --- неожиданность. В середине лампочка скачком поднимает $V$: это скачок $\dd V/\dd t$, всплеск $\delta$ приходится на лампочку. Дальше $V$ растёт ровно так, как требует горизонт, и $\delta=0$; в момент сока награда $r$ и спад $V$ гасят друг друга. Справа $V$ падает без награды --- провал. Это те же всплеск, перенос и провал, что на рис.~\ref{fig:rpe}, но теперь видно, откуда они берутся.}
\label{fig:ctdcases}
\end{figure}

Что нужно для~\eqref{eq:ctd} из схем, которые у нас уже есть? Три вещи.

\emph{Сама оценка $V$.} Её выдаёт группа \nt{ГЛУТ}-нейронов --- критик, --- веса которой учатся тем же правилом~\eqref{eq:threefactor} с тем же $\delta$: если $\delta>0$, ожидание было занижено, и связи, активные перед этим, усиливаются.

\emph{Производная $\dd V/\dd t$.} Её вычисляет любая схема, отвечающая на изменения, а не на уровень: прямое возбуждение вместе с задержанным торможением через \nt{ГАМК}-посредника, как в детекторе направления главы~\ref{sec:gaba}, или истощающийся синапс~\eqref{eq:depression} из главы~\ref{sec:code}.

\emph{Знак в одном проводе.} Ошибка бывает и положительной, и отрицательной, а частота импульсов --- только положительной. Решение то же, что у \nt{СЕРО}-нейрона в главе~\ref{sec:serocomputer}: \nt{ДОФА}-нейрон --- ритмоводитель. Его ровные несколько импульсов в секунду означают $\delta=0$, пачка --- $\delta>0$, пауза --- $\delta<0$. Отрицательной ошибке остаётся меньше места: частоту нельзя опустить ниже нуля. У настоящих \nt{ДОФА}-нейронов провалы действительно мельче всплесков~\cite{Bayer2005}.

Что дофамин ведёт себя как $\delta$, измерено. Как именно мозг вычисляет $\dd V/\dd t$, --- гипотеза.

\section{Актёр и критик}

Соберём всё вместе (рис.~\ref{fig:actorcritic}). Нейроны обстановки возбуждают группы действий, выбирающие победителя через \nt{ГАМК} (утверждение~\ref{prop:wta}), --- это \emph{актёр} --- и критика, выдающего $V$~\cite{SuttonBarto2018}. \nt{ДОФА}-нейрон получает награду $r$ и $V$, вычисляет~\eqref{eq:ctd} и рассылает $\delta$ всем синапсам актёра и критика.

\begin{figure}[t]
\centering
\begin{tikzpicture}[>=Stealth,
  nrn/.style={circle,draw,very thick,fill=blue!6,minimum size=0.9cm,inner sep=0pt,font=\small},
  gab/.style={circle,draw,very thick,fill=red!10,minimum size=0.6cm,inner sep=0pt},
  dofa/.style={circle,draw,very thick,double,double distance=1.2pt,fill=orange!15,minimum size=1.35cm,inner sep=0pt,font=\scriptsize},
  inh/.style={-{Bar[width=7pt]},thick,red!70!black},
  bc/.style={->,thick,densely dashed,orange!85!black}]
\node[nrn] (S1) at (0,2.4) {$x_1$};
\node[nrn] (S2) at (0,1.2) {$x_2$};
\node[nrn] (S3) at (0,0) {$x_3$};
\node[left=0.15cm,align=right,font=\small] at (-0.5,1.2) {обстановка};
\node[nrn] (A) at (4,2.6) {$A$};
\node[nrn] (B) at (4,1.0) {$B$};
\node[gab] (G) at (5.2,1.8) {};
\draw[->,thick] (A) -- (G); \draw[->,thick] (B) -- (G);
\draw[inh] (G) to[bend left=35] (A);
\draw[inh] (G) to[bend right=35] (B);
\node[above,font=\small] at (4.6,3.05) {актёр: выбор действия};
\node[nrn] (V) at (4,-1.1) {$V$};
\node[below,font=\small] at (4,-1.6) {критик};
\foreach \s in {S1,S2,S3}{\foreach \t in {A,B,V}{\draw[->,gray!60!black] (\s) -- (\t);}}
\node[dofa] (D) at (8.0,-1.1) {\nt{ДОФА}};
\draw[->,thick] (V) -- node[above,font=\scriptsize] {$V,\ \dd V/\dd t$} (D);
\draw[->,thick,green!45!black] (10.0,-1.1) node[right,black,font=\small] {награда $r$} -- (D);
\pic[scale=1.1] at (10.6,-0.5) {juice};
\draw[bc] (D.north) -- (8.0,4.0) -- node[above,font=\small,black] {$\delta$ --- всем синапсам актёра и критика} (2.2,4.0) -- (2.2,3.0);
\draw[bc] (D.south) -- (8.0,-2.4) -- (2.2,-2.4) -- (2.2,-0.6);
\end{tikzpicture}
\caption{Сеть, которая учится выбирать действия. Серые стрелки --- \nt{ГЛУТ}-синапсы с меняющимися весами; красный кружок --- \nt{ГАМК}-нейрон; двойной кружок --- \nt{ДОФА}-ритмоводитель, вычисляющий $\delta$ по~\eqref{eq:ctd}; штриховые стрелки --- широковещательная рассылка $\delta$.}
\label{fig:actorcritic}
\end{figure}

Знак действия медиатора выбирает получатель (утверждение~\ref{prop:receptor}), и в области выбора действий часть нейронов несёт дофаминовые рецепторы одного семейства, часть --- другого. В опытах на срезах дофамин вместе с совместной активностью у первых усиливает синапсы, у вторых --- ослабляет~\cite{Shen2008}; в модели это значит, что у вторых знак $\delta$ в~\eqref{eq:threefactor} перевёрнут. Первые учатся тому, что \emph{стоит делать}, вторые --- тому, чего \emph{делать не стоит}.

\section{Итог}

\begin{table}[t]
\centering
\footnotesize
\setlength{\tabcolsep}{4pt}
\renewcommand{\arraystretch}{1.15}
\begin{tabular}{>{\raggedright}p{2.6cm}>{\raggedright}p{2.7cm}>{\raggedright}p{3.1cm}>{\raggedright\arraybackslash}p{5.0cm}}
\toprule
Правило & Что знает синапс & Чему учится сеть & Что известно \\
\midrule
Хебб по частотам~\eqref{eq:oja} & $x$, $y$, свой вес & сжатию входов: главным направлениям & усиление при совместной активности измерено; механизм ограничения весов --- предмет исследований \\
Хебб по времени & порядок импульсов $x$ и $y$ в пределах $\sim20\ms$ & причинным связям, последовательностям & измерено на \nt{ГЛУТ}-синапсах \\
три множителя~\eqref{eq:threefactor} & $x$, $y$, след $e$, общий сигнал $\delta$ & действиям, которые приносят награду & влияние дофамина в течение секунд после активности измерено; как главный механизм обучения выбору --- хорошо обоснованная гипотеза \\
критик~\eqref{eq:ctd} & то же & ожиданию награды $V$ & сходство дофамина с $\delta$ измерено; схема, вычисляющая $\dd V/\dd t$, --- гипотеза \\
обратное распространение ошибки & персональную поправку для каждого веса & всему, но нужны обратные провода и такт & в мозге не найдено в этой форме; ищут приближения к нему \\
\bottomrule
\end{tabular}
\caption{Правила обучения в сети из \nt{ГЛУТ}-, \nt{ГАМК}- и \nt{ДОФА}-нейронов.}
\label{tab:learning}
\end{table}

Правила обучения сведены в табл.~\ref{tab:learning}. \nt{ГЛУТ} несёт данные, \nt{ГАМК} управляет потоком, а \nt{ДОФА} решает, какие изменения в сети оставить.

\section{Задачи}

\begin{exercise}[\lvlA]
\label{ex:06:1}
\emph{Сколько живёт след.} В опыте рис.~\ref{fig:trace} совместная активность $xy=1$ длится $0.2\,\text{с}$, начиная с $e=0$, а дофамин приходит через $0.5\,\text{с}$ после её конца.
(а)~Найдите по~\eqref{eq:threefactor} значение следа в момент прихода дофамина при $\tau_e=1\,\text{с}$ и при $\tau_e=50\ms$ и их отношение.
(б)~При какой наибольшей задержке след при $\tau_e=1\,\text{с}$ ещё не меньше $10\,\%$ своего значения в конце совместной активности?
(в)~При фиксированных длительности активности $T_a$ и задержке $d$ найдите $\tau_e$, при котором след в момент дофамина наибольший. Вычислите его для $T_a=0.2\,\text{с}$ и для $T_a=1\,\text{с}$ при $d=0.5\,\text{с}$. Почему слишком длинный след тоже плох?
\end{exercise}

\begin{exercise}[\lvlA]
\label{ex:06:2}
\emph{Дрейф правила Хебба по времени.} Отправитель и получатель выдают независимые пуассоновские потоки с частотами $\nu_x$ и $\nu_y$; каждая пара импульсов меняет вес по окну рис.~\ref{fig:stdp}.
(а)~Докажите, что средняя скорость изменения веса равна $\nu_x\nu_y\,(A_+\tau_+-A_-\tau_-)$.
(б)~Вычислите её при $\tau_\pm=20\ms$, $A_-=0.6A_+$, $\nu_x=\nu_y=10$ импульсов в секунду. На сколько $A_+$ вырастет вес за час совершенно случайной активности?
(в)~При каком $\tau_-$ (при $A_-=0.6A_+$) дрейф исчезает? Почему это условие нужно выполнять точно, если правило должно учить только причинным связям?
\end{exercise}

\begin{exercise}[\lvlB]
\label{ex:06:3}
\emph{Доказательство утверждения~\ref{prop:oja} до конца.} Усреднённое правило: $\dd\mathbf w/\dd t=\eta\bigl(C\mathbf w-(\mathbf w^{\mathsf T}C\mathbf w)\mathbf w\bigr)$.
(а)~Покажите, что $\dd\|\mathbf w\|^2/\dd t=2\eta\,(\mathbf w^{\mathsf T}C\mathbf w)(1-\|\mathbf w\|^2)$, и выведите, что длина вектора весов стремится к единице.
(б)~Линеаризуйте правило в точке $\mathbf e_k$. Покажите, что отклонение вдоль $\mathbf e_j$, $j\ne k$, растёт как $e^{\eta(\lambda_j-\lambda_k)t}$, а отклонение длины затухает со скоростью $2\eta\lambda_k$.
(в)~Для $C=\begin{pmatrix}2&1\\1&2\end{pmatrix}$ найдите предел $\mathbf w$ и самую медленную постоянную времени сходимости.
(г)~Частоты импульсов положительны, и у входов есть среднее: $\mathbf x=\mathbf m+\mathbf n$, $\langle\mathbf n\rangle=0$, $\langle\mathbf n\mathbf n^{\mathsf T}\rangle=\Sigma$. Пусть $\mathbf m=(\mu,0)$, $\Sigma=\operatorname{diag}(s_1^2,s_2^2)$. При каком условии нейрон выучит направление $\mathbf e_1$? Проверьте его при $\mu=1$, $s_1=0.1$, $s_2=0.5$ и уточните формулировку утверждения~\ref{prop:oja}.
\end{exercise}

\begin{exercise}[\lvlA]
\label{ex:06:4}
\emph{Непрерывный критик как предел шагов.}
(а)~Разбейте время на шаги $\Delta$, награду за шаг возьмите $r(t)\Delta$, а $\gamma=e^{-\Delta/\tau_\gamma}$. Покажите, что $\delta_t/\Delta$ из~\eqref{eq:rpe} при $\Delta\to0$ переходит в~\eqref{eq:ctd}.
(б)~Лампочка зажигается в непредсказуемый момент $t_c$, сок размера $R$ (короткий импульс площади $R$) приходит в момент $t_c+L$. Для выученного ожидания найдите $V(t)$ и $\delta(t)$ на всём опыте, включая случай, когда сок не пришёл. Покажите, что между лампочкой и соком $\delta\equiv0$.
(в)~Вычислите площадь всплеска на лампочку в долях $R$ при $L=1\,\text{с}$ для $\tau_\gamma=2\,\text{с}$ и $\tau_\gamma=4\,\text{с}$. Какое проверяемое предсказание отсюда следует для гипотезы раздела~\ref{sec:horizon}?
\end{exercise}

\begin{exercise}[\lvlB]
\label{ex:06:5}
\emph{Откуда берётся цена широковещания.} Пусть $N$ нейронов дрожат независимо, $\xi_i\sim\mathcal N(0,\sigma^2)$, и при малом шуме $\delta=\sum_ia_i\xi_i$, где $a_i=\partial R/\partial y_i$. Синапс $j$-го нейрона оценивает градиент величиной $\hat g_j=\delta\,\xi_jx_j$.
(а)~Найдите $\langle\hat g_j\rangle$ и $\operatorname{Var}\hat g_j$.
(б)~При одинаковых $a_i$ найдите отношение сигнала к шуму одной попытки и число попыток, нужное для заданной точности. Выведите отсюда утверждение~\ref{prop:broadcastcost}.
(в)~Пусть вместо ошибки рассылается сама награда, $\delta=\bar R+\sum_ia_i\xi_i$. Во сколько раз вырастет дисперсия? Сравните с рассуждением после утверждения~\ref{prop:gradient}.
(г)~Пусть один нейрон ($N=1$) учится задаче за $100$ попыток, а одна попытка занимает $5\,\text{с}$. Оцените время обучения при $N=10^4$ и при $N=10^{10}$.
\end{exercise}

\begin{exercise}[\lvlB]
\label{ex:06:6}
\emph{Проектирование: схема, вычисляющая $\delta$.} \nt{ДОФА}-ритмоводитель с фоновой частотой $\rho_0$ получает награду $r$, оценку $V$ прямым возбуждающим путём с весом $k_1$ и её сглаженную копию $\bar V$, $\tau_d\,\dd\bar V/\dd t=-\bar V+V$, тормозящим путём с весом $k_2$. Его частота $\rho=\rho_0+\lambda\,\hat\delta$, $\hat\delta=r+k_1V-k_2\bar V$, ограничена снизу нулём.
(а)~Выберите $k_1$, $k_2$ так, чтобы для медленных $V$ было $\hat\delta\approx$~\eqref{eq:ctd}.
(б)~Выученное ожидание растёт как $V=V_0e^{t/\tau_\gamma}$, и истинная $\delta=0$. Найдите точно, что выдаёт схема. Какое $\tau_d$ нужно, чтобы ложная ошибка не превышала $5\,\%$ от $V/\tau_\gamma$ при $\tau_\gamma=2\,\text{с}$?
(в)~$V$ скачком выросло на $\Delta V$ (лампочка). Найдите $\hat\delta(t)$ и покажите, что площадь всплеска равна $\Delta V$ при любом $\tau_d$.
(г)~Годится ли в качестве сглаживающего звена истощающийся синапс~\eqref{eq:depression} с $U=0.5$, $\tau_{\text{вос}}=0.8\,\text{с}$ при частоте $20$ импульсов в секунду? Какой будет ложная ошибка?
(д)~Пусть $\rho_0=5$, а наибольшая частота $30$ импульсов в секунду. Каков диапазон представимых $\hat\delta$ в каждую сторону? Если $\hat\delta$ гауссова с нулевым средним и $\lambda\sigma_\delta=\rho_0$, найдите среднее переданного сигнала после обрезки нулём. Чем опасно это смещение для правила~\eqref{eq:threefactor}?
\end{exercise}

\begin{exercise}[\lvlB]
\label{ex:06:7}
\emph{Моделирование: предельная задержка награды.} Работайте с копией \texttt{models/\allowbreak dofa\_learning.py}.
(а)~Добавьте в \texttt{run()} параметр задержки награды $d$ и сделайте длину попытки равной $T_{\text{cue}}+d+T_{\text{rew}}+0.4\,\text{с}$ (при $d=0.5\,\text{с}$ модель должна совпасть с исходной).
(б)~При $d=0.5\,\text{с}$, $\eta=0.05$, $500$ попытках и шести зёрнах постройте долю выборов $A$ в последней сотне попыток в зависимости от $\tau_e\in\{0.05,0.1,0.2,0.5,1,2,4\}\,\text{с}$.
(в)~При $\tau_e=1\,\text{с}$ постройте ту же долю в зависимости от $d\in\{0.25,0.5,1,2,3\}\,\text{с}$.
(г)~Вычислите для каждого прогона множитель $F=(1-e^{-T_{\text{cue}}/\tau_e})\,e^{-d/\tau_e}$ --- долю следа, доживающую до награды (задача~\ref{ex:06:1}), и постройте все точки как функцию $F$. Проверьте утверждение~\ref{prop:threefactor}: где проходит граница обучения и почему доля выборов $A$ убывает и при слишком большом $\tau_e$?
\end{exercise}

\begin{exercise}[\lvlC]
\label{ex:06:8}
\emph{Можно ли обойти цену широковещания.} Модель: $N$ нейронов, выход $y_i=w_i+\xi_i$, $\xi_i\sim\mathcal N(0,\sigma^2)$, награда $R=-\sum_i(y_i-w_i^*)^2$, ошибка $\delta=R-\langle R\rangle$ (ожидание при текущих весах), правило $\Delta w_i=\eta\,\delta\,\xi_i$ после каждой попытки. Пусть $E=\sum_i(w_i-w_i^*)^2$.
(а)~Докажите, что $\langle\Delta E\rangle=-4\eta\sigma^2E+\eta^2\bigl[4\sigma^4(N+2)E+\sigma^6N(2N+8)\bigr]$. Найдите лучшее $\eta$ и покажите, что вдали от шумового дна $E$ убывает в $e$ раз за $N+2$ попытки.
(б)~Пусть в каждой попытке дрожат только $m$ случайно выбранных нейронов из $N$. Покажите, что постоянная времени равна $N(m+2)/m$. Помогает ли разреженный поиск?
(в)~Проверьте (а) и (б) моделью: $\sigma=0.1$, начальные $w_i-w_i^*=1$, $\eta$ --- половина от $1/(\sigma^2(m+2))$; измерьте число попыток до $E<0.1E_0$ при $N=4,16,64,256$ и $m=1,4,N$ и сравните с теорией.
(г)~Открытая часть. Линейный рост с $N$ неизбежен для любого правила, у которого на входе одно число на попытку? Какие изменения условий его ломают? Рассмотрите хотя бы: несколько проводов ошибки (глава~\ref{sec:dofaalt}), местную ошибку предсказания входа, шум в весах вместо шума в нейронах. Оцените, какой из вариантов позволил бы учить $10^{10}$ весов за разумное время.
\end{exercise}

\chapter{Альтернативные современные теории обучения на основе дофамина}
\chaptermark{Другие теории \nt{DOPA}}
\label{sec:dofaalt}
\begin{chaptergoals}
\item как \nt{DOPA}-нейроны с асимметричным ответом на ошибки записывают всё распределение награды;
\item как часть \nt{DOPA}-нейронов сообщает о важности или опасности, а не о знаке ошибки;
\item как один провод несёт быструю ошибку для обучения и медленный уровень, задающий темп действий;
\item почему сигнал \nt{DOPA} --- жгут, а не одно число, и какая теория оспаривает саму ошибку.
\end{chaptergoals}

\section{Что на самом деле идёт по проводу}

В главе~\ref{sec:dofa} \nt{DOPA}-система была у нас одним проводом, по которому идёт одно число --- ошибка предсказания награды $\delta$~\eqref{eq:ctd}. Эта картина объясняет многое, но чем точнее измеряют дофамин, тем больше находят того, что в неё не укладывается. Одни нейроны отвечают на неприятное, как на награду; концентрация дофамина растёт, пока животное бежит к цели, хотя ничего неожиданного не происходит; разные \nt{DOPA}-нейроны ведут себя по-разному.

Для инженера все эти находки --- вопрос об одном: \emph{какой формат у сигнала на широковещательной шине?} Одно число или вектор? Один сигнал на проводе или несколько, разнесённых по частоте? Говорит ли он о будущем, как $\delta$, или о прошлом? Ниже --- шесть современных ответов; для каждого мы отделим измеренное от предполагаемого.

\section{Не одно число, а распределение}

Ошибка~\eqref{eq:ctd} учит критика \emph{среднему} ожиданию награды. Но верные полкапли сока и капля с вероятностью одна вторая имеют одинаковое среднее. Чтобы их различать, нужно всё распределение награды.

Возьмём простейшую задачу: после предвестника приходит награда случайного размера $r$. Пусть \nt{DOPA}-нейронов много, и каждый, $i$-й, обслуживает свою оценку $V_i$, так что в момент награды его ошибка равна $\delta_i=r-V_i$. И пусть положительную ошибку каждый учитывает с весом $\alpha_i^+$, отрицательную --- с весом $\alpha_i^-$. После каждой награды оценка сдвигается на
\begin{equation}
\Delta V_i \;=\; \eta\,\bigl(\alpha_i^+\,[\delta_i]_+ \;-\; \alpha_i^-\,[-\delta_i]_+\bigr).
\label{eq:asymmetric}
\end{equation}
Оценка перестаёт меняться, когда в среднем положительные поправки уравновешивают отрицательные:
\begin{empheq}[box=\keybox]{equation}
\alpha_i^+\,\bigl\langle[r-V_i]_+\bigr\rangle \;=\; \alpha_i^-\,\bigl\langle[V_i-r]_+\bigr\rangle ,
\qquad
\kappa_i \;=\; \frac{\alpha_i^+}{\alpha_i^++\alpha_i^-} .
\label{eq:expectile}
\end{empheq}
При $\kappa_i=1/2$ это обычное среднее. При $\kappa_i>1/2$ нейрон --- «оптимист»: его оценка сдвинута к верхнему краю распределения, при $\kappa_i<1/2$ --- «пессимист».

Пример: сок приходит с вероятностью $1/2$, размер капли $1$. Тогда~\eqref{eq:expectile} даёт $\kappa_i\cdot\tfrac12(1-V_i)=(1-\kappa_i)\cdot\tfrac12V_i$, то есть $V_i=\kappa_i$ (рис.~\ref{fig:expectiles}). Набор нейронов с разными $\kappa_i$ записывает распределение награды так же, как набор квантилей записывает распределение в статистике.

\begin{figure}[t]
\centering
\begin{tikzpicture}[>=Stealth,x=7cm,y=3cm]
\draw[->,gray!70!black] (-0.08,0) -- (1.15,0) node[right,black] {\small награда};
\draw[->,gray!70!black] (-0.08,0) -- (-0.08,0.75) node[left,black] {\small вероятность};
\fill[blue!25] (-0.03,0) rectangle (0.03,0.5);
\fill[blue!25] (0.97,0) rectangle (1.03,0.5);
\node[below,font=\small] at (0,0) {$0$};
\node[below,font=\small] at (1,0) {$1$};
\node[left,font=\scriptsize] at (-0.08,0.5) {$1/2$};
\foreach \k/\c/\lab/\h in {0.2/blue!60!black/{пессимист, $\kappa=0.2$}/0.62, 0.5/gray!50!black/{среднее, $\kappa=0.5$}/0.42, 0.8/red!70!black/{оптимист, $\kappa=0.8$}/0.22}{
  \draw[very thick,\c] (\k,0) -- (\k,\h);
  \node[above,font=\scriptsize,\c] at (\k,\h) {\lab};}
\end{tikzpicture}
\caption{Распределение награды записано набором \nt{DOPA}-нейронов. Награда $0$ или $1$ с вероятностью $1/2$ (синие столбики); нейрон с долей $\kappa$ приходит к оценке $V=\kappa$~\eqref{eq:expectile}.}
\label{fig:expectiles}
\end{figure}

Что измерено: у разных \nt{DOPA}-нейронов разная «точка переворота» --- размер награды, при котором ответ сменяется со всплеска на провал, --- и нейроны с большей асимметрией ответов на положительные и отрицательные ошибки переворачиваются при больших наградах, как и требует~\eqref{eq:expectile}~\cite{Dabney2020}. Из ответов группы нейронов удаётся восстановить форму распределения. Что это главная функция разброса, а не побочный эффект, --- гипотеза.

\begin{keyblock}
\begin{proposition}[Распределение из асимметрии]
\label{prop:expectile}
Если \nt{DOPA}-нейроны различаются долей $\kappa_i$, с которой они учитывают положительные ошибки, их оценки сходятся к разным точкам~\eqref{eq:expectile} распределения награды. Группа нейронов передаёт не среднее, а распределение, и тем самым --- риск.
\end{proposition}
\end{keyblock}

\section{Не только ошибка, но и важность}

По формуле ошибки~\eqref{eq:ctd} неприятное событие --- удар, горький вкус --- должно вызывать провал: вышло хуже, чем ожидалось. Часть \nt{DOPA}-нейронов так и делает, а другая отвечает на неприятное \emph{всплеском}, как на награду~\cite{Matsumoto2009}. Эти нейроны сообщают не знак ошибки, а то, что случилось что-то \emph{важное}: неожиданное, сильное, требующее внимания~\cite{BrombergMartin2010}.

Инженеру удобно думать об этом как о двух сигналах. Первый --- знаковая ошибка $\delta$: куда менять веса. Второй --- её модуль $|\delta|$ или новизна события: \emph{насколько сильно} их менять, то есть скорость обучения; после неожиданного события учиться надо быстрее. По смыслу он близок к норадреналиновому коэффициенту усиления из главы~\ref{sec:neurontypes}.

Есть и третий вариант: отдельный провод для отдельной оси. \nt{DOPA}-нейроны, идущие к одной из областей выбора действий, отвечают не на награду, а на новизну и силу раздражителя и нужны, чтобы животное научилось избегать опасного~\cite{Menegas2018}: по проводу идёт не «лучше или хуже», а «опасно или нет».

Что измерено: разные группы \nt{DOPA}-нейронов отвечают на неприятное и на новое по-разному, и разные выходы учат разному. Как именно мозг пользуется сигналом важности --- как скоростью обучения, как сигналом внимания или как отдельной ошибкой по оси опасности, --- обсуждается.

\section{Не только учить, но и подгонять}

У дофамина есть и немедленное действие: чем его больше, тем охотнее и быстрее животное действует. Как совместить это с обучением?

Инженерный ответ --- \emph{разделение по частоте}. Быстрые всплески и провалы длительностью в сотни миллисекунд несут $\delta$ и учат. Медленный уровень, меняющийся за минуты, несёт другую величину --- среднюю награду в единицу времени $\bar R$~\cite{Niv2007}. Пусть действие, выполненное за время $\tau_a$, требует усилия $C/\tau_a$: чем быстрее, тем дороже. Зато каждая секунда промедления стоит $\bar R$ --- столько награды можно было бы за неё получить. Суммарная цена $C/\tau_a+\bar R\,\tau_a$ минимальна при
\begin{empheq}[box=\keybox]{equation}
\tau_a^{*} \;=\; \sqrt{\frac{C}{\bar R}} .
\label{eq:vigor}
\end{empheq}
Чем богаче обстановка, тем дороже время и тем быстрее выгодно действовать: если $\bar R$ удвоилась, время действия уменьшается в $\sqrt2\approx1.4$ раза. Заметьте, что $\bar R$ --- та самая ожидаемая награда, которую вычитали в главе~\ref{sec:dofa}.

Выброс дофамина в месте назначения может меняться и без изменения частоты импульсов: его подстраивают местные сигналы у самих выходов~\cite{Mohebi2019}. Но медленные составляющие есть и в самих импульсах: в опытах следующего раздела плавный рост при приближении к награде виден и в частоте \nt{DOPA}-нейронов~\cite{Kim2020}. Значит, медленный уровень складывается из двух источников --- частоты импульсов и местной регулировки выброса, --- и какой из них главный, по-видимому, зависит от выхода и от задачи.

\begin{keyblock}
\begin{proposition}[Два сигнала на одном проводе]
\label{prop:multiplex}
\nt{DOPA}-система передаёт по меньшей мере две величины, разнесённые по времени: быструю ошибку $\delta$, которая учит, и медленный уровень, который задаёт цену времени~\eqref{eq:vigor} и тем самым скорость действий. Измерено, что быстрая и медленная составляющие ведут себя по-разному; что медленная равна именно $\bar R$, --- гипотеза.
\end{proposition}
\end{keyblock}

\section{Рост при приближении}

Если животное бежит по коридору к далёкой награде, концентрация дофамина в области выбора действий плавно растёт в течение секунд, пока цель приближается~\cite{Hamid2016}. По главе~\ref{sec:dofa} этого быть не должно: всё идёт по плану, и $\delta$ должна быть нулём.

Первое объяснение: этот рост --- не $\delta$, а само ожидание $V$, сигнал «цель близко, стоит постараться» из предыдущего раздела~\cite{Hamid2016}. Второе объяснение сохраняет $\delta$~\cite{Kim2020}. Если ожидание верно, то при горизонте $\tau_\gamma$ на пути к награде $R$ оно растёт как $V=Re^{-(T-t)/\tau_\gamma}$, и по формуле ошибки~\eqref{eq:ctd} $\dd V/\dd t-V/\tau_\gamma=0$. Но пусть издалека ожидание занижено, а по мере приближения оценка уточняется. Тогда ожидание растёт круче, скажем как $V=Re^{-(T-t)/\tau_v}$ с $\tau_v<\tau_\gamma$, и
\begin{empheq}[box=\keybox]{equation}
\delta(t) \;=\; \frac{\dd V}{\dd t}-\frac{V}{\tau_\gamma} \;=\; V\Bigl(\frac{1}{\tau_v}-\frac{1}{\tau_\gamma}\Bigr) \;>\; 0 .
\label{eq:ramp}
\end{empheq}
Ошибка положительна и растёт вместе с $V$ --- тот самый плавный рост (рис.~\ref{fig:ramp}). При $\tau_\gamma=4\,\text{с}$ и $\tau_v=1\,\text{с}$ получается $\delta=0.75\,V$ в секунду. Эту версию проверяли в виртуальном коридоре, мгновенно перенося животное ближе к цели: дофамин отвечал всплеском, как и положено производной, а не плавному уровню~\cite{Kim2020}.

\begin{figure}[t]
\centering
\begin{tikzpicture}[>=Stealth,x=1.6cm,y=2.6cm]
\draw[->,gray!70!black] (-4.2,0) -- (0.5,0) node[below left,black] {\small $t$};
\draw[->,gray!70!black] (-4.2,0) -- (-4.2,1.2);
\foreach \x/\l in {-4/{$T-4$},-2/{$T-2$},0/{$T$}}{\draw[gray!60!black] (\x,-0.03) -- (\x,0.03); \node[below,font=\scriptsize] at (\x,0) {\l\,с};}
\draw[thick,densely dashed,gray!60!black] plot[domain=-4:0,samples=40] (\x,{exp(\x/4)});
\draw[very thick,blue!60!black] plot[domain=-4:0,samples=60] (\x,{exp(\x)});
\draw[very thick,red!70!black] plot[domain=-4:0,samples=60] (\x,{0.75*exp(\x)});
\node[anchor=south east,font=\small,gray!50!black] at (-1.3,0.77) {$V$ при $\tau_v=\tau_\gamma$: $\delta=0$};
\node[right,font=\small,blue!60!black] at (0.05,1.0) {$V$ при $\tau_v=1\,\text{с}$};
\node[right,font=\small,red!70!black] at (0.05,0.72) {$\delta$ по~\eqref{eq:ramp}};
\end{tikzpicture}
\caption{Рост дофамина при приближении к награде как ошибка предсказания, $\tau_\gamma=4\,\text{с}$. Штриховая линия --- ожидание, растущее ровно как требует горизонт ($\delta=0$); синяя --- ожидание, растущее круче; красная --- ошибка~\eqref{eq:ramp}.}
\label{fig:ramp}
\end{figure}

Что измерено: плавный рост есть --- и в концентрации дофамина, и в частоте импульсов \nt{DOPA}-нейронов~\cite{Kim2020}, --- и мгновенные скачки обстановки вызывают скачки дофамина. Какое из двух объяснений верно --- или оба верны на разных выходах, --- обсуждается.

\section{Взгляд назад}

Все теории выше смотрят \emph{вперёд}. Одна из современных теорий переворачивает направление~\cite{Jeong2022}. Пусть мозг учится не предсказывать награду по предвестнику, а, получив награду, \emph{оглядываться}: какое событие чаще других ей предшествовало? Мера такой связи --- разность двух вероятностей:
\begin{empheq}[box=\keybox]{equation}
M \;=\; P(\text{предвестник перед наградой}) \;-\; P(\text{предвестник в случайном окне}) .
\label{eq:retro}
\end{empheq}
Если предвестник часто бывает и без награды, второе слагаемое велико, и связь слаба. В этой теории дофамин сообщает не ошибку предсказания, а то, насколько данное событие --- вероятная \emph{причина} награды.

Механизм оглядки требует того же, что и правило трёх множителей главы~\ref{sec:dofa}: каждое событие должно оставлять тающий след, чтобы в момент награды можно было прочитать, что было перед ней. Разница не в синапсе, а в том, что передаёт \nt{DOPA}.

В опытах, где теория~\eqref{eq:retro} и ошибка~\eqref{eq:ctd} предсказывают разное, выброс дофамина в ряде случаев следовал первой~\cite{Jeong2022}. Но в других опытах ответ дофамина при обучении постепенно сдвигался от момента награды к моменту предвестника --- как требует обучение по ошибке~\eqref{eq:ctd}~\cite{Amo2022}. Какая теория описывает мозг, пока неизвестно.

\section{Ошибка не только в награде}

Последнее расширение --- о том, \emph{о чём} ошибка. Если заменить ожидаемый сок другим, таким же ценным, но иного вкуса, ценность не изменилась, и ошибка~\eqref{eq:ctd} равна нулю. Однако \nt{DOPA}-нейроны на такую замену отвечают~\cite{Takahashi2017}. А искусственно вызванные всплески дофамина могут научить животное связи между двумя нейтральными раздражителями, вообще без награды~\cite{Sharpe2017}. Похоже, \nt{DOPA} сообщает об ошибке предсказания не только награды, но и того, \emph{что} произойдёт.

Формально это простое обобщение~\eqref{eq:ctd}: вместо одного потока награды $r$ берётся набор признаков происходящего $\varphi_k$ --- вкус, место, звук, --- и для каждого своё ожидание $M_k$ и своя ошибка~\cite{Gardner2018}:
\begin{empheq}[box=\keybox]{equation}
\delta_k(t) \;=\; \varphi_k(t) \;+\; \frac{\dd M_k}{\dd t} \;-\; \frac{M_k}{\tau_\gamma} ,
\qquad k=1,\dots,K .
\label{eq:vectortd}
\end{empheq}
Награда --- лишь один из признаков, и вектор ошибок учит не только тому, что хорошо, но и тому, что за чем следует, --- модели мира.

С этим согласуется неоднородность \nt{DOPA}-нейронов: в одной задаче разные нейроны несут разные величины --- одни связаны с наградой, другие с движением, третьи с тем, куда направлено внимание~\cite{Engelhard2019}.

\begin{keyblock}
\begin{proposition}[Жгут вместо провода]
\label{prop:bundle}
Сигнал \nt{DOPA}-системы --- не одно число. Он разложен по нейронам (распределение~\eqref{eq:expectile}, разные признаки~\eqref{eq:vectortd}, отдельная ось опасности) и по времени (быстрая ошибка и медленный уровень~\eqref{eq:vigor}). Скалярная ошибка~\eqref{eq:ctd} --- первое приближение к этому сигналу, как сумматор с порогом --- первое приближение к нейрону.
\end{proposition}
\end{keyblock}

\section{Итог}

\begin{table}[t]
\centering
\footnotesize
\setlength{\tabcolsep}{4pt}
\renewcommand{\arraystretch}{1.15}
\begin{tabular}{>{\raggedright}p{2.5cm}>{\raggedright}p{3.2cm}>{\raggedright}p{3.6cm}>{\raggedright\arraybackslash}p{4.2cm}}
\toprule
Теория & Что идёт по проводу & Главное измерение & Что известно \\
\midrule
ошибка предсказания (гл.~\ref{sec:dofa}) & скаляр $\delta$~\eqref{eq:ctd} & всплеск, перенос на предвестник, провал & измерено; первое приближение \\
распределение & набор $\delta_i$ с разной асимметрией & разные точки переворота у разных нейронов & согласие с опытом есть; роль разброса --- гипотеза \\
важность & $|\delta|$, новизна; ось опасности & всплески на неприятное у части нейронов & измерено; как используется --- обсуждается \\
цена времени & медленный уровень $\bar R$ & дофамин ускоряет действия; медленная составляющая --- и в выбросе у выходов, и в частоте импульсов & разделение быстрого и медленного измерено; $\bar R$ --- гипотеза \\
рост при приближении & $V$ или $\delta$ при уточнении оценки~\eqref{eq:ramp} & плавный рост, скачки при переносе в коридоре & рост измерен; объяснение обсуждается \\
взгляд назад & мера причинной связи~\eqref{eq:retro} & опыты, где предсказания двух теорий расходятся & результаты противоречат друг другу \\
вектор ошибок & $\delta_k$ по признакам~\eqref{eq:vectortd} & ответ на смену вкуса; обучение без награды & измерено; векторная форма --- гипотеза \\
\bottomrule
\end{tabular}
\caption{Что передаёт \nt{DOPA}-система: стандартная картина главы~\ref{sec:dofa} и её современные расширения.}
\label{tab:dofaalt}
\end{table}

Теории таблицы~\ref{tab:dofaalt} не исключают друг друга: почти все сохраняют ошибку предсказания как основу и расширяют её формат. Всерьёз конкурирует с ней только взгляд назад, и этот спор решат опыты. Для инженера вывод простой: цена широковещания из утверждения~\ref{prop:broadcastcost} падает, если провод на самом деле --- много проводов с разными адресатами.

\section{Задачи}

\begin{exercise}[\lvlA]
\label{ex:07:1}
\emph{Точки распределения для одной капли.} Сок размера $1$ приходит с вероятностью $p=0.2$, иначе награда $0$. Найдите по~\eqref{eq:expectile} $V_i$ для $\kappa_i=0.2$, $0.5$, $0.8$. Чему равна медиана этой награды, и чем точка~\eqref{eq:expectile} отличается от квантиля?
\end{exercise}

\begin{exercise}[\lvlB]
\label{ex:07:2}
\emph{Распределение по двум нейронам.} Награда равна $0$ или $x$, вероятность $x$ равна $p$. Два нейрона с правилом~\eqref{eq:asymmetric} сообщили $V(1/2)=0.25$ и $V(0.8)=0.5$. Восстановите $p$, $x$ и стандартное отклонение награды. Что сообщит нейрон с $\kappa=0.2$?
\end{exercise}

\begin{exercise}[\lvlB]
\label{ex:07:3}
\emph{Моделирование: распределение из асимметрии.} Смоделируйте девять \nt{DOPA}-нейронов с $\kappa_i=0.1,\dots,0.9$ и правилом~\eqref{eq:asymmetric}, $\alpha_i^+=2\kappa_i$, $\alpha_i^-=2(1-\kappa_i)$, при награде, равномерной на $[0,1]$. Как разброс $V_i$ зависит от $\eta$, и какое $\eta$ нужно, чтобы отличить это распределение от двух равновероятных капель $0.25$ и $0.75$?
\end{exercise}

\begin{exercise}[\lvlA]
\label{ex:07:4}
\emph{Цена времени.} (а)~Выведите~\eqref{eq:vigor} и найдите суммарную цену в оптимуме. (б)~Мозг оценил $\bar R$ с ошибкой в $k$ раз. Найдите относительную переплату при $k=2$ и $k=4$. Насколько точно медленный уровень дофамина должен сообщать $\bar R$?
\end{exercise}

\begin{exercise}[\lvlB]
\label{ex:07:5}
\emph{Всплеск при переносе.} Ожидание растёт по~\eqref{eq:ramp} с $\tau_v=1\,\text{с}$, $\tau_\gamma=4\,\text{с}$; животное мгновенно переносят из точки $T-2\,\text{с}$ в $T-1\,\text{с}$. Найдите площадь всплеска $\delta$ в долях $R$ и сколько секунд рампы до переноса он заменяет. Что покажет перенос, если дофамин сообщает само $V$?
\end{exercise}

\begin{exercise}[\lvlB]
\label{ex:07:6}
\emph{Проектирование: два сигнала на проводе.} По проводу \nt{DOPA} идут уровень, растущий до $s=1/60$ импульса в секунду за секунду, и события по $A=3$ импульса. Синапс со следом $\tau_e=1\,\text{с}$ вычитает из частоты её копию, сглаженную с $\tau_f$. При каких $\tau_f$ теряется не больше $10\,\%$ события, а вклад уровня за время следа меньше $10\,\%$ вклада события?
\end{exercise}

\begin{exercise}[\lvlB]
\label{ex:07:7}
\emph{Проектирование опыта: взгляд назад против ошибки.} В окне предвестник бывает с вероятностью $0.1$, после него награда --- с вероятностью $0.8$. Сравните $\Delta P=P(\text{нагр.}\mid\text{предв.})-P(\text{нагр.}\mid\text{нет})$ и $M$ из~\eqref{eq:retro}, если (а)~добавить награды без предвестника, $0.08$ на окно; (б)~удвоить частоту предвестника без наград.
\end{exercise}

\begin{exercise}[\lvlC]
\label{ex:07:8}
\emph{Жгут и цена широковещания.} В модели задачи~\ref{ex:06:8} $N$ нейронов разбиты на $K$ групп со своими проводами, и в каждый провод просачивается доля $\rho$ чужих ошибок. Покажите, что постоянная обучения равна $N/K+2+\rho^2N(1-1/K)$. Сколько попыток она даёт при $K\to\infty$, $N=10^{10}$ и $\rho=0.01$?
\end{exercise}

\chapter{Когда искать, а когда решать: добавляем \nt{НОРА}}
\chaptermark{Добавляем \nt{НОРА}}
\label{sec:nora}

\section{Чего не хватает}

В главе~\ref{sec:dofa} сеть училась благодаря шуму: случайные отклонения активности были поиском, а \nt{ДОФА} сообщал, стало ли от них лучше. Но там же остался незаданным один параметр --- \emph{сколько} искать. Если шум слишком силён, сеть мечется и не пользуется тем, что уже знает. Если слишком слаб, она раз за разом выбирает одно и то же и не замечает, что где-то рядом стало лучше. В обучении с подкреплением это называют выбором между использованием и поиском, и у каждого алгоритма для него есть ручка. В главе~\ref{sec:neurontypes} мы предположили, что в мозге эту ручку крутит \nt{НОРА}.

Норадреналиновых нейронов у человека несколько десятков тысяч (таблица~\ref{tab:types}), почти все они собраны в одной маленькой группе, а их выходы расходятся практически по всему мозгу, хотя разные части этой группы, как выяснилось, отчасти обслуживают разные области~\cite{Poe2020}. Это широковещательный канал, ещё более узкий, чем \nt{ДОФА}. Работает он в двух режимах~\cite{AstonJones2005}. В \emph{фоновом} нейроны выдают импульсы ровно, с частотой от долей до нескольких импульсов в секунду, и эта частота медленно меняется вместе с бодрствованием. В \emph{отрывистом} они отвечают короткой пачкой на важное для текущей задачи событие, примерно в момент принятия решения. Удобная, но неточная контрольная точка --- диаметр зрачка: он меняется вместе с активностью этих нейронов, но также вместе с активностью других областей~\cite{Joshi2016} и холинергических выходов в коре~\cite{Reimer2016}. Зрачок показывает общий уровень возбуждения, а не одну \nt{НОРА}.

\section{Коэффициент усиления}

Измерено вот что: норадреналин подавляет фоновую активность нейронов сильнее, чем их ответ на стимул, и отношение сигнала к фону растёт~\cite{Foote1975NE}. Что крутизна характеристики отдельного нейрона при этом буквально умножается на коэффициент, из опыта не следует. Как и в главе~\ref{sec:neurontypes}, возьмём простую модель, которая передаёт этот эффект: норадреналин меняет крутизну передаточной характеристики получателя~\cite{ServanSchreiber1990}. Слабые, подпороговые входы (фон) тогда дают ещё меньше, а сильные --- ещё больше. Пусть нейрон отвечает частотой
\begin{empheq}[box=\keybox]{equation}
r \;=\; F\bigl(g\,(u-\theta)\bigr), \qquad F(z)=\frac{r_{\max}}{1+e^{-z}} ,
\label{eq:gain}
\end{empheq}
где $u$ --- входной ток, $\theta$ --- порог, а $g$ --- коэффициент усиления, которым в модели управляет \nt{НОРА}. При малом $g$ частота плавно следует за входом в широком диапазоне: нейрон --- \emph{аналоговый усилитель}. При большом $g$ почти любой вход выше порога даёт полную частоту, а ниже --- ноль: нейрон --- \emph{компаратор}, почти цифровой элемент (рис.~\ref{fig:gaincurves}). Одна ручка переключает тот же нейрон между двумя режимами, которые в электронике делают разные схемы.

\begin{figure}[t]
\centering
\begin{tikzpicture}[>=Stealth,x=1.1cm,y=2.4cm]
\draw[->,gray!70!black] (-3.3,0) -- (3.4,0) node[below left,black] {\small $u$};
\draw[->,gray!70!black] (-3.3,0) -- (-3.3,1.2) node[left,black] {\small $r$};
\draw[densely dashed,gray!60!black] (0,0) -- (0,1.1);
\node[below,font=\small] at (0,0) {$\theta$};
\draw[densely dotted,gray!60!black] (-3.3,1) -- (3.2,1) node[right,black,font=\small] {$r_{\max}$};
\draw[very thick,blue!60!black] plot[domain=-3.2:3.2,samples=60] (\x,{1/(1+exp(-0.8*\x))});
\draw[very thick,orange!85!black] plot[domain=-3.2:3.2,samples=60] (\x,{1/(1+exp(-2*\x))});
\draw[very thick,red!70!black] plot[domain=-3.2:3.2,samples=120] (\x,{1/(1+exp(min(-8*\x,9)))});
\node[blue!60!black,font=\small,anchor=west] at (0.5,0.12) {малое $g$: усилитель};
\node[red!70!black,font=\small,anchor=east] at (-0.3,0.85) {большое $g$: компаратор};
\end{tikzpicture}
\caption{Передаточная характеристика~\eqref{eq:gain} при трёх значениях коэффициента усиления $g$.}
\label{fig:gaincurves}
\end{figure}

Помогает ли большое усиление против шума? Не всегда. Пусть два входа отличаются на $\Delta u$, и надо сказать, какой больше. Если шум $\sigma_{\text{до}}$ добавлен к входу \emph{до} усилителя, то усиливается и сигнал, и шум, и их отношение не меняется. Если же шум $\sigma_{\text{после}}$ добавляется \emph{после} усилителя --- от ненадёжных синапсов и случайных импульсов самой сети, как в главе~\ref{sec:code}, --- то отношение сигнала к шуму
\begin{empheq}[box=\keybox]{equation}
d' \;=\; \frac{g\,\Delta u}{\sqrt{g^2\sigma_{\text{до}}^2+\sigma_{\text{после}}^2}}
\label{eq:gainsnr}
\end{empheq}
растёт с $g$, пока $g\sigma_{\text{до}}$ не сравняется с $\sigma_{\text{после}}$~\cite{ServanSchreiber1990}.

\begin{keyblock}
\begin{proposition}[Когда усиление помогает]
\label{prop:gainnoise}
Увеличивая коэффициент усиления, \nt{НОРА} улучшает различение входов только против того шума, который возникает \emph{после} усилителя, в самой сети. Шум, пришедший вместе со входом, усиление не уберёт.
\end{proposition}
\end{keyblock}

\section{Усиление в выборе из двух}

Вернёмся к выбору из главы~\ref{sec:gaba}: две группы \nt{ГЛУТ}-нейронов со входами $I_1$, $I_2$ тормозят друг друга с силой $\beta$. Теперь у каждой группы есть коэффициент усиления $g$: в уравнениях выбора~\eqref{eq:wta} выражение в скобках умножается на $g$. Пока работают обе группы, установившиеся частоты находятся из $r_1=g(I_1-\beta r_2)$, $r_2=g(I_2-\beta r_1)$. Сложив и вычтя эти равенства, получим сумму $r_1+r_2=g(I_1+I_2)/(1+g\beta)$ и разность $r_1-r_2=g(I_1-I_2)/(1-g\beta)$. Их отношение --- насколько резко сеть предпочитает один вход другому:
\begin{empheq}[box=\keybox]{equation}
\frac{r_1-r_2}{r_1+r_2} \;=\; \varepsilon\,\frac{1+g\beta}{1-g\beta}, \qquad \varepsilon=\frac{I_1-I_2}{I_1+I_2}, \quad g\beta<1 .
\label{eq:wtagain}
\end{empheq}
Малая разница входов $\varepsilon$ усиливается в $(1+g\beta)/(1-g\beta)$ раз, и этот множитель неограниченно растёт, когда $g\beta$ подходит к единице. При $g\beta>1$ состояние, где работают обе группы, неустойчиво, как в утверждении~\ref{prop:wta}: побеждает одна, другая молчит (рис.~\ref{fig:pitchfork}).

\begin{figure}[t]
\centering
\begin{tikzpicture}[>=Stealth,x=3.2cm,y=1.5cm]
\draw[->,gray!70!black] (0,0) -- (2.15,0) node[below left,black] {\small $g\beta$};
\draw[->,gray!70!black] (0,-1.25) -- (0,1.35) node[above,black] {\small $(r_1-r_2)/(r_1+r_2)$};
\draw[gray!60!black] (1,-0.04) -- (1,0.04); \node[below right,font=\small] at (1,0) {$1$};
\node[left,font=\scriptsize] at (0,1) {$1$}; \node[left,font=\scriptsize] at (0,-1) {$-1$};
\draw[very thick,blue!60!black] plot[domain=0:0.905,samples=60] (\x,{0.05*(1+\x)/(1-\x)}) -- (1,1) -- (2.05,1);
\draw[very thick,blue!60!black] (1.105,-1) -- (2.05,-1);
\draw[thick,densely dashed,gray!60!black] (1,0) -- (2.05,0);
\node[font=\small,align=center] at (0.45,-0.62) {размышляет:\\обе группы\\работают};
\node[font=\small,align=center] at (1.55,0.45) {решила:\\работает одна};
\node[font=\small,blue!60!black,anchor=south west] at (1.05,-0.97) {слабый вход победил раньше --- держится};
\end{tikzpicture}
\caption{Выбор из двух при разном коэффициенте усиления, входы отличаются на $\varepsilon=0.05$. При $g\beta>1$ устойчивы оба решения (сплошные линии; победа слабого входа --- при $g\beta>(1+\varepsilon)/(1-\varepsilon)\approx1.1$), и сеть держит то, которое уже приняла; симметричное состояние (штриховая линия) неустойчиво.}
\label{fig:pitchfork}
\end{figure}

\begin{keyblock}
\begin{proposition}[Усиление переключает режим выбора]
\label{prop:gainwta}
В сети выбора с взаимным торможением $\beta$ коэффициент усиления $g$ переводит сеть через границу $g\beta=1$. Ниже неё сеть «размышляет»: обе группы работают, а разница входов усиливается по~\eqref{eq:wtagain}. Выше --- сеть «решила»: работает одна группа, и решение держится. Меняя $g$, \nt{НОРА} может переключать сеть между этими режимами, не трогая ни одного веса.
\end{proposition}
\end{keyblock}

\section{Усиление как температура}

Теперь добавим к выбору шум, возникающий в самой сети, после усилителя. Какая группа победит, решает знак величины $g(u_A-u_B)+\eta$, где $u_A-u_B$ --- разница входов, то есть разница выученных весов из главы~\ref{sec:dofa}, а $\eta$ --- шум с масштабом $\sigma$. Если шум распределён по логистическому закону (для гауссова кривая почти та же), вероятность выбрать $A$ равна
\begin{empheq}[box=\keybox]{equation}
P(A) \;=\; \frac{1}{1+e^{-g\,(u_A-u_B)/\sigma}} .
\label{eq:softmax}
\end{empheq}
Программист узнает здесь выбор по softmax с температурой $T=\sigma/g$. Коэффициент усиления --- обратная температура. При малом $g$ даже заметно лучший вариант выбирается ненамного чаще худшего: сеть много ищет. При большом $g$ лучший известный вариант выбирается почти всегда: сеть пользуется тем, что знает.

Уточним, что такое $g$ в~\eqref{eq:wtagain} и~\eqref{eq:softmax}: это \emph{действующее} усиление разницы входов текущей задачи в момент выбора. Два режима \nt{НОРА} действуют на него противоположно. Отрывистая пачка в момент решения повышает его, и сеть закрепляет выбор. Высокая фоновая активность, наоборот, идёт вместе с поиском: у людей перед выбором наугад базовый зрачок шире~\cite{JepmaNieuwenhuis2011}, а у крыс вход \nt{НОРА} в переднюю поясную кору переводит выбор в случайный режим~\cite{Tervo2014stochastic}. Значит, высокой фоновой \nt{НОРА} соответствует \emph{малое} действующее $g$, а пачке --- большое.

\section{Сколько искать}

Какое усиление лучше? Мы проверили это на модели. Сеть выбирает одно из двух действий по~\eqref{eq:softmax}; каждое приносит награду с какой-то вероятностью, и сеть оценивает её по наградам выбранного действия, как в главе~\ref{sec:dofa}. Время от времени мир меняется: обе вероятности разыгрываются заново. Меняться может и выбранное действие, и то, которое сеть давно не пробовала; о втором сеть может узнать только поиском. На рис.~\ref{fig:invertedU} показано, какую долю наилучшей возможной награды получает сеть при разных постоянных $g$.

\begin{figure}[t]
\centering
\begin{tikzpicture}[>=Stealth,x=1.2cm,y=1cm]
\draw[->,gray!70!black] (0,0) -- (7.6,0) node[below left,black] {\small $g$};
\draw[->,gray!70!black] (0,0) -- (0,4.4) node[above,black,font=\small] {доля наилучшей награды};
\foreach \x/\l in {0/0.5,1/1,2/2,3/4,4/8,5/16,6/32,7/64}{\draw[gray!60!black] (\x,-0.06) -- (\x,0.06); \node[below,font=\scriptsize] at (\x,0) {\l};}
\foreach \v/\y in {0.8/0.8,0.9/2.4,1.0/4.0}{\draw[gray!60!black] (-0.06,\y) -- (0.06,\y); \node[left,font=\scriptsize] at (0,\y) {\v};}
\draw[very thick,blue!60!black] plot[mark=*,mark size=1.3pt] coordinates {(0,0.368) (1,0.880) (2,1.728) (3,2.784) (4,3.456) (5,3.616) (6,3.488) (7,3.264)};
\draw[very thick,red!70!black] plot[mark=*,mark size=1.3pt] coordinates {(0,0.368) (1,0.672) (2,1.184) (3,1.840) (4,2.176) (5,1.920) (6,1.456) (7,1.104)};
\node[blue!60!black,font=\small,anchor=south] at (5,3.7) {спокойный мир};
\node[red!70!black,font=\small,anchor=north] at (5.1,1.25) {быстро меняющийся мир};
\draw[->,thick,blue!60!black] (5,3.25) -- (5,3.55);
\draw[->,thick,red!70!black] (4,1.8) -- (4,2.1);
\node[font=\small\itshape,gray!35!black,anchor=west] at (0.15,2.3) {ищет вслепую};
\node[font=\small\itshape,gray!35!black] at (6.45,2.55) {не замечает перемен};
\end{tikzpicture}
\caption{Доля наилучшей возможной награды при постоянном коэффициенте усиления $g$ (шкала по $g$ логарифмическая). Синяя кривая --- мир меняется в среднем раз в тысячу попыток, красная --- раз в двадцать. Стрелки отмечают лучшее $g$: в быстро меняющемся мире оно вдвое меньше. Среднее по 60 прогонам по 4000 попыток.}
\label{fig:invertedU}
\end{figure}

При $g=0.5$ сеть почти не пользуется знанием и получает около трёх четвертей возможного, при очень большом $g$ упускает перемены. Лучшее значение: $g=16$, когда мир меняется раз в тысячу попыток (доля $0.976$), и $g=8$, когда раз в двадцать ($0.886$).

\begin{keyblock}
\begin{proposition}[Перевёрнутое U]
\label{prop:explore}
Награда, которую собирает сеть, зависит от коэффициента усиления как перевёрнутое U: слишком малое $g$ тратит попытки на поиск, слишком большое не замечает перемен. Лучшее $g$ тем меньше, чем быстрее меняется мир.
\end{proposition}
\end{keyblock}

Перевёрнутое U наблюдается и в опытах: животные решают задачу лучше всего при умеренной фоновой активности \nt{НОРА}-нейронов с чёткими отрывистыми ответами, а при высокой фоновой активности отвлекаются и переключаются между задачами~\cite{AstonJones2005}. Это согласуется с различием режимов из предыдущего раздела: отрывистая пачка в момент решения даёт большое действующее $g$ как раз тогда, когда надо выбрать, а высокая фоновая активность без пачек усиливает всё подряд, в том числе посторонние входы, так что действующее $g$ для текущей задачи падает, --- это поиск.

\section{Кто крутит ручку}

Раз лучшее усиление зависит от того, как быстро меняется мир, его хорошо бы подстраивать. Но откуда \nt{НОРА}-нейронам знать, что мир изменился? Естественный признак --- \emph{неожиданность}: ошибки предсказания стали больше обычного. Важно различать два вида неопределённости. Если награда просто случайна, ошибки велики всегда, и это ожидаемо. Если же ошибки вдруг выросли по сравнению с их обычной величиной, значит, что-то поменялось. По одной из гипотез, именно эту неожиданную неопределённость и сообщает \nt{НОРА}, а ожидаемую --- \nt{АЦЕТ}~\cite{YuDayan2005}.

Что с этим сигналом делать? Можно понизить усиление и больше искать. Можно повысить скорость обучения, чтобы быстрее забыть устаревшие оценки: опыты показывают, что после резкой перемены люди учатся быстрее, и зрачок при этом расширяется~\cite{Nassar2012}; напомним, что зрачок --- показатель не только \nt{НОРА}, но и \nt{АЦЕТ}~\cite{Reimer2016}. А можно сделать и то и другое сразу: по ещё одной гипотезе отрывистая пачка \nt{НОРА} работает как \emph{прерывание} --- сигнал, по которому сеть сбрасывает текущее состояние и перестраивается под новую обстановку~\cite{BouretSara2005}, как общая линия прерывания процессора.

Честно скажем, чего мы не нашли. В модели мы пробовали оба простых способа подстройки: понижать $g$ или повышать скорость обучения, когда сглаженная ошибка превышает свою долговременную среднюю. В мире, который то подолгу стоит, то часто меняется, лучшие настройки обоих способов дали $0.926$--$0.927$ от наилучшей награды --- почти столько же, сколько лучшее постоянное $g$ ($0.925$ при $g=8$) или постоянная, но достаточно большая скорость обучения ($0.928$), а неудачные настройки --- меньше. Кривые на рис.~\ref{fig:invertedU} у вершины пологие, и в таком мире подстраивать почти нечего. Подстройка должна окупаться там, где разные режимы требуют очень разных настроек. Какой именно закон управления реализует \nt{НОРА}, пока не установлено.

\section{Итог}

\begin{table}[t]
\centering
\footnotesize
\setlength{\tabcolsep}{4pt}
\renewcommand{\arraystretch}{1.15}
\begin{tabular}{>{\raggedright}p{3.0cm}>{\raggedright}p{4.4cm}>{\raggedright\arraybackslash}p{6.0cm}}
\toprule
Что делает \nt{НОРА} & Как & Что известно \\
\midrule
меняет крутизну ответа нейронов & коэффициент усиления $g$ в~\eqref{eq:gain} & повышение отношения сигнала к шуму измерено; мультипликативная модель --- упрощение \\
переключает выбор & переводит сеть через $g\beta=1$ & следует из модели~\eqref{eq:wtagain}; прямых измерений порога нет \\
задаёт, сколько искать & $g$ --- обратная температура~\eqref{eq:softmax}; фон --- малое действующее $g$ (поиск), пачка --- большое (решение) & перевёрнутое U и два режима работы измерены; связь с поиском --- хорошо обоснованная гипотеза \\
сообщает о переменах & неожиданная неопределённость & зрачок и скорость обучения растут после перемен --- измерено; что это сигнал \nt{НОРА} --- гипотеза \\
прерывает и сбрасывает & отрывистая пачка на неожиданное событие & пачки на важные события измерены; «прерывание» --- гипотеза \\
\bottomrule
\end{tabular}
\caption{Что добавляет \nt{НОРА} к сети из \nt{ГЛУТ}, \nt{ГАМК} и \nt{ДОФА}.}
\label{tab:nora}
\end{table}

\nt{ДОФА} говорит сети, \emph{куда} менять веса. \nt{НОРА} не меняет ни одного веса, но решает, \emph{как} сеть пользуется тем, что уже знает (таблица~\ref{tab:nora}): одна ручка --- коэффициент усиления --- превращает нейрон из усилителя в компаратор, сеть выбора --- из «размышляющей» в «решившую», а выбор --- из поиска в использование. Как \nt{НОРА} узнаёт, насколько быстро меняется мир, --- открытый вопрос.

\section{Задачи}

\begin{exercise}[\lvlA]
\label{ex:08:1}
\emph{Одна ручка на весь мозг.}
(а)~По таблице~\ref{tab:types} оцените, сколько нейронов мозга приходится на один \nt{НОРА}-нейрон и сколько всего мест выброса у \nt{НОРА}-системы.
(б)~Пусть $\Delta u=1$, $\sigma_{\text{до}}=0.5$, $\sigma_{\text{после}}=4$. Вычислите $d'$ по~\eqref{eq:gainsnr} при $g=1$, $g=8$ и $g\to\infty$.
(в)~При каком $g$ достигается $90\,\%$ предельного $d'$? Покажите, что ответ зависит только от отношения $\sigma_{\text{после}}/\sigma_{\text{до}}$, и объясните смысл утверждения~\ref{prop:gainnoise} через это отношение.
\end{exercise}

\begin{exercise}[\lvlB]
\label{ex:08:2}
\emph{Выбор с коэффициентом усиления: вывод.} Система: $\tau\,\dd r_1/\dd t=-r_1+g[I_1-\beta r_2]_+$, $\tau\,\dd r_2/\dd t=-r_2+g[I_2-\beta r_1]_+$.
(а)~Выведите~\eqref{eq:wtagain}. Найдите собственные числа линеаризации в области, где работают обе группы, и время, за которое затухает разность $r_1-r_2$. Вычислите множитель усиления разницы и это время при $g\beta=0.5$ и $0.9$, $\tau=20\ms$.
(б)~Докажите, что состояние, где работают обе группы, существует только при $g\beta<(1-\varepsilon)/(1+\varepsilon)$ (при $\varepsilon>0$). Чему равна эта граница при $\varepsilon=0.05$?
(в)~Докажите утверждение из подписи к рис.~\ref{fig:pitchfork}: победа слабого входа устойчива при $g\beta>(1+\varepsilon)/(1-\varepsilon)$. Что происходит между двумя границами?
(г)~Объясните, почему у границы $g\beta=1$ сеть «размышляет» долго, и чем это полезно и вредно.
\end{exercise}

\begin{exercise}[\lvlA]
\label{ex:08:3}
\emph{Softmax из шума.}
(а)~Покажите, что при логистическом шуме $\eta$ с масштабом $\sigma$, $P(\eta<z)=1/(1+e^{-z/\sigma})$, вероятность выбора равна~\eqref{eq:softmax}.
(б)~Пусть каждая из $K$ групп получает свой независимый шум с распределением Гумбеля, $P(\eta_k<z)=\exp(-e^{-z/\sigma})$, и побеждает группа с наибольшим $gu_k+\eta_k$. Докажите, что $P(k)=e^{gu_k/\sigma}/\sum_le^{gu_l/\sigma}$, и выведите из этого случай $K=2$.
(в)~Вычислите $P(A)$ при $u_A-u_B=0.2$, $\sigma=1$ и $g=1$, $8$, $32$, и найдите $g$, при котором $P(A)=0.9$ при $u_A-u_B=0.1$.
(г)~Для гауссова шума $P(A)=\Phi\bigl(g(u_A-u_B)/s\bigr)$. Подберите $s$ так, чтобы наклоны в нуле совпали с логистической кривой, и найдите наибольшее расхождение кривых. Проверьте слова «кривая почти та же».
\end{exercise}

\begin{exercise}[\lvlB]
\label{ex:08:4}
\emph{Оценка лучшего усиления.} В модели рис.~\ref{fig:invertedU} вероятности наград двух действий после каждой перемены равномерно распределены на $[0,1]$.
(а)~Пусть сеть знает вероятности точно. Какую долю наилучшей награды она получает при $g=0.5$? Сравните с рис.~\ref{fig:invertedU}.
(б)~Пусть лучший вариант лучше худшего на $\Delta$. Сеть пробует худший вариант примерно раз в $e^{g\Delta}$ попыток. Потребуйте, чтобы это время было порядка среднего времени между переменами $1/h$, и получите $g^*\approx\ln(1/h)/\bar\Delta$. Найдите $\bar\Delta=\langle|p_A-p_B|\rangle$.
(в)~Вычислите оценку при $h=0.001$, $0.01$ и $0.05$ и сравните с лучшими $g$ модели. Чего оценка не учитывает?
\end{exercise}

\begin{exercise}[\lvlB]
\label{ex:08:5}
\emph{Моделирование: как лучшее усиление зависит от скорости перемен.} Работайте с копией \texttt{models/\allowbreak nora\_explore.py} в отдельной папке: скрипт пишет файлы \texttt{scan\_h*.txt} в текущую папку.
(а)~Для $h\in\{0.001,0.003,0.01,0.03,0.05,0.1,0.2\}$ и $g=2^{k/2}$, $k=-2,\dots,14$, найдите долю наилучшей награды (60 прогонов по 4000 попыток) и лучшее $g^*(h)$.
(б)~Постройте $g^*$ в зависимости от $\ln(1/h)$ вместе с оценкой задачи~\ref{ex:08:4}. В каком диапазоне оценка верна?
(в)~Объясните, почему при быстрых переменах $g^*$ перестаёт убывать. Указание: сравните время, за которое устаревает добытая поиском информация, со временем $1/\eta$, за которое правило $Q\leftarrow Q+\eta(R-Q)$ её усваивает, и оцените стационарный разброс $Q$ при $\eta=0.2$.
(г)~Как меняется высота вершины перевёрнутого U с $h$ и почему кривые у вершины пологие?
\end{exercise}

\begin{exercise}[\lvlB]
\label{ex:08:6}
\emph{Проектирование: прерывание для сети выбора.} Сеть задачи~\ref{ex:08:2} с $\beta=2$, $\tau=20\ms$, $g=1$ уже решила: $r_1=0.9$, $r_2=0$, а входы теперь $I_1=0.9$, $I_2=1.0$. По утверждению~\ref{prop:wta} старое решение держится, хотя второй вход сильнее. \nt{НОРА} посылает прерывание: на время $T_p$ усиление падает до $g_{\text{lo}}$, затем возвращается к $1$.
(а)~Покажите, что при $g=1$ исход определяется знаком $r_1-r_2-g(I_2-I_1)/(g\beta-1)$ в момент возврата усиления (пока работают обе группы).
(б)~При $g_{\text{lo}}=0$ найдите наименьшее $T_p$ аналитически.
(в)~Найдите наименьшее $T_p$ моделью (Эйлер, шаг $10$~мкс) при $g_{\text{lo}}=0$, $0.1$, $0.25$, $0.4$. Сформулируйте правило проектирования: сколько должна длиться пачка и насколько глубоко опускать усиление.
(г)~Почему прерывание через усиление лучше, чем прямое торможение победителя дополнительным \nt{ГАМК}-входом?
\end{exercise}

\begin{exercise}[\lvlC]
\label{ex:08:7}
\emph{Когда подстройка усиления окупается.} В разделе «Кто крутит ручку» подстройка $g$ по неожиданности почти ничего не дала. Выясните, где предел.
(а)~Оракул знает, спокойная сейчас эпоха или переменчивая, и ставит в каждой своё постоянное $g$. Найдите моделью лучшую пару и долю наилучшей награды в мире главы (эпохи по $1000$ попыток, $h=0.001$ и $0.05$). Сравните с лучшим постоянным $g$ и с подстройкой по неожиданности из \texttt{models/\allowbreak nora\_explore.py}. Какую часть возможного выигрыша подстройка уже забирает?
(б)~Спроектируйте мир, где оракул выигрывает у лучшего постоянного $g$ заметно больше. Используйте задачу~\ref{ex:08:4}: лучшее $g$ зависит и от $h$, и от $\bar\Delta$.
(в)~Проверьте в этом мире подстройку по неожиданности. Какую часть выигрыша оракула она забирает и почему?
(г)~Открытая часть. Какой сигнал должен получать \nt{НОРА}, чтобы приблизиться к оракулу? Сравните сигнал «сглаженная $|\delta|$ выше своей средней» с оценкой частоты перемен $h$ по последовательности наград. Какие измерения зрачка или активности \nt{НОРА}-нейронов могли бы различить эти две стратегии?
\end{exercise}

\chapter{Когда писать, а когда читать: добавляем \nt{ACET}}
\chaptermark{Добавляем \nt{ACET}}
\label{sec:acet}

\section{Чего не хватает}

У микросхемы памяти, кроме проводов адреса и данных, есть ещё одна линия --- разрешение записи. Пока она выключена, память только читается; когда включена, то, что лежит на линиях данных, записывается. Без такой линии память не работает: если каждое чтение одновременно что-то записывает, прочитанное, в том числе прочитанное с ошибкой, тут же записывается обратно.

В мозге эта задача стоит острее, чем в микросхеме. В главе~\ref{sec:glutcomputer} память была группой \nt{GLUT}-нейронов, которую держат включённой её собственные связи (рис.~\ref{fig:bistable}). В главе~\ref{sec:dofa} мы видели, что эти связи учатся правилом Хебба прямо по ходу работы. Значит, одни и те же синапсы и хранят старое, и принимают новое, а отдельной фазы записи, как и часов, нет. Чем сеть отличает момент, когда нужно запомнить то, что видишь, от момента, когда нужно вспомнить то, что знаешь? В главе~\ref{sec:neurontypes} мы предположили, что эту линию держит \nt{ACET}. В этой главе мы проверим, как такая линия должна работать и почему без неё не обойтись.

\section{Три действия одного провода}

Ацетилхолиновых нейронов, работающих внутри мозга, немного, и собраны они в нескольких небольших группах, а их выходы расходятся по всей коре. Это широковещательный канал, такой же, как \nt{DOPA} и \nt{NORA}. Уровень ацетилхолина в коре высок при внимательном бодрствовании и низок в медленном сне~\cite{Hasselmo1999}. Как и у дофамина, смысл сигнала задаёт рецептор получателя (утверждение~\ref{prop:receptor}): у ацетилхолина есть и быстрые рецепторы, открывающие канал, и медленные, запускающие реакции внутри клетки. Для нас важны три действия.

\emph{Первое: ослабить внутренние связи.} Опыты на срезах обонятельной коры показали, что ацетилхолин сильно подавляет передачу по связям между нейронами одной области и почти не трогает входы, приходящие в неё извне~\cite{HasselmoBower1992}.

\emph{Второе: усилить входы.} Ацетилхолин повышает возбудимость нейронов и облегчает передачу по некоторым входным путям; сигнал от органов чувств становится громче собственной активности сети~\cite{Hasselmo2006}.

\emph{Третье: облегчить изменение весов.} При высоком уровне ацетилхолина связи меняются легче~\cite{Hasselmo2006}.

Для инженера все три действия --- одна операция: переключение из режима «читать» в режим «писать». Сеть перестаёт слушать саму себя, начинает слушать вход и запоминает то, что слышит.

\section{Сеть, которая дописывает}

Возьмём $N$ \nt{GLUT}-нейронов, соединённых друг с другом. Сохраним в их связях несколько образов --- наборов по $pN$ одновременно активных нейронов --- правилом Хебба~\cite{Hebb1949}. Если подать на такую сеть половину образа, связи возбудят недостающую половину: сеть \emph{дописывает} образ по части, как группа на рис.~\ref{fig:bistable} поддерживала сама себя. Это ассоциативная память: адресом служит часть содержимого. \nt{GABA}, как в главе~\ref{sec:gaba}, держит общее число активных нейронов около $pN$, чтобы вместе не могли зажечься два образа.

Обозначим уровень ацетилхолина $a$, от $0$ до $1$, и запишем три его действия в модель прямо:
\begin{empheq}[box=\keybox]{equation}
\tau\frac{\dd r_i}{\dd t} = -r_i + F\Bigl(\tfrac{1+a}{2}\,x_i + (1-a)\sum_j W_{ij}\,r_j - g\Bigr),
\qquad
\frac{\dd W_{ij}}{\dd t} = \eta\,a\,(r_i-p)(r_j-p) .
\label{eq:acetnet}
\end{empheq}
Здесь $r_i$ --- частота нейрона $i$, $x_i$ --- его вход от органов чувств, $F$ --- пороговая передаточная характеристика, $g$ --- общее торможение от \nt{GABA}, растущее, когда активных нейронов больше $pN$. Множитель $\frac{1+a}{2}$ --- усиление входа, $1-a$ --- ослабление внутренних связей, $\eta a$ --- скорость обучения. Второе уравнение --- правило Хебба~\eqref{eq:hebb} в форме, удобной для редких образов~\cite{Tsodyks1988}: вес растёт, когда оба нейрона активнее обычного, и уменьшается, когда активен только один. Отрицательная часть весов, как всегда, стоит за \nt{GABA}-посредниками. Часов нет: и нейроны, и веса меняются непрерывно.

\begin{figure}[t]
\centering
\begin{tikzpicture}[>=Stealth,
  nrn/.style={circle,draw,very thick,fill=blue!6,minimum size=0.8cm,inner sep=0pt,font=\small},
  acet/.style={circle,draw,very thick,double,double distance=1.2pt,fill=orange!15,minimum size=1.35cm,inner sep=0pt,font=\scriptsize},
  bc/.style={->,thick,densely dashed,orange!85!black}]
\foreach \i/\y in {1/2.4,2/1.2,3/0}{
  \node[nrn] (N\i) at (3,\y) {$r_\i$};
  \draw[->,thick,blue!60!black] (0,\y) node[left,black] {\small $x_\i$} -- (N\i);}
\node[left,align=right,font=\small] at (-0.5,1.2) {от органов\\чувств};
\draw[->,thick,gray!60!black] (N1) to[bend left=30] (N2);
\draw[->,thick,gray!60!black] (N2) to[bend left=30] (N1);
\draw[->,thick,gray!60!black] (N2) to[bend left=30] (N3);
\draw[->,thick,gray!60!black] (N3) to[bend left=30] (N2);
\draw[->,thick,gray!60!black] (N1.east) .. controls (4.6,2.0) and (4.6,0.4) .. (N3.east);
\node[right,font=\small,gray!40!black,align=left] at (4.7,1.2) {внутренние\\связи $W$};
\node[acet] (A) at (1.2,-1.9) {\nt{ACET}};
\draw[bc] (A) -- (1.6,-0.15);
\node[font=\small,orange!60!black,anchor=east] at (0.95,-0.9) {$+$ вход};
\draw[bc] (A) -- (4.3,0.6);
\node[font=\small,orange!60!black,anchor=west] at (4.3,0.1) {$-$ внутренние связи, $+$ скорость обучения};
\node[right,font=\small,align=left] at (2.2,-1.9) {высокий уровень: «писать»\\низкий уровень: «читать»};
\end{tikzpicture}
\caption{\nt{ACET} как линия разрешения записи. Нейроны $r_i$ получают вход $x_i$ от органов чувств (синие стрелки) и возбуждают друг друга через внутренние связи $W$ (серые), в которых хранятся образы. Ацетилхолин (штриховые стрелки) усиливает вход, ослабляет внутренние связи и облегчает изменение весов, как в~\eqref{eq:acetnet}.}
\label{fig:acetnet}
\end{figure}

\section{Запись и чтение мешают друг другу}

Проверим модель~\eqref{eq:acetnet} на задаче, в которой запись и чтение действительно конфликтуют. Сеть из $200$ нейронов уже хранит пять образов по $20$ активных нейронов. Теперь нужно запомнить шестой, который наполовину совпадает с одним из старых: десять нейронов у них общие. Так бывает постоянно: новая комната похожа на знакомую, новое лицо --- на старое.

\emph{Запись при низком $a$.} Сначала отделим первые два действия ацетилхолина от третьего: оставим скорость обучения полной, а меняться будут только усиление входа и ослабление внутренних связей. При $a=0$ вход зажигает общие десять нейронов, а они через внутренние связи зажигают оставшуюся половину \emph{старого} образа. \nt{GABA} не даёт гореть обоим сразу, и старый образ, поддержанный собственными связями, вытесняет новый, который держится только на входе. Сеть видит не то, что перед ней, а то, что помнит, --- и правило Хебба старательно записывает увиденное.

В итоге новый образ вообще не запомнился: по его отличительной половине сеть ничего не вспоминает. А старый испорчен: вызванный своей отличительной половиной, он тянет за собой в среднем шесть чужих нейронов из десяти. При $a=0.1$ новый образ уже запоминается, но сливается со старым, и порча даже сильнее: каждый из двух, вызванный своей половиной, тянет за собой около восьми чужих нейронов; начиная с $a=0.2$ запись чистая: лишних нейронов при вспоминании меньше одного (усреднение по десяти прогонам).

\emph{Запись с настоящей скоростью обучения.} Теперь пусть ацетилхолин, как в~\eqref{eq:acetnet}, задаёт и скорость обучения. За одно предъявление новый образ запоминается целиком только при $a\ge0.9$; при $a=0.75$ --- на восемь десятых, при $a\le0.5$ не запоминается совсем.

\emph{Чтение.} Подадим на сеть с пятью чисто записанными образами половину образа, а потом уберём и её. При $a\le0.2$ сеть дописывает образ целиком и держит его сама; при $a=0.3$ --- почти целиком; при $a\ge0.4$ внутренние связи слишком ослаблены, и после того как подсказка исчезла, активность гаснет (рис.~\ref{fig:acetsim}).

\begin{figure}[t]
\centering
\begin{tikzpicture}[>=Stealth,x=9cm,y=3.2cm]
\draw[->,gray!70!black] (0,0) -- (1.1,0) node[right,black] {\small $a$};
\draw[->,gray!70!black] (0,0) -- (0,1.2);
\foreach \x in {0,0.2,0.4,0.6,0.8,1.0}{\draw[gray!60!black] (\x,-0.02) -- (\x,0.02); \node[below,font=\scriptsize] at (\x,0) {$\x$};}
\foreach \y in {0,0.5,1}{\draw[gray!60!black] (-0.01,\y) -- (0.01,\y); \node[left,font=\scriptsize] at (0,\y) {$\y$};}
\node[rotate=90,font=\small] at (-0.1,0.6) {доля восстановленного образа};
\draw[very thick,blue!60!black] plot[mark=*,mark size=1.6pt] coordinates {(0,1.00) (0.1,1.00) (0.2,1.00) (0.3,0.96) (0.4,0.04) (0.5,0.00) (0.75,0.00) (1.0,0.00)};
\draw[very thick,red!70!black] plot[mark=*,mark size=1.6pt] coordinates {(0.1,0.00) (0.2,0.00) (0.3,0.00) (0.5,0.00) (0.75,0.80) (0.9,1.00) (1.0,1.00)};
\node[font=\small,blue!60!black,anchor=west,align=left] at (0.03,0.5) {чтение:\\образ\\по половине};
\node[font=\small,red!70!black,anchor=east,align=right] at (1.0,0.35) {запись:\\новый образ\\за раз};
\end{tikzpicture}
\caption{Модель~\eqref{eq:acetnet}: $200$ нейронов, образы по $20$ активных, среднее по десяти прогонам. Синие точки --- чтение: какая доля образа восстановлена по его половине после того, как подсказка убрана. Красные точки --- запись: какая доля нового образа, наполовину похожего на старый, вспоминается после одного предъявления, если ацетилхолин задаёт и скорость обучения. Чтение работает при $a\le0.3$, запись --- при $a\ge0.9$.}
\label{fig:acetsim}
\end{figure}

\begin{keyblock}
\begin{proposition}[Запись и чтение требуют разных режимов]
\label{prop:writeread}
В модели~\eqref{eq:acetnet}, где одни и те же связи хранят старое и учатся новому, а ацетилхолин задаёт и скорость обучения, нет уровня $a$, при котором одинаково хорошо идут и запись, и чтение. Для записи внутренние связи нужно ослабить, иначе сеть запишет не вход, а вспомненное; для чтения их нужно усилить, иначе она не допишет образ по части. Нужен переключатель, и один широковещательный провод может им служить.
\end{proposition}
\end{keyblock}

Если бы ацетилхолин не менял скорость обучения, для обоих дел годилась бы узкая полоса $0.2\le a\le0.3$. Но тогда сеть училась бы и при каждом чтении, то есть заново записывала бы прочитанное вместе с его ошибками. Сама идея --- подавлять внутренние связи во время обучения --- взята из моделей, построенных по измерениям на срезах~\cite{HasselmoAndersonBower1992}. Числа выше относятся к нашей упрощённой модели; где именно лежат границы режимов в мозге, неизвестно.

\section{Сколько верить глазам}

Переключатель «писать/читать» --- крайний случай более общего вопроса: насколько верить входу и насколько --- памяти? Для этого вопроса у инженеров есть точный ответ, и он проливает свет на то, почему один провод управляет сразу и режимом, и скоростью обучения.

Пусть сеть следит за величиной $s(t)$, которая медленно блуждает: за секунду её дисперсия растёт на $q$. Органы чувств сообщают $y(t)=s(t)+$шум с интенсивностью $R$. Лучшая оценка $\hat s$ в непрерывном времени даётся фильтром Калмана--Бьюси~\cite{KalmanBucy1961}:
\begin{empheq}[box=\keybox]{equation}
\frac{\dd\hat s}{\dd t} \;=\; \frac{P}{R}\,\bigl(y-\hat s\bigr),
\qquad
\frac{\dd P}{\dd t} \;=\; q-\frac{P^2}{R} ,
\label{eq:kalman}
\end{empheq}
где $P$ --- дисперсия ошибки оценки, то есть насколько сеть сама не уверена в том, что знает. Первое уравнение --- та же RC-цепь, что у мембраны в модели нейрона~\eqref{eq:glutneuron}, но с постоянной времени $R/P$, которая меняется. Когда сеть не уверена, $P$ велико, постоянная времени мала, и оценка быстро следует за входом: это «писать». Когда уверена, $P$ мало, и оценка почти не слушает вход, держась за память: это «читать».

В установившемся режиме $P=\sqrt{qR}$ и постоянная времени памяти равна $\sqrt{R/q}$: если мир уходит на единицу за сто секунд, $q=0.01$, а шум органов чувств $R=1$, память должна помнить около десяти секунд.

Главное здесь то, что одно число $P/R$ задаёт сразу две вещи: и вес входа по сравнению с памятью, и скорость, с которой меняется хранимая оценка. Это ровно то, что делает ацетилхолин в~\eqref{eq:acetnet}. По гипотезе~\cite{YuDayan2005}, ацетилхолин и сообщает сети величину, подобную $P$, --- \emph{ожидаемую} неопределённость, ту, о которой сеть знает заранее. Медленным этот канал бывает не всегда: часть \nt{ACET}-нейронов отвечает на награду и наказание за два десятка миллисекунд, тем сильнее, чем неожиданнее подкрепление~\cite{Hangya2015}.

\begin{keyblock}
\begin{proposition}[Одна ручка для веса входа и скорости обучения]
\label{prop:kalman}
В оптимальном фильтре~\eqref{eq:kalman} вес, с которым вход смешивается с памятью, и скорость обучения --- одна и та же величина $P/R$. Поэтому разумно управлять ими одним сигналом. При неизменных свойствах мира этот сигнал выходит на постоянный уровень $\sqrt{q/R}$, и подстраивать его нужно, только когда свойства мира меняются.
\end{proposition}
\end{keyblock}

Второе предложение утверждения объясняет результат главы~\ref{sec:nora}, где подстройка скорости обучения по неожиданности не выиграла у лучшей постоянной: в мире с неизменными свойствами перемен лучшая скорость обучения и есть постоянная. Выигрыш появляется, когда мир внезапно становится другим. Тогда оценка вдруг оказывается далеко от входа, а $P$ об этом не знает. Правильное действие --- резко поднять $P$ и тем самым перейти в режим записи. По гипотезе, которую мы обсуждали в главе~\ref{sec:nora}, о такой \emph{неожиданной} перемене сообщает \nt{NORA}~\cite{YuDayan2005}. Разделение труда получается естественным: \nt{NORA} замечает, что мир изменился, \nt{ACET} держит сеть в режиме записи, пока новая обстановка не выучена.

\section{Сон как режим чтения}

Если высокий уровень ацетилхолина --- режим записи, то что делает сеть при низком? В медленном сне уровень ацетилхолина падает, внутренние связи освобождаются от подавления, входа от органов чувств почти нет. Сеть в режиме чтения без подсказки проигрывает то, что в ней записано. Регистрация активности показывает, что нейроны, работавшие днём вместе, в медленном сне снова срабатывают вместе~\cite{WilsonMcNaughton1994}, и даже в том же порядке~\cite{Skaggs1996}. По гипотезе~\cite{Hasselmo1999}, эти повторы переписывают быстро выученное днём в долговременную память (искусственное удлинение коротких вспышек повторов улучшает память~\cite{FernandezRuiz2019}): днём --- запись с входа, ночью --- перезапись изнутри. Для инженера это похоже на запись в быстрый буфер днём и его сброс в медленную постоянную память ночью.

\section{Итог}

\begin{table}[t]
\centering
\footnotesize
\setlength{\tabcolsep}{4pt}
\renewcommand{\arraystretch}{1.15}
\begin{tabular}{>{\raggedright}p{3.2cm}>{\raggedright}p{4.2cm}>{\raggedright\arraybackslash}p{6.0cm}}
\toprule
Что делает \nt{ACET} & Как & Что известно \\
\midrule
ослабляет внутренние связи & множитель $1-a$ в~\eqref{eq:acetnet} & измерено на срезах \\
усиливает вход & множитель $\frac{1+a}{2}$ & повышение возбудимости и передачи по входным путям измерено \\
облегчает обучение & скорость $\eta a$ & измерено \\
переключает «писать/читать» & утверждение~\ref{prop:writeread} & следует из моделей; прямого измерения режимов нет \\
сообщает ожидаемую неопределённость & $P$ в~\eqref{eq:kalman} & гипотеза \\
освобождает сеть для повторов во сне & низкий $a$ в медленном сне & низкий уровень во сне и повторы измерены; перезапись в долговременную память --- гипотеза \\
\bottomrule
\end{tabular}
\caption{Что добавляет \nt{ACET} к сети из \nt{GLUT}, \nt{GABA}, \nt{DOPA} и \nt{NORA}.}
\label{tab:acet}
\end{table}

\nt{DOPA} говорит сети, куда менять веса, \nt{NORA} --- насколько решительно пользоваться тем, что она знает, а \nt{ACET} --- слушать ли ей мир или себя (таблица~\ref{tab:acet}). Это последняя широковещательная линия управления из таблицы~\ref{tab:types}: разрешение записи, без которого память, хранящая образы в тех же связях, которыми она их читает, портила бы себя при каждом чтении.

\section{Задачи}

\begin{exercise}[\lvlA]
\label{ex:09:1}
\emph{Сколько помнить.} Фильтр~\eqref{eq:kalman} с $q=0.01$, $R=1$ помнит около $10$~с. Найдите $P_\infty$ и память $\sqrt{R/q}$, если (а)~шум датчика вырос вдвое по амплитуде; (б)~мир стал блуждать в сто раз быстрее. (в)~Во сколько раз должен вырасти шум, чтобы сеть помнила минуту?
\end{exercise}

\begin{exercise}[\lvlB]
\label{ex:09:2}
\emph{Режим записи после перемены.} Мир изменился, и неопределённость фильтра~\eqref{eq:kalman} с $q=0.01$, $R=1$ подскочила до $P_0\to\infty$. (а)~С какой постоянной времени $P$ возвращается к $P_\infty$? (б)~Через сколько секунд скорость обучения превысит установившуюся не больше чем на $10\,\%$?
\end{exercise}

\begin{exercise}[\lvlB]
\label{ex:09:3}
\emph{Постоянная скорость обучения.} Сеть учится как $\dd\hat s/\dd t=(y-\hat s)/T$; мир блуждает, $\dd s/\dd t=w$, $y=s+n$, где $w$ и $n$ --- белые шумы с интенсивностями $q$ и $R$. Найдите лучшее $T$ и дисперсию ошибки при нём. Во сколько раз она вырастет, если $T$ ошибочно в $2$ и в $10$ раз?
\end{exercise}

\begin{exercise}[\lvlB]
\label{ex:09:4}
\emph{Где граница записи.} В \texttt{models/\allowbreak acet\_memory.py} порог $\theta=0.35$, вес внутри образа $w=0.041$ (в идеале $0.045$). Новый образ делит $m=10$ нейронов со старым. При каком $a$ в~\eqref{eq:acetnet} остальные нейроны старого образа не зажигаются от этих $m$ (порог резкий)?
\end{exercise}

\begin{exercise}[\lvlB]
\label{ex:09:5}
\emph{Моделирование: ёмкость памяти.} Измените в \texttt{models/\allowbreak acet\_memory.py} функцию \texttt{read\_sweep}: при $a=1$ запишите $M$ образов ($M$ от $5$ до $100$), затем при $a=0$ прочитайте первые десять по половине. При каком $M$ появляются ошибки и как предельное $M_{\max}$ зависит от $N$ и $p$?
\end{exercise}

\begin{exercise}[\lvlB]
\label{ex:09:6}
\emph{Проектирование: запись без утечки.} Со скоростью обучения $\eta a$ сеть учится и при чтении. Предложите монотонную $\eta(a)\le\eta$, равную нулю при $a\le0.3$ и не хуже $\eta a$ при $a\ge0.9$. Проверьте: после $20$ чтений при $a=0.2$ с подсказкой, где три чужих нейрона, сколько их вошло в образ?
\end{exercise}

\begin{exercise}[\lvlC]
\label{ex:09:7}
\emph{Как переписать буфер ночью.} (а)~Во сне $a=0.05$, повтор длится $0.1\,\text{с}$ против $0.2\,\text{с}$ днём при $a=1$. Сколько повторов накопят изменение весов одного дневного предъявления при скорости $\eta a$ и при $\eta(a)$ задачи~\ref{ex:09:6}? (б)~Хранилище учится в $100$ раз медленнее буфера и слышит образ $20$ раз за ночь по $0.1\,\text{с}$. Сколько нужно ночей?
\end{exercise}

\chapter{Вся машина}
\chaptermark{Вся машина}
\label{sec:synthesis}

\section{Что мы собрали}

Мы брали по одному типу нейронов из таблицы~\ref{tab:types} и спрашивали, что он добавляет к вычислительной машине. Каждый раз правила были одни и те же: только то, что бывает в живом организме, и никаких часов. Теперь соберём части в одну машину и посмотрим, как они работают вместе.

Картина из главы~\ref{sec:neurontypes} подтвердилась и стала точнее. Огромная адресная сеть \nt{ГЛУТ}-нейронов несёт данные. \nt{ГАМК}-нейроны управляют потоком этих данных. А над сетью висят четыре широковещательных канала, каждый со своей ручкой: \nt{ДОФА} говорит, какие изменения весов оставить, \nt{СЕРО} держит общий уровень и, по гипотезе, горизонт планирования, \nt{НОРА} выбирает между поиском и решением, \nt{АЦЕТ} --- между записью и чтением (рис.~\ref{fig:machine}).

\begin{figure}[t]
\centering
\begin{tikzpicture}[>=Stealth,
  blk/.style={draw,very thick,rounded corners=3pt,align=center,font=\small},
  reg/.style={draw,very thick,double,double distance=1.2pt,circle,minimum size=1.25cm,inner sep=0pt,font=\scriptsize,fill=orange!15},
  bc/.style={->,thick,densely dashed,orange!85!black}]
% sensors, bus, actions
\node[blk,fill=green!8,text width=1.7cm] (S) at (0.9,0) {датчики\\\scriptsize вкл./выкл.};
\draw[very thick,rounded corners=4pt,fill=blue!5] (2.6,-1.35) rectangle (9.0,1.35);
\node[font=\small] at (5.8,0.95) {шина \nt{ГЛУТ}: данные};
\node[font=\scriptsize,align=center] at (4.3,0.25) {совпадения, задержки,\\логика (гл.~\ref{sec:glutcomputer})};
\node[font=\scriptsize,align=center] at (7.4,0.25) {память, код\\(гл.~\ref{sec:code}, \ref{sec:acet})};
\draw[thick,red!70!black,fill=red!8,rounded corners=2pt] (2.8,-1.15) rectangle (8.8,-0.55);
\node[font=\scriptsize,red!60!black] at (5.8,-0.85) {\nt{ГАМК} (гл.~\ref{sec:gaba}): запрет, выбор, деление, ритм};
\node[blk,fill=blue!5,text width=2.0cm] (A) at (10.9,0) {выбор\\действия};
\draw[->,very thick] (S) -- (2.6,0);
\draw[->,very thick] (9.0,0) -- (A);
\draw[->,very thick] (A) -- (12.6,0) node[right,font=\small] {мышцы};
\pic[scale=1.1] at (13.2,-0.5) {spring};
\pic[scale=1.15] at (0.4,0.85) {eyepic};
\pic[scale=1.05] at (1.35,0.85) {speaker};
% registers
\node[reg] (D) at (3.4,3.2) {\nt{ДОФА}};
\node[reg] (C) at (5.6,3.2) {\nt{СЕРО}};
\node[reg] (N) at (7.8,3.2) {\nt{НОРА}};
\node[reg] (H) at (10.0,3.2) {\nt{АЦЕТ}};
\node[font=\scriptsize,align=center,above=1pt] at (D.north) {ошибка $\delta$};
\node[font=\scriptsize,align=center,above=1pt] at (C.north) {уровень, $\tau_\gamma$};
\node[font=\scriptsize,align=center,above=1pt] at (N.north) {усиление $g$};
\node[font=\scriptsize,align=center,above=1pt] at (H.north) {запись $a$};
\draw[bc] (D.south) -- (3.4,1.35);
\draw[bc] (C.south) -- (5.6,1.35);
\draw[bc] (N.south) -- (7.8,1.35);
\draw[bc] (H.south) -- (10.0,2.1) -- (8.8,1.35);
\draw[bc] (H.south east) to[bend left=20] (A.north east);
\draw[->,thick] (2.85,1.35) -- (2.85,2.75) -- (D.west) node[pos=0.25,left,font=\scriptsize] {$V$};
\draw[->,thick,green!45!black] (0.4,3.2) node[left,black,font=\small] {награда $r$} -- (2.2,3.2) -- (D.west);
\pic[scale=0.9] at (-0.45,3.7) {juice};
\draw[->,thick,densely dotted,gray!50!black] ([yshift=0.45cm]D.north) to[bend left=18] node[above,font=\scriptsize,black] {неожиданность} ([yshift=0.45cm]N.north);
\end{tikzpicture}
\caption{Машина целиком. Шина \nt{ГЛУТ} несёт данные от датчиков к выбору действия; \nt{ГАМК} управляет потоком внутри неё. Двойные кружки --- широковещательные каналы; штриховые стрелки --- рассылка их сигналов всей сети, у каждого своя ручка. \nt{ДОФА} получает награду $r$ и ожидание $V$ от критика внутри шины (гл.~\ref{sec:dofa}); точечная стрелка --- гипотеза, что о неожиданности \nt{НОРА} узнаёт по необычно большим ошибкам (гл.~\ref{sec:nora}).}
\label{fig:machine}
\end{figure}

\section{Одна формула для всех ручек}

Ручки можно собрать в одну схематическую формулу. Это не новая модель, а сводка моделей глав~\ref{sec:glutcomputer}--\ref{sec:acet}: каждый символ взят из своей главы.
\begin{empheq}[box=\keybox]{equation}
\begin{aligned}
\tau\,\frac{\dd r_i}{\dd t} \;&=\; -r_i + F\Bigl(g\,\Bigl[\tfrac{1+a}{2}\,x_i + (1-a)\sum_j W_{ij}\,r_j - \sum_k M_{ik}\,q_k - \theta\Bigr]\Bigr),\\
\frac{\dd W_{ij}}{\dd t} \;&=\; \eta\,a\,\delta(t)\,e_{ij}(t),
\qquad \tau_e\,\frac{\dd e_{ij}}{\dd t} = -e_{ij} + r_i\,r_j,
\qquad \delta = r + \frac{\dd V}{\dd t} - \frac{V}{\tau_\gamma} .
\end{aligned}
\label{eq:machine}
\end{empheq}
Здесь $r_i$ --- частоты \nt{ГЛУТ}-нейронов, $x_i$ --- вход от датчиков, $W_{ij}\ge0$ --- их взаимные веса (глава~\ref{sec:glutcomputer}); $q_k$ --- частоты \nt{ГАМК}-нейронов, $M_{ik}\ge0$ --- сила их торможения (глава~\ref{sec:gaba}); $e_{ij}$ --- след, $\delta$ --- ошибка \nt{ДОФА} (глава~\ref{sec:dofa}); $\tau_\gamma$ --- горизонт, по гипотезе задаваемый \nt{СЕРО}; $g$ --- усиление \nt{НОРА} (глава~\ref{sec:nora}); $a$ --- уровень \nt{АЦЕТ} (глава~\ref{sec:acet}).

Посмотрите, чего в~\eqref{eq:machine} нет. Нет такта: всё записано дифференциальными уравнениями, и каждый нейрон меняется сам. Нет отдельной памяти: помнят те же $W_{ij}$, по которым идёт вычисление, и та же активность $r_i$. Нет персональных поправок для огромного большинства синапсов: их обучение --- произведение местного следа и одного общего числа; отдельный провод ошибки к каждому выходу есть лишь в цепи, которая учит движения (глава~\ref{sec:motorlearning}). И управляющих сигналов всего четыре рода: $\delta$, $\tau_\gamma$, $g$, $a$, на восемьдесят миллиардов нейронов. Каждый из них на самом деле не одно число, а небольшой жгут параллельных линий (утверждение~\ref{prop:bundle}), но линий в жгуте единицы.

\begin{keyblock}
\begin{proposition}[Архитектура]
\label{prop:architecture}
Машину мозга, насколько мы её сейчас понимаем, можно описать тремя слоями. Данные идут по адресной шине \nt{ГЛУТ}. Поток данных направляют, ограничивают и нормируют адресные \nt{ГАМК}-нейроны. Режим работы задают несколько широковещательных каналов, каждый --- небольшой жгут параллельных линий, общих для больших участков сети. Первые два слоя установлены твёрдо; распределение ролей между широковещательными каналами --- в основном гипотеза.
\end{proposition}
\end{keyblock}

\section{Все типы в одной таблице}

Таблица~\ref{tab:synthesis} сводит все типы нейронов книги. Четыре типа, которым не досталось своей главы, заняли в ней по строчке; двое из них нашли роль в главах о выходах машины: \nt{ГЛИЦ} --- в петлях спинного мозга (глава~\ref{sec:muscles}), \nt{АДРЕ} --- в химическом выходе (глава~\ref{sec:chemoutput}). Роль остальных двух в вычислении либо простая, либо неизвестна. Вещества, которые не задают тип нейрона, собраны в приложении~\ref{sec:othermediators}.

\begin{table}[t]
\centering
\footnotesize
\setlength{\tabcolsep}{3pt}
\renewcommand{\arraystretch}{1.15}
\begin{tabular}{>{\raggedright}p{1.3cm}>{\raggedright}p{1.0cm}>{\raggedright}p{3.9cm}>{\raggedright}p{1.5cm}>{\raggedright}p{1.8cm}>{\raggedright\arraybackslash}p{3.8cm}}
\toprule
Тип & Глава & Что вычисляет & Время & Шина & Что известно \\
\midrule
\nt{ГЛУТ} & \ref{sec:glutcomputer}, \ref{sec:code}, \ref{sec:sensors}, \ref{sec:motorlearning}, \ref{sec:dendrites} & данные: совпадения, задержки, логика, память, последовательности & мс & адресная & установлено \\
\nt{ГАМК} & \ref{sec:gaba}, \ref{sec:code}, \ref{sec:pain}, \ref{sec:motorlearning} & управление потоком: запрет, выбор, деление, устойчивость, ритм; ворота боли & мс & адресная & установлено; деление частоты --- вместе с фоновым шумом \\
\nt{СЕРО} & \ref{sec:serocomputer}, \ref{sec:pain}, \ref{sec:sleep} & регулятор уровня, сглаживание; горизонт $\tau_\gamma$; громкость боли; режимы сна & секунды & объём & самоторможение измерено; горизонт --- гипотеза \\
\nt{ДОФА} & \ref{sec:dofa}, \ref{sec:dofaalt} & ошибка предсказания $\delta$; медленный уровень --- цена времени & $0.1$\,с и минуты & широковещ. & $\delta$ измерено; расширения --- гл.~\ref{sec:dofaalt} \\
\nt{НОРА} & \ref{sec:nora}, \ref{sec:pain}, \ref{sec:chemoutput}, \ref{sec:sleep} & усиление $g$: искать или решать; неожиданность; «газ» для органов & секунды & широковещ. & рост отношения сигнала к шуму измерен; роль в поиске --- гипотеза \\
\nt{АЦЕТ} & \ref{sec:acet}, \ref{sec:muscles}, \ref{sec:chemoutput}, \ref{sec:sleep} & запись или чтение; скорость обучения; выход на мышцу; «тормоз» для органов & секунды & обе & три эффекта измерены; переключение режимов --- гипотеза \\
\nt{ГЛИЦ} & \ref{sec:muscles}, \ref{sec:pain} & отрицательный вес в спинном мозге и стволе: запрет мышцы-противника & мс & адресная & установлено \\
\nt{ГИСТ} & --- & глобальный сигнал «бодрствуй» & медленно & широковещ. & роль в вычислении не изучена \\
\nt{АДРЕ} & \ref{sec:chemoutput} & «газ» для всего тела через кровь; в мозге неизвестно & медленно & широковещ. & вне мозга установлено; роль в мозге неясна \\
\nt{АСПА} & --- & слабый положительный вес & мс & адресная & роль спорна \\
\bottomrule
\end{tabular}
\caption{Все типы нейронов книги: что вычисляет каждый и насколько это известно.}
\label{tab:synthesis}
\end{table}

\section{Как ручки работают вместе}

До сих пор каждая глава поворачивала одну ручку. Проследим теперь одно событие через всю машину (рис.~\ref{fig:timeline}). Животное давно выучило, что действие $A$ приносит сок. В какой-то момент мир меняется: $A$ перестаёт приносить сок, а приносит его действие $B$, которое животное почти не пробует.

\emph{Ошибка.} Сок после $A$ не пришёл, и \nt{ДОФА}-нейроны отвечают провалом, $\delta<0$ по формуле ошибки~\eqref{eq:ctd}. Синапсы, которые только что выбрали $A$, несут след и ослабевают.

\emph{Неожиданность.} Один провал --- обычная случайность. Но провалы повторяются, и ошибки становятся больше обычного. По гипотезе главы~\ref{sec:nora}, это и есть сигнал \nt{НОРА}: мир изменился. Усиление $g$ падает, выбор становится мягким, и шум чаще приводит к пробе $B$.

\emph{Запись.} По гипотезе главы~\ref{sec:acet}, после перемены \nt{АЦЕТ} поднимается и переводит сеть в режим записи: внутренние связи, хранящие старую привычку, ослаблены, вход от датчиков усилен, веса меняются легко.

\emph{Новая привычка.} Проба $B$ приносит сок, которого никто не ждал: всплеск, $\delta>0$, и синапсы, выбравшие $B$, усиливаются. Несколько таких попыток --- и ожидание $V$ догоняет новый мир, ошибки снова малы.

\emph{Возврат.} Неожиданность прошла, \nt{НОРА} возвращает высокое усиление, выбор снова жёсткий; \nt{АЦЕТ} опускается, сеть в режиме чтения пользуется выученным. Всё это время \nt{СЕРО}, по гипотезе, задаёт горизонт $\tau_\gamma$: насколько далёкий сок ещё стоит ждать.

\begin{figure}[t]
\centering
\begin{tikzpicture}[>=Stealth,x=1.1cm,y=0.95cm]
\foreach \y/\lab in {4/{$\delta$ (\nt{ДОФА})},3/{$g$ (\nt{НОРА})},2/{$a$ (\nt{АЦЕТ})},1/{выбор $B$}}{
  \draw[gray!60!black] (0,\y-0.35) -- (10,\y-0.35);
  \node[left,font=\scriptsize] at (0,\y) {\lab};}
\draw[densely dashed,gray!60!black] (2,0.4) -- (2,4.55) node[above,font=\scriptsize,black] {мир изменился};
\draw[densely dashed,gray!60!black] (5,0.4) -- (5,4.55) node[above,font=\scriptsize,black] {первый сок за $B$};
\pic[scale=0.85] at (2,5.25) {nojuice};
\pic[scale=0.85] at (5,5.25) {juice};
% delta: dips after the change, burst at the first juice for B, then back to zero
\draw[thick,red!70!black] (0,3.65) -- (2.3,3.65) -- (2.3,3.3) -- (2.5,3.3) -- (2.5,3.65) -- (3.2,3.65) -- (3.2,3.3) -- (3.4,3.3) -- (3.4,3.65) -- (4.1,3.65) -- (4.1,3.35) -- (4.3,3.35) -- (4.3,3.65) -- (5.0,3.65) -- (5.0,4.35) -- (5.2,4.35) -- (5.2,3.65) -- (5.8,3.65) -- (5.8,4.1) -- (6.0,4.1) -- (6.0,3.65) -- (6.6,3.65) -- (6.6,3.85) -- (6.8,3.85) -- (6.8,3.65) -- (10,3.65);
% gain: high, drops after surprise, recovers
\draw[thick,blue!60!black] (0,3.2) -- (3.0,3.2) -- (3.4,2.75) -- (5.8,2.75) -- (7.2,3.2) -- (10,3.2);
% ACh: low, rises after the change, falls after learning
\draw[thick,green!45!black] (0,1.75) -- (2.8,1.75) -- (3.3,2.25) -- (6.8,2.25) -- (7.8,1.75) -- (10,1.75);
% choice of B
\draw[thick,black] (0,0.7) -- (3.0,0.7) plot[domain=3:10,samples=40] (\x,{0.7+0.6/(1+exp(-(\x-5.3)*1.8))});
\draw[->,gray!70!black] (0,0.2) -- (10.2,0.2) node[below left,font=\scriptsize,black] {время};
\end{tikzpicture}
\caption{Одно событие, прослеженное через всю машину. Это качественная схема по моделям и гипотезам глав~\ref{sec:dofa}--\ref{sec:acet}, а не результат расчёта. После перемены \nt{ДОФА} отвечает провалами; повторные провалы, по гипотезе, поднимают сигнал неожиданности, \nt{НОРА} снижает усиление, \nt{АЦЕТ} включает запись; первый сок за $B$ даёт всплеск, выбор переходит на $B$, и ручки возвращаются.}
\label{fig:timeline}
\end{figure}

Заметьте, что ни одна ручка не знает, что делают другие. Каждая откликается на свои признаки --- ошибку, её необычную величину, неопределённость, --- и порядок действий складывается сам, как в асинхронной схеме, где каждый блок срабатывает, когда готовы его входы. Такую согласованную работу целиком, насколько нам известно, пока не проверяли: по отдельности ручки измерены, а их совместная работа --- гипотеза.

\section{Чем эта машина не похожа на компьютер}

\begin{table}[t]
\centering
\footnotesize
\setlength{\tabcolsep}{4pt}
\renewcommand{\arraystretch}{1.15}
\begin{tabular}{>{\raggedright}p{2.2cm}>{\raggedright}p{4.2cm}>{\raggedright\arraybackslash}p{6.6cm}}
\toprule
& Обычный компьютер & Мозг \\
\midrule
время & общий такт, около гигагерца & часов нет; каждый нейрон меняется сам, окна задают события и местные ритмы (гл.~\ref{sec:glutcomputer}, \ref{sec:gaba}, \ref{sec:code}) \\
элементы & надёжные двоичные вентили & аналоговые пороговые элементы; синапс теряет большую часть импульсов (гл.~\ref{sec:code}) \\
знак связи & любой & задаётся нейроном-отправителем (гл.~\ref{sec:neurontypes}) \\
память & отдельно от процессора & в тех же связях, что и вычисление, и в самоподдерживающейся активности (гл.~\ref{sec:glutcomputer}, \ref{sec:acet}) \\
обучение & обратное распространение ошибки & местные правила, одно широковещательное число $\delta$ (гл.~\ref{sec:dofa}) и провода ошибки к выходам цепи движений (гл.~\ref{sec:motorlearning}) \\
управление & регистры и шина команд & четыре широковещательных канала, каждый --- жгут из нескольких линий (гл.~\ref{sec:serocomputer}, \ref{sec:dofa}--\ref{sec:acet}) \\
энергия & десятки--сотни ватт на один процессор & около $20$ ватт на весь мозг, в среднем около импульса в секунду на нейрон (гл.~\ref{sec:code}) \\
\bottomrule
\end{tabular}
\caption{Мозг и обычный компьютер.}
\label{tab:computer}
\end{table}

Таблица~\ref{tab:computer} сводит главные отличия. Каждое из них --- не недостаток, который природа не успела исправить, а часть одного решения. Без общего такта нужны датчики «включения---выключения» и детекторы совпадений. Ненадёжным элементам нужна избыточность многих нейронов. Памяти, совмещённой с вычислением, нужен \nt{АЦЕТ}, чтобы запись не портила чтение. Обучению без обратных проводов нужен \nt{ДОФА} и шум. А энергетический предел в один импульс в секунду заставляет весь язык быть скупым и тратить импульсы на новости.

\section{Чего мы не знаем}

Честнее всего закончить эту сводку списком того, чего мы не знаем. Каждый пункт --- задача для инженера.

\emph{Код.} Мы знаем ёмкость импульсного канала и знаем, что получатель выбирает, какую часть сообщения читать (глава~\ref{sec:code}). Но насколько кора пользуется точным временем импульсов и есть ли у нейронов «слова», неизвестно.

\emph{Обучение через много слоёв.} Широковещательный сигнал учит медленно, и медленнее с каждым нейроном, влияющим на результат (утверждение~\ref{prop:broadcastcost}). Как тогда кора настраивает миллиарды весов в многослойных цепях, неизвестно; есть модели, в которых ошибку между слоями разносят пачки импульсов~\cite{Payeur2021}. Возможно, часть ответа --- в вычислениях внутри ветвей самого нейрона, где один нейрон ведёт себя как маленькая многослойная сеть: чтобы повторить ответ модели одного нейрона коры, искусственной сети нужно пять--восемь слоёв~\cite{Beniaguev2021}, а ветви нейронов человека умеют вычислять даже исключающее ИЛИ~\cite{Gidon2020}. Другой кандидат --- самообучение: если кора учится предсказывать собственные входы, ошибка возникает на месте, в каждом слое, и общий сигнал для этого не нужен (глава~\ref{sec:dofa})~\cite{Keller2018}.

\emph{Что на самом деле передают широковещательные каналы.} Для \nt{ДОФА} есть стандартная картина и шесть её расширений (глава~\ref{sec:dofaalt}); для \nt{НОРА} и \nt{АЦЕТ} --- правдоподобные гипотезы (главы~\ref{sec:nora}, \ref{sec:acet}); для \nt{СЕРО} --- в основном догадки.

\emph{Согласование во времени.} Мы видели, как вычислить функцию, отмерить время от события и нарезать поток на окна местным ритмом. Но как огромная сеть без общего такта согласует работу далёких областей, понятно лишь отчасти.

\emph{Энергия.} Двадцать ватт на всё. Мы понимаем, почему код должен быть разреженным (утверждение~\ref{prop:economy}), но построить машину, которая делала бы то же самое при такой же мощности, пока не умеет никто~\cite{MehonicKenyon2022}.

Мозг --- вычислительная машина, собранная не инженером, но по законам, которые любой инженер узнает: RC-цепи, пороги, обратные связи, фильтры, шины и широковещательные регистры. Её схему мы прочитали лишь частично. Дочитывать её будут те, кто умеет читать схемы.

\section{Что дальше в книге}

Эта глава собрала вычислительное ядро машины. Остальные главы выходят за его пределы. Главы~\ref{sec:sensors} и~\ref{sec:recognition} --- о входах: как датчики превращают свет, звук и молекулы в импульсы и как машина узнаёт в них предметы. Главы~\ref{sec:muscles}--\ref{sec:chemoutput} --- о выходах и о том, что их обслуживает: мышцы, боль как сигнал тревоги, обучение движениям по отдельным проводам ошибки и медленный химический выход в тело. Глава~\ref{sec:debugging} --- о том, как машину отлаживают снаружи, в том числе интерфейсами мозг---компьютер. Главы~\ref{sec:eetypes} и~\ref{sec:dendrites} ещё раз сортируют нейроны, теперь по электрической функции и по числу входов. Глава~\ref{sec:sleep} --- о том, что машина делает, когда отключена от мира, а главы~\ref{sec:livingmachine} и~\ref{sec:dishbrain} --- о машинах, собранных из живых нейронов на чипе.

\section{Задачи}

Задачи этой главы соединяют результаты разных глав: каждая опирается по меньшей мере на две из них.

\begin{exercise}[\lvlA]
\label{ex:10:1}
\emph{Шина и регистры.} Сравним потоки данных и управления в машине~\eqref{eq:machine}.
\par\noindent(а)~Оцените сверху поток данных по шине \nt{ГЛУТ}: все $8.6\cdot10^{10}$ нейронов со средней частотой~\eqref{eq:energy}, прочитанные с точностью $1\ms$ по~\eqref{eq:spikeentropy}.
\par\noindent(б)~Оцените поток управления: четыре широковещательных канала, каждый --- жгут примерно из пяти линий (утверждение~\ref{prop:bundle}); значение каждой линии меняется за секунду и различимо на четырёх уровнях. Во сколько раз поток данных больше?
\par\noindent(в)~Сколько времени понадобилось бы каналам управления, чтобы передать столько, сколько шина данных передаёт за одну секунду?
\par\noindent(г)~Пусть широковещательные нейроны \nt{ДОФА}, \nt{СЕРО}, \nt{НОРА} и \nt{АЦЕТ} (таблица~\ref{tab:types}) работают с частотой $5$ импульсов в секунду, а их импульс стоит в $10$--$100$ раз дороже обычного, потому что выходов у них больше. Какую мощность и какую долю энергии сигналов~\eqref{eq:energy} забирает управление? С каким бытовым прибором её можно сравнить?
\end{exercise}

\begin{exercise}[\lvlB]
\label{ex:10:2}
\emph{Две ручки на одной петле.} В~\eqref{eq:machine} усиление $g$ и уровень $a$ умножают одну и ту же скобку. Возьмите петлю \nt{ГЛУТ}--\nt{ГАМК}~\eqref{eq:ei} и поставьте ручки так, как в~\eqref{eq:machine}: $\tau_E\,\dd E/\dd t=-E+g\bigl[(1-a)\,w_{EE}E-w_{EI}I+h\bigr]$, $\tau_I\,\dd I/\dd t=-I+w_{IE}E$; \nt{ГАМК}-нейроны ручками не управляются, передаточная характеристика в рабочей области линейна.
\par\noindent(а)~Выведите, во что превращаются условия утверждения~\ref{prop:ei}.
\par\noindent(б)~Числа рис.~\ref{fig:ei}: $w_{EE}=1.5$, $w_{EI}w_{IE}=2$, $\tau_E=10\ms$, $\tau_I=2\ms$. При $a=0$ найдите усиление $g$, при котором петля начинает колебаться, и частоту колебаний на этой границе.
\par\noindent(в)~Пачка \nt{НОРА} в момент решения поднимает $g$ до $6$. При каком наименьшем $a$ петля остаётся устойчивой? Что будет при $a=0$?
\par\noindent(г)~Что это значит для последовательности на рис.~\ref{fig:timeline}: можно ли крутить ручки \nt{НОРА} и \nt{АЦЕТ} независимо?
\end{exercise}

\begin{exercise}[\lvlB]
\label{ex:10:3}
\emph{Откуда берётся цена широковещания.} Выведите утверждение~\ref{prop:broadcastcost} из~\eqref{eq:perturbation}. Выходы $N$ нейронов $y_k=f(u_k)+\xi_k$ с независимыми гауссовыми $\xi_k$ дисперсии $\sigma^2$ влияют на награду; при малом шуме $R\approx R_0+\sum_kR'_k\xi_k$, и \nt{ДОФА} рассылает $\delta=R-R_0$.
\par\noindent(а)~Покажите, что среднее поправки одного синапса $\delta\,\xi_jx_j$ равно $\sigma^2R'_jx_j$.
\par\noindent(б)~Вычислите её дисперсию.
\par\noindent(в)~Пусть все $R'_k$ одинаковы. Найдите отношение сигнала к шуму одной попытки и число попыток, нужное, чтобы узнать знак поправки с заданной надёжностью. Как оно зависит от $N$?
\par\noindent(г)~Проверьте (в) моделированием на Python при $N=1$, $4$, $16$, $64$ ($5\cdot10^4$ попыток).
\end{exercise}

\begin{exercise}[\lvlB]
\label{ex:10:4}
\emph{Проектирование: связать звук и вспышку.} Событие одновременно звучит и вспыхивает. Звуковой вход приходит на шину через $L_a=5\ms$. Световой --- через $L_v=L_0+t_1$, где $L_0=30\ms$, а $t_1$ --- задержка первого импульса~\eqref{eq:latency} при $\tau=10\ms$ и контрасте от $U=1.2\theta$ до $U=4\theta$. На каждом входе есть и фоновые импульсы, независимые, по $5$ в секунду. Постройте из \nt{ГЛУТ}-нейронов детектор «звук и вспышка от одного события»: линия задержки $d$ на звуковом входе и детектор совпадений~\eqref{eq:window} с порогом $\theta_D$, весом $w$ и постоянной времени $\tau_D$.
\par\noindent(а)~Выберите $d$ и окно так, чтобы связывалось любое событие во всём диапазоне контрастов, а окно было наименьшим.
\par\noindent(б)~Выберите $\tau_D$ при $\theta_D=1.5w$.
\par\noindent(в)~Найдите частоту ложных связываний от фоновых импульсов. Как она зависит от фоновых частот и от разброса задержек? Что это говорит о пользе утверждения~\ref{prop:economy} и о задаче согласования во времени из этой главы?
\end{exercise}

\begin{exercise}[\lvlB]
\label{ex:10:5}
\emph{Моделирование: кто замечает перемену и кто пишет.} Проверим разделение труда из главы~\ref{sec:acet}. Критик следит за ожиданием $s$ по попыткам: в каждой попытке он видит $y=s+\text{шум}$ с дисперсией $R=1$; с вероятностью $h=1/200$ перед попыткой $s$ скачком принимает новое значение из нормального распределения с дисперсией $J^2=4$. Ошибка критика $\delta=y-\hat s$ --- та же ошибка предсказания, что в~\eqref{eq:ctd}. Напишите на Python три критика и для каждого измерьте среднеквадратичную ошибку $\langle(s-\hat s)^2\rangle$ за $10^5$ попыток (десять прогонов).
\par\noindent(а)~Постоянная скорость: $\hat s\leftarrow\hat s+\alpha\delta$. Найдите лучшее $\alpha$.
\par\noindent(б)~Только \nt{АЦЕТ}: фильтр Калмана в попытках, $P\leftarrow P+q$, $K=P/(P+R)$, $\hat s\leftarrow\hat s+K\delta$, $P\leftarrow(1-K)P$, с лучшим постоянным $q$. Объясните результат утверждением~\ref{prop:kalman}.
\par\noindent(в)~\nt{НОРА} и \nt{АЦЕТ}: тот же фильтр при $q=0$, но с детектором неожиданности. Он копит $E\leftarrow\lambda E+\delta^2/(P+R)$ и, когда $E$ превышает порог $c$, поднимает $P$ до $J^2$ (сеть переходит в режим записи) и обнуляет $E$. Подберите $\lambda$ и $c$.
\par\noindent(г)~Найдите нижнюю границу: критик, который точно знает моменты скачков. Что мешает критику (в) до неё дотянуться?
\end{exercise}

\begin{exercise}[\lvlA]
\label{ex:10:6}
\emph{Сколько попыток на смену привычки.} Сеть выучила ожидания $Q_A=0.8$, $Q_B=0.2$ и выбирает по~\eqref{eq:softmax} с $\sigma=1$. Мир меняется: теперь $A$ приносит сок с вероятностью $0.2$, $B$ --- с вероятностью $0.8$. Пока сеть выбирает $A$, ожидание обновляется как $Q_A\leftarrow Q_A+\alpha\,(r-Q_A)$; $B$ почти не пробуется, и $Q_B$ остаётся прежним.
\par\noindent(а)~Найдите среднее $Q_A$ после $n$ выборов $A$.
\par\noindent(б)~При $g=8$ найдите, после скольких выборов $A$ вероятность выбрать $A$ упадёт до $0.6$, при $\alpha=0.1$ и при $\alpha=0.5$ (\nt{АЦЕТ} поднял скорость обучения, глава~\ref{sec:acet}).
\par\noindent(в)~То же, если \nt{НОРА} одновременно опустила $g$ до $2$.
\par\noindent(г)~Почему без проб $B$ вероятность выбрать $A$ никогда не опустится ниже половины, и какая ручка на рис.~\ref{fig:timeline} за это отвечает?
\end{exercise}

\begin{exercise}[\lvlC]
\label{ex:10:7}
\emph{Жгут вместо провода: сколько линий нужно.} В разделе «Чего мы не знаем» остался вопрос, как настраиваются многослойные цепи. Посмотрим, спасает ли дело жгут из $K$ линий (утверждение~\ref{prop:bundle}).
\par\noindent(а)~Модель: $N$ выходов $y_k=w_k+\xi_k$, цель --- $y_k=0$, награда $R=-\sum_ky_k^2$, $w_k(0)=1$. Нейроны разбиты на $K$ групп по $G=N/K$; линия $l$ рассылает своей группе ошибку $\delta_l$ только её выходов, при малом шуме $\delta_l\approx-2\sum_{k\in l}w_k\xi_k$, и $w_k\leftarrow w_k+\eta\,\delta_l\xi_k$. Выведите рекуррентное соотношение для среднего $S=\sum_{k\in l}w_k^2$, найдите лучшее $\eta$ и покажите, что число попыток до $S=\varepsilon S(0)$ равно примерно $(G+2)\ln(1/\varepsilon)$.
\par\noindent(б)~Проверьте моделированием при $N=8$, $32$, $128$ и $K=1$, $4$, $N$, $\sigma=0.1$, $\varepsilon=0.01$.
\par\noindent(в)~Оцените по (а) время обучения цепи из $10^6$ нейронов при жгуте из пяти линий и одной попытке за три секунды.
\par\noindent(г)~Открытый вопрос. Каждой линии жгута нужна своя группа получателей --- это уже адресная, а не широковещательная передача. Сколько линий нужно, чтобы время обучения не зависело от $N$, и откуда могли бы взяться столь многие сигналы ошибки? Сравните с кандидатами из раздела «Чего мы не знаем» и предложите опыт или модель, которые различили бы их.
\end{exercise}

\chapter{Входы машины: как к ней подключены глаза, уши и другие датчики}
\chaptermark{Входы машины}
\label{sec:sensors}

\section{Датчик как преобразователь}

На рис.~\ref{fig:machine} данные входили в машину слева, из блока «датчики». Теперь откроем этот блок. Для инженера каждый орган чувств --- это \emph{преобразователь} с аналоговым входным каскадом и аналого-цифровым преобразователем в конце, только цифры на выходе --- не числа, а импульсы.

Схема у всех органов чувств одна и та же (рис.~\ref{fig:transducer}). Физическая величина --- свет, давление, колебание, молекула --- действует на белок-рецептор в мембране чувствительной клетки. Этот белок работает как затвор: он сам или через короткую цепочку реакций открывает ионный канал. Возникает \emph{рецепторный ток}, а с ним --- аналоговое напряжение на мембране. Дальше напряжение проходит регулировку усиления и попадает в кодер импульсов --- уже знакомый интегрирующий нейрон~\eqref{eq:glutneuron}, который превращает уровень в частоту и моменты импульсов. Выход почти всегда один и тот же: чувствительные нейроны глаза, уха, обоняния и кожи выбрасывают глутамат, так что входы подключаются прямо к шине \nt{ГЛУТ}.

\begin{figure}[t]
\centering
\begin{tikzpicture}[>=Stealth,
  blk/.style={draw,very thick,rounded corners=2pt,align=center,font=\footnotesize,minimum height=1.0cm,text width=1.9cm,fill=blue!5}]
\node[align=center,font=\small] (P) at (0.1,0) {свет, звук,\\давление,\\молекула};
\node[blk] (R) at (3.0,0) {рецептор-\\затвор};
\node[blk] (G) at (6.5,0) {регулировка\\усиления};
\node[blk] (E) at (10.0,0) {кодер\\импульсов};
\node[align=center,font=\small] (B) at (13.4,0) {шина\\\nt{ГЛУТ}};
\draw[->,very thick] (P) -- (R);
\foreach \ic/\xx in {sun/-1.1,speaker/-0.3,press/0.5,molecule/1.3}{\pic[scale=0.85] at (\xx,1.05) {\ic};}
\draw[->,very thick] (R) -- node[above,font=\scriptsize] {ток} (G);
\draw[->,very thick] (G) -- node[above,font=\scriptsize] {уровень} (E);
\draw[->,very thick,red!70!black] (E) -- node[above,font=\scriptsize,black] {импульсы} (B);
\draw[decorate,decoration={brace,amplitude=5pt,mirror},gray!60!black] (1.8,-0.8) -- (7.7,-0.8) node[midway,below=5pt,font=\scriptsize,black] {аналоговая часть};
\draw[decorate,decoration={brace,amplitude=5pt,mirror},gray!60!black] (8.8,-0.8) -- (11.2,-0.8) node[midway,below=5pt,font=\scriptsize,black] {импульсы};
\end{tikzpicture}
\caption{Любой орган чувств как цепь преобразований. Белок-рецептор открывает канал и даёт ток; регулировка усиления приспосабливает его к обстановке; кодер~\eqref{eq:glutneuron} превращает уровень в импульсы, которые уходят на шину \nt{ГЛУТ}. В глазу и ухе первые каскады вообще не выдают импульсов: сигнал остаётся аналоговым, пока не приходится передавать его далеко.}
\label{fig:transducer}
\end{figure}

Одна деталь хорошо знакома электронщику. В глазу и в ухе самые первые клетки импульсов не выдают вовсе: они передают соседям плавное напряжение, как аналоговый каскад на плате. Импульсы появляются только там, где сигнал надо везти далеко. Аналоговый сигнал на расстоянии миллиметров затухает и зашумляется, а импульс, как мы видели в главе~\ref{sec:neurontypes}, доходит до конца провода той же высоты. Оцифровывают там, где начинается длинная линия.

\section{На пределе физики}

Датчики мозга работают у физических пределов чувствительности. Датчик света в темноте отвечает на \emph{один фотон}: поглощение одной частицы света даёт ток около пикоампера, заметный на фоне собственного шума~\cite{Baylor1979}. Датчик звука замечает смещение своего чувствительного пучка меньше чем на нанометр~\cite{Hudspeth1989}.

Когда света мало, точность ограничивает не датчик, а сам свет: фотоны приходят случайно, пуассоновским потоком. Если за время накопления $T$ в датчик попадает в среднем $N=\Phi T$ фотонов, их число дрожит на $\sqrt N$. Чтобы прибавка $\Delta N$ была заметна, она должна быть в несколько раз больше этого дрожания, и различимая доля изменения яркости равна
\begin{empheq}[box=\keybox]{equation}
\frac{\Delta I}{I} \;\approx\; \frac{k}{\sqrt{\Phi\,T}} .
\label{eq:photonnoise}
\end{empheq}
Это та же формула, что и для числа импульсов~\eqref{eq:rateerror}: дробовой шум фотонов и дробовой шум импульсов подчиняются одному закону. В темноте глаз может улучшить точность только одним способом --- дольше копить фотоны, и время накопления действительно растёт до десятых долей секунды. Поэтому в сумерках мы видим медленнее.

\section{Динамический диапазон: регулировка усиления}

Самая трудная инженерная задача у датчиков --- диапазон. Освещённость от звёздной ночи до солнечного снега меняется примерно на десять порядков, интенсивность звука от порога слышимости до болевого --- на двенадцать. А частота импульсов меняется от нуля до нескольких сотен в секунду, меньше чем на три порядка. Линейно такой вход в такой выход не уложить.

Решение то же, что у любого приёмника: \emph{автоматическая регулировка усиления}. Ответы чувствительных клеток глаза хорошо описываются формулой
\begin{empheq}[box=\keybox]{equation}
r \;=\; r_{\max}\,\frac{I}{I+\sigma} ,
\label{eq:nakarushton}
\end{empheq}
где $\sigma$ --- яркость, при которой ответ достигает половины~\cite{NakaRushton1966}. Сама по себе~\eqref{eq:nakarushton} охватывает лишь два-три порядка. Но $\sigma$ не постоянна: при долгой работе на ярком фоне она растёт вслед за средней яркостью. Кривая ответа сдвигается вдоль логарифмической оси яркости и всякий раз ставит свой крутой участок туда, где сейчас находится вход (рис.~\ref{fig:adaptation}). Это то же деление на общую активность, что и шунтирующее торможение главы~\ref{sec:gaba}~\cite{Carandini2012}, только знаменатель здесь --- средняя яркость за последние секунды.

\begin{figure}[t]
\centering
\begin{tikzpicture}[>=Stealth,x=1.25cm,y=2.8cm]
\draw[->,gray!70!black] (-2,0) -- (6.4,0) node[below left,font=\small,black] {$\log_{10} I$};
\draw[->,gray!70!black] (-2,0) -- (-2,1.15) node[left,font=\small,black] {$r/r_{\max}$};
\foreach \u in {-2,0,2,4,6}{\draw[gray!60!black] (\u,-0.02) -- (\u,0.02); \node[below,font=\scriptsize] at (\u,0) {$\u$};}
\draw[densely dashed,gray!60!black] (-2,0.5) -- (6.2,0.5);
\foreach \s/\c/\lab in {0/blue!60!black/{тёмный фон},2/green!45!black/{средний},4/red!70!black/{яркий}}{
  \pgfmathsetmacro{\lo}{max(-2,\s-3)}
  \draw[very thick,\c] (-2,0) -- (\lo,0) plot[domain=\lo:6.2,samples=80] (\x,{1/(1+10^(\s-\x))});
  \node[font=\scriptsize,\c,anchor=west] at (\s+0.15,0.42) {\lab};}
\foreach \s/\ic in {0/moon,2/cloud,4/sun}{\pic at (\s+0.6,0.64) {\ic};}
\end{tikzpicture}
\caption{Регулировка усиления датчика света по~\eqref{eq:nakarushton}. Каждая кривая охватывает два-три порядка яркости; при переходе на более яркий фон $\sigma$ растёт, и кривая сдвигается по логарифмической оси. Вместе сдвигающиеся кривые покрывают весь диапазон.}
\label{fig:adaptation}
\end{figure}

У регулировки есть цена, знакомая по главе~\ref{sec:code}: датчик сообщает не яркость, а яркость \emph{по отношению к фону}. Постоянный вход постепенно перестаёт вызывать ответ, а каждое изменение вызывает. Входы машины с самого начала говорят о переменах, а не об уровнях, и в этом та же экономия импульсов, что и в утверждении~\ref{prop:economy}~\cite{Barlow1961}. Отсюда и включающие и выключающие датчики главы~\ref{sec:glutcomputer}~\cite{Schiller1992}: одни сообщают, что стало ярче, другие --- что темнее.

Инженеры повторили этот приём в \emph{событийных камерах}. Такая камера не снимает кадры по часам: каждый её пиксель сам, без общего такта, выдаёт событие «ярче» или «темнее», когда логарифм яркости изменился на заданную долю~\cite{Lichtsteiner2008}. Это ровно двухпроводный код включения---выключения из главы~\ref{sec:glutcomputer}, воплощённый в кремнии, и он даёт камере диапазон и быстроту, недоступные обычной матрице~\cite{Gallego2022}.

\section{Глаз: сжатие до провода}

В глазу около $10^8$ датчиков света~\cite{Curcio1990}, а выходных нейронов, которые везут изображение в мозг, --- около миллиона. Прежде чем сигнал покинет глаз, его сжимают примерно в сто раз. Главные приёмы сжатия нам знакомы. Каждый выходной нейрон отвечает на разницу между яркостью в маленьком пятне и яркостью вокруг него --- это вычитание соседей через торможение, как в главе~\ref{sec:gaba}. Ровные участки картинки почти ничего не стоят, а края и перемены передаются. Разные выходные нейроны специализированы: одни сообщают о посветлении, другие о потемнении, третьи о движении в определённую сторону (рис.~\ref{fig:direction}); у мыши таких выходных каналов насчитали больше тридцати~\cite{Baden2016}.

Какова пропускная способность получившегося кабеля? У морской свинки по записям выходных нейронов измерили около $875$ тысяч бит в секунду на $\sim10^5$ таких нейронов; пересчёт на миллион выходных нейронов человека даёт порядка $10^7$ бит в секунду~\cite{Koch2006} --- столько же, сколько старый сетевой кабель на десять мегабит. При средней частоте в несколько импульсов в секунду на нейрон это несколько бит на импульс, как и получалось в главе~\ref{sec:code}.

\section{Ухо: банк фильтров}

Ухо решает другую задачу: звук надо разложить по частотам. Инженер поставил бы банк полосовых фильтров, и ухо делает именно это механически. Звуковое колебание бежит по упругой перегородке, свойства которой плавно меняются вдоль её длины, и каждый участок резонирует на свою частоту. Датчики, сидящие на участке, отвечают только на свою полосу (рис.~\ref{fig:filterbank}). Частоты разложены по месту примерно логарифмически: каждой октаве достаётся похожая доля датчиков. Номер сработавшего датчика --- это частота, код местом из главы~\ref{sec:code}.

\begin{figure}[t]
\centering
\begin{tikzpicture}[>=Stealth,x=1.35cm,y=2.2cm]
\draw[->,gray!70!black] (-0.3,0) -- (7.6,0) node[above left,font=\small,black] {частота, Гц};
\draw[->,gray!70!black] (-0.3,0) -- (-0.3,1.15) node[left,font=\small,black] {ответ};
\foreach \k/\f in {0/125,1/250,2/500,3/1000,4/2000,5/4000,6/8000}{
  \draw[gray!60!black] (\k,-0.02) -- (\k,0.02); \node[below,font=\scriptsize] at (\k,0) {\f};
  \draw[thick,blue!60!black] plot[domain={max(\k-1.2,-0.3)}:{\k+0.7},samples=60] (\x,{1/(1+((\x-\k)/(\x<\k?0.45:0.2))^4)});}
% a piano keyboard: seven white keys per octave, like the octaves on the axis
\foreach \i in {0,...,48}{\draw[gray!50!black,fill=white,line width=0.3pt] ({\i/7},-0.44) rectangle ({(\i+1)/7},-0.26);}
\foreach \o in {0,...,6}{\foreach \j in {0,1,3,4,5}{\fill[black] ({\o+(\j+1)/7-0.035},-0.35) rectangle ({\o+(\j+1)/7+0.035},-0.26);}}
\end{tikzpicture}
\caption{Ухо как банк полосовых фильтров, схема. Каждая кривая --- ответ датчиков одного участка на звук разной частоты; центральные частоты идут через октаву, как клавиши на логарифмической шкале. У настоящих фильтров крутой высокочастотный склон и пологий низкочастотный. Внизу --- клавиатура: на ней, как и в ухе, каждой октаве достаётся одинаковое место.}
\label{fig:filterbank}
\end{figure}

Резонаторы уха не пассивны. Часть датчиков сама раскачивает перегородку в такт звуку и усиливает слабые колебания в тысячи раз по амплитуде (на $66$--$76$\,дБ), а с ростом громкости усиление падает~\cite{Ruggero1997,Robles2001}. Для инженера это усилитель с положительной обратной связью, поставленный у самой границы самовозбуждения, --- та же жизнь на грани, что и у сети главы~\ref{sec:gaba}: у границы система особенно чувствительна к малым сигналам. Сжатие громкости здесь делается тем же усилителем: тихое усиливается сильно, громкое --- слабо.

У уха есть и второй код. На низких частотах импульсы нейронов идут в фазе со звуком: нейрон срабатывает на определённом участке каждой волны, хотя и не на каждой волне. У морской свинки привязка к фазе начинает слабеть около $600$\,Гц и ещё заметна примерно до $3.5$\,кГц~\cite{PalmerRussell1986}. Так в моментах импульсов сохраняется точное время звука, а из него --- разница прихода звука в два уха, по которой сеть главы~\ref{sec:glutcomputer} находит направление~\cite{Grothe2010}. На высоких частотах импульсы за волной не успевают, и остаётся только код местом. Одно ухо пользуется обоими кодами главы~\ref{sec:code}: временем там, где нейроны успевают, и местом там, где нет.

\section{Остальные датчики}

\begin{table}[t]
\centering
\footnotesize
\setlength{\tabcolsep}{3pt}
\renewcommand{\arraystretch}{1.15}
\begin{tabular}{>{\raggedright}p{1.9cm}>{\raggedright}p{2.6cm}>{\raggedright}p{4.0cm}>{\raggedright\arraybackslash}p{4.6cm}}
\toprule
Чувство & Что измеряет & Как устроен вход & Приём из книги \\
\midrule
зрение & свет & $\sim10^8$ датчиков, сжатие в $\sim100$ раз, $\sim10^7$ бит/с & регулировка усиления, вычитание соседей, два провода, направление движения \\
слух & давление воздуха & механический банк фильтров, активный усилитель & код местом и код временем, задержки \\
осязание & давление, вибрация & быстро и медленно привыкающие датчики & пара «фильтр верхних частот --- фильтр нижних» \\
температура & тепло & отдельные датчики тепла и холода & двухпроводный код \\
обоняние & молекулы & $\sim400$ типов рецепторов, у каждого датчика один тип & комбинаторный код: запах --- набор сработавших типов \\
вкус & вещества в пище & пять качеств: сладкое, горькое, солёное, кислое, умами & меченые линии; часть клеток выбрасывает АТФ или серотонин, а не глутамат \\
положение тела & растяжение мышц & датчики длины и скорости растяжения & обратная связь для регулятора движения \\
\bottomrule
\end{tabular}
\caption{Как подключены датчики. Все устройства в третьем столбце измерены; правый столбец связывает их с приёмами из предыдущих глав.}
\label{tab:sensors}
\end{table}

Остальные чувства собраны в таблице~\ref{tab:sensors}, и почти в каждой строке узнаётся приём из предыдущих глав. Осязание пользуется парой фильтров: одни датчики давления отвечают на прикосновение и сразу замолкают, другие держат ответ, пока давление есть. Температуру передают два провода, отдельные для тепла и для холода. Самый необычный вход --- обоняние. Типов обонятельных рецепторов у человека около четырёхсот, и каждый обонятельный нейрон несёт только один тип~\cite{Mombaerts2004}. Запах активирует не один тип, а набор, и сообщение --- это \emph{какие} типы сработали. Для инженера это двоичный вектор длиной около четырёхсот, у которого число возможных значений астрономическое.

\begin{keyblock}
\begin{proposition}[Входы машины]
\label{prop:sensors}
Каждое чувство --- преобразователь, который работает у физического предела чувствительности, сжимает огромный динамический диапазон автоматической регулировкой усиления и передаёт по шине \nt{ГЛУТ} прежде всего \emph{перемены}. Для этого датчики пользуются теми же приёмами, что и сама машина: двухпроводным кодом, кодом местом, кодом временем и делением на фон. Устройство датчиков измерено; то, что их коды подстроены под наибольшую информацию на импульс, --- хорошо обоснованная гипотеза.
\end{proposition}
\end{keyblock}

Входы, как и всё остальное, работают асинхронно. Каждый датчик выдаёт импульс, когда у него есть что сообщить, а разные чувства приходят с разными задержками: звук --- через миллисекунды, свет --- через десятки миллисекунд. Как машина сводит их в одну картину мира, если общего такта нет, --- одна из задач согласования во времени из главы~\ref{sec:synthesis}.

\section{Задачи}

\begin{exercise}[\lvlA]
\label{ex:11:1}
\emph{Сколько копить фотоны.} В сумерках в датчик света попадает $\Phi=100$ фотонов в секунду. Прибавка считается заметной, если она в $k=3$ раза больше дрожания, как в~\eqref{eq:photonnoise}.
\par\noindent(а)~Сколько времени нужно одному датчику, чтобы заметить изменение яркости на $10\,\%$?
\par\noindent(б)~Время накопления ограничено $0.2\,\text{с}$. Какое наименьшее изменение яркости заметит один датчик?
\par\noindent(в)~Сколько независимых датчиков нужно сложить, чтобы заметить $10\,\%$ за $0.2\,\text{с}$? Во сколько раз при этом падает разрешение по каждой оси, если датчики складываются квадратным пятном? Какой приём главы~\ref{sec:code} здесь повторяется?
\par\noindent(г)~Днём $\Phi=10^4$ в секунду. Сколько времени нужно одному датчику для тех же $10\,\%$?
\end{exercise}

\begin{exercise}[\lvlA]
\label{ex:11:2}
\emph{Сколько бит в импульсе.} (а)~Глаз передаёт $\sim10^7$ бит/с по $10^6$ выходным нейронам со средней частотой $3$--$5$ импульсов в секунду. Какую долю предельной ёмкости~\eqref{eq:spikeentropy} при $\Delta t=1\ms$ он использует?
\par\noindent(б)~Сколько бит в секунду остаётся на один датчик света после сжатия?
\par\noindent(в)~Обонятельный код: запах включает $20$ из $\approx400$ типов рецепторов. Сколько бит несёт такой набор, если все наборы равновероятны? Сколько общих типов в среднем у двух случайных запахов и какова вероятность, что общих типов три и больше? пять и больше?
\end{exercise}

\begin{exercise}[\lvlB]
\label{ex:11:3}
\emph{Что умеет одна кривая~\eqref{eq:nakarushton}.} (а)~В каком диапазоне яркостей ответ растёт от $10\,\%$ до $90\,\%$ от $r_{\max}$? Сколько это порядков?
\par\noindent(б)~Найдите наибольшую крутизну $\dd r/\dd\log_{10}I$ и яркость, при которой она достигается.
\par\noindent(в)~Сколько положений сдвигающейся кривой нужно, чтобы покрыть десять порядков яркости? двенадцать порядков громкости?
\par\noindent(г)~Пусть регулировка усиления довела $\sigma$ до яркости фона $I_b$. Докажите, что ответ на ступеньку $I_b\to I_b+\Delta I$ зависит только от контраста $\Delta I/I_b$ (закон Вебера). Пусть выход --- пуассоновский счёт импульсов за $T$, и ступенька заметна, когда прибавка счёта равна его разбросу. Найдите наименьший заметный контраст для одного нейрона при $r_{\max}=200$ в секунду и $T=0.1\,\text{с}$. Сколько независимых нейронов нужно для $2\,\%$? Каков предел при корреляции $c=0.1$~\eqref{eq:correlated}, если корреляция похожа на сигнал?
\end{exercise}

\begin{exercise}[\lvlB]
\label{ex:11:4}
\emph{Предел кода временем в ухе.} Разница прихода звука в два уха у человека --- до $0.22\,\text{м}/343\,\text{м/с}$ (глава~\ref{sec:glutcomputer}).
\par\noindent(а)~Импульсы идут в фазе с тоном частоты $f$, поэтому по ним разница прихода известна лишь с точностью до периода $1/f$. Докажите, что направление определяется однозначно, только если $f<1/(2\,\delta t_{\max})$, и вычислите эту частоту. Сравните с частотой, до которой импульсы идут в фазе.
\par\noindent(б)~Сколько детекторов в ряду~\eqref{eq:itd} сработает на тон $500$~Гц и на тон $2$~кГц от одного источника?
\par\noindent(в)~Каждый импульс дрожит на $0.5\ms$, а различать нужно $20$~мкс. Сколько независимых импульсов надо усреднить? Сколько нейронов, если каждый за $50\ms$ выдаёт импульсы с частотой $100$ в секунду?
\end{exercise}

\begin{exercise}[\lvlB]
\label{ex:11:5}
\emph{Проектирование: событийная камера на кабель глаза.} Нужна событийная камера из $10^6$ пикселей с выходом $10^7$ бит/с, как у глаза. Пиксель выдаёт событие «ярче» или «темнее», когда $\ln I$ изменился на $C$ с прошлого события. Событие несёт адрес пикселя и знак; время события --- момент его прихода, отдельно оно не передаётся.
\par\noindent(а)~Сколько бит в событии и сколько событий в секунду пропускает выход?
\par\noindent(б)~Камера должна замечать перепады яркости в $20\,\%$. Выберите $C$. Сколько событий даёт пиксель, через который проходит край с отношением яркостей $4$?
\par\noindent(в)~Край длиной $500$ пикселей движется по кадру. С какой наибольшей скоростью он может двигаться без переполнения выхода? Можно ли, повышая $C$, пропустить край, бегущий $1000$ пикселей в секунду?
\par\noindent(г)~Сравните с обычной камерой, которая через тот же выход передаёт полные кадры по $8$ бит на пиксель: сколько кадров в секунду и на сколько пикселей сдвинется край из (в) между кадрами? Какой приём главы~\ref{sec:glutcomputer} даёт событийной камере преимущество?
\end{exercise}

\begin{exercise}[\lvlB]
\label{ex:11:6}
\emph{Моделирование: регулировка усиления.} Для этой главы файла в \texttt{models/} нет; напишите модель на Python. Датчик отвечает по~\eqref{eq:nakarushton} с $r_{\max}=1$, а $\sigma$ подстраивается под яркость с $\tau=1\,\text{с}$ по одному из двух законов: в логарифмах, $\tau\,\dd\ln\sigma/\dd t=\ln I-\ln\sigma$, или линейно, $\tau\,\dd\sigma/\dd t=I-\sigma$. Фон $I_b=1$ в течение $10\,\text{с}$, затем $100$ в течение $10\,\text{с}$, затем снова $1$. Каждые $100\ms$ подаётся проба: яркость на $10\,\ms$ повышается на $10\,\%$. Ответ на пробу считайте как разность с таким же датчиком без проб.
\par\noindent(а)~Постройте ответы на пробы во времени для обоих законов. Каков ответ после полного привыкания? Сравните с $1.1/2.1-1/2$.
\par\noindent(б)~Каков ответ на первую пробу после скачка фона?
\par\noindent(в)~Через сколько секунд ответ на пробу восстанавливается до половины привыкшего после скачка вверх и после скачка вниз? Выведите эти времена из~\eqref{eq:nakarushton} для обоих законов и сравните.
\end{exercise}

\begin{exercise}[\lvlC]
\label{ex:11:7}
\emph{Под какой мир подстроена кривая.} Утверждение~\ref{prop:sensors} называет гипотезой то, что коды датчиков подстроены под наибольшую информацию.
\par\noindent(а)~Пусть к выходу $r(I)$ добавляется малый шум постоянной величины. Докажите, что информация о яркости наибольшая, когда распределение выхода равномерно, то есть $r(I)=r_{\max}\,\Phi_I(I)$, где $\Phi_I$ --- функция распределения яркостей.
\par\noindent(б)~Для какого распределения яркостей кривая~\eqref{eq:nakarushton} оптимальна в смысле (а)? Какой у него разброс $\log_{10}I$?
\par\noindent(в)~Если шум выхода пуассоновский, его дисперсия пропорциональна $r$. Какая кривая оптимальна тогда?
\par\noindent(г)~Открытый вопрос. Регулировка усиления в главе двигает только $\sigma$, то есть середину кривой. Предложите правило, которое подстраивало бы ещё и крутизну под разброс яркостей сцены, запишите его в непрерывном времени и проверьте на модели задачи~\ref{ex:11:6} со сценами разного контраста. Как опытом отличить датчик, который максимизирует информацию на импульс, от датчика, максимизирующего информацию в секунду?
\end{exercise}

\chapter{Распознавание изображений и видео}
\chaptermark{Распознавание}
\label{sec:recognition}

\section{Задача: одно и то же при разных пикселях}

Глава~\ref{sec:sensors} довела изображение до входа машины: около миллиона выходных нейронов глаза, сжатый поток, в котором передаются в основном края и перемены. Но между датчиком и действием есть шаг, о котором мы пока молчали. Кошка на картинке --- это не набор пикселей. Сдвиньте её на полкадра, поверните, отодвиньте, выключите половину света --- почти все пиксели станут другими, а ответ «кошка» должен остаться тем же. Инженер называет это \emph{инвариантностью}: ответ классификатора не должен меняться при сдвиге, масштабе, повороте и освещении.

Трудность хорошо видна на простом опыте. Возьмём одномерную «сетчатку» из $24$ точек и две фигуры по пять точек, которые могут стоять где угодно. Классификатор, который сравнивает вход с образцом, запомненным в одном месте, угадывает в $49$--$52\%$ случаев --- столько же, сколько монетка. Сдвиг ломает сравнение по пикселям полностью. Нужна другая схема.

\section{Конвейер из ступеней}

Как быстро мозг это делает, мы уже знаем из главы~\ref{sec:code}: животное на мгновенно показанной картинке узнаётся примерно за $150\ms$, и сигнал проходит около десятка ступеней, по $10$--$20\ms$ на каждую~\cite{Thorpe1996}. На каждой ступени нейрон успевает выдать ноль или один импульс. Значит, быстрое распознавание --- это один проход по конвейеру, без долгих раздумий и повторов. Вопрос в том, что делает каждая ступень.

Ответ, который подсказали измерения на первых ступенях зрения~\cite{HubelWiesel1962}, состоит из двух чередующихся операций (рис.~\ref{fig:conveyor}).

\emph{Избирательность.} Нейрон ступени $S$ смотрит на небольшой участок входа и срабатывает, только если там есть определённое сочетание частей: вот этот край, а рядом тот. Это И по частям --- пороговый нейрон главы~\ref{sec:glutcomputer}, у которого порог выше, чем даёт любая одна часть.

\emph{Инвариантность.} Нейрон ступени $C$ собирает выходы многих $S$-нейронов, ищущих одно и то же сочетание в разных местах, и срабатывает, если оно нашлось хоть где-то. Это ИЛИ по положениям --- а точнее, выбор наибольшего.

Для одномерной модели обе операции записываются так:
\begin{empheq}[box=\keybox]{equation}
s_{f,p} \;=\; \Bigl[\sum_i t_{f,i}\,x_{p+i} - \theta\Bigr]_+ ,
\qquad
c_f \;=\; \max_p\, s_{f,p} ,
\label{eq:scstage}
\end{empheq}
где $x$ --- вход, $t_{f,i}$ --- образец признака $f$, $p$ --- положение, $\theta$ --- порог. Первое выражение делает ответ избирательным, второе --- независимым от положения. Наибольшее из многих входов сеть получает из средств главы~\ref{sec:gaba}: взаимное торможение оставляет самый сильный вход (утверждение~\ref{prop:wta}), а деление на общую активность выравнивает ответы при разной яркости~\cite{Carandini2012}.

\begin{figure}[t]
\centering
\begin{tikzpicture}[>=Stealth,
  px/.style={draw,minimum size=0.36cm,inner sep=0pt},
  on/.style={px,fill=blue!55!black},
  snode/.style={circle,draw,very thick,minimum size=0.62cm,inner sep=0pt,font=\scriptsize,fill=blue!6},
  cnode/.style={circle,draw,very thick,minimum size=0.8cm,inner sep=0pt,font=\scriptsize,fill=red!10}]
% two inputs: the same shape at two positions
\foreach \row/\sh/\lab in {0/1/{кадр 1},1/7/{кадр 2}}{
  \begin{scope}[yshift=-\row*1.0cm]
  \node[left,font=\scriptsize] at (-0.25,0) {\lab};
  \foreach \i in {0,...,13}{\node[px] at (\i*0.4,0) {};}
  \foreach \d in {0,1,3,4}{\node[on] at ({(\sh+\d)*0.4},0) {};}
  \end{scope}}
% S stage: detectors of the shape at several positions
\foreach \k in {0,...,5}{
  \node[snode] (S\k) at ({0.8+0.8*\k},1.7) {$S$};}
\node[font=\scriptsize,left] at (-0.25,1.7) {ступень $S$: И по частям};
\draw[->,thick,blue!60!black] (1.2,0.25) -- (S1.south);
\draw[->,thick,orange!85!black] (3.6,0.25) -- (S4.south);
\node[font=\scriptsize,blue!60!black,anchor=east] at (1.25,0.85) {кадр 1};
\node[font=\scriptsize,orange!85!black,anchor=west] at (3.9,0.85) {кадр 2};
% C stage
\node[cnode] (C) at (2.8,3.25) {$C$};
\node[font=\scriptsize,left] at (-0.25,3.25) {ступень $C$: ИЛИ по положениям};
\foreach \k in {0,...,5}{\draw[->,gray!60!black] (S\k.north) -- (C);}
\draw[->,very thick,red!70!black] (C.north) -- ++(0,0.6) node[above,font=\scriptsize,black] {«фигура есть» --- в обоих кадрах};
% next stages
\draw[decorate,decoration={brace,amplitude=5pt},gray!60!black] (6.3,1.4) -- (6.3,-1.3);
\node[right,font=\scriptsize,align=left] at (6.5,0.05) {один и тот же\\объект, другие\\пиксели};
\draw[->,thick,densely dashed,gray!60!black] (3.6,3.25) -- (5.9,3.25) node[right,font=\scriptsize,black,align=left] {следующие пары\\$S$, $C$ --- крупнее,\\сложнее, инвариантнее};
\end{tikzpicture}
\caption{Одна пара ступеней конвейера. Нейроны $S$ ищут одно и то же сочетание частей в разных местах; нейрон $C$ отвечает, если оно нашлось хоть где-то~\eqref{eq:scstage}. Фигура в кадре 1 и в кадре 2 возбуждает разные $S$-нейроны, но один и тот же $C$. Следующие пары ступеней повторяют то же на всё более крупных и сложных сочетаниях.}
\label{fig:conveyor}
\end{figure}

В той же одномерной модели пара ступеней~\eqref{eq:scstage} узнаёт фигуру в любом месте без ошибок, а если каждую точку входа случайно перевернуть с вероятностью $0.1$ --- в $86\%$ случаев, с вероятностью $0.2$ --- в $70\%$. Монетка по-прежнему даёт половину. Сложите несколько таких пар одну над другой: первая ищет края, следующая --- сочетания краёв, ещё выше --- части предметов, наверху --- предметы целиком. С каждой парой сочетания крупнее, а ответ всё меньше зависит от того, где и как предмет лежит~\cite{Riesenhuber1999}.

\begin{keyblock}
\begin{proposition}[Конвейер избирательности и инвариантности]
\label{prop:conveyor}
Быстрое распознавание --- один проход по десятку ступеней. На каждой паре ступеней И по частям~\eqref{eq:scstage} делает ответ избирательным, а ИЛИ по положениям --- независимым от сдвига. Чередование этих двух операций даёт ответ, который различает предметы и почти не зависит от того, где и как они видны. Устройство первых ступеней и скорость прохода измерены; то, что вся цепочка сводится к этим двум операциям, --- упрощение.
\end{proposition}
\end{keyblock}

Если это устройство кажется знакомым, так и должно быть. Именно его чередование --- поиск признака во всех местах, потом объединение по соседним местам --- инженеры перенесли в искусственные \emph{свёрточные} сети~\cite{Fukushima1980}. А недавно пошли обратно: свёрточные сети, обученные распознавать предметы, оказались лучшей из известных моделей ответов нейронов на верхних ступенях зрения~\cite{Yamins2014}. Результат наверху читается просто: по суммарным частотам нескольких сотен нейронов последней ступени линейный решающий блок узнаёт предмет через десятую долю секунды после показа~\cite{Hung2005}. Это частотный код группы из главы~\ref{sec:code}.

\section{Учитель, которого нет: учиться на видео}

Свёрточную сеть учат на миллионах картинок с подписями. Мозгу подписей почти никто не даёт: ребёнок видит предметы часами, а слово «чашка» слышит изредка. Откуда тогда ступени $C$ знают, какие $S$-нейроны объединять? Ведь чтобы объединить «эту фигуру здесь» с «этой фигурой там», надо уже знать, что это одна и та же фигура.

Подсказку даёт само видео. Мир меняется непрерывно: предмет в поле зрения остаётся тем же предметом, пока он сдвигается, поворачивается и меняет освещение, и только иногда взгляд переходит на другой. Всё, что меняется \emph{медленно}, скорее всего относится к предмету, а всё, что меняется быстро, --- к его положению и виду~\cite{WiskottSejnowski2002}. Значит, $C$-нейрону достаточно правила Хебба из главы~\ref{sec:dofa} со следом: связывать то, что шло подряд, а не только то, что шло одновременно~\cite{Foldiak1991}:
\begin{empheq}[box=\keybox]{equation}
\tau_{\text{сл}}\,\frac{\dd\bar y_k}{\dd t} \;=\; -\bar y_k + y_k ,
\qquad
\frac{\dd\mathbf w_k}{\dd t} \;=\; \eta\,\bar y_k\,\bigl(\mathbf x - \mathbf w_k\bigr) .
\label{eq:traceinvariance}
\end{empheq}
Здесь $y_k$ --- активность $C$-нейрона $k$, который выиграл соревнование через торможение, $\bar y_k$ --- её след, та же RC-цепь, что и в правиле~\eqref{eq:threefactor}, а $\mathbf x$ --- выходы $S$-нейронов. Нейрон, выигравший на первом виде предмета, ещё помнит это, когда приходит второй вид, и подтягивает к себе и его.

\begin{figure}[t]
\centering
\begin{tikzpicture}[>=Stealth,x=1.0cm,y=1.0cm]
\foreach \i/\lab in {0/1,1/2,2/3,3/4}{
  \node[draw,thick,fill=blue!8,minimum width=0.9cm,minimum height=0.55cm,font=\scriptsize] at (\i+0.5,2.2) {$A$ в \lab};}
\foreach \i/\lab in {0/4,1/3,2/2,3/1}{
  \node[draw,thick,fill=orange!12,minimum width=0.9cm,minimum height=0.55cm,font=\scriptsize] at (\i+6.5,2.2) {$B$ в \lab};}
\node[font=\scriptsize,gray!40!black] at (5.0,2.2) {пауза};
\draw[->,gray!70!black] (0,0) -- (10.4,0) node[below left,font=\scriptsize,black] {время};
% traces
\draw[thick,blue!60!black] (0,0.05) .. controls (0.4,1.2) and (0.8,1.3) .. (1.2,1.35) -- (4,1.4) .. controls (4.6,0.5) and (5.2,0.15) .. (6,0.08) -- (10,0.05);
\draw[thick,orange!85!black] (0,0.05) -- (6,0.05) .. controls (6.4,1.2) and (6.8,1.3) .. (7.2,1.35) -- (10,1.4);
\node[font=\scriptsize,blue!60!black,anchor=west] at (1.4,1.65) {след $\bar y_1$ --- держится весь проход $A$};
\node[font=\scriptsize,orange!85!black,anchor=west] at (7.2,1.65) {след $\bar y_2$};
\draw[decorate,decoration={brace,amplitude=4pt},gray!60!black] (4.0,2.6) -- (0,2.6);
\draw[decorate,decoration={brace,amplitude=4pt,mirror},gray!60!black] (0,-0.35) -- (4,-0.35) node[midway,below=4pt,font=\scriptsize,black] {все виды $A$ связываются с нейроном 1};
\end{tikzpicture}
\caption{Как видео учит инвариантности, схема. Предмет $A$ проходит через четыре положения, потом пауза, потом предмет $B$. След~\eqref{eq:traceinvariance} нейрона, выигравшего на первом виде $A$, держится весь проход, и все виды $A$ оказываются связаны с одним нейроном. Пауза стирает след, и $B$ достаётся другому нейрону.}
\label{fig:tracevideo}
\end{figure}

Мы проверили это на модели (рис.~\ref{fig:tracevideo}). Два $C$-нейрона соревнуются за входы от $16$ $S$-нейронов: две фигуры в восьми положениях. Фигура проезжает через все положения, по $50\ms$ на каждое, потом полсекунды пауза и следующая случайная фигура. Никаких подписей. Если след короткий, $\tau_{\text{сл}}=5\ms$, нейроны делят вход по \emph{положению}, и ответ совпадает с фигурой не лучше монетки: $52\%$ в среднем по десяти прогонам. Если след длиннее одного положения, $\tau_{\text{сл}}=50$--$200\ms$, во всех десяти прогонах каждый нейрон собирает \emph{все} положения своей фигуры: инвариантность выучена на одном только видео. Пауза между предметами важна: без неё след перетекает на следующий предмет и путает их.

То же самое видно в опыте. Если в эксперименте незаметно подменять предмет, пока глаз перескакивает с одного места на другое, то нейроны верхних ступеней за час-другой начинают отвечать на подменённый предмет как на тот же самый: они связывают то, что шло подряд~\cite{LiDiCarlo2008}. Инвариантность в мозге действительно учится на непрерывности времени. Насколько этим правилом объясняется всё обучение зрения, --- гипотеза.

\section{Видео --- это движение}

Видео --- не только поток кадров для обучения, но и сама задача: что движется, куда и как быстро. Первый шаг нам знаком --- детектор направления главы~\ref{sec:gaba} (рис.~\ref{fig:direction}), в котором запаздывающее торможение гасит движение в одну сторону. Такие детекторы стоят по всему полю зрения, и дальше с ними поступают так же, как с признаками формы: ступени, собирающие много детекторов одного направления в разных местах, отвечают на движение крупного предмета независимо от того, где он. Это та же пара И и ИЛИ~\eqref{eq:scstage}, только признак не «край», а «край, сдвинувшийся за время $d$».

Видео ещё и предсказуемо: следующий кадр почти такой же, как предыдущий. По гипотезе \emph{предсказывающего кодирования} верхние ступени посылают нижним предсказание, а вверх уходит главным образом ошибка --- то, чего предсказание не учло~\cite{RaoBallard1999}. Это та же экономия импульсов, что в утверждении~\ref{prop:economy}: платить за новости, а не за то, что и так известно. Отдельные наблюдения с этой гипотезой согласуются, но как общее устройство конвейера она не доказана.

\section{Мозг и свёрточная сеть}

Насколько далеко заходит сходство? Для быстрого распознавания одного предмета оно действительно глубокое~\cite{Yamins2014}. Но у мозга есть и то, чего нет у простой свёрточной сети.

\emph{Обратные связи.} Конвейер мозга не только прямой: у каждой ступени есть связи назад и вбок. Для трудных картинок --- с перекрытием, в плохом свете --- ответ верхних ступеней складывается позже, и эти поздние ответы лучше описываются сетями с обратными связями, чем прямыми~\cite{Kar2019}. Один проход решает лёгкие случаи, трудные требуют нескольких кругов.

\emph{Устойчивость.} Искусственную сеть можно обмануть незаметной для глаза подделкой пикселей: изменение, которое человек не видит, превращает «панду» в «гиббона»~\cite{Szegedy2014}. Зрение человека к таким подделкам почти нечувствительно. Почему --- вопрос открытый; кандидаты --- шум, деление на общую активность и обратные связи.

\emph{Учитель.} Сети нужны миллионы подписанных примеров, мозгу --- в основном видео без подписей~\eqref{eq:traceinvariance} и немного подсказок от \nt{ДОФА} о том, что важно.

\emph{Энергия.} Весь мозг --- около $20$ ватт (утверждение~\ref{prop:economy}), а распознавание занимает лишь часть этого; обученная свёрточная сеть на графическом ускорителе тратит сотни ватт.

\section{Иллюзии: где машина ошибается и почему}
\label{sec:illusions}

Инженер, который хочет понять чужую схему, подаёт на неё необычные сигналы и смотрит, где она ошибается. Для зрения такие сигналы давно придуманы --- это зрительные иллюзии. Иллюзия --- не поломка. Каждая указывает на определённый механизм машины, который в обычной обстановке работает правильно, а на специально подобранной картинке выдаёт себя.

\emph{Разумные ожидания.} Вход шумный, и машина дополняет его тем, что обычно бывает в мире. Медленные движения встречаются чаще быстрых; когда вход ненадёжен --- например, у картинки мало контраста, --- оценка скорости сдвигается к медленной, и бледный предмет кажется движущимся медленнее яркого~\cite{Weiss2002}. Так же объясняют многие иллюзии яркости: мы видим не яркость пикселя, а то, какой поверхности такая яркость чаще всего соответствовала в прошлом~\cite{Purves2011}. Фотография платья, которое одни видят бело-золотым, а другие сине-чёрным, --- тот же механизм: картинка двусмысленна, и люди расходятся в том, каким они неосознанно считают освещение~\cite{LaferSousa2015,Wallisch2017}. Различия между людьми измерены; то, что мозг здесь сочетает вход с ожиданием по правилам теории вероятностей, --- хорошо обоснованная гипотеза.

\emph{Регулировка усиления.} Посмотрите минуту на водопад, потом на неподвижные камни: камни поплывут вверх. Каналы «вниз» и «вверх» --- пара проводов, как в~\eqref{eq:dualrail}; регулировка усиления~\eqref{eq:nakarushton} приглушила уставший канал, и равновесие пары сдвинулось. Иллюзии наклона и контраста, в которых окружение меняет вид центра, объясняются делением на активность соседей~\cite{Schwartz2009}, как в главе~\ref{sec:gaba}. Есть и поучительный пример осторожности: тёмные пятна на перекрёстках белых полос решётки долго объясняли простым торможением соседей, но опыты с изменённой решёткой показали, что это объяснение не работает~\cite{SchillerCarvey2005}.

\emph{Компенсация задержек.} Путь от глаза до верхних ступеней занимает десятки миллисекунд, и движущийся предмет за это время успевает сдвинуться. Если рядом с ним вспыхнет неподвижная точка, предмет кажется впереди вспышки, хотя они совпадали~\cite{Nijhawan1994}. По одному из объяснений машина продлевает движение вперёд, чтобы скомпенсировать задержку, --- то же, что в главе~\ref{sec:muscles}: при запаздывающей обратной связи приходится предсказывать. Объяснений у этой иллюзии несколько, и спор не решён.

\emph{Ошибка в коде времени.} Неподвижная картинка из колец с резкими цветовыми переходами кажется вращающейся. Контрастные участки обрабатываются быстрее бледных, и разница задержек~\eqref{eq:latency} читается детекторами направления как движение~\cite{Conway2005}. Запускают иллюзию мелкие непроизвольные скачки глаза и моргания: каждый скачок обновляет вход, и детекторы снова видят «движение»~\cite{OteroMillan2012}. Это код временем из главы~\ref{sec:code}, прочитанный неправильно.

\emph{Две картинки в одной.} Каркас куба, который то смотрит на нас одной гранью, то другой, или разные картинки в двух глазах: восприятие перескакивает каждые несколько секунд. Лучшие модели --- взаимное торможение двух толкований, медленная усталость и шум~\cite{MorenoBote2007}, то есть та самая схема~\eqref{eq:halfcentre}, которая в главе~\ref{sec:muscles} порождала ритм шага. Время смены задаёт усталость.

А искусственные сети? Сеть, обученная предсказывать следующий кадр видео, видит вращение на тех же неподвижных кольцах и не видит его на контрольных картинках~\cite{Watanabe2018}. Сети, обученные очищать изображения от шума и повышать их резкость, поддаются тем же иллюзиям цвета и яркости, что и люди~\cite{GomezVilla2020}. Значит, часть иллюзий --- следствие самой задачи, которой учится машина, а не особенностей нервной ткани. Какие иллюзии свойственны только мозгу, --- хорошая проверка для любой модели зрения.

\begin{keyblock}
\begin{proposition}[Иллюзии как отпечатки механизмов]
\label{prop:illusions}
Иллюзия возникает, когда механизм, полезный в обычном мире, получает необычный вход: ожидание побеждает ненадёжный вход, регулировка усиления сдвигает пару каналов, компенсация задержки уводит предмет вперёд, разница задержек читается как движение, взаимное торможение с усталостью перескакивает. Каждая иллюзия --- испытательный сигнал, указывающий на свой механизм. Сами иллюзии измерены; их объяснения --- от хорошо обоснованных до спорных.
\end{proposition}
\end{keyblock}

\section{Итог}

\begin{table}[t]
\centering
\footnotesize
\setlength{\tabcolsep}{4pt}
\renewcommand{\arraystretch}{1.15}
\begin{tabular}{>{\raggedright}p{3.0cm}>{\raggedright}p{4.6cm}>{\raggedright\arraybackslash}p{5.2cm}}
\toprule
Что делает машина & Как & Что известно \\
\midrule
узнаёт за $150\ms$ & один проход по десятку ступеней & скорость измерена \\
делает ответ избирательным & И по частям, ступень $S$~\eqref{eq:scstage} & измерено на первых ступенях \\
делает ответ инвариантным & ИЛИ по положениям, ступень $C$; торможение и деление & измерено на первых ступенях; для верхних --- упрощение \\
выдаёт ответ & частоты сотен нейронов последней ступени, линейное чтение & измерено \\
учится без подписей & правило Хебба со следом~\eqref{eq:traceinvariance} на непрерывном видео & подмена предмета меняет ответы --- измерено; как главное правило --- гипотеза \\
видит движение & детекторы направления, затем та же пара И/ИЛИ & измерено \\
экономит на предсказуемом & вверх идёт ошибка предсказания & гипотеза \\
решает трудные случаи & обратные связи, несколько кругов & поздние ответы и роль обратных связей измерены \\
ошибается предсказуемо & иллюзии: ожидания, регулировка усиления, задержки, взаимное торможение с усталостью & иллюзии измерены; объяснения --- от обоснованных до спорных \\
\bottomrule
\end{tabular}
\caption{Как машина распознаёт изображения и видео.}
\label{tab:recognition}
\end{table}

Распознавание собрано из деталей, которые у нас уже были (таблица~\ref{tab:recognition}): порог даёт И по частям, взаимное торможение --- ИЛИ по положениям, деление --- независимость от яркости, правило Хебба со следом --- учителя, которого нет. Инженеры взяли у зрения чередование избирательности и инвариантности и построили на нём свёрточные сети; обратно мозг пока не отдал главного --- умения учиться на видео без подписей почти без энергии.

\section{Задачи}

\begin{exercise}[\lvlA]
\label{ex:12:1}
\emph{Частота на одной ступени.} Ступень конвейера длится $T=15\ms$, нейроны работают с частотой $r=20$ в секунду. В главе~\ref{sec:code} сто нейронов с независимыми шумами давали за такое время точность около двадцати процентов.
\par\noindent(а)~Какова по~\eqref{eq:rateerror} относительная ошибка частоты одного нейрона за время ступени?
\par\noindent(б)~Теперь шумы коррелированы, $c=0.1$~\eqref{eq:correlated}. Какова точность группы из ста нейронов и каков предел при сколь угодно большой группе? Сколько уровней частоты различимо на одной ступени?
\par\noindent(в)~Что отсюда следует о том, как ступени передают друг другу числа, если корреляция похожа на сигнал? Назовите два выхода из главы~\ref{sec:code}.
\end{exercise}

\begin{exercise}[\lvlA]
\label{ex:12:2}
\emph{Энергия одного узнавания.} Пусть в одном проходе участвуют десять ступеней по $10^7$ нейронов, и каждый выдаёт не больше одного импульса ценой $E_{\text{имп}}\sim10^{-10}$~Дж (глава~\ref{sec:code}).
\par\noindent(а)~Сколько энергии стоит одно узнавание? Какую долю это составляет от энергии всего мозга и от энергии его сигналов~\eqref{eq:energy} за те же $150\ms$?
\par\noindent(б)~С какой средней частотой работают эти нейроны во время прохода? Совместимо ли это со средней частотой~\eqref{eq:energy}, и при каком условии?
\par\noindent(в)~Ускоритель мощностью $300$~Вт узнаёт картинку за $10\ms$. Во сколько раз больше энергии он тратит на одно узнавание?
\end{exercise}

\begin{exercise}[\lvlB]
\label{ex:12:3}
\emph{Пороги ступени $S$.} Одномерная «сетчатка» из $N=24$ точек, фигуры $A=11011$ и $B=10101$ из текста. Образец признака в~\eqref{eq:scstage}: $t_{f,i}=+1$, где у фигуры $1$, и $t_{f,i}=-1$, где $0$ (отрицательный вес --- через \nt{ГАМК}-посредника); $n_f$ --- число единиц фигуры.
\par\noindent(а)~Докажите, что при $n_f-1\le\theta_f<n_f$ нейрон $s_{f,p}$ отвечает тогда и только тогда, когда в окне $p$ стоит в точности фигура $f$. Найдите пороги для $A$ и $B$.
\par\noindent(б)~Найдите пороги, при которых нейрон прощает одну перевёрнутую точку в окне. Проверьте, что при них детектор $A$ не отвечает на чистую $B$ ни в каком положении, и наоборот.
\par\noindent(в)~Докажите, что $c_f$ не меняется при сдвиге фигуры, пока она остаётся внутри «сетчатки».
\par\noindent(г)~Сосчитайте нейроны двух ступеней, если отрицательный вес от каждой точки даёт один общий \nt{ГАМК}-посредник этой точки. Сколько $S$-нейронов понадобится для десяти признаков размером $5\times5$ на картинке $100\times100$?
\end{exercise}

\begin{exercise}[\lvlB]
\label{ex:12:4}
\emph{Сколько длится одна картинка из двух.} Перескоки восприятия описывают схемой~\eqref{eq:halfcentre}.
\par\noindent(а)~Пусть $\tau\ll\tau_a$. Пока побеждает группа~1, $r_2=0$ и $r_1=I-a_1$. Найдите $a_1(t)$ и $a_2(t)$ и выведите условие, при котором группа~2 вырывается вперёд.
\par\noindent(б)~В симметричном цикле усталость каждой группы в начале и в конце её победы одна и та же. Составьте уравнения для длительности победы $T_d$ и решите их численно при $I=1$, $\beta=2$, $b=2$, $\tau_a=0.4\,\text{с}$.
\par\noindent(в)~Смоделируйте~\eqref{eq:halfcentre} при $\tau=20\ms$ и при $\tau=2\ms$ и сравните период с (б) и с рис.~\ref{fig:halfcentre}. Откуда разница?
\par\noindent(г)~Докажите, что в пределе (а) $T_d$ пропорционально $\tau_a$ и не зависит от $I$. Какое $\tau_a$ даёт смену каждые три секунды? Проверьте моделированием.
\end{exercise}

\begin{exercise}[\lvlB]
\label{ex:12:5}
\emph{Моделирование: цена инвариантности.} Используйте функцию \texttt{two\_stage} из \texttt{models/\allowbreak recognition\_models.py} ($4000$ попыток, \texttt{seed}$\,=0$).
\par\noindent(а)~При $N=24$ постройте долю верных ответов обоих классификаторов в зависимости от вероятности переворота точки $p$ от $0$ до $0.5$. При каком $p$ пара ступеней $S$, $C$ ошибается в каждом четвёртом случае?
\par\noindent(б)~При $p=0.1$ измерьте долю верных ответов для $N=8$, $12$, $24$, $48$, $96$, $192$. Как она зависит от $N$?
\par\noindent(в)~Объясните зависимость (б): что делает взятие наибольшего по всё большему числу положений, когда вход шумный? Предложите изменение ступени $C$, которое уменьшило бы эту потерю, и объясните, почему в зрении ступени $C$ собирают лишь соседние положения.
\end{exercise}

\begin{exercise}[\lvlB]
\label{ex:12:6}
\emph{Моделирование: окно следа.} Используйте функцию \texttt{trace\_learning} из \texttt{models/\allowbreak recognition\_models.py} (десять прогонов, \texttt{seed}$\,=0,\dots,9$).
\par\noindent(а)~При паузе $0.3\,\text{с}$ постройте среднюю и наихудшую оценку инвариантности в зависимости от $\tau_{\text{сл}}$ от $5\ms$ до $2\,\text{с}$ (логарифмическая шкала). В каком окне $\tau_{\text{сл}}$ все десять прогонов выучивают инвариантность безошибочно?
\par\noindent(б)~Повторите без паузы.
\par\noindent(в)~Объясните верхнюю границу окна (а) через~\eqref{eq:traceinvariance}: что должно случиться со следом за паузу?
\par\noindent(г)~Проверьте, зависит ли нижняя граница от времени показа одного положения (\texttt{dwell}$\,=0.2\,\text{с}$ вместо $50\ms$), от скорости обучения (\texttt{eta} от $0.5$ до $8$) и от шага интегрирования (\texttt{dt}$\,=1\ms$). Что это говорит о её происхождении?
\end{exercise}

\begin{exercise}[\lvlB]
\label{ex:12:7}
\emph{Проектирование: движение в любом месте.} Ряд из $24$ точек с шагом $\Delta x=0.1^\circ$. На каждой паре соседних точек стоит детектор направления главы~\ref{sec:gaba}: \nt{ГЛУТ}-нейрон $Y$ возбуждается правой точкой $B$ и тормозится левой точкой $A$ через \nt{ГАМК}-посредника с задержкой $d$; торможение гасит возбуждение, если они приходят не дальше чем на $\Delta_{\text{т}}=5\ms$ друг от друга. Над детекторами --- ступень $C$~\eqref{eq:scstage}.
\par\noindent(а)~В какую сторону движения отвечает $Y$? Для какой полосы скоростей одна задержка $d$ гасит противоположное направление?
\par\noindent(б)~Нужно, чтобы $C$ отвечал на движение в свою сторону в любом месте ряда и молчал при движении в противоположную сторону со скоростью от $5$ до $20^\circ$ в секунду. Хватит ли одной задержки? Предложите схему и найдите все допустимые задержки.
\par\noindent(в)~Сосчитайте нейроны схемы. Что произойдёт при скорости $40^\circ$ в секунду в противоположную сторону?
\end{exercise}

\begin{exercise}[\lvlC]
\label{ex:12:8}
\emph{Подделка пикселей.} Раздел «Мозг и свёрточная сеть» оставил открытым вопрос, почему зрение почти нечувствительно к незаметным подделкам.
\par\noindent(а)~В модели \texttt{two\_stage} решает разность $\Delta=c_{\text{верн}}-c_{\text{чуж}}$, где $c_f$ --- наибольшее по окнам $k$ число совпадений фигуры $f$ с окном $k$. Докажите, что для смены ответа нужно не меньше $\lceil(\Delta+1)/2\rceil$ перевёрнутых точек, а для ничьей --- $\lceil\Delta/2\rceil$.
\par\noindent(б)~Найдите $\Delta$ для чистых фигур при $N=24$ и проверьте полным перебором, сколько переворотов на самом деле нужно.
\par\noindent(в)~Классификатор «число единиц на входе больше $3.5$ --- это $A$» тоже инвариантен к сдвигу и безошибочен на чистых картинках. Сколько переворотов нужно, чтобы его обмануть? Какова его доля верных ответов при $p=0.1$ по сравнению с $0.86$ у пары ступеней?
\par\noindent(г)~Открытый вопрос. Что даёт устойчивость к подделке: инвариантность или избирательность? Предложите изменение модели, которое увеличит наименьшее число переворотов, не ухудшив ответа на случайный шум, и проверьте его. Какой из кандидатов текста --- шум, деление на общую активность, обратные связи --- можно проверить в этой модели и как?
\end{exercise}

\chapter{Выходы машины: как мозг управляет мышцами}
\chaptermark{Выходы машины}
\label{sec:muscles}

\section{Быстрый выход}

Всё, что машина делает в мире быстро, она делает мышцами. Слово, взгляд, шаг, улыбка --- это сокращения мышц. Второй, медленный выход --- химия тела --- разобран в главе~\ref{sec:chemoutput}. На рис.~\ref{fig:machine} выход был стрелкой «мышцы»; теперь посмотрим, что за этой стрелкой.

Последнее звено цепи --- \nt{ACET}-нейроны, те самые, у которых в таблице~\ref{tab:types} записано «выход на мышцу». Каждое мышечное волокно получает вход ровно от одного такого нейрона, а один нейрон управляет группой волокон --- от нескольких сотен до пары тысяч~\cite{Feinstein1955mu}. В мышцах, которым нужна тонкая подстройка, единицы мельче, в больших мышцах ног --- крупнее. Нейрон вместе со своими волокнами называют \emph{двигательной единицей}: это наименьший кусок мышцы, который мозг может включить отдельно.

Синапс нейрона на мышце устроен совсем не так, как синапсы коры из главы~\ref{sec:code}. Там большая часть импульсов терялась. Здесь каждый импульс выбрасывает в несколько раз больше ацетилхолина, чем нужно, чтобы волокно сработало~\cite{WoodSlater2001}. Это синапс с \emph{запасом надёжности}: один импульс нейрона --- ровно одно сокращение волокон. Внутри машины можно позволить себе ненадёжные связи и исправлять ошибки избыточностью; на выходе ошибка стоит движения, и надёжность куплена прямо в железе.

\section{Сила: частота и число единиц}

Один импульс даёт одиночное сокращение --- подёргивание: сила нарастает за несколько десятков миллисекунд и спадает. Хорошее приближение для него~\cite{Fuglevand1993}
\begin{equation}
h(t) \;=\; F_1\,\frac{t}{T_c}\,e^{\,1-t/T_c},
\label{eq:twitch}
\end{equation}
где $T_c$ --- время нарастания, порядка $50\ms$, а $F_1$ --- сила одного подёргивания. Если импульсы идут чаще, подёргивания складываются:
\begin{empheq}[box=\keybox]{equation}
F(t) \;=\; \sum_k h\bigl(t-t_k\bigr) .
\label{eq:forcesum}
\end{empheq}
Это тот же фильтр нижних частот, что превращал импульсы в концентрацию в главе~\ref{sec:serocomputer}, только теперь на выходе не концентрация, а сила. Мышца --- цифро-аналоговый преобразователь из импульсов в силу (рис.~\ref{fig:tetanus}). В нашей модели при пяти импульсах в секунду подёргивания почти не перекрываются и сила скачет от $0.2F_1$ до $1.1F_1$; при пятнадцати она уже держится между $1.8F_1$ и $2.2F_1$; при сорока становится почти ровной, около $5.4F_1$. Настоящая мышца при частых импульсах насыщается, так что линейная сумма~\eqref{eq:forcesum} верна лишь до нескольких десятков импульсов в секунду.

\begin{figure}[t]
\centering
\begin{tikzpicture}[>=Stealth,x=20cm,y=0.55cm]
\draw[->,gray!70!black] (0,0) -- (0.66,0) node[right,font=\small,black] {$t$, с};
\draw[->,gray!70!black] (0,0) -- (0,6.4) node[left,font=\small,black] {$F/F_1$};
\foreach \t in {0,0.2,0.4,0.6}{\draw[gray!60!black] (\t,-0.1) -- (\t,0.1); \node[below,font=\scriptsize] at (\t,0) {\t};}
\foreach \f in {1,3,5}{\draw[gray!60!black] (-0.004,\f) -- (0.004,\f); \node[left,font=\scriptsize] at (0,\f) {\f};}
\draw[thick,blue!60!black] plot file {figures/muscle_twitch_5.dat};
\draw[thick,green!45!black] plot file {figures/muscle_twitch_15.dat};
\draw[thick,red!70!black] plot file {figures/muscle_twitch_40.dat};
\node[font=\scriptsize,blue!60!black,anchor=west] at (0.605,0.9) {5 имп/с};
\node[font=\scriptsize,green!45!black,anchor=west] at (0.605,2.1) {15 имп/с};
\node[font=\scriptsize,red!70!black,anchor=west] at (0.605,5.4) {40 имп/с};
\end{tikzpicture}
\caption{Мышца как цифро-аналоговый преобразователь: сила~\eqref{eq:forcesum} при регулярных импульсах одной двигательной единицы, $T_c=50\ms$. Редкие импульсы дают отдельные подёргивания, частые сливаются в ровную силу --- так же, как частые импульсы сливаются в ровную концентрацию на рис.~\ref{fig:lowpass}.}
\label{fig:tetanus}
\end{figure}

У мозга две ручки силы: частота импульсов каждой единицы и число включённых единиц. И вторая ручка устроена очень изобретательно. Единицы в мышце разные: самые слабые в сотню раз слабее самых сильных. Когда сила нужна всё большая, единицы включаются \emph{по порядку размера}, от маленьких к большим~\cite{Henneman1957}. Никто этот порядок не вычисляет: маленькие нейроны просто легче возбуждаются от того же общего входа, так что порядок включения задан самим железом.

Что это даёт? Пусть каждая следующая по порядку единица сильнее предыдущей в $q$ раз, где $q$ немного больше единицы. Сила всех включённых единиц --- сумма геометрической прогрессии, и почти вся она приходится на последние единицы. Тогда очередная включённая единица добавляет
\begin{empheq}[box=\keybox]{equation}
\frac{\Delta F}{F} \;\approx\; q-1 \;=\; \text{const} ,
\label{eq:sizeprinciple}
\end{empheq}
то есть шаг силы пропорционален самой силе. При слабом усилии мозг дозирует его мелкими порциями, при сильном --- крупными. Это \emph{число с плавающей точкой}: порядок задаётся тем, сколько единиц включено, а мантисса --- частотой их импульсов. Та же логарифмическая шкала, что у датчиков главы~\ref{sec:sensors}: вход и выход машины меряют мир в относительных долях.

У плавающей точки есть обратная сторона. Разброс силы тоже растёт пропорционально ей самой: сильное движение неизбежно менее точно, чем слабое. По одной влиятельной гипотезе, именно этот шум, зависящий от сигнала, объясняет, почему наши движения плавные: плавная команда даёт наименьший разброс в конечной точке~\cite{HarrisWolpert1998}.

\section{Тянуть, но не толкать: два провода на сустав}

Мышца умеет только тянуть. Толкать она не может, как не может \nt{GLUT}-нейрон выдать отрицательный сигнал. Решение знакомо по главе~\ref{sec:glutcomputer}: \emph{два провода}. Каждый сустав обслуживает пара мышц, тянущих в противоположные стороны (рис.~\ref{fig:antagonists}). Если их силы $F_1$ и $F_2$, а плечо $\ell$, то вращающий момент и жёсткость сустава равны
\begin{empheq}[box=\keybox]{equation}
M \;=\; \ell\,(F_1-F_2), \qquad K \;\approx\; \kappa_m\,(F_1+F_2),
\label{eq:antagonists}
\end{empheq}
где $\kappa_m$ --- коэффициент, связывающий силу мышцы с её жёсткостью: напряжённая мышца пружинит сильнее расслабленной. Разность сил задаёт момент, сумма --- жёсткость~\cite{Hogan1984}. Двумя проводами мозг управляет двумя независимыми величинами: куда повернуть сустав и насколько твёрдо его держать. Когда надо удержать чашку, не расплескав, мы напрягаем обе мышцы сразу: момент тот же, жёсткость выше, и случайный толчок меньше сбивает руку.

\begin{figure}[t]
\centering
\begin{tikzpicture}[>=Stealth,
  nrn/.style={circle,draw,very thick,minimum size=0.95cm,inner sep=0pt,font=\scriptsize},
  inh/.style={-{Bar[width=7pt]},thick,red!70!black}]
% the joint: a bar on a pivot, two springs pulling down on either side
\fill[gray!30] (4.6,-0.25) -- (5.4,-0.25) -- (5,0.15) -- cycle;
\draw[line width=3pt,gray!50!black] (2.4,0.2) -- (7.6,0.2);
\fill[black] (5,0.2) circle (0.08);
\draw[decorate,decoration={coil,segment length=5pt,amplitude=4pt},very thick,blue!60!black] (3.0,0.2) -- (3.0,-1.6);
\draw[decorate,decoration={coil,segment length=5pt,amplitude=4pt},very thick,orange!85!black] (7.0,0.2) -- (7.0,-1.6);
\draw[very thick] (2.5,-1.6) -- (3.5,-1.6); \draw[very thick] (6.5,-1.6) -- (7.5,-1.6);
\node[font=\small,blue!60!black] at (2.35,-0.75) {$F_1$};
\node[font=\small,orange!85!black] at (7.65,-0.75) {$F_2$};
\draw[->,thick] (5.9,0.75) arc (60:120:1.8) node[above,font=\scriptsize,pos=0.5] {$M=\ell(F_1-F_2)$};
% motor neurons and reflex circuit
\node[nrn,fill=green!12] (M1) at (1.6,-3.0) {\nt{ACET}};
\node[nrn,fill=green!12] (M2) at (8.4,-3.0) {\nt{ACET}};
\node[nrn,fill=red!10] (G) at (5.0,-3.0) {\nt{GLYC}};
\node[nrn,fill=blue!6] (S) at (1.6,-5.0) {\nt{GLUT}};
\draw[->,very thick] (M1) -- (3.0,-1.7);
\draw[->,very thick] (M2) -- (7.0,-1.7);
\draw[->,thick] (S) -- (M1);
\draw[->,thick] (S) -- (G);
\draw[inh] (G) -- (M2);
\draw[->,thick,densely dashed,gray!60!black] (3.15,-1.2) to[bend left=25] node[pos=0.35,right,font=\scriptsize,black] {растяжение} (S.north east);
\draw[->,thick,blue!60!black] (-0.3,-3.0) node[left,font=\scriptsize,black,align=right] {команда\\мозга} -- (M1);
\draw[->,thick,blue!60!black] (10.3,-3.0) node[right,font=\scriptsize,black,align=left] {команда\\мозга} -- (M2);
\end{tikzpicture}
\caption{Две мышцы на один сустав и самая быстрая петля обратной связи. \nt{ACET}-нейроны получают команды мозга и тянут сустав в разные стороны; разность сил поворачивает его, сумма делает жёстче. Датчик растяжения (\nt{GLUT}) левой мышцы возбуждает её собственный \nt{ACET}-нейрон и через \nt{GLYC}-посредника тормозит нейрон мышцы-противника --- запрет из главы~\ref{sec:gaba}.}
\label{fig:antagonists}
\end{figure}

\section{Обратная связь с задержкой}

Мышцы не работают вслепую. В каждой сидят датчики растяжения, о которых шла речь в таблице~\ref{tab:sensors}, и самая короткая петля замыкается прямо в спинном мозге (рис.~\ref{fig:antagonists}). Мышцу неожиданно растянули --- датчик растяжения возбуждает её собственный \nt{ACET}-нейрон, мышца сокращается и возвращает сустав. Одновременно через тормозящий нейрон-посредник выключается мышца-противник, чтобы не мешала. Этим посредником работает \nt{GLYC} --- тот самый тип из таблицы~\ref{tab:types}, которому не досталось своей главы: его место именно здесь, в быстрых петлях спинного мозга.

У каждой петли есть задержка $d$: у спинномозговой около $25$--$30\ms$ для руки, у петли через кору --- в полтора--три раза больше, у петли через глаз --- $100$--$200\ms$. А задержанная обратная связь, как мы видели в утверждении~\ref{prop:ei}, раскачивает систему. Для простейшего регулятора, исправляющего отклонение $x$ со скоростью $G$ по сведениям $d$ секунд давности, $\dd x/\dd t=-G\,x(t-d)$, условие устойчивости таково~\cite{Hayes1950}:
\begin{empheq}[box=\keybox]{equation}
G\,d \;<\; \frac{\pi}{2} ,
\label{eq:delaystable}
\end{empheq}
а на границе система колеблется с частотой $1/(4d)$. Мы проверили это на модели: при $d=30\ms$ отклонение затухает при $G=40$ в секунду, а при $G=55$ в секунду, чуть выше границы $52$, раскачивается.

Отсюда два следствия. Первое: чем длиннее петля, тем медленнее она обязана исправлять. Через глаз с задержкой в полторы сотни миллисекунд руку нельзя поправлять быстрее, чем примерно за десятую долю секунды, иначе она задрожит. Второе: граница устойчивости спинномозговой петли, $1/(4\cdot30\ms)\approx8$ колебаний в секунду, близка к частоте мелкой дрожи, которая есть у каждого здорового человека, --- восемь-двенадцать колебаний в секунду. Петля растяжения --- один из вероятных источников этой дрожи~\cite{McAuleyMarsden2000}.

Как тогда мозг двигается быстро и точно, если обратная связь запаздывает? По распространённой гипотезе --- \emph{предсказывает}: держит внутреннюю модель тела и вычисляет, каким будет результат команды, не дожидаясь датчиков~\cite{WolpertGhahramani2000}. Датчики нужны, чтобы исправлять модель, а не каждое движение. Инженер назвал бы это управлением с упреждением и наблюдателем состояния.

\section{Генераторы ритма движений}

Ходьба, дыхание, жевание --- ритмические движения. Нужны ли для них часы? Нет. В спинном мозге и стволе есть небольшие сети, которые сами порождают ритм, даже если на них не приходит ничего ритмического~\cite{MarderBucher2001}. Простейшая такая сеть собрана из деталей, которые у нас уже есть: две группы \nt{GLUT}-нейронов, которые тормозят друг друга через \nt{GABA}- или \nt{GLYC}-посредников, как на рис.~\ref{fig:wta}, и медленно устают, как память главы~\ref{sec:glutcomputer}:
\begin{empheq}[box=\keybox]{equation}
\tau\,\frac{\dd r_i}{\dd t} = -r_i + \bigl[I - \beta\,r_j - a_i\bigr]_+ ,
\qquad
\tau_a\,\frac{\dd a_i}{\dd t} = -a_i + b\,r_i ,
\label{eq:halfcentre}
\end{empheq}
где $a_i$ --- усталость группы $i$, $j$ --- другая группа. Взаимное торможение с $\beta>1$ выбирает победителя (утверждение~\ref{prop:wta}). Победитель работает и устаёт; когда усталость съедает его вход, проигравший, уже отдохнувший, вырывается вперёд и сам начинает уставать. Группы работают по очереди, и каждая управляет своей мышцей пары (рис.~\ref{fig:halfcentre}).

\begin{figure}[t]
\centering
\begin{tikzpicture}[>=Stealth,x=3.2cm,y=2.4cm]
\draw[->,gray!70!black] (0,0) -- (4.2,0) node[right,font=\small,black] {$t$, с};
\draw[->,gray!70!black] (0,0) -- (0,1.2) node[left,font=\small,black] {$r$};
\foreach \t in {0,1,2,3,4}{\draw[gray!60!black] (\t,-0.03) -- (\t,0.03); \node[below,font=\scriptsize] at (\t,0) {\t};}
\draw[thick,blue!60!black] plot file {figures/halfcentre1.dat};
\draw[thick,orange!85!black] plot file {figures/halfcentre2.dat};
\node[font=\scriptsize,blue!60!black,anchor=west] at (1.6,1.28) {группа 1: «вперёд»};
\node[font=\scriptsize,orange!85!black,anchor=west] at (2.7,1.28) {группа 2: «назад»};
\end{tikzpicture}
\caption{Генератор ритма по модели~\eqref{eq:halfcentre}: $I=1$, $\beta=2$, $b=2$, $\tau=20\ms$, $\tau_a=0.4\,\text{с}$. Две группы со взаимным торможением и усталостью работают по очереди с периодом $0.84\,\text{с}$, хотя на вход подаётся постоянный сигнал.}
\label{fig:halfcentre}
\end{figure}

В нашей модели при постоянном входе группы сменяют друг друга с периодом $0.84\,\text{с}$, примерно вдвое больше времени усталости $\tau_a$. Постоянная команда сверху --- «иди» --- превращается в ритм на месте. Это ещё один местный ритм, как ритм петли \nt{GLUT}--\nt{GABA} в главе~\ref{sec:gaba}: он возникает, только когда сеть работает, и задаёт такт лишь тем мышцам, которыми управляет.

\begin{keyblock}
\begin{proposition}[Выход машины]
\label{prop:muscles}
Мозг управляет мышцами как иерархия регуляторов. Наверху выбирается действие (глава~\ref{sec:dofa}); ниже местные генераторы превращают постоянные команды в ритм, а быстрые петли спинного мозга держат суставы. Сила задаётся в плавающей точке: порядок --- числом единиц, включаемых по размеру, мантисса --- частотой импульсов. Каждый сустав управляется парой тянущих мышц, разность сил которых задаёт момент, а сумма --- жёсткость. Эти устройства измерены; что мозг опирается на внутреннюю модель тела, --- хорошо обоснованная гипотеза.
\end{proposition}
\end{keyblock}

\section{Итог}

\begin{table}[t]
\centering
\footnotesize
\setlength{\tabcolsep}{4pt}
\renewcommand{\arraystretch}{1.15}
\begin{tabular}{>{\raggedright}p{3.0cm}>{\raggedright}p{4.6cm}>{\raggedright\arraybackslash}p{5.2cm}}
\toprule
Что делает выход & Как & Что известно \\
\midrule
превращает импульсы в силу & сумма подёргиваний~\eqref{eq:forcesum}, фильтр нижних частот & измерено; линейная сумма --- упрощение \\
дозирует силу & включение единиц по размеру, плавающая точка~\eqref{eq:sizeprinciple} & порядок включения измерен \\
толкает, умея только тянуть & пара мышц: момент --- разность, жёсткость --- сумма~\eqref{eq:antagonists} & измерено \\
держит сустав & петля растяжения, запрет противника через \nt{GLYC} & измерено \\
не дрожит & $Gd<\pi/2$~\eqref{eq:delaystable} & следует из теории регулирования; вклад в дрожь --- вероятная причина \\
двигается быстро & предсказание по внутренней модели & хорошо обоснованная гипотеза \\
ходит и дышит ритмично & генератор ритма~\eqref{eq:halfcentre} & генераторы измерены; простейшая схема --- модель \\
\bottomrule
\end{tabular}
\caption{Как машина управляет мышцами.}
\label{tab:muscles}
\end{table}

Выход машины устроен теми же приёмами, что и всё остальное (таблица~\ref{tab:muscles}): два провода вместо знака, фильтр нижних частот вместо ЦАП, логарифмическая шкала вместо линейной, взаимное торможение и усталость вместо тактового генератора. Вход (глава~\ref{sec:sensors}) и выход меряют мир в относительных долях, а между ними работает машина, которой посвящена эта книга.

\section{Задачи}

\begin{exercise}[\lvlA]
\label{ex:13:1}
\emph{Плавающая точка в цифрах.} В мышце $N=100$ двигательных единиц с силами $f_n=f_1q^{\,n-1}$; самая сильная в $100$ раз сильнее самой слабой. (а)~Найдите $q$, относительный шаг~\eqref{eq:sizeprinciple} и динамический диапазон в децибелах. (б)~Каков шаг при силе $10f_1$ у «линейного ЦАП» из $100$ одинаковых единиц той же суммарной силы?
\end{exercise}

\begin{exercise}[\lvlA]
\label{ex:13:2}
\emph{Цена запаса надёжности.} На каждый импульс синапс на мышечном волокне выбрасывает пуассоновское число порций со средним $m$, а волокно срабатывает при $n_{\text{пор}}=20$ порциях и больше; запас надёжности $S=m/n_{\text{пор}}$. Сколько сбоев в сутки даёт единица, работающая с частотой $20$ импульсов в секунду, при $S=2$ и при $S=4$?
\end{exercise}

\begin{exercise}[\lvlB]
\label{ex:13:3}
\emph{Мышца как фильтр.} Подёргивание~\eqref{eq:twitch} с $T_c=50\ms$ --- импульсная характеристика линейной системы. (а)~Найдите её передаточную функцию $H(s)$ и частоту среза. (б)~При какой частоте регулярных импульсов размах пульсаций силы составляет $10\%$ от средней силы?
\end{exercise}

\begin{exercise}[\lvlB]
\label{ex:13:4}
\emph{Быстрее задержки не исправить.} Регулятор $\dd x/\dd t=-G\,x(t-d)$. (а)~Покажите, что быстрее всего отклонение затухает без колебаний при $Gd=1/e$, со скоростью $1/d$. (б)~Найдите это $G$ и постоянную времени затухания для спинномозговой петли, $d=30\ms$, и для петли через глаз, $d=150\ms$.
\end{exercise}

\begin{exercise}[\lvlB]
\label{ex:13:5}
\emph{Период генератора ритма.} При $\tau\ll\tau_a$ период $T$ модели~\eqref{eq:halfcentre} задаётся уравнением $(\beta-1)(1-E^{2+b})/(\beta-E)=b\,(1-E^{1+b})/(1+b)$, где $E=e^{-T/(2\tau_a)}$. (а)~Найдите $T$ при $\beta=b=2$, $\tau_a=0.4$~с. (б)~Покажите, что $T$ не зависит от входа $I$, а ритм возможен только при $b>\beta-1$.
\end{exercise}

\begin{exercise}[\lvlB]
\label{ex:13:6}
\emph{Моделирование генератора.} Импортируйте функцию \texttt{halfcentre} из \texttt{models/\allowbreak muscle\_models.py} (сам файл не запускайте) и измерьте период, отбросив первые два цикла, при $b=0.8$, $1.05$, $1.2$, $1.5$, $2$, $4$ и при $I=0.5$, $1$, $2$. Может ли команда «иди быстрее» менять темп через $I$?
\end{exercise}

\begin{exercise}[\lvlB]
\label{ex:13:7}
\emph{Жёсткость против рефлекса.} Сустав охвачен петлёй растяжения: $\dd\theta/\dd t=-a\theta(t)-G\theta(t-d)$, $a=K/B$, $d=30\ms$. Сведите её к задаче~\ref{ex:13:4} и найдите наименьшее $a$, при котором отклонение затухает без колебаний с постоянной времени $15\ms$, и нужное $G$. Как менять $G$ при усилении совместного напряжения мышц?
\end{exercise}

\begin{exercise}[\lvlC]
\label{ex:13:8}
\emph{Ручка темпа.} По задачам~\ref{ex:13:5} и~\ref{ex:13:6} вход $I$ генератора~\eqref{eq:halfcentre} меняет только размах, но не темп. Предложите минимальное изменение модели из деталей книги, при котором ниже некоторого $I_{\min}$ ритма нет, а выше период убывает с ростом $I$, хотя бы вдвое при утроении $I$. Найдите $I_{\min}$ при $\tau\ll\tau_a$ и проверьте моделированием.
\end{exercise}

\chapter{Боль: сигнал тревоги}
\chaptermark{Боль}
\label{sec:pain}
\begin{chaptergoals}
\item почему боль --- приоритетный сигнал тревоги, а не канал измерения;
\item как быстрый и медленный провода сообщают, где повреждение и насколько оно серьёзно;
\item как ворота в спинном мозге дают касанию запретить боль, а широковещательные каналы задают громкость;
\item как сенсибилизация повышает усиление после повреждения и как боль учит и прерывает работу.
\end{chaptergoals}

\section{Тревога, а не измерение}

Все датчики главы~\ref{sec:sensors} измеряли мир: сколько света, какой высоты звук, какое давление. Боль устроена иначе. Она сообщает не «что там», а «опасно, брось всё». Для инженера это не измерительный канал, а \emph{сигнал тревоги}: линия с высшим приоритетом, которая должна пробиться сквозь любую текущую работу машины.

От измерительного канала сигнализация отличается тремя свойствами, и все три у боли есть. Во-первых, её трудно приглушить привыканием: датчики боли привыкают гораздо слабее датчиков касания, а после лёгкого повреждения становятся даже чувствительнее~\cite{LaMotte1982hyper}; ослабление ответа после сильного нагрева измерено лишь у одного их типа~\cite{Treede1995heat}. Датчик света на постоянном фоне замолкает (рис.~\ref{fig:adaptation}), а сигнализация, которая замолчала через минуту, бесполезна. Во-вторых, у неё высокий порог: датчики боли срабатывают только тогда, когда воздействие близко к опасному, например на нагрев --- у разных датчиков пороги от $41^\circ$ до $46$--$48^\circ$, в зависимости от типа~\cite{Treede1995heat, Weidner1999cfib}. В-третьих --- и это главное для этой главы --- громкость тревоги решает не только датчик, но и сама машина. Одна и та же рана болит по-разному в разной обстановке. Посмотрим, из каких уже знакомых деталей это собрано.

Датчики боли --- \nt{GLUT}-нейроны, как и остальные чувствительные нейроны главы~\ref{sec:sensors}. Часть из них вместе с глутаматом выбрасывает вещество~P из приложения~\ref{sec:othermediators}, медленно усиливающее передачу.

\section{Два провода: быстрая и медленная боль}

Ударьте палец --- и вы почувствуете боль дважды: сначала короткую и резкую, а через секунду тупую и долгую. Это два разных провода. Быстрый провод изолирован и проводит импульсы со скоростью $v$ от $5$ до $30$ метров в секунду; медленный тонкий и голый, $0.5$--$2$ метра в секунду. Сигнал от ступни проходит около метра, и время прихода
\begin{empheq}[box=\keybox]{equation}
t \;=\; \frac{L}{v}
\label{eq:paintime}
\end{empheq}
для быстрого провода --- десятки миллисекунд, для медленного --- около секунды. Два провода несут два сообщения. Быстрый говорит \emph{где} и \emph{когда}: его импульсы приходят раньше и точнее, как в коде временем главы~\ref{sec:code}. Медленный говорит \emph{насколько плохо}, и его долгая тупая боль держит машину в режиме тревоги, пока повреждение не прошло.

\section{Ворота в спинном мозге}

Первое, куда приходит сигнал боли, --- спинной мозг, и там он проходит через \emph{ворота}~\cite{MelzackWall1965}; конкретные нейроны этих ворот найдены сравнительно недавно~\cite{Duan2014}. На выходной \nt{GLUT}-нейрон, который передаёт боль дальше в мозг, сходятся провод боли и провод прикосновения. А рядом сидит тормозящий нейрон-посредник, \nt{GABA} или \nt{GLYC}, который прикосновение возбуждает (рис.~\ref{fig:gate}). Прикосновение закрывает ворота. В исходной теории ворот~\cite{MelzackWall1965} боль, наоборот, выключает этого посредника, и сильная боль ворота открывает; это допущение теории, на найденных нейронах ворот оно прямо не показано.

\begin{figure}[t]
\centering
\begin{tikzpicture}[>=Stealth,
  nrn/.style={circle,draw,very thick,minimum size=1.0cm,inner sep=0pt,font=\scriptsize},
  inh/.style={-{Bar[width=7pt]},thick,red!70!black},
  bc/.style={->,thick,densely dashed,orange!85!black}]
\node[nrn,fill=blue!6] (Z) at (6,0) {\nt{GLUT}};
\node[nrn,fill=red!10] (Q) at (3.6,-1.6) {\nt{GABA}};
\draw[->,very thick,red!70!black] (0,0.6) node[left,font=\small,black,align=right] {боль $\nu$} -- (5.45,0.15);
\draw[->,very thick,blue!60!black] (0,-2.6) node[left,font=\small,black,align=right] {касание $\mu$} -- (2.9,-2.0);
\draw[->,thick,blue!60!black] (1.8,-2.35) -- (5.5,-0.35) node[pos=0.8,below right=-2pt,font=\scriptsize] {$w_{z\mu}$};
\draw[inh] (1.6,0.5) -- (3.35,-1.1) node[pos=0.5,left=1pt,font=\scriptsize] {$w_{q\nu}$};
\draw[inh] (Q) -- (Z) node[pos=0.5,above left=-1pt,font=\scriptsize] {$\beta$};
\draw[->,very thick] (Z) -- (8.2,0) node[right,font=\small,align=left] {в мозг: $z$};
\draw[bc] (6,2.1) node[above,font=\scriptsize,black,align=center] {сверху: \nt{SERO}, \nt{NORA}, опиоиды} -- (Z.north) node[pos=0.45,right,font=\scriptsize,black] {усиление $g$};
\end{tikzpicture}
\caption{Ворота боли в спинном мозге по модели~\eqref{eq:gate}. Выходной \nt{GLUT}-нейрон получает боль $\nu$ и слабое возбуждение от касания $\mu$. Тормозящий посредник (\nt{GABA} или \nt{GLYC}) возбуждается касанием и, по допущению теории ворот, выключается болью. Сверху, по широковещательным каналам, приходит усиление $g$.}
\label{fig:gate}
\end{figure}

Запишем это для частот, как в главе~\ref{sec:gaba}. Частота посредника $q$ и выходная частота $z$ равны
\begin{empheq}[box=\keybox]{equation}
q \;=\; \bigl[w_{q\mu}\,\mu + q_0 - w_{q\nu}\,\nu\bigr]_+ ,
\qquad
z \;=\; \bigl[\,g\,(\nu + w_{z\mu}\,\mu) - \beta\,q\,\bigr]_+ ,
\label{eq:gate}
\end{empheq}
где $\nu$ --- вход боли, $\mu$ --- вход касания, $w$ --- веса связей (первый индекс --- кому, второй --- от кого), $\beta$ --- сила торможения, $q_0$ --- собственная фоновая активность посредника, $g$ --- усиление, которое задают сверху (о нём ниже). Это запрет~\eqref{eq:veto} из главы~\ref{sec:gaba}: «боль, но не при касании», --- с одной поправкой, взятой из теории ворот как допущение: сильная боль сама выключает запрещающий нейрон.

Мы посчитали модель при $w_{q\mu}=1$, $q_0=0.2$, $w_{q\nu}=0.5$, $w_{z\mu}=0.3$, $\beta=1.5$ (рис.~\ref{fig:gatecurves}). Без касания боль начинает проходить с $\nu\approx0.18$. С касанием порог поднимается до $0.86$, и при $\nu=1$ выход падает вчетверо, с $1$ до $0.25$. А при сильной боли, $\nu=2$, касание уже ничего не меняет: ворота распахнуты. Так и в жизни: потереть ушиб помогает, а при серьёзной травме --- нет. Само касание, без боли, через ворота не проходит: тревога не поднимается от прикосновения.

\begin{figure}[t]
\centering
\begin{tikzpicture}[>=Stealth,x=5.2cm,y=1.35cm]
\draw[->,gray!70!black] (0,0) -- (2.12,0) node[right,font=\small,black] {боль $\nu$};
\draw[->,gray!70!black] (0,0) -- (0,4.3) node[left,font=\small,black] {$z$};
\foreach \t in {0,0.5,1,1.5,2}{\draw[gray!60!black] (\t,-0.05) -- (\t,0.05); \node[below,font=\scriptsize] at (\t,0) {\t};}
\foreach \f in {1,2,3,4}{\draw[gray!60!black] (-0.01,\f) -- (0.01,\f); \node[left,font=\scriptsize] at (0,\f) {\f};}
\draw[very thick,red!70!black] plot file {figures/pain_gate_notouch.dat};
\draw[very thick,blue!60!black] plot file {figures/pain_gate_touch.dat};
\draw[very thick,orange!85!black,densely dashed] plot file {figures/pain_gate_sens.dat};
\draw[very thick,red!70!black] (0.08,3.9) -- (0.2,3.9) node[right,font=\scriptsize,black] {без касания};
\draw[very thick,blue!60!black] (0.08,3.45) -- (0.2,3.45) node[right,font=\scriptsize,black] {с касанием};
\draw[very thick,orange!85!black,densely dashed] (0.08,3.0) -- (0.2,3.0) node[right,font=\scriptsize,black] {после сенсибилизации, $g=2$};
\end{tikzpicture}
\caption{Выход ворот~\eqref{eq:gate} в зависимости от силы боли. Касание поднимает порог и приглушает слабую и среднюю боль, но не сильную. Сенсибилизация (штриховая линия) вдвое увеличивает усиление: та же боль передаётся вдвое громче, а порог опускается.}
\label{fig:gatecurves}
\end{figure}

\section{Ручка громкости сверху}

Ворота управляются не только снизу. Из ствола мозга к ним идут нисходящие пути, которые меняют усиление $g$ в~\eqref{eq:gate}: \nt{SERO}, \nt{NORA} и опиоидные пептиды из приложения~\ref{sec:othermediators}~\cite{Fields2004}. Это те же широковещательные каналы, что и в главах~\ref{sec:serocomputer}--\ref{sec:acet}, только здесь их ручка --- громкость тревоги. Они умеют крутить её в обе стороны: одни и те же нисходящие цепи могут и приглушать боль, и усиливать её.

Зачем машине убавлять тревогу? Затем, что тревога --- это прерывание, а прерывание не всегда уместно. Животное, убегающее от хищника, не должно останавливаться из-за раненой лапы: сначала бежать, потом зализывать рану. Ожидание облегчения тоже приглушает боль, и часть этого эффекта снимается, если заблокировать опиоидные рецепторы~\cite{Levine1978}: ожидание, вычисленное в мозге, спускается по опиоидному каналу к воротам. Сама по себе эта картина измерена; как именно мозг решает, какой должна быть громкость, --- гипотеза.

\section{Сенсибилизация: тревога, которая учится}

После повреждения выходные нейроны ворот становятся чувствительнее: тот же вход вызывает больший выход, а порог опускается~\cite{Woolf1983}. В модели~\eqref{eq:gate} это рост усиления $g$, и простейшее его описание --- медленная RC-цепь, которую накачивает сама боль:
\begin{empheq}[box=\keybox]{equation}
\tau_s\,\frac{\dd g}{\dd t} \;=\; -(g-1) + k_s\,z ,
\label{eq:sensitization}
\end{empheq}
где $\tau_s$ --- время, за которое чувствительность возвращается к норме; оно больше длительности самой боли, потому что сенсибилизация держится и после того, как повреждающий стимул прекратился~\cite{Woolf1983}, а $k_s$ --- сила накачки. Пока ткань повреждена, выход ворот держит $g$ выше единицы, и всё, что происходит рядом с раной, передаётся громче (штриховая линия на рис.~\ref{fig:gatecurves}: при $g=2$ порог опускается до $0.11$, а выход удваивается). Это разумная настройка сигнализации: пока повреждение не зажило, лучше перестраховаться.

Для инженера это положительная обратная связь с медленной утечкой: боль повышает усиление, усиление --- боль. Пока петля слабая, $g$ возвращается к норме, когда рана зажила. Если же петля сильная, сигнализация может застрять во включённом состоянии и после того, как повреждения уже нет, --- как группа с сильной обратной связью на рис.~\ref{fig:bistable}. Модель~\eqref{eq:sensitization} --- упрощение; то, что центральная сенсибилизация существует, измерено.

\section{Боль как учитель и как прерывание}

У боли в машине две работы помимо сигнализации.

\emph{Учитель.} Боль --- отрицательная награда: в формуле ошибки~\eqref{eq:ctd} она входит как $r<0$. Всё, что говорилось в главе~\ref{sec:dofa}, работает и здесь: предвестник боли начинает вызывать ошибку заранее, и сеть учится избегать того, что к боли ведёт. Опыты на людях показывают, что активность мозга при обучении на боли следует именно ошибке временных разностей~\cite{Seymour2004}, а часть \nt{DOPA}-нейронов отвечает и на неприятные события (глава~\ref{sec:dofaalt}).

\emph{Прерывание.} Боль отрывает внимание от текущей задачи --- это её назначение, а не побочный эффект~\cite{EcclestonCrombez1999}. В главе~\ref{sec:nora} та же роль «прерывания» обсуждалась для отрывистой пачки \nt{NORA}. А самый быстрый ответ даже не доходит до мозга: отдёргивание руки от горячего замыкается прямо в спинном мозге той же парой мышц и тем же запретом противника, что на рис.~\ref{fig:antagonists}.

\begin{keyblock}
\begin{proposition}[Сигнал тревоги]
\label{prop:pain}
Боль --- не измерение, а сигнал тревоги с высоким приоритетом. Она привыкает гораздо слабее измерительных каналов, идёт по двум проводам (быстрому --- «где», медленному --- «насколько плохо»), проходит через ворота, которые закрывает касание (а сильная боль, по допущению теории ворот, открывает), и её громкость задают сверху широковещательные каналы. Эти устройства измерены, кроме открывания ворот болью; простые модели~\eqref{eq:gate} и~\eqref{eq:sensitization} --- упрощения.
\end{proposition}
\end{keyblock}

\section{Итог}

\begin{table}[t]
\centering
\footnotesize
\setlength{\tabcolsep}{4pt}
\renewcommand{\arraystretch}{1.15}
\begin{tabular}{>{\raggedright}p{3.0cm}>{\raggedright}p{4.9cm}>{\raggedright\arraybackslash}p{4.9cm}}
\toprule
Что делает боль & Как & Что известно \\
\midrule
поднимает тревогу & высокий порог, привыкание слабее, чем у касания & порог измерен; сравнение привыкания --- оценка \\
сообщает «где» и «насколько плохо» & два провода с разной скоростью~\eqref{eq:paintime} & измерено \\
проходит через ворота & запрет через \nt{GABA}/\nt{GLYC}~\eqref{eq:gate} & ворота измерены; открывание болью --- допущение теории; модель --- упрощение \\
меняет громкость & нисходящие \nt{SERO}, \nt{NORA}, опиоиды & измерено; правило выбора громкости --- гипотеза \\
учится после повреждения & рост усиления~\eqref{eq:sensitization} & сенсибилизация измерена; модель --- упрощение \\
учит избегать & отрицательная награда в~\eqref{eq:ctd} & сходство с ошибкой временных разностей измерено \\
прерывает работу & захват внимания; рефлекс отдёргивания & измерено \\
\bottomrule
\end{tabular}
\caption{Боль как сигнал тревоги.}
\label{tab:pain}
\end{table}

Боль собрана из деталей, которые мы уже видели (таблица~\ref{tab:pain}): \nt{GLUT}-датчики, два провода, запрет через тормозящего посредника, широковещательная ручка громкости, медленная RC-цепь, ошибка предсказания и прерывание. Новое в ней одно: это единственный вход машины, громкость которого машина задаёт сама, --- сигнализация, которую владелец может и приглушить, и перенастроить.

\section{Задачи}

\begin{exercise}[\lvlA]
\label{ex:14:1}
\emph{Два провода.} Сигнал идёт от ступни, $L=1$~м. (а)~Залп идёт по пучку проводов одного типа, скорости которых заполняют диапазон~\eqref{eq:paintime}. На сколько он растягивается к приходу в спинной мозг для быстрого и медленного типа? (б)~Волны боли по проводам $15$ и $1$~м/с различимы при разнице от $0.1$~с. С какого расстояния они сливаются?
\end{exercise}

\begin{exercise}[\lvlB]
\label{ex:14:2}
\emph{Когда касание болит.} Ворота~\eqref{eq:gate} с параметрами текста, усиление $g$ растёт при сенсибилизации. До какого $g$ касание без боли не проходит через ворота ни при какой силе $\mu$? При каком $g$ начинает проходить касание $\mu=1$?
\end{exercise}

\begin{exercise}[\lvlB]
\label{ex:14:3}
\emph{Устойчивость сенсибилизации.} Входы постоянны, выход ворот линеен по $g$: $z=gA-C$, $A=\nu+w_{z\mu}\mu$, $C=\beta q\ge0$. (а)~Найдите равновесие $g^*$ модели~\eqref{eq:sensitization} и условие его устойчивости. (б)~При $k_s=0.5$ найдите силу боли, выше которой модель уходит вразнос, без касания и с касанием $\mu=1$.
\end{exercise}

\begin{exercise}[\lvlA]
\label{ex:14:4}
\emph{Предвестник боли.} Сигнал в момент $0$ предсказывает боль в момент $T=2$~с, в~\eqref{eq:ctd} $r(t)=-P\,\delta(t-T)$, $\tau_\gamma=2$~с. (а)~Найдите площадь всплеска $\delta$ в момент сигнала. (б)~Действие в момент $t_e=1$~с отменяет боль. Каков всплеск $\delta$ в этот момент и чему он научит сеть?
\end{exercise}

\begin{exercise}[\lvlB]
\label{ex:14:5}
\emph{Сигнализация, которая защёлкивается.} Дополните ворота из \texttt{models/\allowbreak pain\_gate.py} динамикой~\eqref{eq:sensitization} и насыщением $z\le4$. Касание $\mu=1$ есть всегда, повреждение $\nu=2$ длится время $T$, затем $\nu=0$. Моделированием найдите наименьшее $T$ (в единицах $\tau_s$), после которого сигнализация остаётся включённой, при $k_s=4$, $4.6$, $5$, $6$, $8$.
\end{exercise}

\begin{exercise}[\lvlB]
\label{ex:14:6}
\emph{Проектирование ворот.} При $\beta=1.5$, $w_{z\mu}=0.3$, $g=1$ подберите в~\eqref{eq:gate} $q_0$, $w_{q\nu}$, $w_{q\mu}$ так, чтобы порог боли был $0.2$ без касания и $1.0$ с касанием $\mu=1$, а касание переставало приглушать боль при $\nu_\times=2.5$. До какого усиления $g$ касание без боли не проходит через ворота?
\end{exercise}

\begin{exercise}[\lvlC]
\label{ex:14:7}
\emph{Ручка громкости.} Работа приносит ценность $V$ в единицу времени, работа с повреждением $\nu$ стоит $c\nu$, так что прерывать выгодно при $\nu>V/c$; боль прерывает при $z>\theta=0.5$. (а)~Какое усиление $g^*$ ворот~\eqref{eq:gate} без касания даёт этот порог при $V/c=0.25$, $0.5$, $2$? (б)~Из каких сигналов книги машина могла бы оценивать $V$ и $c$?
\end{exercise}

\chapter{Как машина учится двигаться}
\chaptermark{Как машина учится двигаться}
\label{sec:motorlearning}

\section{Три учителя}

В главе~\ref{sec:muscles} мы видели, как машина двигает мышцами, и остановились на гипотезе: быстрые и точные движения требуют внутренней модели тела, которая предсказывает результат команды~\cite{WolpertGhahramani2000}. Откуда берётся эта модель? Её надо выучить, и выучить не так, как мы учились до сих пор.

В книге уже было два учителя. Первый --- \emph{никакой}: правило Хебба (глава~\ref{sec:dofa}) усиливает то, что часто бывает вместе, и сжимает входы. Второй --- \emph{награда}: \nt{DOPA} рассылает всем одно число, «лучше или хуже, чем ожидалось». Движению нужен третий учитель --- \emph{ошибка}. Рука промахнулась мимо чашки на два сантиметра влево: это не «хуже», это вектор, который говорит каждому выходу, куда и насколько он ошибся. Инженер назвал бы это обучением с учителем.

Разница между вторым и третьим учителем --- это разница между одним проводом на всех и отдельным проводом на каждый выход. В утверждении~\ref{prop:broadcastcost} обучение по одному общему числу замедлялось с каждым нейроном, чей шум влиял на результат. Если же каждый выход $i$ получает свою ошибку $e_i$, веса можно менять по простейшему правилу:
\begin{empheq}[box=\keybox]{equation}
\frac{\dd w_{ij}}{\dd t} \;=\; -\eta\,e_i(t)\,x_j(t) .
\label{eq:lms}
\end{empheq}
Это правило наименьших квадратов, давно известное в технике адаптивных фильтров~\cite{WidrowHoff1960}: вес меняется пропорционально произведению своего входа на ошибку своего выхода и в среднем идёт по градиенту квадрата ошибки. Никакого шума-поиска, как в главе~\ref{sec:dofa}, не нужно: направление поправки приходит прямо по проводу ошибки. И скорость не падает с числом выходов, потому что чужие ошибки к синапсу не приходят.

\begin{keyblock}
\begin{proposition}[Три учителя]
\label{prop:teachers}
Правило Хебба учит без учителя тому, что часто; сигнал \nt{DOPA} учит одним числом на всю сеть тому, что выгодно, и медленнее с каждым нейроном; провод ошибки к каждому выходу учит по правилу~\eqref{eq:lms} тому, что точно, и скорость не зависит от числа выходов. Цена третьего способа --- отдельный провод на каждый выход и кто-то, кто знает правильный ответ.
\end{proposition}
\end{keyblock}

\section{Провод ошибки}

Кто может знать правильный ответ для движения? Сами датчики. Если рука ушла влево от цели, об этом сообщают глаз и датчики растяжения из главы~\ref{sec:sensors}. Нужна схема, которая доставит эту ошибку к нужным выходам.

В мозге есть область, которая, по-видимому, устроена ровно под это~\cite{Marr1969,Albus1971}. Её схема необычно правильная, почти кристаллическая, и в ней легко узнать правило~\eqref{eq:lms} (рис.~\ref{fig:errorwire}).

\emph{Расширитель входа.} Сигналы о команде и о состоянии тела приходят на огромный слой маленьких \nt{GLUT}-нейронов. Их около семидесяти миллиардов --- больше, чем во всём остальном мозге, вместе взятом~\cite{Azevedo2009}. Каждый из них смешивает несколько входов, и сигнал раскладывается в очень длинный вектор. Инженеру этот приём знаком: если поднять размерность входа, в новом пространстве почти любую нужную зависимость можно получить простой взвешенной суммой~\cite{Cover1965}.

\emph{Выходные нейроны.} Взвешенную сумму считают около тридцати миллионов больших выходных нейронов~\cite{Andersen1992cereb}, у каждого --- порядка ста тысяч входов от расширителя. И это \nt{GABA}-нейроны: результат обучения передаётся дальше не возбуждением, а торможением, как в запрете главы~\ref{sec:gaba}.

\emph{Провод ошибки.} Каждый выходной нейрон получает ещё один вход, особый: ровно один \nt{GLUT}-провод, который срабатывает редко, около раза в секунду, но каждый раз вызывает в выходном нейроне мощную вспышку~\cite{Ito2001}. Это и есть $e_i$: вероятность вспышки зависит от направления ошибки движения~\cite{Herzfeld2018}. Совпадение вспышки ошибки с активностью входа ослабляет этот вход --- в точности член $-\eta\,e_i x_j$ в~\eqref{eq:lms}.

\begin{figure}[t]
\centering
\begin{tikzpicture}[>=Stealth,
  gran/.style={circle,draw,thick,fill=blue!6,minimum size=0.32cm,inner sep=0pt},
  outn/.style={circle,draw,very thick,fill=red!10,minimum size=1.1cm,inner sep=0pt,font=\scriptsize},
  inh/.style={-{Bar[width=7pt]},very thick,red!70!black}]
% inputs
\foreach \y/\lab in {1.6/{команда},0.4/{состояние тела}}{
  \draw[->,thick,blue!60!black] (-0.6,\y) node[left,font=\scriptsize,black] {\lab} -- (0.35,\y);}
% expansion layer
\foreach \i in {0,...,11}{\node[gran] (g\i) at ({0.8+0.28*mod(\i,2)},{-0.3+0.23*\i}) {};}
\foreach \i in {0,...,11}{\pgfmathsetmacro{\yy}{mod(\i,3)>0 ? 1.6 : 0.4}\draw[gray!60!black] (0.35,\yy) -- (g\i);}
\node[font=\scriptsize,align=center] at (0.95,-1.05) {расширитель\\$\sim7\cdot10^{10}$ \nt{GLUT}};
% parallel wires to the output neuron
\foreach \i in {0,...,11}{\draw[->,gray!70!black] (g\i.east) -- (4.1,{1.0+0.07*(\i-5.5)});}
\node[outn] (P) at (4.6,1.0) {\nt{GABA}};
\node[font=\scriptsize,align=center,above=2pt] at (P.north) {выходной нейрон\\$\sim10^5$ входов $x_j$, веса $w_{ij}$};
\draw[inh] (P) -- (6.8,1.0) node[right,font=\scriptsize,align=left] {поправка\\к движению};
% error wire
\draw[->,very thick,red!70!black,densely dashed] (4.6,-1.4) node[below,font=\scriptsize,black,align=center] {один провод ошибки $e_i$\\\nt{GLUT}, $\sim1$ раз/с} -- (P.south);
\end{tikzpicture}
\caption{Схема обучения с учителем, как её, по-видимому, реализует мозг. Команда и состояние тела раскладываются огромным слоем \nt{GLUT}-нейронов в длинный вектор $x_j$; выходной \nt{GABA}-нейрон складывает его с весами $w_{ij}$. Единственный провод ошибки приносит $e_i$; совпадение ошибки с активным входом ослабляет вес этого входа, как в правиле~\eqref{eq:lms}.}
\label{fig:errorwire}
\end{figure}

Одна трудность знакома нам по главе~\ref{sec:dofa}: ошибка приходит позже, чем вход, который к ней привёл. Промах виден через десятки-сотни миллисекунд после команды. Значит, синапс должен помнить, что был активен, --- нужен след, как в правиле~\eqref{eq:threefactor}. И длина этого следа, судя по опытам, подстроена под задержку обратной связи в каждой задаче: там, где ошибка приходит через сто миллисекунд, сильнее всего ослабляется вход, бывший активным за сто миллисекунд до неё~\cite{Suvrathan2016}.

\begin{keyblock}
\begin{proposition}[Провод ошибки]
\label{prop:errorwire}
В схеме, изображённой на рис.~\ref{fig:errorwire}, каждый выходной нейрон получает ровно один провод ошибки, и совпадение ошибки с входом ослабляет этот вход. Это правило наименьших квадратов~\eqref{eq:lms}, применённое к входу, который расширен до огромной размерности. Устройство схемы и ослабление при совпадении измерены; что вся схема работает как обучение с учителем, --- ведущая гипотеза.
\end{proposition}
\end{keyblock}

\section{Внутренняя модель как адаптивный фильтр}

Чему учится такая схема? Самый естественный ответ --- предсказанию. Пусть на вход приходит команда $u(t)$, пропущенная через набор фильтров с разными постоянными времени, как в расширителе: $x_j(t)=(g_j*u)(t)$. Выход --- их взвешенная сумма, предсказание того, что сообщат датчики:
\begin{empheq}[box=\keybox]{equation}
\hat y(t) \;=\; \sum_j w_j\,(g_j*u)(t), \qquad e(t) \;=\; \hat y(t) - y(t) ,
\label{eq:adaptivefilter}
\end{empheq}
а веса учатся по~\eqref{eq:lms}. Это \emph{адаптивный фильтр}: схема, которая сама подбирает свою передаточную функцию так, чтобы повторить неизвестную систему~\cite{Fujita1982}. Здесь неизвестная система --- собственное тело с его массой, упругостью и задержками. Выученный фильтр и есть внутренняя модель главы~\ref{sec:muscles}: он говорит, что сообщат датчики, не дожидаясь их.

Такая модель полезна дважды. Во-первых, её предсказание можно подать вместо запаздывающей обратной связи, и тогда условие устойчивости~\eqref{eq:delaystable} перестаёт связывать руки: исправлять можно быстро, по модели. Во-вторых, разница между предсказанным и полученным --- это сигнал о \emph{неожиданном}. Всё, что машина сделала сама, модель предсказала и вычла; остаётся то, что сделал мир. Поэтому, по одной из гипотез, мы и не можем пощекотать сами себя~\cite{Blakemore1998}.

\section{Когда мир меняется: быстрое и медленное}

Проверим схему опытом, который легко поставить. Человек тянется к цели, а мир неожиданно меняется: например, на руку действует боковая сила. Первые движения промахиваются, потом промах уменьшается. Если силу убрать, рука сначала промахивается в \emph{другую} сторону --- модель выучила силу и продолжает её компенсировать. Это прямое свидетельство того, что выучена модель, а не просто «стало получаться».

Опыты показывают и более тонкую вещь: учатся одновременно два процесса --- быстрый, который быстро учится и быстро забывает, и медленный, который учится медленно и помнит долго~\cite{Smith2006}. Поправки обновляются после каждого движения, по событию:
\begin{empheq}[box=\keybox]{equation}
e = p - (x_f + x_s), \qquad x_f \leftarrow A_f\,x_f + B_f\,e, \qquad x_s \leftarrow A_s\,x_s + B_s\,e ,
\label{eq:tworate}
\end{empheq}
где $p$ --- возмущение, $x_f$ и $x_s$ --- выученные поправки двух процессов, $A$ --- насколько поправка сохраняется до следующего движения, $B$ --- насколько сильно она учится на ошибке. У быстрого процесса $B$ велико, а $A$ заметно меньше единицы; у медленного --- наоборот.

\begin{figure}[t]
\centering
\begin{tikzpicture}[>=Stealth,x=0.034cm,y=1.9cm]
\draw[->,gray!70!black] (0,0) -- (355,0) node[right,font=\small,black] {движение};
\draw[->,gray!70!black] (0,-1.15) -- (0,1.25);
\foreach \v in {-1,1}{\draw[gray!60!black] (-3,\v) -- (3,\v); \node[left,font=\scriptsize] at (0,\v) {$\v$};}
\foreach \n in {100,200,300}{\draw[gray!60!black] (\n,-0.04) -- (\n,0.04); \node[below,font=\scriptsize] at (\n,0) {\n};}
\draw[thin,gray!70!black] plot file {figures/tworate_p.dat};
\draw[thick,orange!85!black] plot file {figures/tworate_fast.dat};
\draw[thick,blue!60!black] plot file {figures/tworate_slow.dat};
\draw[very thick,black] plot file {figures/tworate_net.dat};
\node[font=\scriptsize,gray!50!black] at (120,1.12) {возмущение $p$};
\node[font=\scriptsize,black,anchor=west] at (238,0.52) {сумма: возврат};
\node[font=\scriptsize,blue!60!black,anchor=west] at (150,0.39) {медленный $x_s$};
\node[font=\scriptsize,orange!85!black,anchor=west] at (60,0.08) {быстрый $x_f$};
\draw[densely dashed,gray!60!black] (226,-1.15) -- (226,1.25);
\node[font=\scriptsize,align=center,anchor=south west] at (228,-1.1) {ошибок\\больше нет};
\end{tikzpicture}
\caption{Два процесса обучения по модели~\eqref{eq:tworate}, $A_f=0.92$, $B_f=0.1$, $A_s=0.996$, $B_s=0.02$. Двести движений с возмущением $p=1$, затем шесть с обратным, $p=-1$, пока сумма не вернётся к нулю. После штриховой линии ошибки перестают поступать, и сумма сама возвращается к старой поправке: быстрый процесс забыл обратное возмущение, медленный помнит прямое.}
\label{fig:tworate}
\end{figure}

Модель~\eqref{eq:tworate} объясняет опыт, который трудно понять без неё (рис.~\ref{fig:tworate}). В нашем расчёте за двести движений с возмущением сумма поправок доходит до $0.84$, причём большую часть держит медленный процесс, $0.63$. Затем шесть движений с обратным возмущением возвращают сумму к нулю: быстрый процесс ушёл в минус, $-0.53$, медленный ещё в плюсе, $0.46$. Если теперь перестать сообщать ошибки, быстрый процесс забывает свою поправку за десятки движений, медленный --- гораздо дольше, и сумма \emph{сама} поднимается до $0.37$ за сорок движений: старый навык возвращается без всякой тренировки. Такое спонтанное восстановление измерено у людей; модель с одним процессом его дать не может --- без ошибок она просто затухает к нулю.

Зачем два процесса? Правдоподобный инженерный ответ даёт оптимальный фильтр главы~\ref{sec:acet}. Тело меняется по разным причинам с разной скоростью: мышцы устают за минуты, растут и слабеют за месяцы. Если помехи бывают быстрыми и медленными, наилучшая оценка состоит из двух фильтров с разными постоянными времени, и каждый ловит свою~\cite{Kording2007}. Быстрый процесс --- это временная поправка, «рука устала»; медленный --- «рука стала другой». Это пока гипотеза, но она объясняет, почему навык, выученный за один день, частично теряется к следующему и почему второй раз он учится быстрее.

\section{Итог}

\begin{table}[t]
\centering
\footnotesize
\setlength{\tabcolsep}{4pt}
\renewcommand{\arraystretch}{1.15}
\begin{tabular}{>{\raggedright}p{3.0cm}>{\raggedright}p{4.6cm}>{\raggedright\arraybackslash}p{5.2cm}}
\toprule
Что делает схема & Как & Что известно \\
\midrule
учится с учителем & провод ошибки к каждому выходу, правило~\eqref{eq:lms} & устройство схемы и ослабление при совпадении измерены; обучение с учителем --- ведущая гипотеза \\
расширяет вход & семьдесят миллиардов \nt{GLUT}-нейронов & число нейронов измерено; роль расширения --- теория \\
учитывает задержку ошибки & след, подстроенный под задержку & измерено \\
строит внутреннюю модель & адаптивный фильтр~\eqref{eq:adaptivefilter} & хорошо обоснованная гипотеза \\
учится двумя темпами & быстрый и медленный процессы~\eqref{eq:tworate} & последействие и спонтанный возврат измерены; два процесса --- модель \\
\bottomrule
\end{tabular}
\caption{Как машина учится двигаться.}
\label{tab:motorlearning}
\end{table}

Теперь у машины три учителя (таблица~\ref{tab:motorlearning}). Хебб сжимает то, что часто; \nt{DOPA} одним числом учит выбирать выгодное; провод ошибки учит точности, и учит быстро, потому что к каждому выходу идёт своя ошибка. Мозг, по-видимому, пользуется всеми тремя, и каждым там, где он дешевле всего: широковещательным сигналом --- для немногих решений, отдельными проводами --- для точной подстройки тела, где правильный ответ сообщают сами датчики.

\section{Задачи}

\begin{exercise}[\lvlA]
\label{ex:15:1}
\emph{Ёмкость схемы с проводом ошибки.} В схеме $3\cdot10^7$ выходных нейронов по $10^5$ входов, у каждого свой провод ошибки ($1$~имп/с). Нейрон с $n$ входами выучивает любую разметку $\approx2n$ векторов~\cite{Cover1965}. (а)~Сколько ассоциаций хранит схема? (б)~Оцените по~\eqref{eq:spikeentropy} при $\Delta t=10\ms$ наименьшее время заполнения ёмкости нейрона.
\end{exercise}

\begin{exercise}[\lvlB]
\label{ex:15:2}
\emph{Скорость правила наименьших квадратов.} Моды правила~\eqref{eq:lms} затухают со скоростями $\eta\lambda_k(C)$. Для пяти фильтров~\eqref{eq:adaptivefilter} на белом шуме $C_{jk}=2\sqrt{\tau_j\tau_k}/(\tau_j+\tau_k)$. Найдите отношение самой быстрой и самой медленной скоростей, если $\tau_j$ растут в $2$ раза ($10$--$160\ms$) и в $4$ раза.
\end{exercise}

\begin{exercise}[\lvlB]
\label{ex:15:3}
\emph{Анализ двух процессов.} Запишите модель~\eqref{eq:tworate} с параметрами рис.~\ref{fig:tworate} как $\mathbf x_{n+1}=M\mathbf x_n+\mathbf b\,p$. (а)~Найдите по собственным числам $M$ постоянные времени в движениях. (б)~Найдите установившиеся $x_f$, $x_s$ и их сумму при постоянном $p=1$.
\end{exercise}

\begin{exercise}[\lvlB]
\label{ex:15:4}
\emph{Моделирование: второй раз быстрее.} По модели~\eqref{eq:tworate} с параметрами из \texttt{models/\allowbreak motor\_learning.py} найдите число движений до $x_f+x_s\ge0.8$ при $p=1$: (а)~у необученной машины; (б)~после $200$ движений с $p=1$ и обратного возмущения до нуля суммы, как на рис.~\ref{fig:tworate}. Может ли так ускориться модель с одним процессом?
\end{exercise}

\begin{exercise}[\lvlB]
\label{ex:15:5}
\emph{Регулятор с внутренней моделью.} Объект $\dd x/\dd t=ku$ виден с задержкой $d$; регулятор $u=-G\hat x$ предсказывает $\hat x(t)=x(t-d)+\hat k\int_{t-d}^{t}u\,\dd s$. (а)~При $\hat k=k=1$, $d=150\ms$ найдите $G$ для постоянной времени $20\ms$ и сравните с пределом~\eqref{eq:delaystable}. (б)~Устойчива ли схема при $\hat k=0.4k$?
\end{exercise}

\begin{exercise}[\lvlB]
\label{ex:15:6}
\emph{Адаптивный фильтр учит мышцу.} Объект --- подёргивание~\eqref{eq:twitch} единичной площади, $T_c=50\ms$; фильтры~\eqref{eq:adaptivefilter} --- RC-цепи с $\tau_j=12.5$, $25$, $50$, $100$, $200\ms$; вход --- новое нормальное число каждые $10\ms$. За сколько секунд правило~\eqref{eq:lms} при $\eta=50$ и $200$ снижает ошибку до $0.5\%$? Почему веса и через $300$~с далеки от наилучших?
\end{exercise}

\begin{exercise}[\lvlC]
\label{ex:15:7}
\emph{Два процесса как оптимальный фильтр.} Возмущение --- сумма процессов $p^{f,s}_{n+1}=A_{f,s}\,p^{f,s}_n+w^{f,s}_n$ с дисперсиями шумов $q_f$, $q_s$, ошибка наблюдается с шумом дисперсии $R$; фильтр Калмана даёт~\eqref{eq:tworate}. При каких $q_f/R$, $q_s/R$ из $A_f=0.92$, $A_s=0.996$ выходит $B_f=0.1$, $B_s=0.02$? Как изменятся $B_f$, $B_s$, если учетверить $q_s$?
\end{exercise}

\chapter{Второй выход: как мозг управляет химией тела}
\chaptermark{Химический выход}
\label{sec:chemoutput}

\section{Быстрый и медленный выходы}

В главе~\ref{sec:muscles} мы видели быстрый выход машины --- мышцы. Но у неё есть и второй, медленный. Мозг держит в нужных пределах температуру тела, частоту сердца, количество воды и сахара, готовит тело к бегу или ко сну. Для этого он не двигает суставы, а выпускает химические сигналы. Путей два. Первый --- нервы, идущие прямо к внутренним органам: к сердцу, сосудам, кишечнику, железам. Второй --- кровь: некоторые нейроны и железы выбрасывают вещество не в щель к соседу, а прямо в кровоток.

Кровь --- самая широковещательная шина из всех, что встречались в книге. Широковещательные каналы глав~\ref{sec:serocomputer}--\ref{sec:acet} рассылали сигнал на большие области мозга; кровь рассылает его каждой клетке тела. Кровь обходит тело примерно за минуту, так что сообщение доходит до всех за десятки секунд и держится минуты и часы, пока вещество не разрушится. Адреса у такой рассылки нет вовсе. Как и в утверждении~\ref{prop:receptor}, смысл сообщения выбирает получатель: откликаются только клетки, у которых есть рецептор к этому веществу.

\section{Газ и тормоз: два провода к органам}

Лучший пример --- сердце. В сердце есть собственный ритмоводитель, как у \nt{СЕРО}-нейрона главы~\ref{sec:serocomputer}: он сам задаёт ритм, и если отключить все нервы, сердце бьётся примерно сто раз в минуту. Мозг ритм не порождает, а только подстраивает, причём двумя проводами противоположного знака (рис.~\ref{fig:gasbrake}). По одному идёт \nt{НОРА}: это \emph{газ}, он ускоряет ритм. По другому --- \nt{АЦЕТ}: это \emph{тормоз}, он замедляет. Грубо,
\begin{empheq}[box=\keybox]{equation}
f \;\approx\; f_0 \;+\; a_S\,S \;-\; a_P\,P ,
\label{eq:gasbrake}
\end{empheq}
где $f_0\approx100$ ударов в минуту --- собственный ритм ритмоводителя, $S$ и $P$ --- активность газа и тормоза. В покое тормоз нажат, и сердце бьётся шестьдесят-семьдесят раз в минуту; при нагрузке тормоз отпускают и нажимают газ.

\begin{figure}[t]
\centering
\begin{tikzpicture}[>=Stealth,
  blk/.style={draw,very thick,rounded corners=2pt,align=center,font=\small},
  pace/.style={circle,draw,very thick,double,double distance=1.2pt,minimum size=1.8cm,inner sep=0pt,font=\scriptsize,fill=red!8,align=center},
  inh/.style={-{Bar[width=7pt]},thick,red!70!black}]
\node[blk,fill=blue!5,text width=1.6cm] (B) at (0,0) {команда\\мозга};
\node[pace] (H) at (6.2,0) {ритмо-\\водитель};
\draw[->,very thick,orange!85!black] (B.north east) to[bend left=20] node[above,font=\scriptsize,black,align=center] {\nt{НОРА}: газ, секунды} (H.north west);
\draw[inh,very thick,blue!60!black] (B.south east) to[bend right=20] node[below,font=\scriptsize,black,align=center] {\nt{АЦЕТ}: тормоз, доли секунды} (H.south west);
\draw[->,very thick] (H) -- (8.4,0) node[right,font=\small] {частота $f$};
\node[blk,fill=orange!10,text width=1.6cm] (A) at (0,-2.6) {\nt{АДРЕ}\\в кровь};
\draw[->,thick,orange!85!black] (B) -- (A);
\foreach \y/\lab in {-2.0/{сердце},-2.6/{сосуды},-3.2/{печень, мышцы}}{
  \draw[->,thick,densely dashed,orange!85!black] (A.east) -- (4.3,\y) node[right,font=\scriptsize,black] {\lab};}
\end{tikzpicture}
\caption{Два провода противоположного знака к ритмоводителю сердца: газ \nt{НОРА} медленный, тормоз \nt{АЦЕТ} быстрый. Внизу --- тот же газ по широковещательной шине: \nt{АДРЕ}-клетки выбрасывают адреналин в кровь, и сигнал получают все органы с подходящим рецептором (штриховые стрелки).}
\label{fig:gasbrake}
\end{figure}

Это двухпроводный код главы~\ref{sec:glutcomputer} и пара мышц главы~\ref{sec:muscles}, только теперь провода управляют не суставом, а темпом ритмоводителя. И у проводов разная скорость. Оба работают как фильтры нижних частот, но тормоз пропускает перемены в несколько раз быстрее газа~\cite{Berger1989}: \nt{АЦЕТ} открывает ионные каналы прямо, а \nt{НОРА} действует через цепочку реакций внутри клетки. Инженер узнает здесь приём с двумя приводами разной скорости: быстрый делает точную мгновенную подстройку, медленный --- крупные долгие перемены.

Здесь же находит роль \nt{АДРЕ}, про который в таблице~\ref{tab:types} было написано, что его роль в мозге неясна. Вне мозга она ясна. Часть клеток газового пути не посылает провод к органу, а выбрасывает адреналин прямо в кровь. Это тот же газ, но по широковещательной шине: за десятки секунд он доходит до сердца, сосудов, печени и мышц сразу и держится минуты. Адресный газ \nt{НОРА} подстраивает один орган; широковещательный газ \nt{АДРЕ} переводит в режим бега всё тело.

\section{Каскад с обратной связью}

Многие гормоны выпускаются не одной железой, а каскадом. Небольшая группа нейронов в основании мозга выбрасывает в кровь первый сигнал; он заставляет железу-посредника выпустить второй; тот --- конечную железу выпустить гормон, который и действует на тело. А конечный гормон возвращается к первым ступеням и тормозит их. Получается усилитель из трёх каскадов, охваченный отрицательной обратной связью, --- регулятор уровня, как у \nt{СЕРО}-системы в формуле~\eqref{eq:feedback}.

Опишем каждую ступень RC-цепью с постоянной времени $\tau$ и усилением $g$, а обратную связь --- коэффициентом $k$:
\begin{equation}
\tau\,\frac{\dd x_1}{\dd t} = -x_1 + g\bigl[u - k\,x_3\bigr]_+ ,\qquad
\tau\,\frac{\dd x_2}{\dd t} = -x_2 + g\,x_1 ,\qquad
\tau\,\frac{\dd x_3}{\dd t} = -x_3 + g\,x_2 ,
\label{eq:cascade}
\end{equation}
где $u$ --- команда мозга, $x_3$ --- уровень конечного гормона. Усиление петли $K=g^3k$. В равновесии $x_3=g^3u/(1+K)$, и, как у любого регулятора с обратной связью, большая петля делает уровень нечувствительным к разбросу усилений ступеней. Но каждая ступень сдвигает фазу: на частоте $\omega=\sqrt3/\tau$ три одинаковые ступени вместе дают сдвиг в полоборота и ослабление в восемь раз. Отсюда классический результат теории регулирования:
\begin{empheq}[box=\keybox]{equation}
K \;<\; 8 ,
\label{eq:cascadestable}
\end{empheq}
иначе регулятор начинает колебаться, с периодом около $2\pi\tau/\sqrt3\approx3.6\,\tau$.

\begin{keyblock}
\begin{proposition}[Точность против устойчивости]
\label{prop:cascade}
У трёхступенчатого каскада с пропорциональной обратной связью ошибка уровня равна $1/(1+K)$, а устойчив он только при $K<8$. Поэтому точнее, чем примерно на десять процентов, такой регулятор уровень не держит; попытка поднять усиление превращает его в генератор колебаний.
\end{proposition}
\end{keyblock}

Мы проверили это на модели~\eqref{eq:cascade} (рис.~\ref{fig:cascade}). При $K=4$ уровень после включения немного колеблется и устанавливается. При $K=7$ он тоже устанавливается, но медленно. При $K=12$ колебания не затухают и не растут без предела: ступень не может выдать отрицательный сигнал, и каскад выходит на ровные пульсации между $0.83$ и $1.24$ от среднего уровня с периодом $3.7$--$3.9\,\tau$.

\begin{figure}[t]
\centering
\begin{tikzpicture}[>=Stealth,x=0.3cm,y=1.2cm]
\draw[->,gray!70!black] (0,0) -- (41.5,0) node[right,font=\small,black] {$t/\tau$};
\draw[->,gray!70!black] (0,0) -- (0,2.8) node[left,font=\small,black] {$x_3/x_3^*$};
\foreach \t in {0,10,20,30,40}{\draw[gray!60!black] (\t,-0.05) -- (\t,0.05); \node[below,font=\scriptsize] at (\t,0) {\t};}
\foreach \f in {1,2}{\draw[gray!60!black] (-0.3,\f) -- (0.3,\f); \node[left,font=\scriptsize] at (0,\f) {\f};}
\draw[densely dashed,gray!60!black] (0,1) -- (40,1);
\draw[thick,blue!60!black] plot file {figures/cascade_K4.dat};
\draw[thick,red!70!black] plot file {figures/cascade_K12.dat};
\node[font=\scriptsize,blue!60!black,anchor=west] at (16,0.55) {$K=4$: устанавливается};
\node[font=\scriptsize,red!70!black,anchor=west] at (16,1.75) {$K=12$: пульсирует};
\end{tikzpicture}
\caption{Каскад из трёх ступеней с обратной связью по~\eqref{eq:cascade}, уровень конечного гормона в долях равновесного $x_3^*$. При усилении петли ниже восьми уровень устанавливается; выше --- каскад сам порождает ровные пульсации.}
\label{fig:cascade}
\end{figure}

Настоящие гормоны каскадов действительно выходят пульсами: уровень гормона стресса, например, колеблется с периодом около часа. По одной из гипотез, эти пульсации порождает сама обратная связь между ступенями с запаздыванием, и отдельный генератор ритма для них не нужен~\cite{Walker2010}. Это та же причина колебаний, что и у петли \nt{ГЛУТ}--\nt{ГАМК} в утверждении~\ref{prop:ei} и у петли растяжения мышцы в условии~\eqref{eq:delaystable}: отрицательная обратная связь, которая запаздывает.

\section{Импульсный код гормонов}

Пульсы бывают не только побочным эффектом, но и кодом. Нейроны, управляющие размножением, выпускают свой сигнал в кровь короткими порциями примерно раз в час. Если железе-получателю подать тот же сигнал не порциями, а непрерывно, она не работает в полную силу, как при пульсах, а замолкает; если снова подавать порциями раз в час, работа восстанавливается~\cite{Belchetz1978}. Значит, получатель читает не уровень, а \emph{частоту} порций.

Для инженера здесь нет загадки. Рецептор, который при долгом действии вещества устаёт и перестаёт отвечать, --- фильтр верхних частот, как истощающийся синапс~\eqref{eq:depression} из главы~\ref{sec:code}. Непрерывный сигнал он глушит, а перемены пропускает. Такому получателю можно сообщить что-то только импульсами, и информация переходит из уровня в частоту и форму порций. Разные частоты порций к тому же по-разному действуют на разные клетки одной и той же железы, так что по одной химической линии можно передать больше одной величины --- как по одному проводу \nt{ДОФА} в главе~\ref{sec:dofaalt} шли быстрая и медленная составляющие.

\section{Суточный ритм: медленный генератор режимов}

Самый медленный ритм тела --- суточный. Его порождает отрицательная обратная связь с запаздыванием внутри клеток: гены производят белки, которые с задержкой в несколько часов подавляют собственное производство~\cite{TakahashiJS2017}. Опять запаздывающая отрицательная обратная связь --- четвёртый раз в этой книге, --- и опять ритм, на этот раз с периодом около суток. Главный генератор --- небольшая группа из десятков тысяч нейронов в основании мозга, ритм которых синхронизирует остальные клетки тела.

Собственный период этого генератора у человека не ровно сутки, а в среднем $24.18$ часа~\cite{Czeisler1999}. Каждый день его подводит свет, приходящий от датчиков глаза (глава~\ref{sec:sensors}). Для электронщика это \emph{фазовая автоподстройка частоты}: генератор со своей частотой $\omega_0$ подстраивается к внешнему сигналу с частотой $\omega$, и разность фаз $\varphi$ подчиняется уравнению
\begin{empheq}[box=\keybox]{equation}
\frac{\dd\varphi}{\dd t} \;=\; (\omega-\omega_0) \;-\; K_\varphi\,\sin\varphi .
\label{eq:pll}
\end{empheq}
Подстройка удерживается, если $|\omega-\omega_0|<K_\varphi$. Генератору с периодом $24.18$ часа нужно каждый день сдвигаться примерно на одиннадцать минут; утреннего света для такого сдвига обычно хватает. В полярную ночь или в пещере без света подстройки нет, и ритм уходит на собственный период.

Суточный генератор --- не тактовый генератор компьютера. Он не отсчитывает шаги вычислений и не синхронизирует импульсы нейронов. Он переключает \emph{режимы}: сон и бодрствование, уровни гормонов, температуру тела, готовность к еде. Это медленный широковещательный регистр того же рода, что и регистры глав~\ref{sec:serocomputer}--\ref{sec:acet}, только его значение меняется по расписанию, подведённому к солнцу.

\begin{keyblock}
\begin{proposition}[Химический выход]
\label{prop:chemoutput}
Медленный выход машины построен из тех же деталей, что и быстрый. Органы получают два провода противоположного знака и разной скорости --- газ \nt{НОРА} и тормоз \nt{АЦЕТ}, --- а весь организм сразу --- широковещательные сигналы по крови, в том числе адреналин \nt{АДРЕ}. Уровни гормонов держат каскады с отрицательной обратной связью, точность которых ограничена устойчивостью; часть гормонов кодируется частотой порций; суточный ритм --- медленный генератор с фазовой подстройкой по свету. Устройство путей и опыты с порциями и с периодом измерены; объяснение пульсаций каскада самой обратной связью --- гипотеза.
\end{proposition}
\end{keyblock}

\section{Итог}

\begin{table}[t]
\centering
\footnotesize
\setlength{\tabcolsep}{4pt}
\renewcommand{\arraystretch}{1.15}
\begin{tabular}{>{\raggedright}p{3.1cm}>{\raggedright}p{4.8cm}>{\raggedright\arraybackslash}p{4.9cm}}
\toprule
Что делает химический выход & Как & Что известно \\
\midrule
подстраивает органы & газ \nt{НОРА} и тормоз \nt{АЦЕТ}~\eqref{eq:gasbrake} & измерено; тормоз быстрее газа \\
переводит всё тело в режим бега & адреналин \nt{АДРЕ} в кровь & измерено \\
держит уровни гормонов & каскад с обратной связью, $K<8$~\eqref{eq:cascadestable} & каскады и обратная связь измерены; пульсации от самой петли --- гипотеза \\
передаёт частотой порций & уставший рецептор как фильтр верхних частот & зависимость от порций измерена \\
переключает режимы за сутки & генератор с фазовой подстройкой по свету~\eqref{eq:pll} & период и подстройка светом измерены \\
\bottomrule
\end{tabular}
\caption{Как машина управляет химией тела.}
\label{tab:chemoutput}
\end{table}

У машины два выхода (таблица~\ref{tab:chemoutput}). Быстрый --- мышцы: миллисекунды, адресно, в каждый сустав. Медленный --- химия тела: секунды, минуты и сутки, по двум проводам к органам и по крови ко всем сразу. Детали у обоих выходов одни и те же --- два провода противоположного знака, ритмоводители, обратная связь и её запаздывание, --- и те же, из которых собрана машина внутри.

\section{Задачи}

\begin{exercise}[\lvlA]
\label{ex:16:1}
\emph{Кровь как фильтр нижних частот.} Гормон выводится из крови с периодом полураспада $t_{1/2}$ и выпускается короткими порциями каждые $T=60$~мин.
(а)~Покажите, что концентрация в крови --- выход RC-цепи с $\tau=t_{1/2}/\ln2$.
(б)~Докажите, что в установившемся режиме отношение максимума концентрации к минимуму равно $2^{T/t_{1/2}}$. Найдите его при $t_{1/2}=10$~мин.
(в)~Во сколько раз цепь ослабляет основную гармонику порций при $t_{1/2}=10$~мин?
(г)~При каком $t_{1/2}$ отношение максимума к минимуму падает до $2$? При каком основная гармоника ослабляется в $10$ раз?
Что это говорит о том, какие гормоны могут передавать сообщение частотой порций?
\end{exercise}

\begin{exercise}[\lvlA]
\label{ex:16:2}
\emph{Фазовая автоподстройка суточного генератора.} Рассмотрите~\eqref{eq:pll}.
(а)~Покажите, что равновесия существуют при $|\omega-\omega_0|<K_\varphi$, найдите устойчивое и скорость приближения к нему.
(б)~Собственный период $24.18$~ч. Пусть свет может сдвигать фазу не больше чем на час в сутки, $K_\varphi=2\pi/24$~рад/сут. Найдите $\omega-\omega_0$ в радианах в сутки и в минутах сдвига в сутки, установившуюся разность фаз $\varphi^*$ в минутах и время релаксации.
(в)~На сколько уйдёт ритм без света за $30$ суток?
(г)~Внешний сигнал скачком сдвинули на $\pm6$~ч. Численно проинтегрировав~\eqref{eq:pll}, найдите, за сколько суток рассогласование станет меньше часа.
\end{exercise}

\begin{exercise}[\lvlB]
\label{ex:16:3}
\emph{Точность против устойчивости.} Рассмотрите линеаризованный каскад~\eqref{eq:cascade}.
(а)~Выведите~\eqref{eq:cascadestable} по критерию Рауса--Гурвица.
(б)~Выведите то же по критерию Найквиста и найдите период колебаний на границе.
(в)~Обобщите на $n$ одинаковых ступеней: найдите $K_c(n)$, период на границе и наименьшую достижимую ошибку $1/(1+K_c)$ для $n=3,\dots,6$ и при $n\to\infty$.
(г)~Замените пропорциональную обратную связь интегральной: вход первой ступени $m$, $\dd m/\dd t=k_I\,(x_{\text{цель}}-x_3)$, ступени с единичным усилением. Покажите, что ошибка в равновесии нулевая, и найдите условие устойчивости и период колебаний на границе. Какова цена?
\end{exercise}

\begin{exercise}[\lvlB]
\label{ex:16:4}
\emph{Проектирование: газ и тормоз.} Пусть в~\eqref{eq:gasbrake} вклады газа $\Delta_S=a_SS$ и тормоза $\Delta_P=a_PP$ (в ударах в минуту) --- выходы RC-цепей с постоянными $\tau_S=2$~с и $\tau_P=0.4$~с, команды на которые ограничены: $0\le\Delta_S^{\text{ком}}\le80$, $0\le\Delta_P^{\text{ком}}\le60$. Собственный ритм $f_0=100$, в покое $\Delta_P=35$, $\Delta_S=0$, $f=65$.
(а)~Нужно как можно быстрее выйти на $f=120$ и удерживать её. Докажите, что до первого достижения цели оптимальны крайние команды, и найдите это время.
(б)~Постройте закон удержания после достижения цели и проверьте, что команды остаются в допустимых пределах.
(в)~Сравните со стратегией «только газ» при нажатом тормозе и со стратегией «только тормоз».
(г)~Объясните с точки зрения управления, зачем в покое держать нажатым быстрый тормоз, а не отпущенным.
\end{exercise}

\begin{exercise}[\lvlB]
\label{ex:16:5}
\emph{Моделирование: каскад с запаздыванием.} Скопируйте функцию \texttt{cascade} из \texttt{models/\allowbreak chem\_models.py} (сам файл не запускайте: он перезаписывает данные рисунков) и добавьте в обратную связь чистую задержку $d$: первая ступень видит $x_3(t-d)$.
(а)~Выведите критическое усиление: $3\arctan(\omega\tau)+\omega d=\pi$, $K_c=(1+\omega^2\tau^2)^{3/2}$. Найдите $K_c$ и период на границе для $d/\tau=0$, $0.5$, $1$, $2$.
(б)~Проверьте моделированием: при $0.9K_c$ колебания затухают, при $1.1K_c$ устанавливаются. Как надёжно отличить медленное затухание от установившихся колебаний?
(в)~Без задержки измерьте период и размах установившихся пульсаций при $K=8.5$, $10$, $12$, $16$, $24$, $40$.
(г)~Какое $\tau$ ступени нужно, чтобы пульсации имели период около часа? Как меняется допустимая точность регулятора при $d=\tau$?
\end{exercise}

\begin{exercise}[\lvlB]
\label{ex:16:6}
\emph{Получатель, который читает порции.} Рецептор устаёт: $y=[c-a]_+$, $\tau_a\,\dd a/\dd t=c-a$.
(а)~Найдите передаточную функцию от $c$ к $c-a$ и покажите, что при постоянном $c$ отклик стремится к нулю.
(б)~Пусть $c$ --- прямоугольные порции амплитуды $c_p$ и длительности $T_p$ с периодом $T$, средняя доза $\bar c=c_pT_p/T$ фиксирована. Докажите, что средний отклик равен $\frac1T\int_{\text{пауза}}a\,\dd t$, и найдите его в замкнутом виде.
(в)~При $\tau_a=30$~мин, $T_p=6$~мин найдите средний отклик для $T=15$, $60$, $120$~мин и предел при $T\to\infty$.
(г)~Пусть отклик насыщается, $y\le y_{\max}$. Найдите численно период порций, дающий наибольший средний отклик, при $y_{\max}=2\bar c$ и $5\bar c$. Как с помощью таких получателей передать по одной линии больше одной величины?
\end{exercise}

\begin{exercise}[\lvlC]
\label{ex:16:7}
\emph{Генератор или обратная связь?} В тексте пульсации гормонов объясняются гипотезой: их порождает сама запаздывающая обратная связь каскада, а не отдельный ритмоводитель. Предложите опыты, которые различат две гипотезы, и получите количественные предсказания на модели~\eqref{eq:cascade}.
(а)~Как должны меняться период и размах пульсаций при изменении усиления петли и при изменении постоянных времени ступеней в каждой гипотезе?
(б)~Что произойдёт, если разомкнуть петлю, подав на первую ступень постоянное значение конечного гормона?
(в)~Как ответит каждая система на короткую добавку конечного гормона в разные фазы цикла?
Оцените, можно ли различить гипотезы по сдвигам периода, если $K$ удаётся изменить лишь в полтора раза, а период измеряется с точностью $5\%$.
\end{exercise}

\chapter{Отладка машины}
\chaptermark{Отладка машины}
\label{sec:debugging}

\section{Как отлаживают машину без схемы}

Инженер, которому дали работающую плату без схемы, отлаживает её четырьмя приёмами. Он ставит \emph{щуп} и смотрит, что идёт по проводу. Он \emph{подаёт тестовый сигнал} и смотрит, что изменилось. Он \emph{отключает} кусок схемы и смотрит, что перестало работать. И он \emph{читает регистры управления}, чтобы узнать, в каком режиме работает плата. Вся эта книга --- попытка прочитать схему мозга, и в этой главе посмотрим, какими из четырёх приёмов инженеры уже умеют пользоваться и что им говорят предыдущие главы о том, чего ждать.

Самый заметный сегодня пример такой отладки --- интерфейсы мозг--компьютер: устройства, которые читают импульсы нейронов и превращают их в команды для внешних машин, или, наоборот, подают в мозг сигналы извне. Им и посвящена большая часть главы, и в первую очередь тому, что делает компания Neuralink.

\section{Щуп: что слышит электрод}

Щуп мозга --- тонкий электрод, вставленный в ткань. Он слышит импульсы нейронов, лежащих в радиусе порядка сотни микрометров, --- обычно от одного до нескольких. Импульс на электроде --- всплеск в десятки--сотни микровольт длительностью около миллисекунды, и по его форме можно отличить соседние нейроны друг от друга. Лучше всего слышны крупные \nt{ГЛУТ}-нейроны: их в коре большинство (таблица~\ref{tab:types}), и их токи сильнее.

Сразу видна главная трудность. Хороший щуп сегодня --- это сотни или тысячи электродов, а нейронов в мозге $8.6\cdot10^{10}$. Мы подслушиваем примерно одну стомиллионную машины. Почему же из такой ничтожной выборки вообще что-то можно понять? Ответ дала глава~\ref{sec:code}. Соседние нейроны шумят согласованно, и усреднение по группе упирается в предел (утверждение~\ref{prop:pooling}): сотня нейронов несёт почти столько же, сколько тысяча. К тому же задача, которую решает двигательная кора, маломерна: у руки немного суставов (глава~\ref{sec:muscles}), и активность сотен нейронов, управляющих ею, укладывается в пространство всего из десятка-другого общих переменных~\cite{Gallego2017}. Щуп на сотню нейронов слышит почти все эти переменные. Если бы мозг хранил независимое сообщение в каждом нейроне, такое подслушивание было бы бесполезно.

\section{Поток данных: сжатие прямо на щупе}

Щуп выдаёт огромный поток. Чтобы увидеть форму импульса, каждый электрод нужно оцифровывать примерно $20$ тысяч раз в секунду с точностью около десяти бит. Тысяча электродов дают
\begin{empheq}[box=\keybox]{equation}
\begin{aligned}
R_{\text{сырой}} \;&=\; N\,f_s\,b \;\approx\; 10^3\cdot2\cdot10^4\cdot10 \;\approx\; 2\cdot10^8\ \text{бит/с},\\
R_{\text{имп}} \;&\approx\; N\,r\,\log_2\frac{e}{r\,\Delta t} \;\approx\; 8\cdot10^4\ \text{бит/с}.
\end{aligned}
\label{eq:probedata}
\end{empheq}
Вторая оценка --- ёмкость потока импульсов~\eqref{eq:spikeentropy} при $r=10$ импульсов в секунду и точности $\Delta t=1\ms$: всё, что в нём на самом деле есть, занимает в две с половиной тысячи раз меньше. Для вживлённого устройства это решающая разница: сырой поток нельзя передать по радио через кожу с тем питанием, которое можно позволить внутри головы, не нагревая ткань. Поэтому вживлённому устройству нужно выделять импульсы прямо на кристалле рядом с электродами и передавать наружу только события --- номер канала и момент импульса. В первом описании системы Neuralink до этого ещё не дошли: кристаллы усиливают и оцифровывают сигналы, весь сырой поток уходит по кабелю, а импульсы выделяет программа снаружи; сжатие на кристалле и связь без проводов названы там планом~\cite{Musk2019}.

Это знакомый приём. Глаз сжимает сигнал в сто раз, прежде чем отправить его по длинному кабелю (глава~\ref{sec:sensors}); событийная камера передаёт не кадры, а изменения. Щуп, чтобы уложиться в провод, вынужден делать то же, что делает сам мозг: говорить импульсами и только о новостях.

\section{Что делает Neuralink}

Устройство Neuralink --- щуп, собранный под эти ограничения~\cite{Musk2019}. Вместо жёсткой решётки иголок --- много гибких нитей тоньше волоса с электродами вдоль каждой: в первом описании системы --- до $3072$ электродов на $96$ нитях, по $32$ на нить, а в устройстве, вживлённом людям, по сообщениям компании, около тысячи каналов на $64$ нитях. Гибкие нити меньше повреждают ткань и лучше переносят то, что мозг внутри черепа всё время слегка сдвигается. Нити вставляет робот, обходя сосуды. В первом описании, как сказано выше, сырой поток~\eqref{eq:probedata} шёл наружу по кабелю. Устройство для людей, по сообщениям компании, устроено иначе: корпус полностью закрыт кожей, импульсы выделяются внутри, данные уходят по радио, питание приходит без проводов.

Задача первого устройства --- чтение. Нити ставят в ту часть коры, которая планирует движения руки, и декодер превращает импульсы в движение курсора на экране. Человек с параличом рук управляет компьютером, представляя движение. Сама идея не нова: интерфейсы на жёстких решётках из сотни электродов давно позволяли управлять курсором~\cite{Hochberg2006}, писать текст, представляя движения руки при письме~\cite{Willett2021}, и даже переводить попытки говорить в текст на экране со скоростью около шестидесяти слов в минуту~\cite{Willett2023}. Новое у Neuralink --- инженерия: число каналов, беспроводность, закрытый корпус и роботизированная вставка, то есть попытка сделать щуп, который можно носить годами вне лаборатории. Насколько нам известно, результаты испытаний на людях публиковались в основном самой компанией, а не в рецензируемых статьях; компания сообщала, в частности, что часть нитей после вставки сместилась и число работающих каналов уменьшилось. Для отладки это обычная неприятность: щуп, который двигается относительно платы, слышит уже другие провода.

Второе направление компании --- запись: устройство для зрения, которое подаёт сигналы в зрительную кору человека, потерявшего зрение. О нём ниже, в разделе о записи.

\section{Чтение: декодер как получатель}

Декодер интерфейса --- это получатель в смысле утверждения~\ref{prop:reader}: он сам решает, какую часть сообщения прочитать. Практически все интерфейсы читают \emph{частоты групп}: считают импульсы каждого канала в окнах по несколько десятков миллисекунд и складывают их с весами, подобранными так, чтобы получалось задуманное движение. Это частотный код группы из раздела~\ref{sec:rate}, и он выбран не случайно: при нескольких импульсах на окно частота одного нейрона не определена, а сумма по сотне уже работает.

Сколько информации удаётся вынуть? Письмо «мысленной рукой» шло со скоростью около девяноста знаков в минуту~\cite{Willett2021}, то есть полтора знака в секунду: даже при пяти битах на знак это меньше восьми бит в секунду. Управление курсором, по сообщениям Neuralink, даёт порядка десяти бит в секунду. А верхняя граница~\eqref{eq:spikeentropy} для \emph{одного} нейрона --- десятки бит в секунду, для сотни --- тысячи. Интерфейс вынимает из подслушанных нейронов малую долю того, что они могли бы нести. Узкое место --- не провод, а то, сколько независимых переменных несёт группа (утверждение~\ref{prop:pooling}) и насколько хорошо декодер понимает код, который, как мы видели в главе~\ref{sec:code}, до конца не расшифрован.

Есть и вторая особенность, знакомая по главе~\ref{sec:motorlearning}. Декодер и мозг учатся друг у друга. Мозг видит, куда поехал курсор, получает ошибку и подстраивает свои команды --- то же обучение движениям по проводу ошибки. А инженеры время от времени заново подбирают веса декодера, потому что нейроны под нитями меняются и связи в сети переучиваются. Получаются два ученика, которые подстраиваются друг под друга; чтобы такая пара не ушла вразнос, скорости их обучения удобно развести: пусть один учится быстро, а другой медленно.

\begin{keyblock}
\begin{proposition}[Почему щуп на тысячу нейронов работает]
\label{prop:probe}
Интерфейс, подслушивающий ничтожную долю нейронов, может управлять движением потому, что активность группы маломерна и её шумы согласованы: несколько сотен нейронов несут почти все общие переменные. По той же причине интерфейс вынимает лишь несколько бит в секунду --- столько, сколько независимых переменных в задаче, а не столько, сколько вмещает поток импульсов. Работа интерфейсов измерена; насколько лучше станет декодер, когда код будет понят, --- открытый вопрос.
\end{proposition}
\end{keyblock}

\section{Запись: труднее, чем чтение}

Подать сигнал в машину труднее, чем прочитать. Электрод, через который пропускают ток, возбуждает не один нейрон, а все нейроны и провода вокруг себя, без разбора типа и знака. Записать им код, в котором важно, \emph{какой} нейрон и \emph{когда} сработал (глава~\ref{sec:code}), нельзя: можно лишь грубо задать, \emph{где} и \emph{насколько сильно}.

Для зрения этого отчасти хватает. Ток через электрод в зрительной коре человек видит как светящуюся точку в определённом месте поля зрения; когда точки зажигали одновременно, до шестнадцати сразу, человек, потерявший зрение, узнавал некоторые буквы и различал границы предметов~\cite{Fernandez2021}. Это код местом, и его записать можно. Но вспомните главу~\ref{sec:sensors}: глаз посылает в мозг не пиксели, а сжатое описание --- края, перемены, посветление и потемнение, по отдельным проводам. Запись «точек света» в кору --- это запись пикселей в машину, которая ждёт на входе совсем другой код. Поэтому искусственное зрение пока грубее, чем можно было бы ждать по числу электродов. Есть и испытательные сигналы, которым электрод не нужен: зрительные иллюзии подают через глаз необычный вход и выводят машину из обычного режима, а каждая ошибка указывает на свой механизм (раздел~\ref{sec:illusions}).

\section{Чего щуп не видит}

Электрод слышит адресную шину --- импульсы \nt{ГЛУТ}-нейронов и хуже, мелких \nt{ГАМК}. Но в главах~\ref{sec:serocomputer}--\ref{sec:acet} мы видели, что режим работы всей машины задают широковещательные регистры: концентрации \nt{СЕРО}, \nt{ДОФА}, \nt{НОРА}, \nt{АЦЕТ}. Их электрод не слышит: нейронов-источников ничтожно мало, и их сигнал --- не импульс рядом с щупом, а концентрация в объёме. Отладчик, который видит шину данных, но не видит регистров управления, всегда будет удивляться: один и тот же вход будет давать разные ответы, потому что где-то изменился $g$, $a$ или $\tau_\gamma$. Концентрации можно мерить химическими датчиками прямо в ткани~\cite{Bunin1998}, но в интерфейсах такие датчики пока не используются. Совсем вне поля зрения щупа остаётся и второй, медленный выход машины --- гормоны в крови (глава~\ref{sec:chemoutput}): их меряют в пробах крови, а не электродом.

Не видит щуп и весов --- а именно в них хранится почти вся память и всё выученное. Их можно лишь угадывать по тому, как меняются ответы. И не видит он боли так, как её чувствует человек: в главе~\ref{sec:pain} мы видели, что сила сигнала тревоги задаётся самой машиной --- воротами в спинном мозге и регуляторами сверху, так что по датчику нельзя прочитать, насколько больно.

\section{Итог: отладка и открытые вопросы}

\begin{table}[t]
\centering
\footnotesize
\setlength{\tabcolsep}{4pt}
\renewcommand{\arraystretch}{1.15}
\begin{tabular}{>{\raggedright}p{2.2cm}>{\raggedright}p{3.6cm}>{\raggedright}p{3.4cm}>{\raggedright\arraybackslash}p{4.0cm}}
\toprule
Приём отладки & Что делает интерфейс & Что объясняет книга & Что известно \\
\midrule
щуп & слышит сотни--тысячи нейронов & маломерность и согласованный шум (гл.~\ref{sec:code}) & измерено \\
сжатие на щупе & передаёт события вместо сырого сигнала & ёмкость потока импульсов~\eqref{eq:spikeentropy} & инженерно решено \\
чтение & частоты групп $\to$ движение, текст, речь & декодер --- получатель, частотный код группы & работает; несколько бит в секунду \\
совместное обучение & мозг и декодер подстраиваются друг под друга & провод ошибки (гл.~\ref{sec:motorlearning}) & наблюдается; правила подстройки --- предмет исследований \\
запись & ток зажигает точки в поле зрения & код местом записать можно, временной --- нет (гл.~\ref{sec:sensors}) & точки измерены; полезное зрение --- пока нет \\
чтение регистров & почти не делается & широковещательные каналы (гл.~\ref{sec:serocomputer}--\ref{sec:acet}) & химические датчики есть в опытах \\
\bottomrule
\end{tabular}
\caption{Отладка машины: что умеют интерфейсы мозг--компьютер и что о них говорят главы книги.}
\label{tab:debugging}
\end{table}

Таблица~\ref{tab:debugging} сводит главу. Интерфейсы мозг--компьютер уже умеют ставить щуп на адресную шину и читать с неё несколько бит в секунду, и то, что это вообще работает, объясняется маломерностью и согласованным шумом из главы~\ref{sec:code}. Записывать в машину они умеют лишь грубо, потому что её входной код сжат и адресован отдельным проводам. А регистры управления и веса --- то, что делает машину обучаемой и настраиваемой, --- щупу пока почти не видны.

Каждый открытый вопрос главы~\ref{sec:synthesis} --- это и задача для отладчика. Какой код читать, решится, когда щупы станут плотнее и декодеры научатся пользоваться временем импульсов, а не только частотами. Как учится многослойная сеть, можно проверить, поставив щупы сразу в несколько слоёв и глядя, как меняются их ответы при обучении. Что передают широковещательные каналы, станет видно, когда к электрическим щупам добавятся химические. Эта книга --- попытка прочитать схему машины по её поведению и по законам, которые знает любой инженер. Отладчики, которые строятся сейчас, дадут возможность проверить эту схему щупом.

\section{Задачи}

\begin{exercise}[\lvlA]
\label{ex:17:1}
\emph{Бюджет радиоканала.} (а)~Пересчитайте~\eqref{eq:probedata} и отношение двух потоков для $N=1000$ и $N=3072$ электродов.
(б)~Пусть передача через кожу стоит $1$~нДж на бит, а на радиопередатчик можно тратить не больше $10$~мВт. Какая мощность нужна для сырого потока и для потока импульсов? Какой энергии на бит потребовал бы сырой поток?
(в)~Импульсы передаются кадрами по $1\ms$: заголовок $16$~бит и $10$-битный номер канала на каждый импульс. Найдите средний поток при $N=1000$, $r=10$~имп/с и сравните его с ёмкостью~\eqref{eq:probedata}.
\end{exercise}

\begin{exercise}[\lvlA]
\label{ex:17:2}
\emph{Сколько бит в сотне нейронов.} Декодер складывает импульсы $N$ нейронов за окно $T=50\ms$, каждый работает с частотой около $r=20$~имп/с, шум пуассоновский.
(а)~По~\eqref{eq:rateerror} и~\eqref{eq:correlated} найдите относительную погрешность суммы при $N=1$, $10$, $100$, $1000$ для $c=0$ и $c=0.1$.
(б)~Пусть управляемая переменная меняет частоту группы со среднеквадратичной амплитудой $30\%$, а шум нормальный. Оцените информацию на окно $\tfrac12\log_2(1+\text{SNR})$ и верхнюю границу скорости, считая окна независимыми.
(в)~Сравните с числами из текста. Какое предположение о коррелированном шуме делает эту оценку пессимистичной?
\end{exercise}

\begin{exercise}[\lvlB]
\label{ex:17:3}
\emph{Скорость интерфейса-клавиатуры.} Интерфейс выбирает один из $N$ знаков; с вероятностью $P$ выбран задуманный, ошибки равновероятны среди остальных $N-1$.
(а)~Докажите, что взаимная информация на выбор при равновероятных знаках равна
$B=\log_2N+P\log_2P+(1-P)\log_2\dfrac{1-P}{N-1}$.
(б)~Докажите, что это ёмкость такого канала.
(в)~Найдите $B$ и скорость в битах в секунду при $N=32$, полутора знаках в секунду и $P=0.99$, $0.94$, $0.9$, $0.8$. Сверьте с оценкой из текста.
(г)~При каком $P$ скорость нулевая? Сколько знаков в секунду нужно при $P=0.94$, чтобы получить $10$~бит/с?
\end{exercise}

\begin{exercise}[\lvlB]
\label{ex:17:4}
\emph{Моделирование: декодер группы.} $N$ нейронов настроены на двумерную скорость $\mathbf v$: частота $\lambda_i=\bigl[b+m\,\mathbf v\cdot\mathbf p_i\bigr]_+$ с $b=20$~имп/с, $m=15$~имп/с, предпочтительные направления $\mathbf p_i$ равномерно по кругу. Каждая компонента $\mathbf v$ --- случайный процесс Орнштейна--Уленбека с постоянной $0.5$~с и стандартным отклонением $0.5$. Импульсы пуассоновские, окна $T=50\ms$, $2000$ окон. Чтобы смоделировать шум, похожий на сигнал, все нейроны получают вместо $\mathbf v$ одно и то же $\mathbf v+\boldsymbol\xi$, где $\boldsymbol\xi$ --- независимый в каждом окне нормальный шум со стандартным отклонением $\sigma_\xi$ по каждой компоненте.
(а)~Декодируйте вектором группы $\hat{\mathbf v}=\frac{2}{NmT}\sum_i(n_i-bT)\,\mathbf p_i$. Выведите дисперсию ошибки: $2b/(Nm^2T)+\sigma_\xi^2$ на компоненту.
(б)~Постройте среднеквадратичную ошибку от $N=10$, $30$, $100$, $300$, $1000$ для $\sigma_\xi=0$ и $0.2$ и сравните с формулой.
(в)~При каком $N$ два вклада в ошибку равны? Сколько бит на компоненту за окно даёт бесконечная группа при $\sigma_\xi=0.2$? Свяжите результат с утверждениями~\ref{prop:pooling} и~\ref{prop:probe}.
\end{exercise}

\begin{exercise}[\lvlB]
\label{ex:17:5}
\emph{Два ученика.} Мозг выдаёт команду $m=b\,v^*$ на задуманную скорость $v^*$ ($\langle v^{*2}\rangle=1$), декодер выдаёт $v=d\,m$. Ошибка $e=v-v^*$. Оба учатся градиентным спуском по $\tfrac12e^2$ со скоростями $\eta_b$ и $\eta_d$.
(а)~Запишите уравнения в непрерывном времени. Покажите, что величина $\eta_d b^2-\eta_b d^2$ сохраняется, и найдите, куда придёт пара, выйдя из $b=1.5$, $d=1$, при $\eta_b=1$, $\eta_d=0.1$ и наоборот.
(б)~Для одновременных обновлений после каждой попытки (оба ученика используют прежние значения $b$ и $d$) линеаризуйте ошибку $\varepsilon=db-1$ вблизи $db=1$ и докажите, что обучение сходится при $\eta_bd^2+\eta_db^2<2$.
(в)~Смоделируйте при $b_0=1.2$, $d_0=1$ и $\eta_b=\eta_d=0.8$, $0.9$, $1.1$, $1.2$, $1.5$, $2.0$. Что происходит, когда каждый ученик по отдельности устойчив, а вместе --- нет?
(г)~Сформулируйте правило выбора скоростей для инженера, который подстраивает декодер.
\end{exercise}

\begin{exercise}[\lvlB]
\label{ex:17:6}
\emph{Проектирование: конвейер данных вживлённого щупа.} Дано: $N=1024$ канала, $20$~тыс. отсчётов в секунду по $10$~бит, средняя частота нейронов $10$~имп/с, радио $1$~нДж/бит с бюджетом $1$~мВт, задержка до декодера не больше $25\ms$, вероятность потери кадра не больше $10^{-6}$. Шум отсчётов белый нормальный с дисперсией $\sigma^2$.
(а)~Выберите порог выделения импульса так, чтобы ложные срабатывания составляли не больше $1\%$ настоящих ($0.1$~Гц на канал). Что будет при пороге $3.5\sigma$?
(б)~Вариант А: кадры по $1\ms$ фиксированной длины на $C$ событий по $10$ бит плюс заголовок $16$ бит. Найдите наименьшее $C$ и мощность передатчика.
(в)~Вариант Б: числа импульсов каждого канала за $20\ms$, по $2$ бита с насыщением на $3$. Найдите поток, мощность и долю насыщенных отсчётов; оцените выигрыш от энтропийного кодирования.
(г)~Сравните варианты с сырым потоком, выберите решение и объясните, что сломается в каждом варианте при синхронной вспышке всей группы.
\end{exercise}

\begin{exercise}[\lvlC]
\label{ex:17:7}
\emph{Предел скорости интерфейса.} Утверждение~\ref{prop:probe} говорит, что скорость ограничена числом независимых переменных в задаче, а не ёмкостью потока импульсов.
(а)~Постройте верхнюю оценку скорости через размерность $D$ общих переменных, полосу $W$ их изменений и отношение сигнал/шум на переменную, $R\le DW\log_2(1+\text{SNR})$. Обоснуйте выбор чисел ссылками на главы книги и оцените $R$ при шуме, ограниченном корреляцией, и при независимом шуме.
(б)~Сравните с измеренными несколькими битами в секунду. Какие факторы объясняют разрыв?
(в)~Открытая часть. Предложите опыт на реальном интерфейсе, который различит два объяснения узкого места: шум, похожий на сигнал (утверждение~\ref{prop:pooling}), и незнание кода, в частности времён импульсов (глава~\ref{sec:code}). Что вы ожидаете увидеть на графике «извлечённая информация --- число используемых каналов» в каждом случае?
\end{exercise}

\chapter{Типы нейронов по электрической функции}
\chaptermark{Нейроны как приборы}
\label{sec:eetypes}

В главе~\ref{sec:neurontypes} мы рассортировали нейроны по тому, какой медиатор выбрасывает их выход. Электронщик рассортировал бы их иначе: по тому, что нейрон делает со своим сигналом, --- каким прибором он работает. Главный прибор книги мы знаем с главы~\ref{sec:glutcomputer}: интегрирующий нейрон с утечкой~\eqref{eq:glutneuron}, RC-цепь с компаратором. Но в мозге есть и другие. В таблице~\ref{tab:eetypes} они собраны по четырём группам, и почти каждый уже встречался в книге.

\begin{center}
\footnotesize
\setlength{\tabcolsep}{3pt}
\renewcommand{\arraystretch}{1.15}
\begin{longtable}{>{\raggedright}p{3.1cm}>{\raggedright}p{3.3cm}>{\raggedright}p{4.7cm}>{\raggedright\arraybackslash}p{2.4cm}}
\caption{Нейроны как приборы: название, аналог в электронике, что прибор делает и где он встречается в книге.}\label{tab:eetypes}\\
\toprule
Нейрон & Прибор & Что делает & Где в книге \\
\midrule
\endfirsthead
\toprule
Нейрон & Прибор & Что делает & Где в книге \\
\midrule
\endhead
\bottomrule
\endfoot
\multicolumn{4}{l}{\emph{Фильтры и интеграторы}} \\
интегрирующий нейрон с утечкой & RC-интегратор с компаратором & складывает входы в окне $\tau$ и срабатывает у порога & гл.~\ref{sec:glutcomputer}, \eqref{eq:glutneuron} \\
детектор совпадений & И с временным окном & срабатывает, только если входы пришли почти вместе; очень малое $\tau$ & гл.~\ref{sec:glutcomputer}, \ref{sec:code}, \eqref{eq:window} \\
фазический нейрон & фильтр верхних частот, детектор фронта & отвечает на изменение входа и замолкает & гл.~\ref{sec:sensors} \\
адаптирующийся нейрон & фильтр верхних частот с автоматической регулировкой усиления & частота падает при постоянном входе & гл.~\ref{sec:sensors}, \eqref{eq:nakarushton} \\
резонатор & колебательный контур, полосовой фильтр & лучше всего отвечает на вход своей частоты & раздел~\ref{sec:resonator} \\
интегратор без утечки & интегратор на операционном усилителе & превращает скорость в положение и держит его & раздел~\ref{sec:perfectint} \\
\midrule
\multicolumn{4}{l}{\emph{Преобразователи}} \\
тонический нейрон & преобразователь напряжение---частота & частота линейно следует за входом и почти не устаёт & гл.~\ref{sec:code} \\
градуальный нейрон (без импульсов) & аналоговый усилитель, буфер & передаёт плавное напряжение на короткое расстояние & гл.~\ref{sec:sensors} \\
выпрямляющий нейрон & однополупериодный выпрямитель & частота не бывает отрицательной, $[u]_+$ & гл.~\ref{sec:glutcomputer}, \ref{sec:gaba} \\
пара «включение---выключение» & пара выпрямителей, двухпроводный код & один отвечает на «больше», другой --- на «меньше» & гл.~\ref{sec:glutcomputer}, \eqref{eq:dualrail} \\
нейрон с шунтирующим торможением & аналоговый делитель, регулировка усиления & делит вход, а не вычитает из него & гл.~\ref{sec:gaba}, \eqref{eq:shunt} \\
\midrule
\multicolumn{4}{l}{\emph{Генераторы и время}} \\
ритмоводитель & релаксационный генератор, мультивибратор & сам выдаёт импульсы с постоянной частотой & гл.~\ref{sec:serocomputer}, \ref{sec:dofa} \\
ритмоводитель с входом & генератор, управляемый напряжением & свой ритм, который вход ускоряет или замедляет & гл.~\ref{sec:chemoutput} \\
пачечный нейрон & стробируемый генератор пачек & выдаёт пачки частых импульсов --- «тире» & гл.~\ref{sec:code}, \eqref{eq:burst} \\
нейрон, привязанный к фазе & фазовый детектор, фазовая автоподстройка & срабатывает на определённой фазе входной волны & гл.~\ref{sec:sensors}, \ref{sec:chemoutput} \\
аксон с задержкой & линия задержки & задерживает сигнал на заданное время & гл.~\ref{sec:glutcomputer}, рис.~\ref{fig:itd} \\
\midrule
\multicolumn{4}{l}{\emph{Память и переключение}} \\
бистабильный нейрон & триггер Шмитта, защёлка & короткий вход переводит его в длительное состояние & раздел~\ref{sec:bistable} \\
нейрон с дендритными ветвями & мультиплексор, маленькая многослойная сеть & ветвь пропускает вход, только если есть разрешающий вход & гл.~\ref{sec:synthesis} \\
\end{longtable}
\end{center}

Одна оговорка важнее всей таблицы: это не разные клетки, а разные \emph{режимы}. Один и тот же нейрон может быть интегратором при медленном входе и детектором совпадений, когда торможение укоротило его $\tau$ (глава~\ref{sec:code}); ритмоводителем в бодрствовании и молчаливым приёмником во сне. Какой прибор получится, решают ионные каналы мембраны и широковещательные каналы, которые их подстраивают.

\section{Четыре прибора, которых в книге ещё не было}

\subsection{Резонатор}
\label{sec:resonator}

У некоторых нейронов есть медленный ток, который противодействует любому изменению напряжения: напряжение растёт --- ток его тянет вниз, падает --- тянет вверх, но с запаздыванием. Для электронщика такой ток ведёт себя как индуктивность, а вместе с ёмкостью мембраны они образуют колебательный контур. Нейрон отвечает сильнее всего на вход определённой частоты, обычно от единиц до десятка-другого герц, и слабее --- на более медленный и более быстрый~\cite{HutcheonYarom2000}. Это полосовой фильтр в одной клетке: из шумного входа он выбирает ритм своей частоты, например местный ритм главы~\ref{sec:gaba}.

\subsection{Интегратор без утечки}
\label{sec:perfectint}

Интегратор с утечкой помнит вход лишь время $\tau$ --- десятки миллисекунд. Но глазу, чтобы стоять неподвижно, нужно помнить своё положение секундами: команда «повернуть» приходит как кратковременная скорость, а держать глаз надо в новом положении. Сеть решает это той же обратной связью, что и память главы~\ref{sec:glutcomputer}. Если группа нейронов с постоянной времени $\tau$ возбуждает сама себя с весом $w$, то $\tau\,\dd r/\dd t=-r+w\,r+I$, и постоянная времени группы равна
\begin{empheq}[box=\keybox]{equation}
\tau_{\text{эфф}} \;=\; \frac{\tau}{1-w} .
\label{eq:perfectint}
\end{empheq}
При $\tau=0.1\,\text{с}$ и $w=0.995$ получается $20\,\text{с}$: почти идеальный интегратор из нейронов, каждый из которых сам по себе забывает за десятую долю секунды~\cite{Seung1996}. Цена --- точная настройка: при $w$ чуть больше единицы интегратор уходит в насыщение, при чуть меньшей --- быстро забывает. Поэтому такому интегратору нужна постоянная подстройка обучением, вроде той, что описана в главе~\ref{sec:motorlearning}.

\subsection{Бистабильный нейрон}
\label{sec:bistable}

В главах~\ref{sec:glutcomputer} и~\ref{sec:gaba} триггер собирался из группы нейронов. Но бывает триггер и в одной клетке. У некоторых нейронов есть ток, который, однажды включившись, сам поддерживает возбуждение: короткий вход переводит нейрон в длительную работу --- \emph{плато}, --- а торможение возвращает в покой. Это триггер Шмитта: у нейрона два устойчивых состояния и гистерезис между ними. Так ведут себя, например, \nt{ACET}-нейроны, управляющие мышцами, и включается это свойство широковещательным каналом: плато появляется, когда до нейрона доходит серотонин~\cite{Hounsgaard1988}. Одна и та же клетка с \nt{SERO} --- защёлка, без него --- обычный интегратор.

\subsection{Интеграторы и резонаторы}

Теория разделяет возбудимые клетки на два класса~\cite{Izhikevich2007}. \emph{Интеграторы} похожи на генератор, управляемый напряжением, который начинает с нуля: чуть выше порога нейрон работает сколь угодно медленно, и частота плавно растёт со входом. Такой нейрон складывает всё, что пришло, и хорошо передаёт частотой величину входа. \emph{Резонаторы} похожи на генератор с минимальной частотой: медленно работать они не умеют, при пороговом входе сразу выдают импульсы с конечной частотой и предпочитают вход своей частоты. Первые --- естественные приборы частотного кода главы~\ref{sec:code}, вторые --- временного.

\begin{keyblock}
\begin{proposition}[Один нейрон --- много приборов]
\label{prop:eetypes}
По электрической функции нейроны работают как интеграторы с утечкой и без неё, детекторы совпадений, фильтры верхних и полосовых частот, преобразователи напряжение---частота, выпрямители, делители, генераторы, линии задержки и триггеры. Какой прибор получится из данной клетки, задают её ионные каналы и широковещательные каналы, которые их подстраивают. Сами режимы измерены; то, насколько мозг использует такую перестройку для вычислений, --- в основном гипотеза.
\end{proposition}
\end{keyblock}

Для инженера это похоже на программируемую логическую матрицу: железо одно, а прибор задаёт конфигурация. Только конфигурацию здесь меняют не при прошивке, а на ходу, --- медленными широковещательными каналами, о которых шла речь в главах~\ref{sec:serocomputer}--\ref{sec:acet}.

\section{Задачи}

\begin{exercise}[\lvlA]
\label{ex:18:1}
\emph{Допуск интегратора без утечки.} Нейроны с $\tau=0.1$~с держат положение глаза по~\eqref{eq:perfectint} в течение $T=1$~с с ошибкой не более $5\%$ в обе стороны. (а)~Найдите допустимый диапазон $w$. (б)~$w$ --- сумма весов $K$ синапсов, каждый с независимой ошибкой $10\%$. При каком $K$ разброс суммы укладывается в допуск?
\end{exercise}

\begin{exercise}[\lvlB]
\label{ex:18:2}
\emph{Резонатор как контур.} Медленный ток раздела~\ref{sec:resonator} опишем переменной $W$: $C\,\dd V/\dd t=-g_LV-g_wW+I$, $\tau_w\,\dd W/\dd t=V-W$. Покажите, что это $RC$-цепь с параллельной ветвью $R_w=1/g_w$, $L_w=\tau_w/g_w$, и при $C/g_L=10\ms$, $\tau_w=100\ms$, $\alpha=g_w/g_L=4$ найдите частоту резонанса и выигрыш $|Z|$ в нём над $|Z(0)|$.
\end{exercise}

\begin{exercise}[\lvlB]
\label{ex:18:3}
\emph{Как интегратор начинает работать.} (а)~Нейрон $\tau\,\dd V/\dd t=-V+\mu$, $\tau=10\ms$, с порогом $\theta$ и сбросом в ноль. При каком $\mu/\theta-1$ частота равна $1$~Гц? (б)~Какую частоту даёт модель $\tau\,\dd v/\dd t=v^2+\mu$ (импульс --- уход $v$ в $+\infty$, сброс в $-\infty$) при $\mu=0.01$?
\end{exercise}

\begin{exercise}[\lvlB]
\label{ex:18:4}
\emph{Моделирование: интегратор и резонатор.} Модель задачи~\ref{ex:18:2} с $g_L=1$, $\tau_m=10\ms$, $\tau_w=100\ms$, порогом $1$ и сбросом $V\to0$ ($W$ не меняется) получает $\mu=1.5\sin2\pi ft$, $f=1$--$40$~Гц. В какой полосе частот нейрон выдаёт импульсы при $\alpha=0$ и при $\alpha=4$? Сравните с условием $1.5\,|Z(f)|>1$.
\end{exercise}

\begin{exercise}[\lvlB]
\label{ex:18:5}
\emph{Защёлка из одной клетки.} Нейрон с током плато: $\tau\,\dd r/\dd t=-r+F(wr+I)$, $F(x)=\min(1,\max(0,x))$, $\tau=10\ms$, $w=1.5$, фон $I=-0.2$. (а)~Какой наименьшей длительности импульс $I=0.3$ переводит нейрон из $r=0$ в плато? (б)~Какой наименьшей амплитуды $B$ длинный импульс $I=-0.2-B$ возвращает его в покой?
\end{exercise}

\begin{exercise}[\lvlB]
\label{ex:18:6}
\emph{Самонастройка интегратора.} При фиксации взгляда вход $I=0$, датчик скольжения даёт скорость ухода $e=\dd r/\dd t$, и $\dd w/\dd t=-\eta\,e\,r$. При $w=0.95$, $\tau=0.1$~с, $\langle r^2\rangle=1$ найдите $\eta$, при котором за $60$~с фиксаций $|1-w|$ войдёт в допуск задачи~\ref{ex:18:1}. Что сломается, если учиться и во время поворотов глаза?
\end{exercise}

\begin{exercise}[\lvlC]
\label{ex:18:7}
\emph{Зачем перестраивать прибор?} Широковещательный канал переключает $\alpha$ модели задачи~\ref{ex:18:2} между $0$ и $4$. Во сколько раз резонатор выигрывает у интегратора в отношении сигнал/шум для слабого синуса в белом шуме тока при $11$~Гц и при $1$~Гц? Придумайте задачу, где выгодно именно переключение режима.
\end{exercise}

\chapter{Нейроны по числу дендритов}
\chaptermark{Число дендритов}
\label{sec:dendrites}

\section{Число входов как параметр схемы}

В главе~\ref{sec:neurontypes} мы сортировали нейроны по тому, что выбрасывает их выход, а в главе~\ref{sec:eetypes} --- по тому, каким прибором они работают. Есть и третий признак, который инженер заметит сразу: \emph{сколько у нейрона входов}. Входы собирают дендриты --- ветви, на которых садятся чужие аксоны (глава~\ref{sec:dofa}, раздел~\ref{sec:weightstore}). У одних нейронов дендритов нет вовсе, у других один или два, у третьих --- дерево, собирающее сотню тысяч входов. Для электронщика это нагрузочная способность по входу, fan-in, и она почти всегда соответствует задаче.

Почему число и размер дендритов так важны, видно из простой оценки. Мембрана --- конденсатор с удельной ёмкостью $c_m\approx1$~мкФ/см$^2$, и чем больше площадь $A$ нейрона вместе с его дендритами, тем больше заряд нужно принести, чтобы поднять напряжение на $\Delta V$. Время, за которое входной ток $I$ это сделает, если не считать утечки,
\begin{empheq}[box=\keybox]{equation}
t \;\approx\; \frac{c_m\,A\,\Delta V}{I} .
\label{eq:dendriteload}
\end{empheq}
Маленький нейрон с несколькими короткими дендритами имеет площадь порядка $2\cdot10^{-5}$~см$^2$ и ёмкость около $20$~пФ. Один крупный синапс с током в несколько наноампер поднимает его на $20\mV$ меньше чем за десятую долю миллисекунды. У нейрона с большим деревом площадь раз в пятнадцать больше и ёмкость около $300$~пФ; обычный синапс с током около $0.1$~нА зарядил бы его за $60\ms$ --- много дольше постоянной времени утечки, то есть не зарядил бы никогда. Такому нейрону нужны десятки одновременных входов. Числа здесь --- порядки величин, но вывод от них не зависит: \emph{мало входов --- быстро и точно, много входов --- медленно и по сумме}.

\begin{figure}[t]
\centering
\begin{tikzpicture}[>=Stealth,
  soma/.style={circle,draw,very thick,fill=blue!6,minimum size=0.55cm,inner sep=0pt},
  ax/.style={->,very thick,red!70!black},
  den/.style={very thick,blue!60!black}]
% 0 dendrites: sensor on a line
\begin{scope}[xshift=0cm]
\draw[den,decorate,decoration={snake,amplitude=1pt,segment length=4pt}] (-0.9,0) -- (0,0);
\node[font=\scriptsize] at (-0.9,0.3) {датчик};
\draw[ax] (0,0) -- (1.0,0);
\draw[thick] (0.1,0) -- (0.1,0.45);
\node[soma,minimum size=0.4cm] at (0.1,0.65) {};
\node[font=\scriptsize,align=center] at (0.05,-0.75) {$0$\\линия с датчиком};
\end{scope}
% 1 dendrite
\begin{scope}[xshift=3.3cm]
\draw[den] (-0.9,0) -- (-0.28,0);
\node[soma] at (0,0) {};
\draw[ax] (0.28,0) -- (0.9,0);
\node[font=\scriptsize,align=center] at (0,-0.75) {$1$\\буфер, повторитель};
\end{scope}
% 2 dendrites: one per ear
\begin{scope}[xshift=6.2cm]
\draw[den] (-0.28,0) -- (-0.9,0); \draw[den] (0.28,0) -- (0.9,0);
\node[soma] at (0,0) {};
\draw[ax] (0,-0.28) -- (0,-0.6);
\node[font=\scriptsize] at (-0.9,0.25) {лев.}; \node[font=\scriptsize] at (0.9,0.25) {прав.};
\node[font=\scriptsize,align=center] at (0,-1.0) {$2$\\«И» по времени};
\end{scope}
% ~4 short dendrites: mixer
\begin{scope}[xshift=9.0cm]
\foreach \a in {45,135,225,315}{\draw[den] (\a:0.28) -- (\a:0.7); \fill[blue!60!black] (\a:0.72) circle (0.06);}
\node[soma] at (0,0) {};
\draw[ax] (0.28,0) -- (1.0,0);
\node[font=\scriptsize,align=center] at (0,-1.0) {$\sim4$\\смеситель};
\end{scope}
% many: summing tree
\begin{scope}[xshift=12.0cm]
\foreach \a in {60,90,120}{\draw[den] (0,0.28) -- ++(\a:0.55) coordinate (b); \foreach \d in {-20,20}{\draw[den] (b) -- ++(\a+\d:0.35);}}
\foreach \a in {-120,-60}{\draw[den] (0,-0.28) -- ++(\a:0.45);}
\node[soma] at (0,0) {};
\draw[ax] (0.28,0) -- (1.0,0);
\node[font=\scriptsize,align=center] at (0,-1.0) {$10^4$--$10^5$\\сумматор};
\end{scope}
\end{tikzpicture}
\caption{Пять классов по числу входов. Синим --- дендриты (входы), красным --- аксон (выход). Чем меньше входов, тем быстрее и точнее нейрон передаёт отдельный сигнал; чем больше, тем лучше он складывает и классифицирует. Схема, не рисунок формы клеток.}
\label{fig:dendriteclasses}
\end{figure}

\section{Ноль дендритов: датчик на конце линии}

У чувствительных нейронов осязания, боли и растяжения мышц (главы~\ref{sec:sensors}, \ref{sec:pain}, \ref{sec:muscles}) на теле клетки нет ни одного дендрита. Аксон выходит из тела одной веточкой и тут же делится надвое: одна ветвь уходит к коже или мышце, и её конец сам служит датчиком; другая идёт в спинной мозг. Импульс бежит от датчика прямо в спинной мозг, минуя тело клетки. Для инженера это линия связи, у которой датчик стоит на дальнем конце, а тело клетки висит сбоку, как блок питания: оно снабжает линию, но в передаче не участвует. Выгода --- скорость и надёжность: на пути нет ни одного места, где сигнал надо было бы суммировать и снова решать, передавать ли его.

Датчики света и звука --- ещё более крайний случай: это не нейроны-сборщики, а сами преобразователи, у которых входом служит не синапс, а белок-рецептор (рис.~\ref{fig:transducer}).

\section{Один дендрит: буфер}

Нейрон с одним дендритом и одним аксоном принимает одну входную линию и передаёт её дальше. Так устроены обонятельные нейроны: единственный дендрит выходит к поверхности носа и несёт единственный тип рецептора~\cite{Mombaerts2004}; так устроены передаточные клетки сетчатки между датчиками света и выходными нейронами глаза; так устроены нейроны, связывающие датчики уха с мозгом. Инженер назвал бы их буфером или повторителем: они не вычисляют, а доставляют один сигнал, иногда меняя его форму --- например, превращая плавное напряжение датчика в импульсы, как в главе~\ref{sec:sensors}.

\section{Два дендрита: «И» по времени}

Детекторы совпадений, определяющие направление на звук (рис.~\ref{fig:itd}), у млекопитающих часто имеют ровно два дендрита, направленных в противоположные стороны: на один приходят входы от левого уха, на другой --- от правого~\cite{Grothe2010}. Зачем разносить входы? Вспомним формулу~\eqref{eq:shunt}: когда на одном участке мембраны открыто много возбуждающих каналов, напряжение там приближается к равновесному потенциалу, и каждый следующий вход добавляет всё меньше. Поэтому много входов от одного уха, собранных на одном дендрите, насыщают его и не могут в одиночку довести нейрон до порога. А по одному входу с каждого дендрита складываются почти линейно. Два дендрита превращают нейрон в более строгое «И»: срабатывай, только если сигналы пришли с обеих сторон и вместе. На моделях показано, что такое разнесение действительно улучшает детекцию совпадений~\cite{AgmonSnir1998}.

\section{Несколько коротких дендритов: смеситель и ретранслятор}

\emph{Смесители.} В цепи обучения движениям из главы~\ref{sec:motorlearning} около семидесяти миллиардов маленьких \nt{ГЛУТ}-нейронов расширяют вход в очень длинный вектор. У каждого из них всего около четырёх коротких дендритов, и каждый заканчивается «когтем», который охватывает окончание одного входа. Каждый смеситель, таким образом, видит лишь около четырёх входов из многих и работает как маленький элемент «И» над своей четвёркой. Из $M$ входных линий можно выбрать $\binom{M}{4}$ разных четвёрок: при $M=100$ это почти четыре миллиона. Зачем столько? Пороговый элемент с $n$ входами может правильно разделить на два класса не больше
\begin{empheq}[box=\keybox]{equation}
P_{\max} \;\approx\; 2n
\label{eq:cover}
\end{empheq}
случайных образов~\cite{Cover1965}. Выходной нейрон той же цепи получает около $10^5$ входов от смесителей и может выучить порядка $2\cdot10^5$ разных соответствий. Смесители с малым числом входов нужны именно для того, чтобы у большого сумматора было много разных, почти независимых входов~\cite{Marr1969,Albus1971}.

\emph{Ретрансляторы.} На пути от уха к детекторам направления есть нейроны с немногими короткими дендритами, которые получают всего несколько огромных синапсов, а некоторые --- по существу один. По~\eqref{eq:dendriteload} это ровно то, что нужно: маленькая ёмкость и большой ток, и нейрон срабатывает через доли миллисекунды после каждого входного импульса, почти не теряя точности времени. Один из таких ретрансляторов получает \nt{ГЛУТ}, а выдаёт \nt{ГЛИЦ}: это инвертор с точностью в десятки микросекунд, поставляющий быстрое торможение детекторам направления~\cite{BorstSoria2012}.

\section{Тысячи входов: сумматор и классификатор}

На другом конце шкалы --- \nt{ГЛУТ}-нейроны коры с десятком тысяч входов и выходные \nt{ГАМК}-нейроны цепи обучения движениям с сотней тысяч. Такой нейрон по~\eqref{eq:dendriteload} медленный и не может откликнуться на отдельный вход: он складывает. Это интегрирующий нейрон главы~\ref{sec:glutcomputer} и тот сумматор, из которого строятся классификаторы. Большое дерево добавляет и новое: его ветви работают как отдельные подсхемы со своими порогами, так что один нейрон ведёт себя как маленькая многослойная сеть~\cite{Beniaguev2021,Gidon2020}.

\section{Дендриты, которые ещё и выход}

Есть нейроны, у которых аксона нет вовсе, а дендриты служат и входом, и выходом: они получают сигналы и сами выбрасывают медиатор. Самый изученный пример --- \nt{ГАМК}-нейроны сетчатки, участвующие в определении направления движения (глава~\ref{sec:gaba}, рис.~\ref{fig:direction}). У такого нейрона каждая ветвь сама по себе отвечает на движение от тела клетки к своему кончику и выдаёт торможение с этого кончика~\cite{Euler2002}. Один нейрон --- это несколько независимых детекторов направления, по одному на ветвь. Для инженера это не один элемент, а микросхема с несколькими каналами в одном корпусе.

\section{Итог}

\begin{table}[t]
\centering
\footnotesize
\setlength{\tabcolsep}{3pt}
\renewcommand{\arraystretch}{1.15}
\begin{tabular}{>{\raggedright}p{1.9cm}>{\raggedright}p{3.4cm}>{\raggedright}p{2.6cm}>{\raggedright}p{1.8cm}>{\raggedright\arraybackslash}p{3.8cm}}
\toprule
Дендритов & Пример & Прибор & Входов & Что известно \\
\midrule
$0$ & датчики осязания, боли, растяжения & линия с датчиком на конце & --- & установлено \\
$1$ & обонятельные нейроны, передаточные клетки сетчатки, нейроны уха & буфер, повторитель & $1$--немного & установлено \\
$2$ & детекторы направления на звук & «И» по времени & два пучка & строение установлено; выигрыш от разнесения входов показан на моделях \\
$\sim4$ & смесители цепи обучения движениям & маленький элемент «И» & $\sim4$ & установлено; роль расширения входа --- хорошо обоснованная гипотеза \\
немного, короткие & ретрансляторы на пути от уха & ретранслятор, инвертор & $1$--несколько огромных & установлено \\
тысячи & \nt{ГЛУТ}-нейроны коры, выходные нейроны обучения движениям & сумматор, классификатор & $10^4$--$10^5$ & установлено; ветви как подсхемы измерены на отдельных примерах \\
дендриты-выходы & \nt{ГАМК}-нейроны направления в сетчатке & многоканальная микросхема & много & измерено \\
\bottomrule
\end{tabular}
\caption{Нейроны по числу дендритов.}
\label{tab:dendrites}
\end{table}

\begin{keyblock}
\begin{proposition}[Число входов следует за задачей]
\label{prop:dendrites}
Где нужна скорость и точность отдельного сигнала --- на датчиках, ретрансляторах и детекторах совпадений, --- у нейронов мало дендритов, мало входов и большие синапсы: малая ёмкость~\eqref{eq:dendriteload} заряжается быстро. Где нужно складывать и классифицировать, у нейронов большие деревья с тысячами входов. Строение классов измерено; вычислительное объяснение хорошо обосновано, но для многих классов остаётся гипотезой.
\end{proposition}
\end{keyblock}

Таблица~\ref{tab:dendrites} дополняет две прежние сортировки: по медиатору (таблица~\ref{tab:types}) и по прибору (таблица~\ref{tab:eetypes}). Все три смотрят на один и тот же элемент с разных сторон: что он выдаёт, что он вычисляет и сколько он слушает.

\section{Задачи}

\begin{exercise}[\lvlA]
\label{ex:19:1}
\emph{Сколько входов нужно большому нейрону.} Нейрон с большим деревом имеет ёмкость $C=300$~пФ и постоянную времени утечки $\tau_m=20\ms$; каждый обычный синапс даёт ток $0.1$~нА (считайте синапсы источниками тока). Порог --- $20\mV$ над покоем.
(а)~Найдите сопротивление мембраны и установившийся сдвиг напряжения от одного синапса.
(б)~Сколько одновременно работающих синапсов нужно, чтобы дойти до порога вообще; за $10\ms$; за $5\ms$?
(в)~Маленький нейрон ($C=20$~пФ, тот же $\tau_m$) получает один синапс в $4$~нА. За какое время он дойдёт до порога и почему здесь утечкой можно пренебречь?
\end{exercise}

\begin{exercise}[\lvlA]
\label{ex:19:2}
\emph{Бюджет смесителей.} (а)~Проверьте, что из $M=100$ входных линий можно выбрать $\binom{100}{4}\approx3.9\cdot10^6$ четвёрок.
(б)~Выходной нейрон с $n=10^5$ входами по~\eqref{eq:cover} запоминает $P_{\max}$ соответствий по одному биту. Сколько бит приходится на один синапс?
(в)~Во сколько раз расширение входа через смесители увеличивает число соответствий, которые может выучить выходной нейрон, по сравнению с нейроном, получающим прямо $M=100$ исходных линий?
\end{exercise}

\begin{exercise}[\lvlB]
\label{ex:19:3}
\emph{Откуда берётся $2n$.} Число способов разбить $P$ точек общего положения в $n$-мерном пространстве на два класса гиперплоскостью, проходящей через начало координат, равно
\begin{equation*}
C(P,n)=2\sum_{k=0}^{n-1}\binom{P-1}{k}.
\end{equation*}
(а)~Покажите, что при $P\le n$ разделимы все $2^P$ разбиений.
(б)~Докажите, что при $P=2n$ разделима ровно половина разбиений.
(в)~Приблизив биномиальное распределение нормальным, найдите долю разделимых разбиений при $n=100$ и $P=180$, $220$. Какова ширина перехода от «почти все» к «почти никакие» и как она зависит от $n$? Почему~\eqref{eq:cover} становится точным порогом при больших $n$?
\end{exercise}

\begin{exercise}[\lvlB]
\label{ex:19:4}
\emph{Два дендрита как строгое «И».} Упростим детектор направления на звук. Каждый дендрит --- отдельный участок с проводимостью утечки $g_L$; синапсы открывают на нём возбуждающую проводимость, и напряжение участка равно~\eqref{eq:shunt} без торможения. Тело клетки складывает напряжения двух участков с коэффициентом связи $\kappa$: $V_{\text{тело}}=\kappa(V_1+V_2)$. Для сравнения точечный нейрон --- один участок, на который приходят все входы.
(а)~Каждый вход открывает проводимость $g_0$. Покажите, что точечный нейрон не отличает $2n$ входов с одной стороны от $n+n$ с двух сторон.
(б)~Найдите для нейрона с двумя дендритами отношение отклика на $n+n$ к отклику на $2n+0$ как функцию $G=ng_0/g_L$ и его предел.
(в)~Выберите порог $\theta=1.1\,\kappa E_E$. Покажите, что никакое число входов с одной стороны не доводит нейрон до порога, и найдите наименьшее $n$ на каждой стороне, при котором он срабатывает, если $g_0=0.5\,g_L$.
\end{exercise}

\begin{exercise}[\lvlB]
\label{ex:19:5}
\emph{Проектирование ретранслятора.} Дрожание момента срабатывания равно $\sigma_t=\sigma_V/(\dd V/\dd t)$, где $\sigma_V$ --- шум напряжения, а производная берётся в момент пересечения порога. Требуется ретранслятор с $\sigma_t\le20$~мкс при $\sigma_V=1\mV$.
(а)~Какой наименьший синаптический ток нужен нейрону с $C=20$~пФ? Утечку на интервале в доли миллисекунды не учитывайте.
(б)~Сколько обычных синапсов по $0.1$~нА должно сработать одновременно, чтобы той же точности достиг нейрон с $C=300$~пФ? Реалистично ли это?
(в)~Каково дрожание нейрона с $C=20$~пФ и одним синапсом в $4$~нА? Сопоставьте с «точностью в десятки микросекунд» из текста.
\end{exercise}

\begin{exercise}[\lvlB]
\label{ex:19:6}
\emph{Моделирование: заряд длинного дендрита.} Пассивный дендрит длины $\ell$ описывается уравнением кабеля
\begin{equation*}
\tau_m\frac{\partial V}{\partial t}=\lambda^2\frac{\partial^2V}{\partial x^2}-V+\frac{i(x,t)}{g},
\end{equation*}
где $\lambda$ --- постоянная длины, $g$ --- проводимость утечки на единицу длины; концы закрыты ($\partial V/\partial x=0$). Перейдите к безразмерным $T=t/\tau_m$, $X=x/\lambda$, $L=\ell/\lambda$.
(а)~Выведите установившееся отношение $V(0)/V(L)$ при постоянном токе, поданном в дальний конец.
(б)~Напишите на Python модель из $40$ участков (Эйлер, шаг $\Delta T\le0.2\,\Delta X^2$). Подайте в дальний конец короткий импульс тока длительностью $0.01\,\tau_m$ и измерьте на ближнем конце ($X=0$) время максимума и его величину в долях значения $Q/(c_m\ell)$, которое дал бы тот же заряд $Q$, мгновенно размазанный по всему дендриту. Сделайте это для $L=0.2$, $1$ и $2$ и проверьте результат (а).
(в)~Когда оценка~\eqref{eq:dendriteload} по полной ёмкости верна, а когда нет? Что это говорит о коротких дендритах ретрансляторов?
\end{exercise}

\begin{exercise}[\lvlC]
\label{ex:19:7}
\emph{Исследование: оптимальное число входов.} Утверждение~\ref{prop:dendrites} говорит, что вычислительное объяснение числа входов для многих классов остаётся гипотезой. Постройте количественную модель компромисса.
(а)~Пусть ёмкость нейрона $C=C_0+nc_s$ растёт с числом входов $n$, одновременно работает $m$ входов по току $I_s$, и нейрон должен запоминать $P\approx2n$ соответствий~\eqref{eq:cover}. Выразите время отклика через $P$ и найдите, что меняется, если работает не фиксированное число входов, а фиксированная доля $m=\varphi n$.
(б)~Предложите критерий стоимости (время, энергия заряда $CV^2$, ёмкость памяти, дрожание из задачи~\ref{ex:19:5}) и найдите оптимальное $n$ для ретранслятора, детектора совпадений и классификатора. Какие строки таблицы~\ref{tab:dendrites} модель объясняет, а какие нет?
(в)~Большое дерево с ветвями-подсхемами ведёт себя как маленькая многослойная сеть. Предложите численный опыт, сравнивающий ёмкость нейрона из $m$ ветвей по $k$ входов с пороговой нелинейностью в каждой ветви и ёмкость точечного нейрона с $n=mk$ входами.
\end{exercise}

\chapter{Что делают нейроны во сне}
\chaptermark{Сон}
\label{sec:sleep}

\section{Сон --- не выключение, а другой режим}

Инженер, увидев выключенный компьютер, скажет: он ничего не делает. Со спящим мозгом так не выходит. Нейроны во сне не молчат: в глубоком сне кора выдаёт импульсы, в быстром сне --- почти столько же, сколько днём. Меняется не то, \emph{работают} ли нейроны, а то, \emph{как} они работают. Сон --- это набор режимов машины, и переключают их те же широковещательные каналы, что и днём.

Для каждого режима можно выписать «состояние регистров» (таблица~\ref{tab:sleepmodes}). Днём \nt{ACET}, \nt{NORA} и \nt{SERO} работают, датчики подключены, мышцы слушаются. В медленном, глубоком сне все три канала затихают, а вход от органов чувств ослаблен~\cite{Hasselmo1999}: выброс ацетилхолина в коре падает, а \nt{NORA}- и \nt{SERO}-нейроны выдают импульсы реже, чем днём~\cite{Marrosu1995,AstonJonesBloom1981,McGintyHarper1976}. В быстром сне картина странная: \nt{ACET} снова поднимается до дневного уровня~\cite{Marrosu1995}, а \nt{NORA} и \nt{SERO} замолкают почти полностью~\cite{AstonJonesBloom1981,McGintyHarper1976,JacobsAzmitia1992}. Мышцы тела в быстром сне выключены: \nt{ACET}-нейроны, управляющие ими, сильно тормозятся через \nt{GABA} и \nt{GLYC}~\cite{BrooksPeever2012} --- тем же запретом из главы~\ref{sec:gaba}, только наложенным на весь выход сразу.

\begin{table}[t]
\centering
\footnotesize
\setlength{\tabcolsep}{4pt}
\renewcommand{\arraystretch}{1.15}
\begin{tabular}{>{\raggedright}p{3.1cm}>{\raggedright}p{2.9cm}>{\raggedright}p{3.2cm}>{\raggedright\arraybackslash}p{3.4cm}}
\toprule
& бодрствование & медленный сон & быстрый сон \\
\midrule
\nt{ACET} (запись, гл.~\ref{sec:acet}) & высокий & низкий & высокий \\
\nt{NORA} (усиление, гл.~\ref{sec:nora}) & средний, вспышки & низкий & почти молчит \\
\nt{SERO} (гл.~\ref{sec:serocomputer}) & средний & низкий & почти молчит \\
вход от датчиков & подключён & ослаблен & ослаблен \\
выход на мышцы & подключён & ослаблен & выключен торможением \\
активность коры & ровная & медленные качели: работа --- тишина & ровная, как днём \\
\bottomrule
\end{tabular}
\caption{Состояние регистров машины в трёх режимах. Все строки измерены; уровни \nt{ACET}, \nt{NORA} и \nt{SERO} --- прямыми записями по стадиям сна~\cite{Marrosu1995,AstonJonesBloom1981,McGintyHarper1976}.}
\label{tab:sleepmodes}
\end{table}

\section{Когда спать: реле с гистерезисом}

Что решает, когда машине переключаться в сон? Хорошо работает простая схема из двух частей~\cite{Borbely1982}. Первая часть --- \emph{давление сна} $S$, которое копится, пока мы бодрствуем, и сбрасывается, пока спим. Это конденсатор, который заряжается днём и разряжается ночью:
\begin{empheq}[box=\keybox]{equation}
\frac{\dd S}{\dd t} \;=\; \frac{1-S}{\tau_r}\ \ \text{(бодрствование)},\qquad
\frac{\dd S}{\dd t} \;=\; -\frac{S}{\tau_d}\ \ \text{(сон)} .
\label{eq:sleeppressure}
\end{empheq}
Вторая часть --- два порога: машина засыпает, когда $S$ доходит до верхнего порога $H(t)$, и просыпается, когда $S$ опускается до нижнего $L(t)$. Оба порога плавно качаются в такт суточному ритму из главы~\ref{sec:chemoutput}. Получилось реле с гистерезисом --- триггер Шмитта из раздела~\ref{sec:bistable}, --- а вместе с конденсатором оно даёт релаксационный генератор, подстраиваемый по фазе суточным ритмом~\eqref{eq:pll}.

\begin{figure}[t]
\centering
\begin{tikzpicture}[>=Stealth,x=0.16cm,y=4.2cm]
\foreach \a/\b in {0/4.3,21.2/29.3,45.4/53.4,69.5/72}{\fill[blue!8] (\a,0) rectangle (\b,1);}
\draw[->,gray!70!black] (0,0) -- (75,0) node[right,font=\small,black] {$t$, ч};
\draw[->,gray!70!black] (0,0) -- (0,1.08) node[left,font=\small,black] {$S$};
\foreach \t in {0,24,48,72}{\draw[gray!60!black] (\t,-0.02) -- (\t,0.02); \node[below,font=\scriptsize] at (\t,0) {\t};}
\draw[thick,densely dashed,red!70!black] plot file {figures/sleep_H.dat};
\draw[thick,densely dashed,green!45!black] plot file {figures/sleep_L.dat};
\draw[very thick,blue!60!black] plot file {figures/sleep_S.dat};
\node[font=\scriptsize,red!70!black,anchor=west] at (73,0.72) {$H$: заснуть};
\node[font=\scriptsize,green!45!black,anchor=west] at (73,0.24) {$L$: проснуться};
\node[font=\scriptsize,blue!60!black] at (25.25,0.93) {сон};
\node[font=\scriptsize,blue!60!black] at (49.4,0.93) {сон};
\end{tikzpicture}
\caption{Давление сна $S$~\eqref{eq:sleeppressure} при $\tau_r=18.2$~ч, $\tau_d=4.2$~ч. Днём $S$ заряжается до верхнего порога $H$, ночью (закрашено) разряжается до нижнего $L$; оба порога качаются вместе с суточным ритмом. Машина сама приходит к режиму около $16$ часов бодрствования и $8$ часов сна.}
\label{fig:twoprocess}
\end{figure}

В нашей модели с обычными значениями $\tau_r=18.2$~ч и $\tau_d=4.2$~ч машина сама выходит на $16.1$ часа бодрствования и $8.0$ часа сна (рис.~\ref{fig:twoprocess}). Модель объясняет и то, что кажется странным: после сорока часов без сна давление доходит до $0.90$, но восстановительный сон в модели длится лишь $8.8$ часа, а не на сутки дольше обычного. Разряд экспоненциальный, и высокий заряд уходит быстро; «долг» возвращается не длительностью, а глубиной сна.

Что физически играет роль $S$? Сильный кандидат --- аденозин, «счётчик нагрузки» из приложения~\ref{sec:othermediators}: его концентрация растёт за время бодрствования и падает во сне~\cite{PorkkaHeiskanen1997,Dunwiddie2001}. Что давление сна --- это именно аденозин и ничего больше, --- гипотеза.

\section{Медленный сон: вся сеть как релаксационный генератор}

В глубоком сне нейроны коры работают по очереди всей группой: реже раза в секунду они вместе включаются на несколько сотен миллисекунд (\emph{состояние работы}) и вместе замолкают (\emph{состояние тишины}): частота этих качелей ниже $1$~Гц, под наркозом чаще всего $0.3$--$0.4$~Гц~\cite{Steriade1993}. Эти медленные качели мы можем собрать из деталей, которые уже есть. Возьмём группу \nt{GLUT}-нейронов с сильной обратной связью на себя --- память из главы~\ref{sec:glutcomputer}, у которой два устойчивых состояния, --- и добавим медленную усталость, как в генераторе ритма главы~\ref{sec:muscles}:
\begin{empheq}[box=\keybox]{equation}
\tau\,\frac{\dd r}{\dd t} = -r + F\bigl(w\,r + I - a\bigr),
\qquad
\tau_a\,\frac{\dd a}{\dd t} = -a + b\,r .
\label{eq:updown}
\end{empheq}
Группа включается, работает и устаёт; усталость $a$ съедает верхнее состояние, и группа падает в тишину; в тишине усталость проходит, нижнее состояние исчезает, и группа снова включается. Получается тот же релаксационный генератор, что и давление сна, только примерно в восемьдесят тысяч раз быстрее.

\begin{figure}[t]
\centering
\begin{tikzpicture}[>=Stealth,x=2.9cm,y=1.0cm]
\begin{scope}[yshift=1.9cm]
\draw[->,gray!70!black] (0,0) -- (4.1,0);
\draw[->,gray!70!black] (0,0) -- (0,1.25) node[left,font=\small,black] {$r$};
\draw[thick,blue!60!black] plot file {figures/updown_sleep.dat};
\node[font=\scriptsize,anchor=west] at (4.15,0.5) {сон: $b=5$};
\end{scope}
\draw[->,gray!70!black] (0,0) -- (4.1,0) node[right,font=\small,black] {$t$, с};
\draw[->,gray!70!black] (0,0) -- (0,1.25) node[left,font=\small,black] {$r$};
\draw[thick,red!70!black] plot file {figures/updown_wake.dat};
\node[font=\scriptsize,anchor=west] at (4.15,0.5) {день: $b=1$};
\foreach \t in {0,1,2,3,4}{\draw[gray!60!black] (\t,-0.05) -- (\t,0.05); \node[below,font=\scriptsize] at (\t,0) {\t};}
\end{tikzpicture}
\caption{Одна и та же группа~\eqref{eq:updown} в двух режимах: $w=5$, $I=2$, $\theta=2.5$, $\tau=10\ms$, $\tau_a=0.3$~с. С сильной усталостью ($b=5$, как в медленном сне) группа качается между работой и тишиной с периодом $1.1$~с (значение модели; в коре период больше секунды, под наркозом --- около $3$~с). Ослабим усталость ($b=1$) --- так действует ацетилхолин --- и та же группа работает ровно.}
\label{fig:updown}
\end{figure}

А что переводит сеть из качелей в ровную дневную работу? Ацетилхолин подавляет в нейронах медленные токи, создающие усталость~\cite{McCormick1992}. В нашей модели при сильной усталости, $b=5$, группа качается с периодом $1.1$~с --- быстрее измеренных качелей, но того же порядка --- и работает около $35\%$ времени, если усреднять по целому числу периодов; при слабой, $b=1$, та же группа работает ровно (рис.~\ref{fig:updown}). Один регистр --- уровень \nt{ACET} --- переключает генератор в усилитель. Поверх медленных качелей в сне идут и более быстрые вспышки ритма $7$--$15$ колебаний в секунду; их порождает петля \nt{GLUT}--\nt{GABA} между корой и промежуточным релейным узлом~\cite{SteriadeMcCormickSejnowski1993}, родственница ритма из главы~\ref{sec:gaba}.

\section{Повторы: перемотка дня}

В главе~\ref{sec:acet} мы видели, что в медленном сне нейроны, работавшие днём вместе, снова срабатывают вместе и в том же порядке. Добавим одну инженерную деталь: повторы идут \emph{с перемоткой}. Последовательность, занимавшая днём секунды, проигрывается во сне за десятые доли секунды~\cite{Nadasdy1999}: в гиппокампе примерно в двадцать раз быстрее~\cite{LeeWilson2002}, в коре --- в шесть-семь раз~\cite{Euston2007}. И проигрывается не когда придётся: короткие вспышки повторов в одной области подстраиваются под качели работы и тишины и под быстрые вспышки ритма в коре. Если искусственно усилить эту согласованность, память улучшается~\cite{Maingret2016} --- это уже причинная связь, а не совпадение.

Зачем машине перематывать день? Инженеру знакома трудность, которую называют \emph{катастрофическим забыванием}: сеть, которая учит новое подряд, одно за другим, портит старое. Спасает \emph{перемешивание}: учить новое вперемешку со старым, небольшими порциями, много раз. Гипотеза о двух системах обучения~\cite{McClelland1995} состоит в том, что быстрая память записывает день целиком, а во сне понемногу, вперемешку и многократно проигрывает его медленной коре. Это та же пара «быстрый буфер --- медленное хранилище», что в главе~\ref{sec:acet}, и та же пара быстрого и медленного учеников, что в главе~\ref{sec:motorlearning}. Повторы и их согласованность измерены; что их назначение --- перемешанное обучение медленной памяти, --- хорошо обоснованная гипотеза.

\section{Перенормировка весов}

Днём правила Хебба из главы~\ref{sec:dofa} в основном усиливают связи. Если только усиливать, веса растут, и вместе с ними растут расходы энергии и падает чувствительность: нейрон, у которого все входы сильные, перестаёт различать их. По гипотезе \emph{синаптического гомеостаза}~\cite{TononiCirelli2014}, сон нужен в том числе для того, чтобы вернуть веса к рабочему уровню: все связи ослабляются в одно и то же число раз, так что их \emph{отношения} --- выученное --- сохраняются. Для инженера это нормировка, как в правиле~\eqref{eq:oja}, только сделанная не на ходу, а в отдельное время.

Прямые измерения этому не противоречат: у мышей площадь контакта на синапсах коры после сна примерно на $18\%$ меньше, чем после бодрствования, уменьшение пропорционально размеру, а самые крупные и устойчивые контакты почти не меняются~\cite{deVivo2017}. Масштабирование весов измерено; что это главная задача сна, --- гипотеза.

\section{Быстрый сон: машина, отключённая от входов и выходов}

В быстром сне кора работает почти как днём, но вход от органов чувств ослаблен, а выход на мышцы выключен торможением. Для инженера это знакомый режим: \emph{испытание на стенде}. Схема работает на полную мощность, но её входы замкнуты на внутренний генератор, а выходы отсоединены от исполнительных механизмов, чтобы ничего не сломать. Сновидения по одной из гипотез --- именно такая внутренняя прогонка модели мира, собранной за день; что именно сеть при этом делает и учится ли, неизвестно. Измерено здесь только то, что записано в таблице~\ref{tab:sleepmodes}: какой канал включён, какой молчит, и что выход заблокирован.

\section{Обслуживание и местный сон}

Долго считалось, что во сне межклеточные щели расширяются и мозг быстрее промывается от продуктов обмена веществ~\cite{Xie2013}. Недавние опыты, измерявшие промывку другим способом, получили обратное: во сне она медленнее~\cite{Miao2024}. Вопрос сейчас открыт, и мы его оставим открытым.

Надёжнее другой факт: сон --- свойство местных сетей, а не всей машины сразу. Если животное долго не давать спать, то отдельные участки его коры, пока животное бодрствует и двигается, ненадолго проваливаются в тишину, как в медленном сне, и в эти моменты животное чаще ошибается~\cite{Vyazovskiy2011}. Каждая сеть устаёт сама и сама переключается; общий режим получается тогда, когда широковещательные каналы переводят все сети сразу.

\begin{keyblock}
\begin{proposition}[Сон как набор режимов]
\label{prop:sleep}
Во сне нейроны не выключаются: широковещательные каналы переводят машину в другие режимы. Когда спать, решает реле с гистерезисом~\eqref{eq:sleeppressure}, подстроенное суточным ритмом. В медленном сне сеть становится релаксационным генератором~\eqref{eq:updown} и проигрывает день с перемоткой; в быстром --- работает, отключённая от входов и выходов. Режимы, повторы и масштабирование весов измерены; их назначение --- перемешанное обучение и перенормировка --- хорошо обоснованные гипотезы.
\end{proposition}
\end{keyblock}

\section{Итог}

\begin{table}[t]
\centering
\footnotesize
\setlength{\tabcolsep}{4pt}
\renewcommand{\arraystretch}{1.15}
\begin{tabular}{>{\raggedright}p{3.2cm}>{\raggedright}p{4.4cm}>{\raggedright\arraybackslash}p{5.2cm}}
\toprule
Что происходит & Как & Что известно \\
\midrule
смена режимов & регистры \nt{ACET}, \nt{NORA}, \nt{SERO} (табл.~\ref{tab:sleepmodes}) & измерено \\
когда спать & конденсатор $S$ и реле с гистерезисом~\eqref{eq:sleeppressure} & модель хорошо описывает опыты; аденозин как $S$ --- гипотеза \\
медленные качели & группа с обратной связью и усталостью~\eqref{eq:updown} & качели измерены; модель --- простейшая схема \\
перемотка дня & повторы в $6$--$20$ раз быстрее, согласованные с качелями & измерено; усиление согласованности улучшает память \\
зачем повторы & перемешанное обучение медленной памяти & хорошо обоснованная гипотеза \\
перенормировка весов & масштабирование связей во сне & уменьшение контактов измерено; главная роль сна --- гипотеза \\
быстрый сон & работа с отключёнными входами и выходом & режим измерен; назначение неизвестно \\
промывка & расширение межклеточных щелей & противоречивые результаты \\
местный сон & отдельные участки проваливаются в тишину & измерено \\
\bottomrule
\end{tabular}
\caption{Что делают нейроны во сне.}
\label{tab:sleep}
\end{table}

Таблица~\ref{tab:sleep} сводит главу. Сон оказался не паузой в работе машины, а её обслуживанием на ходу: переключение режимов теми же широковещательными каналами, генератор из тех же деталей, что ритм шагов, перемотка дня для медленной памяти и перенормировка весов. Днём машина пишет, ночью переписывает и приводит в порядок записанное.

\section{Задачи}

\begin{exercise}[\lvlA]
\label{ex:20:1}
\emph{Реле без суточного ритма.} Пусть пороги постоянны: $H=0.67$, $L=0.17$ (средние значения порогов модели рис.~\ref{fig:twoprocess}). Выведите из~\eqref{eq:sleeppressure} длительности бодрствования и сна и период генератора при $\tau_r=18.2$~ч, $\tau_d=4.2$~ч. Почему модель с качающимися порогами всё же выходит на $16.1+8.0=24$~ч? К какому прибору главы~\ref{sec:chemoutput} относится этот эффект?
\end{exercise}

\begin{exercise}[\lvlA]
\label{ex:20:2}
\emph{Долг сна.} Сон, начатый при давлении $S_0$, длится $t_s=\tau_d\ln(S_0/L)$, $\tau_d=4.2$~ч, $L$ --- нижний порог при пробуждении. (а)~После $40$~ч без сна $S_0=0.90$, а сон длился $8.8$~ч. Каков был $L$? Сколько длился бы сон при обычном $L\approx0.092$? (б)~За ночь площадь синаптических контактов уменьшается на $18\%$. На сколько должны вырасти веса за день?
\end{exercise}

\begin{exercise}[\lvlB]
\label{ex:20:3}
\emph{Проектирование порогов.} Подберите постоянные пороги $H$ и $L$, при которых реле~\eqref{eq:sleeppressure} с $\tau_r=18.2$~ч и $\tau_d=4.2$~ч даёт ровно $16$~ч бодрствования и $8$~ч сна. Чем такая машина хуже машины с качающимися порогами, если однажды ночью её разбудили через $4$~часа?
\end{exercise}

\begin{exercise}[\lvlB]
\label{ex:20:4}
\emph{Период медленных качелей.} В модели~\eqref{eq:updown} $F(x)=1/\bigl(1+e^{-(x-\theta)/k}\bigr)$, $w=5$, $I=2$, $\theta=2.5$, $k=0.25$, $b=5$, $\tau\ll\tau_a=0.3$~с. (а)~Найдите точки поворота $a_{\text{в}}$ и $a_{\text{н}}$ кривой $r=F(wr+I-a)$, где исчезают работа и тишина. (б)~Приняв $r\approx1$ в работе и $r\approx0$ в тишине, найдите период и долю работы.
\end{exercise}

\begin{exercise}[\lvlB]
\label{ex:20:5}
\emph{Когда качели выключаются.} В модели задачи~\ref{ex:20:4} неподвижная точка лежит на пересечении кривой $a=wr+I-F^{-1}(r)$ с прямой $a=br$. Колебания идут, пока пересечение на средней ветви, между точками поворота. При каких $b$ это так? Как, меняя $b$, \nt{ACET} переключает генератор в усилитель?
\end{exercise}

\begin{exercise}[\lvlB]
\label{ex:20:6}
\emph{Моделирование: долг или фаза.} В \texttt{sleep\_models.py} возьмите первое пробуждение после $48$~ч и держите машину бодрствующей $D=24$, $32$, $40$, $48$, $64$~ч (\texttt{stay\_awake}, \texttt{days=8}). Сколько длится восстановительный сон при каждом $D$? Что определяет его длительность?
\end{exercise}

\begin{exercise}[\lvlB]
\label{ex:20:7}
\emph{Моделирование: забывание и перемешивание.} Линейный нейрон $y=\mathbf w\cdot\mathbf x$, $n=40$, учится по правилу $\mathbf w\leftarrow\mathbf w-0.5\,(y-y^*)\,\mathbf x$. Задачи A и B --- по $10$ образов, $x_j\sim\mathcal N(0,1/n)$, $y^*=\pm1$. Какова среднеквадратичная ошибка на A после $200$ эпох A и затем $200$ эпох B? А после $200$ эпох A и B вперемешку?
\end{exercise}

\begin{exercise}[\lvlC]
\label{ex:20:8}
\emph{Исследование: перемешивание или перенормировка?} Нейрон задачи~\ref{ex:20:7} с порогом; две ночные процедуры: (i)~все веса умножаются на $0.82$; (ii)~старые образы повторяются вперемешку с новыми. Что каждая делает с отношениями весов и с ответами нейрона? Как гипотезы различить по размерам синапсов после сна?
\end{exercise}

\chapter{Машина из живых нейронов}
\chaptermark{Машина из живых нейронов}
\label{sec:livingmachine}

\section{Стенд с открытой схемой}

В главе~\ref{sec:debugging} мы отлаживали машину, у которой нет схемы: щуп слышит тысячу нейронов из восьмидесяти миллиардов, и всё остальное приходится угадывать. Инженер в таком положении сделал бы простую вещь --- собрал бы маленький кусок схемы на макетной плате, где видна каждая деталь. С нейронами это тоже можно сделать.

Нейроны коры разделяют и выращивают на чипе с решёткой электродов --- от нескольких десятков до десятков тысяч. Через несколько недель клетки прорастают отростками и соединяются в сеть из тысяч или сотен тысяч нейронов, в основном \nt{ГЛУТ} и примерно на пятую часть \nt{ГАМК}. Каждый электрод умеет и слушать, как щуп главы~\ref{sec:debugging}, и подавать ток, как тестовый сигнал. Вторая разновидность стенда --- \emph{органоиды}: трёхмерные комочки нервной ткани размером около миллиметра, выращенные из стволовых клеток человека~\cite{Lancaster2013}. На таких стендах живые сети работают месяцами; есть площадки, где опыты на органоидах можно ставить удалённо, по сети, больше ста дней подряд~\cite{Jordan2024}. Эту область называют \emph{биологическими вычислениями} или \emph{органоидным интеллектом}~\cite{Smirnova2023}.

У стенда два важных отличия от мозга. Он крошечный: в нём в сто тысяч раз меньше нейронов. И в нём нет широковещательных каналов: нейронов \nt{ДОФА}, \nt{СЕРО}, \nt{НОРА} и \nt{АЦЕТ} в культуре из коры нет. Это машина, у которой есть шина данных и управление потоком, но нет регистров управления. Посмотрим, на что она способна.

\section{Что сеть делает сама}

Оставленная в покое, культура не молчит и не шумит равномерно. Каждые несколько секунд почти вся сеть вспыхивает разом --- \emph{залпом} --- и снова затихает~\cite{Wagenaar2006}. Для инженера это знакомая картина: усилитель с сильной положительной обратной связью и без нагрузки превращается в генератор.

Механизм мы уже знаем. Сильные взаимные связи \nt{ГЛУТ}-нейронов запускают лавину, как в главе~\ref{sec:gaba} при нарушении условий утверждения~\ref{prop:ei}. Лавина быстро истощает запас медиатора в синапсах~\eqref{eq:depression}, и сеть замолкает. Запас восстанавливается с постоянной времени $\tau_{\text{вос}}$, и когда восстановится доля $x^*$, достаточная для новой лавины, следует новый залп. Интервал между залпами
\begin{empheq}[box=\keybox]{equation}
T_{\text{залп}} \;\approx\; \tau_{\text{вос}}\,\ln\frac{1}{1-x^*} .
\label{eq:burstinterval}
\end{empheq}
При $\tau_{\text{вос}}=0.8\,\text{с}$ из главы~\ref{sec:code} и $x^*=0.9$ получается около двух секунд, при $x^*=0.99$ --- около четырёх, то есть те самые несколько секунд. Это генератор релаксационного типа, родственник генератора ритма из главы~\ref{sec:muscles}: там его перезаряжала усталость групп, здесь --- запас медиатора.

Залпы --- то, что сеть делает, когда ей нечего делать. В живом мозге похожая картина бывает в глубоком сне (глава~\ref{sec:sleep}) и в развивающемся мозге, до того как в него пошли настоящие входы. Сходство это наводит на мысль, но одинаковость механизмов --- гипотеза.

\section{Как учить сеть без учителя-дофамина}

Раз \nt{ДОФА} в культуре нет, сигнал «лучше или хуже» приходится подавать самим --- электродами. Опыты показывают, что сеть на стенде действительно меняет ответы, и самое интересное в них --- \emph{чем} её учат.

\emph{Учитель, который замолкает.} Сеть стимулируют через один электрод раз в секунду-три, пока на другом электроде не появится нужный ответ через $50\pm10\ms$ после стимула. Как только он появился, стимуляцию выключают на пять минут. После нескольких таких циклов нужный ответ появляется на стимул сразу~\cite{Shahaf2001}. Награда здесь --- тишина: стимул расшатывает сеть, а прекращение стимула оставляет её в том состоянии, которое дало правильный ответ. Это спуск по градиенту через случайные отклонения из главы~\ref{sec:dofa} (утверждение~\ref{prop:gradient}), где роль ошибки играет сам факт тряски: трясут, пока не станет правильно.

\emph{Сеть, которая играет в теннис.} В опыте на стенде с плотной решёткой электродов около миллиона нейронов подключили к компьютерной игре в теннис с одной ракеткой~\cite{Kagan2022}. Положение мяча подавалось током на «чувствительные» электроды: где мяч --- кодом местом, как далеко --- частотой. Активность в двух «двигательных» областях двигала ракетку вверх или вниз. Правило обучения было необычным. Если ракетка отбила мяч, сеть получала короткий \emph{предсказуемый} стимул; если промахнулась --- несколько секунд \emph{непредсказуемого} случайного тока. За двадцать минут игры серии отбитых мячей стали длиннее, чем у контрольных культур без обратной связи. Подробный разбор этого опыта --- в главе~\ref{sec:dishbrain}.

Авторы объяснили результат гипотезой, что сеть действует так, чтобы её вход был предсказуемым~\cite{Friston2010}. Это та же идея, что экономия импульсов в главе~\ref{sec:code} и предсказание кадра в главе~\ref{sec:recognition}: машина тратит меньше, когда вход ожидаем. Но и сам опыт, и особенно слова авторов о «разумности» культуры, вызвали резкую критику: эффект невелик, а объяснить его можно и проще~\cite{Balci2023}. Честная оценка такая: культура на стенде меняет поведение под действием обратной связи --- это измерено, а что она учится именно предсказывать свой вход --- гипотеза.

\emph{Сеть, которая водит тело.} Похожим образом культуру подключали к простому виртуальному «телу» и учили его двигаться к цели, подбирая пространственно-временной рисунок тренирующих стимулов~\cite{Bakkum2008}. Ни в одном из этих опытов нет дофамина: учителем служит рисунок тока, который выбирает инженер. Что меняется, когда учителем становится программа или сам дофамин, --- в разделах~\ref{sec:dishdopamine}--\ref{sec:cartpole}.

\begin{figure}[t]
\centering
\begin{tikzpicture}[>=Stealth,
  blk/.style={draw,very thick,rounded corners=2pt,align=center,font=\footnotesize,minimum height=0.8cm,fill=blue!5}]
% (a) closed loop
\node[font=\small] at (1.5,4.3) {(а) замкнутая петля};
\node[blk,text width=2.6cm] (G) at (1.5,3.3) {игра: мяч и ракетка};
\draw[very thick,fill=orange!10,rounded corners=3pt] (0.5,0.2) rectangle (2.5,1.6);
\foreach \x in {0.7,1.0,...,2.35}{\foreach \y in {0.4,0.7,1.0,1.3}{\fill[gray!60] (\x,\y) circle (0.04);}}
\draw[->,thick,blue!60!black] (G.west) -| (-0.5,0.9) -- (0.5,0.9);
\node[font=\scriptsize,blue!60!black,anchor=east] at (-0.6,2.1) {мяч $\to$ ток};
\draw[->,thick,red!70!black] (2.5,0.9) -- (3.5,0.9) |- (G.east);
\node[font=\scriptsize,red!70!black,anchor=west] at (3.6,2.1) {импульсы $\to$ ракетка};
\node[font=\scriptsize] at (1.5,-0.2) {нейроны на решётке электродов};
\node[font=\scriptsize,align=center] at (1.5,-0.85) {отбил: предсказуемый стимул\\промах: случайный ток};
% (b) reservoir
\begin{scope}[xshift=7.9cm]
\node[font=\small] at (1.6,4.3) {(б) резервуар};
\node[blk,text width=1.1cm] (X) at (-0.4,1.0) {вход\\$u(t)$};
\draw[very thick,fill=orange!10] (1.6,1.0) circle (0.8);
\foreach \a/\r in {20/0.35,80/0.5,150/0.3,210/0.55,280/0.4,330/0.2,110/0.15}{\fill[red!60!black] ({1.6+\r*cos(\a)},{1.0+\r*sin(\a)}) circle (0.05);}
\node[font=\scriptsize] at (1.6,-0.2) {органоид: не обучается};
\node[blk,text width=1.8cm,fill=green!8] (W) at (4.0,1.0) {считывание\\$W_{\text{вых}}$};
\draw[->,thick] (X) -- (0.8,1.0);
\draw[->,thick] (2.4,1.0) -- node[above,font=\scriptsize] {$\mathbf r(t)$} (W);
\draw[->,thick] (W) -- (4.0,2.3) node[above,font=\scriptsize] {$y(t)$};
\node[font=\scriptsize,green!40!black] at (4.0,0.3) {обучается};
\end{scope}
\end{tikzpicture}
\caption{Два способа заставить живую сеть вычислять. (а) Замкнутая петля: положение мяча подаётся током, импульсы двигательных областей двигают ракетку, а предсказуемый или случайный стимул после удара служит учителем. (б) Резервуар~\eqref{eq:reservoir}: органоид не обучают вовсе, он только расщепляет вход на множество нелинейных откликов с памятью; учится одно линейное считывание снаружи.}
\label{fig:livingmachine}
\end{figure}

\section{Живой резервуар}

Есть и другой способ использовать живую сеть: не учить её вовсе. В технике он называется \emph{вычислением на резервуаре}~\cite{Maass2002,JaegerHaas2004}. Берут готовую нелинейную систему с памятью --- любую, лишь бы её отклик на вход был богатым и зависел от недавнего прошлого. Вход $u(t)$ раскачивает её, получается вектор откликов $\mathbf r(t)$, и обучается только линейное считывание
\begin{empheq}[box=\keybox]{equation}
y(t) \;=\; W_{\text{вых}}\,\mathbf r(t) ,
\label{eq:reservoir}
\end{empheq}
веса которого подбираются правилом наименьших квадратов~\eqref{eq:lms} из главы~\ref{sec:motorlearning}. Резервуар делает то, что делали маленькие \nt{ГЛУТ}-нейроны главы~\ref{sec:motorlearning}: расщепляет вход на длинный вектор, в котором нужные классы становятся разделимы плоскостью (глава~\ref{sec:dendrites}).

Органоид на решётке электродов подошёл на роль такого резервуара. Звуки речи, поданные на него рисунками тока, после обучения одного линейного считывания различались по говорящему с точностью около $78\%$~\cite{Cai2023}. Результат честный и полезный, но важно понимать, где в нём «ум»: органоид даёт нелинейность и затухающую память, а всё обучение сделано снаружи, в кремнии (рис.~\ref{fig:livingmachine}).

\section{Что стенд говорит о машине}

\emph{Энергия.} По оценке~\eqref{eq:energy} импульс стоит около $10^{-10}$~Дж. Культура из миллиона нейронов при частоте в импульс в секунду тратит на передачу сигналов
\begin{empheq}[box=\keybox]{equation}
P \;\approx\; 10^{6}\times1\,\text{с}^{-1}\times10^{-10}\,\text{Дж} \;=\; 10^{-4}\,\text{Вт},
\label{eq:dishpower}
\end{empheq}
десятую долю милливатта. Электроника, которая стимулирует, записывает и обрабатывает эти импульсы, тратит на порядки больше. Как и в главе~\ref{sec:debugging}, узкое место --- не вычисление, а связь с ним.

\emph{Недостающие детали.} Стенд показывает, сколько умеет шина \nt{ГЛУТ} с управлением потоком \nt{ГАМК} без регистров управления: самоорганизоваться, генерировать залпы, менять ответы под действием обратной связи и служить хорошим резервуаром. Чего она без них не умеет, --- сама решить, что считать успехом. Этот сигнал на стенде подаёт инженер, а в мозге, как мы видели в главах~\ref{sec:dofa}--\ref{sec:acet}, --- широковещательные каналы.

\emph{Этика.} Органоиды выращивают из клеток человека, и чем сложнее они становятся, тем серьёзнее вопрос, может ли такая ткань что-то чувствовать. Инженерный взгляд подсказывает частичный ответ: боль в главе~\ref{sec:pain} --- это не сигнал одного датчика, а целая система тревоги с воротами, регулятором громкости и широковещательными каналами, которых на стенде нет. Но это рассуждение, а не доказательство, и сама область предлагает вести этическую оценку вместе с опытами, с самого начала~\cite{Smirnova2023}.

\begin{keyblock}
\begin{proposition}[Живая сеть на стенде]
\label{prop:livingmachine}
Сеть из \nt{ГЛУТ}- и \nt{ГАМК}-нейронов на решётке электродов сама порождает залпы~\eqref{eq:burstinterval}, меняет ответы под действием обратной связи и может служить резервуаром~\eqref{eq:reservoir} для обученного снаружи считывания. Всё это измерено. Что такая сеть учится по какому-то общему правилу, например предсказывать свой вход, --- гипотеза, а громкие слова о её «разумности» большинство специалистов не принимает.
\end{proposition}
\end{keyblock}

\section{Дофамин по вспышке света}
\label{sec:dishdopamine}

Единственный известный нам прямой опыт, где культуре дали дофамин как учителя, поставлен на платформе из органоидов, к которой исследователи подключаются удалённо~\cite{Jordan2024}. В жидкость, омывающую органоиды, добавили дофамин в светочувствительной «клетке»: связанная молекула не действует на рецепторы. Концентрация клеточного дофамина --- $1$~мМ. Когда нужно наградить сеть, на неё светят ультрафиолетом: клетка распадается, и дофамин выходит в объём.

Петля устроена так. Один и тот же стимул подаётся раз за разом; если он вызвал больше импульсов, чем прежде, включается вспышка и высвобождается дофамин. За $240$ повторов число вызванных импульсов в целом росло. Авторы сами пишут, что по одному такому опыту нельзя заключить, что рост вызван дофамином~\cite{Jordan2024}.

Для инженера это первая попытка собрать в чашке правило трёх множителей~\eqref{eq:threefactor}: активность отправителя и получателя оставляет след, а общий сигнал решает, закрепить ли изменение. Но сигнал здесь бывает только положительным: награда есть или её нет, отрицательной ошибки $\delta<0$ установка не выдаёт. Дофамин к тому же расплывается и откачивается медленно, как концентрация в~\eqref{eq:lowpass}, так что частые вспышки могут просто поднять его общий уровень. Чтобы отличить обучение от этого, нужны два контроля: столько же вспышек в случайные моменты, не связанные с ответом сети, и те же вспышки без клеточного дофамина. И нужна проверка после отдыха: осталось ли изменение, когда дофамин ушёл.

\section{Дофамин учит кремний}
\label{sec:biohybrid}

В обратном опыте живые клетки учат электронику~\cite{Keene2020}. Клетки, выделяющие дофамин, выращивают в микроканале над органическим нейроморфным элементом --- транзистором, проводимость канала которого играет роль веса синапса. Дофамин, вышедший из клеток, окисляется у электрода элемента, и это меняет проводимость. Устройство повторяет и механизм уборки медиатора из щели, поэтому вес можно не только изменить надолго, но и вернуть обратно.

Схемотехнически это интегратор с утечкой по химическому входу: вес копит поток дофамина, а «откачка» задаёт, как быстро он забывается, --- та же RC-цепь, что и~\eqref{eq:lowpass}. Главное здесь направление связи. Щуп из главы~\ref{sec:debugging} не видит широковещательных регистров, потому что они химические. Этот элемент читает именно химический регистр и превращает его в электрический вес. Для интерфейсов мозг---компьютер это путь к датчику, который слышит не только шину данных, но и регистры управления.

\section{Органоиды со своими дофаминовыми нейронами}
\label{sec:midbrainorg}

Из стволовых клеток человека умеют выращивать органоиды, похожие на ту область мозга, где сидят \nt{DOPA}-нейроны. В них есть работающие дофаминовые нейроны, которые выделяют дофамин~\cite{Jo2016}. Такие органоиды выращивают в основном для изучения болезней. Опытов, где их дофамин служил бы учителем для другой сети, мы не нашли.

Для машины из живых нейронов это недостающая деталь. Стенд из главы~\ref{sec:livingmachine} --- шина данных и управление потоком без регистров управления; здесь источник широковещательного сигнала уже есть. Не хватает критика: сети, которая вычисляла бы ошибку предсказания~\eqref{eq:ctd} и управляла бы этими нейронами. Соединить органоид коры с таким органоидом так, чтобы дофамин выделялся по ошибке, --- открытая задача, и пока это гипотеза, а не опыт.

\section{Органоид балансирует шестом без дофамина}
\label{sec:cartpole}

Самый полный из недавних опытов обходится без дофамина совсем~\cite{Robbins2026}. Органоиды коры мыши подключили в замкнутую петлю к задаче «шест на тележке»: активность органоида двигает тележку, положение шеста подаётся обратно стимуляцией. Между попытками органоид получает короткие высокочастотные учебные стимулы. Какие именно --- выбирает программа обучения с подкреплением.

Результаты измерены:
\begin{itemize}
\item у большинства органоидов учебные сигналы, выбранные программой, дали лучший результат, чем случайно выбранные сигналы или отсутствие сигнала;
\item после $45$ минут отдыха улучшение не сохранялось;
\item блокада рецепторов глутамата обоих типов, AMPA и NMDA, отменяла улучшение.
\end{itemize}

Последний пункт говорит, где живёт выученное: в \nt{GLUT}-синапсах, и через тот рецептор, который в разделе~\ref{sec:weightstore} работает как детектор совпадения входа и выхода. Второй пункт показывает, чего не хватает: быстрое изменение есть, закрепления нет, как у быстрого ученика из главы~\ref{sec:motorlearning} без медленного. А роль учителя изменилась: инженера, который выбирает рисунок тока, заменила программа, но учитель по-прежнему стоит снаружи чашки.

Среди более ранних опытов обучения культур на электродных матрицах мы также не нашли ни одного, где использовался бы дофамин: учебным сигналом везде служила электрическая стимуляция.

\section{Кто учит культуру}

\begin{table}[t]
\centering
\footnotesize
\setlength{\tabcolsep}{4pt}
\renewcommand{\arraystretch}{1.15}
\begin{tabular}{>{\raggedright}p{2.7cm}>{\raggedright}p{2.5cm}>{\raggedright}p{3.2cm}>{\raggedright\arraybackslash}p{4.4cm}}
\toprule
Опыт & Кто учит & Учебный сигнал & Что известно \\
\midrule
остановка стимуляции при нужном ответе & инженер & конец стимуляции & ответ закрепляется --- измерено \\
DishBrain (гл.~\ref{sec:dishbrain}) & инженер & предсказуемый или случайный ток & значимый рост розыгрышей за сеанс только при полной обратной связи, при тишине эффект меньше --- измерено; от сеанса к сеансу неустойчив; причина обсуждается \\
шест на тележке & программа обучения с подкреплением & выбранные программой высокочастотные стимулы & лучше случайных, но не сохраняется после отдыха --- измерено \\
дофамин по вспышке & правило «больше импульсов --- награда» & дофамин, освобождённый светом & рост импульсов --- наблюдение; роль дофамина не доказана \\
биогибридный синапс & живые клетки & дофамин от клеток & долгое изменение и возврат веса элемента --- измерено \\
органоиды с \nt{DOPA}-нейронами & --- & --- & источник дофамина есть; как учитель не использован \\
\bottomrule
\end{tabular}
\caption{Кто и чем учит живые нейроны на чипе.}
\label{tab:dishteachers}
\end{table}

Во всех опытах, где учится сама культура, учитель пока снаружи (таблица~\ref{tab:dishteachers}): инженер, программа или правило, заданное инженером. Дофамин вошёл в чашку лишь однажды, и его роль ещё предстоит доказать. Собрать в чашке настоящий широковещательный канал --- с критиком, ошибкой и следом, как в главе~\ref{sec:dofa}, --- следующий шаг, которого никто ещё не сделал.


\section{Итог}

\begin{table}[t]
\centering
\footnotesize
\setlength{\tabcolsep}{4pt}
\renewcommand{\arraystretch}{1.15}
\begin{tabular}{>{\raggedright}p{2.8cm}>{\raggedright}p{4.2cm}>{\raggedright\arraybackslash}p{5.8cm}}
\toprule
Опыт & Что делает живая сеть & Что известно \\
\midrule
сеть без входа & залпы каждые несколько секунд~\eqref{eq:burstinterval} & измерено; механизм через истощение синапсов --- общепринятая модель \\
учитель, который замолкает & нужный ответ закрепляется, когда стимуляцию выключают & измерено \\
игра в теннис & серии отбитых мячей удлиняются при предсказуемой обратной связи & изменение поведения измерено; величина эффекта и объяснение предсказанием входа спорны \\
виртуальное тело & движение к цели при подобранном рисунке стимулов & измерено \\
резервуар & нелинейность и память для считывания~\eqref{eq:reservoir} & измерено; обучение --- снаружи, в кремнии \\
энергия & около $10^{-4}$~Вт на миллион нейронов~\eqref{eq:dishpower} & оценка по~\eqref{eq:energy} \\
\bottomrule
\end{tabular}
\caption{Что показали машины из живых нейронов.}
\label{tab:livingmachine}
\end{table}

Машина из живых нейронов (таблица~\ref{tab:livingmachine}) --- это стенд, на котором видно, что умеют шина данных и управление потоком без регистров управления. Они умеют многое, но не умеют сами решить, что хорошо. На стенде эту роль играет инженер со своим рисунком тока; в мозге --- те немногие широковещательные провода, которым посвящена середина книги.

\section{Задачи}

\begin{exercise}[\lvlA]
\label{ex:21:1}
\emph{Интервал между залпами.} (а)~Выведите~\eqref{eq:burstinterval}, считая, что залп истощает запас до нуля, а затем запас $x$ восстанавливается по закону $\tau_{\text{вос}}\,\dd x/\dd t=1-x$. Обобщите формулу на случай, когда залп оставляет долю $x_0$.
(б)~Найдите $T_{\text{залп}}$ при $\tau_{\text{вос}}=0.8$~с и $x^*=0.9$, $0.99$ и чувствительность $\dd T_{\text{залп}}/\dd x^*$ в обоих случаях.
(в)~Пусть $x^*$ от залпа к залпу случайно меняется на $\pm0.005$ около $0.99$. В каком диапазоне лежат интервалы? Что это говорит о регулярности залпов?
\end{exercise}

\begin{exercise}[\lvlA]
\label{ex:21:2}
\emph{Энергия стенда.} (а)~Проверьте~\eqref{eq:dishpower} и, применив ту же оценку ко всему мозгу ($8.6\cdot10^{10}$ нейронов), сравните с~\eqref{eq:energy}.
(б)~Стенд записывает $1024$ электрода с частотой $20$~кГц по $10$ бит на отсчёт. Каков поток данных? Сколько тратит его передача при $1$~нДж на бит и во сколько раз это больше мощности самой культуры?
\end{exercise}

\begin{exercise}[\lvlB]
\label{ex:21:3}
\emph{Проектирование: подавить залпы.} Залпы мешают культуре работать резервуаром: каждый стирает её память. Предлагается подавать через много электродов редкую фоновую стимуляцию, так что каждый синапс выбрасывает медиатор с частотой $r_b$. Используйте модель истощения главы~\ref{sec:code}: $\tau_{\text{вос}}\,\dd x/\dd t=1-x-U r\tau_{\text{вос}}\,x$ с $U=0.5$, $\tau_{\text{вос}}=0.8$~с.
(а)~Найдите установившийся запас $x_{\text{уст}}(r_b)$ и наименьшую $r_b$, при которой новая лавина невозможна, для $x^*=0.9$ и $0.99$.
(б)~Какова цена этого решения для передачи сигнала через синапс?
(в)~Почему стимуляция должна быть распределена по многим электродам, а не подана через один?
\end{exercise}

\begin{exercise}[\lvlB]
\label{ex:21:4}
\emph{Память резервуара не больше числа нейронов.} Резервуар с $N$ выходами $\mathbf r(t)$ возбуждается стационарным случайным входом $u(t)$ с нулевым средним, единичной дисперсией и независимыми значениями в разные моменты. Для каждой задержки $k\ge0$ линейное считывание~\eqref{eq:reservoir} $y_k(t)=\mathbf w_k\cdot\mathbf r(t)$ подбирается так, чтобы максимизировать квадрат коэффициента корреляции $m_k$ между $y_k(t)$ и $u(t-k)$. Докажите, что ёмкость памяти
\begin{equation*}
\mathrm{MC}=\sum_{k=0}^{\infty}m_k\le N .
\end{equation*}
Указание: рассмотрите случайные величины как векторы со скалярным произведением $\langle X,Y\rangle=\mathbb E[XY]$; что можно сказать о системе $\{u(t-k)\}_k$?
\end{exercise}

\begin{exercise}[\lvlB]
\label{ex:21:5}
\emph{Моделирование: резервуар в кремнии.} Напишите на чистом Python резервуар из $N=30$ элементов: $\mathbf r(t+1)=\tanh\bigl(W\mathbf r(t)+\mathbf v\,u(t)+\mathbf c\bigr)$; $W$ --- случайная гауссова матрица, отнормированная на спектральный радиус $0.9$ (степенной метод); $v_i\sim U(-0.5,0.5)$, $c_i\sim U(-0.2,0.2)$. Считывание --- гребневая регрессия (метод наименьших квадратов с регуляризатором $10^{-6}$ на отсчёт) по $1500$ шагам после $100$ шагов разгона, проверка на следующих $1000$.
(а)~При $u(t)\sim U(-1,1)$ найдите $m_k$ для $k=0,\dots,89$ и $\mathrm{MC}$. Повторите без $\tanh$ (линейный резервуар) при спектральном радиусе $0.9$ и $0.99$. Проверьте неравенство задачи~\ref{ex:21:4}.
(б)~При $u(t)=\pm1$ обучите считывание выдавать $y(t)=u(t-1)\,u(t-2)$ (исключающее ИЛИ). Сравните с линейным считыванием прямо с отводов $u(t)$, $u(t-1)$, $u(t-2)$ и с резервуаром без смещений ($\mathbf c=0$). Объясните последний результат.
\end{exercise}

\begin{exercise}[\lvlB]
\label{ex:21:6}
\emph{Проектирование считывания для живого резервуара.} С органоида записывают $1024$ канала, признаки --- счёт импульсов канала за окно. Для обучения есть $P=240$ записей с двумя классами.
(а)~Пользуясь формулой задачи~\ref{ex:19:3}, объясните, почему считывание по всем $1024$ признакам даст стопроцентную точность на обучающих данных при любых, даже случайных, метках.
(б)~Найдите наибольшее число признаков $n$ (после их объединения в группы), при котором случайная разметка $240$ записей разделима с вероятностью не больше $1\%$.
(в)~На $100$ проверочных записях получена точность $78\%$. Насколько она надёжно отличается от случайного угадывания?
\end{exercise}

\begin{exercise}[\lvlC]
\label{ex:21:7}
\emph{Исследование: учитель без широковещательного канала.} По утверждению~\ref{prop:livingmachine} стенд не умеет сам решить, что считать успехом; эту роль играет рисунок тока.
(а)~Предложите протокол стимуляции, который реализует на стенде трёхфакторное правило главы~\ref{sec:dofa}: что играет роль следа, что --- роль ошибки $\delta$, и как подать её всем синапсам сразу через $8$--$1024$ электродов?
(б)~Сравните ожидаемую скорость обучения с «учителем, который замолкает» (опыт с тишиной как наградой), используя~\eqref{eq:perturbation} и оценку шума считывания~\eqref{eq:rateerror}.
(в)~Как отличить на стенде обучение весов в сети от сдвига её текущего состояния? Предложите контроль, в котором считывание заморожено, и назовите результат, который опроверг бы гипотезу об обучении.
\end{exercise}

\chapter{DishBrain}
\chaptermark{DishBrain}
\label{sec:dishbrain}

\section{Сеть в чашке играет в теннис}

DishBrain --- так авторы назвали стенд, на котором культура живых нейронов, выращенная на чипе с электродами, играет в упрощённый теннис с одной ракеткой~\cite{Kagan2022}. Это самый обсуждаемый опыт с машиной из живых нейронов (обзор таких машин --- в главе~\ref{sec:livingmachine}), и для нашей книги он особенно удобен. Здесь маленькая сеть доступна целиком: каждый вход задаёт экспериментатор, каждый выход записывается. Все вопросы предыдущих глав --- как записан сигнал, кто учит, что видит щуп --- можно задать сети, у которой нет ни глаз, ни мышц, ни \nt{ДОФА}-нейронов. Разберём опыт как инженер разбирает чужой стенд: схема, вход, выход, учитель, результат, споры.

\section{Стенд}

Нейроны берут из коры зародыша мыши --- около $800\,000$ клеток на чип --- или выращивают из стволовых клеток человека, около миллиона на чип. Несколько недель они растут на чипе, прорастают друг в друга и начинают сами выдавать импульсы и пачки. На площадке в $8$~мм$^2$ под ними лежит $26\,000$ платиновых электродов; одновременно можно записывать до $1024$ из них, а стимулировать на практике --- через восемь.

Вокруг культуры замкнута петля (рис.~\ref{fig:dishloop}). Компьютер моделирует игру: мяч летит, отражается от стен, ракетка движется вверх и вниз. Положение мяча превращается в ток на восьми «чувствительных» электродах. Импульсы из двух «двигательных» областей считаются окнами по $10\ms$ и двигают ракетку. После каждого удара или промаха культура получает ещё один стимул --- обратную связь. Окна по $10\ms$ --- это такт компьютера стенда, а не культуры: сама сеть по-прежнему работает асинхронно, как и все сети этой книги.

\begin{figure}[t]
\centering
\begin{tikzpicture}[>=Stealth,
  blk/.style={draw,very thick,rounded corners=3pt,align=center,font=\small,minimum height=1.2cm},
  lab/.style={font=\scriptsize,align=center}]
\node[blk,fill=blue!6,text width=3.4cm] (C) at (0,0) {культура на чипе\\\scriptsize $\sim10^6$ нейронов, $26\,000$ электродов};
\node[blk,fill=orange!10,text width=3.0cm] (G) at (7.2,0) {компьютер:\\игра в теннис};
\draw[->,very thick,blue!60!black] (G.north west) to[bend right=25] node[lab,above] {мяч: \emph{место} --- какой из 8 электродов,\\\emph{частота} --- 4--40 имп/с} (C.north east);
\draw[->,very thick,red!70!black] (C.south east) to[bend right=25] node[lab,below] {ракетка: счёт импульсов\\областей «вверх» и «вниз» за $10\ms$} (G.south west);
\draw[->,thick,densely dashed,green!45!black] (G.east) -- ++(0.9,0) |- ++(0,-2.6) -| node[lab,pos=0.25,below] {отбил: предсказуемый стимул; промах: непредсказуемый} (C.south);
\end{tikzpicture}
\caption{Петля DishBrain. Синяя стрелка --- вход: положение мяча записано кодом местом и частотным кодом (глава~\ref{sec:code}). Красная --- выход: разность счётов двух областей двигает ракетку. Зелёная штриховая --- обратная связь после каждого розыгрыша, единственный «учитель» стенда.}
\label{fig:dishloop}
\end{figure}

\section{Вход и выход}

\emph{Вход} записан сразу двумя кодами главы~\ref{sec:code}. Какой из восьми электродов стимулируется, говорит, на какой высоте мяч: это код местом, как у датчиков главы~\ref{sec:sensors}. Как часто идут импульсы стимуляции, от $4$ до $40$ в секунду, говорит, насколько мяч далеко: это частотный код. Культура не «знает» ни одного из этих кодов заранее: код задан экспериментатором, и сеть должна научиться его читать.

\emph{Выход} читается так же, как в интерфейсах главы~\ref{sec:debugging}. Импульсы одной двигательной области толкают ракетку вверх, другой --- вниз; в простейшей записи за каждое окно
\begin{empheq}[box=\keybox]{equation}
\Delta p \;=\; c\,\bigl(n_\uparrow - n_\downarrow\bigr),
\label{eq:dishreadout}
\end{empheq}
где $n_\uparrow$, $n_\downarrow$ --- числа импульсов двух областей за $10\ms$, а $c$ --- коэффициент стенда. Это двухпроводный код главы~\ref{sec:glutcomputer}: знак движения несёт не знак сигнала, а то, какой из двух проводов громче.

Посмотрим на шум такого выхода. Если область выдаёт в сумме $R$ импульсов в секунду, за окно $T=10\ms$ их в среднем $RT$, и по~\eqref{eq:rateerror} относительный разброс равен $1/\sqrt{RT}$. При $R=1000$ импульсов в секунду это около $30$ процентов в каждом окне; при $R=100$ --- больше ста. Ракетка неизбежно дрожит, и обучиться сеть может лишь одним способом --- сдвинуть \emph{среднюю} разность счётов в нужную сторону. Это те же законы, что ограничивают любой интерфейс главы~\ref{sec:debugging}, только теперь по обе стороны петли.

\section{Учитель без дофамина}

Самое интересное в стенде --- учитель. В культуре нейронов коры нет \nt{ДОФА}-нейронов, и никакого сигнала «лучше, чем ожидалось» стенд не посылает. Обратная связь другая. Отбила сеть мяч --- на все восемь чувствительных электродов сразу подаётся короткий \emph{предсказуемый} стимул: $100$ импульсов в секунду в течение $0.1\,\text{с}$. Промахнулась --- четыре секунды идёт стимул, который авторы называют \emph{непредсказуемым}: импульсы по $150$~мВ, $5$ в секунду~\cite{Kagan2022}.

Авторы исходили из гипотезы, что сеть действует так, чтобы её вход становился предсказуемым~\cite{Friston2010}. Для инженера смысл простой: у культуры есть ровно один способ избавиться от четырёх секунд случайного шума --- отбить мяч. Та же идея встречалась нам в экономии импульсов главы~\ref{sec:code} и в предсказании кадра в главе~\ref{sec:recognition}.

Но нужна ли здесь именно предсказуемость? Возможен и более грубый механизм. Пусть непредсказуемый стимул после промаха просто \emph{перемешивает} недавние изменения в сети, а спокойный после удара их не трогает. Опишем настройку сети одним вектором $\boldsymbol\psi$, и пусть в настройке $\boldsymbol\psi$ сеть промахивается с вероятностью $P_{\text{пр}}(\boldsymbol\psi)$. Если толчки симметричны --- из $\boldsymbol\psi$ в $\boldsymbol\psi'$ так же вероятно, как обратно, --- то в установившемся режиме поток уходов из каждой настройки уравновешен потоком приходов, и доля времени, которую сеть проводит в настройке $\boldsymbol\psi$, равна
\begin{empheq}[box=\keybox]{equation}
\rho(\boldsymbol\psi) \;\propto\; \frac{1}{P_{\text{пр}}(\boldsymbol\psi)} .
\label{eq:dishdensity}
\end{empheq}
Настройка, которая промахивается вдесятеро реже, удерживается вдесятеро дольше. Никто не сообщает сети, в какую сторону менять веса: хорошие настройки просто реже разрушаются.

Мы проверили это на модели, сведя культуру к двум числам: ракетка приходит в точку $a\,y+b$, когда мяч приходит на высоту $y$, и отбивает, если промахнулась меньше чем на $0.1$. После промаха оба числа получают случайный толчок, после удара остаются прежними. Доля отбитых мячей за $2000$ розыгрышей выросла с $0.21$ в первых двухстах до $0.38$ в последних пятистах (сорок «культур», рис.~\ref{fig:dishmodel}). Если толчок давать после \emph{каждого} розыгрыша, сведений в обратной связи нет, и доля остаётся около $0.12$; если толчков нет вовсе, она стоит на $0.19$. Опыт на стенде устроен так же: улучшение было только при полной обратной связи, а при тишине вместо шума после промаха и без обратной связи его не было~\cite{Kagan2022}.

\begin{figure}[t]
\centering
\begin{tikzpicture}[x=1cm,y=5cm]
\draw[->,gray!70!black] (0,0) -- (10.6,0);
\draw[->,gray!70!black] (0,0) -- (0,0.5) node[above,font=\small,black] {доля отбитых};
\foreach \v in {0.1,0.2,0.3,0.4}{\draw[gray!60!black] (-0.06,\v) -- (0.06,\v); \node[left,font=\scriptsize] at (0,\v) {\v};}
\foreach \x/\f/\l/\lab in {1.8/0.21/0.38/{шум после\\промаха},5.2/0.13/0.11/{шум после\\каждого розыгрыша},8.6/0.19/0.19/{без шума}}{
  \fill[blue!35] (\x-0.8,0) rectangle (\x-0.05,\f);
  \fill[red!60!black] (\x+0.05,0) rectangle (\x+0.8,\l);
  \node[below,font=\scriptsize,align=center] at (\x,-0.01) {\lab};
  \node[above,font=\scriptsize] at (\x-0.425,\f) {\f};
  \node[above,font=\scriptsize] at (\x+0.425,\l) {\l};}
\fill[blue!35] (7.4,0.47) rectangle (7.7,0.495); \node[right,font=\scriptsize] at (7.75,0.482) {первые 200};
\fill[red!60!black] (7.4,0.41) rectangle (7.7,0.435); \node[right,font=\scriptsize] at (7.75,0.422) {последние 500};
\end{tikzpicture}
\caption{Модель «перемешивание при промахе» (\texttt{dishbrain\_model.py}): доля отбитых мячей в начале и в конце $2000$ розыгрышей, среднее по сорока моделям-культурам. Обучение идёт, только если толчок следует за промахом, но не за ударом, --- как и в опыте, где сеть училась лишь при полной обратной связи.}
\label{fig:dishmodel}
\end{figure}

\begin{keyblock}
\begin{proposition}[Учитель-перемешиватель]
\label{prop:dishscramble}
Если неудача встряхивает сеть, а удача оставляет её в покое, сеть сама дрейфует к настройкам, которые реже ошибаются: время, проведённое в настройке, обратно пропорционально вероятности неудачи~\eqref{eq:dishdensity}. Для такого обучения не нужны ни сигнал награды, ни его знак, ни предсказуемость как таковая --- достаточно, чтобы обратная связь после удачи и после неудачи различалась. Изменение поведения культуры измерено; какой механизм работает в чашке --- предсказание входа или перемешивание, --- не установлено.
\end{proposition}
\end{keyblock}

Сравните с главой~\ref{sec:dofa}. Там шум тоже служил поиском, но у дофамина был знак: ошибка $\delta$ говорила, в какую сторону сдвинуть веса, и шаг шёл по градиенту~\eqref{eq:perturbation}. Здесь знака нет --- есть только «оставить» или «встряхнуть». Такой поиск грубее и медленнее, но требует от сети меньше: ни \nt{ДОФА}-нейронов, ни критика, ни следа, растянутого на секунды.

\section{Что измерили и о чём спорят}

Каждая игра длилась $20$ минут; сравнивали первые $5$ минут с последними $15$. У культур с полной обратной связью серии отбитых мячей стали длиннее, а мгновенных промахов --- меньше; в контролях без клеток, без игры и с компьютерной «ракеткой» ничего подобного не было. Культуры из клеток человека в среднем отбивали дольше мышиных. Улучшение статистически надёжно --- на девяти мышиных и одиннадцати человеческих культурах, в сотне с лишним игр каждой группы, --- но невелико: сеть не стала хорошим игроком, она стала немного лучше случайного~\cite{Kagan2022}. В следующей работе той же группы нашли, что во время игры активность культуры приближается к \emph{критическому} режиму --- границе между затуханием и лавиной, той самой жизни на грани, что в главе~\ref{sec:gaba}, --- и чем ближе, тем лучше игра~\cite{Habibollahi2023}.

Главный спор вызвали не цифры, а слова. Заголовок статьи утверждал, что нейроны в чашке «проявляют разумность» (sentience), и группа исследователей возразила: изменение поведения в замкнутой петле --- ещё не разум и не ощущение, а выводы нужно формулировать скромнее~\cite{Balci2023}. С инженерной точки зрения открыты два вопроса, и оба --- о механизме. Первый: учится ли сеть, то есть меняются ли её веса надолго, или обратная связь лишь сдвигает её текущее состояние? Второй: нужна ли именно предсказуемость, или хватает любого различия между откликом на удачу и на неудачу, как в утверждении~\ref{prop:dishscramble}? Стенд, на котором доступен каждый электрод, позволяет ответить на оба --- и это, пожалуй, его главная ценность.

\section{Спецификация стенда: как его воспроизвести}
\label{sec:dishspec}

Ниже собрано всё, что нужно, чтобы собрать такой стенд заново: оборудование, петля, кодирование входа, чтение выхода, обратная связь и протокол опыта. Каждый параметр помечен источником. «Статья» --- значение взято из текста работы~\cite{Kagan2022}. «Выбор» --- в статье его нет, и для воспроизведения его придётся задать самому; мы предлагаем значение по умолчанию и говорим, как его проверить.

\subsection*{Оборудование и культура}

\begin{center}
\footnotesize
\setlength{\tabcolsep}{4pt}
\renewcommand{\arraystretch}{1.15}
\begin{tabular}{>{\raggedright}p{4.0cm}>{\raggedright}p{6.9cm}>{\raggedright\arraybackslash}p{1.6cm}}
\toprule
Параметр & Значение & Источник \\
\midrule
чип & матрица MaxOne: $26\,000$ платиновых электродов на площадке $8$~мм$^2$ & статья \\
запись & до $1024$ каналов одновременно, $20\,000$ отсчётов в секунду & статья \\
стимуляция & до $32$ электродов технически, в опыте --- $8$ независимых & статья \\
клетки & кора зародыша мыши; нейроны человека из стволовых клеток (два способа выращивания); клетки HEK293T (не нейроны) для контроля & статья \\
посев & около $10^6$ клеток на чип & статья \\
возраст культуры к опыту & мышь: с $\sim10$ суток; человек: $14$--$24$ суток или $\sim80$ суток в зависимости от способа выращивания & статья \\
\bottomrule
\end{tabular}
\end{center}

\subsection*{Петля и время}

Петля работает окнами по $10\ms$ ($200$ отсчётов): за окно считаются импульсы на электродах двигательных областей, игра обновляется, и выдаётся очередная стимуляция. Задержка от импульса в культуре до ответного стимула --- около $5\ms$. Во время игры сигнал каждого электрода проходит фильтр верхних частот Бесселя 2-го порядка с частотой среза $100$~Гц; модуль сигнала сглаживается фильтром нижних частот Бесселя 1-го порядка с частотой $1$~Гц, и порог импульса пропорционален этому сглаженному модулю. Порог в шесть стандартных отклонений фона статья указывает для отдельных измерений активности, а не для игры. Чтобы стимуляция не засчитывалась как ответ культуры, все сигналы гасятся, когда одновременно приходит больше $15$ крупных, выше $75$~мВ, импульсов. Всё это --- статья.

\subsection*{Кодирование мяча}

\begin{center}
\footnotesize
\setlength{\tabcolsep}{4pt}
\renewcommand{\arraystretch}{1.15}
\begin{tabular}{>{\raggedright}p{3.4cm}>{\raggedright}p{7.5cm}>{\raggedright\arraybackslash}p{1.6cm}}
\toprule
Параметр & Значение & Источник \\
\midrule
чувствительная область & $8$ электродов, расположенных так, чтобы соседние электроды означали соседние положения мяча & статья \\
код местом & номер электрода $k=1,\dots,8$ задаёт положение мяча относительно центра ракетки & статья \\
какая координата & вертикальное положение мяча; для воспроизведения --- относительно ракетки, $k=\lceil 8(y-p+\tfrac12)\rceil$ при $y-p\in[-\tfrac12,\tfrac12]$ & статья; формула --- выбор \\
код частотой & частота стимуляции от $4$ до $40$~имп/с несёт расстояние до ракетки & статья \\
направление шкалы & $4$~имп/с, когда мяч у противоположной стены, и линейно до $40$~имп/с, когда он у стены ракетки: $f=4+36\,(1-x/L)$, $x$ --- расстояние до стены ракетки, $L$ --- ширина корта & статья \\
форма импульса & короткий прямоугольный двухфазный: сначала положительная, потом отрицательная фаза & статья \\
амплитуда & $75$~мВ; выбрана так, чтобы не заставлять срабатывать заторможенные клетки & статья \\
длительность фаз & не указана; взять стандартную настройку системы стимуляции и не менять её в течение опыта & выбор \\
\bottomrule
\end{tabular}
\end{center}

\subsection*{Чтение ракетки}

В статье две двигательные области: импульсы в первой двигают ракетку вверх, во второй --- вниз; вклад областей взвешивается, чтобы учесть разную плотность клеток и разную активность~\cite{Kagan2022}. Точная формула не приведена. Для воспроизведения берём правило~\eqref{eq:dishreadout}, $\Delta p=c\,(n_\uparrow-n_\downarrow)$ за каждое окно $10\ms$, с весом, уравнивающим области по их средней активности в покое (выбор): $n_\uparrow$ и $n_\downarrow$ делятся на средний счёт своей области за окно, измеренный перед игрой. Коэффициент $c$ выбираем так, чтобы разница в один средний счёт сдвигала ракетку на $1\%$ высоты корта за окно (выбор); тогда постоянный перевес одной области проводит ракетку через весь корт за $1$~с.

Расположение областей на чипе в тексте статьи координатами не задано; есть только схема. Для воспроизведения достаточно трёх непересекающихся групп электродов: восемь чувствительных в центре и по группе записывающих электродов с двух сторон, одна --- «вверх», другая --- «вниз» (выбор).

\subsection*{Игра}

Размеры корта, длина ракетки и скорость мяча в статье не указаны; сказано только, что мяч летит с постоянной скоростью и отражается от стен. Для воспроизведения (выбор): корт $1\times1$, ракетка на правой стенке длиной $0.2$ (попадание при $|y-p|<0.1$, как в модели~\texttt{models/dishbrain\_model.py}), мяч пересекает корт за $1$~с. Промах без единого отбитого мяча в розыгрыше считается «эйсом»; розыгрыш длиннее трёх ударов подряд --- «длинным» (статья).

\subsection*{Обратная связь}

\begin{center}
\footnotesize
\setlength{\tabcolsep}{4pt}
\renewcommand{\arraystretch}{1.15}
\begin{tabular}{>{\raggedright}p{2.6cm}>{\raggedright}p{8.3cm}>{\raggedright\arraybackslash}p{1.6cm}}
\toprule
Событие & Что подаётся & Источник \\
\midrule
удар & предсказуемый стимул: все $8$ чувствительных электродов одновременно, $75$~мВ, $100$~имп/с, $100\ms$; обычная стимуляция мяча на это время прерывается & статья \\
промах & непредсказуемый стимул: $150$~мВ, $5$~имп/с, $4$~с, в случайные моменты на случайно выбранные из $8$ электродов; затем $4$~с без стимуляции, и мяч стартует в случайном направлении & статья \\
закон случайности & не указан; для воспроизведения --- моменты импульсов как пуассоновский поток со средней частотой $5$~имп/с & выбор \\
\bottomrule
\end{tabular}
\end{center}

\subsection*{Протокол и контроль}

Один сеанс длится $20$~мин; сравнивают первые $5$ и последние $15$~мин (статья). Три меры результата: средняя длина розыгрыша, доля эйсов и доля длинных розыгрышей. Условия опыта (статья):
\begin{itemize}
\item \emph{стимул} --- полная петля, как описано выше;
\item \emph{тишина} --- вместо стимула удара или промаха столько же времени без всякой стимуляции, затем мяч стартует в случайном направлении;
\item \emph{без обратной связи} --- мяч после промаха просто отражается и летит дальше, стимуляция положения мяча продолжается;
\item \emph{покой} --- игра идёт и ракетка движется по активности, но стимуляции нет совсем;
\item \emph{контроли} --- чип только со средой; клетки HEK293T; «культура в компьютере», симуляция вместо клеток.
\end{itemize}
Объём выборки в основной серии: $9$ культур мыши ($101$ сеанс), $11$ культур человека ($138$), $6$ чипов со средой ($80$), $20$ культур в покое ($42$), $3$ симуляции ($38$), всего $399$ сеансов. В серии про обратную связь --- $486$ сеансов за $3$ дня по $3$ сеанса в день, условия «стимул», «тишина» и «без обратной связи» ($n=27$, $17$ и $15$). Рост средней длины розыгрыша от первых $5$ минут к последним $15$ получен только у живых культур и только в условии «стимул»~\cite{Kagan2022}.

\subsection*{Цикл стенда по шагам}

Каждые $10\ms$ компьютер стенда делает следующее.
\begin{enumerate}
\item Считывает число импульсов $n_\uparrow$, $n_\downarrow$ в двух двигательных областях за прошедшее окно и нормирует их на средний счёт области в покое.
\item Сдвигает ракетку на $\Delta p=c\,(n_\uparrow-n_\downarrow)$ и ограничивает её краями корта.
\item Сдвигает мяч на один шаг; отражает его от верхней, нижней и левой стенок.
\item Если мяч дошёл до правой стенки: при $|y-p|<0.1$ --- удар, счёт розыгрыша растёт на один, подаётся стимул удара ($100\ms$); иначе --- промах, розыгрыш записывается, подаётся стимул промаха ($4$~с), затем $4$~с паузы, и мяч стартует заново.
\item Иначе выбирает электрод $k$ по вертикальному положению мяча и частоту $f$ по расстоянию до стены ракетки и подаёт на электрод $k$ двухфазные импульсы $75$~мВ с частотой $f$ до следующего окна.
\end{enumerate}
Этого достаточно, чтобы собрать стенд, совместимый с опубликованным, и проверить главный вопрос главы: учит ли культуру предсказуемость или простая разница между ответом на удар и на промах. Для этого к трём условиям статьи стоит добавить четвёртое, которого в ней нет: после промаха --- \emph{предсказуемый} стимул той же длительности и энергии, что и непредсказуемый. Если культура учится и так, дело в разнице, а не в предсказуемости.


\section{Итог}

\begin{table}[t]
\centering
\footnotesize
\setlength{\tabcolsep}{4pt}
\renewcommand{\arraystretch}{1.15}
\begin{tabular}{>{\raggedright}p{2.1cm}>{\raggedright}p{4.2cm}>{\raggedright}p{1.9cm}>{\raggedright\arraybackslash}p{4.6cm}}
\toprule
Узел стенда & Как устроен & Глава книги & Что известно \\
\midrule
вход & место (какой из 8 электродов) и частота (4--40 имп/с) & \ref{sec:code}, \ref{sec:sensors} & задано экспериментатором \\
выход & разность счётов двух областей за $10\ms$~\eqref{eq:dishreadout} & \ref{sec:glutcomputer}, \ref{sec:debugging} & задано экспериментатором; шум по~\eqref{eq:rateerror} \\
учитель & предсказуемый стимул после удара, непредсказуемый после промаха; дофамина нет & \ref{sec:dofa} & улучшение при полной обратной связи измерено; без неё его нет \\
механизм & предсказание входа или перемешивание при промахе~\eqref{eq:dishdensity} & \ref{sec:dofa}, \ref{sec:recognition} & не установлен; оба объясняют результат \\
режим сети & приближение к критическому режиму во время игры & \ref{sec:gaba} & связь с качеством игры измерена \\
«разумность» & утверждение авторов & --- & не показана; оспорена \\
\bottomrule
\end{tabular}
\caption{DishBrain глазами инженера.}
\label{tab:dishbrain}
\end{table}

DishBrain --- машина из живых нейронов в самом простом виде: вход задан кодом, выход прочитан счётом, учитель --- различие между откликом на удачу и на неудачу (таблица~\ref{tab:dishbrain}). Сеть без глаз, мышц и дофамина чуть-чуть научилась играть, и уже одно это превращает её в испытательный стенд для идей книги. Каким бы ни оказался механизм --- предсказание, перемешивание или что-то третье, --- ответ будет найден так, как советует глава~\ref{sec:debugging}: щупом, тестовым сигналом и выключением частей на машине, которую наконец видно целиком.

\section{Задачи}

\begin{exercise}[\lvlA]
\label{ex:22:1}
\emph{Сколько нужно сдвинуть счёт.} Две двигательные области выдают вместе $R_\uparrow+R_\downarrow=2000$ импульсов в секунду, потоки пуассоновские и независимые.
(а)~Каков относительный разброс счёта одной области за окно $10\ms$ при $R=1000$ и $R=100$ импульсов в секунду?
(б)~Мяч летит до ракетки $1$~с, и за это время смещения~\eqref{eq:dishreadout} складываются. Какой наименьшей разности $R_\uparrow-R_\downarrow$ достаточно, чтобы среднее смещение ракетки вдвое превышало его разброс? Какую долю от суммарной частоты это составляет?
\end{exercise}

\begin{exercise}[\lvlA]
\label{ex:22:2}
\emph{Сколько розыгрышей нужно, чтобы увидеть обучение.} Считайте розыгрыши независимыми, а долю отбитых мячей --- биномиальной.
(а)~Сколько розыгрышей нужно в каждом из двух периодов, чтобы рост доли отбитых с $0.20$ до $0.25$ превышал два стандартных отклонения разности?
(б)~То же для роста с $0.21$ до $0.38$, как в модели рис.~\ref{fig:dishmodel}.
(в)~Почему в опыте розыгрыши одной культуры нельзя считать независимыми и зачем поэтому нужны десятки культур?
\end{exercise}

\begin{exercise}[\lvlB]
\label{ex:22:3}
\emph{Вывод~\eqref{eq:dishdensity}.} Пусть настройка $\boldsymbol\psi$ принимает значения из конечного множества. В каждом розыгрыше сеть промахивается с вероятностью $P(\boldsymbol\psi)>0$ и тогда переходит в $\boldsymbol\psi'$ с вероятностью $K(\boldsymbol\psi,\boldsymbol\psi')=K(\boldsymbol\psi',\boldsymbol\psi)$; при ударе остаётся на месте.
(а)~Докажите, что $\rho(\boldsymbol\psi)\propto1/P(\boldsymbol\psi)$ --- стационарное распределение, и укажите условие его единственности.
(б)~Докажите, что установившаяся доля промахов равна среднему гармоническому $P$ по всем настройкам, и что она не больше доли промахов при отсутствии обратной связи. Когда достигается равенство?
(в)~Пример: две настройки с $P=0.9$ и $P=0.1$. Найдите обе доли промахов.
(г)~Что произойдёт, если $P=0$ в некоторых настройках?
\end{exercise}

\begin{exercise}[\lvlB]
\label{ex:22:4}
\emph{Поглощающая область модели.} В модели рис.~\ref{fig:dishmodel} ракетка приходит в $p=\min\bigl(1,\max(0,ay+b)\bigr)$, мяч --- на высоту $y\sim U(0,1)$, удар --- если $|p-y|<h$, $h=0.1$.
(а)~Покажите, что ограничение $p$ отрезком $[0,1]$ никогда не увеличивает промах.
(б)~Докажите, что при $|b|<h$ и $|a-1+b|<h$ сеть отбивает все мячи, и найдите площадь этой области на плоскости $(a,b)$. Почему такие настройки --- поглощающие?
(в)~Найдите численно среднюю долю отбитых мячей при начальном распределении модели ($a\sim U(-1,1)$, $b\sim U(0,1)$) и сравните со столбцом «без шума» рис.~\ref{fig:dishmodel}.
\end{exercise}

\begin{exercise}[\lvlB]
\label{ex:22:5}
\emph{Моделирование с \texttt{dishbrain\_model.py}.} Работайте в копии модели; усредняйте по $200$ культурам (затравки $0$--$199$).
(а)~Увеличьте число розыгрышей до $20\,000$. Какой будет доля отбитых в последних $500$? Какая доля культур к концу оказалась в поглощающей области задачи~\ref{ex:22:4}? Почему доля отбитых не стремится к единице?
(б)~Ограничьте настройку прямоугольником $a\in[-1,2]$, $b\in[-1,1]$ (обрезайте после каждого толчка) и повторите для $2000$ и $20\,000$ розыгрышей. Объясните разницу с помощью задачи~\ref{ex:22:3}.
(в)~Для $2000$ розыгрышей постройте долю отбитых в последних $500$ как функцию силы толчка \texttt{sigma} от $0.02$ до $0.64$. Где оптимум и почему он порядка $h$?
\end{exercise}

\begin{exercise}[\lvlB]
\label{ex:22:6}
\emph{Проектирование входного кода.} Ракетка отбивает мяч, если ошибка по высоте меньше $h=0.1$. Сеть успевает прочитать вход за $T=0.25$~с; частота стимуляции ограничена $4$--$40$ импульсами в секунду; по главе~\ref{sec:code} до частоты $r$ за время $T$ различимо около $\sqrt{rT}$ уровней.
(а)~Сколько уровней высоты нужно различать? Хватает ли кода местом на восьми электродах?
(б)~Можно ли было бы передать высоту частотой на одном электроде? Какая предельная частота для этого понадобилась бы?
(в)~Сколько бит о расстоянии до мяча несёт частотный код стенда? Предложите распределение двух величин между кодами при том же числе импульсов стимуляции и обоснуйте его.
\end{exercise}

\begin{exercise}[\lvlC]
\label{ex:22:7}
\emph{Исследование: предсказание или перемешивание?} Глава оставляет открытым, что работает в чашке.
(а)~Обобщите модель на $d$ настраиваемых чисел (например, ракетка $p=\sum_{i=0}^{d-1}a_iy^i$) и численно найдите, как число розыгрышей до доли отбитых $0.5$ растёт с $d$ при перемешивании. Сравните с правилом, в котором знак ошибки известен (глава~\ref{sec:dofa}, \eqref{eq:perturbation}). Какая зависимость от $d$ ожидается в каждом случае?
(б)~Предложите опыт на стенде, в котором две гипотезы дают противоположные предсказания. Например: после промаха подавать \emph{предсказуемый}, но сильный стимул, а после удара --- непредсказуемый, но слабый. Что предскажет каждая гипотеза?
(в)~Оцените по задаче~\ref{ex:22:2}, сколько культур и розыгрышей нужно, чтобы опыт (б) дал надёжный ответ при ожидаемой разнице долей около $0.05$.
\end{exercise}

\appendix
\renewcommand{\chaptername}{\appendixname}
\chapter{Вещества, которые не задают тип нейрона}
\label{sec:othermediators}

В главе~\ref{sec:neurontypes} мы рассортировали нейроны по главному медиатору и получили десять типов (таблица~\ref{tab:types}). Теперь пора добавить осложнения. Нейрон выбрасывает не только главный медиатор, а в мозге кроме главных медиаторов работают и другие вещества. Они меняют то, как сеть передаёт и хранит сигналы.

Кроме адресной и широковещательной шин (раздел~\ref{sec:buses}) есть и третья: \emph{обратный канал}. Несколько веществ выбрасывает не отправитель, а \emph{получатель}, и они идут через щель назад, к выходу отправителя, и меняют, сколько медиатора тот выбросит в следующий раз. В сетях связи это похоже на подтверждение приёма, которое отправитель использует, чтобы настроить свою передачу. Именно этим каналом нейроны пользуются, когда меняют веса своих связей: по нему выход отправителя узнаёт об активности получателя, то есть о втором множителе правил обучения главы~\ref{sec:dofa}.

Полный список известных медиаторов длинен: больше сотни веществ, большинство из них --- короткие белковые цепочки, нейропептиды. Но все они укладываются в несколько классов, и внутри каждого класса вещества ведут себя сходно. Вещества, которые не бывают главным медиатором, собраны в таблице~\ref{tab:othermediators}, с теми же тремя вопросами, которые задал бы инженер: по какой шине идёт сигнал, быстро ли он действует и какой у него обычно знак. Тип нейрона они не задают: они выбрасываются вместе с главным медиатором, или выбрасываются не нейронами, или идут в обратную сторону.

\begin{table}[t]
\centering
\footnotesize
\setlength{\tabcolsep}{4pt}
\renewcommand{\arraystretch}{1.12}
\begin{tabular}{>{\raggedright}p{2.25cm}>{\raggedright}p{2.8cm}>{\raggedright}p{1.8cm}>{\raggedright}p{1.5cm}>{\centering}p{0.8cm}>{\raggedright\arraybackslash}p{4.3cm}}
\toprule
Класс & Вещество & Шина & Скорость & Знак & Роль в вычислении \\
\midrule
Аминокислоты & D-серин & локальная & медленно & И & второй ключ для рецептора глутамата: условие «И» \\
\midrule
Моноамины & следовые амины & широковещ. & медленно & $\pm$ & тонкая настройка моноаминовых каналов \\
\midrule
\multirow{2}{2.25cm}{Пурины}
 & АТФ & адресная & быстро и медленно & $+$ & сопутствующий медиатор; сигнал между нейронами и глией \\
 & аденозин & объёмная & медленно & $-$ & накапливается при активности: счётчик нагрузки~\cite{Dunwiddie2001} \\
\midrule
\multirow{8}{2.25cm}{Нейропептиды (больше сотни)}
 & энкефалины, эндорфины, динорфин & объёмная & медленно & $-$ & подавление передачи \\
 & вещество P & объёмная & медленно & $+$ & медленное усиление \\
 & нейропептид Y & объёмная & медленно & $-$ & медленное торможение \\
 & соматостатин & объёмная & медленно & $-$ & медленное торможение, сопутствует GABA \\
 & вазоактивный интестинальный пептид (ВИП) & объёмная & медленно & $+$ & растормаживание: торможение тормозящих нейронов \\
 & холецистокинин & объёмная & медленно & $\pm$ & сопутствует GABA в части тормозящих нейронов \\
 & орексин & широковещ. & медленно & $+$ & стабилизатор режима бодрствования \\
 & окситоцин, вазопрессин & широковещ. & медленно & $\pm$ & глобальные настройки поведения \\
\midrule
\multirow{2}{2.25cm}{Газы}
 & оксид азота NO & обратная, объёмная & медленно & $\pm$ & проходит сквозь мембраны; обратный сигнал при изменении весов~\cite{Garthwaite2008} \\
 & CO, H$_2$S & объёмная & медленно & $\pm$ & роль изучена мало \\
\midrule
Липиды & эндоканнабиноиды (анандамид, 2-АГ) & обратная & медленно & $-$ & получатель уменьшает выброс у отправителя: обратная связь по весу~\cite{WilsonNicoll2001} \\
\bottomrule
\end{tabular}
\caption{Вещества, которые не задают тип нейрона. \emph{Шина}, \emph{скорость} и \emph{знак} --- как в таблице~\ref{tab:types}; кроме того, объёмная шина --- медиатор расплывается в межклеточном пространстве вокруг места выброса, обратная --- идёт от получателя назад к отправителю, локальная --- действует рядом с местом выброса. «И» --- разрешающий сигнал: сам по себе ток не вызывает, но без него не срабатывает другой вход, как второй вход логического элемента «И». Нейропептиды почти всегда выбрасываются вместе с главным медиатором и в основном при высокой частоте импульсов~\cite{vandenPol2012}.}
\label{tab:othermediators}
\end{table}

\chapter{Обозначения}
\label{sec:notation}

Типы нейронов обозначены первыми четырьмя буквами английского названия главного медиатора: \nt{GLUT} --- глутамат, \nt{GABA} --- гамма-аминомасляная кислота, \nt{GLYC} --- глицин, \nt{ACET} --- ацетилхолин, \nt{DOPA} --- дофамин, \nt{SERO} --- серотонин, \nt{HIST} --- гистамин, \nt{NORA} --- норадреналин, \nt{ADRE} --- адреналин, \nt{ASPA} --- аспартат (таблица~\ref{tab:types}).

Букв меньше, чем величин, и в разных главах одна буква иногда означает разное. В таблице ниже у каждого значения своя строка; в правом столбце --- формула, где оно введено.

\begin{center}
\footnotesize
\setlength{\tabcolsep}{4pt}
\renewcommand{\arraystretch}{1.12}
\begin{longtable}{>{\raggedright}p{2.0cm}>{\raggedright}p{9.6cm}>{\raggedright\arraybackslash}p{2.4cm}}
\toprule
Символ & Значение & Где введено \\
\midrule
\endfirsthead
\toprule
Символ & Значение & Где введено \\
\midrule
\endhead
\bottomrule
\endfoot
$a$ & крутизна линейной передаточной характеристики, $f(u)=au$ & \eqref{eq:feedback} \\
$a$ & уровень \nt{ACET}: $0$ --- чтение, $1$ --- запись & \eqref{eq:acetnet} \\
$a_i$, $a$ & усталость группы: в генераторе ритма; в качелях медленного сна & \eqref{eq:halfcentre}, \eqref{eq:updown} \\
$a_S$, $a_P$ & чувствительность ритма сердца к газу и к тормозу & \eqref{eq:gasbrake} \\
$A(r)$, $A_0$ & ответ синапса на импульс при частоте $r$; у отдохнувшего синапса & \eqref{eq:depression} \\
$A_1$ & ответ получателя на одну порцию медиатора & \eqref{eq:quantal} \\
$A$ & площадь мембраны нейрона вместе с дендритами & \eqref{eq:dendriteload} \\
$A_f$, $A_s$ & доля поправки, сохраняемая до следующего движения, у быстрого и медленного процессов & \eqref{eq:tworate} \\
$b_i$ & собственный приток ритмоводителя & \eqref{eq:seroneuron} \\
$b$ & как быстро работа утомляет группу в генераторе ритма или в качелях сна & \eqref{eq:halfcentre}, \eqref{eq:updown} \\
$b$ & число бит на один отсчёт оцифровки & \eqref{eq:probedata} \\
$B_f$, $B_s$ & сила, с которой поправка учится на ошибке, у быстрого и медленного процессов & \eqref{eq:tworate} \\
$c$ & концентрация медиатора в объёме & \eqref{eq:seroneuron} \\
$c$ & коэффициент корреляции шумов соседних нейронов & \eqref{eq:correlated} \\
$c$ & коэффициент, с которым разность счётов двигает ракетку & \eqref{eq:dishreadout} \\
$c_f$ & ответ $C$-нейрона: наибольший из $s_{f,p}$ по всем положениям & \eqref{eq:scstage} \\
$c_m$ & удельная ёмкость мембраны & \eqref{eq:dendriteload} \\
$C$ & цена усилия: действие за время $\tau_a$ стоит $C/\tau_a$ & \eqref{eq:vigor} \\
$d'$ & различимость двух входов: разница, делённая на шум & \eqref{eq:gainsnr} \\
$d$, $d_j$ & задержка на пути сигнала & \eqref{eq:glutneuron}, \eqref{eq:vstar} \\
$d$ & задержка петли обратной связи с мышцей & \eqref{eq:delaystable} \\
$D$ & коэффициент диффузии & \eqref{eq:diffusionlength} \\
$D$ & знаменатель в частотах \nt{GLUT}--\nt{GABA}-сети & \eqref{eq:isn} \\
$D$ & время ожидания отложенной награды & \eqref{eq:patience} \\
$e$, $e_{ij}$ & след в синапсе & \eqref{eq:threefactor} \\
$e$, $e_i$ & ошибка выхода или движения, приходящая по проводу ошибки & \eqref{eq:lms}, \eqref{eq:adaptivefilter}, \eqref{eq:tworate} \\
$E_{\text{имп}}$ & энергия одного импульса вместе с работой синапсов & \eqref{eq:energy} \\
$E$ & частота \nt{GLUT}-группы & \eqref{eq:ei} \\
$E_E$ & равновесный потенциал возбуждающего синапса & \eqref{eq:shunt} \\
$f$, $F$ & передаточная характеристика: частота как функция входа & \eqref{eq:ratenet}, \eqref{eq:gain} \\
$f$, $f_0$ & частота ударов сердца; собственный ритм ритмоводителя & \eqref{eq:gasbrake} \\
$f_s$ & частота оцифровки одного электрода & \eqref{eq:probedata} \\
$F(t)$, $F_1$ & сила мышцы; сила одного подёргивания & \eqref{eq:forcesum} \\
$F_1$, $F_2$ & силы двух мышц, тянущих сустав в разные стороны & \eqref{eq:antagonists} \\
$g$ & чувствительность рецептора к концентрации & \eqref{eq:seroneuron} \\
$g_L$, $g_E$, $g_I$ & проводимости утечки, возбуждения и торможения & \eqref{eq:shunt} \\
$g$ & коэффициент усиления, задаваемый \nt{NORA} & \eqref{eq:gain}, \eqref{eq:machine} \\
$g$ & общее торможение от \nt{GABA} в сети памяти & \eqref{eq:acetnet} \\
$g$ & усиление ворот боли, задаваемое сверху & \eqref{eq:gate} \\
$g$ & усиление одной ступени гормонального каскада & \eqref{eq:cascade} \\
$g_j$ & фильтр $j$ со своей постоянной времени на входе адаптивного фильтра & \eqref{eq:adaptivefilter} \\
$G$ & усиление петли обратной связи & \eqref{eq:feedback} \\
$G$ & скорость исправления в петле с задержкой & \eqref{eq:delaystable} \\
$h$, $h_I$ & внешний вход \nt{GLUT}-группы и \nt{GABA}-группы & \eqref{eq:ei}, \eqref{eq:isn} \\
$h(t)$ & одиночное подёргивание мышцы & \eqref{eq:twitch} \\
$H(t)$ & верхний порог давления сна: дойдя до него, машина засыпает & \eqref{eq:sleeppressure} \\
$I$, $I_i$ & внешний приток к нейрону или группе & \eqref{eq:ratenet}, \eqref{eq:wta}, \eqref{eq:updown} \\
$I$ & частота \nt{GABA}-группы & \eqref{eq:ei} \\
$I$ & яркость или интенсивность раздражителя & \eqref{eq:photonnoise}, \eqref{eq:nakarushton} \\
$I$ & входной ток, заряжающий мембрану & \eqref{eq:dendriteload} \\
$k$ & число импульсов в пачке & \eqref{eq:burst} \\
$k$ & во сколько раз прибавка должна превышать шум & \eqref{eq:photonnoise} \\
$k$ & коэффициент обратной связи гормонального каскада & \eqref{eq:cascade} \\
$k_s$ & сила, с которой боль накачивает сенсибилизацию & \eqref{eq:sensitization} \\
$K$ & число признаков результата & \eqref{eq:vectortd} \\
$K$ & жёсткость сустава & \eqref{eq:antagonists} \\
$K$ & усиление петли гормонального каскада, $K=g^3k$ & \eqref{eq:cascade}, \eqref{eq:cascadestable} \\
$K_\varphi$ & сила подстройки суточного генератора к внешнему сигналу & \eqref{eq:pll} \\
$L$ & длина линии задержек & \eqref{eq:itd} \\
$L$ & длина провода от датчика боли & \eqref{eq:paintime} \\
$L(t)$ & нижний порог давления сна: дойдя до него, машина просыпается & \eqref{eq:sleeppressure} \\
$M$ & мера причинной связи предвестника и награды & \eqref{eq:retro} \\
$M_k$ & ожидание признака $k$ & \eqref{eq:vectortd} \\
$M_{ik}$ & сила торможения от \nt{GABA}-нейрона $k$ & \eqref{eq:machine} \\
$M$ & вращающий момент сустава & \eqref{eq:antagonists} \\
$M$ & число входных линий смесителей & \eqref{eq:cover} \\
$n$, $\bar n$ & число импульсов за окно; его среднее & \eqref{eq:rateerror} \\
$n$ & число входов порогового элемента & \eqref{eq:cover} \\
$n_\uparrow$, $n_\downarrow$ & счёт импульсов областей «вверх» и «вниз» за окно & \eqref{eq:dishreadout} \\
$N$ & число нейронов & \eqref{eq:correlated}, \eqref{eq:energy} \\
$N$ & число электродов щупа & \eqref{eq:probedata} \\
$N_s$ & число мест выброса между двумя нейронами & \eqref{eq:quantal} \\
$p$ & вероятность выброса медиатора на импульс & \eqref{eq:burst} \\
$p$ & доля активных нейронов в сети памяти & \eqref{eq:acetnet} \\
$p$ & возмущение, которое учатся компенсировать & \eqref{eq:tworate} \\
$P$ & мощность, затрачиваемая на сигналы & \eqref{eq:energy}, \eqref{eq:dishpower} \\
$P$ & неопределённость хранимой оценки & \eqref{eq:kalman} \\
$P$ & активность «тормоза»: \nt{ACET}-провода к сердцу & \eqref{eq:gasbrake} \\
$P_{\max}$ & сколько случайных образов пороговый элемент может разделить на два класса & \eqref{eq:cover} \\
$P_{\text{пр}}$ & вероятность промаха & \eqref{eq:dishdensity} \\
$q$ & скорость, с которой меняется мир & \eqref{eq:kalman} \\
$q_k$ & частота \nt{GABA}-нейрона $k$ & \eqref{eq:machine} \\
$q_0$ & фоновая активность тормозящего посредника & \eqref{eq:gate} \\
$q$ & частота тормозящего посредника в воротах боли & \eqref{eq:gate} \\
$q$ & во сколько раз следующая двигательная единица сильнее предыдущей & \eqref{eq:sizeprinciple} \\
$r$, $r_i$ & частота импульсов нейрона & \eqref{eq:ratenet} \\
$r$, $r_t$ & награда (поток награды) & \eqref{eq:rpe}, \eqref{eq:ctd} \\
$\mathbf r(t)$ & вектор откликов резервуара & \eqref{eq:reservoir} \\
$R$, $\bar R$ & полученная награда; её ожидание или среднее в единицу времени & \eqref{eq:perturbation}, \eqref{eq:vigor} \\
$R$ & дисперсия шума на входе фильтра & \eqref{eq:kalman} \\
$R$, $r_0$ & отложенная и немедленная награды & \eqref{eq:patience} \\
$R_{\max}$ & предельная скорость передачи, бит/с & \eqref{eq:spikeentropy} \\
$r_{\max}$ & наибольшая частота импульсов & \eqref{eq:gain}, \eqref{eq:nakarushton} \\
$R_{\text{сырой}}$, $R_{\text{имп}}$ & поток данных щупа: сырой и только импульсы, бит/с & \eqref{eq:probedata} \\
$s_j$ & знак выходов нейрона $j$ & \eqref{eq:dalesign} \\
$s_i$ & знак рецептора получателя & \eqref{eq:seroneuron} \\
$s$, $\hat s$ & состояние мира; его оценка & \eqref{eq:rpe}, \eqref{eq:kalman} \\
$s_{f,p}$ & ответ $S$-нейрона на признак $f$ в положении $p$ & \eqref{eq:scstage} \\
$S$ & активность «газа»: \nt{NORA}-провода к сердцу & \eqref{eq:gasbrake} \\
$S$ & давление сна & \eqref{eq:sleeppressure} \\
$t_1$ & момент первого импульса & \eqref{eq:latency} \\
$t_k$ & момент $k$-го импульса & \eqref{eq:forcesum} \\
$t_{f,i}$ & образец признака $f$ & \eqref{eq:scstage} \\
$T$ & окно счёта импульсов; время накопления фотонов & \eqref{eq:rateerror}, \eqref{eq:photonnoise} \\
$T_c$ & время нарастания подёргивания мышцы & \eqref{eq:twitch} \\
$T_{\text{залп}}$ & интервал между залпами культуры & \eqref{eq:burstinterval} \\
$u$ & суммарный вход нейрона & \eqref{eq:gain} \\
$u$, $u(t)$ & команда мозга: на входе адаптивного фильтра; гормональному каскаду & \eqref{eq:adaptivefilter}, \eqref{eq:cascade} \\
$u(t)$ & вход резервуара & \eqref{eq:reservoir} \\
$U$ & уровень постоянного входа & \eqref{eq:latency} \\
$U$ & доля запаса медиатора, выбрасываемая на импульс & \eqref{eq:depression} \\
$v$ & скорость сигнала в проводе; скорость движения пятна & \eqref{eq:itd}, \eqref{eq:vstar} \\
$V$ & напряжение на мембране & \eqref{eq:glutneuron}, \eqref{eq:shunt} \\
$V$ & ожидание будущей награды & \eqref{eq:rpe}, \eqref{eq:ctd} \\
$w$, $w_{ij}$, $W_{ij}$ & вес связи & \eqref{eq:glutneuron}, \eqref{eq:ratenet}, \eqref{eq:acetnet}, \eqref{eq:lms}, \eqref{eq:adaptivefilter} \\
$\mathbf w_k$ & веса входов $C$-нейрона $k$ & \eqref{eq:traceinvariance} \\
$w_{q\mu}$, $w_{q\nu}$, $w_{z\mu}$ & веса связей в воротах боли & \eqref{eq:gate} \\
$W_{\text{вых}}$ & веса линейного считывания резервуара & \eqref{eq:reservoir} \\
$x$, $y$ & логические входы и выход & \eqref{eq:dualrail}, \eqref{eq:veto} \\
$x$, $y$ & активность отправителя и получателя & \eqref{eq:hebb} \\
$x$ & место детектора на линии задержек & \eqref{eq:itd} \\
$x_i$ & вход от органов чувств & \eqref{eq:acetnet} \\
$x_p$ & вход в положении $p$ & \eqref{eq:scstage} \\
$\mathbf x$ & выходы $S$-нейронов, вход $C$-нейрона & \eqref{eq:traceinvariance} \\
$x_j$ & вход $j$ адаптивного фильтра: команда после фильтра $g_j$ & \eqref{eq:lms}, \eqref{eq:adaptivefilter} \\
$x_f$, $x_s$ & поправки быстрого и медленного процессов & \eqref{eq:tworate} \\
$x_1$, $x_2$, $x_3$ & уровни на трёх ступенях гормонального каскада; $x_3$ --- конечный гормон & \eqref{eq:cascade} \\
$x^*$ & доля восстановленного запаса медиатора, при которой начинается новая лавина & \eqref{eq:burstinterval} \\
$y$ & зашумлённый вход фильтра & \eqref{eq:kalman} \\
$y$, $\hat y$ & сигнал датчиков; его предсказание адаптивным фильтром & \eqref{eq:adaptivefilter} \\
$y_k$, $\bar y_k$ & активность $C$-нейрона $k$; её след & \eqref{eq:traceinvariance} \\
$y(t)$ & выход считывания резервуара & \eqref{eq:reservoir} \\
$z$ & частота выхода ворот боли & \eqref{eq:gate} \\
\midrule
$\alpha_i^\pm$ & вес положительных и отрицательных ошибок & \eqref{eq:asymmetric} \\
$\beta$ & сила взаимного торможения & \eqref{eq:wta} \\
$\beta$ & сила торможения выхода ворот боли & \eqref{eq:gate} \\
$\gamma$ & коэффициент дисконтирования & \eqref{eq:rpe} \\
$\delta(\cdot)$ & дельта-функция: мгновенный скачок & \eqref{eq:glutneuron} \\
$\delta$, $\delta_t$ & ошибка предсказания награды & \eqref{eq:rpe}, \eqref{eq:ctd} \\
$\delta t$ & разница времён прихода звука в два уха & \eqref{eq:itd} \\
$\Delta$, $\Delta_{\max}$ & интервал между импульсами; окно совпадения & \eqref{eq:window} \\
$\Delta t$ & точность, с которой получатель различает время & \eqref{eq:spikeentropy} \\
$\Delta F$ & прибавка силы от очередной двигательной единицы & \eqref{eq:sizeprinciple} \\
$\Delta I$ & различимая прибавка яркости & \eqref{eq:photonnoise} \\
$\Delta p$ & сдвиг ракетки за окно & \eqref{eq:dishreadout} \\
$\Delta u$ & разница входов двух групп & \eqref{eq:gainsnr} \\
$\Delta V$ & на сколько нужно поднять напряжение мембраны & \eqref{eq:dendriteload} \\
$\Delta x$ & расстояние между соседними датчиками & \eqref{eq:vstar} \\
$\varepsilon$ & относительная разница входов & \eqref{eq:wtagain} \\
$\eta$ & скорость обучения & \eqref{eq:plasticity}, \eqref{eq:lms}, \eqref{eq:traceinvariance} \\
$\eta$ & шум в сети после усилителя & \eqref{eq:softmax} \\
$\mu$ & вход прикосновения к воротам боли & \eqref{eq:gate} \\
$\nu$ & вход боли & \eqref{eq:gate} \\
$\theta$ & порог срабатывания & \eqref{eq:window}, \eqref{eq:scstage} \\
$\kappa$ & количество медиатора на один импульс & \eqref{eq:seroneuron} \\
$\kappa_i$ & отношение весов положительных и отрицательных ошибок & \eqref{eq:expectile} \\
$\kappa_m$ & связь силы мышцы с её жёсткостью & \eqref{eq:antagonists} \\
$\ell$ & расстояние, на которое расплывается медиатор & \eqref{eq:diffusionlength} \\
$\ell$ & плечо, на котором мышца тянет сустав & \eqref{eq:antagonists} \\
$\xi$ & случайное дрожание выхода нейрона & \eqref{eq:perturbation} \\
$\rho(\boldsymbol\psi)$ & доля времени, проведённого в настройке $\boldsymbol\psi$ & \eqref{eq:dishdensity} \\
$\sigma$ & разброс, масштаб шума & \eqref{eq:rateerror}, \eqref{eq:softmax} \\
$\sigma$ & яркость, при которой ответ равен половине наибольшего & \eqref{eq:nakarushton} \\
$\sigma_{\text{до}}$, $\sigma_{\text{после}}$ & шум до усилителя и после него & \eqref{eq:gainsnr} \\
$\tau$ & постоянная времени мембраны & \eqref{eq:glutneuron} \\
$\tau$ & постоянная времени ступени гормонального каскада & \eqref{eq:cascade} \\
$\tau_{\text{эфф}}$ & постоянная времени группы, возбуждающей сама себя & \eqref{eq:perfectint} \\
$\tau_c$ & время жизни медиатора в объёме & \eqref{eq:seroneuron} \\
$\tau_E$, $\tau_I$ & постоянные времени \nt{GLUT}- и \nt{GABA}-группы & \eqref{eq:ei} \\
$\tau_e$ & время жизни следа & \eqref{eq:threefactor} \\
$\tau_{\text{сл}}$ & время жизни следа активности $C$-нейрона & \eqref{eq:traceinvariance} \\
$\tau_s$ & время возврата чувствительности после повреждения & \eqref{eq:sensitization} \\
$\tau_r$, $\tau_d$ & время накопления давления сна при бодрствовании и его сброса во сне & \eqref{eq:sleeppressure} \\
$\tau_{\text{вос}}$ & время восстановления запаса медиатора & \eqref{eq:depression} \\
$\tau_\gamma$ & горизонт планирования & \eqref{eq:ctd} \\
$\tau_v$ & скорость, с которой уточняется ожидание & \eqref{eq:ramp} \\
$\tau_a$ & время усталости в генераторе ритма или в качелях сна & \eqref{eq:halfcentre}, \eqref{eq:updown} \\
$\tau_a^*$ & лучшее время выполнения действия & \eqref{eq:vigor} \\
$\Phi$ & правая часть правила обучения & \eqref{eq:plasticity} \\
$\Phi$ & поток фотонов & \eqref{eq:photonnoise} \\
$\varphi_k$ & признак $k$ полученного результата & \eqref{eq:vectortd} \\
$\varphi$ & разность фаз суточного генератора и внешнего сигнала & \eqref{eq:pll} \\
$\boldsymbol\psi$ & настройка сети в модели DishBrain & \eqref{eq:dishdensity} \\
$\omega$, $\omega_0$ & частота внешнего сигнала (света); собственная частота суточного генератора & \eqref{eq:pll} \\
\end{longtable}
\end{center}

\chapter{Ответы}
\label{sec:answers}

Краткие ответы к задачам в конце глав. Полные решения --- в отдельном пособии для преподавателей.

\answer{ex:01:1}{(а)~$8\cdot10^{10}$ \nt{ГЛУТ} --- $10$ населений Земли; широковещательных $\approx1.9\cdot10^6$ --- город размером с Вену. (б)~$\approx8\cdot10^{14}$ окончаний, из них широковещательных $\approx3\cdot10^{11}$, доля $\approx4\cdot10^{-4}$ --- в $\approx16$ раз больше доли этих нейронов ($2.2\cdot10^{-5}$).}
\answer{ex:01:2}{(а)~$\tau_c\ln(c_0/c_*)\approx4.6\ms$. (б)~Линия свободна при $T\ge\tau_c\ln101$: до $\approx22$ выбросов в секунду вместо $\approx220$.}
\answer{ex:01:3}{$V(s_k)=\gamma^{K-1-k}r$; $\delta_{\text{лампа}}=\gamma^Kr=re^{-T/\tau_\gamma}\approx0.60\,r$ при любом $\Delta$, $\tau_\gamma\approx1.95\,\text{с}$.}
\answer{ex:01:4}{$n^*=93$, $47$, $25$, $12$: $n^*\approx5/\alpha$ при малых $\alpha$.}
\answer{ex:01:5}{\nt{ГЛУТ}: $1\to10\to10^4\to10^7$ --- $\approx10^7$ нейронов, $4$ звена, $4\ms$. \nt{ДОФА}: $1\to10\to3.2\cdot10^5$ --- $\approx3\cdot10^5$ нейронов, $3$ звена, в $30$ раз меньше; тот же порядок, что число \nt{ДОФА}-нейронов.}
\answer{ex:01:6}{$\lambda=f_EJ_E-f_IJ_I$, откуда $J_I=37.5$; отношение токов $f_IJ_I/f_EJ_E=0.94$; лавина ($\lambda\ge1$) при ослаблении торможения всего на $6.7\,\%$.}
\answer{ex:01:7}{$n\ge57$ предъявлений на каждую задержку в каждом состоянии. Тот же спад дала бы, например, неточность внутренней оценки времени, растущая с $T$.}

\answer{ex:02:1}{(а)~$\delta t_{\max}=0.64\ms$, $L=v\,\delta t_{\max}=3.2$~мм, шаг $v\cdot10\,\text{мкс}/2=25\um$, $\approx129$ детекторов. (б)~Срабатывает полоса шириной $v\,\tau\ln2=1.7$~мм, $\approx69$ детекторов; направление --- центр полосы.}
\answer{ex:02:2}{(а)~$-\tau\ln\frac{w_2}{\theta-w_1}\le\delta\le\tau\ln\frac{w_1}{\theta-w_2}$. (б)~$c=\frac\tau2\ln\frac{w_1(\theta-w_1)}{w_2(\theta-w_2)}$; окно $[-0.235,\,0.178]\ms$, $c=-28$~мкс: раньше должен прийти слабый вход. (в)~Направление смещается на $28$~мкс в сторону сильного уха; $c=0$ только при $\theta=w_1+w_2$, где окно нулевой ширины.}
\answer{ex:02:3}{(а)~$\sum_i\max(-z_i,0)$ удовлетворяет неравенству Грёнвалля с нулевым начальным значением и тождественно равно нулю. (в)~Нужны $s_i$ с $s_iw_{ij}s_j\ge0$; они существуют тогда и только тогда, когда циклы чётны. Взаимное торможение --- да (монотонна, устойчивых колебаний нет), петля \nt{ГЛУТ}--\nt{ГАМК} --- нет (может колебаться).}
\answer{ex:02:4}{(б)~Шесть нейронов $g_k=[x+y+z\ge k]$ и $\bar g_k=[\bar x+\bar y+\bar z\ge4-k]$, $k=1,2,3$; $C=g_2$, $\bar C=\bar g_2$, $S=[g_1+\bar g_2+g_3\ge2]$, $\bar S=[\bar g_1+g_2+\bar g_3\ge2]$: восемь нейронов. (г)~$12$ нейронов.}
\answer{ex:02:5}{$r_0=0.00197$; $T_{\min}(0.6)=0.718\,\tau$; граница $\tau\ln[2(1-r_0)]=0.691\,\tau$; $I_c=0.225$, $T_{\min}\propto(I-I_c)^{-1/2}$.}
\answer{ex:02:6}{(а)~$500$ звеньев, $5\cdot10^4$ нейронов и импульсов. (б)~$4.5\ms$ ($0.45\,\%$) за $1\,\text{с}$, $1.4\ms$ ($1.4\,\%$) за $100\ms$: относительная ошибка $\propto T^{-1/2}$. (в)~Общий для всех звеньев случайный множитель задержки с разбросом $5\,\%$.}
\answer{ex:02:7}{(а)~Освежать раз в $\tau_F\ln2=0.69\,\text{с}$: $4.3$ импульса в секунду на нейрон против $20$; $4.3\cdot10^3$ импульсов ($4\cdot10^{-7}$~Дж) против $2\cdot10^4$ ($2\cdot10^{-6}$~Дж). (б)~Триггер теряет бит; конденсатор сохраняет его, если глушение началось не позже чем через $0.49\,\text{с}$ после освежения (вероятность $0.71$), а наверняка --- при освежении раз в $0.49\,\text{с}$, $6.1$ импульса в секунду.}

\answer{ex:03:1}{(а)~$35\to17.5\mV$ при $g_E=g_L$; $56\to40\mV$ при $g_E=4g_L$ (вычитание дало бы $38.5\mV$); шунт равносилен делению $g_E$ на $3$. (б)~$\tau_{\text{эфф}}=10\to5\ms$, окно $6.9\to3.5\ms$. (в)~$v^*=2$~мм/с.}
\answer{ex:03:2}{(а)~$\operatorname{tr}=\frac{w_{EE}-1}{\tau_E}-\frac1{\tau_I}$, $\det=\frac{1-w_{EE}+w_{EI}w_{IE}}{\tau_E\tau_I}$. (б)~$\omega=\sqrt{\det}$; граница $\tau_I=20\ms$, период $73\ms$; $\lambda=(-0.225\pm0.156i)\,\text{мс}^{-1}$ при $\tau_I=2\ms$ и $(0.0083\pm0.070i)\,\text{мс}^{-1}$ при $30\ms$ (период $89\ms$), рост ограничивает $[\cdot]_+$. (в)~$-1/3$; при $w_{EE}>1$.}
\answer{ex:03:3}{(б)~$\omega=\sqrt{G^2-1}/\tau_E$, $d_c=\tau_E\arccos(-1/G)/\sqrt{G^2-1}$, период $2\pi\tau_E/\sqrt{G^2-1}\in(2d_c,4d_c)$. (в)~$G=2$: $d_c=12.1\ms$, период $36\ms$; $G=5$: $d_c=3.6\ms$, период $12.8\ms$.}
\answer{ex:03:4}{(а)~$0.73$ и $0.53$; $1$ и $0$; после перестановки $0.9$ и $0$. (б)~Переключения при $I_2=\beta I_1=2$ и $I_2=I_1/\beta=0.5$, ширина петли $1.5$. (в)~$+\tau\ln10/(\beta-1)=2.3\tau$ на каждый десятичный порядок $\varepsilon$.}
\answer{ex:03:5}{(а)~Три посредника, $d_k=5$, $10$, $15\ms$; $k\,T_I\ge15\ms$. (в)~Вне полосы детектор отвечает в обоих направлениях; при $T_I=10\ms$ хватает двух.}
\answer{ex:03:6}{(а)~$n+M+1$ нейронов ($M$ --- число наборов с $f=1$); для $x\oplus y\oplus z$ --- $8$. (б)~\nt{ГАМК} $=[x+y+z\ge2]$, выход $=[x+y+z-2\,\nt{ГАМК}\ge1]$. (в)~Прямое возбуждение надо задержать на время пути через \nt{ГАМК}; при неизвестном разбросе моментов --- нельзя.}
\answer{ex:03:7}{(а)~Один посредник, возбуждаемый всеми и тормозящий всех, плюс самовозбуждение $P=\beta I$. (б)~$\binom84=70<100\le\binom94=126$: $9$ посредников. (в)~$100$: ранг $\beta(J-I)$ равен $n$.}
\answer{ex:03:8}{(а)~$h_c=\frac1{2w_{EE}}+\frac{w_{EI}w_{IE}-w_{EE}}{4w_{EE}^2}=\frac7{18}\approx0.39$; условие~(ii) нарушается при $h=4$. (б)~Смена знака между $h=0.38$ и $0.40$; при резком включении $h\gtrsim2.4$ из тишины активность уходит в разнос.}

\answer{ex:04:1}{(а)~$\approx1.1\cdot10^{11}$ бит/Дж. (б)~$\log_2\sqrt{10}\approx1.7$ бита против $81$: в $\approx50$ раз.}
\answer{ex:04:2}{$r^*=(P_0/E_{\text{имп}})\ln\bigl(1/(r^*\Delta t)\bigr)$, $P_0/E_{\text{имп}}=1.16\,\text{с}^{-1}$: $r^*\approx6$ импульсов в секунду, $7.4\cdot10^{10}$ бит/Дж.}
\answer{ex:04:3}{$\mathcal I=\frac{N}{1+(N-1)c}=9.9$ (предел $10$) против $\mathcal I=\frac{N}{1-c}=1111$: в $\approx110$ раз; корреляция вредна, только когда похожа на сигнал.}
\answer{ex:04:4}{$\theta\approx68.7$ и $19.9$ (в единицах $w$), $m=19$ и $m=10$: быстрому получателю нужно почти вдвое меньше совпадающих входов, хотя его средний вход в пять раз меньше.}
\answer{ex:04:5}{$U\tau_{\text{вос}}\ge0.725\,\text{с}$ и $\tau_{\text{вос}}(1-2U)\le0.05\,\text{с}$, что требует $U\ge0.483$; наименьшее $\tau_{\text{вос}}=0.725\,\text{с}$ при $U=1$ ($1.45\,\text{с}$ при $U=0.5$).}
\answer{ex:04:6}{(а)~$\theta\,e/(e-1)\approx1.58\,\theta$ при любых $\tau$ и $\theta$. (б)~$14$ импульсов.}
\answer{ex:04:7}{(а)~$1.62$ бита на слово: $0.77$ --- число импульсов, $0.84$ --- их расположение. (б)~$1.13$ и $1.59$ бита: по $30$ словам оценка теряет почти треть.}

\answer{ex:05:1}{$\ell\approx17\um$; $\approx100$~суток: рассылку по мозгу делает ветвление провода, диффузия лишь размывает каждое место выброса на $\sim20\um$.}
\answer{ex:05:2}{$r=ab/(1+G)$, $S=1/(1+G)$ точно; $+1.7\,\%$ вместо $+20\,\%$.}
\answer{ex:05:3}{$T^*=1.21\,\text{с}$ при $G=2$ и $0.19\,\text{с}$ при $G=9$: реально $G\sim2$--$4$, подавление в $3$--$5$ раз.}
\answer{ex:05:4}{$\mathrm{CV}=1/\sqrt{2\lambda\tau_c}\approx9\cdot10^{-4}$ при суммарной частоте $\lambda$; одному нейрону нужно $\tau_c=12.5\,\text{с}$.}
\answer{ex:05:5}{$c\approx0.40$ в обоих случаях; $\mathrm{CV}=0.050$ против $0.158$, в $\sqrt{10}$ раз. Частоты отдельных нейронов остаются пропорциональны $a_i$: регулируется только сумма.}
\answer{ex:05:6}{$N_1=\overline{A\vee B}$, $N_2=\overline{A\vee N_1}$, $N_3=\overline{B\vee N_1}$, $N_4=\overline{N_2\vee N_3}$, XOR $=N_5=\overline{N_4}$; каждое переключение занимает $\tau_c\ln2$, путь из $4$ вентилей --- $2.8\,\text{с}$ на один XOR.}
\answer{ex:05:7}{(а)~$\tau_\gamma=6/\ln3\approx5.5\,\text{с}$. (б)~$\bar D<4\,\text{с}$ против $D<2.2\,\text{с}$: непредсказуемость задержки делает терпеливее.}
\answer{ex:05:8}{$k_2=1/\tau_d=10\,\text{с}^{-1}$, $k_1=1/\tau_d-1/\tau_\gamma=9.5\,\text{с}^{-1}$; $\tau_\gamma=1/(k_2-mk_1)$: удвоение при $+2.6\,\%$ (а при $+5.3\,\%$ горизонт бесконечен).}

\answer{ex:06:1}{$e=(1-e^{-T_a/\tau_e})e^{-d/\tau_e}$. (а)~$0.110$ против $4.5\cdot10^{-5}$, в $\approx2.5\cdot10^3$ раз. (б)~$\tau_e^*=T_a/\ln(1+T_a/d)=0.59\,\text{с}$.}
\answer{ex:06:2}{Скорость $\nu_x\nu_y(A_+\tau_+-A_-\tau_-)=0.8A_+$ в секунду, $\approx2900A_+$ за час; $\tau_-=\tau_+/0.6=33\ms$.}
\answer{ex:06:3}{$\mathbf e_1$ при $\mu^2+s_1^2>s_2^2$; здесь $1.01>0.25$, и нейрон выучит $\mathbf e_1$, хотя разброс вдоль $\mathbf e_2$ в $25$ раз больше: правило находит главный вектор матрицы корреляций, а не ковариаций.}
\answer{ex:06:4}{$V=Re^{-(t_c+L-t)/\tau_\gamma}$ между лампочкой и соком, так что $\dot V=V/\tau_\gamma$; всплеск площади $Re^{-L/\tau_\gamma}$: $0.61R$ и $0.78R$.}
\answer{ex:06:5}{Отношение сигнала к шуму $1/(N+1)$, попыток $\propto N+1$: $5\cdot10^5$ попыток $\approx1$~мес. и $5\cdot10^{11}$ попыток $\approx8\cdot10^4$~лет.}
\answer{ex:06:6}{(а)~$k_2=1/\tau_d$, $k_1=1/\tau_d-1/\tau_\gamma$. (б)~Ложная ошибка $-(V/\tau_\gamma)\,\varepsilon/(1+\varepsilon)$, $\varepsilon=\tau_d/\tau_\gamma$, откуда $\tau_d\le0.105\,\text{с}$.}
\answer{ex:06:7}{$0.46$, $0.90$, $0.99$, $0.95$ и $0.70$ при $d=3\,\text{с}$. Точки ложатся на одну кривую от доли следа $F=(1-e^{-T_{\text{cue}}/\tau_e})e^{-d/\tau_e}$: обучение при $F\gtrsim0.1$, случайный выбор при $F<10^{-2}$.}
\answer{ex:06:8}{$\eta^*=1/\bigl(2\sigma^2(N+2)\bigr)$, за $N+2$ попытки; при $m$ дрожащих из $N$ --- за $N(m+2)/m\ge N+2$: разреженность не помогает.}

\answer{ex:07:1}{$V=\kappa p/\bigl(\kappa p+(1-\kappa)(1-p)\bigr)$: $0.059$, $0.20$, $0.50$; медиана равна $0$, а точка~\eqref{eq:expectile} меняется с $\kappa$ плавно.}
\answer{ex:07:2}{$p=1/3$, $x=0.75$, $\sigma_r=0.35$, $V(0.2)=0.083$.}
\answer{ex:07:3}{$V=\sqrt\kappa/(\sqrt\kappa+\sqrt{1-\kappa})$, от $0.25$ до $0.75$; разброс $\propto\sqrt\eta$; у капель $V=0.25+0.5\kappa$, распределения различаются на $0.05$ у крайних нейронов, нужно $\eta\lesssim0.01$.}
\answer{ex:07:4}{(а)~$J^*=2\sqrt{C\bar R}$. (б)~Переплата $\bigl(k^{1/2}+k^{-1/2}\bigr)/2-1$: $6\,\%$ при $k=2$ и $25\,\%$ при $k=4$ --- хватает точности «в пределах двух раз».}
\answer{ex:07:5}{Всплеск площади $R(e^{-1}-e^{-2})\approx0.23R$ --- как $2.3\,\text{с}$ рампы; если дофамин сообщает $V$, всплеска нет, уровень просто ступенькой вырастает в $e$ раз.}
\answer{ex:07:6}{Событие сдвигает вес $\propto A\tau_f/(\tau_e+\tau_f)$, уровень даёт ошибку $s\tau_f$: $9\,\text{с}\le\tau_f\le17\,\text{с}$.}
\answer{ex:07:7}{Исходно $\Delta P=0.80$, $M=0.90$. (а)~$\Delta P=0.74$ ($-8\,\%$), $M=0.43$ ($-52\,\%$); (б)~$\Delta P=0.40$ ($-50\,\%$), $M=0.80$ ($-11\,\%$): двойная диссоциация.}
\answer{ex:07:8}{$N_{\text{эфф}}\to2+\rho^2N\approx10^6$ попыток: в $10^4$ раз меньше, чем с одним проводом, но утечка в $1\,\%$ всё ещё стоит миллиона попыток.}

\answer{ex:08:1}{(а)~$\approx1.7\cdot10^6$. (б)~$g=\frac{0.9}{\sqrt{0.19}}\,\sigma_{\text{после}}/\sigma_{\text{до}}\approx16.5$: ответ зависит только от отношения шумов.}
\answer{ex:08:2}{(а)~Собственные числа $-(1\pm g\beta)/\tau$; разность затухает за $\tau/(1-g\beta)$: $40\ms$ при $g\beta=0.5$ (разница входов усилена втрое), $200\ms$ при $0.9$ (в $19$ раз). (б)~$0.905$ и $1.105$; между границами побеждает только сильный вход.}
\answer{ex:08:3}{$P(k)=e^{gu_k/\sigma}/\sum_le^{gu_l/\sigma}$, то есть~\eqref{eq:softmax}; $g=10\ln9\approx22$.}
\answer{ex:08:4}{$\bar\Delta=1/3$, $g^*\approx3\ln(1/h)$: $21$, $14$, $9$ против $16$--$23$, $11$ и $8$ в модели.}
\answer{ex:08:5}{$g^*$ от $16$--$23$ при $h=0.001$ до $6$--$8$ при $h=0.2$; оценка $3\ln(1/h)$ верна в пределах полутора раз при $h\le0.05$; при $h\gtrsim\eta/10$ правило $Q\leftarrow Q+\eta(R-Q)$ не успевает усвоить добытое поиском, и $g^*$ застревает около $8$.}
\answer{ex:08:6}{$T_p=\tau\ln9=44\ms$ при $g_{\text{lo}}=0$ (модель $44\ms$) и $70\ms$ при $g_{\text{lo}}=0.25$: пачка в $2$--$3\tau$ при усилении около нуля.}
\answer{ex:08:7}{(а)~Лучшее постоянное $0.925$ ($g=8$), оракул $0.929$ ($16/8$) --- выигрывать почти нечего. (б)~Например, спокойные эпохи с редкими наградами ($p\in[0,0.3]$) и переменчивые с $h=0.1$: $0.850$ против $0.872$; подстройка забирает около четверти при $\eta=0.2$ и почти всё при $\eta=0.05$.}

\answer{ex:09:1}{(а)~$R=4$: $P_\infty=0.2$, ошибка $0.45$, скорость $0.05\,\text{с}^{-1}$, память $20\,\text{с}$. (б)~$q=1$: $1$, $1$, $1\,\text{с}^{-1}$, $1\,\text{с}$. (в)~Память пропорциональна амплитуде шума: минута --- при шуме вшестеро сильнее. (г)~\nt{АЦЕТ} ниже в (а), выше в (б).}
\answer{ex:09:2}{(б)~$\hat s(t)=\frac1t\int_0^ty\,\dd t'$. (в)~$\frac12\sqrt{R/q}$ --- вдвое короче памяти. (г)~$t=\frac12\ln21\,\sqrt{R/q}\approx15\,\text{с}$.}
\answer{ex:09:3}{$\dd e/\dd t=-e/T+w-n/T$, дисперсия $qT/2+R/(2T)$; $T^*=\sqrt{R/q}$, дисперсия $\sqrt{qR}=P_\infty$; рост в $(\kappa+1/\kappa)/2$ раз: $1.25$ и $5.05$.}
\answer{ex:09:4}{(а)~$a_w=1-\theta/(mw)$: $0.22$ при $w=0.045$, $0.15$ при $w=0.041$. (б)~$a_c\approx0.38$ и $0.31$. (в)~$a<1-\theta/((2m-1)w)$: $0.59$ и $0.55$; границу чтения задаёт дописывание.}
\answer{ex:09:5}{Ошибки появляются при $M\approx40$--$60$, при $M=80$ --- $1.9$ лишних нейрона, при $M=100$ --- доля образа $0.86$; разброс $h$ совпадает с $0.92\sqrt{(M-1)p(1-p)/N}$ (при $M=80$: $0.172$ против $0.173$); $M_{\max}\approx80$, $M_{\max}\propto N/p$.}
\answer{ex:09:6}{Например, $\eta(a)=\eta\min\bigl(1,\max(0,(a-0.3)/0.6)\bigr)$. С $\eta a$ после $20$ чтений при $a=0.2$ все три чужих нейрона вошли в образ, при $a=0.3$ доля образа упала до $0.6$; с $\eta(a)$ --- ни одного, запись при $a\ge0.9$ та же (доля $1.00$).}
\answer{ex:09:7}{(а)~$40$ повторов; с $\eta(a)$ --- никогда. (б)~Ночью $a_B$ низкий, $a_S$ высокий: $S$ пишет, слушая $B$ как датчик. (в)~$200$ повторов, $10$ ночей; за день --- не больше $\approx M_{\max}/10$ образов.}

\answer{ex:10:1}{(а)~$\approx11.4$ бит на импульс, $\approx10^{12}$ бит/с. (б)~$\approx40$ бит/с, в $\sim10^{10}$ раз меньше. (в)~$2.5\cdot10^{10}\,\text{с}$ --- около восьмисот лет. (г)~$10$--$100$ мВт, $0.1$--$1\,\%$ энергии сигналов --- светодиод индикатора.}
\answer{ex:10:2}{(а)~(i)~$g\,w_{EI}w_{IE}>g(1-a)w_{EE}-1$; (ii)~$\tau_I<\tau_E/\bigl(g(1-a)w_{EE}-1\bigr)$. (б)~$g=4$, $\approx62$~Гц. (в)~$a>1/3$; при $a=0$ нарастающие колебания $\approx67$~Гц. (г)~Нет: устойчивость задаёт произведение $g(1-a)$.}
\answer{ex:10:3}{(а)~$\sigma^2R'_jx_j$. (б)~$x_j^2\sigma^4\bigl(\sum_{k\ne j}R'^2_k+2R'^2_j\bigr)$. (в)~$1/(N+1)$ на попытку, нужно $n\approx z^2(N+1)\propto N$ попыток. (г)~$0.51$, $0.197$, $0.059$, $0.016$ против $0.5$, $0.2$, $0.059$, $0.015$.}
\answer{ex:10:4}{(а)~$L_v\in[32.9,\,47.9]\ms$, $d=35.4\ms$, $\Delta_{\max}=7.5\ms$. (б)~$\tau_D=\Delta_{\max}/\ln2=10.8\ms$. (в)~$2\Delta_{\max}r_ar_v\approx0.38$ в секунду, пропорционально $r_ar_v$ и разбросу задержек.}
\answer{ex:10:5}{(а)~$\alpha=0.2$, ошибка $0.180$. (б)~$q=0.03$, $0.180$ --- то же. (в)~$\lambda=0.9$, $c=20$: $0.098$, на $45\,\%$ меньше. (г)~$0.025$; мешают запаздывание обнаружения и пропуск малых скачков.}
\answer{ex:10:6}{(а)~$0.2+0.6(1-\alpha)^n$. (б)~$n=\ln\bigl(\ln1.5/(0.6g)\bigr)/\ln(1-\alpha)$: $24$ и $4$. (в)~$11$ и $2$. (г)~$Q_A\to Q_B$ и $P(A)\to\frac12$; узнать о $B$ можно только пробой --- ручка \nt{НОРА}.}
\answer{ex:10:7}{(а)~$S'=S\bigl(1-4\eta\sigma^2+4\eta^2\sigma^4(G+2)\bigr)$, лучшее $\eta\sigma^2=1/(2(G+2))$, $n\approx(G+2)\ln(1/\varepsilon)$. (б)~При $N=128$: $595$, $151$, $12$ попыток против $599$, $157$, $14$. (в)~$\approx10^6$ попыток, около месяца. (г)~Нужно $K\propto N$ --- по сути свой сигнал ошибки у каждого нейрона.}

\answer{ex:11:1}{(а)~$T=k^2/\bigl(\Phi(\Delta I/I)^2\bigr)=9\,\text{с}$. (б)~$45$ датчиков, пятно $\approx7\times7$, разрешение хуже в $\approx7$ раз.}
\answer{ex:11:2}{(а)~$2$--$3.3$ бита на импульс, $22$--$34\,\%$ границы ($9.1$--$9.8$ бита). (б)~$\log_2\binom{400}{20}\approx111$ бит; в среднем $1$ общий тип.}
\answer{ex:11:3}{(а)~$\sigma/9\le I\le9\sigma$, $1.9$ порядка. (б)~$6$ и $7$.}
\answer{ex:11:4}{(а)~$f<343/0.44\approx780$~Гц. (б)~$625$ импульсов, $125$ нейронов.}
\answer{ex:11:5}{$21$ бит на событие, $4.8\cdot10^5$ событий/с, $7$ событий на пиксель края: $136$ пикселей/с. Обычная камера --- $1.25$ кадра/с.}
\answer{ex:11:6}{В логарифмах $\approx1.0\,\text{с}$ в обе стороны (теория $0.96\tau$); линейно $0.2\,\text{с}$ вверх и $3.0\,\text{с}$ вниз (теория $0.18\tau$ и $3.0\tau$).}
\answer{ex:11:7}{(а)~$\ln I$ распределён логистически с медианой $\ln\sigma$ и масштабом $1$: разброс $\pi/\sqrt3$ в натуральных логарифмах, $0.79$ порядка. (б)~$r=r_{\max}\Phi_I(I)^2$.}

\answer{ex:12:1}{$0.60$ (без корреляции было бы $0.18$), предел $\sqrt{c/(rT)}\approx0.58$: сто нейронов различают лишь «есть/нет».}
\answer{ex:12:2}{(а)~$10^8$ импульсов, $\approx10\,\text{мДж}$: $0.3\,\%$ от $3$~Дж всего мозга. (б)~$3$~Дж, в $\approx300$ раз больше.}
\answer{ex:12:3}{(а)~$\theta_A\in[2,3)$, $\theta_B\in[1,2)$; чужая фигура даёт не больше $2$ и $1$. (б)~$10\cdot96^2\approx9.2\cdot10^4$.}
\answer{ex:12:4}{(а)~$T_d=0.85\,\tau_a=0.34\,\text{с}$. (б)~$T_d\propto\tau_a$ и не зависит от $I$: $\tau_a\approx3.5\,\text{с}$ (моделирование: $3.1\,\text{с}$).}
\answer{ex:12:5}{$p\approx0.17$. Доля падает медленно: $0.90$, $0.88$, $0.86$, $0.84$, $0.82$, $0.79$ при $N=8$, $12$, $24$, $48$, $96$, $192$.}
\answer{ex:12:6}{Все десять прогонов безошибочны при $\tau_{\text{сл}}=20$, $50$ и $200\ms$ (при $30\ms$ один прогон из десяти сорвался); при $5$--$10\ms$ --- около $0.5$, при $0.5$--$2\,\text{с}$ качество падает до $0.86$. Сверху: след должен угаснуть за паузу, $e^{-0.3\,\text{с}/\tau_{\text{сл}}}\ll1$.}
\answer{ex:12:7}{Одной задержки мало: окно $\pm5\ms$ не покрывает $\Delta x/v$ от $5$ до $20\ms$. Две ветви торможения с $5<d_1\le10\ms$ и $15\le d_2\le d_1+10\ms$.}
\answer{ex:12:8}{\texttt{two\_stage}: $\Delta=2$, ничья от одного переворота, смена ответа --- от двух. Счётчик единиц: один переворот; $0.57$ при $p=0.1$.}

\answer{ex:13:1}{$q=100^{1/99}\approx1.048$, шаг $\approx4.8\%$; диапазон $\approx2.2\cdot10^3$ ($67$~дБ); $\approx15$ включений на удвоение; у линейного ЦАП шаг $22f_1$, т.\,е. $220\%$ при $10f_1$; точность $10\%$ достигается при $n\ge53$.}
\answer{ex:13:2}{$P_{\text{сбоя}}\approx0.47$, $1.8\cdot10^{-4}$, $2.8\cdot10^{-16}$; при $S=2$ около $300$ сбоев в сутки, при $S=4$ около $5\cdot10^{-10}$ в сутки; оценка Чернова даёт $\le10^{-14}$.}
\answer{ex:13:3}{$H(s)=eF_1T_c/(1+sT_c)^2$, частота среза $1/(2\pi T_c)\approx3.2$~Гц; $\bar F=e\,rT_cF_1$ ($2.0F_1$ и $5.4F_1$); при $5$~имп/с $F$ от $0.21F_1$ до $1.10F_1$; размах $10\%$ при $r\approx22$~имп/с (по первой гармонике $\approx20$).}
\answer{ex:13:4}{$Gd<\pi/2$, $f=1/(4d)$; быстрее всего без колебаний при $Gd=1/e$, скорость затухания $1/d$; $d=30\ms$: $G\approx12$~с$^{-1}$, $\tau=30\ms$; $d=150\ms$: $G\approx2.5$~с$^{-1}$, $\tau=150\ms$.}
\answer{ex:13:5}{$E\approx0.427$, $T\approx1.70\,\tau_a=0.68$~с (модель с $\tau=20\ms$ даёт $0.84$~с); ритм существует при $b>\beta-1$; $T$ не зависит от $I$.}
\answer{ex:13:6}{$T\approx0.48$, $0.84$, $1.55$~с при $\tau_a=0.2$, $0.4$, $0.8$~с ($0.39$, $0.73$, $1.42$~с при $\tau=5\ms$); при любом $I$ $T=0.840$~с; при $b=0.8$ ритма нет, при $b=1.05$, $1.2$, $1.5$, $2$, $4$ $T\approx2.73$, $1.74$, $1.19$, $0.84$, $0.45$~с.}
\answer{ex:13:7}{$a\ge33.3$~с$^{-1}$, $F_1\approx58$~Н, $F_2\approx8.3$~Н, $G=e^{-2}/d\approx4.5$~с$^{-1}$; без мышцы-противника $\tau\approx17\ms$, требование не выполнено; предел $\tau\approx9.2\ms$; чем сильнее совместное напряжение, тем меньше надо брать $G$.}
\answer{ex:13:8}{Открытая задача. Например, порог в накачке усталости $b[r_i-r_0]_+$ даёт $I_{\min}=\beta b r_0/(b-\beta+1)\approx0.53$ при $b=4$, $r_0=0.2$ и период от $1.8$~с при $I=0.6$ до $0.52$~с при $I=3$.}

\answer{ex:14:1}{(а)~На $0.17$~с и на $1.5$~с. (б)~С расстояния меньше $11$~см.}
\answer{ex:14:2}{Касание без боли не проходит ни при каком $\mu$, пока $g\le\beta w_{q\mu}/w_{z\mu}=5$; касание $\mu=1$ проходит при $g>6$: сильная сенсибилизация делает болью обычное прикосновение.}
\answer{ex:14:3}{(а)~$g^*=(1-k_sC)/(1-k_sA)$, устойчиво при $k_sA<1$. (б)~Без касания при $\nu\ge2$, с касанием --- уже при $\nu\ge1.7$.}
\answer{ex:14:4}{$V(t)=-Pe^{-(T-t)/\tau_\gamma}$. (а)~$-Pe^{-T/\tau_\gamma}\approx-0.37P$. (б)~$+Pe^{-(T-t_e)/\tau_\gamma}\approx+0.61P$; он растёт с $t_e$ до $+P$ и учит избегать боли.}
\answer{ex:14:5}{При $k_s=4$ защёлкивания нет (второе устойчивое состояние есть при $k_s\ge4.58$); $T_{\min}\approx4.35$, $1.40$, $0.65$, $0.32\,\tau_s$ при $k_s=4.6$, $5$, $6$, $8$.}
\answer{ex:14:6}{$w_{q\nu}=4/9\approx0.444$, $q_0=2/9\approx0.222$, $w_{q\mu}=49/45\approx1.089$; касание без боли безопасно до $g=49/9\approx5.4$.}
\answer{ex:14:7}{(а)~Порог $\nu_c(g)=\theta/g$ при $\nu_c\ge q_0/w_{q\nu}$, иначе $(\theta+\beta q_0)/(g+\beta w_{q\nu})$; $g^*\approx2.45$, $1.0$, $0.25$. (б)~Открытый вопрос.}

\answer{ex:15:1}{(а)~$2\cdot10^5$ на нейрон, $6\cdot10^{12}$ на схему. (б)~$\approx8$~бит/с на провод, не меньше $2.5\cdot10^4$~с $\approx7$~ч.}
\answer{ex:15:2}{Шаг $2$: соседняя корреляция $0.943$, $\lambda$ от $4.18$ до $7.3\cdot10^{-4}$, отношение $\approx5.7\cdot10^3$; шаг $4$: корреляция $0.8$, отношение $\approx94$.}
\answer{ex:15:3}{(а)~$\lambda=0.988$ и $0.808$: $82$ и $4.7$ движения. (б)~$x_s^*=0.690$, $x_f^*=0.172$, сумма $0.862$.}
\answer{ex:15:4}{(а)~$116$ движений. (б)~$19$. Модель с одним процессом при нулевой сумме экономии не даёт.}
\answer{ex:15:5}{(а)~$G=50$~с$^{-1}$ против предела $10.5$~с$^{-1}$ без модели. (б)~Нет: при $G=50$~с$^{-1}$ критическая задержка $d_c\approx103\ms$; при $d=150\ms$ нужно $\hat k>0.445k$ (при любых $G$ и $d$ --- $\hat k>k/2$).}
\answer{ex:15:6}{Ошибка падает ниже $0.5\%$ за $6$--$30$~с, пол $\approx(6$--$9)\cdot10^{-4}$ при наилучших весах $7.5\cdot10^{-4}$; веса при $\eta=50$ ещё отличаются от наилучших на величину до $0.2$ вдоль плохо обусловленных направлений $C$ (задача~\ref{ex:15:2}).}
\answer{ex:15:7}{$\mathbf B=A\mathbf K$; $q_f/R\approx0.035$, $q_s/R\approx8.6\cdot10^{-4}$; при $4q_s$ $B_f\approx0.088$, $B_s\approx0.047$: медленный процесс учится более чем вдвое быстрее, а быстрый --- медленнее.}

\answer{ex:16:1}{Кровь --- RC-цепь с $\tau=t_{1/2}/\ln2=14.4$~мин; максимум к минимуму $2^{T/t_{1/2}}=64$; основная гармоника ослаблена до $0.55$; отношение $2$ при $t_{1/2}=T=60$~мин.}
\answer{ex:16:2}{$\varphi^*=\arcsin\bigl((\omega-\omega_0)/K_\varphi\bigr)\approx41$~мин ($\omega-\omega_0\approx0.047$~рад/сут), время релаксации $1/(K_\varphi\cos\varphi^*)\approx3.9$~сут; после сдвига на $+6$ и $-6$~ч $7.7$ и $7.9$~сут.}
\answer{ex:16:3}{$K_c(n)=\sec^n(\pi/n)$: $8$ и $2.37$, ошибка $11\%$ и $30\%$; при $n\to\infty$ $K_c\to1$, ошибка $\to50\%$.}
\answer{ex:16:4}{Крайние команды (газ $80$, тормоз $0$) дают $f=120$ через $0.76$~с; «только газ» --- $2.33$~с, втрое дольше.}
\answer{ex:16:5}{$3\arctan(\omega\tau)+\omega d=\pi$, $K_c=(1+\omega^2\tau^2)^{3/2}$: $K_c=8$, $3.54$, $2.50$, $1.76$ и период $3.63$, $5.46$, $6.86$, $9.27\,\tau$; при $0.9K_c$ колебания затухают, при $1.1K_c$ устанавливаются.}
\answer{ex:16:6}{$0.60\bar c$, $0.87\bar c$, $0.90\bar c$, предел $0.91\bar c$; при постоянной концентрации отклик нулевой: такой получатель читает только порции.}
\answer{ex:16:7}{Открытая задача. При обратной связи ритм есть только при $K>8$, период пропорционален $\tau$ ступеней и почти не зависит от $K$ ($3.64\,\tau$ при $K=8.5$, $4.23\,\tau$ при $K=40$): изменение $K$ в полтора раза сдвигает период на $2$--$4\%$, меньше погрешности. Решает размыкание петли (постоянный конечный гормон на входе первой ступени убивает такой ритм) или ответ на короткую добавку гормона в разные фазы цикла.}

\answer{ex:17:1}{$2\cdot10^8$ и $8.1\cdot10^4$~бит/с (в $2.5\cdot10^3$ раза); при $3072$ каналах $6.1\cdot10^8$ и $2.5\cdot10^5$; $0.2$~Вт против $81$~мкВт, для сырого потока нужно не больше $50$~пДж/бит; кадры дают $1.16\cdot10^5$~бит/с, в $1.4$ раза больше ёмкости.}
\answer{ex:17:2}{При $c=0.1$ погрешность $1$, $0.44$, $0.33$, $0.32$ (предел $\sqrt c=0.32$), при $c=0$ $1$, $0.32$, $0.10$, $0.03$; скорость при $c=0.1$ $1.2$, $5.6$, $8.7$, $9.2$~бит/с (предел $9.3$), при $c=0$ до $65$~бит/с при $N=1000$.}
\answer{ex:17:3}{$B=4.87$, $4.37$, $4.04$, $3.29$~бит на знак, то есть $7.3$, $6.6$, $6.1$, $4.9$~бит/с; $B=0$ при $P=1/N$; для $10$~бит/с нужно $2.3$~знака в секунду.}
\answer{ex:17:4}{Ошибка $0.61$, $0.34$, $0.19$, $0.11$, $0.060$ при $\sigma_\xi=0$ и $0.64$, $0.40$, $0.27$, $0.23$, $0.21$ при $\sigma_\xi=0.2$ (формула совпадает до $2\%$); вклады равны при $N\approx90$; предел $\tfrac12\log_2(1+0.25/0.04)\approx1.4$~бита на компоненту за окно.}
\answer{ex:17:5}{$\dot b=-\eta_bd\varepsilon$, $\dot d=-\eta_db\varepsilon$; при $\eta_b=1$, $\eta_d=0.1$ приход в $(1.06,\,0.94)$, при обратных скоростях в $(1.48,\,0.67)$; сходимость при $\eta_bd^2+\eta_db^2<2$; при $\eta=1.1$ и $1.2$ устойчивые колебания периода $2$, при $1.5$ непериодические колебания, при $2$ разнос; скорости надо разводить.}
\answer{ex:17:6}{Порог $\approx4.4\sigma$ (при $3.5\sigma$ около $4.7$~Гц ложных); А: $C=29$, $306$~кбит/с, $0.31$~мВт; Б: $102$~кбит/с, $0.10$~мВт, насыщено $5.7\cdot10^{-5}$ отсчётов, энтропия $0.77$~бит на отсчёт ($40$~кбит/с); сырой поток $205$~мВт.}
\answer{ex:17:7}{Открытая задача. Например, $D=10$, $W=5$~Гц: около $46$~бит/с при шуме, ограниченном корреляцией ($\text{SNR}\approx0.9$), и около $330$~бит/с при независимом шуме от $1000$ нейронов; измерено порядка $10$~бит/с.}

\answer{ex:18:1}{(а)~$0.9949\le w\le1.0049$, то есть $|1-w|\lesssim0.005$. (б)~$K\ge400$ синапсов.}
\answer{ex:18:2}{$(\omega_R\tau_w)^2=\sqrt{\alpha(\alpha+2+2\varrho)}/\varrho-1$, $\varrho=\tau_m/\tau_w$; $f_R\approx11.1$~Гц; $|Z(\omega_R)|/|Z(0)|=4.6$.}
\answer{ex:18:3}{(а)~$f=1/\bigl(\tau\ln\frac{\mu}{\mu-\theta}\bigr)$; для $1$~Гц нужно $\mu/\theta-1\approx3.7\cdot10^{-44}$. (б)~$T=\pi\tau/\sqrt\mu$, $f\approx3.2$~Гц: $f\propto\sqrt\mu$.}
\answer{ex:18:4}{Интегратор работает при $f\lesssim17.7$~Гц (фильтр нижних частот), резонатор --- только в полосе $5.5$--$22$~Гц (полосовой фильтр).}
\answer{ex:18:5}{(а)~$T_{\min}=\frac{\tau}{w-1}\ln\frac{0.5}{0.3}\approx10.2\ms$. (б)~$B>w-1-0.2=0.3$.}
\answer{ex:18:6}{Постоянная времени подстройки $\tau/(\eta\langle r^2\rangle)$, $\eta=\tau\ln10/60\,\text{с}\approx3.8\cdot10^{-3}$; при $I\neq0$ вес смещается, из $e$ нужно вычесть копию команды $I/\tau$.}
\answer{ex:18:7}{Открытая задача. Резонатор выигрывает лишь в $1.35$ раза при $11$~Гц и проигрывает в $17$ раз при $1$~Гц; основной выигрыш от перестройки --- в пороговом выборе полосы (задача~\ref{ex:18:4}).}

\answer{ex:19:1}{(а)~Один синапс даёт $6.7\mV$ ($R=66.7$~МОм), так что до порога вообще нужно не меньше $4$ (трёх хватает лишь асимптотически); $8$ за $10\ms$; $14$ за $5\ms$. (б)~$0.1\ms\ll\tau_m$, поправка от утечки $\sim0.25\%$.}
\answer{ex:19:2}{(а)~$1.75\cdot10^{10}$ (четверть смесителей отвечает на всё подряд), $4.4\cdot10^9$ и $0.06$ (при $k=40$ не отвечает почти никогда). (б)~В $10^3$ раз: $2\cdot10^5$ против $200$.}
\answer{ex:19:3}{(а)~$C(2n,n)=2^{2n-1}$ по симметрии биномиальных коэффициентов. (б)~$0.93$ и $0.09$; ширина перехода $\sim\sqrt{2n}$ по $P$, относительная $\sim1/\sqrt n$.}
\answer{ex:19:4}{С одного дендрита $V_{\text{тело}}<\kappa E_E<\theta$ при любом числе входов; при $g_0=0.5g_L$ нужно $n\ge3$ на каждом дендрите.}
\answer{ex:19:5}{$I\ge C\sigma_V/\sigma_t=15$~нА, то есть $150$ синхронных синапсов --- нереалистично; маленькому нейрону хватило бы $1$~нА, а с синапсом в $4$~нА дрожание $5$~мкс.}
\answer{ex:19:6}{Максимум на ближнем конце: $0.03\,\tau_m$ и $0.97$ ($L=0.2$); $0.32\,\tau_m$ и $0.66$ ($L=1$); $0.79\,\tau_m$ и $0.33$ ($L=2$). Оценка~\eqref{eq:dendriteload} по полной ёмкости верна только при $L\ll1$.}
\answer{ex:19:7}{Открытая задача. При фиксированном $m$ время отклика растёт линейно с запоминаемым $P$; при фиксированной доле $m=\varphi n$ оно перестаёт зависеть от $n$, и медленность большого нейрона --- следствие суммирования многих входов, а не размера как такового.}

\answer{ex:20:1}{$t_{\text{б}}=\tau_r\ln\frac{1-L}{1-H}=16.8$~ч, $t_{\text{с}}=\tau_d\ln\frac HL=5.8$~ч, период $22.5$~ч; к $24$~ч генератор подтягивает суточное качание порогов --- фазовая автоподстройка~\eqref{eq:pll}.}
\answer{ex:20:2}{(а)~$L=0.111$; при $L=0.092$ сон длился бы $9.6$~ч. (б)~$S_0$ больше в $1.44$ раза, сон длиннее в $1.10$ раза: скорость разряда пропорциональна $S$. (в)~На $22\%$ ($1/0.82=1.22$).}
\answer{ex:20:3}{$H=\frac{p-1}{p-1/q}=0.623$, $L=H/q=0.093$, где $p=e^{16/\tau_r}$, $q=e^{8/\tau_d}$. После побудки через $4$~ч фаза сдвигается навсегда (засыпание на $7.2$~ч раньше); с качающимися порогами она возвращается за двое--трое суток.}
\answer{ex:20:4}{(а)~$r_\pm=\frac12\bigl(1\pm\sqrt{1-4k/w}\bigr)=0.947,\ 0.053$; $a_{\text{в}}=3.51$, $a_{\text{н}}=0.49$. (б)~$T_{\text{раб}}\approx0.33$~с, $T_{\text{тиш}}\approx0.59$~с, $T\approx0.93$~с, доля работы $0.36$. (в)~Модель $1.11$~с из-за конечного $\tau/\tau_a$; при росте $\tau_a$ отношение $T/\tau_a$ стремится к $3.2$. Быстрее суточного в $7.8\cdot10^4$ раз.}
\answer{ex:20:5}{(а)~$a_{\text{в}}/r_+<b<a_{\text{н}}/r_-$: $3.71<b<9.20$. (б)~$r\approx1$, $a\approx1$, $F'\approx0$ --- устойчиво. (в)~\nt{АЦЕТ} уменьшает $b$ ниже $3.71$, и пересечение уходит на верхнюю ветвь.}
\answer{ex:20:6}{(а)~$16.79$ и $5.76$~ч. (б)~При $D=24$, $32$, $40$, $48$, $64$ сон длится $4.9$, $6.8$, $8.8$, $5.5$, $8.9$~ч и совпадает с $\tau_d\ln(S_0/L)$ до $0.01$~ч; длительность задаёт суточная фаза засыпания, а не долг. (в)~Колебания при $3.8\le b\le9.0$ (нет при $3.7$ и $9.2$); период $1.61$~с при $b=3.8$, $1.11$~с при $b=5$, минимум $0.98$~с около $b=8$.}
\answer{ex:20:7}{(а)~Ошибка на A после B в среднем $0.51$ (от $0.19$ до $1.08$), ошибка нулевого ответа --- $1$. (б)~Обе ошибки меньше $10^{-4}$. (в)~Изменения лежат в линейной оболочке образов B; общее решение есть, пока $P_A+P_B\le n$.}
\answer{ex:20:8}{Открытая задача. Равномерное масштабирование сохраняет отношения весов, но сохраняет ответы порогового нейрона лишь при масштабировании порога в то же число раз; перемешанные повторы меняют отношения. Неизменность крупных контактов противоречит чисто равномерному масштабированию.}

\answer{ex:21:1}{(а)~$T=\tau_{\text{вос}}\ln\frac{1-x_0}{1-x^*}$. (б)~$1.84$~с и $3.68$~с; $\dd T/\dd x^*=\tau_{\text{вос}}/(1-x^*)=8$~с и $80$~с. (в)~От $3.36$ до $4.24$~с: у $x^*\to1$ залпы нерегулярны.}
\answer{ex:21:2}{(а)~$10^{-4}$~Вт; для мозга $8.6$~Вт, как в~\eqref{eq:energy}. (б)~$205$~Мбит/с, $0.2$~Вт --- в $2\cdot10^3$ раз больше культуры.}
\answer{ex:21:3}{(а)~$x_{\text{уст}}=1/(1+Ur_b\tau_{\text{вос}})$; $r_b>(1/x^*-1)/(U\tau_{\text{вос}})=0.28$~Гц при $x^*=0.9$ и $0.025$~Гц при $x^*=0.99$. (б)~Сила синапса падает до $x_{\text{уст}}$, то есть на $1$--$10\%$. (в)~Залп запускается там, где запас восстановлен.}
\answer{ex:21:4}{$u(t-k)$ ортонормированы, $m_k=\|P_Vu(t-k)\|^2$, где $P_V$ --- проектор на $N$-мерную оболочку выходов; по неравенству Бесселя $\sum_k m_k\le N$.}
\answer{ex:21:5}{(а)~$\mathrm{MC}\approx9$--$11$ с $\tanh$, $15$--$18$ и $17$--$21$ для линейного при $0.9$ и $0.99$ (три затравки); всегда $\le30$. (б)~Резервуар: точность $0.985$--$1.0$, $r^2=0.84$--$0.91$; отводы и резервуар без смещений: $r^2<0.01$ --- нечётная нелинейность не создаёт чётных произведений.}
\answer{ex:21:6}{(а)~При $P\le n$ разделимы все $2^P$ разметок. (б)~$n\le102$. (в)~На $5.6$ стандартных отклонения выше $50\%$.}
\answer{ex:21:7}{Открытая задача. Ключевой контроль: заморозить считывание и проверить, сохраняется ли улучшение после паузы без стимуляции; если нет, изменилось состояние, а не веса.}

\answer{ex:22:1}{(а)~$32\%$ и $100\%$. (б)~$R_\uparrow-R_\downarrow\ge2\sqrt{2000\,\text{Гц}/1\,\text{с}}\approx89$~Гц, около $4.5\%$ суммарной частоты.}
\answer{ex:22:2}{(а)~$n\ge4(p_1q_1+p_2q_2)/\Delta^2\approx560$. (б)~$\approx56$. (в)~Удары одной культуры коррелированы через её настройку; единица статистики --- культура.}
\answer{ex:22:3}{(а)~Детальный баланс $\rho P K$ симметричен; единственно при неприводимом $K$. (б)~Доля промахов $=1/\langle1/P\rangle\le\langle P\rangle$, равенство при $P=\text{const}$. (в)~$0.18$ против $0.5$. (г)~Такие настройки поглощающие; вся масса собирается в них.}
\answer{ex:22:4}{(б)~Параллелограмм площади $4h^2=0.04$ (с учётом ограничения $p$ --- численно $0.045$). (в)~$\approx0.18$, в согласии с $0.19$ на рис.~\ref{fig:dishmodel}.}
\answer{ex:22:5}{(а)~$0.49$; в поглощающей области $44\%$ культур (после $2000$ розыгрышей --- $32\%$); остальные уходят в область, где промах почти всегда. (б)~$0.80$ и $1.00$. (в)~$0.25$, $0.32$, $0.38$, $0.36$, $0.32$, $0.23$ при $\sigma=0.02$--$0.64$ ($\pm0.03$); оптимум при $\sigma\approx0.1\approx h$.}
\answer{ex:22:6}{(а)~$\ge5$ уровней; $8$ электродов хватает. (б)~Нет: нужно $r_{\max}\ge25/T=100$~Гц. (в)~$\sqrt{10}\approx3$ уровня, $\approx1.7$~бита; высоту --- местом, расстояние --- частотой.}
\answer{ex:22:7}{Открытая задача. Время случайного поиска растёт с $d$ примерно как отношение объёмов $(L/h)^d$, у поиска со знаком ошибки --- линейно по $d$; опыт с перестановкой предсказуемости и силы стимулов разделяет гипотезы, если в каждой группе несколько десятков культур.}

\answer{ex:23:1}{Ответ падает как $(1-f)^n$; вдвое --- после $n=\ln2/[-\ln(1-f)]\approx14$, $7$ и $1$ запаха (модель при $f=5\%$: $0.50$ после $14$). Без забывания банк быстро «забивается»; $\tau_c$ освобождает строки.}
\answer{ex:23:2}{$K=5$: $\binom{2000}{2}/\binom{51}{5}\approx0.85$ пары (при $K=4$ --- $8$, при $K=7$ --- $0.017$). Доля общих входов $K/51$: $0.14$ при $K=7$ и $0.49$ при $K=25$ --- клетки становятся полукопиями друг друга.}
\answer{ex:23:3}{$\ln2/[-\ln(1-\eta)]$: $6.6$ постукивания ($66$~с) без еды и $13.5$ ($135$~с) на еде. Конечный уровень $A_\infty$ тот же: флаг меняет скорость, а не цель.}
\answer{ex:23:4}{$\Delta w/(\eta\tau)=e^{\Delta/\tau}(1-e^{-T/\tau})$ при $\Delta<0$; $e^{-\Delta/\tau}-e^{-(T-\Delta)/\tau}$ при $0\le\Delta\le T$; $-(1-e^{-T/\tau})\,e^{-(\Delta-T)/\tau}$ при $\Delta>T$. Ноль --- при $\Delta=T/2$.}
\answer{ex:23:5}{Одно сочетание: $0.95$, $0.17$, $0.034$; десять: $1.40$, $0.50$, $0.27$. На четвёртый день память почти целиком в медленном банке, и десять сочетаний дают её в восемь раз больше, чем одно.}
\answer{ex:23:6}{С неподвижным порогом загорается $\approx65\%$ клеток, и выученный запах избегается лишь на $\approx0.08$; при $1.5$-кратной силе --- $36\%$ и $0.14$. С отбором $k$ победителей код не меняется, отвращение $1.00$.}
\answer{ex:23:7}{Открытая задача. Выход, выученный на сахар, должен возбуждать учителя наказания, а сахар --- гасить это возбуждение; тогда учитель наказания несёт отрицательную часть $\delta$ из~\eqref{eq:ctd}, учитель награды --- положительную, и пара линий заменяет одно число со знаком.}


\chapter{Экспериментальное приложение}
\label{sec:supplement}

Для каждой главы здесь собраны её главные утверждения и то, на чём они держатся: опыт, в котором это измерено, что именно измерено и где это опубликовано. Статус в правом столбце говорит, насколько утверждение твёрдо: \emph{измерено} --- есть прямой опыт; \emph{теория} --- математическое следствие, выведенное в главе или в классической работе; \emph{модель} --- результат расчёта по модели книги (скрипты в папке \texttt{models/}); \emph{оценка} --- число порядка величины без прямого измерения; \emph{гипотеза} --- правдоподобное объяснение, которое опытом ещё не доказано. Где опыта нет, так и сказано.

\section{Глава~\ref{sec:neurontypes}. Типы нейронов}

Числа главы держатся на прямых подсчётах клеток в мозге человека там, где такие подсчёты есть, правило «один нейрон --- один медиатор» --- на клеточных атласах и опытах, которые нашли и исключения, роль \nt{ДОФА} --- на записях импульсов, а роли остальных широковещательных каналов --- на гипотезах.

{\footnotesize
\setlength{\tabcolsep}{3pt}\renewcommand{\arraystretch}{1.15}
\begin{longtable}{>{\raggedright}p{0.5cm}>{\raggedright}p{4.0cm}>{\raggedright}p{5.6cm}>{\raggedright}p{2.3cm}>{\raggedright\arraybackslash}p{1.7cm}}
\toprule
№ & Утверждение & Опыт и что измерено & Источник & Статус \\
\midrule\endfirsthead
\toprule
№ & Утверждение & Опыт и что измерено & Источник & Статус \\
\midrule\endhead
1 & В мозге человека около $8.6\cdot10^{10}$ нейронов, и почти все \nt{ГЛУТ}-нейроны --- мелкие нейроны мозжечка (табл.~\ref{tab:types}) & Изотропный фракционатор: ткань растворяют в однородную взвесь ядер и считают ядра с нейронным маркером NeuN. У взрослых мужчин $86.1\pm8.1$ млрд нейронов; в коре больших полушарий из них лишь $19\,\%$, основная масса --- в мозжечке & \cite{Azevedo2009} & измерено \\
2 & У почти каждого нейрона один главный медиатор (утверждение~\ref{prop:dale}), но есть исключения & Транскриптомный атлас мозга мыши: около $4\cdot10^6$ клеток, $5322$ кластера; кластерам приписан медиатор по генам синтеза и упаковки. Исключения измерены прямо: \nt{ДОФА}-нейроны в срезах выбрасывают ещё и ГАМК через тот же везикулярный транспортёр и тормозят нейроны стриатума; у части \nt{ДОФА}-аксонов глутамат и дофамин выходят из разных окончаний одного провода (электронная микроскопия и оптогенетика) & \cite{Strata1999,Yao2023,Tritsch2012,Zhang2015} & измерено \\
3 & Весь столбец матрицы весов одного знака~\eqref{eq:dalesign} & Вывод в главе из утверждения~\ref{prop:dale} и того, что рецепторы глутамата почти все возбуждающие, а ГАМК --- тормозные. Прямой проверки «все получатели одного нейрона получают вес одного знака» как общего правила нет; контрпримеры --- совместный выброс ГАМК \nt{ДОФА}-нейронами (строка~2) и получатели с рецепторами разного знака (строка~11) & вывод в главе & теория \\
4 & От трёх четвертей до четырёх пятых нейронов коры --- \nt{ГЛУТ}, от пятой части до четверти --- \nt{ГАМК} & Иммуноокраска на ГАМК в столбиках шириной $50$\,мкм через всю толщу коры в $10$ областях коры обезьяны: в $8$ областях ГАМК-нейроны составляют около $25\,\%$ всех нейронов, в первичной зрительной коре --- меньше $20\,\%$ & \cite{Hendry1987} & измерено \\
5 & У \nt{ГЛУТ}-нейрона около $10^4$ выходов & Стереология на вскрытии: в новой коре молодых мужчин $1.64\cdot10^{14}$ синапсов (электронная микроскопия, дисектор) и около $2.3\cdot10^{10}$ нейронов. Отношение даёт $\sim7\cdot10^3$ синапсов на нейрон; в среднем число входов равно числу выходов & \cite{Tang2001synapses,Pakkenberg1997} & измерено \\
6 & Выход \nt{ДОФА}-нейрона ветвится на $10^5$--$10^6$ окончаний; это оценка по охвату мишени, а не подсчёт & Вирусная метка мембраны отдельных \nt{ДОФА}-нейронов крысы и полная реконструкция аксона: один нейрон покрывает от $0.45$ до $5.7\,\%$ (в среднем $2.7\,\%$) объёма стриатума. Измерен охват, а не число окончаний: оно выведено из охвата по порядку величины, прямого счёта окончаний нет. Для \nt{СЕРО} и \nt{НОРА} числа окончаний --- грубые оценки & \cite{Matsuda2009} & измерено \\
7 & Широковещательных нейронов мало: \nt{ДОФА} $\sim5\cdot10^5$ (главная из двух групп, нижняя граница), \nt{СЕРО} $\sim3\cdot10^5$ (сумма двух частичных подсчётов), \nt{ГИСТ} $\sim6\cdot10^4$ (табл.~\ref{tab:types}, рис.~\ref{fig:typepie}) & Стереологические подсчёты у человека: $550\,000$ пигментированных нейронов в чёрной субстанции (без вентральной покрышки); $235\,000$ нейронов в дорсальном ядре шва (все нейроны ядра) и ещё $125\,000$ серотониновых нейронов в покрышке моста; около $64\,000$ гистаминовых нейронов гипоталамуса. Для \nt{НОРА}, \nt{ГЛИЦ}, \nt{АЦЕТ} и \nt{АДРЕ} прямого подсчёта нет: в таблице это грубые оценки по порядку величины & \cite{Pakkenberg1991,Baker1990,Baker1991,Airaksinen1991} & измерено \\
8 & Глутамат в щели взлетает примерно до миллимоля и спадает за $\sim1\ms$ & По вытеснению быстро отсоединяющегося конкурентного блокатора с рецепторов NMDA во время синаптической передачи (культура нейронов гиппокампа): пик $1.1$\,мМ, постоянная спада $1.2\ms$ & \cite{Clements1992} & измерено \\
9 & В работающей коре возбуждающий и тормозящий токи почти уравновешены, нейрон сидит у порога & Теория: сеть с сильными связями сама приходит в уравновешенный режим. Опыт: в префронтальной коре хорька in vivo во время «вверх»-состояний потенциал реверсии синаптического тока держится около $-37$\,мВ сотни миллисекунд, хотя проводимость меняется на $\sim21$\,нСм: возбуждение и торможение растут и падают вместе & \cite{vanVreeswijk1996,Haider2006} & измерено \\
10 & \nt{ДОФА}-нейроны сообщают ошибку предсказания награды~\eqref{eq:rpe} (рис.~\ref{fig:rpe}) & Импульсы \nt{ДОФА}-нейронов обезьяны: пачка на неожиданный сок; после обучения --- пачка на сигнал и нет ответа на предсказанный сок; провал частоты, когда обещанный сок не дан. В количественном опыте частота пропорциональна разности текущей награды и взвешенного среднего прошлых, но только для положительных ошибок. Само уравнение~\eqref{eq:rpe} --- алгоритм временных разностей & \cite{Schultz1997,Bayer2005,SuttonBarto2018} & измерено \\
11 & Знак и скорость задаёт рецептор: ионотропные --- доли миллисекунды, метаботропные --- гораздо медленнее (утверждение~\ref{prop:receptor}) & Быстрые импульсы глутамата на кусочки мембраны нейронов гиппокампа: ток через ионотропные рецепторы нарастает (от $20$ до $80\,\%$) за $0.2$--$0.6\ms$ и спадает за $2$--$3\ms$. Медленный тормозный потенциал пирамидных нейронов снимается избирательным блокатором метаботропного рецептора ГАМК. Два семейства рецепторов дофамина с противоположным действием; в стриатуме нейронам с рецептором D1 дофамин нужен для усиления синапсов, нейронам с D2 --- для ослабления & \cite{Colquhoun1992,DutarNicoll1988,Beaulieu2011,Shen2008} & измерено \\
12 & Широковещательный канал --- не один провод, а жгут из немногих линий (раздел~\ref{sec:buses}) & Вирусно-генетическая разметка серотониновых нейронов дорсального ядра шва мыши: нейроны, посылающие выход в миндалину и в лобную кору, лежат в разных местах ядра, ветвятся в дополняющие друг друга области и противоположно отвечают на неприятный раздражитель. Обзор: подгруппы нейронов голубого пятна (\nt{НОРА}) кодируют разные процессы & \cite{Ren2018,Poe2020} & измерено \\
13 & \nt{НОРА} задаёт коэффициент усиления $g$ & Гипотеза адаптивного усиления; модель сети, в которой катехоламины повышают крутизну передаточной характеристики. Прямого измерения «$g$ растёт с норадреналином» во всей сети нет; измерено, что диаметр зрачка следует за импульсами нейронов голубого пятна обезьяны & \cite{AstonJones2005,ServanSchreiber1990,Joshi2016} & гипотеза \\
14 & \nt{АЦЕТ} переключает «запись/чтение» и задаёт скорость обучения & Гипотеза. Опора: в срезах обонятельной коры крысы агонисты ацетилхолина ($100$\,мкМ) ослабляют внутренние, «вспоминающие» синапсы больше чем на $60\,\%$, а входные --- меньше чем на $15\,\%$ & \cite{Hasselmo2006,HasselmoBower1992} & гипотеза \\
15 & \nt{СЕРО} задаёт коэффициент дисконтирования $\gamma$; серотониновые нейроны работают ровно, несколько импульсов в секунду, и почти молчат в одной из фаз сна & Гипотеза «серотонин $=\gamma$». Сильнейший довод: оптогенетическая активация серотониновых нейронов дорсального ядра шва мыши уменьшает число преждевременных отказов от ожидания отложенной награды. Частота импульсов и молчание во сне с быстрыми движениями глаз измерены (обзор записей) & \cite{Doya2002,Miyazaki2014,JacobsAzmitia1992} & гипотеза \\
\bottomrule
\end{longtable}}

\section{Глава~\ref{sec:glutcomputer}. Что вычисляют \nt{GLUT}-нейроны}

Логические утверждения главы --- теоремы о схемах из неотрицательных весов, выведенные в самой главе; опыт подтверждает модель нейрона и то, что схемы, собранные по этим теоремам (детектор совпадений, линия задержки, самоподдерживающаяся группа, цепочка), в мозге действительно встречаются.

{\footnotesize
\setlength{\tabcolsep}{3pt}\renewcommand{\arraystretch}{1.15}
\begin{longtable}{>{\raggedright}p{0.5cm}>{\raggedright}p{4.0cm}>{\raggedright}p{5.6cm}>{\raggedright}p{2.3cm}>{\raggedright\arraybackslash}p{1.7cm}}
\toprule
№ & Утверждение & Опыт и что измерено & Источник & Статус \\
\midrule\endfirsthead
\toprule
№ & Утверждение & Опыт и что измерено & Источник & Статус \\
\midrule\endhead
1 & Нейрон --- RC-цепь с порогом и сбросом~\eqref{eq:glutneuron} & В тело пирамидных нейронов коры крысы (срезы) подавали шумовой ток, похожий на синаптический поток в живом мозге. Зависимость средней частоты от среднего и разброса тока описывается интегрирующим нейроном с утечкой, если добавить медленную адаптацию частоты. Диапазоны $\tau$ и задержек приведены в главе без ссылки на опыт & \cite{Rauch2003} & измерено \\
2 & Модель~\eqref{eq:glutneuron} --- упрощение: ветви входов сами выдают местные импульсы & Запись прямо с дендритов пирамидных нейронов зрительной коры мыши in vivo: зрительный раздражитель вызывает местные дендритные импульсы, зависящие от рецепторов NMDA; их подавление расширяет настройку нейрона на ориентацию & \cite{Smith2013} & измерено \\
3 & Окно совпадения $\Delta_{\max}=\tau\ln\frac{w}{\theta-w}$~\eqref{eq:window}; при $\theta=1.5w$ оно равно $0.7\tau$ & Следствие модели~\eqref{eq:glutneuron}. Прямое измерение узкого окна суммирования есть для нейронов с быстрым торможением (глава~\ref{sec:gaba}, строка~3 её таблицы) & вывод в главе & теория \\
4 & Сеть из одних \nt{GLUT}-нейронов вычисляет только монотонные функции входов (утверждение~\ref{prop:monotone}) & Доказательство в главе: итерация от тишины с неубывающей правой частью. Это утверждение о модели; опыта, который проверял бы его на живой сети без торможения, нет & вывод в главе & теория \\
5 & Органы чувств дают пару проводов: одни клетки отвечают на усиление раздражителя, другие --- на ослабление & Сетчатка: уже на биполярных клетках сигнал колбочек расходится на включающий и выключающий каналы. Фармакологический блок включающих биполярных клеток у обезьяны: каналы идут раздельно до коры; после блока ухудшается обнаружение усиления света, но не ослабления & \cite{Schiller1992} & измерено \\
6 & С двухпроводным кодом сеть из \nt{GLUT}-нейронов вычисляет любую логическую функцию асинхронно, без общего такта, ценой вдвое большего числа нейронов (утверждение~\ref{prop:dualrail}, \eqref{eq:dualrail}, рис.~\ref{fig:dualxor}) & Двухпроводный код и определение готовности по самому сигналу --- стандартная техника асинхронных схем, нечувствительных к задержкам. Схема исключающего ИЛИ из шести нейронов собрана в главе. Опыта, показывающего, что мозг вычисляет логические функции именно двухпроводным кодом, нет & \cite{Sparso2001} & теория \\
7 & Наибольшая разница времени прихода звука в уши у человека $\approx0.6\ms$, а мозг различает около десяти микросекунд & Слушатели в опыте с двумя вариантами ответа: порог различения разницы времени (75\,\% правильных) --- $9$\,мкс для шума в полосе $150$--$1700$\,Гц, $11$\,мкс для тона $1$\,кГц, $28$\,мкс для щелчка. Предел $0.6\ms$ --- расчёт главы, $0.22\,\text{м}/343\,\text{м/с}$ & \cite{KlumppEady1956} & измерено \\
8 & У птиц разница времён превращается в место через линии задержки и детекторы совпадений~\eqref{eq:itd} (рис.~\ref{fig:itd}) & Сипуха: аксоны от двух ушей входят в ядро детекторов совпадений с противоположных сторон; время прихода импульсов вдоль аксона меняется упорядоченно с глубиной, и разброс задержек через ядро ($160$\,мкс на $700$\,мкм) совпадает с диапазоном межушных задержек & \cite{Carr1990} & измерено \\
9 & У млекопитающих ряда задержек не нашли; ту же задачу решают детекторы совпадений с точно рассчитанным торможением от \nt{GLYC}-нейронов & Запись in vivo в медиальной верхней оливе песчанки: ответы не образуют карты направлений; подведение глицина и его блокатора через микропипетку показывает, что точно рассчитанное во времени \emph{глициновое} торможение задаёт рабочий диапазон задержек. Торможение здесь от \nt{GLYC}, а не от \nt{GABA} & \cite{Brand2002,Grothe2010} & измерено \\
10 & Группа с сильной обратной связью --- триггер; самоподдерживающаяся активность держится секунды после раздражителя (рис.~\ref{fig:bistable}) & Префронтальная кора обезьяны в задаче с отложенным движением глаз к запомненной точке (задержка $1$--$6$\,с): $87$ из $170$ нейронов, связанных с задачей, меняют частоту во время задержки, у $79\,\%$ из них это зависит от направления на запомненную точку. Что эту активность держит обратная связь с двумя устойчивыми состояниями, --- теория & \cite{Funahashi1989,Wang2001} & измерено \\
11 & Бит записывается, если приток длится дольше $\approx0.7\tau$ (рис.~\ref{fig:recurrentgroup}б) & Расчёт уравнения $\tau\,\dd r/\dd t=-r+f(wr+I)$ методом Эйлера при $w=1$, $I=0.6$; скрипта в \texttt{models/} нет. Опыта нет & расчёт в главе & модель \\
12 & Память можно хранить в облегчении синапсов, почти без импульсов & Теория: остаточный кальций в окончаниях держит запись около секунды. Опора: в медиальной префронтальной коре хорька (запись с нескольких нейронов сразу) есть подсеть пирамидных нейронов, связанных облегчающимися синапсами с долгим последействием. Обзор наблюдений «тихих» промежутков при удержании в памяти. Какой способ кора использует когда --- открытый вопрос & \cite{Mongillo2008,WangMarkram2006,Stokes2015} & гипотеза \\
13 & Память стирается усталостью: запас медиатора истощается & Парные записи пирамидных нейронов коры: ответ синапса падает при частых импульсах, скорость падения задаётся вероятностью выброса. Что именно это гасит самоподдерживающуюся группу, прямо не показано & \cite{Tsodyks1997} & измерено \\
14 & Цепочка групп --- таймер, запускаемый событием; у певчих птиц нейроны срабатывают строго по очереди & Запись опознанных нейронов высшего вокального центра зебровой амадины во сне и при пении: каждый нейрон, посылающий выход в моторное ядро, выдаёт одну короткую пачку в один точный момент песни, и нейроны срабатывают последовательно & \cite{Hahnloser2002} & измерено \\
\bottomrule
\end{longtable}}

\section{Глава~\ref{sec:gaba}. \nt{ГЛУТ} и \nt{ГАМК} вместе}

Логические и динамические утверждения главы выведены в ней самой из линейного анализа и закона Ома; опыт показывает, что схемы, на которые они опираются, --- запрет с задержкой, прямое торможение вслед за возбуждением, стабилизация торможением, ритм петли \nt{ГЛУТ}--\nt{ГАМК}, --- в мозге действительно работают.

{\footnotesize
\setlength{\tabcolsep}{3pt}\renewcommand{\arraystretch}{1.15}
\begin{longtable}{>{\raggedright}p{0.5cm}>{\raggedright}p{4.0cm}>{\raggedright}p{5.6cm}>{\raggedright}p{2.3cm}>{\raggedright\arraybackslash}p{1.7cm}}
\toprule
№ & Утверждение & Опыт и что измерено & Источник & Статус \\
\midrule\endfirsthead
\toprule
№ & Утверждение & Опыт и что измерено & Источник & Статус \\
\midrule\endhead
1 & \nt{ГАМК}-нейронов в коре около пятой части & Иммуноокраска на ГАМК в $10$ областях коры обезьяны: около $25\,\%$ нейронов в $8$ областях, меньше $20\,\%$ в первичной зрительной коре. Обзор разнообразия тормозящих нейронов коры & \cite{Hendry1987,Markram2004} & измерено \\
2 & Что делает торможение, во многом решает место посадки: дендрит, тело, место рождения импульса, окончание чужого провода & Обзор: классы \nt{ГАМК}-нейронов коры различаются мишенью на получателе. Одновременная запись с тела и дендрита пирамидного нейрона: прямое торможение гораздо сильнее на теле, чем на дендритах. Торможение на окончании чужого провода в спинном мозге уменьшает выброс медиатора одной входной линии (обзор измерений) & \cite{Markram2004,Pouille2001,Rudomin1999} & измерено \\
3 & Смена знака через посредника стоит одной задержки, около миллисекунды; торможение, пришедшее вслед за возбуждением, сужает окно совпадения & Пирамидные нейроны гиппокампа крысы в срезах: торможение через одного посредника приходит вслед за прямым возбуждением с короткой задержкой, и окно, в котором два возбуждающих входа суммируются до порога, меньше $2\ms$. На дендритах, где это торможение слабее, окно шире & \cite{Pouille2001} & измерено \\
4 & Запрет $a\wedge\bar b$~\eqref{eq:veto} с ИЛИ и постоянной единицей функционально полон (утверждение~\ref{prop:gabacomplete}) & Доказательство в главе. Классический результат: сеть пороговых элементов с тормозящими входами реализует любое логическое выражение. Утверждение о модели, опыта нет & \cite{McCullochPitts1943} & теория \\
5 & Запрет с задержкой даёт детектор направления движения с оптимальной скоростью $v^*=\Delta x/d$~\eqref{eq:vstar} & Сетчатка кролика: избирательность к направлению создаётся торможением, которое при движении в «нулевом» направлении гасит возбуждение. Раздельная запись возбуждающего и тормозящего входов избирательных клеток показала, что звёздчатые амакриновые клетки (\nt{ГАМК}) тормозят их несимметрично --- только отростками, направленными в «нулевую» сторону. Формула $v^*$ --- вывод главы & \cite{Barlow1965,Fried2002} & измерено \\
6 & Шунтирующее торможение делит напряжение, а не вычитает~\eqref{eq:shunt} (рис.~\ref{fig:shunt}) & Закон Ома для делителя из трёх проводимостей при равновесном потенциале хлора около покоя; для нейрона без импульсов & вывод в главе & теория \\
7 & На частоту импульсов шунт действует вычитанием, а деление получается от торможения вместе с уравновешенным фоновым шумом & Модель: у срабатывающего нейрона ток через шунт почти постоянен, и эффект на частоту вычитающий. Опыт: в пирамидные нейроны соматосенсорной коры крысы (срезы) через динамический фиксатор вводили фоновые возбуждающие и тормозящие проводимости; одновременный уравновешенный рост обеих делит крутизну ответа без заметного сдвига частоты. Обзор: деление на общую активность встречается во многих отделах мозга & \cite{HoltKoch1997,Chance2002,Carandini2012} & измерено \\
8 & При $\beta>1$ взаимное торможение выбирает одного победителя (утверждение~\ref{prop:wta}, \eqref{eq:wta}) & Собственные числа $-(1\pm\beta)/\tau$ --- вывод в главе. Частоты $0.73$ и $0.53$ при $\beta=0.5$ --- неподвижная точка $r_{1,2}=(I_{1,2}-\beta I_{2,1})/(1-\beta^2)$; скрипта в \texttt{models/} нет. Прямого опыта в главе не приводится & вывод в главе & теория \\
9 & Сброс памяти сигналом через \nt{ГАМК}; две группы с взаимным торможением --- защёлка & Следует из рис.~\ref{fig:bistable} и утверждения~\ref{prop:wta}. Опыта нет & вывод в главе & теория \\
10 & Равновесие петли \nt{ГЛУТ}--\nt{ГАМК} устойчиво, если торможение достаточно сильное и достаточно быстрое (утверждение~\ref{prop:ei}) & Модель двух популяций~\eqref{eq:ei}; условия (i) и (ii) --- определитель и след её матрицы, выведены в главе & \cite{WilsonCowan1972} & теория \\
11 & При $w_{EE}=1.5$ и $w_{EI}w_{IE}=2$ быстрое торможение ($\tau_I=2\ms$) держит $E^*=0.67h$, медленное ($30\ms$) раскачивает колебания (рис.~\ref{fig:ei}) & Численное решение~\eqref{eq:ei}, скрипт \texttt{figures/make\_ei.py} & расчёт & модель \\
12 & Кора работает в режиме стабилизации торможением; подпись --- подтолкнёшь \nt{ГАМК}-нейроны, и их частота упадёт~\eqref{eq:isn} & Теория предсказала парадоксальный ответ. Опыт: внутриклеточные записи в первичной зрительной коре кошки --- при подавлении ответа раздражителем вокруг тормозящая проводимость не растёт, а падает вместе с возбуждающей. Оптогенетическое возбуждение PV-нейронов мыши в зрительной, соматосенсорной и моторной коре снижает частоту тормозящих нейронов, в том числе без раздражителя и при лёгком наркозе. Коэффициент $-1/3$ при числах рис.~\ref{fig:ei} --- вывод главы & \cite{Tsodyks1997isn,Ozeki2009,Sanzeni2020} & измерено \\
13 & Петля «\nt{ГЛУТ} возбуждает \nt{ГАМК}, \nt{ГАМК} тормозит \nt{ГЛУТ}» порождает быстрый ритм & Оптогенетическая стимуляция быстрых \nt{ГАМК}-нейронов в соматосенсорной коре мыши на частотах $8$--$200$\,Гц избирательно усиливает гамма-ритм ($20$--$80$\,Гц, период $12$--$50\ms$); стимуляция \nt{ГЛУТ}-нейронов усиливает только более медленные ритмы. В главе период назван $10$--$30\ms$, то есть верхняя часть этой полосы & \cite{Cardin2009,Buzsaki2012} & измерено \\
\bottomrule
\end{longtable}}

\section{Глава~\ref{sec:code}. Азбука Морзе}

Предельные оценки главы --- теория информации и счёт; измерено, где в канале возникает шум, как быстро работает мозг и сколько стоит импульс, а вопрос о том, каким кодом кора пользуется на самом деле, остаётся открытым.

{\footnotesize
\setlength{\tabcolsep}{3pt}\renewcommand{\arraystretch}{1.15}
\begin{longtable}{>{\raggedright}p{0.5cm}>{\raggedright}p{4.0cm}>{\raggedright}p{5.6cm}>{\raggedright}p{2.3cm}>{\raggedright\arraybackslash}p{1.7cm}}
\toprule
№ & Утверждение & Опыт и что измерено & Источник & Статус \\
\midrule\endfirsthead
\toprule
№ & Утверждение & Опыт и что измерено & Источник & Статус \\
\midrule\endhead
1 & Частоты \nt{ГЛУТ}-нейронов коры --- от долей до десятков импульсов в секунду, и большую часть времени нейроны работают у нижней границы & Запись через прилегающую к клетке пипетку (выбор нейрона не зависит от его активности) в слуховой коре бодрствующей крысы: в каждый момент частоту выше $20$ импульсов в секунду дают меньше $5\,\%$ нейронов; у части нейронов фоновая частота ниже $0.01$ в секунду; распределение частот логнормальное & \cite{Hromadka2008} & измерено \\
2 & Ёмкость импульсного канала $R_{\max}\approx r\log_2\frac{e}{r\Delta t}$~\eqref{eq:spikeentropy}, почти вся она --- во времени импульсов (утверждение~\ref{prop:capacity}) & Энтропия двоичной строки из клеток длины $\Delta t$. Числа главы: $8$ бит на импульс при $r=10$ в секунду и $\Delta t=1\ms$, около $5$ бит при $\Delta t=10\ms$, меньше $2$ бит в секунду у счётчика импульсов & \cite{MacKay1952} & теория \\
3 & Чувствительные нейроны передают от одного до нескольких бит на импульс, в лучших случаях до половины предела & Восстановление раздражителя по импульсам чувствительных нейронов, для которых раздражитель известен; сводка измерений в книге & \cite{Rieke1997} & измерено \\
4 & Генератор импульсов точен; точное время задают быстрые перепады входа & Нейроны коры крысы в срезах: на повторяемый ток с быстрыми колебаниями импульсы повторяются с точностью меньше $1\ms$, на постоянный ток моменты расползаются & \cite{Mainen1995} & измерено \\
5 & Синапс ненадёжен: больше половины импульсов не доходят до получателя из-за случайности выброса & Минимальная стимуляция входов пирамидных нейронов гиппокампа (срезы): больше половины импульсов не вызывают ответа. Исключены колебания порога стимуляции, сбои проведения на ветвлениях аксона и низкая температура опыта; остаётся вероятностный выброс & \cite{AllenStevens1994} & измерено \\
6 & Поток импульсов в живой коре выглядит случайным: интервалы разбросаны как у пуассоновского потока, дисперсия числа импульсов близка к среднему; часть «шума» --- сообщение о движениях животного & Записи в зрительных областях V1 и MT бодрствующей обезьяны: коэффициент вариации интервалов около $1$, отношение дисперсии числа импульсов к среднему --- как у случайного процесса; модель, суммирующая независимые малые входы, даёт разброс гораздо меньше. В зрительной коре мыши значительная часть разброса предсказывается по видеозаписи движений самого животного & \cite{Softky1993,Stringer2019} & измерено \\
7 & Частотный код медленен: ошибка $1/\sqrt{rT}$~\eqref{eq:rateerror}, а мозг узнаёт картинку за $\sim150\ms$ & Формула --- пуассоновская статистика, вывод главы. Опыт: люди решали, есть ли животное на незнакомой фотографии, показанной на $20\ms$; потенциалы мозга на «да» и «нет» расходятся примерно через $150\ms$. Разбиение этого времени на десяток ступеней по $10$--$20\ms$ --- оценка главы, не измерение & \cite{Thorpe1996} & измерено \\
8 & Усреднение по группе ограничено корреляцией: ошибка падает не больше чем в $1/\sqrt c$ раз~\eqref{eq:correlated} (утверждение~\ref{prop:pooling}) & Пары нейронов области MT обезьяны, записанные одновременно: коэффициент корреляции чисел импульсов в среднем $0.12$, и теоретический анализ показывает, что даже такая корреляция ограничивает выигрыш от группы. Теория: предел ставит только та часть корреляции, что похожа на сигнал. Модель противоположной позиции: частота группы из $50$--$100$ нейронов читается за один межимпульсный интервал, $10$--$50\ms$, и время отдельных импульсов в коре не используется & \cite{Zohary1994,MorenoBote2014,ShadlenNewsome1998} & измерено \\
9 & Число можно писать во время первого импульса~\eqref{eq:latency}, в порядке импульсов группы или в фазе местного ритма & Сетчатка саламандры: пространственная структура коротко показанной картинки записана в относительных задержках первых импульсов выходных нейронов, почти независимо от контраста. Клетки места гиппокампа крысы: при проходе через «своё» место импульсы сдвигаются всё раньше в цикле тета-ритма ($7$--$12$\,Гц), сдвиг от $100^\circ$ до $355^\circ$. Насколько кора пользуется временем импульсов --- открытый вопрос (строка~8) & \cite{Gollisch2008,OKeefe1993} & измерено \\
10 & Какой код прочитан, решает постоянная времени получателя~\eqref{eq:reader} (утверждение~\ref{prop:reader}); в активной коре нейроны ближе к детекторам совпадений & Формула~\eqref{eq:reader} --- вывод главы. Обзор записей in vivo: в активной коре проводимость мембраны в несколько раз выше, чем в покое, и постоянная времени соответственно короче. Что кора переключает код именно так, прямо не показано & \cite{Destexhe2003} & гипотеза \\
11 & Истощающийся синапс передаёт изменение частоты, по существу $\Delta r/r$~\eqref{eq:depression} (рис.~\ref{fig:depression}) & Парные записи пирамидных нейронов коры: скорость истощения задаётся вероятностью выброса; при медленном истощении ответ отражает частоту, при быстром --- совпадения. Модель на основе измеренного истощения в коре крысы: одинаковые относительные изменения частоты на быстрых и медленных входах дают одинаковый ответ, то есть синапс передаёт $\Delta r/r$. Числа $1.7\to2.2$, $6.7$ и $90\ms$ --- расчёт главы & \cite{Tsodyks1997,Abbott1997depression} & измерено \\
12 & Пачка --- «тире»: её вероятность потеряться $(1-p)^k$~\eqref{eq:burst} мала, облегчение уменьшает её ещё & Обзор: на ненадёжных синапсах одиночные импульсы часто теряются, а пачки передаются надёжно благодаря облегчению; пачка вызывает долговременные изменения синапсов. Что мозг читает пачку как отдельный знак, --- гипотеза & \cite{Lisman1997,AllenStevens1994} & гипотеза \\
13 & Быстрые \nt{ГАМК}-нейроны задают ритм и «пунктуацию»; ещё один класс тормозит тормозящих и открывает ворота (утверждение~\ref{prop:punctuation}) & Обзор свойств классов \nt{ГАМК}-нейронов коры. Оптогенетика: ритмическое возбуждение быстрых \nt{ГАМК}-нейронов вызывает гамма-ритм, и ответ на раздражитель зависит от фазы цикла. В зрительной коре мыши PV-нейроны тормозят друг друга, SST-нейроны --- все остальные классы, VIP-нейроны --- преимущественно SST. Возбуждение VIP-нейронов у бодрствующей мыши снимает торможение с \nt{ГЛУТ}-нейронов; их включает сигнал подкрепления. Разделение ролей как общее правило --- гипотеза & \cite{Tremblay2016,Cardin2009,Pfeffer2013,Pi2013} & гипотеза \\
14 & Энергия ограничивает среднюю частоту величиной порядка одного импульса в секунду~\eqref{eq:energy} & Энергетический бюджет серого вещества грызуна по анатомическим и физиологическим данным: импульсы и синаптические токи --- $47\,\%$ и $34\,\%$ энергии сигнализации; одновременно могут быть активны не больше $\sim15\,\%$ нейронов. Для коры человека: не больше $\sim1\,\%$ нейронов могут быть сильно активны одновременно. Число $\sim10^9$ АТФ на импульс и мощность $\sim10$\,Вт на сигнализацию --- оценка главы на основе этих расчётов & \cite{Attwell2001,Lennie2003} & теория \\
15 & Код разрежен и подстроен под наибольшее число бит на джоуль; кора меняет точность на энергию (утверждение~\ref{prop:economy}) & Гипотеза экономного кода. Опора: у мышей при ограничении пищи в зрительной коре падает проводимость AMPA-рецепторов и синаптический расход АТФ ($-29\,\%$), частота импульсов сохраняется, а настройка на ориентацию расширяется на $32\,\%$ & \cite{Barlow1961,Padamsey2022} & гипотеза \\
\bottomrule
\end{longtable}}

\section{Глава~\ref{sec:serocomputer}. \nt{СЕРО}-нейроны}

Свойства самих \nt{СЕРО}-нейронов (ритм, автоторможение, знак по рецептору, подсистемы) измерены прямо; вычисления главы --- регулятор, фильтр, логика --- следуют из её уравнений; $\tau_c$, коэффициент диффузии серотонина и число мест выброса --- оценки; роль петли как регулятора уровня в мозге --- гипотеза (в ядре шва автоторможение точечное и даёт лишь короткие паузы), а роль \nt{СЕРО} как горизонта --- гипотеза с поведенческими опытами в её пользу.

{\footnotesize
\setlength{\tabcolsep}{3pt}\renewcommand{\arraystretch}{1.15}
\begin{longtable}{>{\raggedright}p{0.5cm}>{\raggedright}p{4.0cm}>{\raggedright}p{5.6cm}>{\raggedright}p{2.3cm}>{\raggedright\arraybackslash}p{1.7cm}}
\toprule
№ & Утверждение & Опыт и что измерено & Источник & Статус \\
\midrule\endfirsthead
\toprule
№ & Утверждение & Опыт и что измерено & Источник & Статус \\
\midrule\endhead
1 & Знак действия серотонина выбирает рецептор получателя (утверждение~\ref{prop:receptor}) & Рецепторы серотонина разделены на семейства по фармакологии и механизму: 5-HT$_{1A}$ через G-белок открывает калиевые каналы и тормозит клетку, 5-HT$_{2A}$ и ионотропный 5-HT$_3$ возбуждают. Один медиатор даёт противоположные ответы в разных клетках. & \cite{BarnesSharp1999} & измерено \\
2 & \nt{СЕРО}-нейрон в бодрствовании ровно выдаёт несколько импульсов в секунду & Запись нейронов дорсального ядра шва у свободно движущихся кошек: медленный ритмичный разряд $2.8\pm0.2$ имп/с в спокойном бодрствовании; в медленном сне частота падает на $34$--$68\,\%$, в быстром --- на $98\,\%$. Под наркозом у крыс --- $1.7\pm0.5$ имп/с, очень регулярно. & \cite{TrulsonJacobs1979, GagnonParent2014, JacobsAzmitia1992} & измерено \\
3 & \nt{СЕРО}-нейрон --- ритмоводитель: толчок на каждый импульс не нужен, но нужен ровный фоновый приток (в срезе --- норадреналин через $\alpha_1$-рецепторы, в организме --- от \nt{НОРА}) & В срезе мозга клетки разряжаются медленно и ровно, после каждого импульса --- следовая гиперполяризация на $200$--$800\ms$. Спонтанно работающих клеток в срезе меньше, чем в организме: ровному ритму нужен тонический вход норадреналина через $\alpha_1$-рецепторы, блокада которых его гасит. & \cite{VandermaelenAghajanian1983} & измерено \\
4 & Почти все рецепторы метаботропные; ответ медленный: нарастает и спадает в пределах примерно секунды & Все семейства рецепторов, кроме 5-HT$_3$, связаны с G-белками. В срезе вызванный выброс серотонина даёт через 5-HT$_{1A}$ тормозящий ток, который нарастает и спадает в пределах секунды. & \cite{BarnesSharp1999, CourtneyFord2016} & измерено \\
5 & В областях назначения серотонин выходит в объём, а не только в щель; в самом ядре шва торможение соседей идёт точка-точка & Быстрая циклическая вольтамперометрия в срезах: выброс на один импульс одинаков в области с синапсами и в ядре шва, где синапсов почти нет; не зависит от блокады захвата и рецепторов; молекул на окончание намного меньше, чем переносчиков и рецепторов, а внеклеточная концентрация совпадает со сродством главного рецептора. В ядре шва торможение через 5-HT$_{1A}$ идёт точка-точка: переносчик не даёт серотонину накопиться вне щели. & \cite{Bunin1998, CourtneyFord2016} & измерено \\
6 & Концентрация по~\eqref{eq:seroneuron} убывает из-за откачки; $\tau_c$ порядка секунды --- оценка & Вольтамперометрия в живом мозге: внеклеточный серотонин ограничивают прежде всего обратный захват и разрушение, а не синтез. Прямого измерения $\tau_c$ в цитируемых работах нет; порядок величины --- оценка главы. & \cite{Hashemi2012serotonin} & измерено; $\tau_c$ --- оценка \\
7 & НЕ и ИЛИ-НЕ на одном ритмоводителе; полнота \nt{СЕРО}-логики (утверждение~\ref{prop:serocomplete}, рис.~\ref{fig:serogates}) & Следует из~\eqref{eq:seroneuron}: тормозимый ритмоводитель --- ИЛИ-НЕ, а ИЛИ-НЕ функционально полон. Опыта, где из \nt{СЕРО}-нейронов собрана логика, нет; условие раздельных объёмов в мозге не выполняется. & вывод в главе & теория \\
8 & Серотонин тормозит сам \nt{СЕРО}-нейрон через рецепторы на нём & У наркотизированных крыс агонисты 5-HT$_{1A}$ внутривенно и ионофоретически тормозят разряд нейронов дорсального шва с той же зависимостью доза--ответ, что и сам серотонин; агонисты 5-HT$_{1B}$ почти не действуют. В срезе эндогенный серотонин соседних клеток даёт через 5-HT$_{1A}$ ток и короткую паузу в разряде, не меняя средней частоты. & \cite{Sprouse1987, CourtneyFord2016} & измерено \\
9 & В модели петля через общий объём --- регулятор уровня: разброс и постоянная времени уменьшаются в $1+G$ раз~\eqref{eq:feedback}; что петля так работает в мозге, --- гипотеза & Классическая формула усилителя с отрицательной обратной связью; в главе выведена заново. Усиление петли $G$ у \nt{СЕРО}-системы не измерено. Измерено: в срезе автоторможение идёт точка-точка, даёт короткую паузу и не меняет средней частоты. & \cite{Black1934feedback, CourtneyFord2016} & теория; в мозге --- гипотеза \\
10 & Концентрация --- скользящее среднее частоты, фильтр нижних частот~\eqref{eq:lowpass}, рис.~\ref{fig:lowpass} & Решение линейного уравнения~\eqref{eq:seroneuron}; рисунок --- расчёт по этой формуле при $\tau_c=1\,\text{с}$. Отдельной модели в \texttt{models/} нет. & вывод в главе & теория \\
11 & Извилистость щелей замедляет диффузию малой молекулы примерно в $2.6$ раза & Метод точечного источника: ионофорез тетраметиламмония и запись его концентрации ионоселективным электродом в срезах и в живом мозге. Извилистость $\lambda\approx1.6$, то есть эффективный $D$ меньше свободного в $\lambda^2\approx2.6$ раза; доля межклеточного объёма около $0.2$. Коэффициент диффузии самого серотонина в ткани в этих работах не измерялся; в главе он --- оценка. & \cite{Sykova2008} & измерено \\
12 & Длина независимого «провода» $\ell\sim\sqrt{D\tau}\approx20$~мкм~\eqref{eq:diffusionlength}; число проводов ограничено числом таких областей (утверждение~\ref{prop:volumewires}) & Закон диффузии и подстановка оценок $D$ и $\tau$ (строки~6 и~11); утверждение~\ref{prop:volumewires} --- следствие. Прямого опыта, считающего независимые каналы объёмной передачи, нет. & вывод в главе & теория \\
13 & Выход одного \nt{СЕРО}-нейрона раскинут по огромным участкам мозга; мест выброса по оценке $\sim10^5$ & Полная реконструкция отдельных нейронов дорсального шва у крысы: все восходящие аксоны сильно ветвятся, общая длина аксона одной клетки до $18.7$~см, ветви одной клетки доходят до стриатума, префронтальной коры и миндалины. Число мест выброса на клетку прямо не подсчитано. & \cite{GagnonParent2014} & измерено; $10^5$ --- оценка \\
14 & \nt{СЕРО}-нейроны делятся на несколько подсистем с разными выходами и ответами & Вирусно-генетическое мечение у мышей: нейроны, идущие к миндалине и к лобной коре, почти не пересекаются по ветвлению, получают разные входы и противоположно отвечают на неприятный стимул; включение и выключение каждой группы по-разному меняет поведение. & \cite{Ren2018} & измерено \\
15 & Условие терпения $D<\tau_\gamma\ln(R/r_0)$~\eqref{eq:patience} при экспоненциальном обесценивании; при измеренном, гиперболическом, --- $D<\tau_\gamma(R/r_0-1)$ & Следует из обесценивания $e^{-D/\tau_\gamma}$ или $1/(1+D/\tau_\gamma)$; вывод в главе. Обезьяны, выбор на задержках в секунды: обесценивание ближе к гиперболе, чем к экспоненте; ответ \nt{ДОФА}-нейронов на предвестник отложенной награды падает с задержкой тоже по гиперболе. & вывод в главе; \cite{KobayashiSchultz2008} & теория \\
16 & \nt{СЕРО} задаёт горизонт $\tau_\gamma$ (раздел~\ref{sec:horizon}) & Оптогенетическое включение \nt{СЕРО}-нейронов дорсального шва у мышей удлиняет ожидание отложенной награды и повышает упорство в поиске награды, ставшей редкой. У людей понижение серотонина (диета без триптофана) делает обесценивание отложенной награды круче. Как концентрация превращается в $\tau_\gamma$, неизвестно. & \cite{Doya2002, Miyazaki2014, Lottem2018, Schweighofer2008} & гипотеза \\
\bottomrule
\end{longtable}}

\section{Глава~\ref{sec:dofa}. Обучение с \nt{ДОФА}}

Механизмы синапса (порции, детектор совпадений, кальций, окна пластичности) и поведение дофамина как ошибки предсказания измерены прямо; то, что правила ведут к градиенту и чего стоит широковещание, --- теоремы; итог обучения выбору --- расчёт по модели книги; схема, вычисляющая ошибку, --- гипотеза.

{\footnotesize
\setlength{\tabcolsep}{3pt}\renewcommand{\arraystretch}{1.15}
\begin{longtable}{>{\raggedright}p{0.5cm}>{\raggedright}p{4.0cm}>{\raggedright}p{5.6cm}>{\raggedright}p{2.3cm}>{\raggedright\arraybackslash}p{1.7cm}}
\toprule
№ & Утверждение & Опыт и что измерено & Источник & Статус \\
\midrule\endfirsthead
\toprule
№ & Утверждение & Опыт и что измерено & Источник & Статус \\
\midrule\endhead
1 & Обратного распространения ошибки в той форме, что в искусственных сетях, в мозге не найдено & Прямого опыта нет: это утверждение об отсутствии. Обзоры разбирают, чего требует алгоритм (обратные связи, повторяющие прямые, отдельная фаза) и какие биологические приближения к нему предложены. & \cite{Lillicrap2020, WhittingtonBogacz2019} & гипотеза \\
2 & Вес раскладывается на число мест выброса, вероятность выброса и ответ на порцию: $w=N_s\,p\,A_1$~\eqref{eq:quantal} & Нервно-мышечный синапс лягушки при сниженном выбросе: ответ получателя колеблется ступенями, кратными спонтанному миниатюрному ответу на одну порцию; число порций на импульс подчиняется закону Пуассона. & \cite{delCastillo1954} & измерено \\
3 & Рецептор получателя --- детектор совпадений; кальций запускает изменение веса & Патч-клэмп нейронов мыши: ионы Mg$^{2+}$ закупоривают канал, открытый глутаматом, при потенциале покоя и уходят при деполяризации. Срезы гиппокампа: хелатор кальция в получателе блокирует долговременное усиление, а освобождение кальция светом внутри клетки само усиливает передачу. & \cite{Nowak1984, Malenka1988calcium} & измерено \\
4 & Решение исполняет любая сторона: растёт $A_1$ у получателя, меняется $p$ у отправителя & Глутамат, освобождаемый светом над одним шипиком, вызывает быстрый и стойкий рост шипика и тока через его рецепторы; усиление сопровождается встраиванием AMPA-рецепторов. Деполяризация получателя на время подавляет выброс у отправителя; эффект переносят эндоканнабиноиды, идущие назад через щель (показано на \nt{ГАМК}-входах). Кратковременное истощение зависит от вероятности выброса у отправителя. & \cite{Matsuzaki2004, Malinow2002, WilsonNicoll2001, Tsodyks1997} & измерено \\
5 & Синапс усиливается, когда отправитель и получатель активны вместе~\eqref{eq:hebb} & Кролик под наркозом: после коротких серий импульсов по входному пути гиппокампа ответ усилен на время от $30$~мин до $10$~ч. В срезе гиппокампа импульсы получателя усиливают только те синапсы, которые в тот же момент активны. & \cite{Bliss1973LTP, Kelso1986hebb} & измерено \\
6 & Правило~\eqref{eq:oja} находит главный собственный вектор корреляций входов (утверждение~\ref{prop:oja}) & Теорема; доказательство в главе. Опыта, где нейрон выучил главную компоненту своих входов, нет; механизм ограничения роста весов в синапсе не установлен. & \cite{Oja1982} & теория \\
7 & Хебб по времени: окно около $\pm20\ms$ (рис.~\ref{fig:stdp}); из повторяемых последовательностей вырастают цепочки & Пары нейронов гиппокампа в культуре: импульс получателя в пределах $20\ms$ после импульса отправителя даёт стойкое усиление, в пределах $20\ms$ до --- ослабление; оба зависят от NMDA-рецепторов. Отношение $A_-=0.6A_+$ на рисунке --- иллюстрация. Что так в сети вырастают цепочки, прямо не измерено. & \cite{BiPoo1998} & измерено \\
8 & След~\eqref{eq:threefactor}: дофамин, пришедший через $1$--$2\,\text{с}$ после совместной активности, усиливает связь, позже --- нет (утверждение~\ref{prop:threefactor}) & Срезы стриатума, раздельная оптогенетическая стимуляция глутаматных и дофаминовых входов: шипик растёт, только если дофамин приходит через $0.3$--$2\,\text{с}$ после глутаматного входа; окно задаётся быстрым подъёмом и распадом цАМФ. В коре спаривание оставляет раздельные следы для усиления и ослабления, которые моноаминовые рецепторы превращают в долговременные изменения в окне порядка секунд. & \cite{Yagishita2014, He2015eligibility} & измерено \\
9 & Окна длиной в секунды бывают и без дофамина: третий сигнал --- сильное возбуждение получателя & Живая мышь, поле CA1: одно плато-событие в получателе усиливает входы, пришедшие за секунды до и после него, и за один проход создаёт поле места. & \cite{Bittner2017} & измерено \\
10 & Средний шаг правила трёх множителей идёт по градиенту награды~\eqref{eq:perturbation} (утверждение~\ref{prop:gradient}) & Теорема о градиенте для обучения по случайным отклонениям; в главе выведена для малого шума. & \cite{Williams1992} & теория \\
11 & Шум --- это поиск: сеть учится на своих случайных отклонениях & Молодые амадины: выключение ядра, выходящего из цепи базальных ганглиев, резко и обратимо убирает изменчивость песни. Взрослые амадины: помеха подаётся на исполнения слога с высотой по одну сторону от средней, и птицы быстро смещают высоту в другую сторону --- учатся по собственному разбросу. & \cite{Olveczky2005, TumerBrainard2007} & измерено \\
12 & Выбор из двух учится при $\tau_e=1\,\text{с}$ и $\delta=R-\bar R$; не учится при $\tau_e=50\ms$; без вычитания веса растут без предела, а при большой скорости обучения сеть может закрепить худший выбор & \texttt{models/dofa\_learning.py}, $\eta=0.05$, $500$ попыток, $6$ зёрен. $\tau_e=1\,\text{с}$: доля выборов $A$ $0.65$--$0.79\to0.96$--$1.00$, веса $0.85$--$1.33$. $\tau_e=50\ms$: $0.38$--$0.55$, веса не меняются. $\delta=R$: вес победителя $860$--$1190$. $\delta=R$, $\eta=0.5$, $3$ зерна: одно закрепилось на $B$ ($0.04\to0$). Скрипт печатает все четыре случая; пересчёт совпал с главой. & \texttt{models/\allowbreak dofa\_\allowbreak learning.py} & модель \\
13 & Время обучения по одному общему сигналу растёт пропорционально числу шумящих нейронов (утверждение~\ref{prop:broadcastcost}) & Кривые обучения линейной сети: возмущение нейронов медленнее точного градиента во столько раз, сколько выходных нейронов, возмущение весов --- ещё во столько раз, сколько входов. & \cite{Werfel2005} & теория \\
14 & Выходы \nt{ДОФА} идут прежде всего к областям выбора действий & Полная реконструкция аксонов отдельных нигростриатных \nt{ДОФА}-нейронов крысы: ветви одной клетки покрывают $0.45$--$5.7\,\%$ объёма стриатума (в среднем $2.7\,\%$), вне стриатума ветвей почти нет. & \cite{Matsuda2009} & измерено \\
15 & Кора учится предсказывать свой вход, ошибка считается на месте & Обзор: в сенсорной коре есть ответы на рассогласование ожидаемого и пришедшего входа. Что это общее правило обучения коры, не доказано. & \cite{Keller2018} & гипотеза \\
16 & Дофамин ведёт себя как ошибка~\eqref{eq:ctd}: всплеск, перенос на предвестник, провал (рис.~\ref{fig:ctdcases}) & \nt{ДОФА}-нейроны обезьяны: всплеск на неожиданный сок; после обучения --- всплеск на предвестник и нет ответа на обещанный сок; провал, когда сок не пришёл. Непрерывная форма~\eqref{eq:ctd} --- теория. & \cite{Schultz1997, Doya2000} & измерено \\
17 & Схема, вычисляющая $\dd V/\dd t$ и вычитающая ожидание & Мыши, запись и оптогенетика в вентральной покрышке: \nt{ДОФА}-нейроны вычитают ожидание из награды; соседние \nt{ГАМК}-нейроны тормозят их, когда награда ожидается, и их возбуждение или торможение меняет ошибку. Как вычисляется производная, не показано. & \cite{Eshel2015arithmetic} & гипотеза \\
18 & \nt{ДОФА}-нейрон --- ритмоводитель; провалы мельче всплесков & В срезе \nt{ДОФА}-нейроны разряжаются спонтанно и очень регулярно, на медленном пейсмекерном потенциале. У обезьяны частота количественно следует положительной ошибке, а отрицательную передаёт лишь слабо. & \cite{GraceOnn1989, Bayer2005} & измерено \\
19 & Актёр и критик (рис.~\ref{fig:actorcritic}) & Алгоритм --- из теории обучения с подкреплением. У людей (фМРТ) брюшной стриатум ведёт себя как критик и при выборе, и без него, спинной --- только когда надо выбирать действия. & \cite{SuttonBarto2018, ODoherty2004actor} & гипотеза \\
20 & Знак $\delta$ в правиле перевёрнут у нейронов со вторым семейством рецепторов & Срезы стриатума: у нейронов с D1-рецепторами спаривание даёт усиление, которому нужен D1; у нейронов с D2 --- ослабление, которому нужен D2. & \cite{Shen2008} & измерено \\
\bottomrule
\end{longtable}}

\section{Глава~\ref{sec:dofaalt}. Другие теории \nt{DOPA}}

Каждая из расширенных картин опирается на свои прямые измерения дофамина (разброс точек переворота, ответы на неприятное и новое, медленный рост, ответы на смену вкуса), но формулы, которые их объясняют, --- теории, и между двумя из них опыты пока противоречат друг другу.

{\footnotesize
\setlength{\tabcolsep}{3pt}\renewcommand{\arraystretch}{1.15}
\begin{longtable}{>{\raggedright}p{0.5cm}>{\raggedright}p{4.0cm}>{\raggedright}p{5.6cm}>{\raggedright}p{2.3cm}>{\raggedright\arraybackslash}p{1.7cm}}
\toprule
№ & Утверждение & Опыт и что измерено & Источник & Статус \\
\midrule\endfirsthead
\toprule
№ & Утверждение & Опыт и что измерено & Источник & Статус \\
\midrule\endhead
1 & Асимметричное правило~\eqref{eq:asymmetric} сходится к точке~\eqref{eq:expectile}; при $\kappa=1/2$ это среднее, для капли с вероятностью $1/2$ --- $V=\kappa$ (рис.~\ref{fig:expectiles}) & Точка~\eqref{eq:expectile} --- экспектиль, решение задачи наименьших квадратов с разными весами положительных и отрицательных отклонений; для примера главы вывод в самой главе. & \cite{NeweyPowell1987}; вывод в главе & теория \\
2 & У разных \nt{DOPA}-нейронов разные точки переворота, и они упорядочены по асимметрии (утверждение~\ref{prop:expectile}) & Мыши, оптогенетически помеченные \nt{DOPA}-нейроны вентральной покрышки, награды разного размера: точка, где всплеск сменяется провалом, у разных нейронов разная; у нейронов с большей асимметрией ответов она выше; по ответам группы восстанавливается форма распределения наград. Что это главная функция разброса --- гипотеза. & \cite{Dabney2020} & измерено \\
3 & Часть \nt{DOPA}-нейронов отвечает на неприятное всплеском & Обезьяны: одни \nt{DOPA}-нейроны возбуждаются на предвестник награды и тормозятся на предвестник струи воздуха в глаз, другие возбуждаются на оба; группы расположены в разных частях среднего мозга. & \cite{Matsumoto2009, BrombergMartin2010} & измерено \\
4 & Модуль ошибки или новизна задаёт скорость обучения & Прямого опыта, где сигнал важности у \nt{DOPA}-нейронов меняет скорость обучения, в цитируемых работах нет; обзор предлагает роль сигнала «внимание, важно». & \cite{BrombergMartin2010} & гипотеза \\
5 & Отдельный провод для оси опасности & Мыши: \nt{DOPA}-нейроны, идущие к заднему концу стриатума, отвечают на новые и сильные раздражители, а не на награду; их разрушение ослабляет избегание нового предмета. & \cite{Menegas2018} & измерено \\
6 & Оптимальное время действия $\tau_a^*=\sqrt{C/\bar R}$~\eqref{eq:vigor} & Минимум суммы $C/\tau_a+\bar R\tau_a$; вывод в главе. В исходной модели скорость действий задаёт цена времени, равная средней награде. & \cite{Niv2007}; вывод в главе & теория \\
7 & Медленный уровень дофамина несёт $\bar R$ и задаёт скорость действий (утверждение~\ref{prop:multiplex}) & Крысы, микродиализ и вольтамперометрия в прилежащем ядре: уровень дофамина от минуты к минуте следует за частотой наград и за скоростью действий. У людей L-DOPA усиливает зависимость времени реакции от средней частоты наград. Что медленный уровень равен именно $\bar R$, не доказано. & \cite{Hamid2016, Beierholm2013vigor} & гипотеза \\
8 & Медленный уровень складывается из двух источников: выброс меняется у выходов без изменения частоты импульсов, но медленный рост бывает и в самой частоте & Крысы, одна задача: выброс в прилежащем ядре следует за меняющимся ожиданием награды, а частота импульсов опознанных \nt{DOPA}-нейронов вентральной покрышки --- нет. Мыши в виртуальном коридоре (строка~10): плавный рост при приближении к награде есть и в импульсах опознанных \nt{DOPA}-нейронов. & \cite{Mohebi2019, Kim2020} & измерено \\
9 & Дофамин плавно растёт, пока животное приближается к далёкой награде & Крысы в лабиринте, вольтамперометрия в стриатуме: концентрация дофамина нарастает за секунды по мере приближения к цели, крутизна зависит от расстояния и размера награды. & \cite{Howe2013ramp, Hamid2016} & измерено \\
10 & Рост --- ошибка $\delta$ при уточнении оценки~\eqref{eq:ramp}; скачок в коридоре даёт всплеск (рис.~\ref{fig:ramp}) & Формула и число $\delta=0.75\,V$ в секунду --- вывод в главе. Опыт: мыши в виртуальном коридоре, мгновенный перенос ближе к цели; импульсы тел, кальций тел и выходов и концентрация в стриатуме отвечают как ошибка, а не как ожидание $V$. Какое из двух объяснений верно на всех выходах, обсуждается. & \cite{Kim2020} & измерено \\
11 & Взгляд назад: дофамин сообщает меру причинной связи~\eqref{eq:retro} & Мыши, выброс дофамина в прилежащем ядре в опытах, где эта теория и ошибка~\eqref{eq:ctd} предсказывают разное: в ряде опытов выброс следовал первой. & \cite{Jeong2022} & гипотеза \\
12 & Ответ дофамина при обучении постепенно сдвигается от награды к предвестнику & Мыши, сигнал выходов \nt{DOPA}-нейронов в брюшном стриатуме по ходу обучения: всплеск постепенно переходит с момента награды на момент предвестника, как требует обучение по~\eqref{eq:ctd}. & \cite{Amo2022} & измерено \\
13 & \nt{DOPA}-нейроны отвечают на смену вкуса награды при той же ценности & Крысы, запись нейронов вентральной покрышки: замена сока другим, равным по ценности, вызывает ответ, хотя ошибка~\eqref{eq:ctd} равна нулю. & \cite{Takahashi2017} & измерено \\
14 & Искусственные всплески учат связи двух нейтральных раздражителей & Крысы, опыт с предварительным сочетанием двух раздражителей без награды: оптогенетическое включение \nt{DOPA}-нейронов на переходе от первого ко второму создаёт связь, выключение мешает её выучить. & \cite{Sharpe2017} & измерено \\
15 & Вектор ошибок по признакам~\eqref{eq:vectortd} & Теория ошибки предсказания признаков; прямой проверки векторной формы нет. & \cite{Gardner2018} & гипотеза \\
16 & Разные \nt{DOPA}-нейроны в одной задаче несут разные величины & Мыши в виртуальном коридоре, двухфотонная съёмка \nt{DOPA}-нейронов: одни связаны с наградой, другие со скоростью и поворотом, третьи с предвестниками и точностью выбора. & \cite{Engelhard2019} & измерено \\
17 & Жгут вместо провода (утверждение~\ref{prop:bundle}) & Сводка строк~2--16: каждое разложение измерено по отдельности; единой картины, что все они --- части одного сигнала, опытом нет. & --- & гипотеза \\
\bottomrule
\end{longtable}}

\section{Глава~\ref{sec:nora}. \nt{НОРА}}

Устройство \nt{НОРА}-системы, её два режима, связь со зрачком и перевёрнутое U в поведении измерены прямо; то, что усиление переключает режим выбора и работает как обратная температура, --- выводы главы; выгода от подстройки усиления --- расчёт по модели книги; «\nt{НОРА} --- ручка поиска» и «\nt{НОРА} --- прерывание» --- гипотезы.

{\footnotesize
\setlength{\tabcolsep}{3pt}\renewcommand{\arraystretch}{1.15}
\begin{longtable}{>{\raggedright}p{0.5cm}>{\raggedright}p{4.0cm}>{\raggedright}p{5.6cm}>{\raggedright}p{2.3cm}>{\raggedright\arraybackslash}p{1.7cm}}
\toprule
№ & Утверждение & Опыт и что измерено & Источник & Статус \\
\midrule\endfirsthead
\toprule
№ & Утверждение & Опыт и что измерено & Источник & Статус \\
\midrule\endhead
1 & \nt{НОРА}-нейронов у человека несколько десятков тысяч, почти все в одной группе & Серийные срезы ствола мозга человека, окраска на тирозингидроксилазу: $53\,900$ клеток в голубом пятне и ещё $6\,260$ в соседнем подъядре. & \cite{Baker1989LC} & измерено \\
2 & Выходы расходятся почти по всему мозгу, но части группы отчасти обслуживают разные области & Вирусно-генетическое мечение у мышей: \nt{НОРА}-нейроны голубого пятна получают входы из широких областей мозга и шлют выходы в широкие области мозга и спинного мозга. У крыс: группа, идущая к миндалине, усиливает выучивание страха, а отдельная группа, идущая к префронтальной коре, --- его угашение. & \cite{Schwarz2015LC, Uematsu2017LC, Poe2020} & измерено \\
3 & Фоновый режим: ровные импульсы с частотой от долей до нескольких в секунду, меняющейся с бодрствованием & Крысы, запись нейронов голубого пятна в цикле сна и бодрствования: частота наибольшая в бодрствовании, ниже в медленном сне, почти ноль в быстром; изменения частоты опережают смену состояния. & \cite{AstonJonesBloom1981} & измерено \\
4 & Отрывистый режим: короткая пачка на важное для задачи событие, примерно в момент решения & Обезьяны, задача найти редкую цель: все нейроны голубого пятна отвечают пачкой только на цель, через $\sim90\ms$ после неё и за $\sim200\ms$ до ответа рукой; при плохой работе пачки слабее. В задаче с выбором из двух пачка точнее привязана к ответу, чем к стимулу, и есть и при верных, и при ошибочных ответах. & \cite{AstonJones1994vigilance, Clayton2004LC} & измерено \\
5 & Диаметр зрачка --- контрольная точка \nt{НОРА} & Обезьяны: импульсы голубого пятна, вплоть до отдельных импульсов одной клетки, предсказывают расширение зрачка; похожая, но более поздняя связь есть в холмиках и поясной коре. Мыши: колебания зрачка следуют за активностью и норадреналиновых, и холинергических выходов в коре --- зрачок отражает не одну \nt{НОРА}. & \cite{Joshi2016, Reimer2016} & измерено \\
6 & \nt{НОРА} меняет крутизну передаточной характеристики~\eqref{eq:gain} & Бодрствующие обезьяны, слуховая кора, ионофорез норадреналина: фоновая активность нейронов подавляется сильнее, чем ответ на голоса сородичей, --- отношение сигнала к шуму растёт. Мультипликативная сигмоида~\eqref{eq:gain} --- упрощение, взятое из сетевой модели. & \cite{Foote1975NE, ServanSchreiber1990} & измерено \\
7 & Усиление повышает $d'$ только против шума после усилителя~\eqref{eq:gainsnr} (утверждение~\ref{prop:gainnoise}) & Вывод в главе; в сетевой модели усиление улучшает обнаружение сигнала. Опыта, который отделял бы шум до и после усилителя и проверял бы это различие, нет. & вывод в главе; \cite{ServanSchreiber1990} & теория \\
8 & Граница $g\beta=1$ переключает сеть выбора из «размышляет» в «решила»~\eqref{eq:wtagain} (утверждение~\ref{prop:gainwta}, рис.~\ref{fig:pitchfork}) & Линейный анализ двух тормозящих друг друга групп; вывод в главе. Прямых измерений этого порога в мозге нет. & вывод в главе & теория \\
9 & Коэффициент усиления --- обратная температура выбора~\eqref{eq:softmax} & При логистическом шуме после усилителя вероятность выбора --- softmax с $T=\sigma/g$; вывод в главе. & вывод в главе & теория \\
10 & \nt{НОРА} задаёт, сколько искать & Люди: перед выбором наугад базовый зрачок шире, чем перед выбором лучшего известного; у людей с более широким зрачком склонность искать выше. Крысы: переход в режим случайного выбора управляется входом \nt{НОРА} в переднюю поясную кору. Направление: больше тонической \nt{НОРА} --- больше поиска, то есть меньше действующее $g$; отождествление \nt{НОРА} с $g$ держится на различии фонового и отрывистого режимов. & \cite{JepmaNieuwenhuis2011, Tervo2014stochastic, AstonJones2005} & гипотеза \\
11 & Награда зависит от $g$ как перевёрнутое U; лучшее $g$ меньше в быстро меняющемся мире (утверждение~\ref{prop:explore}, рис.~\ref{fig:invertedU}) & \texttt{models/nora\_explore.py}, $60$ прогонов по $4000$ попыток. Перемена раз в $1000$ попыток: лучшее $g=16$, доля $0.976$; раз в $20$: $g=8$, доля $0.886$; при $g=0.5$ --- $0.77$. Пересчёт совпал с главой. & \texttt{models/\allowbreak nora\_\allowbreak explore.py} & модель \\
12 & Перевёрнутое U в поведении: лучше всего при умеренной фоновой активности с чёткими пачками & Обезьяны, та же задача, что в строке~4: в периоды хорошей работы фоновая частота умеренная, а пачки на цель чёткие; при высокой фоновой частоте пачек нет и ложных ответов больше. Изменения фоновой и вызванной активности следуют за колебаниями качества работы. & \cite{Usher1999LC, AstonJones2005} & измерено \\
13 & \nt{НОРА} сообщает неожиданную неопределённость, \nt{АЦЕТ} --- ожидаемую & Теория байесовского вывода с двумя видами неопределённости; прямого опыта, разделяющего эти сигналы по медиаторам, в цитируемых работах нет. & \cite{YuDayan2005} & гипотеза \\
14 & После резкой перемены люди учатся быстрее, и зрачок расширяется & Люди, задача предсказания с внезапными переменами: короткие изменения зрачка следуют за оценкой вероятности перемены и вместе с базовым зрачком предсказывают, насколько сильно новые данные меняют оценку (скорость обучения); изменение зрачка, не связанное со светом и задачей, меняет эту скорость. Что зрачок здесь показывает именно \nt{НОРА}, --- вывод по косвенным данным строки~5. & \cite{Nassar2012} & измерено \\
15 & Отрывистая пачка работает как прерывание: сеть сбрасывает состояние & Прямого опыта, где пачка \nt{НОРА} сбрасывает состояние сети, нет; измерены только пачки на важные события (строка~4). & \cite{BouretSara2005} & гипотеза \\
16 & Подстройка $g$ или скорости обучения по неожиданности почти не лучше лучших постоянных настроек & \texttt{models/nora\_explore.py}, мир с эпохами по $1000$ попыток (перемена раз в $1000$ и раз в $20$). Лучшее постоянное $g=8$: $0.925$; подстройка $g$: до $0.927$; подстройка скорости обучения: до $0.926$; постоянная скорость $0.4$: $0.928$. Пересчёт совпал с главой. & \texttt{models/\allowbreak nora\_\allowbreak explore.py} & модель \\
\bottomrule
\end{longtable}}

\section{Глава~\ref{sec:acet}. Добавляем \nt{АЦЕТ}}

Три действия ацетилхолина измерены на срезах и в живом мозге, конфликт записи и чтения показан на модели книги, а то, что \nt{АЦЕТ} служит линией разрешения записи и сообщает ожидаемую неопределённость, --- гипотеза с частичной опорой на опыт.

{\footnotesize
\setlength{\tabcolsep}{3pt}\renewcommand{\arraystretch}{1.15}
\begin{longtable}{>{\raggedright}p{0.5cm}>{\raggedright}p{4.0cm}>{\raggedright}p{5.6cm}>{\raggedright}p{2.3cm}>{\raggedright\arraybackslash}p{1.7cm}}
\toprule
№ & Утверждение & Опыт и что измерено & Источник & Статус \\
\midrule\endfirsthead
\toprule
№ & Утверждение & Опыт и что измерено & Источник & Статус \\
\midrule\endhead
1 & \nt{АЦЕТ}-нейронов мозга немного, они собраны в нескольких группах, а выходы расходятся по всей коре: широковещательный канал & Окраска антителами к ферменту синтеза ацетилхолина и ретроградная метка у крысы: кора, гиппокамп, миндалина, обонятельная луковица и таламус получают ацетилхолин в основном от шести компактных групп клеток базального переднего мозга и ствола & \cite{Mesulam1983} & измерено \\
2 & Уровень ацетилхолина в коре высок при бодрствовании и низок в медленном сне & Микродиализ в коре и гиппокампе свободно движущихся кошек с записью ЭЭГ: выброс ацетилхолина в спокойном бодрствовании на $\sim100\%$, в активном --- на $\sim175\%$ выше, чем в медленном сне; в быстром сне в коре --- как при бодрствовании & \cite{Marrosu1995,Hasselmo1999} & измерено \\
3 & Смысл сигнала задаёт рецептор: у ацетилхолина есть быстрые рецепторы-каналы и медленные, запускающие реакции в клетке (утверждение~\ref{prop:receptor}) & Обзор измерений: действие ацетилхолина на возбудимость, передачу и пластичность зависит от места выброса, подтипа рецептора (никотиновые --- ионные каналы, мускариновые --- через G-белки) и клетки-получателя & \cite{Picciotto2012} & измерено \\
4 & Первое действие: ослабить внутренние связи, почти не трогая входы извне (множитель $1-a$ в~\eqref{eq:acetnet}) & Срезы обонятельной коры крысы: при $100$\,мкМ агонистов ацетилхолина потенциалы внутренних связей (слой Ib) падают более чем на $60\%$, внешнего входа (слой Ia) --- менее чем на $15\%$, в той же клетке; подавление пресинаптическое. Срезы гиппокампа (CA1): $89\%$ против $40\%$ для внутреннего и входного путей & \cite{HasselmoBower1992,HasselmoSchnell1994} & измерено \\
5 & Второе действие: усилить вход (множитель $\frac{1+a}{2}$) & Срезы неокортекса: никотиновые рецепторы усиливают синапсы входа из таламуса и не действуют на внутрикорковые; мускариновые ослабляют оба типа. Возбудимость нейронов ацетилхолин повышает (обзор) & \cite{Gil1997,Hasselmo2006} & измерено \\
6 & Третье действие: облегчить изменение весов (скорость $\eta a$) & Крыса: электрическая стимуляция базального ядра (источника ацетилхолина для коры) в паре с тоном вызывает у взрослого животного сильную перестройку карты слуховой коры под предъявляемый звук & \cite{KilgardMerzenich1998,Hasselmo2006} & измерено \\
7 & Правило Хебба для редких образов во втором уравнении~\eqref{eq:acetnet} даёт ассоциативную память, дописывающую образ по части & Расчёт ёмкости сети с разреженными образами и ковариационным правилом: при малой доле активных нейронов $p$ сеть хранит заметно больше образов, чем при $p=1/2$ & \cite{Tsodyks1988} & теория \\
8 & Запись при низком $a$ портит память: сеть записывает вспомненное, а не увиденное; при $a=0.1$ порча не слабее, чем при $a=0$ & \texttt{models/\allowbreak acet\_memory.py}: $200$ нейронов, пять образов по $20$, новый образ делит с одним из них $10$ нейронов, $10$ прогонов. При $a=0$ новый образ не запомнен, старый тянет $5.9$ чужих нейронов из $10$; при $a=0.1$ новый образ вспоминается ($0.97$), но оба сливаются: новый тянет $7.3$ чужих нейрона, старый --- $7.9$; с $a=0.2$ --- меньше одного & вывод в главе & модель \\
9 & Утверждение~\ref{prop:writeread}: нет уровня $a$, при котором одинаково хороши запись и чтение (рис.~\ref{fig:acetsim}) & Тот же скрипт, скорость обучения $\eta a$: новый образ запоминается за одно предъявление целиком при $a\ge0.9$, на $0.8$ при $a=0.75$, не запоминается при $a\le0.5$. Чтение по половине образа: $1.0$ при $a\le0.2$, $0.96$ при $a=0.3$, $0.04$ при $a=0.4$ & вывод в главе & модель \\
10 & В мозге один широковещательный провод --- \nt{АЦЕТ} --- переключает запись и чтение & Идея взята из моделей, построенных по измерениям на срезах. Самый прямой опыт --- у человека: блокатор мускариновых рецепторов скополамин ухудшает заучивание новых пар слов, сильнее всего пересекающихся со старыми, но не вспоминание уже выученных пар. Прямого измерения двух режимов сети нет & \cite{HasselmoAndersonBower1992,Atri2004} & гипотеза \\
11 & Фильтр~\eqref{eq:kalman} оптимален; вес входа и скорость обучения --- одна величина $P/R$, в установившемся режиме $P=\sqrt{qR}$ и память $\sqrt{R/q}$ (утверждение~\ref{prop:kalman}) & Опыта не требует: оптимальность фильтра Калмана--Бьюси --- классический результат теории оценивания; установившийся режим выведен в главе & \cite{KalmanBucy1961} & теория \\
12 & \nt{АЦЕТ} сообщает сети величину, подобную $P$, --- ожидаемую неопределённость & Предложено в вероятностной модели внимания. Прямого измерения того, что уровень ацетилхолина следует за неопределённостью оценки, нет & \cite{YuDayan2005} & гипотеза \\
13 & Часть \nt{АЦЕТ}-нейронов отвечает на награду и наказание за два десятка миллисекунд, тем сильнее, чем неожиданнее подкрепление & Мышь, оптогенетически опознанные холинергические нейроны базального переднего мозга: ответ на награду и наказание через $18\pm3$\,мс, величина растёт с неожиданностью подкрепления, ответы схожи у нейронов двух ядер & \cite{Hangya2015} & измерено \\
14 & \nt{НОРА} замечает неожиданную перемену мира, \nt{АЦЕТ} держит сеть в режиме записи & Разделение труда предложено в модели. Косвенно: у человека диаметр зрачка, связанный с \nt{НОРА}, отслеживает перемены и скорость обучения. Прямой проверки пары \nt{НОРА}--\nt{АЦЕТ} нет & \cite{YuDayan2005,Nassar2012} & гипотеза \\
15 & В медленном сне сеть в режиме чтения проигрывает записанное днём, и эти повторы переписывают его в долговременную память & Записи ансамблей нейронов гиппокампа крысы: клетки, работавшие вместе днём, чаще срабатывают вместе в последующем медленном сне и в том же порядке --- измерено. Для перезаписи: подавление рябей гиппокампа после обучения ухудшает пространственную память, а удлинение спонтанных рябей её улучшает; что причина --- низкий уровень \nt{АЦЕТ}, не доказано & \cite{WilsonMcNaughton1994,Skaggs1996,Girardeau2009,FernandezRuiz2019,Hasselmo1999} & гипотеза \\
\bottomrule
\end{longtable}}

\section{Глава~\ref{sec:synthesis}. Вся машина}

Глава-сводка новых опытов не приводит: каждая её строка опирается на главу, где утверждение разобрано, и на те же источники; твёрдо установлены шина \nt{ГЛУТ} и управление потоком через \nt{ГАМК}, а роли широковещательных каналов и их совместная работа --- в основном гипотеза.

{\footnotesize
\setlength{\tabcolsep}{3pt}\renewcommand{\arraystretch}{1.15}
\begin{longtable}{>{\raggedright}p{0.5cm}>{\raggedright}p{4.0cm}>{\raggedright}p{5.6cm}>{\raggedright}p{2.3cm}>{\raggedright\arraybackslash}p{1.7cm}}
\toprule
№ & Утверждение & Опыт и что измерено & Источник & Статус \\
\midrule\endfirsthead
\toprule
№ & Утверждение & Опыт и что измерено & Источник & Статус \\
\midrule\endhead
1 & На $8.6\cdot10^{10}$ нейронов приходится всего четыре рода управляющих сигналов~\eqref{eq:machine} & Подсчёт ядер клеток в гомогенате мозга (изотропный фракционатор): у взрослого мужчины $86.1\pm8.1$ млрд нейронов & \cite{Azevedo2009} & измерено \\
2 & Утверждение~\ref{prop:architecture}, первые два слоя: данные идут по адресной шине \nt{ГЛУТ}, знак связи задаёт отправитель, поток направляют и нормируют \nt{ГАМК}-нейроны (главы~\ref{sec:neurontypes}, \ref{sec:gaba}) & Один главный медиатор на нейрон (обзор измерений); прямое торможение сужает окно, в котором пирамидный нейрон складывает входы; деление на общую активность измерено во многих областях & \cite{Strata1999,Pouille2001,Carandini2012} & измерено \\
3 & Элементы ненадёжны: синапс теряет большую часть импульсов (табл.~\ref{tab:computer}, глава~\ref{sec:code}) & Срезы гиппокампа, минимальная стимуляция: больше половины импульсов не вызывают ответа в CA1; сбои порога и проведения по аксону исключены, причина --- вероятностный выброс медиатора & \cite{AllenStevens1994} & измерено \\
4 & Мозг работает на $\sim20$ ваттах, в среднем около импульса в секунду на нейрон (глава~\ref{sec:code}); такой машины инженеры пока не построили & Оценка~\eqref{eq:energy} из измеренных затрат: около $10^9$ молекул АТФ на импульс вместе с работой синапсов (кора грызунов); подробная оценка для коры даёт ещё меньшую частоту. Состояние техники --- по обзору & \cite{Attwell2001,Lennie2003,MehonicKenyon2022} & теория \\
5 & Обучение большинства синапсов --- произведение местного следа и одного общего числа $\delta$ (второе уравнение~\eqref{eq:machine}, глава~\ref{sec:dofa}) & \nt{ДОФА}-нейроны обезьяны отвечают на ошибку предсказания награды; в срезах полосатого тела дофамин увеличивает шипик, только если приходит через $0.3$--$2$\,с после глутаматного входа --- след измерен & \cite{Schultz1997,Yagishita2014} & измерено \\
6 & Отдельный провод ошибки к каждому выходу есть лишь в цепи обучения движениям (глава~\ref{sec:motorlearning}) & Мозжечок: совпадение сигнала лазящего волокна со входом ослабляет этот вход; у обезьяны сложные импульсы клеток Пуркинье несут направление ошибки движения и вызывают поправку & \cite{Ito2001,Herzfeld2018} & измерено \\
7 & Каждый широковещательный канал --- жгут из нескольких линий (утверждение~\ref{prop:bundle}) & Для \nt{ДОФА}: в одной задаче разные нейроны несут награду, движение и признаки стимула; быстрая ошибка и медленный уровень дофамина различны. Для \nt{СЕРО}, \nt{НОРА}, \nt{АЦЕТ} такая раскладка изучена хуже & \cite{Engelhard2019,Hamid2016,Mohebi2019} & измерено \\
8 & \nt{СЕРО}: регулятор уровня и, по гипотезе, горизонт $\tau_\gamma$ (глава~\ref{sec:serocomputer}) & Самоторможение: агонисты собственных рецепторов 5-HT$_{1A}$ тормозят серотониновые нейроны шва --- измерено. Горизонт предложен в теории; сильнейший опыт --- оптогенетическая активация этих нейронов у мыши удлиняет ожидание отложенной награды & \cite{Sprouse1987,Doya2002,Miyazaki2014} & гипотеза \\
9 & \nt{НОРА}: усиление $g$ выбирает между поиском и решением (глава~\ref{sec:nora}) & Рост отношения сигнала к шуму: у бодрствующей обезьяны норадреналин в слуховой коре подавляет фоновую активность сильнее, чем ответ на голоса сородичей --- измерено; мультипликативное усиление --- модель; роль в выборе между поиском и решением --- теория & \cite{Foote1975NE,ServanSchreiber1990,AstonJones2005} & гипотеза \\
10 & \nt{АЦЕТ}: переключает запись и чтение (глава~\ref{sec:acet}) & Три действия (ослабление внутренних связей, усиление входа, облегчение пластичности) измерены; переключение режимов косвенно поддержано опытом со скополамином у человека & \cite{HasselmoBower1992,Gil1997,Atri2004} & гипотеза \\
11 & Формула~\eqref{eq:machine} описывает машину целиком & Это сводка моделей глав~\ref{sec:glutcomputer}--\ref{sec:acet}, каждая часть опирается на свою главу; как целое опытом не проверялась & вывод в главе & гипотеза \\
12 & Ручки работают согласованно без общего управления: ошибка, неожиданность, запись, возврат (рис.~\ref{fig:timeline}) & Опыта нет: схема качественная. Отдельные звенья --- теория разделения труда \nt{НОРА} и \nt{АЦЕТ} и связь зрачка со скоростью обучения у человека & \cite{YuDayan2005,Nassar2012} & гипотеза \\
13 & Обучение по одному общему сигналу медленнее, чем больше нейронов влияют на результат (утверждение~\ref{prop:broadcastcost}) & Расчёт кривых обучения линейной сети: обучение по возмущению выходов медленнее градиентного спуска во столько раз, сколько выходов & \cite{Werfel2005} & теория \\
14 & Как кора настраивает многослойные цепи, неизвестно; кандидаты --- пачки импульсов, вычисления в ветвях нейрона, самообучение предсказанием & Пачки как носитель ошибки --- модель; повторить ответ подробной модели нейрона коры смогла лишь сеть в $5$--$8$ слоёв --- модель; кальциевые импульсы в ветвях нейронов коры человека, дающие исключающее ИЛИ, --- измерено в срезах; самообучение --- обзор свидетельств & \cite{Payeur2021,Beniaguev2021,Gidon2020,Keller2018} & гипотеза \\
\bottomrule
\end{longtable}}

\section{Глава~\ref{sec:sensors}. Входы машины}

Устройство датчиков измерено прямо --- записями тока одиночных клеток, колебаний улитки и подсчётом клеток сетчатки; предел чувствительности и формула дробового шума следуют из физики, а то, что коды датчиков подстроены под наибольшую информацию, --- гипотеза с сильной опорой.

{\footnotesize
\setlength{\tabcolsep}{3pt}\renewcommand{\arraystretch}{1.15}
\begin{longtable}{>{\raggedright}p{0.5cm}>{\raggedright}p{4.0cm}>{\raggedright}p{5.6cm}>{\raggedright}p{2.3cm}>{\raggedright\arraybackslash}p{1.7cm}}
\toprule
№ & Утверждение & Опыт и что измерено & Источник & Статус \\
\midrule\endfirsthead
\toprule
№ & Утверждение & Опыт и что измерено & Источник & Статус \\
\midrule\endhead
1 & Датчик --- преобразователь: рецептор даёт ток, а выход чувствительных клеток глаза и уха --- глутамат на шину \nt{ГЛУТ}; первые клетки импульсов не выдают (рис.~\ref{fig:transducer}) & Палочки и колбочки черепахи выбрасывают возбуждающую аминокислоту: её ловят глутаматные каналы на отсечённом кусочке мембраны. Внутренние волосковые клетки улитки крысы: токи в окончаниях нервных волокон идут через глутаматные AMPA-рецепторы, частота выбросов растёт с плавной деполяризацией клетки до $150$\,Гц & \cite{CopenhagenJahr1989,GlowatzkiFuchs2002} & измерено \\
2 & Датчик света в темноте отвечает на один фотон током около пикоампера & Всасывающий электрод на наружном сегменте палочки жабы: ответы на тусклые вспышки квантованы, событие $\approx1$\,пА ($5\%$ темнового тока) с разбросом $0.2$\,пА, эффективность $0.5$. Человек сообщает о попадании одного фотона чаще случайного & \cite{Baylor1979,Tinsley2016} & измерено \\
3 & Датчик звука замечает смещение меньше нанометра & Метод Мёссбауэра, основная мембрана улитки морской свинки на месте $16$--$19$\,кГц: на пороге ответа слухового нерва скорость мембраны около $0.04$\,мм/с, то есть смещение $\approx0.35$\,нм (пересчёт наш). Мерили мембрану, а не пучок датчика & \cite{Sellick1982,Hudspeth1989} & измерено \\
4 & Точность в темноте ограничена дробовым шумом фотонов~\eqref{eq:photonnoise}; при слабом свете накопление длиннее & Формула следует из пуассоновской статистики. В палочке шум при тусклом свете совпадает с суммой независимых квантовых событий; на ярком фоне ответ на вспышку меньше и короче. Число «десятые доли секунды» здесь отдельно не проверено & \cite{Baylor1979} & теория \\
5 & Диапазон входа --- десять порядков для света и двенадцать для звука, а частота импульсов --- меньше трёх порядков & Для звука: ответ основной мембраны на характеристической частоте растёт с крутизной до $0.2$\,дБ/дБ в полосе $40$--$80$\,дБ, сжатие видно и выше $100$\,дБ. Сами числа порядков --- стандартные оценки, в главе без источника & \cite{Ruggero1997} & измерено \\
6 & Ответ датчика~\eqref{eq:nakarushton}, а $\sigma$ следует за средней яркостью, сдвигая кривую по логарифмической оси (рис.~\ref{fig:adaptation}) & Формула впервые подобрана к S-потенциалам сетчатки рыбы. Колбочки черепахи: ответы подчиняются $V=I V_m/(I+\sigma)$ и охватывают $\sim3.5$ порядка яркости; после $2$--$3$ минут фона кривая сдвигается по горизонтали, сохраняя ширину. Связь с делением на общую активность --- обзор & \cite{NakaRushton1966,NormannPerlman1979,Carandini2012} & измерено \\
7 & Датчики передают перемены, а их коды подстроены под наибольшую информацию на импульс (утверждение~\ref{prop:sensors}) & Принцип предложен теоретически. Сильнейшее свидетельство: у мухи характеристика «контраст $\to$ ответ» нейронов первого слоя совпадает с накопленным распределением контрастов природных сцен --- так достигается наибольшая информация & \cite{Barlow1961,Laughlin1981} & гипотеза \\
8 & Включающие и выключающие датчики --- двухпроводный код; событийная камера повторяет его в кремнии & Обзор опытов: включающий и выключающий каналы сетчатки разделяются уже на первом синапсе и выключаются по отдельности. Камера: $128\times128$ пикселей без кадров, диапазон $120$\,дБ, задержка $15$\,мкс & \cite{Schiller1992,Lichtsteiner2008,Gallego2022} & измерено \\
9 & В глазу $\sim10^8$ датчиков света и $\sim10^6$ выходных нейронов: сжатие в $\sim100$ раз & Подсчёт на целых препаратах сетчатки человека: в среднем $92$ млн палочек и $4.6$ млн колбочек; выходных (ганглиозных) клеток от $0.7$ до $1.5$ млн в разных глазах & \cite{Curcio1990,CurcioAllen1990} & измерено \\
10 & У мыши выходных каналов глаза больше тридцати & Мышь: двухфотонная кальциевая съёмка более $11\,000$ выходных клеток и кластеризация ответов --- заметно больше $30$ функциональных типов. Для человека число не измерено & \cite{Baden2016} & измерено \\
11 & Глаз морской свинки передаёт $\sim875\,000$ бит/с; пересчёт на глаз человека даёт $\sim10^7$ бит/с & Морская свинка, многоэлектродная матрица, движение в природных сценах: $6$--$13$ бит/с на клетку, $\sim875\,000$ бит/с на $\sim10^5$ выходных клеток. Для человека ($\sim10^6$ клеток) $\sim10^7$ бит/с --- пересчёт авторов & \cite{Koch2006} & измерено \\
12 & Ухо --- банк полосовых фильтров с почти логарифмической картой частот (рис.~\ref{fig:filterbank}) & Место наибольшего колебания основной мембраны зависит от частоты; функция «частота--место» почти экспоненциальна, одной формой описывает улитки человека и живых животных & \cite{Greenwood1990,Robles2001} & измерено \\
13 & Резонаторы активны: часть датчиков усиливает слабые колебания в тысячи раз по амплитуде ($66$--$76$\,дБ), с ростом громкости усиление падает & Лазерная виброметрия основания улитки шиншиллы: при слабых тонах на характеристической частоте усиление $66$--$76$\,дБ относительно стремечка; смерть снижает чувствительность на $60$--$81$\,дБ (в $10^3$--$10^4$ раз по амплитуде) и убирает сжатие. Источник силы --- наружные волосковые клетки & \cite{Ruggero1997,Robles2001} & измерено \\
14 & Усилитель уха стоит у границы самовозбуждения & Расчёт: у точки бифуркации Хопфа неизбежны сжатие диапазона, острая настройка и комбинационные тоны --- все известные нелинейности слуха. Что улитка настроена именно так, прямо не доказано & \cite{Eguiluz2000} & гипотеза \\
15 & На низких частотах импульсы идут в фазе со звуком (у морской свинки привязка слабеет от $\sim600$\,Гц и заметна до $\sim3.5$\,кГц), и по ним находится направление; на высоких остаётся код местом & Слуховой нерв морской свинки: привязка к фазе начинает падать около $600$\,Гц и исчезает выше $3.5$\,кГц; у кошки и обезьяны граница примерно на октаву выше. Разница времени прихода в два уха --- обзор & \cite{PalmerRussell1986,Grothe2010} & измерено \\
16 & Остальные датчики (табл.~\ref{tab:sensors}): у обоняния $\sim400$ типов рецепторов, по одному на нейрон, и комбинаторный код; вкус частично через АТФ и серотонин; осязание --- быстро и медленно привыкающие датчики; тепло и холод --- отдельные & Мышь: один рецептор узнаёт много запахов, один запах --- много рецепторов, разные запахи --- разные наборы. Без рецепторов АТФ нерв вкуса не отвечает; вкусовые почки выбрасывают серотонин. Запись одиночных волокон кожи руки человека: четыре типа по скорости привыкания. Холодовой канал отличен от тепловых & \cite{Mombaerts2004,Malnic1999,Finger2005,Huang2005,JohanssonVallbo1979,McKemy2002} & измерено \\
\bottomrule
\end{longtable}}

\section{Глава~\ref{sec:recognition}. Распознавание}

Скорость распознавания, устройство первых ступеней, линейное чтение ответа и обучение на непрерывности времени измерены на человеке и обезьяне; сведение всей цепочки к двум операциям и обучение на видео проверены на моделях книги, а объяснения иллюзий --- от хорошо обоснованных до спорных.

{\footnotesize
\setlength{\tabcolsep}{3pt}\renewcommand{\arraystretch}{1.15}
\begin{longtable}{>{\raggedright}p{0.5cm}>{\raggedright}p{4.0cm}>{\raggedright}p{5.6cm}>{\raggedright}p{2.3cm}>{\raggedright\arraybackslash}p{1.7cm}}
\toprule
№ & Утверждение & Опыт и что измерено & Источник & Статус \\
\midrule\endfirsthead
\toprule
№ & Утверждение & Опыт и что измерено & Источник & Статус \\
\midrule\endhead
1 & Сравнение с образцом по пикселям при сдвиге угадывает не лучше монетки; пара ступеней~\eqref{eq:scstage} узнаёт фигуру в любом месте & \texttt{models/\allowbreak recognition\_models.py}, «сетчатка» из $24$ точек, $4000$ проб: образец --- $0.49$, $0.51$, $0.52$ при доле перевёрнутых точек $0$, $0.1$, $0.2$; пара $S$ и $C$ --- $1.00$, $0.86$, $0.70$ & вывод в главе & модель \\
2 & Животное на картинке узнаётся за $\sim150$\,мс, один проход по десятку ступеней по $10$--$20$\,мс & Человек, задача «есть животное или нет» на фотографиях, показанных на $20$\,мс: вызванные потенциалы для двух ответов расходятся через $\sim150$\,мс. Обезьяна: время первого ответа растёт от ступени к ступени (V1 раньше всех, затем V2 и V4). Число ступеней и «ноль или один импульс на ступень» --- вывод из этих чисел & \cite{Thorpe1996,Schmolesky1998} & измерено \\
3 & Ступень $S$ --- И по частям, ступень $C$ --- наибольшее по положениям (утверждение~\ref{prop:conveyor}) & Кошка, первая зрительная кора: простые клетки отвечают на край определённой ориентации в определённом месте, сложные --- на ту же ориентацию в большей области. Внутриклеточная запись сложных клеток: ответ на две полоски у многих клеток близок к большему из двух ответов, в среднем ближе всего к операции max. Наибольшее через взаимное торможение --- утверждение~\ref{prop:wta}, деление на общую активность --- обзор & \cite{HubelWiesel1962,Lampl2004,Carandini2012} & измерено \\
4 & Выше по цепочке сочетания крупнее и сложнее, ответ всё меньше зависит от положения & Обезьяна, упрощение стимулов до «критического признака»: клетки V4 и задней нижневисочной коры отвечают на сложные сочетания формы, цвета и текстуры при малом поле; в передней нижневисочной коре поле больше. Чередование $S$ и $C$ в модели воспроизводит эту картину & \cite{KobatakeTanaka1994,Riesenhuber1999} & измерено \\
5 & Свёрточные сети повторяют это чередование и лучше всех предсказывают ответы верхних ступеней; предмет читается линейно по частотам группы нейронов & Сеть, обученная распознавать, без подгонки к мозгу: верхний слой предсказывает ответы нейронов нижневисочной коры обезьяны, средние --- ответы V4. Линейный классификатор по активности $\sim100$ случайно выбранных клеток за окна от $12.5$\,мс узнаёт предмет и категорию при разном положении и размере & \cite{Fukushima1980,Yamins2014,Hung2005} & измерено \\
6 & Правило Хебба со следом~\eqref{eq:traceinvariance} учит инвариантности на видео без подписей, если след не короче $\sim20$\,мс & Идея медленных признаков и правила со следом --- теория. \texttt{models/\allowbreak recognition\_models.py}, два $C$-нейрона, $16$ входов, $10$ прогонов: пауза между предметами $0.3$\,с. При $\tau_{\text{сл}}=5$\,мс совпадение с фигурой $0.52$ в среднем, при $10$\,мс --- $0.56$; при $20$, $50$ и $200$\,мс --- $1.00$ во всех прогонах (при $30$\,мс один прогон из десяти ошибается) & \cite{Foldiak1991,WiskottSejnowski2002} & модель \\
7 & В мозге инвариантность учится на непрерывности времени & Обезьяна: предмет незаметно подменяли во время скачка глаза в определённое место поля; позиционная инвариантность нейронов нижневисочной коры менялась в соответствии с подменой, значимо уже через час. Что это главное правило обучения зрения --- гипотеза & \cite{LiDiCarlo2008} & измерено \\
8 & Движение: детекторы направления, затем та же пара И/ИЛИ по большим областям & Детекторы направления в сетчатке кролика; в области MT обезьяны все $110$ записанных нейронов избирательны к направлению, ответ на движение сильнее, чем в V1 & \cite{Barlow1965,Albright1984} & измерено \\
9 & Верхние ступени посылают вниз предсказание, вверх идёт ошибка & Модель объясняет внеклассические эффекты полей первой зрительной коры; отдельные наблюдения согласуются (обзор), как общее устройство конвейера не доказано & \cite{RaoBallard1999,Keller2018} & гипотеза \\
10 & Трудные картинки требуют обратных связей и нескольких кругов & Обезьяна: для «трудных» картинок решение в нижневисочной коре складывается на $\sim30$\,мс позже; эти поздние ответы плохо предсказывает прямая сеть и лучше --- сети с обратными связями или более глубокие & \cite{Kar2019} & измерено \\
11 & Незаметная подделка пикселей обманывает сеть; зрение человека гораздо устойчивее, но не неуязвимо & На сетях: малое специально подобранное возмущение меняет ответ; пример «панда $\to$ гиббон» --- из второй работы. У человека при коротком показе подделки, хорошо переносимые между сетями, сдвигают выбор & \cite{Szegedy2014,Goodfellow2015,Elsayed2018} & измерено \\
12 & Иллюзии ожидания: ненадёжный вход уступает ожиданию --- бледное движется «медленнее», платье видят по-разному (раздел~\ref{sec:illusions}) & Решётки меньшего контраста кажутся движущимися медленнее. Опрос $1401$ человека: платье видят в трёх устойчивых вариантах, подсказка об освещении переключает цвет; у $\sim13\,000$ наблюдателей решает предположение об освещении. Байесовская модель с ожиданием медленных движений объясняет ряд иллюзий движения & \cite{Thompson1982,LaferSousa2015,Wallisch2017,Weiss2002,Purves2011} & гипотеза \\
13 & Регулировка усиления: водопад, наклон и контраст; решётка Германа --- не торможение соседей & Детекторы направления сетчатки кролика при долгом движении снижают частоту, а после остановки падают ниже фона. Иллюзия наклона воспроизводится моделью с делением на активность соседей. Изменения решётки, не меняющие ответ выходных нейронов сетчатки, убирают пятна --- объяснение соседним торможением отвергнуто & \cite{BarlowHill1963,Schwartz2009,SchillerCarvey2005} & измерено \\
14 & Компенсация задержек: движущийся предмет кажется впереди вспышки & Иллюзия измерена. Психофизика противоречит и продлению движения вперёд, и разнице задержек; предложено объяснение задним числом --- по событиям $\sim80$\,мс после вспышки. Спор не решён & \cite{Nijhawan1994,EaglemanSejnowski2000} & гипотеза \\
15 & Неподвижные «змеи» вращаются: разница задержек~\eqref{eq:latency} читается детекторами направления & Иллюзия работает и без цвета; пары соседних элементов порождают её поодиночке; нейроны коры обезьяны, избирательные к направлению, дают направленный ответ на эти неподвижные пары. У человека вращение запускают микроскачки глаз и моргания & \cite{Conway2005,OteroMillan2012} & измерено \\
16 & Двойственные картинки: взаимное торможение, усталость и шум; смены вызывает главным образом шум & Модель конкурирующих групп нейронов: смены в ней вызывает \emph{шум}, без шума их нет; она объясняет рост частоты смен с силой стимула. Усталость здесь играет меньшую роль & \cite{MorenoBote2007} & гипотеза \\
17 & Часть иллюзий --- следствие задачи, которой учится машина & Сеть, обученная предсказывать следующий кадр видео, «видит» вращение неподвижных колец и не видит его на контрольных картинках; сети для очистки и повышения резкости изображений поддаются иллюзиям цвета и яркости, как человек & \cite{Watanabe2018,GomezVilla2020} & измерено \\
\bottomrule
\end{longtable}}

\section{Глава~\ref{sec:muscles}. Выход на мышцы}

Устройство выхода --- двигательные единицы, запас надёжности синапса, включение по размеру, пара мышц, петля растяжения и генераторы ритма --- измерено прямо; условие устойчивости петли с задержкой --- теорема; внутренняя модель тела --- хорошо обоснованная гипотеза; числа рисунков --- расчёт по модели книги.

{\footnotesize
\setlength{\tabcolsep}{3pt}\renewcommand{\arraystretch}{1.15}
\begin{longtable}{>{\raggedright}p{0.5cm}>{\raggedright}p{4.0cm}>{\raggedright}p{5.6cm}>{\raggedright}p{2.3cm}>{\raggedright\arraybackslash}p{1.7cm}}
\toprule
№ & Утверждение & Опыт и что измерено & Источник & Статус \\
\midrule\endfirsthead
\toprule
№ & Утверждение & Опыт и что измерено & Источник & Статус \\
\midrule\endhead
1 & Двигательная единица: один \nt{АЦЕТ}-нейрон управляет группой волокон, от нескольких сотен до пары тысяч; в мышцах тонкой подстройки единицы мельче & Мышцы человека, подсчёт двигательных аксонов и мышечных волокон: в среднем около $340$ волокон на нейрон в межкостной мышце кисти, около $410$ в плечелучевой мышце, около $1930$ в медиальной икроножной. Точного числа для самых мелких единиц глава не приводит. & \cite{Feinstein1955mu} & измерено \\
2 & Синапс на мышце выбрасывает в несколько раз больше медиатора, чем нужно для срабатывания волокна & Обзор измерений на нервно-мышечных синапсах: запас надёжности (отношение выброшенного медиатора к пороговому) у взрослых млекопитающих обычно $3$--$5$; у человека запас держится больше на складках мембраны получателя, чем на величине выброса. & \cite{WoodSlater2001} & измерено \\
3 & Подёргивание~\eqref{eq:twitch} нарастает за время $T_c$ порядка $50\ms$ & Одиночные двигательные единицы кисти человека при произвольном усилии, сила усреднена по импульсам единицы: время нарастания $30$--$100\ms$, у $80\,\%$ единиц меньше $70\ms$. Функция $(t/T_c)e^{1-t/T_c}$ --- аппроксимация из модели пула единиц. & \cite{MilnerBrown1973rec, Fuglevand1993} & измерено \\
4 & Сумма подёргиваний~\eqref{eq:forcesum}: при $5$, $15$ и $40$ имп/с сила $0.2$--$1.1F_1$, $1.8$--$2.2F_1$ и около $5.4F_1$ (рис.~\ref{fig:tetanus}) & Расчёт по \texttt{models/\allowbreak muscle\_\allowbreak models.py}, $T_c=50\ms$: $0.21$--$1.10$, $1.76$--$2.19$ и $5.28$--$5.49\,F_1$ на последних $0.3\,\text{с}$. Линейное сложение --- упрощение: насыщение настоящей мышцы в модели нет. & \texttt{models/\allowbreak muscle\_\allowbreak models.py} & модель \\
5 & Единицы включаются по размеру; самые слабые в сотню раз слабее самых сильных & Вентральные корешки кошки: при нарастающем рефлекторном входе первыми разряжаются нейроны с мелкими импульсами в корешке, то есть мелкие. Межкостная мышца кисти человека: сила подёргивания единиц от $0.1$ до $10$~г и растёт почти линейно с силой, при которой единица включается. & \cite{Henneman1957, MilnerBrown1973rec} & измерено \\
6 & Шаг силы пропорционален силе, $\Delta F/F\approx q-1$~\eqref{eq:sizeprinciple}: плавающая точка & Выведено в главе из геометрической прогрессии сил. Согласуется с опытом: при больших усилиях на то же приращение силы включается гораздо меньше новых единиц. & вывод в главе; \cite{MilnerBrown1973rec} & теория \\
7 & Разброс силы растёт пропорционально самой силе & Произвольное изометрическое усилие мышцы большого пальца: стандартное отклонение силы линейно растёт со средней силой; при том же усилии, вызванном электрической стимуляцией, разброс постоянен. Модель пула: линейный рост даёт включение единиц по порядку силы. & \cite{Jones2002sdn} & измерено \\
8 & Плавность движений --- следствие шума, зависящего от сигнала & Модель: траектория, минимизирующая разброс конечной точки при таком шуме, воспроизводит измеренные траектории глаза и руки. Прямого опыта, отличающего эту причину от других, нет. & \cite{HarrisWolpert1998} & гипотеза \\
9 & Разность сил пары мышц задаёт момент, сумма --- жёсткость~\eqref{eq:antagonists} & Теория управления импедансом совместным напряжением. Рука человека: жёсткость каждого сустава, измеренная малыми толчками, примерно пропорциональна моменту в нём; разные соотношения совместного напряжения меняют форму и направление эллипса жёсткости кисти. & \cite{Hogan1984, GomiOsu1998stiff} & измерено \\
10 & Датчик растяжения возбуждает свой \nt{АЦЕТ}-нейрон и через тормозящего посредника выключает противника (рис.~\ref{fig:antagonists}) & Спинной мозг кошки: одиночный посредник, возбуждаемый волокнами датчиков растяжения, даёт моносинаптические тормозные потенциалы в \nt{АЦЕТ}-нейронах мышцы-противника, синаптическая задержка $0.28$--$0.42\ms$. Медиатор посредника (глицин) в этом опыте не определялся. & \cite{Jankowska1972Ia} & измерено \\
11 & Задержка петли: спинномозговой $25$--$30\ms$, через кору в полтора--три раза больше, через глаз $100$--$200\ms$ & Рука человека, толчок по суставу: короткий ответ мышц через $20$--$45\ms$, длинный $45$--$105\ms$, произвольный $120$--$180\ms$. Смещение видимого положения пальца во время движения: поправка через $160\ms$ в среднем. & \cite{Pruszynski2008rapid, SaundersKnill2003vis} & измерено \\
12 & $\dd x/\dd t=-Gx(t-d)$ устойчиво при $Gd<\pi/2$~\eqref{eq:delaystable}; на границе частота $1/(4d)$ & Теорема о корнях уравнения с запаздыванием. Проверка по \texttt{models/\allowbreak muscle\_\allowbreak models.py}, $d=30\ms$: к $3\,\text{с}$ амплитуда $3\cdot10^{-6}$ при $G=40$, $0.13$ при $G=50$, $38$ при $G=55$ в секунду (граница $52$). & \cite{Hayes1950} & теория \\
13 & Физиологическая дрожь $8$--$12$~Гц; петля растяжения --- один из её источников & Обзор измерений: мелкая дрожь в диапазоне около $10$~Гц есть у здоровых людей; её происхождение смешанное --- центральные колебания, свойства разрядов единиц, механический резонанс и резонанс рефлекторной петли, и в разных условиях преобладает разное. & \cite{McAuleyMarsden2000} & гипотеза \\
14 & Мозг предсказывает результат команды по внутренней модели тела & Человек держит груз щипком и двигает рукой: при инерционной, вязкой и упругой нагрузке сила сжатия меняется одновременно с нагрузкой, а не с запаздыванием петли, то есть нагрузка предсказана. Обзор сводит такие опыты; прямого наблюдения самой модели в нейронах нет. & \cite{FlanaganWing1997grip, WolpertGhahramani2000} & гипотеза \\
15 & Ритм ходьбы, дыхания, жевания порождают местные сети при постоянном входе & Обзор опытов: изолированная нервная система (спинной мозг, ганглии) выдаёт ритмический двигательный узор без ритмического входа и без обратной связи от мышц. & \cite{MarderBucher2001} & измерено \\
16 & Взаимное торможение с усталостью~\eqref{eq:halfcentre} даёт ритм с периодом $0.84\,\text{с}\approx2\tau_a$ (рис.~\ref{fig:halfcentre}) & Теорема: сеть со взаимным торможением и адаптацией без устойчивого состояния при постоянном входе обязательно колеблется. Расчёт по \texttt{models/\allowbreak muscle\_\allowbreak models.py} ($I=1$, $\beta=2$, $b=2$, $\tau=20\ms$, $\tau_a=0.4\,\text{с}$): период $0.84\,\text{с}$. Что настоящие генераторы устроены именно так, не показано. & \cite{Matsuoka1985osc} & модель \\
\bottomrule
\end{longtable}}

\section{Глава~\ref{sec:pain}. Боль}

Датчики боли, два провода, ворота в спинном мозге, нисходящая ручка громкости, сенсибилизация и захват внимания измерены прямо; формулы ворот и сенсибилизации --- упрощённые модели книги; правило, по которому машина выбирает громкость тревоги, --- гипотеза.

{\footnotesize
\setlength{\tabcolsep}{3pt}\renewcommand{\arraystretch}{1.15}
\begin{longtable}{>{\raggedright}p{0.5cm}>{\raggedright}p{4.0cm}>{\raggedright}p{5.6cm}>{\raggedright}p{2.3cm}>{\raggedright\arraybackslash}p{1.7cm}}
\toprule
№ & Утверждение & Опыт и что измерено & Источник & Статус \\
\midrule\endfirsthead
\toprule
№ & Утверждение & Опыт и что измерено & Источник & Статус \\
\midrule\endhead
1 & Высокий порог: датчики боли срабатывают при нагреве выше примерно $43^\circ$ & Запись одиночных волокон. Кожа обезьяны: медианный порог медленных (C) датчиков $41^\circ$, быстрых датчиков второго типа $46^\circ$, первого типа --- выше $53^\circ$. Кожа человека: C-датчики, отвечающие и на давление, $40.7^\circ$, не отвечающие на давление --- $48^\circ$. Число $43^\circ$ --- середина этого диапазона, а не единый порог. & \cite{Treede1995heat, Weidner1999cfib} & измерено \\
2 & После повреждения датчики не глохнут, а становятся чувствительнее & Лёгкий ожог кожи ($50^\circ$, $100\,\text{с}$): через $5$--$10$~мин порог боли у людей и порог C-датчиков у обезьяны снижены на $2$--$6^\circ$; у других кожных датчиков такого сдвига нет. Слабое привыкание прямо не сравнивалось; у быстрых датчиков второго типа ответ после сильного нагрева, наоборот, подавлен. & \cite{LaMotte1982hyper, Treede1995heat} & измерено \\
3 & Два провода: быстрый $5$--$30$~м/с, медленный $0.5$--$2$~м/с; время прихода $t=L/v$~\eqref{eq:paintime} & Скорость проведения: быстрые датчики боли обезьяны в среднем $15$ и $25$~м/с (два типа); C-датчики человека $0.8$--$1.0$~м/с. При $L=1$~м это десятки миллисекунд и около секунды. & \cite{Treede1995heat, Weidner1999cfib} & измерено \\
4 & Боль приходит дважды: сначала быстрая и точная, потом тупая и долгая & Время реакции на нагрев предплечья человека: от $1100$ до $400\ms$ с ростом температуры; сильный нагрев волосистой кожи даёт двойную боль, ладони --- нет. Первую боль могут нести только волокна быстрее $6$~м/с --- по латентности ей соответствуют быстрые датчики второго типа, которых в ладони нет. Раздел ролей («где» и «насколько плохо») --- толкование. & \cite{CampbellLaMotte1983first, Treede1995heat} & измерено \\
5 & Ворота: выходной нейрон спинного мозга получает боль и касание, тормозящий посредник касанием возбуждается (рис.~\ref{fig:gate}) & Схема ворот предложена как теория. Мыши, генетическая метка и удаление групп нейронов: возбуждающие нейроны с соматостатином нужны для механической боли; тормозящие нейроны с динорфином не дают входу от датчиков касания дойти до них --- без этих нейронов лёгкое касание вызывает боль. & \cite{MelzackWall1965, Duan2014} & измерено \\
6 & Касание приглушает боль & Человек, лазерные импульсы боли на руке: одновременное касание снижает оценку боли и способность различать энергию импульсов; эффект слабеет с расстоянием между точкой касания и точкой боли внутри одного кожного сегмента. & \cite{Mancini2014touch} & измерено \\
7 & Сильная боль сама выключает тормозящего посредника, и ворота распахиваются & Это часть исходной схемы ворот. В приведённых опытах прямо не показано: работа на мышах описывает, как ворота не пускают касание в канал боли, а не выключение посредника болью. & \cite{MelzackWall1965} & гипотеза \\
8 & Ворота~\eqref{eq:gate}: порог $0.18$ без касания, $0.86$ с касанием, $0.11$ при $g=2$; при $\nu=1$ касание снижает выход с $1$ до $0.25$, при $\nu=2$ не действует; касание само не проходит (рис.~\ref{fig:gatecurves}) & Расчёт по \texttt{models/\allowbreak pain\_\allowbreak gate.py} при весах из текста воспроизводит все эти числа. & \texttt{models/\allowbreak pain\_\allowbreak gate.py} & модель \\
9 & Нисходящие пути из ствола меняют усиление ворот в обе стороны & Крыса, продолговатый мозг, запись при нагреве хвоста: одни нейроны разряжаются перед отдёргиванием («вкл»), другие замолкают («выкл»); слабая стимуляция этой области подавляет рефлекс. Обзор: те же цепи и облегчают, и тормозят передачу боли, опиоиды действуют через них. & \cite{Fields1983rvm, Fields2004} & измерено \\
10 & Ожидание облегчения приглушает боль через опиоидный канал & Люди с острой болью: у тех, кому ожидание облегчения снизило боль, блокада опиоидных рецепторов вернула боль к уровню остальных; у остальных блокада боль не меняла. & \cite{Levine1978} & измерено \\
11 & Громкость тревоги мозг выбирает по выгоде прерывания & Опыта, в котором измерено правило выбора громкости, нет. & --- & гипотеза \\
12 & Сенсибилизация: после повреждения выход ворот растёт, порог падает, отчасти за счёт спинного мозга & Крыса: после повреждения кожи порог сгибательного рефлекса падает, а ответ растёт; электрофизиологический анализ показывает, что часть роста возникает в самом спинном мозге. Время спада «от часов до дней» в этой работе не измерялось. & \cite{Woolf1983} & измерено \\
13 & Сенсибилизация~\eqref{eq:sensitization} --- положительная обратная связь; при сильной петле тревога может защёлкнуться после заживления & Следует из модели~\eqref{eq:sensitization} (задача в конце главы). Опыта, показывающего, что стойкая боль после заживления --- защёлкнувшаяся петля именно такого вида, нет. & вывод в главе & гипотеза \\
14 & Боль --- отрицательная награда, и обучение на ней следует ошибке временных разностей & Томограф, люди, цепочки предвестников боли: активность вентрального стриатума и передней островковой коры совпадает с сигналами обучения по временным разностям. Обезьяна: часть \nt{ДОФА}-нейронов возбуждается от неприятного обдува и его предвестника, другая часть тормозится. & \cite{Seymour2004, Matsumoto2009} & измерено \\
15 & Боль прерывает текущую работу & Люди, электрический болевой стимул и звуковые пробы: ответ на пробу через $100\ms$ после начала боли заметно замедлен, через $1.5\,\text{с}$ разницы с контролем уже нет. Обзор сводит такие опыты в модель захвата внимания. & \cite{Crombez1996disrupt, EcclestonCrombez1999} & измерено \\
16 & Отдёргивание замыкается в спинном мозге & Крыса: нейроны глубоких слоёв заднего рога повторяют пространственный узор рефлекса отдёргивания отдельных мышц; их не удаётся возбудить антидромно из шейного отдела, то есть это спинальные посредники. & \cite{Schouenborg1995wr} & измерено \\
\bottomrule
\end{longtable}}

\section{Глава~\ref{sec:motorlearning}. Обучение движениям}

Правило наименьших квадратов и выгода отдельного провода ошибки --- теоремы; устройство схемы с проводом ошибки, ослабление при совпадении, подстройка следа под задержку, последействие и спонтанный возврат навыка измерены; то, что вся схема работает как обучение с учителем и строит внутреннюю модель, --- ведущая гипотеза; числа двухпроцессной модели --- расчёт по модели книги.

{\footnotesize
\setlength{\tabcolsep}{3pt}\renewcommand{\arraystretch}{1.15}
\begin{longtable}{>{\raggedright}p{0.5cm}>{\raggedright}p{4.0cm}>{\raggedright}p{5.6cm}>{\raggedright}p{2.3cm}>{\raggedright\arraybackslash}p{1.7cm}}
\toprule
№ & Утверждение & Опыт и что измерено & Источник & Статус \\
\midrule\endfirsthead
\toprule
№ & Утверждение & Опыт и что измерено & Источник & Статус \\
\midrule\endhead
1 & Правило~\eqref{eq:lms} в среднем идёт по градиенту квадрата ошибки & Результат теории адаптивных фильтров: поправка, пропорциональная входу и ошибке выхода, есть стохастический спуск по градиенту среднего квадрата ошибки. & \cite{WidrowHoff1960} & теория \\
2 & С проводом ошибки на каждый выход скорость не зависит от числа выходов; с одним общим числом --- падает с числом шумящих нейронов (утверждение~\ref{prop:teachers}) & Кривые обучения линейной сети: обучение по случайным возмущениям нейронов медленнее точного градиента во столько раз, сколько возмущаемых нейронов. & \cite{Werfel2005} & теория \\
3 & Схема с проводом ошибки работает как обучение с учителем (утверждение~\ref{prop:errorwire}) & Теории, выведенные из строения схемы. Опыты строк 7--10 согласуются с ними; что вся схема работает именно так, не доказано. & \cite{Marr1969, Albus1971} & гипотеза \\
4 & Расширитель: около $7\cdot10^{10}$ мелких \nt{GLUT}-нейронов, больше, чем во всём остальном мозге & Подсчёт ядер во взвеси растворённой ткани: в мозжечке человека $69\cdot10^9$ нейронов из $86\cdot10^9$ во всём мозге. Стереология: $101\cdot10^9$ мелких нейронов в мозжечке пожилых мужчин. & \cite{Azevedo2009, Andersen1992cereb} & измерено \\
5 & Подъём размерности входа делает нужную зависимость линейной & Теорема о числе разбиений: взвешенная сумма $n$ входов с порогом реализует произвольную разметку примерно $2n$ векторов общего положения. Что мозжечок пользуется именно этим, --- часть гипотезы строки~3. & \cite{Cover1965} & теория \\
6 & Около $3\cdot10^7$ выходных \nt{GABA}-нейронов, у каждого порядка $10^5$ входов & Стереология мозжечка человека: $30.5\cdot10^6$ выходных нейронов. Электронная микроскопия, крыса: около $1.75\cdot10^5$ входов от расширителя на один выходной нейрон. Тормозной медиатор выходных нейронов --- в обзоре. & \cite{Andersen1992cereb, NapperHarvey1988, Ito2001} & измерено \\
7 & Ровно один провод ошибки на выходной нейрон, около раза в секунду, каждый раз мощная вспышка & Обзор записей: каждый выходной нейрон получает один лазящий \nt{GLUT}-вход, который вызывает сложный импульс с частотой порядка $1$~Гц. & \cite{Ito2001} & измерено \\
8 & Вероятность вспышки зависит от направления ошибки движения & Обезьяна, глазодвигательный отдел мозжечка, адаптация саккад: направление зрительной ошибки задаёт вероятность сложного импульса, величина --- его время; вспышка меняет простые импульсы в следующей попытке и сдвигает движение вдоль предпочтительной ошибки клетки. & \cite{Herzfeld2018} & измерено \\
9 & Совпадение вспышки ошибки с активным входом ослабляет вход ($-\eta e_i x_j$) & Обзор опытов: совместная стимуляция входа от расширителя и провода ошибки вызывает долговременное ослабление этого входа; стимуляция одного входа --- нет. & \cite{Ito2001} & измерено \\
10 & Длина следа подстроена под задержку ошибки: при задержке $100\ms$ сильнее всего ослабляется вход, бывший активным за $100\ms$ & Срезы и живые мыши: в отделе, отвечающем за глазодвигательное обучение, ослабление узко настроено на интервал около $120\ms$ между входом и ошибкой --- столько идёт сигнал ошибки в эту задачу; в другом отделе разные клетки настроены на разные интервалы. & \cite{Suvrathan2016} & измерено \\
11 & Схема --- адаптивный фильтр~\eqref{eq:adaptivefilter}, выучивающий внутреннюю модель тела & Модель, выведенная из строения схемы. Прямого опыта, где выученный фильтр измерен как передаточная функция тела, нет. & \cite{Fujita1982} & гипотеза \\
12 & Предсказание вычитает следствия собственных движений, поэтому себя не пощекотать & Люди: собственное касание кажется менее щекотным, чем то же чужое; в томографе соматосенсорная кора отвечает на собственное касание слабее, а мозжечок меньше активен, когда движение касается кожи. Объяснение через вычитание предсказания --- гипотеза. & \cite{Blakemore1998} & гипотеза \\
13 & При внешней силе промах уменьшается, а после её снятия рука промахивается в другую сторону & Люди тянутся рукоятью робота в силовом поле: траектории возвращаются к прямым; после снятия поля последействие --- почти зеркальное отражение первых искажений; оно переносится и в не тренированные части пространства. & \cite{ShadmehrMussaIvaldi1994} & измерено \\
14 & Учатся два процесса~\eqref{eq:tworate}; после обратного возмущения навык возвращается сам & Люди, силовое поле, затем короткое обратное поле и попытки с зажатой ошибкой: поправка сама возвращается к старой; модель из двух процессов это воспроизводит, модель из одного --- нет. & \cite{Smith2006} & измерено \\
15 & Числа рис.~\ref{fig:tworate}: $0.84$ (медленный $0.63$) за $200$ движений; $6$ обратных; быстрый $-0.53$, медленный $0.46$; возврат до $0.37$ за $40$ движений & Расчёт по \texttt{models/\allowbreak motor\_\allowbreak learning.py}, $A_f=0.92$, $B_f=0.1$, $A_s=0.996$, $B_s=0.02$: $0.840$ ($0.634$ и $0.205$), $6$ движений, $-0.531$ и $0.457$, пик $0.371$ на $40$-м движении. & \texttt{models/\allowbreak motor\_\allowbreak learning.py} & модель \\
16 & Два процесса нужны, потому что тело меняется с разными скоростями & Теория: оптимальная оценка при помехах с несколькими временами жизни даёт несколько фильтров с разными постоянными времени и объясняет спонтанный возврат и ускоренное повторное обучение. & \cite{Kording2007} & гипотеза \\
\bottomrule
\end{longtable}}

\section{Глава~\ref{sec:chemoutput}. Химический выход}

Два провода к сердцу и их разная скорость, отрицательная обратная связь в гормональных каскадах, код частотой порций, суточный генератор и его подстройка светом измерены прямо; граница $K<8$ --- теорема теории регулирования, выведенная в главе; пульсации, порождаемые самой обратной связью каскада, --- гипотеза с сильным опытом в её пользу.

{\footnotesize
\setlength{\tabcolsep}{3pt}\renewcommand{\arraystretch}{1.15}
\begin{longtable}{>{\raggedright}p{0.5cm}>{\raggedright}p{4.0cm}>{\raggedright}p{5.6cm}>{\raggedright}p{2.3cm}>{\raggedright\arraybackslash}p{1.7cm}}
\toprule
№ & Утверждение & Опыт и что измерено & Источник & Статус \\
\midrule\endfirsthead
\toprule
№ & Утверждение & Опыт и что измерено & Источник & Статус \\
\midrule\endhead
1 & Ритмоводитель сердца сам бьётся около $100$ раз в минуту; в покое тормоз нажат, и частота обычно порядка $60$--$70$ & Люди, полное выключение обоих проводов к сердцу: собственная частота выше частоты покоя и убывает с возрастом, у взрослых порядка $100$ ударов в минуту. Значит, в покое преобладает тормоз. Частота покоя $60$--$70$ --- типичное значение, отдельно не цитировано. & \cite{JoseCollison1970ihr} & измерено \\
2 & Газ \nt{NORA} и тормоз \nt{ACET} --- фильтры нижних частот, тормоз быстрее~\eqref{eq:gasbrake} & Собака, частотно-модулированная стимуляция блуждающего или симпатического нерва сердца и передаточная функция к частоте предсердий: оба пути --- фильтры нижних частот, у симпатического частота среза много ниже и добавлена чистая задержка около $1.7\,\text{с}$. Характеристики зависят от среднего тонуса, так что линейная формула~\eqref{eq:gasbrake} --- приближение. & \cite{Berger1989} & измерено \\
3 & Тормоз быстр, потому что калиевый канал \nt{ACET} открывает без второго посредника, а \nt{NORA} --- через цепочку реакций & Мембранные лоскуты клеток предсердий: калиевый канал, открываемый ацетилхолином, открывают $\beta\gamma$-субъединицы G-белка, связанного с рецептором, без второго посредника внутри клетки. Путь \nt{NORA} через цАМФ здесь не цитирован. & \cite{Logothetis1987gbg} & измерено \\
4 & Кровь обходит тело за время порядка минуты; \nt{ADRE} из кровотока доходит до всех органов с рецептором & Общая оценка порядка величины; в главе так и сформулирована, источник не цитирован. & --- & оценка \\
5 & Гормональный каскад охвачен отрицательной обратной связью: конечный гормон тормозит первые ступени & Обзор опытов: гормоны коры надпочечников подавляют выброс АКТГ гипофизом и рилизинг-гормона гипоталамусом в трёх временных диапазонах --- быстром, промежуточном и медленном. & \cite{KellerWood1984fb} & измерено \\
6 & Три одинаковые ступени устойчивы при $K<8$~\eqref{eq:cascadestable}, на границе период $2\pi\tau/\sqrt3\approx3.6\tau$ & Выведено в главе: на частоте $\sqrt3/\tau$ три ступени дают сдвиг фазы $180^\circ$ и ослабление в $8$ раз. & вывод в главе & теория \\
7 & Ошибка уровня $1/(1+K)$, поэтому точнее примерно $10\,\%$ такой регулятор не держит (утверждение~\ref{prop:cascade}) & Следствие строки~6 для пропорциональной обратной связи. Настоящая ось имеет обратную связь на нескольких ступенях и в нескольких временных диапазонах (строка~5), так что граница относится к модели~\eqref{eq:cascade}, а не измерена. & вывод в главе & теория \\
8 & При $K=4$ и $7$ уровень устанавливается, при $K=12$ --- ровные пульсации с периодом $3.7$--$3.9\,\tau$ (рис.~\ref{fig:cascade}) & Расчёт по \texttt{models/\allowbreak chem\_\allowbreak models.py}: при $K=12$ последовательные периоды $3.9$, $3.7$, $3.8$, $3.7\,\tau$; при $K=4$ размах в конце счёта $0.003$. & \texttt{models/\allowbreak chem\_\allowbreak models.py} & модель \\
9 & Гормон стресса выходит пульсами около раза в час; пульсы порождает сама обратная связь, без отдельного генератора & Крысы, постоянное вливание рилизинг-гормона: АКТГ и кортикостерон колеблются с физиологической, почти часовой частотой --- ритмичный вход от гипоталамуса для пульсов не нужен. Модель гипофиз--надпочечник с запаздывающей обратной связью даёт такие колебания. & \cite{Walker2012glc, Walker2010} & гипотеза \\
10 & Получатель читает частоту порций, а не уровень & Обезьяны без собственного выброса рилизинг-гормона гонадотропинов: постоянное вливание не восстанавливает выброс гормонов гипофиза, порции раз в час восстанавливают; переход на постоянное вливание снова глушит ответ. Утомление рецептора как фильтр верхних частот --- толкование. & \cite{Belchetz1978} & измерено \\
11 & Разные частоты порций по-разному действуют на разные клетки одной железы & Те же обезьяны: учащение порций до $2$--$5$ в час снижает оба гормона гипофиза; урежение до одной за $3$~часа снижает один (ЛГ) и повышает другой (ФСГ), так что их отношение задаётся частотой. & \cite{Wildt1981gnrh} & измерено \\
12 & Суточный ритм порождает внутриклеточная отрицательная обратная связь с задержкой в несколько часов & Обзор: белки часовых генов с задержкой подавляют транскрипцию собственных генов; петля из транскрипции и трансляции даёт период около суток. & \cite{TakahashiJS2017} & измерено \\
13 & Главный генератор --- группа из десятков тысяч нейронов, синхронизирующая остальное тело & Обзор: надперекрёстное ядро --- главный суточный генератор; отдельные его нейроны в культуре --- самостоятельные генераторы, сеть их синхронизирует; у мыши около $10^4$ нейронов с каждой стороны. & \cite{Welsh2010scn} & измерено \\
14 & Собственный период у человека $24.18$~часа & Люди в тусклом свете при расписании, оторванном от суток: период ритмов мелатонина, температуры и кортизола в среднем $24.18$~часа у молодых и пожилых, с узким разбросом. & \cite{Czeisler1999} & измерено \\
15 & Свет подводит генератор как фазовая автоподстройка~\eqref{eq:pll}; нужен сдвиг около $11$~мин в сутки & Люди, одиночный яркий свет $6.7$~часа в разное время суток: кривая фазового ответа первого типа с размахом $5$~часов --- задержка до минимума температуры тела, опережение после. Уравнение~\eqref{eq:pll} --- стандартная модель; $11$~мин $=0.18$~часа --- арифметика. & \cite{Khalsa2003prc} & измерено \\
16 & Без света подстройки нет, и ритм уходит на собственный период & Полностью слепые люди без восприятия света, живущие в обычном обществе: примерно у половины ритмы мелатонина и кортизола идут со своим периодом, не совпадающим с сутками. & \cite{Sack1992blind} & измерено \\
\bottomrule
\end{longtable}}

\section{Глава~\ref{sec:debugging}. Отладка машины}

Что слышит электрод, маломерность активности, работа интерфейсов на жёстких решётках, совместное обучение мозга и декодера и зрительные точки от стимуляции измерены в рецензируемых опытах; оценки потоков данных --- арифметика главы; о вживлённом людям устройстве Neuralink, насколько удалось проверить, публиковала данные в основном сама компания.

{\footnotesize
\setlength{\tabcolsep}{3pt}\renewcommand{\arraystretch}{1.15}
\begin{longtable}{>{\raggedright}p{0.5cm}>{\raggedright}p{4.0cm}>{\raggedright}p{5.6cm}>{\raggedright}p{2.3cm}>{\raggedright\arraybackslash}p{1.7cm}}
\toprule
№ & Утверждение & Опыт и что измерено & Источник & Статус \\
\midrule\endfirsthead
\toprule
№ & Утверждение & Опыт и что измерено & Источник & Статус \\
\midrule\endhead
1 & Электрод слышит импульсы нейронов в радиусе порядка сотни микрометров, обычно от одного до нескольких; их различают по форме & Крыса, поле CA1, одновременная внутриклеточная и внеклеточная запись: форма внеклеточного импульса повторяет ширину и амплитуду внутриклеточного. Оценка: тетрод мог бы различать $60$--$100$ нейронов, а на деле выделяют около шести. & \cite{Henze2000extra} & измерено \\
2 & В мозге $8.6\cdot10^{10}$ нейронов; щуп слышит примерно одну стомиллионную & Подсчёт ядер во взвеси растворённой ткани: $86\cdot10^9$ нейронов в мозге человека. Доля --- арифметика главы. & \cite{Azevedo2009} & измерено \\
3 & Активность сотен нейронов двигательной коры укладывается в десяток-другой общих переменных & Обзор записей в двигательной коре обезьян: активность популяции описывается немногими скрытыми переменными, общими для многих нейронов. & \cite{Gallego2017} & измерено \\
4 & Шум соседей согласован, и сотня нейронов несёт почти столько же, сколько тысяча (утверждение~\ref{prop:pooling}) & Зрительная кора обезьяны: коэффициент корреляции шума соседних нейронов около $0.1$, и выигрыш от усреднения насыщается при сотне нейронов. Теория ограничивающих информацию корреляций. & \cite{Zohary1994, MorenoBote2014} & измерено \\
5 & Сырой поток $2\cdot10^8$~бит/с против $8\cdot10^4$~бит/с в импульсах~\eqref{eq:probedata} & Арифметика главы по ёмкости потока импульсов~\eqref{eq:spikeentropy}. Параметры оцифровки реальны: кристалл Neuralink --- $19.3$~кГц, $10$~бит на канал. & вывод в главе; \cite{Rieke1997, Musk2019} & теория \\
6 & Импульсы выделяют на кристалле у электродов и передают наружу только события; корпус закрыт кожей, связь и питание без проводов & Система, описанная в статье компании, передаёт полный сырой поток по кабелю USB-C, импульсы выделяет программа в реальном времени вне кристалла; сжатие на борту, питание и передача без проводов названы там планом для клинических устройств. Для вживлённого людям устройства --- только сообщения компании. & \cite{Musk2019}; отчёты Neuralink, рецензируемой статьи нет & измерено (отчёт компании) \\
7 & До $3072$ электродов на $96$ нитях по $32$; нити вставляет робот, обходя сосуды & Статья компании: $3072$ электрода на $96$ нитях по $32$; робот вставляет $6$ нитей ($192$ электрода) в минуту; до $70\,\%$ электродов у крыс с долговременно вживлённой решёткой слышат импульсы. & \cite{Musk2019} & измерено (статья компании) \\
8 & У людей около тысячи каналов на $64$ нитях; часть нитей сместилась; курсор порядка $10$~бит/с & Только сообщения компании. Поиск в PubMed не нашёл рецензируемой статьи с результатами на людях; это согласуется с оговоркой в тексте главы. & отчёты Neuralink; рецензируемой статьи нет & измерено (отчёт компании) \\
9 & Интерфейс на решётке из сотни электродов управляет курсором & Человек с параличом, $96$ электродов в первичной двигательной коре: намерение двигать рукой меняет импульсы через три года после травмы; декодер даёт «нейронный курсор» (почта, телевизор), управление протезом кисти и роботом-рукой. & \cite{Hochberg2006} & измерено \\
10 & Письмо «мысленной рукой»: около $90$ знаков в минуту, меньше $8$~бит/с & Человек с параличом кисти, попытки писать от руки, рекуррентный декодер: $90$ знаков в минуту при точности $94.1\,\%$ в реальном времени, выше $99\,\%$ после автокоррекции. Оценка в битах --- арифметика главы. & \cite{Willett2021} & измерено \\
11 & Попытки говорить переводятся в текст со скоростью около $60$ слов в минуту & Человек, потерявший внятную речь, решётки в коре: $62$ слова в минуту; $9.1\,\%$ ошибок в словах при словаре из $50$ слов, $23.8\,\%$ при словаре из $125\,000$. & \cite{Willett2023} & измерено \\
12 & Интерфейс вынимает малую долю ёмкости: у одного нейрона десятки бит в секунду, у сотни --- тысячи (утверждение~\ref{prop:probe}) & Верхняя граница~\eqref{eq:spikeentropy} --- теория; сравнение со строками~8, 10 --- вывод главы. Что узкое место --- маломерность и незнание кода, а не провод, прямым опытом не проверено. & \cite{Rieke1997} & теория \\
13 & Мозг и декодер учатся друг у друга & Обезьяны управляют курсором; декодер подстраивается в замкнутой петле, а нейроны одновременно переучиваются: точность растёт, навык сохраняется и меньше страдает от движений собственной руки. Совет развести скорости обучения --- инженерное правило, отдельного опыта в главе нет. & \cite{Orsborn2014coad} & измерено \\
14 & Ток через электрод возбуждает окружающие нейроны и провода без разбора типа & Кора мыши и кошки, двухфотонная запись кальция: даже слабый ток возбуждает редкие нейроны, разбросанные на расстояние до миллиметров, через прямое возбуждение аксонов в объёме диаметром в десятки микрометров. То есть возбуждение неизбирательно, но не сплошное. & \cite{Histed2009stim} & измерено \\
15 & Стимуляция зрительной коры зажигает точки; до $16$ точек сразу позволяют узнать некоторые буквы и границы предметов & Слепой человек, $96$ электродов в зрительной коре на $6$ месяцев: порог одного электрода $66.8\pm36.5$~мкА; одновременная стимуляция групп (до $16$ электродов --- предел стимулятора) снижает порог и даёт различимые узоры, некоторые буквы и границы предметов узнаются. & \cite{Fernandez2021} & измерено \\
16 & Второе направление Neuralink --- устройство, подающее сигналы в зрительную кору & Только сообщения компании о разработке; рецензируемых данных о его работе не найдено. & отчёты Neuralink & гипотеза \\
17 & Концентрации широковещательных медиаторов можно мерить химическим датчиком в ткани & Срезы мозга крысы, быстрая циклическая вольтамперометрия на углеродном волокне: измерены выброс и обратный захват серотонина; вещество выходит за пределы синапса. & \cite{Bunin1998} & измерено \\
\bottomrule
\end{longtable}}

\section{Глава~\ref{sec:eetypes}. Нейроны как приборы}

Режимы работы отдельных клеток измерены прямо, током через электрод и записью напряжения; математика интегратора и деления на классы --- классическая теория, а то, что мозг пользуется перестройкой режимов для вычислений, остаётся гипотезой.

{\footnotesize
\setlength{\tabcolsep}{3pt}\renewcommand{\arraystretch}{1.15}
\begin{longtable}{>{\raggedright}p{0.5cm}>{\raggedright}p{4.0cm}>{\raggedright}p{5.6cm}>{\raggedright}p{2.3cm}>{\raggedright\arraybackslash}p{1.7cm}}
\toprule
№ & Утверждение & Опыт и что измерено & Источник & Статус \\
\midrule\endfirsthead
\toprule
№ & Утверждение & Опыт и что измерено & Источник & Статус \\
\midrule\endhead
1 & Резонатор: медленный ток, противодействующий изменению напряжения, делает нейрон полосовым фильтром с максимумом на единицах--десятках герц (раздел~\ref{sec:resonator}) & В срезах в \nt{GLUT}-нейроны вводили синусоидальный ток с плавно растущей частотой и мерили импеданс $|Z(f)|$: максимум на $2$--$5$~Гц при $33^\circ$C (около $7$~Гц при $38^\circ$C); резонанс создают медленные калиевый и катионный токи и усиливает постоянный натриевый ток. Обзор даёт тот же приём для многих классов клеток & \cite{HuVervaekeStorm2002,HutcheonYarom2000} & измерено \\
2 & Медленный ток ведёт себя как индуктивность, вместе с ёмкостью мембраны --- колебательный контур & Подпороговый отклик аксона на ток измерен и описан линеаризованными уравнениями проводимостей; эквивалентная схема содержит «феноменологическую» индуктивность, величина которой зависит от напряжения & \cite{Mauro1970} & теория \\
3 & Постоянная времени группы с обратной связью $\tau_{\text{эфф}}=\tau/(1-w)$~\eqref{eq:perfectint}; при $\tau=0.1$~с и $w=0.995$ --- $20$~с & Выводится в главе из линейного уравнения группы; как механизм удержания положения глаза предложено в теоретической работе & вывод в главе; \cite{Seung1996} & теория \\
4 & Интегратор положения глаза держит вход сетью, а не свойством отдельной клетки & Внутриклеточная запись у рыбы в узле, держащем положение глаза: после быстрого поворота частота нейрона ступенькой меняется и держится; короткий ток в одну клетку вызывает лишь кратковременный сдвиг, а ступеньки напряжения сохраняются и ниже порога --- их поддерживают синаптические входы & \cite{Aksay2001} & измерено \\
5 & Интегратору без утечки нужна постоянная подстройка обучением; при расстройке он либо забывает, либо уходит в насыщение & Рыбу обучали зрительной обратной связью, пропорциональной $\pm$ положению глаза: удержание становилось неустойчивым или с утечкой, постоянная времени нейронов падала больше чем на порядок, до $<1$~с; обычная зрительная обратная связь постепенно возвращала устойчивость & \cite{Major2004} & измерено \\
6 & Бистабильный нейрон: короткий вход включает длительное плато, торможение выключает; свойство включается серотонином (раздел~\ref{sec:bistable}) & Внутриклеточная запись \nt{ACET}-нейронов, управляющих мышцами, у кошки: короткое синаптическое возбуждение переводит нейрон в длительную работу, короткое торможение сбрасывает; у спинальной кошки состояние появляется после введения предшественника серотонина & \cite{Hounsgaard1988} & измерено \\
7 & Два класса: интеграторы (частота растёт от нуля) и резонаторы (при пороге сразу конечная частота) & Классы выведены из теории бифуркаций. Опыт: в срезах коры \nt{GLUT}-нейроны начинают работать со сколь угодно малой частотой (тип~1), быстрые \nt{GABA}-нейроны --- скачком с конечной частоты (тип~2) & \cite{Izhikevich2007,Tateno2004} & измерено \\
8 & Один нейрон --- интегратор или детектор совпадений: торможение укорачивает окно сложения & В срезах гиппокампа окно, в котором возбуждающие входы складываются до порога, короче $2$~мс; его укорачивает прямое торможение, приходящее сразу за возбуждением; в дендритах, где торможение слабее, окно шире & \cite{Pouille2001} & измерено \\
9 & Линия задержки из аксонов (таблица~\ref{tab:eetypes}) & У сипухи измерены задержки в аксонах, идущих к детекторам совпадений: разная длина аксона даёт разную задержку, и детекторы настроены на разность времён прихода звука & \cite{Carr1990} & измерено \\
10 & Ритмоводитель в бодрствовании и молчаливая клетка во сне & \nt{SERO}-нейроны работают почти регулярно днём, замедляются в медленном сне и почти молчат в быстром (подробнее --- глава~\ref{sec:sleep}) & \cite{JacobsAzmitia1992,McGintyHarper1976} & измерено \\
11 & Прибор задают ионные каналы и широковещательные каналы, которые их подстраивают (утверждение~\ref{prop:eetypes}) & Показано на малых сетях: вещества-регуляторы меняют собственные свойства нейронов и силу синапсов, так что одна и та же схема выдаёт разные ритмы & \cite{Marder2012} & измерено \\
12 & Мозг использует перестройку режимов для вычислений (утверждение~\ref{prop:eetypes}) & Прямого опыта, где перестройка режима клетки была бы нужна для решения задачи, нет; есть лишь опыты, где регуляторы меняют выход схемы & \cite{Marder2012} & гипотеза \\
\bottomrule
\end{longtable}}

\section{Глава~\ref{sec:dendrites}. Число дендритов}

Строение классов и числа входов измерены анатомически и электрическими записями; оценка~\eqref{eq:dendriteload} --- простая физика конденсатора, а вычислительное объяснение числа входов опирается на теорию и модели.

{\footnotesize
\setlength{\tabcolsep}{3pt}\renewcommand{\arraystretch}{1.15}
\begin{longtable}{>{\raggedright}p{0.5cm}>{\raggedright}p{4.0cm}>{\raggedright}p{5.6cm}>{\raggedright}p{2.3cm}>{\raggedright\arraybackslash}p{1.7cm}}
\toprule
№ & Утверждение & Опыт и что измерено & Источник & Статус \\
\midrule\endfirsthead
\toprule
№ & Утверждение & Опыт и что измерено & Источник & Статус \\
\midrule\endhead
1 & Удельная ёмкость мембраны $c_m\approx1$~мкФ/см$^2$ & С тела нейрона снимали мембрану на пипетку, подавали ступеньку напряжения и по спаду ёмкостного тока находили полную ёмкость, затем делили на площадь: $0.9$~мкФ/см$^2$ у трёх классов нейронов & \cite{Gentet2000} & измерено \\
2 & Время заряда $t\approx c_mA\Delta V/I$~\eqref{eq:dendriteload}: большому нейрону нужны десятки одновременных входов, маленькому хватает одного крупного синапса & Оценка по закону заряда конденсатора; числа ($20$ и $300$~пФ, $0.1$ и несколько нА) --- порядки величин & вывод в главе & теория \\
3 & Один огромный синапс даёт ток в несколько наноампер и сразу доводит маленький нейрон до порога & Одновременная запись окончания и получателя в срезе: ток одного синапса $2$--$13$~нА, каждый импульс входа вызывает надпороговый ответ, задержка около $0.4$~мс при $36^\circ$C; получатель выдаёт \nt{ГЛИЦ} & \cite{Borst1995,BorstSoria2012} & измерено \\
4 & Ноль дендритов: у чувствительных нейронов осязания и боли импульс идёт от датчика в спинной мозг мимо тела клетки & Строение (аксон, делящийся надвое у тела) установлено анатомически. То, что проведение не требует возбудимости тела, показано на проверенной модели такого нейрона: без возбудимого тела импульс проходит развилку так же надёжно & \cite{AmirDevor2003} & теория \\
5 & Один дендрит: обонятельный нейрон несёт один тип рецептора и передаёт одну входную линию & Обзор опытов с генами рецепторов: каждый обонятельный нейрон выражает один ген рецептора из большого семейства & \cite{Mombaerts2004} & измерено \\
6 & Два дендрита: у детекторов направления на звук входы от левого и правого уха садятся на разные дендриты & Анатомия и записи у млекопитающих, собранные в обзоре: входы от двух ушей разнесены по двум противоположным дендритам & \cite{Grothe2010} & измерено \\
7 & Разнесение входов по двум дендритам делает «И» строже & Показано на модели: насыщение~\eqref{eq:shunt} на одном дендрите не даёт входам с одной стороны дойти до порога, и детекция совпадений улучшается. Прямого опыта, где бы меняли раскладку входов, нет & \cite{AgmonSnir1998} & теория \\
8 & Около $7\cdot10^{10}$ маленьких \nt{ГЛУТ}-нейронов цепи обучения движениям, по $\sim4$ входа у каждого & Подсчёт клеток изотропным фракционированием: в этой части мозга человека около $69$ млрд нейронов. Запись одной такой клетки у живой крысы: на касание приходят пачки дискретных токов от немногих входов, и клетка срабатывает, когда они складываются & \cite{Azevedo2009,Chadderton2004} & измерено \\
9 & Пороговый элемент с $n$ входами разделяет не больше $P_{\max}\approx2n$ случайных образов~\eqref{eq:cover} & Классический результат комбинаторной геометрии. Оговорка: при одних лишь положительных весах и с запасом на шум ёмкость меньше; для выходного нейрона цепи обучения движениям оценено $\sim4\cdot10^4$ соответствий, а не $2\cdot10^5$ & \cite{Cover1965,Brunel2004} & теория \\
10 & Смесители с малым числом входов нужны, чтобы расширить вход; «четыре» близко к оптимуму & Теория: размерность представления, которую даёт слой смесителей, максимальна примерно при том числе входов, которое наблюдается анатомически; исходная гипотеза расширения --- в классических работах & \cite{Marr1969,Albus1971,LitwinKumar2017} & гипотеза \\
11 & Большие деревья: $10^4$--$10^5$ входов & Подсчёт синапсов по электронным снимкам: $\approx1.7\cdot10^5$ синапсов на одном выходном \nt{ГАМК}-нейроне цепи обучения движениям у крысы; около $30\,000$ возбуждающих и $1\,700$ тормозящих входов на \nt{ГЛУТ}-нейроне гиппокампа & \cite{NapperHarvey1988,Megias2001} & измерено \\
12 & Ветви большого дерева --- отдельные подсхемы со своими порогами; нейрон --- маленькая многослойная сеть & Точечная стимуляция двух мест на тонких дендритах: входы на одной ветви складываются по сигмоиде, на разных ветвях --- линейно. В дендритах нейронов коры человека найдены градуальные кальциевые импульсы, позволяющие одной клетке решать линейно неразделимую задачу. Модель: чтобы воспроизвести один нейрон, нужна глубокая сеть & \cite{Polsky2004,Gidon2020,Beniaguev2021} & измерено \\
13 & Дендриты-выходы: каждая ветвь \nt{ГАМК}-нейрона сетчатки --- отдельный детектор направления & Двухфотонная запись кальция в отдельных ветвях: сигнал в ветви избирателен к движению от тела клетки к кончику, а напряжение тела --- нет & \cite{Euler2002} & измерено \\
14 & Число входов следует за задачей (утверждение~\ref{prop:dendrites}) & Строение классов измерено (строки 3--13); что число входов выбрано ради скорости или классификации, показано лишь теорией и моделями для отдельных классов & \cite{LitwinKumar2017,AgmonSnir1998} & гипотеза \\
\bottomrule
\end{longtable}}

\section{Глава~\ref{sec:sleep}. Сон}

Режимы сна, повторы и уменьшение синапсов измерены записями нейронов, микродиализом и электронной микроскопией; расписание сна и медленные качели описаны моделями книги, а назначение сна --- в основном гипотезы.

{\footnotesize
\setlength{\tabcolsep}{3pt}\renewcommand{\arraystretch}{1.15}
\begin{longtable}{>{\raggedright}p{0.5cm}>{\raggedright}p{4.0cm}>{\raggedright}p{5.6cm}>{\raggedright}p{2.3cm}>{\raggedright\arraybackslash}p{1.7cm}}
\toprule
№ & Утверждение & Опыт и что измерено & Источник & Статус \\
\midrule\endfirsthead
\toprule
№ & Утверждение & Опыт и что измерено & Источник & Статус \\
\midrule\endhead
1 & Во сне нейроны коры не молчат; меняется режим работы & Долгие записи нейронов коры у крыс: в медленном сне частота остаётся ненулевой, периоды работы чередуются с периодами тишины; после долгого бодрствования частота выше во всех состояниях, после сна --- ниже & \cite{Vyazovskiy2009} & измерено \\
2 & \nt{ACET} высок днём и в быстром сне, низок в медленном (таблица~\ref{tab:sleepmodes}) & Микродиализ в коре у свободно двигающихся кошек: в бодрствовании выброс ацетилхолина на $100$--$175\%$ выше, чем в медленном сне, а в быстром сне --- примерно на уровне спокойного бодрствования & \cite{Marrosu1995} & измерено \\
3 & \nt{NORA} и \nt{SERO} затихают в медленном сне и почти молчат в быстром & Запись отдельных \nt{NORA}-нейронов у крыс и \nt{SERO}-нейронов у кошек по стадиям сна: частота максимальна в бодрствовании, ниже в медленном сне и почти нулевая в быстром & \cite{AstonJonesBloom1981,McGintyHarper1976} & измерено \\
4 & В быстром сне мышцы выключены торможением через \nt{GABA} и \nt{GLYC} & У крыс мерили активность \nt{ACET}-нейронов жевательной мышцы и подавали блокаторы прямо к ним: паралич снимается, только если заблокированы и \nt{GABA}-рецепторы обоих типов, и \nt{GLYC}-рецепторы & \cite{BrooksPeever2012} & измерено \\
5 & Давление сна $S$ --- конденсатор с двумя порогами~\eqref{eq:sleeppressure}, $\tau_r=18.2$~ч, $\tau_d=4.2$~ч & Модель двух процессов; постоянные времени получены из того, как мощность медленных волн ЭЭГ растёт с бодрствованием и падает во сне, а форма порогов --- из времени спонтанного пробуждения & \cite{Borbely1982,Daan1984} & теория \\
6 & Машина сама выходит на $16.1$~ч бодрствования и $8.0$~ч сна (рис.~\ref{fig:twoprocess}) & Расчёт по~\eqref{eq:sleeppressure} с качающимися порогами, скрипт \texttt{models/sleep\_models.py}: $16.1$~ч бодрствования и $7.9$--$8.0$~ч сна & --- & модель \\
7 & После $40$~ч без сна $S=0.90$, но сон длится лишь $8.8$~ч: долг возвращается глубиной, а не длительностью & Модель: \texttt{models/sleep\_models.py} даёт $S=0.90$ и $8.8$~ч. Опыт: после $40.5$~ч без сна у восьми человек в первую восстановительную ночь мощность медленных волн ЭЭГ значимо выше обычной & \cite{Borbely1981} & модель \\
8 & Физически $S$ --- аденозин & Микродиализ у кошек: внеклеточный аденозин в одной из областей, управляющих бодрствованием, растёт за время бодрствования, растёт дальше при лишении сна и падает во время сна. Что $S$ --- только аденозин, не показано & \cite{PorkkaHeiskanen1997,Dunwiddie2001} & гипотеза \\
9 & В медленном сне кора качается: работа несколько сотен мс --- тишина, реже раза в секунду ($<1$~Гц, под наркозом чаще $0.3$--$0.4$~Гц) & Внутриклеточная запись у кошек под наркозом: деполяризация $0.8$--$1.5$~с и долгая гиперполяризация, ритм $<1$~Гц, чаще всего $0.3$--$0.4$~Гц; это же чередование работы и тишины видно в естественном сне (строка~1) & \cite{Steriade1993,Vyazovskiy2009} & измерено \\
10 & Качели --- группа с обратной связью и усталостью~\eqref{eq:updown}; при $b=5$ период $1.1$~с (значение модели) и работа около $35\%$ времени, при $b=1$ работа ровная (рис.~\ref{fig:updown}) & Расчёт, скрипт \texttt{models/sleep\_models.py}: период $1.11$~с, доля времени работы $0.35$ (среднее по целому числу периодов); при $b=1$ доля работы $1.0$. Период модели короче измеренного (строка~9) & --- & модель \\
11 & \nt{ACET} выключает качели и переводит кору в ровную работу & Стимуляция ядра \nt{ACET}-нейронов у кошки блокировала медленный ритм в $79\%$ нейронов коры и переводила их в ровную работу, в основном за счёт исчезновения долгих гиперполяризаций. Ацетилхолин подавляет медленные калиевые токи, создающие усталость & \cite{SteriadeAmzica1993,McCormick1992} & измерено \\
12 & Вспышки ритма $7$--$15$~Гц порождает петля \nt{GLUT}--\nt{GABA} между корой и релейным узлом & В срезах релейного узла зрения у хорька вспышки создаются ответными пачками релейных клеток на торможение от соседнего \nt{GABA}-ядра, которое они сами возбуждают & \cite{vonKrosigk1993,SteriadeMcCormickSejnowski1993} & измерено \\
13 & Повторы дня идут с перемоткой: в гиппокампе примерно в $20$ раз быстрее, в коре --- в $6$--$7$ раз & В медленном сне повторяющиеся последовательности нейронов гиппокампа идут сжатыми во времени. После бега по дорожке последовательности нейронов места повторялись в том же порядке короткими вспышками по $\sim100$~мс, сжатыми примерно в $20$ раз; в коре сжатие в $6$--$7$ раз & \cite{Nadasdy1999,LeeWilson2002,Euston2007} & измерено \\
14 & Усиление согласованности повторов с корой улучшает память (причинная связь) & У крыс после обучения стимуляцией, привязанной к повторам, усиливали их связь с медленными волнами и вспышками ритма в коре: на следующий день память была высокой, у контрольных --- на уровне случая & \cite{Maingret2016} & измерено \\
15 & Назначение повторов --- перемешанное обучение медленной памяти & Теория двух систем обучения и расчёты на искусственных сетях: последовательное обучение портит старое, перемешанное --- нет. Прямого опыта на мозге, что повторы нужны именно для этого, нет & \cite{McClelland1995} & гипотеза \\
16 & После сна синапсы уменьшаются пропорционально размеру, крупные почти не меняются & Трёхмерная электронная микроскопия $6920$ синапсов коры мыши: площадь контакта после сна на $\sim18\%$ меньше, чем после бодрствования, уменьшение пропорционально размеру, крупные устойчивые синапсы не меняются & \cite{deVivo2017} & измерено \\
17 & Главная задача сна --- перенормировка весов & Гипотеза синаптического гомеостаза; строки 1 и 16 с ней согласуются, но не доказывают, что это главная функция & \cite{TononiCirelli2014} & гипотеза \\
18 & Быстрый сон --- испытание модели мира на стенде & Режим (таблица~\ref{tab:sleepmodes}) измерен; что сеть прогоняет модель мира, --- теоретическое предположение, опыта нет & \cite{Hobson2009} & гипотеза \\
19 & Промывка мозга во сне: вопрос открыт & Один опыт: во сне межклеточное пространство на $60\%$ шире и уборка быстрее. Другой, с иным способом измерения: во сне и под наркозом уборка заметно медленнее. Результаты противоречат друг другу & \cite{Xie2013,Miao2024} & гипотеза \\
20 & Местный сон: участки коры бодрствующего животного ненадолго проваливаются в тишину & У крыс после долгого бодрствования нейроны отдельного участка коры коротко замолкают, как в медленном сне, с медленной волной в местной ЭЭГ; в такие моменты крыса чаще ошибается в задаче & \cite{Vyazovskiy2011} & измерено \\
\bottomrule
\end{longtable}}

\section{Глава~\ref{sec:livingmachine}. Машина из живых нейронов}

Поведение культур на решётках электродов измерено в прямых опытах с записью и стимуляцией; механизм залпов опирается на модель истощения синапсов и опыты на микрокультурах, а утверждения о дофамине в чашке --- на единичные опыты, часть которых сами авторы считают лишь демонстрацией.

{\footnotesize
\setlength{\tabcolsep}{3pt}\renewcommand{\arraystretch}{1.15}
\begin{longtable}{>{\raggedright}p{0.5cm}>{\raggedright}p{4.0cm}>{\raggedright}p{5.6cm}>{\raggedright}p{2.3cm}>{\raggedright\arraybackslash}p{1.7cm}}
\toprule
№ & Утверждение & Опыт и что измерено & Источник & Статус \\
\midrule\endfirsthead
\toprule
№ & Утверждение & Опыт и что измерено & Источник & Статус \\
\midrule\endhead
1 & Культура на чипе: в основном \nt{ГЛУТ}, примерно пятая часть \nt{ГАМК} & Сеть из тысяч нейронов на решётке электродов --- стандартный опыт (строка~3). Для доли \nt{ГАМК} «около пятой части» глава ссылки не даёт, и проверенного источника для этого числа в культурах коры мы не нашли; доля зависит от препарата: в культурах гиппокампа \nt{ГАМК}-нейронов около $6\%$ & \cite{Benson1994} & измерено \\
2 & Органоиды: трёхмерная нервная ткань размером около миллиметра из стволовых клеток человека & Из стволовых клеток вырастили трёхмерные органоиды с несколькими областями мозга, в том числе корой со слоями предшественников и зрелыми нейронами & \cite{Lancaster2013} & измерено \\
3 & Без входа культура каждые несколько секунд вспыхивает залпом & $58$ культур коры плотностью от $3\,000$ до $50\,000$ нейронов записывали пять недель ($963$ получасовые записи): после двух недель почти во всех культурах преобладают залпы всей сети; интервалы между залпами --- от $1$ до $300$~с и сильно зависят от препарата & \cite{Wagenaar2006} & измерено \\
4 & Залп кончается истощением медиатора, интервал $T_{\text{залп}}\approx\tau_{\text{вос}}\ln\frac{1}{1-x^*}$~\eqref{eq:burstinterval} & Формула выведена в главе. Модель: сеть с истощающимися синапсами сама порождает залпы всей сети. Опыт на микрокультурах из $4$--$30$ нейронов: длительность залпа зависит от времени после предыдущего, истощение пузырьков медиатора отменяет залпы; вывод авторов --- залп ограничен запасом медиатора. Восстановление там --- десятки секунд, а не $0.8$~с & вывод в главе; \cite{Tsodyks2000,CohenSegal2011,Tsodyks1997} & теория \\
5 & «Учитель, который замолкает»: стимуляция выключается, когда появился нужный ответ, и ответ закрепляется & Культуру стимулировали с частотой $0.3$--$1$~Гц, пока через $50\pm10$~мс после стимула не появится выбранный ответ, затем стимуляцию выключали на $5$~мин; после нескольких циклов нужный ответ вызывался сразу; строили кривые обучения & \cite{Shahaf2001} & измерено \\
6 & Культура играет в теннис; серии отбитых мячей удлиняются только при обратной связи & Подробно --- в главе~\ref{sec:dishbrain} и приложении к ней & \cite{Kagan2022} & измерено \\
7 & Что культура учится предсказывать свой вход --- гипотеза; громкие слова о «разумности» оспорены & Принцип свободной энергии --- теоретическое предложение; прямого опыта, отличающего его от более простых объяснений, нет. Тридцать исследователей в ответной статье указали, что изменение поведения в петле не показывает ни разума, ни ощущений & \cite{Friston2010,Balci2023} & гипотеза \\
8 & Культура водит виртуальное тело к цели при подобранном рисунке стимулов & Программа сама подбирала обучающие рисунки стимуляции по текущему результату; за десятки минут сеть достигала заданных состояний, пластичность держалась $80$~мин после $2$~ч обучения. Те же стимулы, поданные без связи с поведением, не действовали & \cite{Bakkum2008} & измерено \\
9 & Резервуар: обучается только линейное считывание $y=W_{\text{вых}}\mathbf r$~\eqref{eq:reservoir} & Теория вычисления на резервуаре: любая нелинейная система с затухающей памятью плюс обученный линейный выход & \cite{Maass2002,JaegerHaas2004} & теория \\
10 & Органоид как резервуар различает говорящих с точностью около $78\%$ & Звуки речи нескольких говорящих подавали органоиду рисунками тока на решётке электродов; обученное считывание различало говорящего с точностью около $78\%$. Оговорка: авторы сообщают, что при тренировке менялась и функциональная связность органоида, и называют это обучением без учителя в самом органоиде, --- обучение не целиком снаружи & \cite{Cai2023} & измерено \\
11 & Миллион нейронов тратит на импульсы около $10^{-4}$~Вт~\eqref{eq:dishpower} & Оценка: цена импульса $10^{-10}$~Дж из энергетического бюджета коры~\eqref{eq:energy}, умноженная на число импульсов; прямого измерения мощности культуры нет & вывод в главе; \cite{Attwell2001} & теория \\
12 & Дофамин по вспышке света: $1$~мМ клеточного дофамина, $240$ повторов, число вызванных импульсов росло & На удалённой платформе с органоидами дофамин в светочувствительной клетке ($1$~мМ) освобождали ультрафиолетом ($365$~нм, $0.8$~с), когда стимул вызывал больше импульсов; за $240$ шагов (около $5$~ч) число импульсов в целом росло. Авторы пишут, что по одному опыту связь с дофамином установить нельзя; контролей нет & \cite{Jordan2024} & гипотеза \\
13 & Дофамин от живых клеток меняет вес органического транзистора надолго и обратимо & Клетки, выделяющие дофамин, выращены над органическим нейроморфным элементом; окисление дофамина у электрода меняет проводимость канала; показаны долговременное изменение веса и его возврат & \cite{Keene2020} & измерено \\
14 & Органоиды с работающими \nt{ДОФА}-нейронами существуют; как учитель их дофамин не использовали & В органоидах из стволовых клеток человека найдены электрически активные зрелые \nt{ДОФА}-нейроны и выработка дофамина. Опытов, где этот дофамин учит другую сеть, в литературе мы не нашли & \cite{Jo2016} & измерено \\
15 & Органоид балансирует шестом: учебные стимулы от программы лучше случайных, без закрепления, через \nt{ГЛУТ}-синапсы & Органоиды коры мыши в замкнутой петле с задачей «шест на тележке»: у большинства органоидов сигналы, выбранные обучением с подкреплением, дали лучший результат, чем случайные или никакие; после $45$~мин отдыха улучшение не сохранялось; блокада AMPA- и NMDA-рецепторов его отменяла & \cite{Robbins2026} & измерено \\
16 & Без регистров управления сеть на стенде не может сама решить, что считать успехом (утверждение~\ref{prop:livingmachine}) & Во всех опытах строк 5--15 учебный сигнал задаёт экспериментатор или программа; опыта, где культура сама выработала бы критерий успеха, нет --- это вывод из отсутствия, а не опыт & --- & гипотеза \\
\bottomrule
\end{longtable}}

\section{Глава~\ref{sec:dishbrain}. DishBrain}

Устройство стенда и результаты игры взяты из одной экспериментальной работы и её продолжения; формулы шума и дрейфа к хорошим настройкам выведены в главе и проверены моделью книги, а механизм обучения в чашке не установлен.

{\footnotesize
\setlength{\tabcolsep}{3pt}\renewcommand{\arraystretch}{1.15}
\begin{longtable}{>{\raggedright}p{0.5cm}>{\raggedright}p{4.0cm}>{\raggedright}p{5.6cm}>{\raggedright}p{2.3cm}>{\raggedright\arraybackslash}p{1.7cm}}
\toprule
№ & Утверждение & Опыт и что измерено & Источник & Статус \\
\midrule\endfirsthead
\toprule
№ & Утверждение & Опыт и что измерено & Источник & Статус \\
\midrule\endhead
1 & Стенд: около $800\,000$ клеток коры мыши или около $10^6$ клеток человека на чипе; $26\,000$ электродов на $8$~мм$^2$, запись до $1024$, стимуляция через $8$ & Методы работы: чип MaxOne, $26\,000$ платиновых электродов на $8$~мм$^2$, запись $1024$ каналов по $20\,000$ отсчётов в секунду; стимулировать технически можно до $32$ электродов, независимо удалось $8$; мышиные культуры использовали примерно с $10$~суток & \cite{Kagan2022} & измерено \\
2 & Петля с тактом $10$~мс: импульсы двигательных областей двигают ракетку (рис.~\ref{fig:dishloop}) & Импульсы считаются окнами по $10$~мс; задержка от импульса до ответного стимула около $5$~мс, её задаёт буфер аппаратуры & \cite{Kagan2022} & измерено \\
3 & Вход: код местом (какой из $8$ электродов) и частотой $4$--$40$~имп/с & В основном опыте частота стимуляции растёт линейно от $4$~Гц, когда мяч у дальней стенки, до $40$~Гц у ракетки; амплитуда импульсов мяча --- $75$~мВ; импульсы двухфазные прямоугольные & \cite{Kagan2022} & измерено \\
4 & Выход $\Delta p=c\,(n_\uparrow-n_\downarrow)$~\eqref{eq:dishreadout}, шум счёта за окно $1/\sqrt{RT}$ & Правило разности двух областей описано в работе (с весами областей по активности), точная формула не дана. Оценка шума --- следствие пуассоновского счёта, выведена в главе & \cite{Kagan2022}; вывод в главе & теория \\
5 & Учитель: после удара --- предсказуемый стимул $100$~имп/с, $0.1$~с; после промаха --- непредсказуемый, $150$~мВ, $5$~имп/с, $4$~с & Методы: после удара $75$~мВ, $100$~Гц, $100$~мс на все $8$ электродов сразу; после промаха $150$~мВ, $5$~Гц в случайные электроды в случайные моменты в течение $4$~с. Дофамина в культуре нет & \cite{Kagan2022} & измерено \\
6 & Сеть действует так, чтобы её вход стал предсказуемым & Принцип свободной энергии --- теоретическое предложение, из которого авторы выводили протокол; опыт с ним согласуется, но не отличает его от более простых механизмов (строки~7--8) & \cite{Friston2010,Kagan2022} & гипотеза \\
7 & Если неудача встряхивает сеть, а удача нет, время в настройке $\rho\propto1/P_{\text{пр}}$~\eqref{eq:dishdensity} (утверждение~\ref{prop:dishscramble}) & Следует из баланса потоков в марковской цепи с симметричными толчками, выведено в главе. Прямой проверки на культуре нет & вывод в главе & теория \\
8 & Модель «перемешивание при промахе» учится: доля ударов $0.21\to0.38$; при толчке после каждого розыгрыша $\approx0.12$, без толчков $0.19$ (рис.~\ref{fig:dishmodel}) & Расчёт, скрипт \texttt{models/dishbrain\_model.py}, $40$ моделей-культур по $2000$ розыгрышей: $0.21\to0.38$; $0.13\to0.11$; $0.19\to0.19$ & --- & модель \\
9 & Серии отбитых мячей удлиняются только у живых культур с обратной связью; контроли не учатся & $399$ сеансов по $20$~мин: $9$ культур мыши ($101$ сеанс), $11$ человека ($138$), $6$ чипов со средой ($80$), $20$ культур в покое ($42$), $3$ симуляции ($38$). Средняя длина розыгрыша за последние $15$~мин больше, чем за первые $5$, только у живых культур; у клеток, не являющихся нейронами (HEK293T), и у чипов со средой эффекта нет & \cite{Kagan2022} & измерено \\
10 & Без обратной связи и при тишине вместо стимула улучшения нет & $486$ сеансов за $3$ дня: «стимул» ($n=27$), «тишина» ($17$), «без обратной связи» ($15$). Рост во времени значим только в условии «стимул». Оговорка: в конце сеанса условие «тишина» всё же превосходило «без обратной связи», хоть и с меньшим эффектом & \cite{Kagan2022} & измерено \\
11 & Эффект невелик; учится ли сеть надолго --- открытый вопрос & Внутри сеанса улучшение надёжно, но, по словам самих авторов, обучение между сеансами за несколько дней устойчиво не наблюдалось & \cite{Kagan2022} & измерено \\
12 & Во время игры сеть приближается к критическому режиму, и чем ближе, тем лучше игра & Анализ лавин активности тех же культур: структурированный вход приближает сеть к критическому режиму, качество игры коррелирует с близостью к нему; но одной критичности без сведений о последствиях действий для обучения мало & \cite{Habibollahi2023} & измерено \\
13 & «Разумность» культуры не показана & Ответная статья тридцати исследователей: изменение поведения в замкнутой петле не доказывает ни разума, ни ощущений; опыта, который это показывал бы, нет & \cite{Balci2023} & гипотеза \\
14 & Предложенный четвёртый контроль: предсказуемый стимул после промаха той же длительности и энергии (раздел~\ref{sec:dishspec}) & Такой опыт не поставлен; это предложение книги & --- & гипотеза \\
\bottomrule
\end{longtable}}


\begin{thebibliography}{99}
\bibitem{Azevedo2009} F. A. C. Azevedo et al., \emph{Equal numbers of neuronal and nonneuronal cells make the human brain an isometrically scaled-up primate brain}, J. Comp. Neurol. \textbf{513}, 532 (2009).
\bibitem{Schultz1997} W. Schultz, P. Dayan, and P. R. Montague, \emph{A neural substrate of prediction and reward}, Science \textbf{275}, 1593 (1997).
\bibitem{SuttonBarto2018} R. S. Sutton and A. G. Barto, \emph{Reinforcement Learning: An Introduction}, 2nd ed. (MIT Press, Cambridge, MA, 2018).
\bibitem{Doya2002} K. Doya, \emph{Metalearning and neuromodulation}, Neural Networks \textbf{15}, 495 (2002).
\bibitem{Strata1999} P. Strata and R. Harvey, \emph{Dale's principle}, Brain Res. Bull. \textbf{50}, 349 (1999).
\bibitem{vanVreeswijk1996} C. van Vreeswijk and H. Sompolinsky, \emph{Chaos in neuronal networks with balanced excitatory and inhibitory activity}, Science \textbf{274}, 1724 (1996).
\bibitem{AstonJones2005} G. Aston-Jones and J. D. Cohen, \emph{An integrative theory of locus coeruleus-norepinephrine function: adaptive gain and optimal performance}, Annu. Rev. Neurosci. \textbf{28}, 403 (2005).
\bibitem{Beaulieu2011} J.-M. Beaulieu and R. R. Gainetdinov, \emph{The physiology, signaling, and pharmacology of dopamine receptors}, Pharmacol. Rev. \textbf{63}, 182 (2011).
\bibitem{Sparso2001} J. Spars{\o} and S. Furber (eds.), \emph{Principles of Asynchronous Circuit Design: A Systems Perspective} (Kluwer, Boston, 2001).
\bibitem{Carr1990} C. E. Carr and M. Konishi, \emph{A circuit for detection of interaural time differences in the brain stem of the barn owl}, J. Neurosci. \textbf{10}, 3227 (1990).
\bibitem{Wang2001} X.-J. Wang, \emph{Synaptic reverberation underlying mnemonic persistent activity}, Trends Neurosci. \textbf{24}, 455 (2001).
\bibitem{Hahnloser2002} R. H. R. Hahnloser, A. A. Kozhevnikov, and M. S. Fee, \emph{An ultra-sparse code underlies the generation of neural sequences in a songbird}, Nature \textbf{419}, 65 (2002).
\bibitem{Barlow1965} H. B. Barlow and W. R. Levick, \emph{The mechanism of directionally selective units in rabbit's retina}, J. Physiol. \textbf{178}, 477 (1965).
\bibitem{Carandini2012} M. Carandini and D. J. Heeger, \emph{Normalization as a canonical neural computation}, Nat. Rev. Neurosci. \textbf{13}, 51 (2012).
\bibitem{Pouille2001} F. Pouille and M. Scanziani, \emph{Enforcement of temporal fidelity in pyramidal cells by somatic feed-forward inhibition}, Science \textbf{293}, 1159 (2001).
\bibitem{Tsodyks1997isn} M. V. Tsodyks, W. E. Skaggs, T. J. Sejnowski, and B. L. McNaughton, \emph{Paradoxical effects of external modulation of inhibitory interneurons}, J. Neurosci. \textbf{17}, 4382 (1997).
\bibitem{Buzsaki2012} G. Buzs\'aki and X.-J. Wang, \emph{Mechanisms of gamma oscillations}, Annu. Rev. Neurosci. \textbf{35}, 203 (2012).
\bibitem{Rieke1997} F. Rieke, D. Warland, R. de Ruyter van Steveninck, and W. Bialek, \emph{Spikes: Exploring the Neural Code} (MIT Press, Cambridge, MA, 1997).
\bibitem{Mainen1995} Z. F. Mainen and T. J. Sejnowski, \emph{Reliability of spike timing in neocortical neurons}, Science \textbf{268}, 1503 (1995).
\bibitem{Softky1993} W. R. Softky and C. Koch, \emph{The highly irregular firing of cortical cells is inconsistent with temporal integration of random EPSPs}, J. Neurosci. \textbf{13}, 334 (1993).
\bibitem{Thorpe1996} S. Thorpe, D. Fize, and C. Marlot, \emph{Speed of processing in the human visual system}, Nature \textbf{381}, 520 (1996).
\bibitem{Zohary1994} E. Zohary, M. N. Shadlen, and W. T. Newsome, \emph{Correlated neuronal discharge rate and its implications for psychophysical performance}, Nature \textbf{370}, 140 (1994).
\bibitem{MorenoBote2014} R. Moreno-Bote et al., \emph{Information-limiting correlations}, Nat. Neurosci. \textbf{17}, 1410 (2014).
\bibitem{Gollisch2008} T. Gollisch and M. Meister, \emph{Rapid neural coding in the retina with relative spike latencies}, Science \textbf{319}, 1108 (2008).
\bibitem{OKeefe1993} J. O'Keefe and M. L. Recce, \emph{Phase relationship between hippocampal place units and the EEG theta rhythm}, Hippocampus \textbf{3}, 317 (1993).
\bibitem{Destexhe2003} A. Destexhe, M. Rudolph, and D. Par\'e, \emph{The high-conductance state of neocortical neurons in vivo}, Nat. Rev. Neurosci. \textbf{4}, 739 (2003).
\bibitem{Tsodyks1997} M. V. Tsodyks and H. Markram, \emph{The neural code between neocortical pyramidal neurons depends on neurotransmitter release probability}, Proc. Natl. Acad. Sci. USA \textbf{94}, 719 (1997).
\bibitem{Lisman1997} J. E. Lisman, \emph{Bursts as a unit of neural information: making unreliable synapses reliable}, Trends Neurosci. \textbf{20}, 38 (1997).
\bibitem{Attwell2001} D. Attwell and S. B. Laughlin, \emph{An energy budget for signaling in the grey matter of the brain}, J. Cereb. Blood Flow Metab. \textbf{21}, 1133 (2001).
\bibitem{Lennie2003} P. Lennie, \emph{The cost of cortical computation}, Curr. Biol. \textbf{13}, 493 (2003).
\bibitem{Yagishita2014} S. Yagishita et al., \emph{A critical time window for dopamine actions on the structural plasticity of dendritic spines}, Science \textbf{345}, 1616 (2014).
\bibitem{Werfel2005} J. Werfel, X. Xie, and H. S. Seung, \emph{Learning curves for stochastic gradient descent in linear feedforward networks}, Neural Comput. \textbf{17}, 2699 (2005).
\bibitem{Doya2000} K. Doya, \emph{Reinforcement learning in continuous time and space}, Neural Comput. \textbf{12}, 219 (2000).
\bibitem{Oja1982} E. Oja, \emph{Simplified neuron model as a principal component analyzer}, J. Math. Biol. \textbf{15}, 267 (1982).
\bibitem{BiPoo1998} G.-Q. Bi and M.-M. Poo, \emph{Synaptic modifications in cultured hippocampal neurons: dependence on spike timing, synaptic strength, and postsynaptic cell type}, J. Neurosci. \textbf{18}, 10464 (1998).
\bibitem{Williams1992} R. J. Williams, \emph{Simple statistical gradient-following algorithms for connectionist reinforcement learning}, Mach. Learn. \textbf{8}, 229 (1992).
\bibitem{Bayer2005} H. M. Bayer and P. W. Glimcher, \emph{Midbrain dopamine neurons encode a quantitative reward prediction error signal}, Neuron \textbf{47}, 129 (2005).
\bibitem{Shen2008} W. Shen, M. Flajolet, P. Greengard, and D. J. Surmeier, \emph{Dichotomous dopaminergic control of striatal synaptic plasticity}, Science \textbf{321}, 848 (2008).
\bibitem{Dabney2020} W. Dabney, Z. Kurth-Nelson, N. Uchida, C. K. Starkweather, D. Hassabis, R. Munos, and M. Botvinick, \emph{A distributional code for value in dopamine-based reinforcement learning}, Nature \textbf{577}, 671 (2020).
\bibitem{Matsumoto2009} M. Matsumoto and O. Hikosaka, \emph{Two types of dopamine neuron distinctly convey positive and negative motivational signals}, Nature \textbf{459}, 837 (2009).
\bibitem{BrombergMartin2010} E. S. Bromberg-Martin, M. Matsumoto, and O. Hikosaka, \emph{Dopamine in motivational control: rewarding, aversive, and alerting}, Neuron \textbf{68}, 815 (2010).
\bibitem{Menegas2018} W. Menegas, K. Akiti, R. Amo, N. Uchida, and M. Watabe-Uchida, \emph{Dopamine neurons projecting to the posterior striatum reinforce avoidance of threatening stimuli}, Nat. Neurosci. \textbf{21}, 1421 (2018).
\bibitem{Niv2007} Y. Niv, N. D. Daw, D. Joel, and P. Dayan, \emph{Tonic dopamine: opportunity costs and the control of response vigor}, Psychopharmacology \textbf{191}, 507 (2007).
\bibitem{Mohebi2019} A. Mohebi et al., \emph{Dissociable dopamine dynamics for learning and motivation}, Nature \textbf{570}, 65 (2019).
\bibitem{Hamid2016} A. A. Hamid et al., \emph{Mesolimbic dopamine signals the value of work}, Nat. Neurosci. \textbf{19}, 117 (2016).
\bibitem{Kim2020} H. R. Kim, A. N. Malik et al., \emph{A unified framework for dopamine signals across timescales}, Cell \textbf{183}, 1600 (2020).
\bibitem{Jeong2022} H. Jeong et al., \emph{Mesolimbic dopamine release conveys causal associations}, Science \textbf{378}, eabq6740 (2022).
\bibitem{Amo2022} R. Amo et al., \emph{A gradual temporal shift of dopamine responses mirrors the progression of temporal difference error in machine learning}, Nat. Neurosci. \textbf{25}, 1082 (2022).
\bibitem{Takahashi2017} Y. K. Takahashi et al., \emph{Dopamine neurons respond to errors in the prediction of sensory features of expected rewards}, Neuron \textbf{95}, 1395 (2017).
\bibitem{Sharpe2017} M. J. Sharpe et al., \emph{Dopamine transients are sufficient and necessary for acquisition of model-based associations}, Nat. Neurosci. \textbf{20}, 735 (2017).
\bibitem{Gardner2018} M. P. H. Gardner, G. Schoenbaum, and S. J. Gershman, \emph{Rethinking dopamine as generalized prediction error}, Proc. R. Soc. B \textbf{285}, 20181645 (2018).
\bibitem{Engelhard2019} B. Engelhard et al., \emph{Specialized coding of sensory, motor and cognitive variables in VTA dopamine neurons}, Nature \textbf{570}, 509 (2019).
\bibitem{Clements1992} J. D. Clements, R. A. J. Lester, G. Tong, C. E. Jahr, and G. L. Westbrook, \emph{The time course of glutamate in the synaptic cleft}, Science \textbf{258}, 1498 (1992).
\bibitem{Hasselmo2006} M. E. Hasselmo, \emph{The role of acetylcholine in learning and memory}, Curr. Opin. Neurobiol. \textbf{16}, 710 (2006).
\bibitem{JacobsAzmitia1992} B. L. Jacobs and E. C. Azmitia, \emph{Structure and function of the brain serotonin system}, Physiol. Rev. \textbf{72}, 165 (1992).
\bibitem{Schiller1992} P. H. Schiller, \emph{The ON and OFF channels of the visual system}, Trends Neurosci. \textbf{15}, 86 (1992).
\bibitem{Grothe2010} B. Grothe, M. Pecka, and D. McAlpine, \emph{Mechanisms of sound localization in mammals}, Physiol. Rev. \textbf{90}, 983 (2010).
\bibitem{Markram2004} H. Markram et al., \emph{Interneurons of the neocortical inhibitory system}, Nat. Rev. Neurosci. \textbf{5}, 793 (2004).
\bibitem{Joshi2016} S. Joshi, Y. Li, R. M. Kalwani, and J. I. Gold, \emph{Relationships between pupil diameter and neuronal activity in the locus coeruleus, colliculi, and cingulate cortex}, Neuron \textbf{89}, 221 (2016).
\bibitem{ServanSchreiber1990} D. Servan-Schreiber, H. Printz, and J. D. Cohen, \emph{A network model of catecholamine effects: gain, signal-to-noise ratio, and behavior}, Science \textbf{249}, 892 (1990).
\bibitem{YuDayan2005} A. J. Yu and P. Dayan, \emph{Uncertainty, neuromodulation, and attention}, Neuron \textbf{46}, 681 (2005).
\bibitem{Nassar2012} M. R. Nassar et al., \emph{Rational regulation of learning dynamics by pupil-linked arousal systems}, Nat. Neurosci. \textbf{15}, 1040 (2012).
\bibitem{BouretSara2005} S. Bouret and S. J. Sara, \emph{Network reset: a simplified overarching theory of locus coeruleus noradrenaline function}, Trends Neurosci. \textbf{28}, 574 (2005).
\bibitem{Hebb1949} D. O. Hebb, \emph{The Organization of Behavior: A Neuropsychological Theory} (Wiley, New York, 1949).
\bibitem{MacKay1952} D. M. MacKay and W. S. McCulloch, \emph{The limiting information capacity of a neuronal link}, Bull. Math. Biophys. \textbf{14}, 127 (1952).
\bibitem{WilsonCowan1972} H. R. Wilson and J. D. Cowan, \emph{Excitatory and inhibitory interactions in localized populations of model neurons}, Biophys. J. \textbf{12}, 1 (1972).
\bibitem{AllenStevens1994} C. Allen and C. F. Stevens, \emph{An evaluation of causes for unreliability of synaptic transmission}, Proc. Natl. Acad. Sci. USA \textbf{91}, 10380 (1994).
\bibitem{Barlow1961} H. B. Barlow, \emph{Possible principles underlying the transformations of sensory messages}, in \emph{Sensory Communication}, ed. W. A. Rosenblith (MIT Press, Cambridge, MA, 1961), pp.~217--234.
\bibitem{Cardin2009} J. A. Cardin et al., \emph{Driving fast-spiking cells induces gamma rhythm and controls sensory responses}, Nature \textbf{459}, 663 (2009).
\bibitem{Lillicrap2020} T. P. Lillicrap, A. Santoro, L. Marris, C. J. Akerman, and G. Hinton, \emph{Backpropagation and the brain}, Nat. Rev. Neurosci. \textbf{21}, 335 (2020).
\bibitem{Fremaux2016} N. Fr\'emaux and W. Gerstner, \emph{Neuromodulated spike-timing-dependent plasticity, and theory of three-factor learning rules}, Front. Neural Circuits \textbf{9}, 85 (2016).
\bibitem{Haider2006} B. Haider, A. Duque, A. R. Hasenstaub, and D. A. McCormick, \emph{Neocortical network activity in vivo is generated through a dynamic balance of excitation and inhibition}, J. Neurosci. \textbf{26}, 4535 (2006).
\bibitem{HoltKoch1997} G. R. Holt and C. Koch, \emph{Shunting inhibition does not have a divisive effect on firing rates}, Neural Comput. \textbf{9}, 1001 (1997).
\bibitem{Chance2002} F. S. Chance, L. F. Abbott, and A. D. Reyes, \emph{Gain modulation from background synaptic input}, Neuron \textbf{35}, 773 (2002).
\bibitem{Ozeki2009} H. Ozeki, I. M. Finn, E. S. Schaffer, K. D. Miller, and D. Ferster, \emph{Inhibitory stabilization of the cortical network underlies visual surround suppression}, Neuron \textbf{62}, 578 (2009).
\bibitem{HasselmoBower1992} M. E. Hasselmo and J. M. Bower, \emph{Cholinergic suppression specific to intrinsic not afferent fiber synapses in rat piriform (olfactory) cortex}, J. Neurophysiol. \textbf{67}, 1222 (1992).
\bibitem{HasselmoAndersonBower1992} M. E. Hasselmo, B. P. Anderson, and J. M. Bower, \emph{Cholinergic modulation of cortical associative memory function}, J. Neurophysiol. \textbf{67}, 1230 (1992).
\bibitem{Hasselmo1999} M. E. Hasselmo, \emph{Neuromodulation: acetylcholine and memory consolidation}, Trends Cogn. Sci. \textbf{3}, 351 (1999).
\bibitem{Tsodyks1988} M. V. Tsodyks and M. V. Feigel'man, \emph{The enhanced storage capacity in neural networks with low activity level}, Europhys. Lett. \textbf{6}, 101 (1988).
\bibitem{KalmanBucy1961} R. E. Kalman and R. S. Bucy, \emph{New results in linear filtering and prediction theory}, J. Basic Eng. \textbf{83}, 95 (1961).
\bibitem{WilsonMcNaughton1994} M. A. Wilson and B. L. McNaughton, \emph{Reactivation of hippocampal ensemble memories during sleep}, Science \textbf{265}, 676 (1994).
\bibitem{BarnesSharp1999} N. M. Barnes and T. Sharp, \emph{A review of central 5-HT receptors and their function}, Neuropharmacology \textbf{38}, 1083 (1999).
\bibitem{Bunin1998} M. A. Bunin and R. M. Wightman, \emph{Quantitative evaluation of 5-hydroxytryptamine (serotonin) neuronal release and uptake: an investigation of extrasynaptic transmission}, J. Neurosci. \textbf{18}, 4854 (1998).
\bibitem{Sprouse1987} J. S. Sprouse and G. K. Aghajanian, \emph{Electrophysiological responses of serotoninergic dorsal raphe neurons to 5-HT$_{1A}$ and 5-HT$_{1B}$ agonists}, Synapse \textbf{1}, 3 (1987).
\bibitem{Sykova2008} E. Syková and C. Nicholson, \emph{Diffusion in brain extracellular space}, Physiol. Rev. \textbf{88}, 1277 (2008).
\bibitem{WilsonNicoll2001} R. I. Wilson and R. A. Nicoll, \emph{Endogenous cannabinoids mediate retrograde signalling at hippocampal synapses}, Nature \textbf{410}, 588 (2001).
\bibitem{Garthwaite2008} J. Garthwaite, \emph{Concepts of neural nitric oxide-mediated transmission}, Eur. J. Neurosci. \textbf{27}, 2783 (2008).
\bibitem{vandenPol2012} A. N. van den Pol, \emph{Neuropeptide transmission in brain circuits}, Neuron \textbf{76}, 98 (2012).
\bibitem{Dunwiddie2001} T. V. Dunwiddie and S. A. Masino, \emph{The role and regulation of adenosine in the central nervous system}, Annu. Rev. Neurosci. \textbf{24}, 31 (2001).
\bibitem{Skaggs1996} W. E. Skaggs and B. L. McNaughton, \emph{Replay of neuronal firing sequences in rat hippocampus during sleep following spatial experience}, Science \textbf{271}, 1870 (1996).
\bibitem{Baylor1979} D. A. Baylor, T. D. Lamb, and K.-W. Yau, \emph{Responses of retinal rods to single photons}, J. Physiol. \textbf{288}, 613 (1979).
\bibitem{Hudspeth1989} A. J. Hudspeth, \emph{How the ear's works work}, Nature \textbf{341}, 397 (1989).
\bibitem{NakaRushton1966} K. I. Naka and W. A. H. Rushton, \emph{S-potentials from colour units in the retina of fish (Cyprinidae)}, J. Physiol. \textbf{185}, 536 (1966).
\bibitem{Lichtsteiner2008} P. Lichtsteiner, C. Posch, and T. Delbruck, \emph{A $128\times128$ 120 dB 15 $\mu$s latency asynchronous temporal contrast vision sensor}, IEEE J. Solid-State Circuits \textbf{43}, 566 (2008).
\bibitem{Koch2006} K. Koch et al., \emph{How much the eye tells the brain}, Curr. Biol. \textbf{16}, 1428 (2006).
\bibitem{Robles2001} L. Robles and M. A. Ruggero, \emph{Mechanics of the mammalian cochlea}, Physiol. Rev. \textbf{81}, 1305 (2001).
\bibitem{Mombaerts2004} P. Mombaerts, \emph{Genes and ligands for odorant, vomeronasal and taste receptors}, Nat. Rev. Neurosci. \textbf{5}, 263 (2004).
\bibitem{WoodSlater2001} S. J. Wood and C. R. Slater, \emph{Safety factor at the neuromuscular junction}, Prog. Neurobiol. \textbf{64}, 393 (2001).
\bibitem{Henneman1957} E. Henneman, \emph{Relation between size of neurons and their susceptibility to discharge}, Science \textbf{126}, 1345 (1957).
\bibitem{HarrisWolpert1998} C. M. Harris and D. M. Wolpert, \emph{Signal-dependent noise determines motor planning}, Nature \textbf{394}, 780 (1998).
\bibitem{Hogan1984} N. Hogan, \emph{Adaptive control of mechanical impedance by coactivation of antagonist muscles}, IEEE Trans. Autom. Control \textbf{29}, 681 (1984).
\bibitem{McAuleyMarsden2000} J. H. McAuley and C. D. Marsden, \emph{Physiological and pathological tremors and rhythmic central motor control}, Brain \textbf{123}, 1545 (2000).
\bibitem{WolpertGhahramani2000} D. M. Wolpert and Z. Ghahramani, \emph{Computational principles of movement neuroscience}, Nat. Neurosci. \textbf{3}, 1212 (2000).
\bibitem{MarderBucher2001} E. Marder and D. Bucher, \emph{Central pattern generators and the control of rhythmic movements}, Curr. Biol. \textbf{11}, R986 (2001).
\bibitem{Miyazaki2014} K. W. Miyazaki et al., \emph{Optogenetic activation of dorsal raphe serotonin neurons enhances patience for future rewards}, Curr. Biol. \textbf{24}, 2033 (2014).
\bibitem{Fuglevand1993} A. J. Fuglevand, D. A. Winter, and A. E. Patla, \emph{Models of recruitment and rate coding organization in motor-unit pools}, J. Neurophysiol. \textbf{70}, 2470 (1993).
\bibitem{Curcio1990} C. A. Curcio, K. R. Sloan, R. E. Kalina, and A. E. Hendrickson, \emph{Human photoreceptor topography}, J. Comp. Neurol. \textbf{292}, 497 (1990).
\bibitem{Hayes1950} N. D. Hayes, \emph{Roots of the transcendental equation associated with a certain difference-differential equation}, J. London Math. Soc. \textbf{25}, 226 (1950).
\bibitem{MelzackWall1965} R. Melzack and P. D. Wall, \emph{Pain mechanisms: a new theory}, Science \textbf{150}, 971 (1965).
\bibitem{Fields2004} H. Fields, \emph{State-dependent opioid control of pain}, Nat. Rev. Neurosci. \textbf{5}, 565 (2004).
\bibitem{Levine1978} J. D. Levine, N. C. Gordon, and H. L. Fields, \emph{The mechanism of placebo analgesia}, Lancet \textbf{312}, 654 (1978).
\bibitem{Woolf1983} C. J. Woolf, \emph{Evidence for a central component of post-injury pain hypersensitivity}, Nature \textbf{306}, 686 (1983).
\bibitem{Seymour2004} B. Seymour et al., \emph{Temporal difference models describe higher-order learning in humans}, Nature \textbf{429}, 664 (2004).
\bibitem{EcclestonCrombez1999} C. Eccleston and G. Crombez, \emph{Pain demands attention: a cognitive-affective model of the interruptive function of pain}, Psychol. Bull. \textbf{125}, 356 (1999).
\bibitem{WidrowHoff1960} B. Widrow and M. E. Hoff, \emph{Adaptive switching circuits}, IRE WESCON Convention Record, part 4, 96 (1960).
\bibitem{Marr1969} D. Marr, \emph{A theory of cerebellar cortex}, J. Physiol. \textbf{202}, 437 (1969).
\bibitem{Albus1971} J. S. Albus, \emph{A theory of cerebellar function}, Math. Biosci. \textbf{10}, 25 (1971).
\bibitem{Cover1965} T. M. Cover, \emph{Geometrical and statistical properties of systems of linear inequalities with applications in pattern recognition}, IEEE Trans. Electron. Comput. \textbf{EC-14}, 326 (1965).
\bibitem{Ito2001} M. Ito, \emph{Cerebellar long-term depression: characterization, signal transduction, and functional roles}, Physiol. Rev. \textbf{81}, 1143 (2001).
\bibitem{Suvrathan2016} A. Suvrathan, H. L. Payne, and J. L. Raymond, \emph{Timing rules for synaptic plasticity matched to behavioral function}, Neuron \textbf{92}, 959 (2016).
\bibitem{Fujita1982} M. Fujita, \emph{Adaptive filter model of the cerebellum}, Biol. Cybern. \textbf{45}, 195 (1982).
\bibitem{Blakemore1998} S.-J. Blakemore, D. M. Wolpert, and C. D. Frith, \emph{Central cancellation of self-produced tickle sensation}, Nat. Neurosci. \textbf{1}, 635 (1998).
\bibitem{Smith2006} M. A. Smith, A. Ghazizadeh, and R. Shadmehr, \emph{Interacting adaptive processes with different timescales underlie short-term motor learning}, PLoS Biol. \textbf{4}, e179 (2006).
\bibitem{Kording2007} K. P. K\"ording, J. B. Tenenbaum, and R. Shadmehr, \emph{The dynamics of memory as a consequence of optimal adaptation to a changing body}, Nat. Neurosci. \textbf{10}, 779 (2007).
\bibitem{Yao2023} Z. Yao et al., \emph{A high-resolution transcriptomic and spatial atlas of cell types in the whole mouse brain}, Nature \textbf{624}, 317 (2023).
\bibitem{Sanzeni2020} A. Sanzeni, B. Akitake, H. C. Goldbach, C. E. Leedy, N. Brunel, and M. H. Histed, \emph{Inhibition stabilization is a widespread property of cortical networks}, eLife \textbf{9}, e54875 (2020).
\bibitem{Padamsey2022} Z. Padamsey, D. Katsanevaki, N. Dupuy, and N. L. Rochefort, \emph{Neocortex saves energy by reducing coding precision during food scarcity}, Neuron \textbf{110}, 280 (2022).
\bibitem{Stringer2019} C. Stringer, M. Pachitariu, N. Steinmetz, C. B. Reddy, M. Carandini, and K. D. Harris, \emph{Spontaneous behaviors drive multidimensional, brainwide activity}, Science \textbf{364}, eaav7893 (2019).
\bibitem{Bittner2017} K. C. Bittner, A. D. Milstein, C. Grienberger, S. Romani, and J. C. Magee, \emph{Behavioral time scale synaptic plasticity underlies CA1 place fields}, Science \textbf{357}, 1033 (2017).
\bibitem{WhittingtonBogacz2019} J. C. R. Whittington and R. Bogacz, \emph{Theories of error back-propagation in the brain}, Trends Cogn. Sci. \textbf{23}, 235 (2019).
\bibitem{Reimer2016} J. Reimer et al., \emph{Pupil fluctuations track rapid changes in adrenergic and cholinergic activity in cortex}, Nat. Commun. \textbf{7}, 13289 (2016).
\bibitem{Beniaguev2021} D. Beniaguev, I. Segev, and M. London, \emph{Single cortical neurons as deep artificial neural networks}, Neuron \textbf{109}, 2727 (2021).
\bibitem{Gidon2020} A. Gidon et al., \emph{Dendritic action potentials and computation in human layer 2/3 cortical neurons}, Science \textbf{367}, 83 (2020).
\bibitem{Gallego2022} G. Gallego et al., \emph{Event-based vision: a survey}, IEEE Trans. Pattern Anal. Mach. Intell. \textbf{44}, 154 (2022).
\bibitem{Berger1989} R. D. Berger, J. P. Saul, and R. J. Cohen, \emph{Transfer function analysis of autonomic regulation. I. Canine atrial rate response}, Am. J. Physiol. \textbf{256}, H142 (1989).
\bibitem{Walker2010} J. J. Walker, J. R. Terry, and S. L. Lightman, \emph{Origin of ultradian pulsatility in the hypothalamic--pituitary--adrenal axis}, Proc. R. Soc. B \textbf{277}, 1627 (2010).
\bibitem{Belchetz1978} P. E. Belchetz, T. M. Plant, Y. Nakai, E. J. Keogh, and E. Knobil, \emph{Hypophysial responses to continuous and intermittent delivery of hypothalamic gonadotropin-releasing hormone}, Science \textbf{202}, 631 (1978).
\bibitem{Czeisler1999} C. A. Czeisler et al., \emph{Stability, precision, and near-24-hour period of the human circadian pacemaker}, Science \textbf{284}, 2177 (1999).
\bibitem{TakahashiJS2017} J. S. Takahashi, \emph{Transcriptional architecture of the mammalian circadian clock}, Nat. Rev. Genet. \textbf{18}, 164 (2017).
\bibitem{Musk2019} E. Musk and Neuralink, \emph{An integrated brain-machine interface platform with thousands of channels}, J. Med. Internet Res. \textbf{21}, e16194 (2019).
\bibitem{Hochberg2006} L. R. Hochberg et al., \emph{Neuronal ensemble control of prosthetic devices by a human with tetraplegia}, Nature \textbf{442}, 164 (2006).
\bibitem{Willett2021} F. R. Willett et al., \emph{High-performance brain-to-text communication via handwriting}, Nature \textbf{593}, 249 (2021).
\bibitem{Willett2023} F. R. Willett et al., \emph{A high-performance speech neuroprosthesis}, Nature \textbf{620}, 1031 (2023).
\bibitem{Gallego2017} J. A. Gallego, M. G. Perich, L. E. Miller, and S. A. Solla, \emph{Neural manifolds for the control of movement}, Neuron \textbf{94}, 978 (2017).
\bibitem{Fernandez2021} E. Fern\'andez et al., \emph{Visual percepts evoked with an intracortical 96-channel microelectrode array inserted in human occipital cortex}, J. Clin. Invest. \textbf{131}, e151331 (2021).
\bibitem{Tritsch2012} N. X. Tritsch, J. B. Ding, and B. L. Sabatini, \emph{Dopaminergic neurons inhibit striatal output through non-canonical release of GABA}, Nature \textbf{490}, 262 (2012).
\bibitem{Zhang2015} S. Zhang et al., \emph{Dopaminergic and glutamatergic microdomains in a subset of rodent mesoaccumbens axons}, Nat. Neurosci. \textbf{18}, 386 (2015).
\bibitem{Mongillo2008} G. Mongillo, O. Barak, and M. Tsodyks, \emph{Synaptic theory of working memory}, Science \textbf{319}, 1543 (2008).
\bibitem{Stokes2015} M. G. Stokes, \emph{`Activity-silent' working memory in prefrontal cortex: a dynamic coding framework}, Trends Cogn. Sci. \textbf{19}, 394 (2015).
\bibitem{Smith2013} S. L. Smith, I. T. Smith, T. Branco, and M. H\"ausser, \emph{Dendritic spikes enhance stimulus selectivity in cortical neurons in vivo}, Nature \textbf{503}, 115 (2013).
\bibitem{Baden2016} T. Baden et al., \emph{The functional diversity of retinal ganglion cells in the mouse}, Nature \textbf{529}, 345 (2016).
\bibitem{Lottem2018} E. Lottem et al., \emph{Activation of serotonin neurons promotes active persistence in a probabilistic foraging task}, Nat. Commun. \textbf{9}, 1000 (2018).
\bibitem{Payeur2021} A. Payeur, J. Guerguiev, F. Zenke, B. A. Richards, and R. Naud, \emph{Burst-dependent synaptic plasticity can coordinate learning in hierarchical circuits}, Nat. Neurosci. \textbf{24}, 1010 (2021).
\bibitem{MehonicKenyon2022} A. Mehonic and A. J. Kenyon, \emph{Brain-inspired computing needs a master plan}, Nature \textbf{604}, 255 (2022).
\bibitem{Poe2020} G. R. Poe et al., \emph{Locus coeruleus: a new look at the blue spot}, Nat. Rev. Neurosci. \textbf{21}, 644 (2020).
\bibitem{Hangya2015} B. Hangya, S. P. Ranade, M. Lorenc, and A. Kepecs, \emph{Central cholinergic neurons are rapidly recruited by reinforcement feedback}, Cell \textbf{162}, 1155 (2015).
\bibitem{FernandezRuiz2019} A. Fern\'andez-Ruiz et al., \emph{Long-duration hippocampal sharp wave ripples improve memory}, Science \textbf{364}, 1082 (2019).
\bibitem{Herzfeld2018} D. J. Herzfeld, Y. Kojima, R. Soetedjo, and R. Shadmehr, \emph{Encoding of error and learning to correct that error by the Purkinje cells of the cerebellum}, Nat. Neurosci. \textbf{21}, 736 (2018).
\bibitem{Duan2014} B. Duan et al., \emph{Identification of spinal circuits transmitting and gating mechanical pain}, Cell \textbf{159}, 1417 (2014).
\bibitem{Tremblay2016} R. Tremblay, S. Lee, and B. Rudy, \emph{GABAergic interneurons in the neocortex: from cellular properties to circuits}, Neuron \textbf{91}, 260 (2016).
\bibitem{Ren2018} J. Ren et al., \emph{Anatomically defined and functionally distinct dorsal raphe serotonin sub-systems}, Cell \textbf{175}, 472 (2018).
\bibitem{Pfeffer2013} C. K. Pfeffer, M. Xue, M. He, Z. J. Huang, and M. Scanziani, \emph{Inhibition of inhibition in visual cortex: the logic of connections between molecularly distinct interneurons}, Nat. Neurosci. \textbf{16}, 1068 (2013).
\bibitem{Pi2013} H.-J. Pi et al., \emph{Cortical interneurons that specialize in disinhibitory control}, Nature \textbf{503}, 521 (2013).
\bibitem{Keller2018} G. B. Keller and T. D. Mrsic-Flogel, \emph{Predictive processing: a canonical cortical computation}, Neuron \textbf{100}, 424 (2018).
\bibitem{HubelWiesel1962} D. H. Hubel and T. N. Wiesel, \emph{Receptive fields, binocular interaction and functional architecture in the cat's visual cortex}, J. Physiol. \textbf{160}, 106 (1962).
\bibitem{Riesenhuber1999} M. Riesenhuber and T. Poggio, \emph{Hierarchical models of object recognition in cortex}, Nat. Neurosci. \textbf{2}, 1019 (1999).
\bibitem{Fukushima1980} K. Fukushima, \emph{Neocognitron: a self-organizing neural network model for a mechanism of pattern recognition unaffected by shift in position}, Biol. Cybern. \textbf{36}, 193 (1980).
\bibitem{Yamins2014} D. L. K. Yamins et al., \emph{Performance-optimized hierarchical models predict neural responses in higher visual cortex}, Proc. Natl. Acad. Sci. USA \textbf{111}, 8619 (2014).
\bibitem{Hung2005} C. P. Hung, G. Kreiman, T. Poggio, and J. J. DiCarlo, \emph{Fast readout of object identity from macaque inferior temporal cortex}, Science \textbf{310}, 863 (2005).
\bibitem{WiskottSejnowski2002} L. Wiskott and T. J. Sejnowski, \emph{Slow feature analysis: unsupervised learning of invariances}, Neural Comput. \textbf{14}, 715 (2002).
\bibitem{Foldiak1991} P. F\"oldi\'ak, \emph{Learning invariance from transformation sequences}, Neural Comput. \textbf{3}, 194 (1991).
\bibitem{LiDiCarlo2008} N. Li and J. J. DiCarlo, \emph{Unsupervised natural experience rapidly alters invariant object representation in visual cortex}, Science \textbf{321}, 1502 (2008).
\bibitem{RaoBallard1999} R. P. N. Rao and D. H. Ballard, \emph{Predictive coding in the visual cortex: a functional interpretation of some extra-classical receptive-field effects}, Nat. Neurosci. \textbf{2}, 79 (1999).
\bibitem{Kar2019} K. Kar, J. Kubilius, K. Schmidt, E. B. Issa, and J. J. DiCarlo, \emph{Evidence that recurrent circuits are critical to the ventral stream's execution of core object recognition behavior}, Nat. Neurosci. \textbf{22}, 974 (2019).
\bibitem{Szegedy2014} C. Szegedy et al., \emph{Intriguing properties of neural networks}, in \emph{Proc. ICLR} (2014), arXiv:1312.6199.
\bibitem{Weiss2002} Y. Weiss, E. P. Simoncelli, and E. H. Adelson, \emph{Motion illusions as optimal percepts}, Nat. Neurosci. \textbf{5}, 598 (2002).
\bibitem{Purves2011} D. Purves, W. T. Wojtach, and R. B. Lotto, \emph{Understanding vision in wholly empirical terms}, Proc. Natl. Acad. Sci. USA \textbf{108}, Suppl.~3, 15588 (2011).
\bibitem{LaferSousa2015} R. Lafer-Sousa, K. L. Hermann, and B. R. Conway, \emph{Striking individual differences in color perception uncovered by `the dress' photograph}, Curr. Biol. \textbf{25}, R545 (2015).
\bibitem{Wallisch2017} P. Wallisch, \emph{Illumination assumptions account for individual differences in the perceptual interpretation of a profoundly ambiguous stimulus in the color domain: ``The dress''}, J. Vis. \textbf{17}(4), 5 (2017).
\bibitem{Schwartz2009} O. Schwartz, T. J. Sejnowski, and P. Dayan, \emph{Perceptual organization in the tilt illusion}, J. Vis. \textbf{9}(4), 19 (2009).
\bibitem{SchillerCarvey2005} P. H. Schiller and C. E. Carvey, \emph{The Hermann grid illusion revisited}, Perception \textbf{34}, 1375 (2005).
\bibitem{Nijhawan1994} R. Nijhawan, \emph{Motion extrapolation in catching}, Nature \textbf{370}, 256 (1994).
\bibitem{Conway2005} B. R. Conway, A. Kitaoka, A. Yazdanbakhsh, C. C. Pack, and M. S. Livingstone, \emph{Neural basis for a powerful static motion illusion}, J. Neurosci. \textbf{25}, 5651 (2005).
\bibitem{OteroMillan2012} J. Otero-Millan, S. L. Macknik, and S. Martinez-Conde, \emph{Microsaccades and blinks trigger illusory rotation in the ``rotating snakes'' illusion}, J. Neurosci. \textbf{32}, 6043 (2012).
\bibitem{MorenoBote2007} R. Moreno-Bote, J. Rinzel, and N. Rubin, \emph{Noise-induced alternations in an attractor network model of perceptual bistability}, J. Neurophysiol. \textbf{98}, 1125 (2007).
\bibitem{Watanabe2018} E. Watanabe, A. Kitaoka, K. Sakamoto, M. Yasugi, and K. Tanaka, \emph{Illusory motion reproduced by deep neural networks trained for prediction}, Front. Psychol. \textbf{9}, 345 (2018).
\bibitem{GomezVilla2020} A. G\'omez-Villa, A. Mart\'in, J. Vazquez-Corral, M. Bertalm\'io, and J. Malo, \emph{Color illusions also deceive CNNs for low-level vision tasks: analysis and implications}, Vision Res. \textbf{176}, 156 (2020).
\bibitem{delCastillo1954} J. del Castillo and B. Katz, \emph{Quantal components of the end-plate potential}, J. Physiol. \textbf{124}, 560 (1954).
\bibitem{Nowak1984} L. Nowak, P. Bregestovski, P. Ascher, A. Herbet, and A. Prochiantz, \emph{Magnesium gates glutamate-activated channels in mouse central neurones}, Nature \textbf{307}, 462 (1984).
\bibitem{Malinow2002} R. Malinow and R. C. Malenka, \emph{AMPA receptor trafficking and synaptic plasticity}, Annu. Rev. Neurosci. \textbf{25}, 103 (2002).
\bibitem{Matsuzaki2004} M. Matsuzaki, N. Honkura, G. C. R. Ellis-Davies, and H. Kasai, \emph{Structural basis of long-term potentiation in single dendritic spines}, Nature \textbf{429}, 761 (2004).
\bibitem{Rudomin1999} P. Rudomin and R. F. Schmidt, \emph{Presynaptic inhibition in the vertebrate spinal cord revisited}, Exp. Brain Res. \textbf{129}, 1 (1999).
\bibitem{HutcheonYarom2000} B. Hutcheon and Y. Yarom, \emph{Resonance, oscillation and the intrinsic frequency preferences of neurons}, Trends Neurosci. \textbf{23}, 216 (2000).
\bibitem{Seung1996} H. S. Seung, \emph{How the brain keeps the eyes still}, Proc. Natl. Acad. Sci. USA \textbf{93}, 13339 (1996).
\bibitem{Hounsgaard1988} J. Hounsgaard, H. Hultborn, B. Jespersen, and O. Kiehn, \emph{Bistability of alpha-motoneurones in the decerebrate cat and in the acute spinal cat after intravenous 5-hydroxytryptophan}, J. Physiol. \textbf{405}, 345 (1988).
\bibitem{Izhikevich2007} E. M. Izhikevich, \emph{Dynamical Systems in Neuroscience: The Geometry of Excitability and Bursting} (MIT Press, Cambridge, MA, 2007).
\bibitem{AgmonSnir1998} H. Agmon-Snir, C. E. Carr, and J. Rinzel, \emph{The role of dendrites in auditory coincidence detection}, Nature \textbf{393}, 268 (1998).
\bibitem{BorstSoria2012} J. G. G. Borst and J. Soria van Hoeve, \emph{The calyx of Held synapse: from model synapse to auditory relay}, Annu. Rev. Physiol. \textbf{74}, 199 (2012).
\bibitem{Euler2002} T. Euler, P. B. Detwiler, and W. Denk, \emph{Directionally selective calcium signals in dendrites of starburst amacrine cells}, Nature \textbf{418}, 845 (2002).
\bibitem{AstonJonesBloom1981} G. Aston-Jones and F. E. Bloom, \emph{Activity of norepinephrine-containing locus coeruleus neurons in behaving rats anticipates fluctuations in the sleep-waking cycle}, J. Neurosci. \textbf{1}, 876 (1981).
\bibitem{BrooksPeever2012} P. L. Brooks and J. H. Peever, \emph{Identification of the transmitter and receptor mechanisms responsible for REM sleep paralysis}, J. Neurosci. \textbf{32}, 9785 (2012).
\bibitem{Borbely1982} A. A. Borb\'ely, \emph{A two process model of sleep regulation}, Hum. Neurobiol. \textbf{1}, 195 (1982).
\bibitem{PorkkaHeiskanen1997} T. Porkka-Heiskanen et al., \emph{Adenosine: a mediator of the sleep-inducing effects of prolonged wakefulness}, Science \textbf{276}, 1265 (1997).
\bibitem{Steriade1993} M. Steriade, A. N\'u\~nez, and F. Amzica, \emph{A novel slow ($<1$ Hz) oscillation of neocortical neurons in vivo: depolarizing and hyperpolarizing components}, J. Neurosci. \textbf{13}, 3252 (1993).
\bibitem{McCormick1992} D. A. McCormick, \emph{Neurotransmitter actions in the thalamus and cerebral cortex and their role in neuromodulation of thalamocortical activity}, Prog. Neurobiol. \textbf{39}, 337 (1992).
\bibitem{SteriadeMcCormickSejnowski1993} M. Steriade, D. A. McCormick, and T. J. Sejnowski, \emph{Thalamocortical oscillations in the sleeping and aroused brain}, Science \textbf{262}, 679 (1993).
\bibitem{Nadasdy1999} Z. N\'adasdy et al., \emph{Replay and time compression of recurring spike sequences in the hippocampus}, J. Neurosci. \textbf{19}, 9497 (1999).
\bibitem{Maingret2016} N. Maingret et al., \emph{Hippocampo-cortical coupling mediates memory consolidation during sleep}, Nat. Neurosci. \textbf{19}, 959 (2016).
\bibitem{McClelland1995} J. L. McClelland, B. L. McNaughton, and R. C. O'Reilly, \emph{Why there are complementary learning systems in the hippocampus and neocortex}, Psychol. Rev. \textbf{102}, 419 (1995).
\bibitem{TononiCirelli2014} G. Tononi and C. Cirelli, \emph{Sleep and the price of plasticity: from synaptic and cellular homeostasis to memory consolidation and integration}, Neuron \textbf{81}, 12 (2014).
\bibitem{deVivo2017} L. de Vivo et al., \emph{Ultrastructural evidence for synaptic scaling across the wake/sleep cycle}, Science \textbf{355}, 507 (2017).
\bibitem{Xie2013} L. Xie et al., \emph{Sleep drives metabolite clearance from the adult brain}, Science \textbf{342}, 373 (2013).
\bibitem{Miao2024} A. Miao et al., \emph{Brain clearance is reduced during sleep and anesthesia}, Nat. Neurosci. \textbf{27}, 1046 (2024).
\bibitem{Vyazovskiy2011} V. V. Vyazovskiy et al., \emph{Local sleep in awake rats}, Nature \textbf{472}, 443 (2011).
\bibitem{Lancaster2013} M. A. Lancaster et al., \emph{Cerebral organoids model human brain development and microcephaly}, Nature \textbf{501}, 373 (2013).
\bibitem{Jordan2024} F. D. Jordan et al., \emph{Open and remotely accessible Neuroplatform for research in wetware computing}, Front. Artif. Intell. \textbf{7}, 1376042 (2024).
\bibitem{Smirnova2023} L. Smirnova et al., \emph{Organoid intelligence (OI): the new frontier in biocomputing and intelligence-in-a-dish}, Front. Sci. \textbf{1}, 1017235 (2023).
\bibitem{Wagenaar2006} D. A. Wagenaar, J. Pine, and S. M. Potter, \emph{An extremely rich repertoire of bursting patterns during the development of cortical cultures}, BMC Neurosci. \textbf{7}, 11 (2006).
\bibitem{Shahaf2001} G. Shahaf and S. Marom, \emph{Learning in networks of cortical neurons}, J. Neurosci. \textbf{21}, 8782 (2001).
\bibitem{Kagan2022} B. J. Kagan et al., \emph{In vitro neurons learn and exhibit sentience when embodied in a simulated game-world}, Neuron \textbf{110}, 3952 (2022).
\bibitem{Friston2010} K. Friston, \emph{The free-energy principle: a unified brain theory?}, Nat. Rev. Neurosci. \textbf{11}, 127 (2010).
\bibitem{Balci2023} F. Balci et al., \emph{A response to claims of emergent intelligence and sentience in a dish}, Neuron \textbf{111}, 604 (2023).
\bibitem{Bakkum2008} D. J. Bakkum, Z. C. Chao, and S. M. Potter, \emph{Spatio-temporal electrical stimuli shape behavior of an embodied cortical network in a goal-directed learning task}, J. Neural Eng. \textbf{5}, 310 (2008).
\bibitem{Maass2002} W. Maass, T. Natschl\"ager, and H. Markram, \emph{Real-time computing without stable states: a new framework for neural computation based on perturbations}, Neural Comput. \textbf{14}, 2531 (2002).
\bibitem{JaegerHaas2004} H. Jaeger and H. Haas, \emph{Harnessing nonlinearity: predicting chaotic systems and saving energy in wireless communication}, Science \textbf{304}, 78 (2004).
\bibitem{Cai2023} H. Cai et al., \emph{Brain organoid reservoir computing for artificial intelligence}, Nat. Electron. \textbf{6}, 1032 (2023).
\bibitem{Habibollahi2023} F. Habibollahi, B. J. Kagan, A. N. Burkitt, and C. French, \emph{Critical dynamics arise during structured information presentation within embodied in vitro neuronal networks}, Nat. Commun. \textbf{14}, 5287 (2023).
\bibitem{Robbins2026} A. Robbins et al., \emph{Goal-directed learning in cortical organoids}, Cell Rep. \textbf{45}, 116984 (2026).
\bibitem{Keene2020} S. T. Keene et al., \emph{A biohybrid synapse with neurotransmitter-mediated plasticity}, Nat. Mater. \textbf{19}, 969 (2020).
\bibitem{Jo2016} J. Jo et al., \emph{Midbrain-like organoids from human pluripotent stem cells contain functional dopaminergic and neuromelanin-producing neurons}, Cell Stem Cell \textbf{19}, 248 (2016).
\bibitem{Hendry1987} S. H. Hendry, H. D. Schwark, E. G. Jones, and J. Yan, \emph{Numbers and proportions of GABA-immunoreactive neurons in different areas of monkey cerebral cortex}, J. Neurosci. \textbf{7}, 1503 (1987).
\bibitem{Tang2001synapses} Y. Tang, J. R. Nyengaard, D. M. G. De Groot, and H. J. G. Gundersen, \emph{Total regional and global number of synapses in the human brain neocortex}, Synapse \textbf{41}, 258 (2001).
\bibitem{Pakkenberg1997} B. Pakkenberg and H. J. G. Gundersen, \emph{Neocortical neuron number in humans: effect of sex and age}, J. Comp. Neurol. \textbf{384}, 312 (1997).
\bibitem{Matsuda2009} W. Matsuda, T. Furuta, K. C. Nakamura, H. Hioki, F. Fujiyama, R. Arai, and T. Kaneko, \emph{Single nigrostriatal dopaminergic neurons form widely spread and highly dense axonal arborizations in the neostriatum}, J. Neurosci. \textbf{29}, 444 (2009).
\bibitem{Pakkenberg1991} B. Pakkenberg, A. M{\o}ller, H. J. G. Gundersen, A. Mouritzen Dam, and H. Pakkenberg, \emph{The absolute number of nerve cells in substantia nigra in normal subjects and in patients with Parkinson's disease estimated with an unbiased stereological method}, J. Neurol. Neurosurg. Psychiatry \textbf{54}, 30 (1991).
\bibitem{Baker1990} K. G. Baker, G. M. Halliday, and I. T\"ork, \emph{Cytoarchitecture of the human dorsal raphe nucleus}, J. Comp. Neurol. \textbf{301}, 147 (1990).
\bibitem{Baker1991} K. G. Baker, G. M. Halliday, P. Halasz, J.-P. Hornung, L. B. Geffen, R. G. H. Cotton, and I. T\"ork, \emph{Cytoarchitecture of serotonin-synthesizing neurons in the pontine tegmentum of the human brain}, Synapse \textbf{7}, 301 (1991).
\bibitem{Airaksinen1991} M. S. Airaksinen, A. Paetau, L. Palj\"arvi, K. Reinikainen, P. Riekkinen, R. Suomalainen, and P. Panula, \emph{Histamine neurons in human hypothalamus: anatomy in normal and Alzheimer diseased brains}, Neuroscience \textbf{44}, 465 (1991).
\bibitem{Colquhoun1992} D. Colquhoun, P. Jonas, and B. Sakmann, \emph{Action of brief pulses of glutamate on AMPA/kainate receptors in patches from different neurones of rat hippocampal slices}, J. Physiol. \textbf{458}, 261 (1992).
\bibitem{DutarNicoll1988} P. Dutar and R. A. Nicoll, \emph{A physiological role for GABA$_B$ receptors in the central nervous system}, Nature \textbf{332}, 156 (1988).

\bibitem{Rauch2003} A. Rauch, G. La Camera, H.-R. L\"uscher, W. Senn, and S. Fusi, \emph{Neocortical pyramidal cells respond as integrate-and-fire neurons to in vivo-like input currents}, J. Neurophysiol. \textbf{90}, 1598 (2003).
\bibitem{KlumppEady1956} R. G. Klumpp and H. R. Eady, \emph{Some measurements of interaural time difference thresholds}, J. Acoust. Soc. Am. \textbf{28}, 859 (1956).
\bibitem{Brand2002} A. Brand, O. Behrend, T. Marquardt, D. McAlpine, and B. Grothe, \emph{Precise inhibition is essential for microsecond interaural time difference coding}, Nature \textbf{417}, 543 (2002).
\bibitem{Funahashi1989} S. Funahashi, C. J. Bruce, and P. S. Goldman-Rakic, \emph{Mnemonic coding of visual space in the monkey's dorsolateral prefrontal cortex}, J. Neurophysiol. \textbf{61}, 331 (1989).
\bibitem{WangMarkram2006} Y. Wang, H. Markram, P. H. Goodman, T. K. Berger, J. Ma, and P. S. Goldman-Rakic, \emph{Heterogeneity in the pyramidal network of the medial prefrontal cortex}, Nat. Neurosci. \textbf{9}, 534 (2006).

\bibitem{McCullochPitts1943} W. S. McCulloch and W. Pitts, \emph{A logical calculus of the ideas immanent in nervous activity}, Bull. Math. Biophys. \textbf{5}, 115 (1943).
\bibitem{Fried2002} S. I. Fried, T. A. M\"unch, and F. S. Werblin, \emph{Mechanisms and circuitry underlying directional selectivity in the retina}, Nature \textbf{420}, 411 (2002).

\bibitem{Hromadka2008} T. Hrom\'adka, M. R. DeWeese, and A. M. Zador, \emph{Sparse representation of sounds in the unanesthetized auditory cortex}, PLoS Biol. \textbf{6}, e16 (2008).
\bibitem{Abbott1997depression} L. F. Abbott, J. A. Varela, K. Sen, and S. B. Nelson, \emph{Synaptic depression and cortical gain control}, Science \textbf{275}, 221 (1997).
\bibitem{ShadlenNewsome1998} M. N. Shadlen and W. T. Newsome, \emph{The variable discharge of cortical neurons: implications for connectivity, computation, and information coding}, J. Neurosci. \textbf{18}, 3870 (1998).

\bibitem{TrulsonJacobs1979} M. E. Trulson and B. L. Jacobs, \emph{Raphe unit activity in freely moving cats: correlation with level of behavioral arousal}, Brain Res. \textbf{163}, 135 (1979).
\bibitem{GagnonParent2014} D. Gagnon and M. Parent, \emph{Distribution of VGLUT3 in highly collateralized axons from the rat dorsal raphe nucleus as revealed by single-neuron reconstructions}, PLoS One \textbf{9}, e87709 (2014).
\bibitem{VandermaelenAghajanian1983} C. P. Vandermaelen and G. K. Aghajanian, \emph{Electrophysiological and pharmacological characterization of serotonergic dorsal raphe neurons recorded extracellularly and intracellularly in rat brain slices}, Brain Res. \textbf{289}, 109 (1983).
\bibitem{CourtneyFord2016} N. A. Courtney and C. P. Ford, \emph{Mechanisms of 5-HT$_{1A}$ receptor-mediated transmission in dorsal raphe serotonin neurons}, J. Physiol. \textbf{594}, 953 (2016).
\bibitem{Hashemi2012serotonin} P. Hashemi et al., \emph{Brain dopamine and serotonin differ in regulation and its consequences}, Proc. Natl. Acad. Sci. USA \textbf{109}, 11510 (2012).
\bibitem{Black1934feedback} H. S. Black, \emph{Stabilized feedback amplifiers}, Bell Syst. Tech. J. \textbf{13}, 1 (1934).
\bibitem{KobayashiSchultz2008} S. Kobayashi and W. Schultz, \emph{Influence of reward delays on responses of dopamine neurons}, J. Neurosci. \textbf{28}, 7837 (2008).
\bibitem{Schweighofer2008} N. Schweighofer et al., \emph{Low-serotonin levels increase delayed reward discounting in humans}, J. Neurosci. \textbf{28}, 4528 (2008).

\bibitem{Malenka1988calcium} R. C. Malenka, J. A. Kauer, R. S. Zucker, and R. A. Nicoll, \emph{Postsynaptic calcium is sufficient for potentiation of hippocampal synaptic transmission}, Science \textbf{242}, 81 (1988).
\bibitem{Bliss1973LTP} T. V. P. Bliss and T. L{\o}mo, \emph{Long-lasting potentiation of synaptic transmission in the dentate area of the anaesthetized rabbit following stimulation of the perforant path}, J. Physiol. \textbf{232}, 331 (1973).
\bibitem{Kelso1986hebb} S. R. Kelso, A. H. Ganong, and T. H. Brown, \emph{Hebbian synapses in hippocampus}, Proc. Natl. Acad. Sci. USA \textbf{83}, 5326 (1986).
\bibitem{He2015eligibility} K. He et al., \emph{Distinct eligibility traces for LTP and LTD in cortical synapses}, Neuron \textbf{88}, 528 (2015).
\bibitem{Olveczky2005} B. P. \"Olveczky, A. S. Andalman, and M. S. Fee, \emph{Vocal experimentation in the juvenile songbird requires a basal ganglia circuit}, PLoS Biol. \textbf{3}, e153 (2005).
\bibitem{TumerBrainard2007} E. C. Tumer and M. S. Brainard, \emph{Performance variability enables adaptive plasticity of `crystallized' adult birdsong}, Nature \textbf{450}, 1240 (2007).
\bibitem{Eshel2015arithmetic} N. Eshel et al., \emph{Arithmetic and local circuitry underlying dopamine prediction errors}, Nature \textbf{525}, 243 (2015).
\bibitem{GraceOnn1989} A. A. Grace and S.-P. Onn, \emph{Morphology and electrophysiological properties of immunocytochemically identified rat dopamine neurons recorded in vitro}, J. Neurosci. \textbf{9}, 3463 (1989).
\bibitem{ODoherty2004actor} J. O'Doherty et al., \emph{Dissociable roles of ventral and dorsal striatum in instrumental conditioning}, Science \textbf{304}, 452 (2004).

\bibitem{NeweyPowell1987} W. K. Newey and J. L. Powell, \emph{Asymmetric least squares estimation and testing}, Econometrica \textbf{55}, 819 (1987).
\bibitem{Beierholm2013vigor} U. Beierholm et al., \emph{Dopamine modulates reward-related vigor}, Neuropsychopharmacology \textbf{38}, 1495 (2013).
\bibitem{Howe2013ramp} M. W. Howe, P. L. Tierney, S. G. Sandberg, P. E. M. Phillips, and A. M. Graybiel, \emph{Prolonged dopamine signalling in striatum signals proximity and value of distant rewards}, Nature \textbf{500}, 575 (2013).

\bibitem{Baker1989LC} K. G. Baker, I. T\"ork, J.-P. Hornung, and P. Halasz, \emph{The human locus coeruleus complex: an immunohistochemical and three dimensional reconstruction study}, Exp. Brain Res. \textbf{77}, 257 (1989).
\bibitem{Schwarz2015LC} L. A. Schwarz et al., \emph{Viral-genetic tracing of the input--output organization of a central noradrenaline circuit}, Nature \textbf{524}, 88 (2015).
\bibitem{Uematsu2017LC} A. Uematsu et al., \emph{Modular organization of the brainstem noradrenaline system coordinates opposing learning states}, Nat. Neurosci. \textbf{20}, 1602 (2017).
\bibitem{AstonJones1994vigilance} G. Aston-Jones, J. Rajkowski, P. Kubiak, and T. Alexinsky, \emph{Locus coeruleus neurons in monkey are selectively activated by attended cues in a vigilance task}, J. Neurosci. \textbf{14}, 4467 (1994).
\bibitem{Clayton2004LC} E. C. Clayton, J. Rajkowski, J. D. Cohen, and G. Aston-Jones, \emph{Phasic activation of monkey locus ceruleus neurons by simple decisions in a forced-choice task}, J. Neurosci. \textbf{24}, 9914 (2004).
\bibitem{Foote1975NE} S. L. Foote, R. Freedman, and A. P. Oliver, \emph{Effects of putative neurotransmitters on neuronal activity in monkey auditory cortex}, Brain Res. \textbf{86}, 229 (1975).
\bibitem{JepmaNieuwenhuis2011} M. Jepma and S. Nieuwenhuis, \emph{Pupil diameter predicts changes in the exploration--exploitation trade-off: evidence for the adaptive gain theory}, J. Cogn. Neurosci. \textbf{23}, 1587 (2011).
\bibitem{Tervo2014stochastic} D. G. R. Tervo et al., \emph{Behavioral variability through stochastic choice and its gating by anterior cingulate cortex}, Cell \textbf{159}, 21 (2014).
\bibitem{Usher1999LC} M. Usher, J. D. Cohen, D. Servan-Schreiber, J. Rajkowski, and G. Aston-Jones, \emph{The role of locus coeruleus in the regulation of cognitive performance}, Science \textbf{283}, 549 (1999).

\bibitem{Mesulam1983} M.-M. Mesulam, E. J. Mufson, B. H. Wainer, and A. I. Levey, \emph{Central cholinergic pathways in the rat: an overview based on an alternative nomenclature (Ch1--Ch6)}, Neuroscience \textbf{10}, 1185 (1983).
\bibitem{Marrosu1995} F. Marrosu et al., \emph{Microdialysis measurement of cortical and hippocampal acetylcholine release during sleep-wake cycle in freely moving cats}, Brain Res. \textbf{671}, 329 (1995).
\bibitem{Picciotto2012} M. R. Picciotto, M. J. Higley, and Y. S. Mineur, \emph{Acetylcholine as a neuromodulator: cholinergic signaling shapes nervous system function and behavior}, Neuron \textbf{76}, 116 (2012).
\bibitem{HasselmoSchnell1994} M. E. Hasselmo and E. Schnell, \emph{Laminar selectivity of the cholinergic suppression of synaptic transmission in rat hippocampal region CA1: computational modeling and brain slice physiology}, J. Neurosci. \textbf{14}, 3898 (1994).
\bibitem{Gil1997} Z. Gil, B. W. Connors, and Y. Amitai, \emph{Differential regulation of neocortical synapses by neuromodulators and activity}, Neuron \textbf{19}, 679 (1997).
\bibitem{KilgardMerzenich1998} M. P. Kilgard and M. M. Merzenich, \emph{Cortical map reorganization enabled by nucleus basalis activity}, Science \textbf{279}, 1714 (1998).
\bibitem{Atri2004} A. Atri et al., \emph{Blockade of central cholinergic receptors impairs new learning and increases proactive interference in a word paired-associate memory task}, Behav. Neurosci. \textbf{118}, 223 (2004).
\bibitem{Girardeau2009} G. Girardeau, K. Benchenane, S. I. Wiener, G. Buzs\'aki, and M. B. Zugaro, \emph{Selective suppression of hippocampal ripples impairs spatial memory}, Nat. Neurosci. \textbf{12}, 1222 (2009).

\input{supplement/bib_10}
\bibitem{CopenhagenJahr1989} D. R. Copenhagen and C. E. Jahr, \emph{Release of endogenous excitatory amino acids from turtle photoreceptors}, Nature \textbf{341}, 536 (1989).
\bibitem{GlowatzkiFuchs2002} E. Glowatzki and P. A. Fuchs, \emph{Transmitter release at the hair cell ribbon synapse}, Nat. Neurosci. \textbf{5}, 147 (2002).
\bibitem{Tinsley2016} J. N. Tinsley et al., \emph{Direct detection of a single photon by humans}, Nat. Commun. \textbf{7}, 12172 (2016).
\bibitem{Sellick1982} P. M. Sellick, R. Patuzzi, and B. M. Johnstone, \emph{Measurement of basilar membrane motion in the guinea pig using the M\"ossbauer technique}, J. Acoust. Soc. Am. \textbf{72}, 131 (1982).
\bibitem{Ruggero1997} M. A. Ruggero, N. C. Rich, A. Recio, S. S. Narayan, and L. Robles, \emph{Basilar-membrane responses to tones at the base of the chinchilla cochlea}, J. Acoust. Soc. Am. \textbf{101}, 2151 (1997).
\bibitem{NormannPerlman1979} R. A. Normann and I. Perlman, \emph{The effects of background illumination on the photoresponses of red and green cones}, J. Physiol. \textbf{286}, 491 (1979).
\bibitem{Laughlin1981} S. Laughlin, \emph{A simple coding procedure enhances a neuron's information capacity}, Z. Naturforsch. C \textbf{36}, 910 (1981).
\bibitem{CurcioAllen1990} C. A. Curcio and K. A. Allen, \emph{Topography of ganglion cells in human retina}, J. Comp. Neurol. \textbf{300}, 5 (1990).
\bibitem{Greenwood1990} D. D. Greenwood, \emph{A cochlear frequency-position function for several species --- 29 years later}, J. Acoust. Soc. Am. \textbf{87}, 2592 (1990).
\bibitem{Eguiluz2000} V. M. Egu\'iluz, M. Ospeck, Y. Choe, A. J. Hudspeth, and M. O. Magnasco, \emph{Essential nonlinearities in hearing}, Phys. Rev. Lett. \textbf{84}, 5232 (2000).
\bibitem{PalmerRussell1986} A. R. Palmer and I. J. Russell, \emph{Phase-locking in the cochlear nerve of the guinea-pig and its relation to the receptor potential of inner hair-cells}, Hear. Res. \textbf{24}, 1 (1986).
\bibitem{Malnic1999} B. Malnic, J. Hirono, T. Sato, and L. B. Buck, \emph{Combinatorial receptor codes for odors}, Cell \textbf{96}, 713 (1999).
\bibitem{Finger2005} T. E. Finger et al., \emph{ATP signaling is crucial for communication from taste buds to gustatory nerves}, Science \textbf{310}, 1495 (2005).
\bibitem{Huang2005} Y.-J. Huang et al., \emph{Mouse taste buds use serotonin as a neurotransmitter}, J. Neurosci. \textbf{25}, 843 (2005).
\bibitem{JohanssonVallbo1979} R. S. Johansson and A. B. Vallbo, \emph{Tactile sensibility in the human hand: relative and absolute densities of four types of mechanoreceptive units in glabrous skin}, J. Physiol. \textbf{286}, 283 (1979).
\bibitem{McKemy2002} D. D. McKemy, W. M. Neuhausser, and D. Julius, \emph{Identification of a cold receptor reveals a general role for TRP channels in thermosensation}, Nature \textbf{416}, 52 (2002).

\bibitem{Schmolesky1998} M. T. Schmolesky et al., \emph{Signal timing across the macaque visual system}, J. Neurophysiol. \textbf{79}, 3272 (1998).
\bibitem{Lampl2004} I. Lampl, D. Ferster, T. Poggio, and M. Riesenhuber, \emph{Intracellular measurements of spatial integration and the MAX operation in complex cells of the cat primary visual cortex}, J. Neurophysiol. \textbf{92}, 2704 (2004).
\bibitem{KobatakeTanaka1994} E. Kobatake and K. Tanaka, \emph{Neuronal selectivities to complex object features in the ventral visual pathway of the macaque cerebral cortex}, J. Neurophysiol. \textbf{71}, 856 (1994).
\bibitem{Albright1984} T. D. Albright, \emph{Direction and orientation selectivity of neurons in visual area MT of the macaque}, J. Neurophysiol. \textbf{52}, 1106 (1984).
\bibitem{Goodfellow2015} I. J. Goodfellow, J. Shlens, and C. Szegedy, \emph{Explaining and harnessing adversarial examples}, in \emph{Proc. ICLR} (2015), arXiv:1412.6572.
\bibitem{Elsayed2018} G. F. Elsayed et al., \emph{Adversarial examples that fool both computer vision and time-limited humans}, in \emph{Advances in Neural Information Processing Systems 31} (2018), arXiv:1802.08195.
\bibitem{Thompson1982} P. Thompson, \emph{Perceived rate of movement depends on contrast}, Vision Res. \textbf{22}, 377 (1982).
\bibitem{BarlowHill1963} H. B. Barlow and R. M. Hill, \emph{Evidence for a physiological explanation of the waterfall phenomenon and figural after-effects}, Nature \textbf{200}, 1345 (1963).
\bibitem{EaglemanSejnowski2000} D. M. Eagleman and T. J. Sejnowski, \emph{Motion integration and postdiction in visual awareness}, Science \textbf{287}, 2036 (2000).

\bibitem{Feinstein1955mu} B. Feinstein, B. Lindeg{\aa}rd, E. Nyman, and G. Wohlfart, \emph{Morphologic studies of motor units in normal human muscles}, Acta Anat. \textbf{23}, 127 (1955).
\bibitem{MilnerBrown1973rec} H. S. Milner-Brown, R. B. Stein, and R. Yemm, \emph{The orderly recruitment of human motor units during voluntary isometric contractions}, J. Physiol. \textbf{230}, 359 (1973).
\bibitem{Jones2002sdn} K. E. Jones, A. F. de C. Hamilton, and D. M. Wolpert, \emph{Sources of signal-dependent noise during isometric force production}, J. Neurophysiol. \textbf{88}, 1533 (2002).
\bibitem{GomiOsu1998stiff} H. Gomi and R. Osu, \emph{Task-dependent viscoelasticity of human multijoint arm and its spatial characteristics for interaction with environments}, J. Neurosci. \textbf{18}, 8965 (1998).
\bibitem{Jankowska1972Ia} E. Jankowska and W. J. Roberts, \emph{Synaptic actions of single interneurones mediating reciprocal Ia inhibition of motoneurones}, J. Physiol. \textbf{222}, 623 (1972).
\bibitem{Pruszynski2008rapid} J. A. Pruszynski, I. Kurtzer, and S. H. Scott, \emph{Rapid motor responses are appropriately tuned to the metrics of a visuospatial task}, J. Neurophysiol. \textbf{100}, 224 (2008).
\bibitem{SaundersKnill2003vis} J. A. Saunders and D. C. Knill, \emph{Humans use continuous visual feedback from the hand to control fast reaching movements}, Exp. Brain Res. \textbf{152}, 341 (2003).
\bibitem{FlanaganWing1997grip} J. R. Flanagan and A. M. Wing, \emph{The role of internal models in motion planning and control: evidence from grip force adjustments during movements of hand-held loads}, J. Neurosci. \textbf{17}, 1519 (1997).
\bibitem{Matsuoka1985osc} K. Matsuoka, \emph{Sustained oscillations generated by mutually inhibiting neurons with adaptation}, Biol. Cybern. \textbf{52}, 367 (1985).

\bibitem{Treede1995heat} R.-D. Treede, R. A. Meyer, S. N. Raja, and J. N. Campbell, \emph{Evidence for two different heat transduction mechanisms in nociceptive primary afferents innervating monkey skin}, J. Physiol. \textbf{483}, 747 (1995).
\bibitem{Weidner1999cfib} C. Weidner et al., \emph{Functional attributes discriminating mechano-insensitive and mechano-responsive C nociceptors in human skin}, J. Neurosci. \textbf{19}, 10184 (1999).
\bibitem{CampbellLaMotte1983first} J. N. Campbell and R. H. LaMotte, \emph{Latency to detection of first pain}, Brain Res. \textbf{266}, 203 (1983).
\bibitem{LaMotte1982hyper} R. H. LaMotte, J. G. Thalhammer, H. E. Torebj{\"o}rk, and C. J. Robinson, \emph{Peripheral neural mechanisms of cutaneous hyperalgesia following mild injury by heat}, J. Neurosci. \textbf{2}, 765 (1982).
\bibitem{Mancini2014touch} F. Mancini, T. Nash, G. D. Iannetti, and P. Haggard, \emph{Pain relief by touch: a quantitative approach}, Pain \textbf{155}, 635 (2014).
\bibitem{Fields1983rvm} H. L. Fields, J. Bry, I. Hentall, and G. Zorman, \emph{The activity of neurons in the rostral medulla of the rat during withdrawal from noxious heat}, J. Neurosci. \textbf{3}, 2545 (1983).
\bibitem{Crombez1996disrupt} G. Crombez, C. Eccleston, F. Baeyens, and P. Eelen, \emph{The disruptive nature of pain: an experimental investigation}, Behav. Res. Ther. \textbf{34}, 911 (1996).
\bibitem{Schouenborg1995wr} J. Schouenborg, H.-R. Weng, J. Kalliom{\"a}ki, and H. Holmberg, \emph{A survey of spinal dorsal horn neurones encoding the spatial organization of withdrawal reflexes in the rat}, Exp. Brain Res. \textbf{106}, 19 (1995).

\bibitem{Andersen1992cereb} B. B. Andersen, L. Korbo, and B. Pakkenberg, \emph{A quantitative study of the human cerebellum with unbiased stereological techniques}, J. Comp. Neurol. \textbf{326}, 549 (1992).
\bibitem{ShadmehrMussaIvaldi1994} R. Shadmehr and F. A. Mussa-Ivaldi, \emph{Adaptive representation of dynamics during learning of a motor task}, J. Neurosci. \textbf{14}, 3208 (1994).

\bibitem{JoseCollison1970ihr} A. D. Jose and D. Collison, \emph{The normal range and determinants of the intrinsic heart rate in man}, Cardiovasc. Res. \textbf{4}, 160 (1970).
\bibitem{Logothetis1987gbg} D. E. Logothetis, Y. Kurachi, J. Galper, E. J. Neer, and D. E. Clapham, \emph{The $\beta\gamma$ subunits of GTP-binding proteins activate the muscarinic K$^+$ channel in heart}, Nature \textbf{325}, 321 (1987).
\bibitem{KellerWood1984fb} M. E. Keller-Wood and M. F. Dallman, \emph{Corticosteroid inhibition of ACTH secretion}, Endocr. Rev. \textbf{5}, 1 (1984).
\bibitem{Walker2012glc} J. J. Walker et al., \emph{The origin of glucocorticoid hormone oscillations}, PLoS Biol. \textbf{10}, e1001341 (2012).
\bibitem{Wildt1981gnrh} L. Wildt et al., \emph{Frequency and amplitude of gonadotropin-releasing hormone stimulation and gonadotropin secretion in the rhesus monkey}, Endocrinology \textbf{109}, 376 (1981).
\bibitem{Welsh2010scn} D. K. Welsh, J. S. Takahashi, and S. A. Kay, \emph{Suprachiasmatic nucleus: cell autonomy and network properties}, Annu. Rev. Physiol. \textbf{72}, 551 (2010).
\bibitem{Khalsa2003prc} S. B. S. Khalsa, M. E. Jewett, C. Cajochen, and C. A. Czeisler, \emph{A phase response curve to single bright light pulses in human subjects}, J. Physiol. \textbf{549}, 945 (2003).
\bibitem{Sack1992blind} R. L. Sack, A. J. Lewy, M. L. Blood, L. D. Keith, and H. Nakagawa, \emph{Circadian rhythm abnormalities in totally blind people: incidence and clinical significance}, J. Clin. Endocrinol. Metab. \textbf{75}, 127 (1992).

\bibitem{Henze2000extra} D. A. Henze et al., \emph{Intracellular features predicted by extracellular recordings in the hippocampus in vivo}, J. Neurophysiol. \textbf{84}, 390 (2000).
\bibitem{Histed2009stim} M. H. Histed, V. Bonin, and R. C. Reid, \emph{Direct activation of sparse, distributed populations of cortical neurons by electrical microstimulation}, Neuron \textbf{63}, 508 (2009).
\bibitem{Orsborn2014coad} A. L. Orsborn et al., \emph{Closed-loop decoder adaptation shapes neural plasticity for skillful neuroprosthetic control}, Neuron \textbf{82}, 1380 (2014).

\bibitem{HuVervaekeStorm2002} H. Hu, K. Vervaeke, and J. F. Storm, \emph{Two forms of electrical resonance at theta frequencies, generated by M-current, h-current and persistent Na$^+$ current in rat hippocampal pyramidal cells}, J. Physiol. \textbf{545}, 783 (2002).
\bibitem{Mauro1970} A. Mauro, F. Conti, F. Dodge, and R. Schor, \emph{Subthreshold behavior and phenomenological impedance of the squid giant axon}, J. Gen. Physiol. \textbf{55}, 497 (1970).
\bibitem{Aksay2001} E. Aksay et al., \emph{In vivo intracellular recording and perturbation of persistent activity in a neural integrator}, Nat. Neurosci. \textbf{4}, 184 (2001).
\bibitem{Major2004} G. Major et al., \emph{Plasticity and tuning of the time course of analog persistent firing in a neural integrator}, Proc. Natl. Acad. Sci. USA \textbf{101}, 7745 (2004).
\bibitem{Tateno2004} T. Tateno, A. Harsch, and H. P. C. Robinson, \emph{Threshold firing frequency-current relationships of neurons in rat somatosensory cortex: type 1 and type 2 dynamics}, J. Neurophysiol. \textbf{92}, 2283 (2004).
\bibitem{Marder2012} E. Marder, \emph{Neuromodulation of neuronal circuits: back to the future}, Neuron \textbf{76}, 1 (2012).
\bibitem{McGintyHarper1976} D. J. McGinty and R. M. Harper, \emph{Dorsal raphe neurons: depression of firing during sleep in cats}, Brain Res. \textbf{101}, 569 (1976).

\bibitem{Gentet2000} L. J. Gentet, G. J. Stuart, and J. D. Clements, \emph{Direct measurement of specific membrane capacitance in neurons}, Biophys. J. \textbf{79}, 314 (2000).
\bibitem{Borst1995} J. G. G. Borst, F. Helmchen, and B. Sakmann, \emph{Pre- and postsynaptic whole-cell recordings in the medial nucleus of the trapezoid body of the rat}, J. Physiol. \textbf{489}, 825 (1995).
\bibitem{AmirDevor2003} R. Amir and M. Devor, \emph{Electrical excitability of the soma of sensory neurons is required for spike invasion of the soma, but not for through-conduction}, Biophys. J. \textbf{84}, 2181 (2003).
\bibitem{Chadderton2004} P. Chadderton, T. W. Margrie, and M. H\"ausser, \emph{Integration of quanta in cerebellar granule cells during sensory processing}, Nature \textbf{428}, 856 (2004).
\bibitem{LitwinKumar2017} A. Litwin-Kumar et al., \emph{Optimal degrees of synaptic connectivity}, Neuron \textbf{93}, 1153 (2017).
\bibitem{NapperHarvey1988} R. M. A. Napper and R. J. Harvey, \emph{Number of parallel fiber synapses on an individual Purkinje cell in the cerebellum of the rat}, J. Comp. Neurol. \textbf{274}, 168 (1988).
\bibitem{Megias2001} M. Meg\'{\i}as, Z. Emri, T. F. Freund, and A. I. Guly\'as, \emph{Total number and distribution of inhibitory and excitatory synapses on hippocampal CA1 pyramidal cells}, Neuroscience \textbf{102}, 527 (2001).
\bibitem{Polsky2004} A. Polsky, B. W. Mel, and J. Schiller, \emph{Computational subunits in thin dendrites of pyramidal cells}, Nat. Neurosci. \textbf{7}, 621 (2004).
\bibitem{Brunel2004} N. Brunel et al., \emph{Optimal information storage and the distribution of synaptic weights: perceptron versus Purkinje cell}, Neuron \textbf{43}, 745 (2004).

\bibitem{Vyazovskiy2009} V. V. Vyazovskiy et al., \emph{Cortical firing and sleep homeostasis}, Neuron \textbf{63}, 865 (2009).
\bibitem{Daan1984} S. Daan, D. G. Beersma, and A. A. Borb\'ely, \emph{Timing of human sleep: recovery process gated by a circadian pacemaker}, Am. J. Physiol. \textbf{246}, R161 (1984).
\bibitem{Borbely1981} A. A. Borb\'ely et al., \emph{Sleep deprivation: effect on sleep stages and EEG power density in man}, Electroencephalogr. Clin. Neurophysiol. \textbf{51}, 483 (1981).
\bibitem{SteriadeAmzica1993} M. Steriade, F. Amzica, and A. Nu\~nez, \emph{Cholinergic and noradrenergic modulation of the slow ($\approx0.3$ Hz) oscillation in neocortical cells}, J. Neurophysiol. \textbf{70}, 1385 (1993).
\bibitem{vonKrosigk1993} M. von Krosigk, T. Bal, and D. A. McCormick, \emph{Cellular mechanisms of a synchronized oscillation in the thalamus}, Science \textbf{261}, 361 (1993).
\bibitem{LeeWilson2002} A. K. Lee and M. A. Wilson, \emph{Memory of sequential experience in the hippocampus during slow wave sleep}, Neuron \textbf{36}, 1183 (2002).
\bibitem{Euston2007} D. R. Euston, M. Tatsuno, and B. L. McNaughton, \emph{Fast-forward playback of recent memory sequences in prefrontal cortex during sleep}, Science \textbf{318}, 1147 (2007).
\bibitem{Hobson2009} J. A. Hobson, \emph{REM sleep and dreaming: towards a theory of protoconsciousness}, Nat. Rev. Neurosci. \textbf{10}, 803 (2009).

\bibitem{CohenSegal2011} D. Cohen and M. Segal, \emph{Network bursts in hippocampal microcultures are terminated by exhaustion of vesicle pools}, J. Neurophysiol. \textbf{106}, 2314 (2011).
\bibitem{Tsodyks2000} M. Tsodyks, A. Uziel, and H. Markram, \emph{Synchrony generation in recurrent networks with frequency-dependent synapses}, J. Neurosci. \textbf{20}, RC50 (2000).
\bibitem{Benson1994} D. L. Benson, F. H. Watkins, O. Steward, and G. Banker, \emph{Characterization of GABAergic neurons in hippocampal cell cultures}, J. Neurocytol. \textbf{23}, 279 (1994).

\input{supplement/bib_22}
\bibitem{White1986} J. G. White, E. Southgate, J. N. Thomson, and S. Brenner, \emph{The structure of the nervous system of the nematode Caenorhabditis elegans}, Phil. Trans. R. Soc. Lond. B \textbf{314}, 1 (1986).
\bibitem{Cook2019} S. J. Cook et al., \emph{Whole-animal connectomes of both Caenorhabditis elegans sexes}, Nature \textbf{571}, 63 (2019).
\bibitem{Dorkenwald2024} S. Dorkenwald et al., \emph{Neuronal wiring diagram of an adult brain}, Nature \textbf{634}, 124 (2024).
\bibitem{Goodman1998} M. B. Goodman, D. H. Hall, L. Avery, and S. R. Lockery, \emph{Active currents regulate sensitivity and dynamic range in C.~elegans neurons}, Neuron \textbf{20}, 763 (1998).
\bibitem{Sulston1975} J. Sulston, M. Dew, and S. Brenner, \emph{Dopaminergic neurons in the nematode Caenorhabditis elegans}, J. Comp. Neurol. \textbf{163}, 215 (1975).
\bibitem{Sawin2000} E. R. Sawin, R. Ranganathan, and H. R. Horvitz, \emph{C.~elegans locomotory rate is modulated by the environment through a dopaminergic pathway and by experience through a serotonergic pathway}, Neuron \textbf{26}, 619 (2000).
\bibitem{Chase2004} D. L. Chase, J. S. Pepper, and M. R. Koelle, \emph{Mechanism of extrasynaptic dopamine signaling in Caenorhabditis elegans}, Nat. Neurosci. \textbf{7}, 1096 (2004).
\bibitem{Hills2004} T. Hills, P. J. Brockie, and A. V. Maricq, \emph{Dopamine and glutamate control area-restricted search behavior in Caenorhabditis elegans}, J. Neurosci. \textbf{24}, 1217 (2004).
\bibitem{Sanyal2004} S. Sanyal et al., \emph{Dopamine modulates the plasticity of mechanosensory responses in Caenorhabditis elegans}, EMBO J. \textbf{23}, 473 (2004).
\bibitem{Kindt2007} K. S. Kindt et al., \emph{Dopamine mediates context-dependent modulation of sensory plasticity in C.~elegans}, Neuron \textbf{55}, 662 (2007).
\bibitem{Hukema2008} R. K. Hukema, S. Rademakers, and G. Jansen, \emph{Gustatory plasticity in C.~elegans involves integration of negative cues and NaCl taste mediated by serotonin, dopamine, and glutamate}, Learn. Mem. \textbf{15}, 829 (2008).
\bibitem{Voglis2008} G. Voglis and N. Tavernarakis, \emph{A synaptic DEG/ENaC ion channel mediates learning in C.~elegans by facilitating dopamine signalling}, EMBO J. \textbf{27}, 3288 (2008).
\bibitem{Caron2013} S. J. Caron, V. Ruta, L. F. Abbott, and R. Axel, \emph{Random convergence of olfactory inputs in the Drosophila mushroom body}, Nature \textbf{497}, 113 (2013).
\bibitem{Honegger2011} K. S. Honegger, R. A. Campbell, and G. C. Turner, \emph{Cellular-resolution population imaging reveals robust sparse coding in the Drosophila mushroom body}, J. Neurosci. \textbf{31}, 11772 (2011).
\bibitem{Lin2014} A. C. Lin, A. M. Bygrave, A. de Calignon, T. Lee, and G. Miesenb\"ock, \emph{Sparse, decorrelated odor coding in the mushroom body enhances learned odor discrimination}, Nat. Neurosci. \textbf{17}, 559 (2014).
\bibitem{Dasgupta2017} S. Dasgupta, C. F. Stevens, and S. Navlakha, \emph{A neural algorithm for a fundamental computing problem}, Science \textbf{358}, 793 (2017).
\bibitem{Aso2014} Y. Aso et al., \emph{The neuronal architecture of the mushroom body provides a logic for associative learning}, eLife \textbf{3}, e04577 (2014).
\bibitem{AsoMBON2014} Y. Aso et al., \emph{Mushroom body output neurons encode valence and guide memory-based action selection in Drosophila}, eLife \textbf{3}, e04580 (2014).
\bibitem{ClaridgeChang2009} A. Claridge-Chang et al., \emph{Writing memories with light-addressable reinforcement circuitry}, Cell \textbf{139}, 405 (2009).
\bibitem{Liu2012} C. Liu et al., \emph{A subset of dopamine neurons signals reward for odour memory in Drosophila}, Nature \textbf{488}, 512 (2012).
\bibitem{Hige2015} T. Hige, Y. Aso, M. N. Modi, G. M. Rubin, and G. C. Turner, \emph{Heterosynaptic plasticity underlies aversive olfactory learning in Drosophila}, Neuron \textbf{88}, 985 (2015).
\bibitem{Levin1992} L. R. Levin et al., \emph{The Drosophila learning and memory gene rutabaga encodes a Ca$^{2+}$/calmodulin-responsive adenylyl cyclase}, Cell \textbf{68}, 479 (1992).
\bibitem{Tomchik2009} S. M. Tomchik and R. L. Davis, \emph{Dynamics of learning-related cAMP signaling and stimulus integration in the Drosophila olfactory pathway}, Neuron \textbf{64}, 510 (2009).
\bibitem{Handler2019} A. Handler et al., \emph{Distinct dopamine receptor pathways underlie the temporal sensitivity of associative learning}, Cell \textbf{178}, 60 (2019).
\bibitem{Tanimoto2004} H. Tanimoto, M. Heisenberg, and B. Gerber, \emph{Event timing turns punishment to reward}, Nature \textbf{430}, 983 (2004).
\bibitem{AsoRubin2016} Y. Aso and G. M. Rubin, \emph{Dopaminergic neurons write and update memories with cell-type-specific rules}, eLife \textbf{5}, e16135 (2016).
\bibitem{Li2020} F. Li et al., \emph{The connectome of the adult Drosophila mushroom body provides insights into function}, eLife \textbf{9}, e62576 (2020).
\bibitem{Felsenberg2017} J. Felsenberg, O. Barnstedt, P. Cognigni, S. Lin, and S. Waddell, \emph{Re-evaluation of learned information in Drosophila}, Nature \textbf{544}, 240 (2017).

\end{thebibliography}


\end{document}
